\PassOptionsToPackage{hyperfootnotes=false}{hyperref}
\documentclass[oneside]{amsart}

\usepackage[all]{xy}
\usepackage{tikz-cd}
\usepackage[T1]{fontenc}
\usepackage{xstring}
\usepackage{xparse}
\usepackage{xcolor}
\definecolor{brightmaroon}{rgb}{0.76, 0.13, 0.28}
\usepackage[linktocpage=true,colorlinks=true,hyperindex,hypertexnames=false,citecolor=blue,linkcolor=brightmaroon]{hyperref}
\usepackage[a4paper,left=1.25in,right=1.25in,top=0.75in,bottom=0.75in]{geometry}
%\usepackage[charter,greekfamily=didot]{mathdesign}
%\usepackage{Baskervaldx}
\usepackage{amssymb}
\usepackage{mathrsfs}
\usepackage{mathpazo}
\linespread{1.05}

\usepackage[nobottomtitles]{titlesec}
\usepackage{marginnote}
\usepackage[shortlabels]{enumitem}
\usepackage{longtable}
\usepackage{multicol}
\usepackage{aurical}
\usepackage{microtype}

\newtheoremstyle{ega-env-style}%
  {}{}{\rmfamily}{}{\bfseries}{.}{ }{\thmnote{(#3)}}%

\newtheoremstyle{ega-thm-env-style}%
  {}{}{\itshape}{}{\bfseries}{. --- }{ }{\thmname{#1}\thmnote{ (#3)}}%

\newtheoremstyle{ega-defn-env-style}%
  {}{}{\rmfamily}{}{\bfseries}{. --- }{ }{\thmname{#1}\thmnote{ (#3)}}%

\theoremstyle{ega-env-style}
\newtheorem*{env}{---}

\theoremstyle{ega-thm-env-style}
\newtheorem*{theorem}{Theorem}
\newtheorem*{proposition}{Proposition}
\newtheorem*{lemma}{Lemma}
\newtheorem*{corollary}{Corollary}

\theoremstyle{ega-defn-env-style}
\newtheorem*{definition}{Definition}
\newtheorem*{example}{Example}
\newtheorem*{examples}{Examples}
\newtheorem*{remark}{Remark}
\newtheorem*{remarks}{Remarks}
\newtheorem*{scholium}{Scholium}
\newtheorem*{notation}{Notation}

% indent subsections, see https://tex.stackexchange.com/questions/177290/.
% also make section titles bigger.
% also add § to \thesection, https://tex.stackexchange.com/questions/119667/ and https://tex.stackexchange.com/questions/308737/.
\makeatletter
\def\l@subsection{\@tocline{2}{0pt}{2.5pc}{2.2pc}{}}
\def\section{\@startsection{section}{1}%
  \z@{.7\linespacing\@plus\linespacing}{.5\linespacing}%
  {\normalfont\bfseries\Large\scshape\centering}}
\renewcommand{\@seccntformat}[1]{%
  \ifnum\pdfstrcmp{#1}{section}=0\textsection\fi%
  \csname the#1\endcsname.~}
\makeatother

%\allowdisplaybreaks[1]
%\binoppenalty=9999
%\relpenalty=9999

% for Chapter 0, Chapter I, etc.
% credit for ZeroRoman https://tex.stackexchange.com/questions/211414/
% added into scripts/make_book.py
%\newcommand{\ZeroRoman}[1]{\ifcase\value{#1}\relax 0\else\Roman{#1}\fi}
%\renewcommand{\thechapter}{\ZeroRoman{chapter}}

\def\mathcal{\mathscr}
\def\sh{\mathcal}                   % sheaf font
\def\bb{\mathbf}                    % bold font
\def\cat{\mathtt}                   % category font
\def\leq{\leqslant}                 % <=
\def\geq{\geqslant}                 % >=
\def\setmin{-}                      % set minus
\def\rad{\mathfrak{r}}              % radical
\def\nilrad{\mathfrak{N}}           % nilradical
\def\emp{\varnothing}               % empty set
\def\vphi{\phi}                     % for switching \phi and \varphi, change if needed
\def\HH{\mathrm{H}}                 % cohomology H
\def\CHH{\check{\HH}}               % Čech cohomology H
\def\RR{\mathrm{R}}                 % right derived R
\def\LL{\mathrm{L}}                 % left derived L
\def\dual#1{{#1}^\vee}              % dual
\def\kres{k}                        % residue field k
\def\C{\cat{C}}                     % category C
\def\op{^\cat{op}}                  % opposite category
\def\Set{\cat{Set}}                 % category of sets
\def\CHom{\cat{Hom}}                % functor category
\def\supertilde{{\,\widetilde{\,}\,}}   % use \supertilde instead of ^\sim
\def\GL{\bb{GL}}
\def\bbGamma{\boldsymbol\Gamma}
\def\red{\mathrm{red}}
\def\rg{\operatorname{rg}}
\def\gr{\operatorname{gr}}
\def\Gr{\operatorname{Gr}}
\def\Hom{\operatorname{Hom}}
\def\Proj{\operatorname{Proj}}
\def\Tor{\operatorname{Tor}}
\def\Ext{\operatorname{Ext}}
\def\Supp{\operatorname{Supp}}
\def\Ker{\operatorname{Ker}}
\def\Im{\operatorname{Im}}
\def\Coker{\operatorname{Coker}}
\def\Spec{\operatorname{Spec}}
\def\Spm{\operatorname{Spm}}
\def\Spf{\operatorname{Spf}}
\def\grad{\operatorname{grad}}
\def\dimc{\operatorname{dimc}}
\def\codim{\operatorname{codim}}
\def\id{\operatorname{id}}
\def\Der{\operatorname{Der}}
\def\Diff{\operatorname{Diff}}
\def\Hyp{\operatorname{Hyp}}
\def\Tr{\operatorname{Tr}}

\renewcommand{\to}{\mathchoice{\longrightarrow}{\rightarrow}{\rightarrow}{\rightarrow}}
\newcommand{\from}{\mathchoice{\longleftarrow}{\leftarrow}{\leftarrow}{\leftarrow}}
\let\mapstoo\mapsto
\renewcommand{\mapsto}{\mathchoice{\longmapsto}{\mapstoo}{\mapstoo}{\mapstoo}}
\def\isoto{\simeq} % isomorphism
\def\simto{\xrightarrow{\sim}}

\renewcommand{\hat}[1]{\widehat{#1}}
\renewcommand{\tilde}[1]{\widetilde{#1}}


\def\shHom{\sh{H}\textup{\kern-2.2pt{\Fontauri\slshape om}}}   % sheaf Hom
\def\shProj{\sh{P}\textup{\kern-2.2pt{\Fontauri\slshape roj}}} % sheaf Proj
\def\shExt{\sh{E}\textup{\kern-2.2pt{\Fontauri\slshape xt}}}   % sheaf Ext
\def\shTor{\sh{T}\textup{\kern-2.2pt{\Fontauri\slshape or}}}   % sheaf Tor
\def\shGr{\sh{G}\textup{\kern-2.2pt{\Fontauri\slshape r}}}     % sheaf Gr
\def\shDer{\sh{D}\,\textup{\kern-2.2pt{\Fontauri\slshape er}}} % sheaf Der
\def\shDiff{\sh{D}\,\textup{\kern-2.2pt{\Fontauri\slshape if{}f}}\,} % sheaf Diff
\def\shHomcont{\sh{H}\textup{\kern-2.2pt{\Fontauri\slshape om.\,cont}}}   % sheaf Hom.cont
\def\shAut{\sh{A}\textup{\kern-2.2pt{\Fontauri\slshape ut}}}   % sheaf Aut

% if unsure of a translation
%\newcommand{\unsure}[2][]{\hl{#2}\marginpar{#1}}
%\newcommand{\completelyunsure}{\unsure{[\ldots]}}
\def\unsure#1{#1 {\color{red}(?)}}
\def\completelyunsure{{\color{red}(???)}}

% use to mark where original page starts
\newcommand{\oldpage}[2][]{{\marginnote{\normalfont{\textbf{#1}~|~#2}}}\ignorespaces}
\def\sectionbreak{\begin{center}***\end{center}}

% Give every source label its own zero-size PDF destination.  The inherited
% unnumbered theorem environments otherwise leave many labels pointing to the
% preceding subsection or display anchor rather than to their exact statement.
\let\EGAOriginalLabel\label
\renewcommand{\label}[1]{\phantomsection\EGAOriginalLabel{#1}}

% for referencing environments.
% use as \sref{chapter-number.x.y.z}, with optional args
% for volume and indices, e.g. \sref[volume]{chapter-number.x.y.z}[i].
\makeatletter
\newcommand{\EGAUnavailableHyperref}[2]{%
  \begingroup
  \let\@latex@warning\@gobble
  \let\@latex@warning@no@line\@gobble
  \let\G@refundefinedtrue\relax
  \hyperref[#1]{#2}%
  \endgroup
}
\NewDocumentCommand{\sref}{o m o}{%
  \@ifundefined{r@#2}%
    {\IfNoValueTF{#1}%
      {\IfNoValueTF{#3}%
        {\EGAUnavailableHyperref{#2}{\normalfont{(\StrBehind{#2}{.})}}}%
        {\EGAUnavailableHyperref{#2}{\normalfont{(\StrBehind{#2}{.},~{#3})}}}}%
      {\IfNoValueTF{#3}%
        {\EGAUnavailableHyperref{#2}{\normalfont{(\textbf{#1},~\StrBehind{#2}{.})}}}%
        {\EGAUnavailableHyperref{#2}{\normalfont{(\textbf{#1},~\StrBehind{#2}{.},~{#3})}}}}}%
    {\IfNoValueTF{#1}%
      {\IfNoValueTF{#3}%
        {\hyperref[#2]{\normalfont{(\StrBehind{#2}{.})}}}%
        {\hyperref[#2]{\normalfont{(\StrBehind{#2}{.},~{#3})}}}}%
      {\IfNoValueTF{#3}%
        {\hyperref[#2]{\normalfont{(\textbf{#1},~\StrBehind{#2}{.})}}}%
        {\hyperref[#2]{\normalfont{(\textbf{#1},~\StrBehind{#2}{.},~{#3})}}}}}%
}
\makeatother

% for marking changes following errata
% use as \erratum[volume]{correction} to say where the erratum is given
\newcommand{\erratum}[2][]{{#2}\marginpar{\footnotesize\textbf{Err}\textsubscript{#1}}}

% for referencing equations, use as \eref{eq:chapter-number.x.y.z}.
\newcommand{\eref}[1]{\hyperref[#1]{\normalfont{(\StrBehind{#1}{.})}}}

\ifdefined\pdfinfoomitdate\pdfinfoomitdate=1\fi
\ifdefined\pdftrailerid\pdftrailerid{}\fi
\ifdefined\pdfsuppressptexinfo\pdfsuppressptexinfo=15\fi
\setlength{\emergencystretch}{3em}

\makeatletter
\let\EGAReaderOriginalHyperref\hyperref
\renewcommand{\hyperref}[2][]{%
  \@ifundefined{r@#1}{#2}{\EGAReaderOriginalHyperref[#1]{#2}}}
\makeatother
\hypersetup{
  pdftitle={Elements of Algebraic Geometry III-1 — corrected English edition},
  pdfauthor={Alexandre Grothendieck and Jean Dieudonné; English translation},
  pdfsubject={Cohomological study of coherent sheaves; corrected English edition}
}

\begin{document}

\title{Elements of Algebraic Geometry\\
III. Cohomological Study of Coherent Sheaves\\
First Part}
\author{Alexandre Grothendieck and Jean Dieudonné}
\date{}
\maketitle

% BEGIN SCHOLARLY TRANSLATION ATTRIBUTION

\begingroup
\small
\begin{quote}
\textbf{Translation sources and contributors.}
The original work is by Alexandre Grothendieck and Jean Dieudonn\'e.
The EGA English edition series draws on the human community translation
\texttt{github.com/ryankeleti/ega}, source version \texttt{3d7afe28}.
Its named contributors are Mateo Carmona, Tim Holzschuh, Tim Hosgood,
Ryan Keleti, Thiago Solovera \& Nery, and Wayne Yang.
The inherited English text includes Chapter 0, sections 8--10 and portions of section 11, and Chapter III, sections 1 and 3.
The additional translation is AI-translated; subsequent corrections, editing,
typesetting and assembly are AI-produced. The publisher did not translate
the text. Recorded AI work includes OpenAI Codex,
GPT-6 Astra, Ultra; the III-1 sources also include GPT-5.6 Sol, Ultra.
The accompanying \texttt{ATTRIBUTION.md} gives source links, contribution
scope and rights information. Attribution does not assert endorsement
or a licence grant from the upstream contributors.

\medskip
\textbf{Rights and availability.}
Rights in the original works and in any protected contributions to these
derivative translations and transcriptions remain with their respective
rights holders. This project does not dedicate those works or contributions
to the public domain or grant a licence over them. Existing third-party
notices and any permissions already granted remain applicable.

These editions are available free of charge for non-commercial, purely
educational purposes, on the presumption that the respective rights holders
would not object to these translations and transcriptions being freely
available alongside the French originals on NUMDAM
(\texttt{www.numdam.org/journals/PMIHES/}). This presumption is not a statement
that permission has been obtained; free access to the originals does not
itself grant permission to translate, adapt or republish them.

Rights holders who object to the availability of their material may open
an issue at \texttt{github.com/KokunoYumeto/ega-en/issues}. The publisher
will take down the affected material at the rights holder's request.
\end{quote}
\endgroup
% END SCHOLARLY TRANSLATION ATTRIBUTION


\begin{center}
\large Corrected English edition
\end{center}

\begin{quote}\small
Source preparation: OpenAI Codex --- GPT-5.6 Sol and GPT-6 Astra,
Ultra effort.
\end{quote}

\phantomsection
\label{section:ega0}
\pdfbookmark[0]{Chapter 0 (continued)}{chapter.0.continued}
\addcontentsline{toc}{part}{Chapter 0 (continued)}
\section*{CHAPTER 0 (continued)}
\setcounter{section}{7}
\setcounter{subsection}{0}
\setcounter{equation}{0}

\section{Representable functors}
\label{section:0.8}

\subsection{Representable functors}
\label{subsection:0.8.1}

\oldpage[0\textsubscript{III}]{5}
\begin{env}[8.1.1]
\label{0.8.1.1}
We denote by $\Set$ the category of sets.
Let $\C$ be a category; for two objects $X$, $Y$ of $\C$, we set $h_X(Y)=\Hom(Y,X)$; for each morphism $u:Y\to Y'$ in $\C$, we denote by $h_X(u)$ the map $v\mapsto vu$ from $\Hom(Y',X)$ to $\Hom(Y,X)$.
It is immediate that with these definitions, $h_X:\C\to\Set$ is a \emph{contravariant functor}, i.e., an object of the category $\CHom(\C\op,\Set)$, of covariant functors from the category $\C\op$ (the dual of the category $\C$) to the category $\Set$ (T, 1.7, (d) and \cite{III-29}).
\end{env}

\begin{env}[8.1.2]
\label{0.8.1.2}
Now let $w:X\to X'$ be a morphism in $\C$; for each $Y\in\C$ and each $v\in\Hom(Y,X)=h_X(Y)$,
we have $wv\in\Hom(Y,X')=h_{X'}(Y)$; we denote by $h_w(Y)$ the map $v\mapsto wv$ from
$h_X(Y)$ to $h_{X'}(Y)$. It is immediate that for each morphism $u:Y\to Y'$ in $\C$, the
diagram
\[
  \xymatrix{
    h_X(Y')\ar[r]^{h_X(u)}\ar[d]_{h_w(Y')} &
    h_X(Y)\ar[d]^{h_w(Y)}\\
    h_{X'}(Y')\ar[r]_{h_{X'}(u)} &
    h_{X'}(Y)
  }
\]
is commutative; in other words, $h_w$ is a \emph{natural transformation (or functorial morphism)
$h_X\to h_{X'}$} (T, 1.2), also a morphism in the category $\CHom(\C\op,\Set)$ (T, 1.7, (d)). The
definitions of $h_X$ and of $h_w$ therefore constitute the definition of a \emph{canonical
covariant functor}
\[
  h:\C\to\CHom(\C\op,\Set),\quad X\mapsto h_X.
  \tag{8.1.2.1}
\]
\end{env}

\begin{env}[8.1.3]
\label{0.8.1.3}
Let $X$ be an object in $\C$, $F$ a contravariant functor from $\C$ to $\Set$
(an object of $\CHom(\C\op,\Set)$). Let $g:h_X\to F$ be a \emph{natural transformation}: for all $Y\in\C$,
\oldpage[0\textsubscript{III}]{6}
$g(Y)$ is thus a map $h_X(Y)\to F(Y)$ such that for each morphism $u:Y\to Y'$ in $\C$,
the diagram
\[
  \xymatrix{
    h_X(Y')\ar[r]^{h_X(u)}\ar[d]_{g(Y')} &
    h_X(Y)\ar[d]^{g(Y)}\\
    F(Y')\ar[r]_{F(u)} &
    F(Y)
  }
  \tag{8.1.3.1}\label{0.8.1.3.1}
\]
is commutative. In particular, we have a map $g(X):h_X(X)=\Hom(X,X)\to F(X)$, hence an element
\[
  \alpha(g)=(g(X))(1_X)\in F(X)
  \tag{8.1.3.2}
\]
and as a result a canonical map
\[
  \alpha:\Hom(h_X,F)\to F(X).
  \tag{8.1.3.3}
\]

Conversely, consider an element $\xi\in F(X)$; for each morphism $v:Y\to X$ in $\C$, $F(v)$ is a
map $F(X)\to F(Y)$; consider the map
\[
  v\mapsto(F(v))(\xi)
  \tag{8.1.3.4}
\]
from $h_X(Y)$ to $F(Y)$; if we denote by $(\beta(\xi))(Y)$ this map,
\[
  \beta(\xi):h_X\to F
  \tag{8.1.3.5}
\]
is a \emph{natural transformation}, since for each morphism $u:Y\to Y'$ in $\C$ we have
$(F(vu))(\xi)=(F(u)\circ F(v))(\xi)$, which makes (\hyperref[0.8.1.3.1]{8.1.3.1}) commutative for $g=\beta(\xi)$.
We have thus defined a canonical map
\[
  \beta:F(X)\to\Hom(h_X,F).
  \tag{8.1.3.6}
\]
\end{env}

\begin{proposition}[8.1.4]
\label{0.8.1.4}
The maps $\alpha$ and $\beta$ are the inverse bijections of each other.
\end{proposition}

\begin{proof}
We calculate $\alpha(\beta(\xi))$ for $\xi\in F(X)$; for each $Y\in\C$, $(\beta(\xi))(Y)$ is a map
$g_1(Y):v\mapsto(F(v))(\xi)$ from $h_X(Y)$ to $F(Y)$. We thus have
\[
  \alpha(\beta(\xi))=(g_1(X))(1_X)=(F(1_X))(\xi)=1_{F(X)}(\xi)=\xi.
\]
We now calculate $\beta(\alpha(g))$ for $g\in\Hom(h_X,F)$; for each $Y\in\C$, $(\beta(\alpha(g)))(Y)$
is the map $v\mapsto(F(v))((g(X))(1_X))$; according to the commutativity of (\hyperref[0.8.1.3.1]{8.1.3.1}), this map
is none other than $v\mapsto(g(Y))((h_X(v))(1_X))=(g(Y))(v)$ by definition of $h_X(v)$, in other words,
it is equal to $g(Y)$, which finishes the proof.
\end{proof}

\begin{env}[8.1.5]
\label{0.8.1.5}
Recall that a \emph{subcategory} $\C'$ of a category $\C$ is defined by the condition that its objects
are objects of $\C$, and that if $X'$, $Y'$ are two objects of $\C'$, then the set
$\Hom_{\C'}(X',Y')$ of morphisms $X'\to Y'$ \emph{in $\C'$} is a subset of the set $\Hom_\C(X',Y')$ of
morphisms $X'\to Y'$ \emph{in $\C$}, the canonical map of ``composition of morphisms''
\[
  \Hom_{\C'}(X',Y')\times\Hom_{\C'}(Y',Z')\to\Hom_{\C'}(X',Z')
\]
\oldpage[0\textsubscript{III}]{7}
being the restriction of the canonical map
\[
  \Hom_\C(X',Y')\times\Hom_\C(Y',Z')\to\Hom_\C(X',Z').
\]

We say that $\C'$ is a \emph{full} subcategory of $\C$ if $\Hom_{\C'}(X',Y')=\Hom_\C(X',Y')$ for every
pair of objects in $\C'$. The subcategory $\C''$ of $\C$ consisting of the objects of $\C$ isomorphic to
objects of $\C'$ is then again a full subcategory of $\C$, \emph{equivalent} (T, 1.2) to $\C'$ as we
verify easily.

A covariant functor $F:\C_1\to\C_2$ is called \emph{fully faithful} if for every pair of objects
$X_1$, $Y_1$ of $\C_1$, the map $u\mapsto F(u)$ from $\Hom(X_1,Y_1)$ to $\Hom(F(X_1),F(Y_1))$ is
\emph{bijective}; this implies that the subcategory $F(\C_1)$ of $\C_2$ is \emph{full}. In addition,
if two objects $X_1$, $X_1'$ have the same image $X_2$, then there exists a unique isomorphism
$u:X_1\to X_1'$ such that $F(u)=1_{X_2}$. For each object $X_2$ of $F(\C_1)$, let $G(X_2)$ be one of
the objects $X_1$ of $\C_1$ such that $F(X_1)=X_2$ ($G$ is defined by means of the axiom of choice); for
each morphism $v:X_2\to Y_2$ in $F(\C_1)$, $G(v)$ will be the unique morphism $u:G(X_2)\to G(Y_2)$ such
that $F(u)=v$; $G$ is then a \emph{functor} from $F(\C_1)$ to $\C_1$; $FG$ is the identity functor on
$F(\C_1)$, and the above shows that there exists an isomorphism of functors $\vphi:1_{\C_1}\to GF$
such that $F$, $G$, $\vphi$, and the identity $1_{F(\C_1)}\to FG$ define an \emph{equivalence} between
the category $\C_1$ and the full subcategory $F(\C_1)$ of $\C_2$ (T, 1.2).
\end{env}

\begin{env}[8.1.6]
\label{0.8.1.6}
We apply Proposition \sref{0.8.1.4} to the case where $F$ is $h_{X'}$, $X'$ being any
object of $\C$; the map $\beta:\Hom(X,X')\to\Hom(h_X,h_{X'})$ is none other than the map
$w\mapsto h_w$ defined in \sref{0.8.1.2}; this map being \emph{bijective}, we see
with the terminology of \sref{0.8.1.5} that:
\end{env}

\begin{proposition}[8.1.7]
\label{0.8.1.7}
The canonical functor $h:\C\to\CHom(\C\op,\Set)$ is fully faithful.
\end{proposition}

\begin{env}[8.1.8]
\label{0.8.1.8}
Let $F$ be a contravariant functor from $\C$ to $\Set$; we say that $F$ is \emph{representable} if there
exists an object $X\in\C$ such that $F$ is \emph{isomorphic} to $h_X$; it follows from
Proposition \sref{0.8.1.7} that the data of an $X\in\C$ and an isomorphism of functors
$g:h_X\to F$ determines $X$ up to unique isomorphism. Proposition \sref{0.8.1.7} then
implies that $h$ defines an \emph{equivalence} between $\C$ and the full subcategory of
$\CHom(\C\op,\Set)$ consisting of the \emph{contravariant representable functors}. It follows from
Proposition \sref{0.8.1.4} that the data of a natural transformation $g:h_X\to F$ is
equivalent to that of an element $\xi\in F(X)$; to say that $g$ is an \emph{isomorphism} is
equivalent to the following condition on $\xi$: \emph{for every object $Y$ of $\C$ the map
$v\mapsto(F(v))(\xi)$ from $\Hom(Y,X)$ to $F(Y)$ is bijective}. When $\xi$ satisfies this condition,
we say that the pair $(X,\xi)$ \emph{represents} the representable functor $F$. By abuse of language,
we also say that the object $X\in\C$ represents $F$ if there exists a $\xi\in F(X)$ such that
$(X,\xi)$ represents $F$, in other words if $h_X$ is isomorphic to $F$.

Let $F$, $F'$ be two contravariant representable functors from $\C$ to $\Set$, $h_X\to F$ and
$h_{X'}\to F'$ two isomorphisms of functors. Then it follows from \sref{0.8.1.6} that
there is a canonical bijective correspondence between $\Hom(X,X')$ and the set $\Hom(F,F')$ of
natural transformations $F\to F'$.
\end{env}

\begin{env}[8.1.9]
\label{0.8.1.9}
\textit{Example I. Projective limits}.
The notion of a contravariant representable functor covers in particular the ``dual'' notion of the usual notion of a ``solution to a universal problem''.
\oldpage[0\textsubscript{III}]{8}
More generally, we will see that the notion of the \emph{projective limit} is a special case of the notion of a representable functor.
Recall that in a category $\C$, we define a \emph{projective system} by the data of a preordered set $I$, a family $(A_\alpha)_{\alpha\in I}$ of objects of $\C$, and for every pair of indices $(\alpha,\beta)$ such that $\alpha\leq\beta$, a morphism $u_{\alpha\beta}:A_\beta\to A_\alpha$.
A \emph{projective limit} of this system in $\C$ consists of an object $B$ of $\C$ (denoted $\varprojlim A_\alpha$), and for each $\alpha\in I$, a morphism $u_\alpha:B\to A_\alpha$ such that: 1\textsuperscript{st}. $u_\alpha=u_{\alpha\beta}u_\beta$ for $\alpha\leq\beta$; 2\textsuperscript{nd}. for every object $X$ of $\C$ and every family $(v_\alpha)_{\alpha\in I}$ of morphisms $v_\alpha:X\to A_\alpha$ such that $v_\alpha=u_{\alpha\beta}v_\beta$ for $\alpha\leq\beta$, there exists a unique morphism $v:X\to B$ (denoted $\varprojlim v_\alpha$) such that $v_\alpha=u_\alpha v$ for all $\alpha\in I$ (T, 1.8).
This can be interpreted in the following way: the $u_{\alpha\beta}$ canonically define maps
\[
  \overline{u}_{\alpha\beta}:\Hom(X,A_\beta)\to\Hom(X,A_\alpha)
\]
which define a \emph{projective system} of sets $(\Hom(X,A_\alpha),\overline{u}_{\alpha\beta})$, and $(v_\alpha)$ is by definition an element of the set $\varprojlim\Hom(X,A_\alpha)$; it is clear that $X\mapsto\varprojlim\Hom(X,A_\alpha)$ is a \emph{contravariant functor} from $\C$ to $\Set$, and the existence of the projective limit $B$ is equivalent to saying that $(v_\alpha)\mapsto\varprojlim v_\alpha$ is an \emph{isomorphism} of functors in $X$
\[
  \varprojlim\Hom(X,A_\alpha)\isoto\Hom(X,B),
  \tag{8.1.9.1}
\]
in other words, that the functor $X\mapsto\varprojlim\Hom(X,A_\alpha)$ is \emph{representable}.
\end{env}

\begin{env}[8.1.10]
\label{0.8.1.10}
\textit{Example II. Final objects}.
Let $\C$ be a category, $\{a\}$ a singleton set.
Consider the contravariant functor $F:\C\to\Set$ which sends every object $X$ of $\C$ to the set $\{a\}$, and every morphism $X\to X'$ in $\C$ to the unique map $\{a\}\to\{a\}$.
To say that this functor is \emph{representable} means that there exists an object $e\in\C$ such that for every $Y\in\C$, $\Hom(Y,e)=h_e(Y)$ is a \emph{singleton set}; we say that $e$ is an \emph{final object} of $\C$, and it is clear that two final objects of $\C$ are isomorphic (which allows us to define, in general with the axiom of choice, \emph{one} final object of $\C$ which we then denote $e_\C$).
For example, in the category $\Set$, the final objects are the singleton sets; in the category of \emph{augmented algebras} over a field $K$ (where the morphisms are the algebra homomorphisms compatible with the augmentation), $K$ is a final object; in the category of \emph{$S$-preschemes} \sref[I]{I.2.5.1}, $S$ is a final object.
\end{env}

\begin{env}[8.1.11]
\label{0.8.1.11}
For two objects $X$ and $Y$ of a category $\C$, set $h_X'(Y)=\Hom(X,Y)$, and for every morphism $u:Y\to Y'$, let $h_X'(u)$ be the map $v\mapsto uv$ from $\Hom(X,Y)$ to $\Hom(X,Y')$; $h_X'$ is then a \emph{covariant functor $\C\to\Set$}, so we deduce as in \sref{0.8.1.2} the definition of a canonical covariant functor $h':\C\op\to\CHom(\C,\Set)$; a \emph{covariant} functor $F$ from $\C$ to $\Set$, in other words an object of $\CHom(\C,\Set)$, is then \emph{representable} if there exists an object $X\in\C$ (necessarily unique up to unique isomorphism) such that $F$ is \emph{isomorphic to $h_X'$}; we leave it to the reader to develop the ``dual'' notions of the above, which this time cover the notion of an \emph{inductive limit}, and in particular the usual notion of a ``solution to a universal problem''.
\end{env}

\subsection{Algebraic structures in categories}
\label{subsection:0.8.2}

\begin{env}[8.2.1]
\label{0.8.2.1}
\oldpage[0\textsubscript{III}]{9}
Given two contravariant functors $F$ and $F'$ from $\C$ to $\Set$, recall that for every object $Y\in\C$, we set $(F\times F')(Y)=F(Y)\times F'(Y)$, and for every morphism $u:Y\to Y'$ in $\C$, we set $(F\times F')(u)=F(u)\times F'(u)$, which is the map $(t,t')\mapsto(F(u)(t),F'(u)(t'))$ from $F(Y')\times F'(Y')$ to $F(Y)\times F'(Y)$; $F\times F':\C\to\Set$ is thus a \emph{contravariant functor} (which is none other than the \emph{product} of the objects $F$ and $F'$ in the category $\CHom(\C\op,\Set)$).
Given an object $X\in\C$, we call an \emph{internal composition law} on $X$ a \emph{natural transformation}
\[
\label{0.8.2.1.1}
  \gamma_X:h_X\times h_X\to h_X.
  \tag{8.2.1.1}
\]

In other words (T, 1.2), for every object $Y\in\C$, $\gamma_X(Y)$ is a map $h_X(Y)\times h_X(Y)\to h_X(Y)$ (thus by definition an \emph{internal composition law} on the set $h_X(Y)$) with the condition that for every morphism $u:Y\to Y'$ in $\C$, the diagram
\[
  \xymatrix{
    h_X(Y')\times h_X(Y')\ar[rr]^{h_X(u)\times h_X(u)}\ar[dd]_{\gamma_X(Y')} & &
    h_X(Y)\times h_X(Y)\ar[dd]^{\gamma_X(Y)}\\\\
    h_X(Y')\ar[rr]_{h_X(u)} & &
    h_X(Y)
  }
\]
is commutative; this implies that for the composition laws $\gamma_X(Y)$ and $\gamma_X(Y')$, $h_X(u)$ is a \emph{homomorphism} from $h_X(Y')$ to $h_X(Y)$.

In a similar way, given two objects $Z$ and $X$ of $\C$, we call an \emph{external composition law on $X$, with $Z$ as its domain of operators} a natural transformation
\[
\label{0.8.2.1.2}
  \omega_{X,Z}:h_Z\times h_X\to h_X.
  \tag{8.2.1.2}
\]

We see as above that for every $Y\in\C$, $\omega_{X,Z}(Y)$ is an external composition law on $h_X(Y)$, with $h_Z(Y)$ as its domain of operators and such that for every morphism $u:Y\to Y'$, $h_X(u)$ and $h_Z(u)$ form a \emph{di-homomorphism} from $(h_Z(Y'),h_X(Y'))$ to $(h_Z(Y),h_X(Y))$.
\end{env}

\begin{env}[8.2.2]
\label{0.8.2.2}
Let $X'$ be a second object of $\C$, and suppose we are given an internal composition law $\gamma_{X'}$ on $X'$; we say that a morphism $w:X\to X'$ in $\C$ is a \emph{homomorphism} for the composition laws if for every $Y\in\C$, $h_w(Y):h_X(Y)\to h_{X'}(Y)$ is a \emph{homomorphism} for the composition laws $\gamma_X(Y)$ and $\gamma_{X'}(Y)$.
If $X''$ is a third
\oldpage[0\textsubscript{III}]{10}
object of $\C$ equipped with an internal composition law $\gamma_{X''}$ and $w':X'\to X''$ is a morphism in $\C$ which is a homomorphism for $\gamma_{X'}$ and $\gamma_{X''}$, then it is clear that the morphism $w'w:X\to X''$ is a homomorphism for the composition laws $\gamma_X$ and $\gamma_{X''}$.
An isomorphism $w:X\isoto X'$ in $\C$ is called an \emph{isomorphism for the composition laws $\gamma_X$ and $\gamma_{X'}$} if $w$ is a homomorphism for these composition laws, and if its inverse morphism $w^{-1}$ is a homomorphism for the composition laws $\gamma_{X'}$ and $\gamma_X$.

We define in a similar way the \emph{di-homomorphisms} for pairs of objects of $\C$ equipped with external composition laws.
\end{env}

\begin{env}[8.2.3]
\label{0.8.2.3}
When an internal composition law $\gamma_X$ on an object $X\in\C$ is such that $\gamma_X(Y)$ is a \emph{group law} on $h_X(Y)$ for \emph{every} $Y\in\C$, we say that $X$, equipped with this law, is a \emph{$\C$-group} or a \emph{group object in $\C$}.
We similarly define \emph{$\C$-rings}, \emph{$\C$-modules}, etc.
\end{env}

\begin{env}[8.2.4]
\label{0.8.2.4}
Suppose that the \emph{product} $X\times X$ of an object $X\in\C$ by itself exists in $\C$; by definition, we then have $h_{X\times X}=h_X\times h_X$ up to canonical isomorphism, since it is a particular case of the projective limit \sref{0.8.1.9}; an internal composition law on $X$ can thus be considered as a functorial morphism $\gamma_X:h_{X\times X}\to h_X$, and thus canonically determines \sref{0.8.1.6} an element $c_X\in\Hom(X\times X,X)$ such that $h_{c_X}=\gamma_X$; in this case, the data of an internal composition law on $X$ is equivalent to the data of a morphism $X\times X\to X$; when $\C$ is the category $\Set$, we recover the classical notion of an internal composition law on a set.
We have an analogous result for an external composition law when the product $Z\times X$ exists in $\C$.
\end{env}

\begin{env}[8.2.5]
\label{0.8.2.5}
With the above notation, suppose that in addition $X\times X\times X$ exists in $\C$; the characterization of the product as an object representing a functor \sref{0.8.1.9} implies the existence of canonical isomorphisms
\[
  (X\times X)\times X\isoto X\times X\times X\isoto X\times(X\times X);
\]
if we canonically identify $X\times X\times X$ with $(X\times X)\times X$, then the map $\gamma_X(Y)\times 1_{h_X(Y)}$ identifies with $h_{c_X\times 1_X}(Y)$ for all $Y\in\C$.
As a result, it is equivalent to say that for every $Y\in\C$, the internal law $\gamma_X(Y)$ is associative, or that the diagram of maps
\[
  \xymatrix{
    h_X(Y)\times h_X(Y)\times h_X(Y)\ar[rr]^{\gamma_X(Y)\times 1}\ar[dd]_{1\times\gamma_X(Y)} & &
    h_X(Y)\times h_X(Y)\ar[dd]^{\gamma_X(Y)}\\\\
    h_X(Y)\times h_X(Y)\ar[rr]_{\gamma_X(Y)} & &
    h_X(Y)
  }
\]
\oldpage[0\textsubscript{III}]{11}
is commutative, or that the diagram of morphisms
\[
  \xymatrix{
    X\times X\times X\ar[rr]^{c_X\times 1_X}\ar[dd]_{1_X\times c_X} & &
    X\times X\ar[dd]^{c_X}\\\\
    X\times X\ar[rr]_{c_X} & &
    X
  }
\]
is commutative.
\end{env}

\begin{env}[8.2.6]
\label{0.8.2.6}
Under the hypotheses of \sref{0.8.2.5}, if we want to express, for every $Y\in\C$, the internal law $\gamma_X(Y)$ as a \emph{group law}, then it is first necessary that it is associative, and second that there exists a map $\alpha_X(Y):h_X(Y)\to h_X(Y)$ having the properties of the \emph{inverse} operation of a group; as for every morphism $u:Y\to Y'$ in $\C$, we have seen that $h_X(u)$ must be a group homomorphism $h_X(Y')\to h_X(Y)$, we first see that $\alpha_X:h_X\to h_X$ must be a \emph{natural transformation}.
On the other hand, one can express the characteristic properties of the inverse $s\mapsto s^{-1}$ of a group $G$ without involving the identity element: it suffices to check that the two composite maps
\[
  (s,t)\mapsto(s,s^{-1},t)\mapsto(s,s^{-1}t)\mapsto s(s^{-1}t),
\]
\[
  (s,t)\mapsto(s,s^{-1},t)\mapsto(s,ts^{-1})\mapsto(ts^{-1})s
\]
are equal to the second projection $(s,t)\mapsto t$ from $G\times G$ to $G$.
By \sref{0.8.1.3}, we have $\alpha_X=h_{a_X}$, where $a_X\in\Hom(X,X)$; the first condition above then expresses that the composite morphism
\[
  X\times X\xrightarrow{(1_X,a_X)\times 1_X}X\times X\times X\xrightarrow{1_X\times c_X}X\times X\xrightarrow{c_X}X
\]
is the second projection $X\times X\to X$ in $\C$, and the second condition is similar.
\end{env}

\begin{env}[8.2.7]
\label{0.8.2.7}
Now suppose that there exists a \emph{final object} $e$ \sref{0.8.1.10} in $\C$.
Let us always assume that $\gamma_X(Y)$ is a group law on $h_X(Y)$ for every $Y\in\C$, and denote by $\eta_X(Y)$ the identity element of $\gamma_X(Y)$.
As, for every morphism $u:Y\to Y'$ in $\C$, $h_X(u)$ is a group homomorphism, we have $\eta_X(Y)=(h_X(u))(\eta_X(Y'))$; taking in particular $Y'=e$, in which case $u$ is the unique element $\varepsilon$ of $\Hom(Y,e)$, we see that the element $\eta_X(e)$ completely determines $\eta_X(Y)$ for every $Y\in\C$.
Set $e_X=\eta_X(X)$, the identity element of the group $h_X(X)=\Hom(X,X)$; the commutativity of the diagram
\[
  \xymatrix{
    h_X(e)\ar[r]^{h_X(\varepsilon)}\ar[d]_{h_{e_X}(e)} &
    h_X(Y)\ar[d]^{h_{e_X}(Y)}\\
    h_X(e)\ar[r]_{h_X(\varepsilon)} &
    h_X(Y)
  }
\]
(cf. \sref{0.8.1.2}) shows that, on the set $h_X(Y)$, the map $h_{e_X}(Y)$ is none other than
\oldpage[0\textsubscript{III}]{12}
$s\mapsto\eta_X(Y)$ sending every element to the identity element.
We then verify that the fact that $\eta_X(Y)$ is the identity element of $\gamma_X(Y)$ for every $Y\in\C$ is equivalent to saying that the composite morphism
\[
  X\xrightarrow{(1_X,1_X)}X\times X\xrightarrow{1_X\times e_X}X\times X\xrightarrow{c_X}X,
\]
and the analog in which we swap $1_X$ and $e_X$, are both \emph{equal} to $1_X$.
\end{env}

\begin{env}[8.2.8]
\label{0.8.2.8}
One could of course easily extend the examples of algebraic structures in categories.
The example of groups was treated with enough detail, but later on we will usually leave it to the reader to develop analogous notions for the examples of algebraic structures we will encounter.
\end{env}

\section{Constructible sets}
\label{section:0.9}

\subsection{Constructible sets}
\label{subsection:0.9.1}

\begin{definition}[9.1.1]
\label{0.9.1.1}
We say that a continuous map $f:X\to Y$ is \emph{quasi-compact} if for every quasi-compact open subset $U$ of $Y$, $f^{-1}(U)$ is quasi-compact.
We say that a subset $Z$ of a topological space $X$ is \emph{retrocompact} in $X$ if the canonical injection $Z\to X$ is quasi-compact, in other words, if for every quasi-compact open subset $U$ of $X$, $U\cap Z$ is quasi-compact.
\end{definition}

A \emph{closed} subset of $X$ is retrocompact in $X$, but a quasi-compact subset of $X$ is not necessarily retrocompact in $X$.
If $X$ is quasi-compact, every retrocompact open subset of $X$ is quasi-compact.
It is clear that every \emph{finite} union of retrocompact sets in $X$ is retrocompact in $X$, as every finite union of quasi-compact sets is quasi-compact.
Every finite intersection of retrocompact open sets in $X$ is a retrocompact open set in $X$.
In a \emph{locally Noetherian} space $X$, every quasi-compact set is a Noetherian subspace, and as a result \emph{every subset} of $X$ is retrocompact in $X$.

\begin{definition}[9.1.2]
\label{0.9.1.2}
Given a topological space $X$, we say that a subset of $X$ is \emph{constructible} if it belongs to the smallest set of subsets $\mathfrak{F}$ of $X$ containing all the retrocompact open subsets of $X$ and stable under finite intersections and complements (which implies that $\mathfrak{F}$ is also stable under finite unions).
\end{definition}

\begin{proposition}[9.1.3]
\label{0.9.1.3}
For a subset of $X$ to be constructible, it is necessary and sufficient for it to be a finite union of sets of the form $U\cap\complement{V}$, where $U$ and $V$ are retrocompact open sets in $X$.
\end{proposition}

\begin{proof}
It is clear that the condition is sufficient.
To see that it is necessary, consider the set $\mathfrak{G}$ of finite unions of sets of the form $U\cap\complement{V}$, where $U$ and $V$ are retrocompact open sets in $X$; it suffices to see that every complement of a set in $\mathfrak{G}$ is in $\mathfrak{G}$.
Let $Z=\bigcup_{i\in I}(U_i\cap\complement{V_i})$, where $I$ is finite, $U_i$ and $V_i$ retrocompact open sets in $X$; we have
$\complement{Z}=\bigcap_{i\in I}(V_i\cup\complement{U_i})$, so $\complement{Z}$ is a finite union of sets which are intersections of a certain number of the $V_i$ and of a certain number of the $\complement{U_i}$, thus of the form
\oldpage[0\textsubscript{III}]{13}
$V\cap\complement{U}$, where $U$ is the union of a certain number of the $U_i$ and $V$ is the intersection of a certain number of the $V_i$; but we have noted above that finite unions and intersections of retrocompact open sets in $X$ are retrocompact open sets in $X$, hence the conclusion.
\end{proof}

\begin{corollary}[9.1.4]
\label{0.9.1.4}
Every constructible subset of $X$ is retrocompact in $X$.
\end{corollary}

\begin{proof}
It suffices to show that if $U$ and $V$ are retrocompact open sets in $X$, then $U\cap\complement{V}$ is retrocompact in $X$; if $W$ is a quasi-compact open set in $X$, then $W\cap U\cap\complement{V}$ is closed in the quasi-compact space $W\cap U$, hence it is quasi-compact.
\end{proof}

In particular:
\begin{corollary}[9.1.5]
\label{0.9.1.5}
For an open subset $U$ of $X$ to be constructible, it is necessary and sufficient for it to be retrocompact in $X$.
For a closed subset $F$ of $X$ to be constructible, it is necessary and sufficient for the open set $\complement{F}$ to be retrocompact.
\end{corollary}

\begin{env}[9.1.6]
\label{0.9.1.6}
An important case is when every quasi-compact open subset of $X$ is retrocompact, in other words, when the intersection of two quasi-compact open subsets of $X$ is quasi-compact (cf.~\sref[I]{I.5.5.6}).
When $X$ is also quasi-compact, this implies that the retrocompact open subsets of $X$ are identical to the quasi-compact open subsets of $X$, and the constructible subsets of $X$ are finite unions of sets of the form $U\cap\complement{V}$, where $U$ and $V$ are quasi-compact open sets.
\end{env}

\begin{corollary}[9.1.7]
\label{0.9.1.7}
For a subset of a Noetherian space $X$ to be constructible, it is necessary and sufficient for it to be a finite union of locally closed subsets of $X$.
\end{corollary}

\begin{proposition}[9.1.8]
\label{0.9.1.8}
Let $X$ be a topological space, $U$ an open subset of $X$.
\begin{enumerate}
  \item[{\rm(i)}] If $T$ is a constructible subset of $X$, then $T\cap U$ is a constructible subset of $U$.
  \item[{\rm(ii)}] In addition, suppose that $U$ is retrocompact in $X$. For a subset $Z$ of $U$ to be constructible in $X$, it is necessary and sufficient for it to be constructible in $U$.
\end{enumerate}
\end{proposition}

\begin{proof}
\medskip\noindent
\begin{enumerate}
  \item[(i)] Using Proposition \sref{0.9.1.3}, we reduce to showing that if $T$ is a retrocompact open set in $X$, then $T\cap U$ is a retrocompact open set in $U$, in other words, for every quasi-compact open $W\subset U$, $T\cap U\cap W=T\cap W$ is quasi-compact, which immediately follows from the hypothesis.
  \item[(ii)] The condition is necessary by (i), so it remains to show that it is sufficient.
    By Proposition \sref{0.9.1.3}, it suffices to consider the case where $Z$ is a retrocompact open set \emph{in $U$}, because it will then follow that $U\setmin Z$ is constructible in $X$, and if $Z$ and $Z'$ are two retrocompact opens \emph{in $U$}, then $Z\cap(U\setmin Z')$ will be constructible in $X$.
    If $W$ is a quasi-compact open set in $X$ and $Z$ a retrocompact open set in $U$, then we have $Z\cap W=Z\cap(W\cap U)$, and by hypothesis $W\cap U$ is a quasi-compact open set in $U$; so $W\cap Z$ is quasi-compact, and as a result $Z$ is a retrocompact open set in $X$, and \emph{a fortiori} constructible in $X$.
\end{enumerate}
\end{proof}

\begin{corollary}[9.1.9]
\label{0.9.1.9}
Let $X$ be a topological space, $(U_i)_{i\in I}$ a finite cover of $X$ by retrocompact open sets in $X$.
For a subset $Z$ of $X$ to be constructible in $X$, it is necessary and sufficient for $Z\cap U_i$ to be constructible in $U_i$ for all $i\in I$.
\end{corollary}

\begin{env}[9.1.10]
\label{0.9.1.10}
In particular, suppose that $X$ is quasi-compact and every point
\oldpage[0\textsubscript{III}]{14}
of $X$ admits a fundamental system of retrocompact open neighbourhoods in $X$ (and \emph{a fortiori} quasi-compact); then the condition for a subset $Z$ of $X$ to be constructible in $X$ is of a \emph{local} nature, in other words, it is necessary and sufficient that for every $x\in X$, there exists an open neighbourhood $V$ of $x$ such that $V\cap Z$ is constructible in $V$.
Indeed, if this condition is satisfied, then there exists for every $x\in X$ an open neighbourhood $V$ of $x$ which is \emph{retrocompact in $X$} and such that $V\cap Z$ is constructible in $V$, by the hypotheses on $X$ and by Proposition \sref{0.9.1.8}[i]; it then suffices to cover $X$ by a finite number of these neighbourhoods and to apply Corollary \sref{0.9.1.9}.
\end{env}

\begin{definition}[9.1.11]
\label{0.9.1.11}
Let $X$ be a topological space.
We say that a subset $T$ of $X$ is \emph{locally constructible} in $X$ if for every $x\in X$ there exists an open neighbourhood $V$ of $x$ such that $T\cap V$ is constructible in $V$.
\end{definition}

It follows from Proposition \sref{0.9.1.8}[i] that if $V$ is such that $V\cap T$ is constructible in $V$, then for every open $W\subset V$, $W\cap T$ is constructible in $W$.
If $T$ is locally constructible in $X$, then for every open set $U$ in $X$, $T\cap U$ is locally constructible in $U$, as a result of the above remark.
The same remark shows that the set of locally constructible subsets of $X$ is stable under finite unions and finite intersections; on the other hand, it is clear that it is also stable under taking complements.

\begin{proposition}[9.1.12]
\label{0.9.1.12}
Let $X$ be a topological space.
Every constructible set in $X$ is locally constructible in $X$.
The converse is true if $X$ is quasi-compact and if its topology admits a basis formed by the retrocompact sets in $X$.
\end{proposition}

\begin{proof}
The first assertion follows from Definition \sref{0.9.1.11} and the second from \sref{0.9.1.10}.
\end{proof}

\begin{corollary}[9.1.13]
\label{0.9.1.13}
Let $X$ be a topological space whose topology admits a basis formed by the retrocompact sets in $X$.
Then every locally constructible subset $T$ of $X$ is retrocompact in $X$.
\end{corollary}

\begin{proof}
Let $U$ be a quasi-compact open set in $X$; $T\cap U$ is locally constructible in $U$, hence constructible in $U$ by Proposition~\sref{0.9.1.12}, and as a result quasi-compact by Corollary \sref{0.9.1.4}.
\end{proof}

\subsection{Constructible subsets of Noetherian spaces}
\label{subsection:0.9.2}

\begin{env}[9.2.1]
\label{0.9.2.1}
We have seen \sref{0.9.1.7} that in a Noetherian space $X$, the constructible subsets of $X$ are the \emph{finite unions of locally closed subsets of $X$}.

The inverse image of a constructible set in $X$ by a continuous map from a Noetherian space $X'$ to $X$ is constructible in $X'$.
If $Y$ is a constructible subset of a Noetherian space $X$, then the subsets of $Y$ which are constructible as subspaces of $Y$ are identical to those which are constructible as subspaces of $X$.
\end{env}

\begin{proposition}[9.2.2]
\label{0.9.2.2}
Let $X$ be an irreducible Noetherian space, $E$ a constructible subset of $X$.
For $E$ to be everywhere dense in $X$, it is necessary and sufficient for $E$ to contain a nonempty open subset of $X$.
\end{proposition}

\begin{proof}
\oldpage[0\textsubscript{III}]{15}
The condition is evidently sufficient, as every nonempty open set is dense in $X$.
Conversely, let $E=\bigcup_{i=1}^n(U_i\cap F_i)$ be a constructible subset of $X$, the $U_i$ being nonempty open sets and the $F_i$ closed in $X$; we then have $\overline{E}\subset\bigcup_i F_i$.
As a result, if $\overline{E}=X$, then $X$ is equal to one of the $F_i$, hence $E\supset U_i$, which finishes the proof.
\end{proof}

When $X$ admits a generic point $x$ \sref[0]{0.2.1.2}, the condition of Proposition \sref{0.9.2.2} is equivalent to the relation $x\in E$.

\begin{proposition}[9.2.3]
\label{0.9.2.3}
Let $X$ be a Noetherian space.
For a subset $E$ of $X$ to be constructible, it is necessary and sufficient that, for every irreducible closed subset $Y$ of $X$, $E\cap Y$ is rare in $Y$ or contains a nonempty open subset of $Y$.
\end{proposition}

\begin{proof}
The necessity of the condition follows from the fact that $E\cap Y$ must be a constructible subset of $Y$ and from Proposition \sref{0.9.2.2}, since a nondense subset of $Y$ is necessarily rare in the irreducible space $Y$ \sref[0]{0.2.1.1}.
To prove that the condition is sufficient, apply the principle of Noetherian induction \sref[0]{0.2.2.2} to the set $\mathfrak{F}$ of closed subsets $Y$ of $X$ such that $Y\cap E$ is constructible (with respect to $Y$ or with respect to $X$, which are equivalent): we can thus assume that for every closed subset $Y\neq X$ of $X$, $E\cap Y$ is constructible.
First suppose that $X$ is not irreducible, and let $X_i$ ($1\leq i\leq m$) be its irreducible components, necessarily of finite number \sref[0]{0.2.2.5}; by hypothesis the $E\cap X_i$ are constructible, hence their union $E$ is as well.
Suppose now that $X$ is irreducible; then by hypothesis, if $E$ is rare, then $\overline{E}\neq X$ and $E=E\cap\overline{E}$ is constructible; if $E$ contains a nonempty open subset $U$ of $X$, then it is the union of $U$ and $E\cap(X\setmin U)$; but $X\setmin U$ is a closed set distinct from $X$, so $E\cap(X\setmin U)$ is constructible; as a result, $E$ is itself constructible, which finishes the proof.
\end{proof}

\begin{corollary}[9.2.4]
\label{0.9.2.4}
Let $X$ be a Noetherian space, $(E_\alpha)$ an increasing filtered family of constructible subsets of $X$, such that
\begin{enumerate}
  \item[\emph{(1st)}] $X$ is the union of the family $(E_\alpha)$.
  \item[\emph{(2nd)}] Every irreducible closed subset of $X$ is contained in the closure of one of the $E_\alpha$.
\end{enumerate}

Then there exists an index $\alpha$ such that $X=E_\alpha$.

When every irreducible closed subset of $X$ admits a generic point, the hypothesis \emph{(2nd)} can be omitted.
\end{corollary}

\begin{proof}
We apply the principle of Noetherian induction \sref[0]{0.2.2.2} to the set $\mathfrak{M}$ of closed subsets of $X$ contained in at least one of the $E_\alpha$; we can thus suppose that every closed subset $Y\neq X$ of $X$ is contained in one of the $E_\alpha$.
The proposition is evident if $X$ is not irreducible, because each of the irreducible components $X_i$ of $X$ ($1\leq i\leq m$) is contained in an $E_{\alpha_i}$, and there exists an $E_\alpha$ containing all of the $E_{\alpha_i}$.
Now suppose that $X$ is irreducible.
By hypothesis, there exists a $\beta$ such that $X=\overline{E_\beta}$, so \sref{0.9.2.2} $E_\beta$ contains a nonempty open subset $U$ of $X$.
But then the closed set $X\setmin U$ is contained in an $E_\gamma$, and it suffices to take an $E_\alpha$ containing $E_\beta$ and $E_\gamma$.
When every irreducible closed subset $Y$ of $X$
\oldpage[0\textsubscript{III}]{16}
admits a generic point $y$, there exists $\alpha$ such that $y\in E_\alpha$, so $Y=\overline{\{y\}}\subset\overline{E_\alpha}$, and condition (2nd) is therefore a consequence of (1st).
\end{proof}

\begin{proposition}[9.2.5]
\label{0.9.2.5}
Let $X$ be a Noetherian space, $x$ a point of $X$, and $E$ a constructible subset of $X$.
For $E$ to be a neighbourhood of $x$, it is necessary and sufficient that for every irreducible closed subset $Y$ of $X$ containing $x$, $E\cap Y$ is dense in $Y$ (if there exists a generic point $y$ of $Y$, this also means \sref{0.9.2.2} that $y\in E$).
\end{proposition}

\begin{proof}
The condition is evidently necessary; we will prove that it is sufficient.
Applying the principle of Noetherian induction to the set $\mathfrak{M}$ of closed subsets $Y$ of $X$ containing $x$ and such that $E\cap Y$ is a neighbourhood of $x$ in $Y$, we can assume that every closed subset $Y\neq X$ of $X$ containing $x$ belongs to $\mathfrak{M}$.
If $X$ is not irreducible, then each of the irreducible components $X_i$ of $X$ containing $x$ is distinct from $X$, hence $E\cap X_i$ is a neighbourhood of $x$ with respect to $X_i$; as a result, $E$ is a neighbourhood of $x$ in the union of the irreducible components of $X$ containing $x$, and as this union is a neighbourhood of $x$ in $X$, so is $E$.
If $X$ is irreducible, then $E$ is dense in $X$ by hypothesis, so it contains a nonempty open subset $U$ of $X$ \sref{0.9.2.2}; the proposition is then evident if $x\in U$; otherwise, by hypothesis, $x$ is an interior point of $E\cap(X\setmin U)$ relative to $X\setmin U$, so the closure of $X\setmin E$ in $X$ does not contain $x$, and the complement of this closure is a neighbourhood of $x$ contained in $E$, which finishes the proof.
\end{proof}

\begin{corollary}[9.2.6]
\label{0.9.2.6}
Let $X$ be a Noetherian space, $E$ a subset of $X$.
For $E$ to be an open set in $X$, it is necessary and sufficient that for every irreducible closed subset $Y$ of $X$ intersecting $E$, $E\cap Y$ contains a nonempty open subset of $Y$.
\end{corollary}

\begin{proof}
The condition is evidently necessary; conversely, if it is satisfied, then it implies that $E$ is constructible by Proposition \sref{0.9.2.3}.
In addition, Proposition \sref{0.9.2.5} shows that $E$ is then a neighbourhood of each of its points, hence the conclusion.
\end{proof}

\subsection{Constructible functions}
\label{subsection:0.9.3}

\begin{definition}[9.3.1]
\label{0.9.3.1}
Let $h$ be a map from a topological space $X$ to a set $T$.
We say that $h$ is \emph{constructible} if $h^{-1}(t)$ is constructible for every $t\in T$, and empty except for finitely many values of $t$; then for every subset $S$ of $T$, $h^{-1}(S)$ is constructible.
We say that $h$ is \emph{locally constructible} if every $x\in X$ has an open neighbourhood $V$ such that $h|V$ is constructible.
\end{definition}

Every constructible function is locally constructible; the converse is true when $X$ is quasi-compact and admits a basis formed by the retrocompact open sets in $X$ (in particular, when $X$ is Noetherian).

\begin{proposition}[9.3.2]
\label{0.9.3.2}
Let $h$ be a map from a Noetherian space $X$ to a set $T$.
For $h$ to be constructible, it is necessary and sufficient that for every irreducible closed subset $Y$ of $X$, there exists a nonempty subset $U$ of $Y$, open with respect to $Y$, in which $h$ is constant.
\end{proposition}

\begin{proof}
The condition is necessary: indeed, by hypothesis, $h$ takes only finitely many values $t_i$ on $Y$, and each of the sets $h^{-1}(t_i)\cap Y$ is constructible in $Y$ \sref{0.9.2.1}; as they can not all be rare subsets of the space $Y$, at least one of them contains a nonempty open set \sref{0.9.2.3}.

\oldpage[0\textsubscript{III}]{17}
To see that the condition is sufficient, we apply the principle of Noetherian induction on the set $\mathfrak{M}$ of closed subsets $Y$ of $X$ such that the restriction $h|Y$ is constructible; we can thus assume that for every closed subset $Y\neq X$ of $X$,
$h|Y$ is constructible.
If $X$ is not irreducible, then the restriction of $h$ to each of the (finitely many) irreducible components $X_i$ of $X$ is constructible, and it then follows immediately from Definition \sref{0.9.3.1} that $h$ is constructible.
If $X$ is irreducible, then there exists by hypothesis a nonempty open subset $U$ of $X$ on which $h$ is constant; on the other hand, the restriction of $h$ to $X\setmin U$ is constructible by hypothesis, and it follows immediately that $h$ is constructible.
\end{proof}

\begin{corollary}[9.3.3]
\label{0.9.3.3}
Let $X$ be a Noetherian space in which every irreducible closed subset admits a generic point.
If $h$ is a map from $X$ to a set $T$ such that, for every $t\in T$, $h^{-1}(t)$ is constructible, then $h$ is constructible.
\end{corollary}

\begin{proof}
If $Y$ is an irreducible closed subset of $X$ and $y$ its generic point, then $Y\cap h^{-1}(h(y))$ is constructible and contains $y$, hence \sref{0.9.2.2} this set contains a nonempty open subset of $Y$, and it suffices to apply Proposition \sref{0.9.3.2}.
\end{proof}

\begin{proposition}[9.3.4]
\label{0.9.3.4}
Let $X$ be a Noetherian space in which every irreducible closed subset admits a generic point, $h$ a constructible map from $X$ to an ordered set.
For $h$ to be upper semi-continuous on $X$, it is necessary and sufficient that for every $x\in X$ and every generization \sref[0]{0.2.1.2} $x'$ of $x$, we have $h(x')\leq h(x)$.
\end{proposition}

\begin{proof}
The function $h$ takes only finitely many values; therefore, to say that it is upper semi-continuous means that for every $x\in X$, the set $E$ of the $y\in X$ such that $h(y)\leq h(x)$ is a neighbourhood of $x$.
By hypothesis, $E$ is a constructible subset of $X$; on the other hand, to say that an irreducible closed subset $Y$ of $X$ contains $x$ means that its generic point $y$ is a generization of $x$; the conclusion then follows from Proposition \sref{0.9.2.5}.
\end{proof}

\section{Supplement on flat modules}
\label{section:0.10}

For any proofs missing in \hyperref[subsection:0.10.1]{(10.1)} and \hyperref[subsection:0.10.2]{(10.2)}, we refer the reader to Bourbaki,~\emph{Alg.~comm.}, chap.~II and~III.

\subsection{Relations between flat modules and free modules}
\label{subsection:0.10.1}

\begin{env}[10.1.1]
\label{0.10.1.1}
Let $A$ be a ring, $\mathfrak{J}$ an ideal of $A$, and $M$ an $A$-module;
for every integer $p\geq0$, we have a canonical homomorphism of $(A/\mathfrak{J})$-modules
\[
\label{0.10.1.1.1}
  \vphi_p:(M/\mathfrak{J}M)\otimes_{A/\mathfrak{J}}(\mathfrak{J}^p/\mathfrak{J}^{p+1})\to\mathfrak{J}^pM/\mathfrak{J}^{p+1}M,
  \tag{10.1.1.1}
\]
which is evidently \emph{surjective}.
We denote by $\gr(A)=\oplus_{p\geq0}\mathfrak{J}^p/\mathfrak{J}^{p+1}$ the graded ring associated to $A$ filtered by the $\mathfrak{J}^p$, and by $\gr(M)=\oplus_{p\geq0}\mathfrak{J}^pM/\mathfrak{J}^{p+1}M$ the graded $\gr(A)$-module associated to $M$ filtered by the $\mathfrak{J}^pM$;
we then have $\gr_p(A)=\mathfrak{J}^p/\mathfrak{J}^{p+1}$, and $\gr_p(M)=\mathfrak{J}^pM/\mathfrak{J}^{p+1}M$;
the $\vphi$ define a \emph{surjective} homomorphism of graded $\gr(A)$-modules
\[
\label{0.10.1.1.2}
  \vphi:\gr_0(M)\otimes_{\gr_0(A)}\gr(A)\to\gr(M).
  \tag{10.1.1.2}
\]
\end{env}

\oldpage[0\textsubscript{III-1}]{18}
\begin{env}[10.1.2]
\label{0.10.1.2}
Suppose that \emph{one} of the following hypotheses is satisfied:
\begin{enumerate}
  \item[(i)] $\mathfrak{J}$ is nilpotent;
  \item[(ii)] $A$ is Noetherian, $\mathfrak{J}$ is contained in the radical of $A$, and $M$ is of finite type.
\end{enumerate}
Then the following properties are equivalent.
\begin{enumerate}
  \item[(a)] $M$ is a free $A$-module.
  \item[(b)] $M/\mathfrak{J}M=M\otimes_A(A/\mathfrak{J})$ is a free $(A/\mathfrak{J})$-module, and $\Tor_1^A(M,A/\mathfrak{J})=0$.
  \item[(c)] $M/\mathfrak{J}M$ is a free $(A/\mathfrak{J})$-module, and the canonical homomorphism \sref{0.10.1.1.2} is injective (and thus bijective).
\end{enumerate}
\end{env}

\begin{env}[10.1.3]
\label{0.10.1.3}
Suppose that $A/\mathfrak{J}$ is a \emph{field} (in other words, that $\mathfrak{J}$ is maximal), and that one of the hypotheses, (i) and (ii), of \sref{0.10.1.2} is satisfied.
Then the following properties are equivalent.
\begin{enumerate}
  \item[(a)] $M$ is a free $A$-module.
  \item[(b)] $M$ is a projective $A$-module.
  \item[(c)] $M$ is a flat $A$-module.
  \item[(d)] $\Tor_1^A(M,A/\mathfrak{J})=0$.
  \item[(e)] The canonical homomorphism \sref{0.10.1.1.2} is bijective.
\end{enumerate}

This result can be applied, in particular, to the following two cases:
\begin{enumerate}
  \item[(i)] $M$ is an \emph{arbitrary} module, over a local ring $A$ whose maximal ideal $\mathfrak{J}$ is \emph{nilpotent} (for example, a local Artinian ring);
  \item[(ii)] $M$ is a module \emph{of finite type} over a \emph{local Noetherian} ring.
\end{enumerate}
\end{env}

\subsection{Local flatness criteria}
\label{subsection:0.10.2}

\begin{env}[10.2.1]
\label{0.10.2.1}
With the hypotheses and notation of \sref{0.10.1.1}, consider the following conditions.
\begin{enumerate}
  \item[(a)] $M$ is a flat $A$-module.
  \item[(b)] $M/\mathfrak{J}M$ is a flat $(A/\mathfrak{J})$-module, and $\Tor_1^A(M,A/\mathfrak{J})=0$.
  \item[(c)] $M/\mathfrak{J}M$ is a flat $(A/\mathfrak{J})$-module, and the canonical homomorphism \sref{0.10.1.1.2} is bijective.
  \item[(d)] For all $n\geq1$, $M/\mathfrak{J}^nM$ is a flat $(A/\mathfrak{J}^n)$-module.
\end{enumerate}

Then we have the implications
\[
  \text{(a)}\implies\text{(b)}\implies\text{(c)}\implies\text{(d)},
\]
and, if $\mathfrak{J}$ is \emph{nilpotent}, then the four conditions are \emph{equivalent}.
This is also the case if $A$ is Noetherian and $M$ is \emph{ideally separated}, that is to say, for every ideal $\mathfrak{a}$ of $A$, the $A$-module $\mathfrak{a}\otimes_A M$ is \emph{separated} for the $\mathfrak{J}$-preadic topology.
\end{env}

\begin{env}[10.2.2]
\label{0.10.2.2}
Let $A$ be a Noetherian ring, $B$ a commutative Noetherian $A$-algebra, $\mathfrak{J}$ an ideal of $A$ such that $\mathfrak{J}B$ is contained in the radical of $B$, and $M$ a $B$-\emph{module} of finite type.
Then, when $M$ is considered as an $A$-\emph{module}, the four conditions of \sref{0.10.2.1} are \emph{equivalent}.
\oldpage[0\textsubscript{III-1}]{19}
This result applies first and foremost in the case where $A$ and $B$ are \emph{local} Noetherian rings, with the homomorphism $A\to B$ being a \emph{local} homomorphism.
More specifically, if $\mathfrak{J}$ is then the \emph{maximal} ideal of $A$, we can, in conditions~\emph{(b)} and \emph{(c)}, remove the hypothesis that $M/\mathfrak{J}M$ is flat, since it is automatically satisfied, and condition~\emph{(d)} implies that the modules $M/\mathfrak{J}^nM$ are \emph{free} over the $A/\mathfrak{J}^n$.
\end{env}

\begin{env}[10.2.3]
\label{0.10.2.3}
With the hypotheses on $A$, $B$, $\mathfrak{J}$, and $M$ from the start of \sref{0.10.2.2}, let $\widehat{A}$ be the separated completion of $A$ for the $\mathfrak{J}$-preadic topology, and $\widehat{M}$ the separated completion of $M$ for the $\mathfrak{J}B$-preadic topology.
Then, for $M$ to be a flat $A$-module, it is necessary and sufficient for $\widehat{M}$ to be a flat $\widehat{A}$-module.
\end{env}

\begin{env}[10.2.4]
\label{0.10.2.4}
Let $\rho:A\to B$ be a local homomorphism of local Noetherian rings, $k$ the residue field of $A$, and $M$ and $N$ both $B$-modules of finite type, with $N$ assumed to be \emph{$A$-flat}.
Let $u:M\to N$ be a $B$-homomorphism.
Then the following conditions are equivalent.
\begin{enumerate}
  \item[(a)] $u$ is injective, and $\Coker(u)$ is a flat $A$-module.
  \item[(b)] $u\otimes1:M\otimes_A k\to N\otimes_A k$ is injective.
\end{enumerate}
\end{env}

\begin{env}[10.2.5]
\label{0.10.2.5}
Let $\rho:A\to B$ and $\sigma:B\to C$ be local homomorphisms of local Noetherian rings, $k$ the residue field of $A$, and $M$ a $C$-module of finite type.
Suppose that $B$ is a \emph{flat} $A$-module.
Then the following conditions are equivalent.
\begin{enumerate}
  \item[(a)] $M$ is a flat $B$-module.
  \item[(b)] $M$ is a flat $A$-module, and $M\otimes_A k$ is a flat $(B\otimes_A k)$-module.
\end{enumerate}
\end{env}

\begin{proposition}[10.2.6]
\label{0.10.2.6}
Let $A$ and $B$ be local Noetherian rings, $\rho:A\to B$ a local homomorphism, $\mathfrak{J}$ an ideal of $B$ contained in the maximal ideal, and $M$ a $B$-module of finite type.
Suppose that, for all $n\geq0$, $M_n=M/\mathfrak{J}^{n+1}M$ is a flat $A$-module.
Then $M$ is a flat $A$-module.
\end{proposition}

\begin{proof}
We have to prove that, for every injective homomorphism $u:N'\to N$ of $A$-modules of finite type, $v=1\otimes u:M\otimes_A N'\to M\otimes_A N$ is injective.
But $M\otimes_A N'$ and $M\otimes_A N$ are $B$-modules of finite type, and thus separated for the $\mathfrak{J}$-preadic topology \sref[0\textsubscript{I}]{0.7.3.5};
it thus suffices to prove that the homomorphism $\widehat{v}:\widehat{M\otimes_A N'}\to\widehat{M\otimes_A N}$ of the separated completions is injective.
But $\widehat{v}=\varprojlim v_n$, where $v_n$ is the homomorphism $1\otimes u:M_n\otimes_A N'\to M_n\otimes_A N$;
since, by hypothesis, $M_n$ is $A$-flat, $v_n$ is injective for all $n$, and thus so too is $\widehat{v}$, because the functor $\varprojlim$ is left exact.
\end{proof}

\begin{corollary}[10.2.7]
\label{0.10.2.7}
Let $A$ be a Noetherian ring, $B$ a local Noetherian ring, $\rho:A\to B$ a homomorphism, $f$ an element of the maximal ideal of $B$, and $M$ a $B$-module of finite type.
Suppose that the homothety $f_M:x\to fx$ on $M$ is injective, and that $M/fM$ is a flat $A$-module.
Then $M$ is a flat $A$-module.
\end{corollary}

\begin{proof}
Let $M_i=f^iM$ for $i\geq0$;
since $f_M$ is injective, $M_i/M_{i+1}$ is isomorphic to $M/fM$, and thus $A$-flat for all $i\geq0$;
the exact sequence
\[
  0\to M_i/M_{i+1}\to M/M_{i+1}\to M/M_i\to 0
\]
gives us, by induction on $i$, that $M/M_i$ is \emph{$A$-flat} for all $i\geq0$ \sref[0\textsubscript{I}]{0.6.1.2};
we can thus apply \sref{0.10.2.6}.
We can also argue directly as follows:
for every $A$-module $N$ of finite type, $M\otimes_A N$ is a $B$-module of finite type;
since $f$ belongs to the radical $\mathfrak{n}$ of $B$, the $(f)$-adic topology on $M\otimes_A N$ is finer than the
\oldpage[0\textsubscript{III-1}]{20}
$\mathfrak{n}$-adic topology, and we know that the latter is \emph{separated} \sref[0\textsubscript{I}]{0.7.3.5}.
Now, since $M/M_i$ is $A$-flat, we have that
\[
  f^i(M\otimes_A N)=\Im(M_i\otimes_A N\to M\otimes_A N)=\Ker(M\otimes_A N\to(M/M_i)\otimes_A N)
\]
by \sref[0\textsubscript{I}]{0.6.1.2}.
So let $N$ be an $A$-module of finite type, and $N'$ a submodule of $N$, with canonical injection $j:N'\to N$; in the commutative diagram
\[
  \xymatrix{
    M\otimes_A N'\ar[r]\ar[d]_{1_M\otimes j} &
    (M/M_i)\otimes_A N'\ar[d]^{1_{M/M_i}\otimes j}\\
    M\otimes_A N\ar[r] &
    (M/M_i)\otimes_A N
  }
\]
$1_{M/M_i}\otimes j$ is injective, because $M/M_i$ is $A$-flat; we thus conclude that
\[
  \Ker(M\otimes_A N'\to M\otimes_A N)\subset\Ker(M\otimes_A N'\to(M/M_i)\otimes_A N')
\]
for any $i$; since the intersection of the latter kernels is $0$, as we saw above, so too is the former kernel, and so $M$ is $A$-flat.
\end{proof}

\begin{proposition}[10.2.8]
\label{0.10.2.8}
Let $A$ be a reduced Noetherian ring, and $M$ an $A$-module of finite type.
Suppose that, for every $A$-algebra $B$ which is a discrete valuation ring, $M\otimes_A B$ is a flat $B$-module (and thus free \sref{0.10.1.3}).
Then $M$ is a flat $A$-module.
\end{proposition}

\begin{proof}
We know that, for $M$ to be flat, it is necessary and sufficient for $M_\mathfrak{m}$ to be a flat $A_\mathfrak{m}$-module for every maximal ideal $\mathfrak{m}$ of $A$ \sref[0\textsubscript{I}]{0.6.3.3};
we can thus restrict to the case where $A$ is \emph{local} \sref[0\textsubscript{I}]{0.1.2.8}.
So let $\mathfrak{m}$ be the maximal ideal of $A$, $\mathfrak{p}_i$ ($1\leq i\leq r$) the minimal prime ideals of $A$, and $k$ the residue field $A/\mathfrak{m}$.
We know \sref[II]{II.7.1.7} that there exists, for each $i$, a discrete valuation ring $B_i$ that has the same field of fractions $K_i$ as the integral domain $A/\mathfrak{p}_i$, and that, further, dominates $A/\mathfrak{p}_i$.
Let $M_i=M\otimes_A B_i$.
By hypothesis, $M_i$ is free over $B_i$, and so, denoting by $k_i$ the residue field of $B_i$, we have
\[
\label{0.10.2.8.1}
  \rg_{k_i}(M_i\otimes_{B_i}k_i)=\rg_{K_i}(M_i\otimes_{B_i}K_i).
  \tag{10.2.8.1}
\]

But it is clear that the composite homomorphism $A\to A/\mathfrak{p}_i\to B_i$ is local, and so $k_i$ is an extension of $k$, and that we have $M_i\otimes_{B_i}k_i = M\otimes_A k_i = (M\otimes_A k)\otimes_k k_i$, and also that $M_i\otimes_{B_i}K_i = M\otimes_A K_i$.
Equation~\sref{0.10.2.8.1} thus implies that
\[
  \rg_k(M\otimes_A k)=\rg_{K_i}(M\otimes_A K_i)\quad\mbox{for $1\leq i\leq r$}
\]
and since $A$ is reduced, we know that this condition implies that $M$ is a \emph{free} $A$-module (Bourbaki, \emph{Alg. comm.}, chap.~II, \textsection~3, n\textsuperscript{o}~2, prop.~7).
\end{proof}

\subsection{Existence of flat extensions of local rings}
\label{subsection:0.10.3}

\begin{proposition}[10.3.1]
\label{0.10.3.1}
Let $A$ be a local Noetherian ring, with maximal ideal $\mathfrak{J}$, and residue field $k=A/\mathfrak{J}$.
Let $K$ be a field extension of $k$.
Then there exists a local homomorphism from $A$ to a local Noetherian ring $B$, such that $B/\mathfrak{J}B$ is $k$-isomorphic to $K$, and such that $B$ is a flat $A$-module.
\end{proposition}

The rest of this section is devoted to proving this proposition, step-by-step.
\begin{env}[10.3.1.1]
\label{0.10.3.1.1}
First suppose that $K=k(T)$, where $T$ is an indeterminate.
In the ring of polynomials $A'=A[T]$, consider the prime ideal $\mathfrak{J}'=\mathfrak{J}A'$, consisting of the
\oldpage[0\textsubscript{III-1}]{21}
polynomials that have coefficients in the ideal $\mathfrak{J}$;
it is clear that $A'/\mathfrak{J}'$ is canonically isomorphic to $k[T]$.
We will show that the ring of fractions $B=A'_{\mathfrak{J}'}$ is that for which we are searching (that is, a ring which satisfies the conditions of the conclusion of the proposition);
it is clearly a local Noetherian ring, with maximal ideal $\mathfrak{L}=\mathfrak{J}B$.
Further, $B/\mathfrak{L}=(A'/\mathfrak{J}')_{\mathfrak{J}'}=(k[T])_{\mathfrak{J}'}$ is exactly the field of fractions $K$ of $k[T]$.
Finally, $B$ is a flat $A'$-module, and $A'$ is a free $A$-module, so $B$ is a flat $A$-module \sref[0\textsubscript{I}]{0.6.2.1}.
\end{env}

\begin{env}[10.3.1.2]
\label{0.10.3.1.2}
Now suppose that $K=k(t)=k[t]$, where $t$ is algebraic over $k$;
let $f\in k[T]$ be the minimal polynomial of $t$;
there exists a monic polynomial $F\in A[T]$ whose canonical image in $k[T]$ is $f$.
So let $A'=A[T]$, and let $\mathfrak{J}'$ be the ideal $\mathfrak{J}A'+(F)$ in $A'$.
We will see that the quotient ring $B=A'/(F)$ is that for which we are searching.
First of all, it is clear that $B$ is a \emph{free} $A$-module, and thus flat.
The ring $A'/\mathfrak{J}'$ is isomorphic to
\[
  (A'/\mathfrak{J}A')/\big((\mathfrak{J}A'+(F))/\mathfrak{J}A'\big)=k[T]/(f)=K;
\]
the image $\mathfrak{L}$ of $\mathfrak{J}'$ in $B$ is thus maximal, and we evidently have that $\mathfrak{L}=\mathfrak{J}B$.
Finally, $B$ is a semi-local ring, because it is an $A$-module of finite type (Bourbaki, \emph{Alg. comm.}, chap.~IV, \textsection~2, n\textsuperscript{o}~5, cor.~3 of prop.~9), and its maximal ideals are in bijective correspondence with those of $B/\mathfrak{J}B$ (\cite[vol.~I, p.~259]{I-13});
the previous arguments then prove that $B$ is a local ring.
\end{env}

\begin{lemma}[10.3.1.3]
\label{0.10.3.1.3}
Let $(A_\lambda,f_{\mu\lambda})$ be a filtered inductive system of local rings, such that the $f_{\mu\lambda}$ are local homomorphisms;
let $\mathfrak{m}_\lambda$ be the maximal ideal of $A_\lambda$, and let $K_\lambda=A_\lambda/\mathfrak{m}_\lambda$.
Then $A'=\varinjlim A_\lambda$ is a local ring, with maximal ideal $\mathfrak{m}'=\varinjlim\mathfrak{m}_\lambda$, and residue field $K=\varinjlim K_\lambda$.
Further, if $\mathfrak{m}_\mu=\mathfrak{m}_\lambda A_\mu$ with $\lambda<\mu$, then we have $\mathfrak{m}'=\mathfrak{m}_\lambda A'$ for all $\lambda$.
If, further, for $\lambda<\mu$, $A_\mu$ is a flat $A_\lambda$-module, and if all the $A_\lambda$ are Noetherian, then $A'$ is Noetherian and is a flat $A_\lambda$-module for all $\lambda$.
\end{lemma}

\begin{proof}
Since, by hypothesis, $f_{\mu\lambda}(\mathfrak{m}_\lambda)\subset\mathfrak{m}_\mu$ for $\lambda<\mu$, the $\mathfrak{m}_\lambda$ form an inductive system, and its limit $\mathfrak{m}'$ is evidently an ideal of $A'$.
Further, if $x'\not\in\mathfrak{m}'$, there exists a $\lambda$ such that $x'=f_\lambda(x_\lambda)$ for some $x_\lambda\in A_\lambda$ (where $f_\lambda:A_\lambda\to A'$ denotes the canonical homomorphism);
because $x'\not\in\mathfrak{m}'$, we necessarily have that $x_\lambda\not\in\mathfrak{m}_\lambda$, and so $x_\lambda$ admits an inverse $y_\lambda$ in $A_\lambda$, and $y'=f_\lambda(y_\lambda)$ is the inverse of $x'$ in $A'$, which proves that $A'$ is a local ring with maximal ideal $\mathfrak{m}'$;
the claim about $K$ follows immediately from the fact that $\varinjlim$ is an exact functor.
The hypothesis that $\mathfrak{m}_\mu=\mathfrak{m}_\lambda A_\mu$ implies that the canonical map $\mathfrak{m}_\lambda\otimes_{A_\lambda}A_\mu\to\mathfrak{m}_\mu$ is surjective;
the equality $\mathfrak{m}'=\mathfrak{m}_\lambda A'$ then follows from, again, the fact that the functor $\varinjlim$ is exact and commutes with the tensor product.

Now suppose that, for $\lambda<\mu$, we have $\mathfrak{m}_\mu=\mathfrak{m}_\lambda A_\mu$, and that $A_\mu$ is a flat $A_\lambda$-module.
Then $A'$ is a flat $A_\lambda$-module for all $\lambda$, by \sref[0\textsubscript{I}]{0.6.2.3};
since $A'$ and $A_\lambda$ are local rings, and since $\mathfrak{m}'=\mathfrak{m}_\lambda A'$, $A'$ is even a \emph{faithfully flat} $A_\lambda$-module \sref[0\textsubscript{I}]{0.6.6.2}.
Finally, suppose further that the $A_\lambda$ are \emph{Noetherian};
the $\mathfrak{m}_\lambda$-adic topologies are then separated \sref[0\textsubscript{I}]{0.7.3.5};
we now show that, from this, it follows that the $\mathfrak{m}'$-adic topology on $A'$ is \emph{separated}.
Indeed, if $x'\in A'$ belongs to all the $\mathfrak{m}^{'n}$ ($n>0$), then it is the image of some $x_\mu\in A_\mu$ for a specific index $\mu$, and since the inverse image in $A_\mu$ of $\mathfrak{m}^{'n}=\mathfrak{m}_\mu^n A'$ is $\mathfrak{m}_\mu^n$ \sref[0\textsubscript{I}]{0.6.6.1}, $x_\mu$ belongs to all the $\mathfrak{m}_\mu^n$, so $x_\mu=0$, by hypothesis, and so $x'=0$.
Denote by $\widehat{A}'$ the completion of $A'$ for the $\mathfrak{m}'$-adic topology;
the above shows that we have $A'\subset\widehat{A}'$.
We will now show that $\widehat{A}'$ is \emph{Noetherian} and \emph{$A_\lambda$-flat} for all $\lambda$;
from this,
\oldpage[0\textsubscript{III-1}]{22}
it will follow that $\widehat{A}'$ is $A'$-flat \sref[0\textsubscript{I}]{0.6.2.3}, and since $\mathfrak{m}'\widehat{A}'\neq\widehat{A}'$, that $\widehat{A}'$ is a faithfully flat $A'$-module \sref[0\textsubscript{I}]{0.6.6.2}, whence the final conclusion that $A'$ is \emph{Noetherian} \sref[0\textsubscript{I}]{0.6.5.2}, which will finish the proof of the lemma.

We have $\widehat{A}'=\varprojlim_n A'/\mathfrak{m}^{\prime n}$;
by the fact that $A'$ is $A_\lambda$-flat, we have that
\[
  \mathfrak{m}^{\prime n}/\mathfrak{m}^{\prime n+1}=(\mathfrak{m}_\lambda^{n}/\mathfrak{m}_\lambda^{n+1})\otimes_{A_\lambda}A'=(\mathfrak{m}_\lambda^{n}/\mathfrak{m}_\lambda^{n+1})\otimes_{K_\lambda}(K_\lambda\otimes_{A_\lambda}A')=(\mathfrak{m}_\lambda^{n}/\mathfrak{m}_\lambda^{n+1})\otimes_{K_\lambda}K;
\]
since $\mathfrak{m}_\lambda^{n}/\mathfrak{m}_\lambda^{n+1}$ is a $K_\lambda$-vector space of finite dimension, $\mathfrak{m}^{\prime n}/\mathfrak{m}^{\prime n+1}$ is a $K$-vector space of finite dimension for all $n\geq 0$.
It thus follows from \sref[0\textsubscript{I}]{0.7.2.12} and \sref[0\textsubscript{I}]{0.7.2.8} that $\widehat{A}'$ is \emph{Noetherian}.
We further know that the maximal ideal of $\widehat{A}'$ is $\mathfrak{m}'\widehat{A}'$, and that $\widehat{A}'/\mathfrak{m}^{\prime n}\widehat{A}'$ is isomorphic to $A'/\mathfrak{m}^{\prime n}$;
since $A'/\mathfrak{m}^{\prime n}=(A_\lambda/\mathfrak{m}_\lambda^n)\otimes_{A_\lambda}A'$, we see that $A'/\mathfrak{m}^{\prime n}$ is a flat $(A_\lambda/\mathfrak{m}_\lambda^n)$-module \sref[0\textsubscript{I}]{0.6.2.1};
criterion~\sref{0.10.2.2} is thus applicable to the Noetherian $A_\lambda$-algebra $\widehat{A}'$, and shows that $\widehat{A}'$ is \emph{$A_\lambda$-flat}.
\end{proof}

\begin{env}[10.3.1.4]
\label{0.10.3.1.4}
We now treat the general case.
There exists an ordinal $\gamma$ and, for every ordinal $\lambda\leq\gamma$, a subfield $k_\lambda$ of $K$ that contains $k$, such that (i) for all $\lambda<\gamma$, $k_{\lambda+1}$ is an extension of $k_\lambda$ generated by a single element; (ii) for every limit ordinal $\mu$, $k_\mu=\bigcup_{\lambda<\mu}k_\lambda$; and (iii) $K=k_\gamma$.
In fact, it suffices to consider a bijection $\xi\mapsto t_\xi$ from the set of ordinals $\xi\leq\beta$ (for some suitable $\beta$) to $K$, and to define $k_\lambda$ by transfinite induction (for $\lambda\leq\beta$) as the union of the $k_\mu$ for $\mu<\lambda$ if $\lambda$ is a limit ordinal, and as $k_\nu(t_\xi)$ if $\lambda=\nu+1$, where $\xi$ is the smallest ordinal such that $t_\xi\not\in k_\nu$;
$\gamma$ is then, by definition, the smallest ordinal $\leq\beta$ such that $k_\gamma=K$.

With this in mind, we will define, by transfinite induction, a family of local Noetherian rings $A_\lambda$ for $\lambda\leq\gamma$, and local homomorphisms $f_{\mu\lambda}:A_\lambda\to A_\mu$ for $\lambda\leq\mu$, satisfying the following conditions:
\begin{enumerate}
  \item[(i)] $(A_\lambda,f_{\mu\lambda})$ is an inductive system, and $A_0=A$;
  \item[(ii)] for all $\lambda$, we have a $k$-isomorphism $A_\lambda/\mathfrak{J}A_\lambda\isoto k_\lambda$;
  \item[(iii)] for $\lambda\leq\mu$, $A_\mu$ is a flat $A_\lambda$-module.
\end{enumerate}

So suppose that the $A_\lambda$ and the $f_{\mu\lambda}$ are defined for $\lambda<\mu<\xi$, and suppose, first of all, that $\xi=\zeta+1$, so that $k_\xi=k_\zeta(t)$.
If $t$ is transcendental over $k_\zeta$, we define $A_\xi$, following the procedure of \sref{0.10.3.1.1}, to be equal to $(A_\zeta[t])_{\mathfrak{J}A_\zeta[t]}$;
the canonical map is $f_{\xi\zeta}$, and, for $\lambda<\zeta$, we take $f_{\xi\lambda}=f_{\xi\zeta}\circ f_{\zeta\lambda}$;
the verification of conditions~(i) to (iii) is then immediate, given that what we have shown in \sref{0.10.3.1.1}.
So now suppose that $t$ is algebraic, and let $h$ be its minimal polynomial in $k_\zeta[T]$, and $H$ a monic polynomial in $A_\zeta[T]$ whose image in $k_\zeta[T]$ is $h$;
we then take $A_\xi$ to be equal to $A_\zeta[T]/(H)$, with the $f_{\xi\lambda}$ being defined as before;
the verification of conditions~(i) to (iii) then follows from what we have shown in \sref{0.10.3.1.2}.

Now suppose that $\xi$ has no predecessor;
we then take $A_\xi$ to be the inductive limit of the inductive system of local rings $(A_\lambda,f_{\mu\lambda})$ for $\lambda<\xi$;
we define $f_{\xi\lambda}$ as the canonical map for $\lambda<\xi$.
The fact that $A_\xi$ is local and Noetherian, that the $f_{\xi\lambda}$ are local homomorphisms, and that conditions~(i) to (iii) are satisfied for $\lambda\leq\xi$
\oldpage[0\textsubscript{III-1}]{23}
then follows from the induction hypothesis, and from Lemma~\sref{0.10.3.1.3}.
With this construction, it is clear that the ring $B=A_\gamma$ satisfies the conditions of \sref{0.10.3.1}.\qed
\end{env}

We note that, by \sref{0.10.2.1}[c], we have a canonical isomorphism
\[
\label{0.10.3.1.5}
  \gr(A)\otimes_k K\simto\gr(B).
  \tag{10.3.1.5}
\]

We can also replace $B$ by its $\mathfrak{J}B$-adic completion $\widehat{B}$ without changing the conclusions of \sref{0.10.3.1}, because $\widehat{B}$ is a flat $B$-module \sref[0\textsubscript{I}]{0.7.3.3}, and thus a flat $A$-module \sref[0\textsubscript{I}]{0.6.2.1}.

We have also shown the following:
\begin{corollary}[10.3.2]
\label{0.10.3.2}
If $K$ is an extension of finite degree, then we can assume that $B$ is a finite $A$-algebra.
\end{corollary}

\section{Supplement on homological algebra}
\label{section:0.11}

\subsection{Review of spectral sequences}
\label{subsection:0.11.1}

\begin{env}[11.1.1]
\label{0.11.1.1}
In the following, we use a more general notion of a spectral sequence than that defined in (T, 2.4); keeping the notation of (T, 2.4), we call a \emph{spectral sequence} in an abelian category $\cat{C}$ a system $E$ consisting of the following parts:
\begin{enumerate}
  \item[(a)] A family $(E_r^{pq})$ of objects of $\cat{C}$ defined for $p,q\in\bb{Z}$ and $r\geq 2$.
  \item[(b)] A family of morphisms $d_r^{pq}:E_r^{pq}\to E_r^{p+r,q-r+1}$ such that $d_r^{p+r,q-r+1}d_r^{pq}=0$.
    We set $Z_{r+1}(E_r^{pq})=\Ker(d_r^{pq})$ and $B_{r+1}(E_r^{pq})=\Im(d_r^{p-r,q+r-1})$, so that
    \[
      B_{r+1}(E_r^{pq})\subset Z_{r+1}(E_r^{pq})\subset E_r^{pq}.
    \]
  \item[(c)] A family of isomorphisms $\alpha_r^{pq}:Z_{r+1}(E_r^{pq})/B_{r+1}(E_r^{pq})\isoto E_{r+1}^{pq}$.

  We then define for $k\geq r+1$, by induction on $k$, the subobjects $B_k(E_r^{pq})$ and $Z_k(E_r^{pq})$ as the inverse images, under the canonical morphism $E_r^{pq}\to E_r^{pq}/B_{r+1}(E_r^{pq})$ of the subobjects of this quotient identified via $\alpha_r^{pq}$ with the subobjects $B_k(E_{r+1}^{pq})$ and $Z_k(E_{r+1}^{pq})$ respectively.
It is clear that we then have, up to isomorphism,
  \[
  \label{0.11.1.1.1}
    Z_k(E_r^{pq})/B_k(E_r^{pq})=E_k^{pq}\text{ for }k\geq r+1,
    \tag{11.1.1.1}
  \]
  and, if we set $B_r(E_r^{pq})=0$ and $Z_r(E_r^{pq})=E_r^{pq}$, then we have the inclusion relations
  \[
  \label{0.11.1.1.2}
    0=B_r(E_r^{pq})\subset B_{r+1}(E_r^{pq})\subset B_{r+2}(E_r^{pq})\subset\cdots\subset Z_{r+2}(E_r^{pq})\subset Z_{r+1}(E_r^{pq})\subset Z_r(E_r^{pq})=E_r^{pq}.
    \tag{11.1.1.2}
  \]
  The other parts of the data of $E$ are then:
  \item[(d)] Two subobjects $B_\infty(E_2^{pq})$ and $Z_\infty(E_2^{pq})$ of $E_2^{pq}$ such that we have $B_\infty(E_2^{pq})\subset Z_\infty(E_2^{pq})$ and, for every $k\geq 2$,
    \[
      B_k(E_2^{pq})\subset B_\infty(E_2^{pq})\text{ and }Z_\infty(E_2^{pq})\subset Z_k(E_2^{pq}).
    \]
    We set
    \[
    \label{0.11.1.1.3}
      E_\infty^{pq}=Z_\infty(E_2^{pq})/B_\infty(E_2^{pq}).
      \tag{11.1.1.3}
    \]
  \item[(e)]
\oldpage[0\textsubscript{III}]{24}
    A family $(E^n)$ of objects of $\cat{C}$, each equipped with a \emph{decreasing filtration $(F^p(E^n))_{p\in\bb{Z}}$}.
    As usual, we denote by $\gr(E^n)$ the graded object associated to the filtered object $E^n$, the direct sum of the $\gr_p(E^n)=F^p(E^n)/F^{p+1}(E^n)$.
  \item[(f)] For every pair $(p,q)\in\bb{Z}\times\bb{Z}$, an isomorphism $\beta^{pq}:E_\infty^{pq}\isoto\gr_p(E^{p+q})$.
\end{enumerate}

The family $(E^n)$, without the filtrations, is called the \emph{abutment} (or \emph{limit}) of the spectral sequence $E$.

Suppose that the category $\cat{C}$ admits infinite direct sums, or that for every $r\geq 2$ and every $n\in\bb{Z}$, there are finitely many pairs $(p,q)$ such that $p+q=n$ and $E_r^{pq}\neq 0$ (it suffices for it to hold for $r=2$).
Then the $E_r^{(n)}=\sum_{p+q=n}E_r^{pq}$ are defined, and if we denote by $d_r^{(n)}$ the morphism $E_r^{(n)}\to E_r^{(n+1)}$ whose restriction to $E_r^{pq}$ is $d_r^{pq}$ for every pair $(p,q)$ such that $p+q=n$, then $d_r^{(n+1)}\circ d_r^{(n)}=0$, in other words, $(E_r^{(n)})_{n\in\bb{Z}}$ is a \emph{complex $E_r^{(\bullet)}$} in $\cat{C}$, with differential of degree $+1$, and it follows from (c) that $\HH^n(E_r^{(\bullet)})$ is \emph{isomorphic to $E_{r+1}^{(n)}$} for every $r\geq 2$.
\end{env}

\begin{env}[11.1.2]
\label{0.11.1.2}
A \emph{morphism $u:E\to E'$} from a spectral sequence $E$ to a spectral sequence $E'=(E_r^{\prime pq},E^{\prime n})$ consists of systems of morphisms $u_r^{pq}:E_r^{pq}\to E_r^{\prime pq}$ and $u^n:E^n\to E^{\prime n}$, the $u^n$ compatible with the filtrations on $E^n$ and $E^{\prime n}$, and the diagrams
\[
  \xymatrix{
    E_r^{pq}\ar[r]^-{d_r^{pq}}\ar[d]_{u_r^{pq}} &
    E_r^{p+r,q-r+1}\ar[d]^{u_r^{p+r,q-r+1}}\\
    E_r^{\prime pq}\ar[r]_{d_r^{\prime pq}} &
    E_r^{\prime p+r,q-r+1}
  }
\]
being commutative; in addition, by passing to quotients, $u_r^{pq}$ gives a morphism $\overline{u}_r^{pq}:Z_{r+1}(E_r^{pq})/B_{r+1}(E_r^{pq})\to Z_{r+1}(E_r^{\prime pq})/B_{r+1}(E_r^{\prime pq})$ and we must have $\alpha_r^{\prime pq}\circ\overline{u}_r^{pq}=u_{r+1}^{pq}\circ\alpha_r^{pq}$; finally, we must have $u_2^{pq}(B_\infty(E_2^{pq}))\subset B_\infty(E_2^{\prime pq})$ and $u_2^{pq}(Z_\infty(E_2^{pq}))\subset Z_\infty(E_2^{\prime pq})$; by passing to quotients, $u_2^{pq}$ then gives a morphism $u_\infty^{pq}:E_\infty^{pq}\to E_\infty^{\prime pq}$, and the diagram
\[
  \xymatrix{
    E_\infty^{pq}\ar[r]^-{u_\infty^{pq}}\ar[d]_{\beta^{pq}} &
    E_\infty^{\prime pq}\ar[d]^{\beta^{\prime pq}}\\
    \gr_p(E^{p+q})\ar[r]_{\gr_p(u^{p+q})} &
    \gr_p(E^{\prime p+q})
  }
\]
must be commutative.

The above definitions show, by induction on $r$, that if the $u_2^{pq}$ are \emph{isomorphisms}, then so are the $u_r^{pq}$ for $r\geq 2$; \emph{if in addition we know that $u_2^{pq}(B_\infty(E_2^{pq}))=B_\infty(E_2^{\prime pq})$ and $u_2^{pq}(Z_\infty(E_2^{pq}))=Z_\infty(E_2^{\prime pq})$ and the $u^n$ are isomorphisms}, then we can conclude that $u$ is an \emph{isomorphism}.
\end{env}

\begin{env}[11.1.3]
\label{0.11.1.3}
\oldpage[0\textsubscript{III}]{25}
Recall that if $(F^p(X))_{p\in\bb{Z}}$ is a (decreasing) \emph{filtration} of an object $X\in\cat{C}$, then we say that this filtration is \emph{separated} if $\inf(F^p(X))=0$, \emph{discrete} if there exists a $p$ such that $F^p(X)=0$, \emph{exhaustive} (or \emph{coseparated}) if $\sup(F^p(X))=X$, \emph{codiscrete} if there exists a $p$ such that $F^p(X)=X$.

We say that a spectral sequence $E=(E_r^{pq},E^n)$ is \emph{weakly convergent} if we have $B_\infty(E_2^{pq})=\sup_k(B_k(E_2^{pq}))$ and $Z_\infty(E_2^{pq})=\inf_k(Z_k(E_2^{pq}))$ (in other words, the objects $B_\infty(E_2^{pq})$ and $Z_\infty(E_2^{pq})$ are determined from the data (a)--(c) of the spectral sequence $E$).
We say that the spectral sequence $E$ is \emph{regular} if it is weakly convergent and if, in addition:
\begin{enumerate}
  \item[(1st)] For every pair $(p,q)$, the decreasing sequence $(Z_k(E_2^{pq}))_{k\geq 2}$ is \emph{stable}; the hypothesis that $E$ is weakly convergent then implies that $Z_\infty(E_2^{pq})=Z_k(E_2^{pq})$ for $k$ large enough (depending on $p$ and $q$).
  \item[(2nd)] For every $n$, the filtration $(F^p(E^n))_{p\in\bb{Z}}$ of $E^n$ is \emph{discrete} and \emph{exhaustive}.
\end{enumerate}

We say that the spectral sequence $E$ is \emph{coregular} if it is weakly convergent and if, in addition:
\begin{enumerate}
  \item[(3rd)] For every pair $(p,q)$, the increasing sequence $(B_k(E_2^{pq}))_{k\geq 2}$ is \emph{stable}, which implies that $B_\infty(E_2^{pq})=B_k(E_2^{pq})$, and as a result, $E_\infty^{pq}=\inf E_k^{pq}$.
  \item[(4th)] For every $n$, the filtration of $E^n$ is \emph{codiscrete}.
\end{enumerate}

Finally, we say that $E$ is \emph{biregular} if it is both regular and coregular, in other words if we have the following conditions:
\begin{enumerate}
  \item[(a)] For every pair $(p,q)$, the sequences $(B_k(E_2^{pq}))_{k\geq 2}$ and $(Z_k(E_2^{pq}))_{k\geq 2}$ are \emph{stable} and we have $B_\infty(E_2^{pq})=B_k(E_2^{pq})$ and $Z_\infty(E_2^{pq})=Z_k(E_2^{pq})$ for $k$ large enough (which implies that $E_\infty^{pq}=E_k^{pq}$).
  \item[(b)] For every $n$, the filtration $(F^p(E^n))_{p\in\bb{Z}}$ is \emph{discrete} and \emph{codiscrete} (which we also call \emph{finite}).
\end{enumerate}

The spectral sequences defined in (T, 2.4) are thus biregular spectral sequences.
\end{env}

\begin{env}[11.1.4]
\label{0.11.1.4}
Suppose that in the category $\cat{C}$, filtered inductive limits exist and the functor $\varinjlim$ is \emph{exact} (which is equivalent to saying that the axiom (AB~5) of (T, 1.5) is satisfied (cf.~T, 1.8)).
The condition that the filtration $(F^p(X))_{p\in\bb{Z}}$ of an object $X\in\cat{C}$ is exhaustive is then expressed as $\varinjlim_{p\to-\infty}F^p(X)=X$.
If a spectral sequence $E$ is weakly convergent, then we have $B_\infty(E_2^{pq})=\varinjlim_{k\to\infty}B_k(E_2^{pq})$; if in addition $u:E\to E'$ is a morphism from $E$ to a weakly convergent spectral sequence $E'$ in $\cat{C}$, then we have $u_2^{pq}(B_\infty(E_2^{pq}))=B_\infty(E_2^{\prime pq})$, by the exactness of $\varinjlim$.
In addition:
\end{env}

\begin{proposition}[11.1.5]
\label{0.11.1.5}
Let $\cat{C}$ be an abelian category in which filtered inductive limits are exact, $E$ and $E'$ two regular spectral sequences in $\cat{C}$, $u:E\to E'$ a morphism of spectral sequences.
If the $u_2^{pq}$ are isomorphisms, then so is $u$.
\end{proposition}

\begin{proof}
We already know \sref{0.11.1.2} that the $u_r^{pq}$ are isomorphisms and that
\[
  u_2^{pq}(B_\infty(E_2^{pq}))=B_\infty(E_2^{\prime pq});
\]
\oldpage[0\textsubscript{III}]{26}
the hypothesis that $E$ and $E'$ are regular also implies that $u_2^{pq}(Z_\infty(E_2^{pq}))=Z_\infty(E_2^{\prime pq})$, and as $u_2^{pq}$ is an isomorphism, so is $u_\infty^{pq}$; we thus conclude that $\gr_p(u^{p+q})$ is also an isomorphism.
But as the filtrations of the $E^n$ and the $E^{\prime n}$ are discrete and exhaustive, this implies that the $u^n$ are also isomorphisms (Bourbaki, \emph{Alg. comm.}, chap.~III, \textsection2, n\textsuperscript{o}8, th.~1).
\end{proof}

\begin{env}[11.1.6]
\label{0.11.1.6}
It follows from (\hyperref[0.11.1.1.2]{11.1.1.2}) and the definition (\hyperref[0.11.1.1.3]{11.1.1.3}) that if, for a spectral sequence $E$, we have $E_r^{pq}=0$, then we have $E_k^{pq}=0$ for $k\geq r$ and $E_\infty^{pq}=0$.
We say that a spectral sequence \emph{degenerates} if there exists an integer $r\geq 2$ and, for every integer $n\in\bb{Z}$, an integer $q(n)$ such that \emph{$E_r^{n-q,q}=0$ for every $q\neq q(n)$}.
We first deduce from the previous remark that we also have $E_k^{n-q,q}=0$ for $k\geq r$ (including $k=\infty$) and $q\neq q(n)$.
In addition, the definition of $E_{r+1}^{pq}$ shows that we have $E_{r+1}^{n-q(n),q(n)}=E_r^{n-q(n),q(n)}$; if $E$ is \emph{weakly convergent}, then we also have $E_\infty^{n-q(n),q(n)}=E_r^{n-q(n),q(n)}$; in other words, for every $n\in\bb{Z}$, $\gr_p(E^n)=0$ for $p\neq n-q(n)$ and $\gr_{n-q(n)}(E^n)=E_r^{n-q(n),q(n)}$.
If in addition the filtration of $E^n$ is \emph{discrete} and \emph{exhaustive}, then the spectral sequence $E$ is \emph{regular}, and we have $E^n=E_r^{n-q(n),q(n)}$ up to isomorphism.
\end{env}

\begin{env}[11.1.7]
\label{0.11.1.7}
Suppose that filtered inductive limits exist and are exact in the category $\cat{C}$, and let $(E_\lambda,u_{\mu\lambda})$ be an inductive system (over a filtered set of indices) of spectral sequences in $\cat{C}$.
Then the \emph{inductive limit} of this inductive system exists in the additive category of spectral sequences of objects of $\cat{C}$: to see this, it suffices to define $E_r^{pq}$, $d_r^{pq}$, $\alpha_r^{pq}$, $B_\infty(E_2^{pq})$, $Z_\infty(E_2^{pq})$, $E^n$, $F^p(E^n)$, and $\beta^{pq}$ as the respective inductive limits of the $E_{r,\lambda}^{pq}$, $d_{r,\lambda}^{pq}$, $\alpha_{r,\lambda}^{pq}$, $B_\infty(E_{2,\lambda}^{pq})$, $Z_\infty(E_{2,\lambda}^{pq})$, $E_\lambda^n$, $F^p(E_\lambda^n)$, and $\beta_\lambda^{pq}$; the verification of the conditions of \sref{0.11.1.1} follows from the exactness of the functor $\varinjlim$ on $\cat{C}$.
\end{env}

\begin{remark}[11.1.8]
\label{0.11.1.8}
Suppose that the category $\cat{C}$ is the category of $A$-modules over a \emph{Noetherian} ring $A$ (resp. a ring $A$).
Then the definitions of \sref{0.11.1.1} show that if, for a given $r$, the $E_r^{pq}$ are $A$-modules \emph{of finite type} (resp. \emph{of finite length}), then so are each of the modules $E_s^{pq}$ for $s\geq r$, hence so is $E_\infty^{pq}$.
If in addition the filtration of the abutment $(E^n)$ is \emph{discrete} and \emph{codiscrete} for all $n$, then we conclude that each of the $E^n$ is also an $A$-module \emph{of finite type} (resp. \emph{of finite length}).
\end{remark}

\begin{env}[11.1.9]
\label{0.11.1.9}
We will have to consider conditions which ensure that a spectral sequence $E$ is biregular in a ``uniform'' way in $p+q=n$.
We will then use the following lemma:
\end{env}

\begin{lemma}[11.1.10]
\label{0.11.1.10}
Let $(E_r^{pq})$ be a family of objects of $\cat{C}$ related by the data of \emph{(a)}, \emph{(b)}, and \emph{(c)} of \sref{0.11.1.1}.
For a fixed integer $n$, the following properties are equivalent:
\begin{enumerate}
  \item[{\rm(a)}] There exists an integer $r(n)$ such that for $r\geq r(n)$, $p+q=n$ or $p+q=n-1$, the morphisms $d_r^{pq}$ are zero.
  \item[{\rm(b)}] There exists an integer $r(n)$ such that for $p+q=n$ or $p+q=n+1$, we have $B_r(E_2^{pq})=B_s(E_2^{pq})$ for $s\geq r\geq r(n)$.
  \item[{\rm(c)}] There exists an integer $r(n)$ such that for $p+q=n$ or $p+q=n-1$, we have $Z_r(E_2^{pq})=Z_s(E_2^{pq})$ for $s\geq r\geq r(n)$.
  \item[{\rm(d)}] There exists an integer $r(n)$ such that for $p+q=n$, we have $B_r(E_2^{pq})=B_s(E_2^{pq})$ and $Z_r(E_2^{pq})=Z_s(E_2^{pq})$ for $s\geq r\geq r(n)$.
\end{enumerate}
\end{lemma}

\begin{proof}
\oldpage[0\textsubscript{III}]{27}
According to the conditions (a), (b), and (c) of \sref{0.11.1.1}, we have that saying $Z_{r+1}(E_2^{pq})=Z_r(E_2^{pq})$ is equivalent to saying that $d_r^{pq}=0$ and that saying $B_r(E_2^{p+r,q-r+1})=B_{r+1}(E_2^{p+r,q-r+1})$ is equivalent to saying that $d_r^{pq}=0$; the lemma immediately follows from this remark.
\end{proof}

\subsection{The spectral sequence of a filtered complex}
\label{subsection:0.11.2}

\begin{env}[11.2.1]
\label{0.11.2.1}
Given an abelian category $\cat{C}$, we will agree to denote by notation such as $K^\bullet$ the \emph{complexes $(K^i)_{i\in\bb{Z}}$} of objects of $\cat{C}$ whose differential is of degree $+1$, and by the notation such as $K_\bullet$ the complexes $(K_i)_{i\in\bb{Z}}$ of objects of $\cat{C}$ whose differential is of degree $-1$.
To each complex $K^\bullet=(K^i)$ whose differential $d$ is of degree $+1$, we can associate a complex $K_\bullet'=(K_i')$ by setting $K_i'=K^{-i}$, the differential $K_i'\to K_{i-1}'$ being the operator $d:K^{-i}\to K^{-i+1}$; and \emph{vice versa}, which, depending on the circumstances, will allow one to consider either one of the types of complexes and translate any result from one type into results for the other.
We similarly denote by notation such as $K^{\bullet\bullet}=(K^{ij})$ (resp. $K_{\bullet\bullet}=(K_{ij})$) the \emph{bicomplexes} (or \emph{double complexes}) of objects of $\cat{C}$ in which the \emph{two} differentials are of degree $+1$ (resp. $-1$); we can still pass from one type to the other by changing the signs of the indices, and we have similar notation and remarks for any multicomplexes.
The notation $K^\bullet$ and $K_\bullet$ will also be used for \emph{$\bb{Z}$-graded objects} of $\cat{C}$, which are not necessarily complexes (they can be considered as such for the \emph{zero} differentials); for example, we write $\HH^\bullet(K^\bullet)=(\HH^i(K^\bullet))_{i\in\bb{Z}}$ for the \emph{cohomology} of a complex $K^\bullet$ whose differential is of degree $+1$, and $\HH_\bullet(K_\bullet)=(\HH_i(K_\bullet))_{i\in\bb{Z}}$ for the \emph{homology} of a complex $K_\bullet$ whose differential is of degree $-1$; when we pass from $K^\bullet$ to $K_\bullet'$ by the method described above, we have $\HH_i(K_\bullet')=\HH^{-i}(K^\bullet)$.

Recall in this case that for a complex $K^\bullet$ (resp. $K_\bullet$), we will write in general $Z^i(K^\bullet)=\Ker(K^i\to K^{i+1})$ (``object of cocycles'') and $B^i(K^\bullet)=\Im(K^{i-1}\to K^i)$ (``object of coboundaries'') (resp. $Z_i(K_\bullet)=\Ker(K_i\to K_{i-1})$ (``object of cycles'') and $B_i(K_\bullet)=\Im(K_{i+1}\to K_i)$ (``object of boundaries'')) so that $\HH^i(K^\bullet)=Z^i(K^\bullet)/B^i(K^\bullet)$ (resp. $\HH_i(K_\bullet)=Z_i(K_\bullet)/B_i(K_\bullet)$).

If $K^\bullet=(K^i)$ (resp. $K_\bullet=(K_i)$) is a complex in $\cat{C}$ and $T:\cat{C}\to\cat{C}'$ a functor from $\cat{C}$ to an abelian category $\cat{C}'$, then we denote by $T(K^\bullet)$ (resp. $T(K_\bullet)$) the complex $(T(K^i))$ (resp. $(T(K_i))$) in $\cat{C}'$.

We will not review the definition of the \emph{$\partial$-functors} (T, 2.1), except to note that we \emph{also} say $\partial$-functor in place of $\partial^*$-functor when the morphism $\partial$ decreases the degree by one, the context clarifying this point if there is cause for confusion.

Finally, we say that a \emph{graded object $(A_i)_{i\in\bb{Z}}$} of $\cat{C}$ is \emph{bounded below} (resp. \emph{above}) if there exists an $i_0$ such that $A_i=0$ for $i<i_0$ (resp. $i>i_0$).
\end{env}

\begin{env}[11.2.2]
\label{0.11.2.2}
Let $K^\bullet$ be a complex in $\cat{C}$ whose differential $d$ is of degree $+1$, and suppose it is equipped with a \emph{filtration $F(K^\bullet)=(F^p(K^\bullet))_{p\in\bb{Z}}$} consisting of \emph{graded} subobjects
\oldpage[0\textsubscript{III}]{28}
of $K^\bullet$, in other words, $F^p(K^\bullet)=(K^i\cap F^p(K^\bullet))_{i\in\bb{Z}}$; in addition, we assume that $d(F^p(K^\bullet))\subset F^p(K^\bullet)$ for every $p\in\bb{Z}$.
Let us quickly recall how one \emph{functorially} defines a spectral sequence $E(K^\bullet)$ from $K^\bullet$ (M,~XV,~4 and G,~I,~4.3).
For $r\geq 2$, the canonical morphism $F^p(K^\bullet)/F^{p+r}(K^\bullet)\to F^p(K^\bullet)/F^{p+1}(K^\bullet)$ defines a morphism in cohomology
\[
  \HH^{p+q}(F^p(K^\bullet)/F^{p+r}(K^\bullet))\to\HH^{p+q}(F^p(K^\bullet)/F^{p+1}(K^\bullet)).
\]

We denote by $Z_r^{pq}(K^\bullet)$ the image of this morphism.
Similarly, from the exact sequence
\[
  0\to F^p(K^\bullet)/F^{p+1}(K^\bullet)\to F^{p-r+1}(K^\bullet)/F^{p+1}(K^\bullet)\to F^{p-r+1}(K^\bullet)/F^p(K^\bullet)\to 0,
\]
we deduce from the exact sequence in cohomology a morphism
\[
  \HH^{p+q-1}(F^{p-r+1}(K^\bullet)/F^p(K^\bullet))\to\HH^{p+q}(F^p(K^\bullet)/F^{p+1}(K^\bullet)),
\]
and we denote by $B_r^{pq}(K^\bullet)$ the image of this morphism; we show that $B_r^{pq}(K^\bullet)\subset Z_r^{pq}(K^\bullet)$ and we take $E_r^{pq}(K^\bullet)=Z_r^{pq}(K^\bullet)/B_r^{pq}(K^\bullet)$; we will not specify the definition of the $d_r^{pq}$ or the $\alpha_r^{pq}$.

We note here that all the $Z_r^{pq}(K^\bullet)$ and $B_r^{pq}(K^\bullet)$, for $p$ and $q$ fixed, are subobjects of the same object $\HH^{p+q}(F^p(K^\bullet)/F^{p+1}(K^\bullet))$, which we denote by $Z_1^{pq}(K^\bullet)$; we set $B_1^{pq}(K^\bullet)=0$, so that the above definitions of $Z_r^{pq}(K^\bullet)$ and $B_r^{pq}(K^\bullet)$ also work for $r=1$; we still set $E_1^{pq}(K^\bullet)=Z_1^{pq}(K^\bullet)$.
We define $d_1^{pq}$ and $\alpha_1^{pq}$ such that the conditions of \sref{0.11.1.1} are satisfied for $r=1$.
On the other hand, we define the subobjects $Z_\infty^{pq}(K^\bullet)$ as the image of the morphism
\[
  \HH^{p+q}(F^p(K^\bullet))\to\HH^{p+q}(F^p(K^\bullet)/F^{p+1}(K^\bullet))=E_1^{pq}(K^\bullet),
\]
and $B_\infty^{pq}(K^\bullet)$ as the image of the morphism
\[
  \HH^{p+q-1}(K^\bullet/F^p(K^\bullet))\to\HH^{p+q}(F^p(K^\bullet)/F^{p+1}(K^\bullet))=E_1^{pq}(K^\bullet),
\]
induced as above from the exact sequence in cohomology.
We set $Z_\infty(E_2^{pq}(K^\bullet))$ and $B_\infty(E_2^{pq}(K^\bullet))$ to be the canonical images in $E_2^{pq}(K^\bullet)$ of $Z_\infty^{pq}(K^\bullet)$ and $B_\infty^{pq}(K^\bullet)$.

Finally, we denote by $F^p(\HH^n(K^\bullet))$ the image in $\HH^n(K^\bullet)$ of the morphism $\HH^n(F^p(K^\bullet))\to\HH^n(K^\bullet)$ induced from the canonical injection $F^p(K^\bullet)\to K^\bullet$; by the exact sequence in cohomology, this is also the kernel of the morphism $\HH^n(K^\bullet)\to\HH^n(K^\bullet/F^p(K^\bullet))$.
This defines a filtration on $E^n(K^\bullet)=\HH^n(K^\bullet)$; we will not give here the definition of the isomorphisms $\beta^{pq}$.
\end{env}

\begin{env}[11.2.3]
\label{0.11.2.3}
The \emph{functorial} nature of $E(K^\bullet)$ is understood in the following way: given two \emph{filtered} complexes $K^\bullet$ and $K^{\prime\bullet}$ in $\cat{C}$ and a morphism of complexes $u:K^\bullet\to K^{\prime\bullet}$ that is \emph{compatible with the filtrations}, we induce in an evident way the morphisms $u_r^{pq}$ (for $r\geq 1$) and $u^n$, and we show that these morphisms are compatible with the $d_r^{pq}$, $\alpha_r^{pq}$, and $\beta^{pq}$ in the sense of \sref{0.11.1.2}, and thus define a well-defined morphism $E(u):E(K^\bullet)\to E(K^{\prime\bullet})$ of spectral sequences.
In addition, we show that if $u$ and $v$ are morphisms $K^\bullet\to K^{\prime\bullet}$ of the above type, \emph{homotopic of order $\leq k$}, then $u_r^{pq}=v_r^{pq}$ for $r>k$ and $u^n=v^n$ for all $n$ (M,~XV,~3.1).
\end{env}

\begin{env}[11.2.4]
\label{0.11.2.4}
\oldpage[0\textsubscript{III}]{29}
Suppose that filtered inductive limits in $\cat{C}$ are exact.
Then if the filtration $(F^p(K^\bullet))$ of $K^\bullet$ is \emph{exhaustive}, then so is the filtration $(F^p(\HH^n(K^\bullet)))$ for all $n$, since by hypothesis we have $K^\bullet=\varinjlim_{p\to-\infty}F^p(K^\bullet)$ and since the hypothesis on $\cat{C}$ implies that cohomology commutes with inductive limits.
In addition, for the same reason, we have $B_\infty(E_2^{pq}(K^\bullet))=\sup_k B_k(E_2^{pq}(K^\bullet))$.
We say that the filtration $(F^p(K^\bullet))$ of $K^\bullet$ is \emph{regular} if for every $n$ there exists an integer $u(n)$ such that $\HH^n(F^p(K^\bullet))=0$ for $p>u(n)$.
This is particularly the case when the filtration of $K^\bullet$ is \emph{discrete}.
When the filtration of $K^\bullet$ is regular and exhaustive, and filtered inductive limits are exact in $\cat{C}$, we have (M,~XV,~4) that the spectral sequence $E(K^\bullet)$ is \emph{regular}.
\end{env}

\subsection{The spectral sequences of a bicomplex}
\label{subsection:0.11.3}

\begin{env}[11.3.1]
\label{0.11.3.1}
With regard to the conventions for bicomplexes, we follow those of~(T,~2.4) rather than those of~(M), the two differentials $d'$ and $d''$ (of degree~$+1$) of such a bicomplex $K^{\bullet\bullet}=(K^{ij})$ thus being assumed to be \emph{permutable}.
Suppose that \emph{one} of the following two conditions is satisfied:
\begin{enumerate}
   \item \emph{Infinite direct sums} exist in $\cat{C}$;
   \item For all $n\in\bb{Z}$, there is only a \emph{finite} number of pairs $(p,q)$ such that $p+q=n$ and $K^{pq}\neq 0$.
 \end{enumerate}
Then the bicomplex $K^{\bullet\bullet}$ defines a (simple) \emph{complex} $(K^{\prime n})_{n\in\bb{Z}}$ with $K^{\prime n}=\sum_{i+j=n}K^{ij}$, the differential $d$ (of degree~$+1$) of this complex being given by $dx=d'x+(-1)^i d''x$ for $x\in K^{ij}$.
When we later speak of the (simple) \emph{complex defined by a bicomplex $K^{\bullet\bullet}$}, it will always be understood that one of the above conditions is satisfied.
We adopt the analogous conventions for multicomplexes.

We denote by $K^{i,\bullet}$ (resp.~$K^{\bullet,j}$) the simple complex $(K^{ij})_{j\in\bb{Z}}$ (resp.~$(K^{ij})_{i\in\bb{Z}}$), by $Z_\mathrm{II}^p(K^{i,\bullet})$, $B_\mathrm{II}^p(K^{i,\bullet})$, $\HH_\mathrm{II}^p(K^{i,\bullet})$ (resp.~$Z_\mathrm{I}^p(K^{\bullet,j})$, $B_\mathrm{I}^p(K^{\bullet,j})$, $\HH_\mathrm{I}^p(K^{\bullet,j})$) its $p$ objects of cocycles, of coboundaries, and of cohomology, respectively;
the differential $d':K^{i,\bullet}\to K^{i+1,\bullet}$ is a morphism of complexes, which thus gives an operator on the cocycles, coboundaries, and cohomology,
\begin{align*}
  d'&:Z_\mathrm{II}^p(K^{i,\bullet})\to Z_\mathrm{II}^p(K^{i+1,\bullet})
\\d'&:B_\mathrm{II}^p(K^{i,\bullet})\to B_\mathrm{II}^p(K^{i+1,\bullet})
\\d'&:\HH_\mathrm{II}^p(K^{i,\bullet})\to\HH_\mathrm{II}^p(K^{i+1,\bullet})
\end{align*}
and it is clear that for these operators, $(Z_\mathrm{II}^p(K^{i,\bullet}))_{i\in\bb{Z}}$, $(B_\mathrm{II}^p(K^{i,\bullet}))_{i\in\bb{Z}}$, and $(\HH_\mathrm{II}^p(K^{i,\bullet}))_{i\in\bb{Z}}$ are complexes;
we denote the complex $(\HH_\mathrm{II}^p(K^{i,\bullet}))_{i\in\bb{Z}}$ by $\HH_\mathrm{II}^p(K^{\bullet\bullet})$, and its $q$ objects of cocycles, coboundaries, and cohomology by $Z_\mathrm{I}^q(\HH_\mathrm{II}^p(K^{\bullet\bullet}))$, $B_\mathrm{I}^q(\HH_\mathrm{II}^p(K^{\bullet\bullet}))$, and $\HH_\mathrm{I}^q(\HH_\mathrm{II}^p(K^{\bullet\bullet}))$.
We similarly define the complexes $\HH_\mathrm{I}^p(K^{\bullet\bullet})$ and their cohomology objects $\HH_\mathrm{II}^q(\HH_\mathrm{I}^p(K^{\bullet\bullet}))$.
Recall, however, that $\HH^n(K^{\bullet\bullet})$ denotes the $n$-th object of the cohomology of the \emph{(simple)} complex defined by $K^{\bullet\bullet}$.
\end{env}

\begin{env}[11.3.2]
\label{0.11.3.2}
On the complex defined by a bicomplex $K^{\bullet\bullet}$, we can consider two canonical filtrations, $(F_\mathrm{I}^p(K^{\bullet\bullet}))$ and $(F_\mathrm{II}^p(K^{\bullet\bullet}))$, given by
\[
\label{0.11.3.2.1}
  F_\mathrm{I}^p(K^{\bullet\bullet})
  = \left(\sum_{i+j=n,i\geq p}K^{ij}\right)_{n\in\bb{Z}}
  \quad\text{and}\quad
  F_\mathrm{II}^p(K^{\bullet\bullet})
  = \left(\sum_{i+j=n,j\geq p}K^{ij}\right)_{n\in\bb{Z}}
  \tag{11.3.2.1}
\]
\oldpage[0\textsubscript{III}]{30}
which, by definition, are graded subobjects of the (simple) complex defined by $K^{\bullet\bullet}$, and thus make this complex a filtered complex;
moreover, it is clear that these filtrations are \emph{exhaustive} and \emph{separated}.

There corresponds to each of these filtrations a spectral sequence~\sref{0.11.2.2};
we denote by $'E(K^{\bullet\bullet})$ and $''E(K^{\bullet\bullet})$ the spectral sequences corresponding to $(F_\mathrm{I}^p(K^{\bullet\bullet}))$ and $(F_\mathrm{II}^p(K^{\bullet\bullet}))$ respectively, called the \emph{spectral sequence of the bicomplex $K^{\bullet\bullet}$}, and both having as their abutment the cohomology $(\HH^n(K^{\bullet\bullet}))$.
We can further show~(M,~XV,~6) that we have
\[
\label{0.11.3.2.2}
  'E_2^{pq}(K^{\bullet\bullet})
  = \HH_\mathrm{I}^p(\HH_\mathrm{II}^q(K^{\bullet\bullet}))
  \quad\text{and}\quad
  ''E_2^{pq}(K^{\bullet\bullet})
  = \HH_\mathrm{II}^q(\HH_\mathrm{I}^p(K^{\bullet\bullet})).
  \tag{11.3.2.2}
\]

Every morphism $u:K^{\bullet\bullet}\to K^{\prime\bullet\bullet}$ of bicomplexes is \emph{ipso facto} compatible with the filtrations of the same type of $K^{\bullet\bullet}$ and $K^{\prime\bullet\bullet}$, thus defines a morphism for each of the two spectral sequences;
in addition, two \emph{homotopic} morphisms define a homotopy \emph{of order $\leq 1$} of the corresponding filtered (simple) complexes, thus the \emph{same} morphism for each of the two spectral sequences~(M,~XV,~6.1).
\end{env}

\begin{proposition}[11.3.3]
\label{0.11.3.3}
Let $K^{\bullet\bullet}=(K^{ij})$ be a bicomplex in an abelian category $\cat{C}$.
\begin{enumerate}
  \item[{\rm(i)}] If there exist $i_0$ and $j_0$ such that $K^{ij}=0$ for $i<i_0$ or $j<j_0$ (resp.~$i>i_0$ or $j>j_0$), then the two spectral sequences $'E(K^{\bullet\bullet})$ and $''E(K^{\bullet\bullet})$ are biregular.
  \item[{\rm(ii)}] If there exist $i_0$ and $i_1$ such that $K^{ij}=0$ for $i<i_0$ or $i>i_1$ (resp.~if there exist $j_0$ and $j_1$ such that $K^{ij}=0$ for $j<j_0$ or $j>j_1$), then the two spectral sequences $'E(K^{\bullet\bullet})$ and $''E(K^{\bullet\bullet})$ are biregular.
\end{enumerate}
Suppose in addition that filtered inductive limits exist and are exact in $\cat{C}$. Then:
\begin{enumerate}
  \item[{\rm(iii)}] If there exists $i_0$ such that $K^{ij}=0$ for $i>i_0$ (resp.~if there exists $j_0$ such that $K^{ij}=0$ for $j<j_0$), then the spectral sequence $'E(K^{\bullet\bullet})$ is regular.
  \item[{\rm(iv)}] If there exists $i_0$ such that $K^{ij}=0$ for $i<i_0$ (resp.~if there exists $j_0$ such that $K^{ij}=0$ for $j>j_0$), then the spectral sequence $''E(K^{\bullet\bullet})$ is regular.
\end{enumerate}
\end{proposition}

\begin{proof}
The proposition follows immediately from the definitions~\sref{0.11.1.3} and from~\sref{0.11.2.4}, as well as from the following observations relating to the filtration $F_\mathrm{I}$ (and similar observations that we can deduce for $F_\mathrm{II}$ by exchanging the roles of the two indices in $K^{\bullet\bullet}$):
\begin{enumerate}
  \item[{$1^{\circ}$}] If there exists $i_0$ such that $K^{ij}=0$ for $i>i_0$, then the filtration $F_\mathrm{I}(K^{\bullet\bullet})$ is \emph{discrete}.
  \item[{$2^{\circ}$}] If there exists $i_0$ such that $K^{ij}=0$ for $i<i_0$, then the filtration $F_\mathrm{I}(K^{\bullet\bullet})$ is \emph{co-discrete}.
    We can then immediately deduce that it is the same for the corresponding filtration $F_\mathrm{I}(\HH^n(K^{\bullet\bullet}))$ for all $n$.
    Furthermore, the definition of the $B_{r}^{pq}$ corresponding to the filtration $F_\mathrm{I}(K^{\bullet\bullet})$ \sref{0.11.2.2} shows that for any pair $(p,q)$, the sequence $(B_{r}^{pq})_{r\geq 2}$ is stable.
  \item[{$3^{\circ}$}] If there exists $j_0$ such that $K^{ij}=0$ for $j<j_0$, then we have 
    \[
      F_\mathrm{I}^{p+r}(K^{\bullet\bullet})
      \cap (\sum_{i+j=n}K^{ij})
      = 0
    \]
    as soon as $p+r+j_{0}>n$, so $Z_{r}^{pq} = Z_\infty(E_2^{pq})$ for $r>q-j_{0}+1$.
    On the other hand, $\HH^{n}(F_\mathrm{I}^{p}(K^{\bullet\bullet})) = 0$ for $p>n-j_{0}+1$.
  \item[{$4^{\circ}$}] If there exists $j_0$ such that $K^{ij}=0$ for $j>j_0$, then we have 
    \[
      F_\mathrm{I}^{p-r+1}(K^{\bullet\bullet})
      \cap (\sum_{i+j=n}K^{ij})
      = \sum_{i+j=n}K^{ij}
    \]
    \oldpage[0\textsubscript{III}]{31}
    as soon as $p-r+1+j_{0}<n$, so $B_{r}^{pq} = B_\infty(E_2^{pq})$ for $r>j_{0}-q+1$.
    On the other hand, $\HH^{n}(F_\mathrm{I}^{p}(K^{\bullet\bullet})) = \HH^{n}(K^{\bullet\bullet})$ for $p+j_{0}<n-1$.
\end{enumerate}
\end{proof}

\begin{env}[11.3.4]
\label{0.11.3.4}
Suppose that the bicomplex $K^{\bullet\bullet}=(K^{ij})$ is such that $K^{ij}=0$ for $i<0$ or $j<0$.
We know that we can define, for all $p\in\bb{Z}$, a canonical ``edge-homomorphism''.
\[
\label{0.11.3.4.1}
  'E_2^{p0}(K^{\bullet\bullet})
  \rightarrow \HH^{p}(K^{\bullet\bullet})
\tag{11.3.4.1}
\]
(M,~XV,~6).
Recall that this is due to, on the one hand, the spectral sequence
$'E(K^{\bullet\bullet})$, since $Z_r^{p0}=Z_{\mathrm{I}}^{p}(Z_{\mathrm{II}}^{0}(K^{\bullet\bullet}))$ for $2\leq r \leq +\infty$, and on the other hand, that $\HH^{p}(F_\mathrm{I}^{p+1}(K^{\bullet\bullet})) = 0$, so that the isomorphism $\beta^{p0}:\ 'E_\infty^{p0} \simto \HH^{p}(F_\mathrm{I}^p)/\HH^{p}(F_\mathrm{I}^{p+1})$ gives a homomorphism $'E_\infty^{p0} \rightarrow \HH^{p}(F_\mathrm{I}^p(K^{\bullet\bullet})) \rightarrow \HH^{p}(K^{\bullet\bullet})$.
The equality of all the $Z_r^{p0}$ allows us to define canonical homomorphisms $'E_r^{p0} \rightarrow 'E_s^{p0}$ for $r \leq s$, and in particular a homomorphism $'E_2^{p0} \rightarrow 'E_\infty^{p0}$, whence by composition the edge-homomorphism $'E_2^{p0}(K^{\bullet\bullet}) \rightarrow \HH^{p}(K^{\bullet\bullet})$.
Furthermore, we can immediately verify that the edge-homomorphism thus defined sends the class modulo $B_2^{p0}$ of an element $z\in Z_\mathrm{II}^0(K^{\bullet\bullet})\subset K^{p0}$ such that $d'z=0$ to the class of $z$ modulo $B_\infty^{p0}$ in $'E_\infty^{p0}$, and then to the cohomology class of $z$ in $\HH^p(K^{\bullet\bullet})$.
We therefore finally see that the \emph{edge-homomorphism} \sref{0.11.3.4.1} \emph{comes from, by passing to cohomology, the canonical injection} $Z_\mathrm{II}^0(K^{\bullet\bullet}) \rightarrow K^{\bullet\bullet}$  (where $K^{\bullet\bullet}$ is considered as simple complex).
We naturally interpret the edge-homomorphism in the same way 
\[
\label{0.11.3.4.2}
  ''E_2^{p0}(K^{\bullet\bullet})
  \rightarrow \HH^{p}(K^{\bullet\bullet})
\tag{11.3.4.2}
\]
as coming from the canonical injection $Z_\mathrm{I}^0(K^{\bullet\bullet}) \rightarrow K^{\bullet\bullet}$.
\end{env}

\begin{env}[11.3.5]
\label{0.11.3.5}
Now let $K_{\bullet\bullet} = (K_{ij})$ be a bicomplex in $\cat{C}$ whose two differential operators are of degrees~$-1$.
We then write $K_{i,\bullet}$ (resp. $K_{\bullet,j}$) to mean the simple complex $(K_{ij})_{j\in\bb{Z}}$ (resp. $(K_{ij})_{i\in\bb{Z}}$), $\HH_{p}^\mathrm{II}(K_{i, \bullet})$ (resp. $\HH_{p}^\mathrm{I}(K_{\bullet, j})$) to mean the $p$-th homology object, $\HH_{p}^\mathrm{II}(K_{\bullet\bullet})$ (resp. $\HH_{p}^\mathrm{I}(K_{\bullet\bullet})$) to mean the complex $(\HH_{p}^\mathrm{II}(K_{i, \bullet}))_{i \in\bb{Z}}$ (resp. $(\HH_{p}^\mathrm{I}(K_{\bullet, j}))_{j \in\bb{Z}}$), and $\HH_q^\mathrm{I}(\HH_{p}^\mathrm{II}(K_{\bullet\bullet}))$ (resp. $\HH_q^\mathrm{II}(\HH_{p}^\mathrm{I}(K_{\bullet\bullet}))$) to mean the $q$-th homology object;
we use analogous notation for the objects of cycles and objects of boundaries;
finally, we will denote by $\HH_n(K_{\bullet\bullet})$ the $n$-th homology object (when it exists) of the simple complex (with a differential operator of degree~$-1$) defined by $K_{\bullet\bullet}$.
Let $K'^{\bullet\bullet}=(K'^{ij})$, with $K'^{ij}=K_{-i,-j}$, be the bicomplex with differential operators of degrees~$+1$ associated to $K_{\bullet\bullet}$.
By definition, the \emph{spectral sequences} of $K_{\bullet\bullet}$ are those of $K'^{\bullet\bullet}$, that we write as $'E(K_{\bullet\bullet})$ and $''E(K_{\bullet\bullet})$, where, however, we change the notation, by putting
\[
  'E_{pq}^{r}(K_{\bullet\bullet})
  = 'E_{r}^{-p,-q}(K'^{\bullet\bullet})
  \quad\text{and}\quad
  ''E_{pq}^{r}(K_{\bullet\bullet})
  = ''E_{r}^{-p,-q}(K'^{\bullet\bullet}), 
\]
for $2\leq r \leq \infty$.
With this notation, we have
\[
  'E_{pq}^{2}(K_{\bullet\bullet})
  = \HH_{p}^\mathrm{I}(\HH_{q}^\mathrm{II}(K_{\bullet\bullet}))
  \quad\text{and}\quad
  ''E_{pq}^{2}(K_{\bullet\bullet})
  = \HH_{q}^\mathrm{II}(\HH_{p}^\mathrm{I}(K_{\bullet\bullet})).
\]
To avoid sign errors, it will be preferable in general to return to the complex $K'^{\bullet\bullet}$ for the relations between these spectral sequences and their abutments.
We note, however, the criteria corresponding to \sref{0.11.3.3}:
\end{env}

\begin{env}[11.3.6]
\label{0.11.3.6}
The spectral sequences $'E(K_{\bullet\bullet})$ and $''E(K_{\bullet\bullet})$ are \emph{biregular} in the following cases:
\begin{enumerate}
  \item[(a)] There exist $i_0$ and $j_0$ such that $K_{ij}=0$ for $i>i_0$ or $j>j_0$ (resp. for $i<i_0$
    \oldpage[0\textsubscript{III}]{32}
    or $j<j_0$);
  \item[(b)] There exist $i_0$ and $i_1$ such that $K_{ij}=0$ for $i<i_0$ and $i>i_1$;
  \item[(c)] There exist $j_0$ and $j_1$ such that $K_{ij}=0$ for $j<j_0$ and $j>j_1$.
\end{enumerate}

The sequence $'E(K_{\bullet\bullet})$ is \emph{regular} if there exists $i_0$ such that $K_{ij}=0$ for $i<i_0$, or if there exists $j_0$ such that $K_{ij}=0$ for $j>j_0$.

The sequence $''E(K_{\bullet\bullet})$ is \emph{regular} if there exists $i_0$ such that $K_{ij}=0$ for $i>i_0$, or if there exists $j_0$ such that $K_{ij}=0$ for $j<j_0$.
\end{env}


\subsection{Hypercohomology of a functor with respect to a complex $K^{\bullet}$}
\label{subsection:0.11.4}

\begin{env}[11.4.1]
\label{0.11.4.1}
Let $\cat{C}$ be an abelian category.
Recall that we defined the \emph{right resolution} (or \emph{cohomological resolution}) of an object $A$ in $\cat{C}$ to be a complex of objects in $\cat{C}$, whose differential operator is of degree $+1$,
\[
  0 \rightarrow L^0 \rightarrow L^1 \rightarrow L^2 \rightarrow \ldots
\]
endowed with a morphism $\varepsilon: A \rightarrow L^0$, called the \emph{augmentation} of the resolution, which we can consider as a morphism of complexes
\[
  \xymatrix{
    0 \ar[r]
    & A \ar[r] \ar[d]_{\varepsilon}
    & 0 \ar[r] \ar[d]
    & 0 \ar[r] \ar[d]
    & \ldots
  \\0 \ar[r]
    & L^0 \ar[r]
    & L^1 \ar[r]
    & L^2 \ar[r]
    & \ldots
  }
\]
such that the sequence
\[
  0 \rightarrow A \xrightarrow{\varepsilon} L^0 \rightarrow L^1 \longrightarrow \ldots
\]
is exact.
Similarly, a \emph{left resolution} (or \emph{homological resolution}) of $A$ is a complex $0 \leftarrow L_0 \leftarrow L_1 \leftarrow L_2 \leftarrow \ldots$ of objects in $\cat{C}$, whose differential operator is of degree~$-1$, endowed with an \emph{augmentation} $L_0 \xrightarrow{\varepsilon} A$ such that the sequence
\[
  0 \leftarrow A \xleftarrow{\varepsilon} L_0 \leftarrow L_1 \leftarrow \ldots
\]
is exact.

When the right resolution $(L^i)_{i\geq 0}$ of an object $A$ is such that $L^i=0$ for $i\geq n+1$, we say that this resolution is \emph{of length $\leq n$}.
We define a left resolution of length $\leq n$ similarly.
A resolution which is of length $\leq n$ for some integer $n$ is said to be \emph{finite}.

A resolution of $A$ is called \emph{projective} (resp. \emph{injective}) if the objects of $\cat{C}$, apart from $A$, of which it is composed are \emph{projective} (resp. \emph{injective}). 
When $\cat{C}$ is the category of (left, say) modules over a ring, we say that a resolution of $A$ is \emph{flat} (resp. \emph{free}) when the modules, apart from $A$, of which it is composed are \emph{flat} (resp. \emph{free}).
\end{env}

\begin{env}[11.4.2]
\label{0.11.4.2}
Let $K^\bullet = (K^i)_{i\in\bb{Z}}$ be a complex of objects of $\cat{C}$, whose differential operator is of degree~$+1$.
We define the \emph{right Cartan--Eilenberg resolution} of $K^\bullet$ to be a pair consisting of a bicomplex $L^{\bullet\bullet} = (L^{ij})$ with differential operator of degree~$+1$ and such that $L^{ij}=0$ for $j<0$, and a morphism of simple complexes $\varepsilon: K^\bullet \rightarrow L^{\bullet, 0}$, such that the following conditions are fulfilled:
\oldpage[0\textsubscript{III}]{33}
\begin{enumerate}
  \item[(i)] For each index $i$, the sequences
    \[
      \begin{aligned}
        0 \rightarrow
        & K^i
        \xrightarrow{\varepsilon} L^{i0}
        \rightarrow L^{i1}
        \rightarrow \ldots
      \\0 \rightarrow
        & B^i(K^\bullet)
        \xrightarrow{\varepsilon} B_\mathrm{I}^i(L^{\bullet,0})
        \rightarrow B_\mathrm{I}^i(L^{\bullet,1})
        \rightarrow \ldots
      \\0 \rightarrow
        & Z^i(K^\bullet)
        \xrightarrow{\varepsilon} Z_\mathrm{I}^i(L^{\bullet,0})
        \rightarrow Z_\mathrm{I}^i(L^{\bullet,1})
        \rightarrow \ldots
      \\0 \rightarrow
        & \HH^i(K^\bullet)
        \xrightarrow{\varepsilon} \HH_\mathrm{I}^i(L^{\bullet,0})
        \rightarrow \HH_\mathrm{I}^i(L^{\bullet,1})
        \rightarrow \ldots
      \end{aligned}
    \]
    are exact, or, in other words, $(L^{i,\bullet}), (B_\mathrm{I}^i(L^{\bullet\bullet})), (Z_\mathrm{I}^i(L^{\bullet\bullet}))$ and $(\HH_\mathrm{I}^i(L^{\bullet\bullet}))$ are \emph{resolutions} of $K^i, B^i(K^\bullet), Z^i(K^\bullet)$ and $\HH^i(K^\bullet)$ (respectively).
  \item[(ii)] For each $j$, the simple complex $L^{\bullet, j}$ is \emph{split}, or, in other words, the exact sequences
    \[
      \label{0.11.4.2.1}
      0 \rightarrow B_\mathrm{I}^i(L^{\bullet,j}) \rightarrow Z_\mathrm{I}^i(L^{\bullet,j}) \rightarrow \HH_\mathrm{I}^i(L^{\bullet,j}) \rightarrow 0
      \tag{11.4.2.1}
    \]
    \[
      \label{0.11.4.2.2}
      0 \rightarrow Z_\mathrm{I}^i(L^{\bullet,j}) \rightarrow L^{ij} \rightarrow B_\mathrm{I}^{i+1}(L^{\bullet,j}) \rightarrow 0
      \tag{11.4.2.2}
    \]
    are \emph{split}.
\end{enumerate}

We can prove (M,~XVII,~1.2) that if every object of $\cat{C}$ is a subobject of an injective object, then every complex $K^\bullet$ of $\cat{C}$ admits an \emph{injective} Cartan--Eilenberg resolution, i.e. consisting of injective objects $L^{ij}$ (the condition (ii) above then implying that the $(B_\mathrm{I}^i(L^{\bullet,j})), (Z_\mathrm{I}^i(L^{\bullet,j}))$ and $(\HH_\mathrm{I}^i(L^{\bullet,j}))$ also consist of injective objects).
Furthermore, for any morphism $f: K^\bullet \rightarrow K'^\bullet$ of complexes of $\cat{C}$, 
any Cartan--Eilenberg resolution $L^{\bullet\bullet}$ of $K^\bullet$, and any \emph{injective} Cartan--Eilenberg resolution $L'^{\bullet\bullet}$ of $K'^\bullet$, there exists a morphism of bicomplexes $F: L^{\bullet\bullet} \rightarrow L'^{\bullet\bullet}$ that is compatible with $f$ and its augmentation, and if $f$ and $g$ are homotopic morphisms from $K^\bullet$ to $K'^\bullet$, then the corresponding morphisms from $L^{\bullet\bullet}$ to $L'^{\bullet\bullet}$ are also homotopic \emph{(loc. cit.)}.

When $K^{\bullet}$ is \emph{bounded below} (resp. \emph{bounded above}), we can take $L^{\bullet\bullet}$ to be such that $L^{ij}=0$ for $i<i_0$ (resp. $i>i_0$) if $K^i=0$ for $i<i_0$ (resp. $i>i_0$) (M,~XVII,~1.3).

Suppose, on the other hand, that there exists an integer $n$ such that every object of $\cat{C}$ admits a \emph{injective resolution of length $\leq n$};
then we can assume that we have $L^{ij}=0$ for $j>n$ (M,~XVII,~1.4).
\end{env}


\begin{env}[11.4.3]
\label{0.11.4.3}
Now let $T$ be an \emph{additive covariant functor} from $\cat{C}$ to an abelian category $\cat{C'}$.
Given a complex $K^\bullet$ of $\cat{C}$ and an \emph{injective} Cartan--Eilenberg resolution $L^{\bullet\bullet}$ of $K^{\bullet}$, suppose that the (simple) complex defined by the bicomplex $T(L^{\bullet\bullet})$ exists \sref{0.11.3.1};
then the two spectral sequences $'E(T(L^{\bullet\bullet}))$ and $''E(T(L^{\bullet\bullet}))$ of this bicomplex are called the \emph{hypercohomology spectral sequences} of $T$ with respect to $K^\bullet$;
by \sref{0.11.4.2} and \sref{0.11.3.2}, they depend only on the complex $K^\bullet$, not on the choice of the injective Cartan--Eilenberg resolution $L^{\bullet\bullet}$;
furthermore, they depend \emph{functorially} on $K^\bullet$.
They have the same abutment $\HH^\bullet(T(L^{\bullet\bullet}))$, also called the \emph{hypercohomology} of $T$ with respect to $K^\bullet$, and denoted by $\RR^\bullet T(K^\bullet)$.
We can show that the $E_2$ terms of the two spectral sequences above are given by 
\[
\label{0.11.4.3.1}
  'E_2^{pq}
  = \HH^p(\RR^q T(K^\bullet))
\tag{11.4.3.1}
\]
\[
\label{0.11.4.3.2}
  ''E_2^{pq}
  = \RR^p T(\HH^q(K^\bullet))
\tag{11.4.3.2}
\]
\oldpage[0\textsubscript{III}]{34}
where $\RR^p T$ denotes, as usual, the $p$-th \emph{derived functor} of $T$ for $p\in\bb{Z}$;
$\RR^q T(K^\bullet)$ denotes the complexes $(\RR^q T(K^i))_{i\in\bb{Z}}$.
\phantomsection\label{0.11.4.4}
Unless explicitly otherwise mentioned, we will henceforth assume that every object of $\cat{C}$ is a subobject of an injective object of $\cat{C}$, so that injective Cartan--Eilenberg resolutions exist for any complex of $\cat{C}$.
Since $L^{ij}=0$ for $j<0$, the criteria of \sref{0.11.3.3} show that the \emph{two} hypercohomology spectral sequences of $T$ with respect to $K^\bullet$ exist and are \emph{biregular} in each of the following two cases:
\begin{enumerate}
  \item $K^\bullet$ is bounded below;
  \item Every object of $\cat{C}$ admits an injective resolution of length at most $n$, for some integer $n$ independent of the object in question.
\end{enumerate}
Indeed, we can suppose, in the first case, by \sref{0.11.4.2}, that there exists $i_0$ such that $L^{ij}=0$ for $i<i_0$, and, in the second case, that there exists $j_1$ such that $L^{ij}=0$ for $j>j_1$;
in each of these two cases, it is furthermore clear that, for any given $n$, there exist only a finite number of pairs $(i,j)$ such that $L^{ij}\neq 0$ and $i+j=n$, which proves our claims.

If we assume that filtrant inductive limits exist in $\cat{C}'$ and are exact (which implies, in particular, the existence of infinite direct sums in $\cat{C}'$), then the complex defined by the bicomplex $T(L^{\bullet\bullet})$ exists, and criterion \sref{0.11.3.3} shows that the sequence $'E(T(L^{\bullet\bullet}))$ is always \emph{regular}.

If $K^\bullet$ is a complex such that all the $K^i$ are zero except for a single $K^{i_0}$, then
$\RR^n T(K^\bullet)$ is isomorphic to $\RR^{n-i_0} T(K^{i_0})$, as follows immediately from the definitions by taking a Cartan--Eilenberg resolution $L^{\bullet\bullet}$ such that $L^{ij}=0$ for $i\neq i_0$.

If $K^\bullet$ and $K'^\bullet$ are two complexes of $\cat{C}$, and $f$ and $g$ are homotopic morphisms from $K^\bullet$ to $K'^\bullet$, then the morphisms $\RR^\bullet T(K^\bullet) \rightarrow \RR^\bullet T(K'^\bullet)$ induced by $f$ and $g$ are identical, and the same is true for the morphisms of the cohomology spectral sequences.
\end{env}

\begin{proposition}[11.4.5]
\label{0.11.4.5}
Suppose that filtrant inductive limits exist in $\cat{C}'$ and are exact.
If $\RR^n T(K^i) = 0$ for all $n>0$ and all $i\in\bb{Z}$, then we have functorial isomorphisms
\[
\label{0.11.4.5.1}
  \RR^i T(K^\bullet) \simto \HH^i(T(K^\bullet))
\tag{11.4.5.1}
\]
for all $i\in\bb{Z}$.
\end{proposition}

\begin{proof}
Indeed, the only non-zero $E_2$ terms of the first spectral sequence \sref{0.11.4.3.1} are then $'E_2^{p0} = \HH^p(T(K^\bullet))$;
in other words, this sequence is \emph{degenerate};
since it is regular \sref{0.11.4.4}, the conclusion follows from \sref{0.11.1.6}.
\end{proof}

\begin{env}[11.4.6]
\label{0.11.4.6}
Now consider, for example, a covariant bifunctor $(M, N) \rightarrow T(M, N)$ from $\cat{C} \times \cat{C}'$ to $\cat{C}''$, where $\cat{C}$, $\cat{C}'$, and $\cat{C}''$ are three abelian categories;
we assume, for simplicity, that $T$ is additive in each of its arguments, and furthermore that every object of $\cat{C}$ and every object of $\cat{C}'$ are subobjects of an injective object, 
and that filtered inductive limits exist in $\cat{C}''$ and are exact.
We then define the \emph{hypercohomology} of $T$ with respect to the complexes $K^\bullet$ and $K'^\bullet$ of $\cat{C}$ and $\cat{C}'$ (respectively), with differential operators of degree~$+1$, by replacing $K^\bullet$ (resp. $K'^\bullet$) by an injective Cartan--Eilenberg resolution $L^{\bullet\bullet}$ (resp. $L'^{\bullet\bullet}$);
then $T(L^{\bullet\bullet}, L'^{\bullet\bullet})$ is a quadricomplex of $\cat{C}''$, which we consider as a \emph{bicomplex} of $\cat{C}''$, with the degree of $T(L^{ij}, L'^{hk})$ being the integers $i+h$ and $j+k$.
The \emph{hypercohomology} of $T$ with respect to $K^\bullet$ and $K'^\bullet$ is by definition the cohomology $\HH^\bullet(T(L^{\bullet\bullet}, L'^{\bullet\bullet}))$ of this bicomplex  (in other words, that of the associated simple complex) and is denoted by $\RR^\bullet T(K^\bullet, K'^\bullet)$;
\oldpage[0\textsubscript{III}]{35}
it is the abutment of two spectral sequences whose $E_2$ terms are given by
\[
\label{0.11.4.6.1}
  'E_2^{pq}
  = \HH^p(\RR^q T(K^\bullet, K'^\bullet))
\tag{11.4.6.1}
\]
\[
\label{0.11.4.6.2}
  ''E_2^{pq}
  = \sum_{q'+q''=q}\RR^p T(\HH^{q'}(K^\bullet),\HH^{q''}(K'^\bullet))
\tag{11.4.6.2}
\]
(cf.~M,~XVII,~2).

Here $\RR^q T(K^\bullet, K'^\bullet)$ is the bicomplex $(\RR^q T(K^i, K'^j))_{(i,j)\in \bb{Z}\times \bb{Z}}$, and the right-hand side of \sref{0.11.4.6.1} is its cohomology when we consider it as a simple complex.

Furthermore, the first spectral sequence is always regular, and the two spectral sequences are biregular when there exists $n$ such that every object of $\cat{C}$ and every object of $\cat{C}'$ admit an injective resolution of length $\leq n$, or when $K^\bullet$ and $K'^\bullet$ are bounded below;
in the latter case, we can furthermore omit the hypothesis that inductive limits exist in $\cat{C}''$.

If $K_1^\bullet$ and ${K'_1}^\bullet$ are two other complexes of $\cat{C}$ and $\cat{C}'$ (respectively), $f$ and $g$ \emph{homotopic} morphisms from $K^\bullet$ to $K_1^\bullet$, 
and $f'$ and $g'$ \emph{homotopic} morphisms from $K'^\bullet$ to ${K'_1}^\bullet$, then the morphisms $\RR^\bullet  T(K^\bullet, K'^\bullet) \rightarrow \RR^\bullet  T(K_1^\bullet, {K'_1}^\bullet)$ induced by $f$ and $f'$ on the one hand, and by $g$ and $g'$ on the other hand, are identical, and the same is true for the morphisms of the hypercohomology spectral sequences.

We can generalize easily to any additive covariant multifunctor.
\end{env}

\begin{proposition}[11.4.7]
\label{0.11.4.7}
Suppose that for any injective object $I$ of $\cat{C}$ (resp. $I'$ of $\cat{C}'$), $A'\mapsto T(I, A')$ (resp. $A\mapsto T(A, I')$) is an exact functor.
Then, with the notation of \sref{0.11.4.6}, we have canonical isomorphisms
\[
  \label{0.11.4.7.1} \RR^\bullet T(K^\bullet, K'^\bullet) \simto \HH^\bullet(T(L^{\bullet\bullet}, K'^\bullet)) \simto \HH^\bullet(T(K^\bullet, L'^{\bullet\bullet})) \tag{11.4.7.1} 
\]
where the last two terms are the cohomology of simple complexes defined by the tricomplexes $T(L^{\bullet\bullet}, K'^\bullet)$ and $T(K^\bullet, L'^{\bullet\bullet})$ (respectively).
\end{proposition}

\begin{proof}
Let us define, for example, the first of these isomorphisms.
The quadricomplex $T(L^{\bullet\bullet}, L'^{\bullet\bullet})$ can be considered as a bicomplex, where the degrees of $T(L^{ij}, L'^{hk})$ are the integers $i+j+h$ and $k$.
Since, for each $h$, ${L'}^{h,\bullet}$ is a \emph{resolution} of ${K'}^h$, we have, for this bicomplex, by virtue of the hypotheses on $T$, $\HH_\mathrm{II}^q(T(L^{\bullet\bullet}, L'^{\bullet\bullet}))=0$ for $q\neq 0$, and $\HH_\mathrm{II}^0(T(L^{\bullet\bullet}, L'^{\bullet\bullet})) = T(L^{\bullet\bullet}, K'^\bullet)$;
the first spectral sequence of this bicomplex is thus \emph{degenerate};
since $L'^{hk}=0$ for $k<0$, this sequence is also regular \sref{0.11.3.3}, and the conclusion then follows from \sref{0.11.1.6}.
\end{proof}

We have similar results for a covariant multifunctor of any number $n$ of arguments: in the calculation of the hypercohomology, it is not necessary to replace \emph{all} the complexes 
by a Cartan--Eilenberg resolution, but only $n-1$ of them, provided that, when we fix $n-1$ arbitrary arguments by taking them to be \emph{injective} objects, the covariant functor in the remaining argument is \emph{exact}.


\subsection{Passage to the inductive limit in the hypercohomology}
\label{subsection:0.11.5}

\begin{lemma}[11.5.1]
\label{0.11.5.1}
Let $K^\bullet=(K^i)_{i\in\bb{Z}}$ be a complex of $\cat{C}$, and, for every integer $r\in\bb{Z}$, let $K^\bullet_{(r)}$ be the complex such that $K^i_{(r)}=0$ for $i<r$ and $K^i_{(r)}=K^i$ for $i\geq r$.
Let $T$ be an additive covariant functor
\oldpage[0\textsubscript{III}]{36}
from $\cat{C}$ to $\cat{C}'$ that commutes with inductive limits; we suppose that filtered inductive limits exist and are exact in $\cat{C}$ and $\cat{C}'$.
Then $\RR^\bullet T(K^\bullet)$ is isomorphic to the inductive limit
\[
  \varinjlim_{r\to-\infty}\RR^\bullet T(K^\bullet_{(r)}).
\]
\end{lemma}

An injective Cartan--Eilenberg resolution of $K^\bullet$ is constructed by choosing arbitrarily, for every $i$, an injective resolution $(X_{\mathrm{B}}^{ij})_{j\geq 0}$ of $\mathrm{B}^i(K^\bullet)$ and an injective resolution $(X_{\mathrm{H}}^{ij})_{j\geq 0}$ of $\HH^i(K^\bullet)$.
Once this is done, the construction shows that the injective resolution $(L^{ij})_{j\geq 0}$ of $K^i$ and the differential operators $L^{i,j}\to L^{i+1,j}$ depend only on the resolutions $(X_{\mathrm{B}}^{i,\bullet})$, $(X_{\mathrm{H}}^{i,\bullet})$, and $(X_{\mathrm{B}}^{i+1,\bullet})$ (M,~XVII,~1.2).

Now it is clear that $\mathrm{B}^i(K^\bullet_{(r)})=\mathrm{B}^i(K^\bullet)$ and $\HH^i(K^\bullet_{(r)})=\HH^i(K^\bullet)$ for $i\geq r+1$.
On the other hand, for every $r$ there is a canonical injection
\[
  \varphi^\bullet_{r-1,r}:K^\bullet_{(r)}\longrightarrow K^\bullet_{(r-1)},
\]
where $\varphi^i_{r-1,r}$ is the identity for $i\neq r-1$.
The preceding observation shows that, if $L^{\bullet\bullet}=(L^{ij})$ is an injective Cartan--Eilenberg resolution of $K^\bullet$, then for every $r$ one can define an injective Cartan--Eilenberg resolution $L^{\bullet\bullet}_{(r)}=(L^{ij}_{(r)})$ of $K^\bullet_{(r)}$ such that $L^{ij}_{(r)}=0$ for $i<r$ and $L^{ij}_{(r)}=L^{ij}$ for $i\geq r+1$.
One can also define a morphism of bicomplexes
\[
  \Phi^{\bullet\bullet}_{r-1,r}:L^{\bullet\bullet}_{(r)}
  \longrightarrow L^{\bullet\bullet}_{(r-1)}
\]
corresponding to $\varphi^\bullet_{r-1,r}$; the construction of this morphism (\emph{loc.\ cit.}) again shows that it can be chosen so that $\Phi^{ij}_{r-1,r}$ is the identity for $i\neq r$ and $i\neq r-1$.

We have thus defined an inductive system $(L^{\bullet\bullet}_{(r)})$ of bicomplexes of $\cat{C}$ whose inductive limit, as $r$ tends to $-\infty$, is evidently $L^{\bullet\bullet}$.
Since direct sums commute with inductive limits, the simple complex associated with $L^{\bullet\bullet}$ is likewise the inductive limit of the simple complexes associated with the $L^{\bullet\bullet}_{(r)}$.
By hypothesis $T$ commutes with inductive limits, and so does cohomology, because the functor $\varinjlim$ is exact.
Consequently, up to isomorphism,
\[
  \HH^\bullet(T(L^{\bullet\bullet}))
  =\varinjlim_{r\to-\infty}\HH^\bullet(T(L^{\bullet\bullet}_{(r)})).
\]

Lemma~\sref{0.11.5.1} allows one, by passage to the inductive limit, to extend to arbitrary complexes $K^\bullet$ properties that hold for complexes \emph{bounded below}.
As a first example, we prove the following.

\begin{proposition}[11.5.2]
\label{0.11.5.2}
Under the hypotheses of \sref{0.11.5.1} concerning $\cat{C}$, $\cat{C}'$, and $T$, the functor $\RR^\bullet T(K^\bullet)$ is a cohomological functor on the abelian category of complexes of $\cat{C}$.
\end{proposition}

\begin{proof}
We first show that it is enough to treat complexes bounded below.
Given an exact sequence
\[
  0\longrightarrow K'^\bullet\longrightarrow K^\bullet
  \longrightarrow K''^\bullet\longrightarrow 0
\]
of complexes, for every $r$ it induces an exact sequence
\[
  0\longrightarrow K'^\bullet_{(r)}\longrightarrow K^\bullet_{(r)}
  \longrightarrow K''^\bullet_{(r)}\longrightarrow 0.
\]
By hypothesis, the latter gives an exact sequence
\[
  \cdots\longrightarrow\RR^nT(K'^\bullet_{(r)})
  \longrightarrow\RR^nT(K^\bullet_{(r)})
  \longrightarrow\RR^nT(K''^\bullet_{(r)})
  \xrightarrow{\partial}\RR^{n+1}T(K'^\bullet_{(r)})
  \longrightarrow\cdots .
\]
These exact sequences form an inductive system.
Lemma~\sref{0.11.5.1} and the exactness of $\varinjlim$ therefore give an exact sequence
\[
  \cdots\longrightarrow\RR^nT(K'^\bullet)
  \longrightarrow\RR^nT(K^\bullet)
  \longrightarrow\RR^nT(K''^\bullet)
  \xrightarrow{\partial}\RR^{n+1}T(K'^\bullet)
  \longrightarrow\cdots .
\]
\end{proof}

To treat complexes bounded below, it is enough to consider those for which $K^i=0$ when $i<0$; these plainly form an abelian category, which we denote by $\cat{K}$.

\begin{lemma}[11.5.2.1]
\label{0.11.5.2.1}
In $\cat{K}$, let $\mathfrak{J}$ be the class of complexes $Q^\bullet=(Q^i)_{i\geq0}$ having the
\oldpage[0\textsubscript{III}]{37}
following properties:
\begin{enumerate}
  \item[{\rm(1)}] every $Q^i$ is an injective object of $\cat{C}$;
  \item[{\rm(2)}] for every $i\geq0$, $\mathrm{Z}^i(Q^\bullet)=\mathrm{B}^i(Q^\bullet)$, and $\mathrm{Z}^i(Q^\bullet)$ is a direct summand of $Q^i$.
\end{enumerate}
Then:
\begin{enumerate}
  \item[{\rm(i)}] every $Q^\bullet\in\mathfrak{J}$ is an injective object of $\cat{K}$;
  \item[{\rm(ii)}] every object of $\cat{K}$ is isomorphic to a subcomplex of a complex belonging to $\mathfrak{J}$.
\end{enumerate}
\end{lemma}

\begin{proof}
{\rm(i)}
Let $A^\bullet=(A^i)$ be an object of $\cat{K}$, let $A'^\bullet=(A'^i)$ be a subobject of $A^\bullet$, let $Q^\bullet=(Q^i)$ belong to $\mathfrak{J}$, and suppose that a morphism
\[
  f=(f^i):A'^\bullet\longrightarrow Q^\bullet
\]
is given.
We must extend it to a morphism $g=(g^i):A^\bullet\to Q^\bullet$.
For simplicity, we use the language of the category of modules (cf.~[27]).

Identify $Q^i$ with $\mathrm{B}^i(Q^\bullet)\oplus\mathrm{B}^{i+1}(Q^\bullet)$.
We proceed by induction on $i$.
Suppose that $g^j$ has been defined for $j<i$, compatibly with the differentials $d^j:A^j\to A^{j+1}$ and $d'^j:Q^j\to Q^{j+1}$ for $j<i-1$, and that in addition:
\begin{enumerate}
  \item[{\rm(1)}] $g^{i-1}(\mathrm{Z}^{i-1}(A^\bullet))\subset\mathrm{Z}^{i-1}(Q^\bullet)$;
  \item[{\rm(2)}] if $C^j=(d^j)^{-1}(A'^{j+1})$ for every $j$, then $d'^{i-1}\circ g^{i-1}$ coincides with $f^i\circ d^{i-1}$ on $C^{i-1}$.
\end{enumerate}
By composing $f^i:A'^i\to Q^i$ with the two projections, we obtain morphisms
\[
  f'^i:A'^i\longrightarrow\mathrm{B}^i(Q^\bullet),
  \qquad
  f''^i:A'^i\longrightarrow\mathrm{B}^{i+1}(Q^\bullet).
\]
Since $d'^{i-1}\circ g^{i-1}$ maps $A^{i-1}$ into $\mathrm{B}^i(Q^\bullet)$ and vanishes on $\mathrm{Z}^{i-1}(A^\bullet)$, it defines a morphism
\[
  h^i:\mathrm{B}^i(A^\bullet)\longrightarrow\mathrm{B}^i(Q^\bullet).
\]
Condition~{\rm(2)} says that $h^i$ coincides with $f'^i$ on $\mathrm{B}^i(A^\bullet)\cap A'^i$.
Because $\mathrm{B}^i(Q^\bullet)$ is a direct summand of $Q^i$, it is injective; hence there is a morphism
\[
  g'^i:A^i\longrightarrow\mathrm{B}^i(Q^\bullet)
\]
that coincides with $h^i$ on $\mathrm{B}^i(A^\bullet)$ and with $f'^i$ on $A'^i$.

On the other hand, consider the morphism
\[
  f'^{i+1}\circ d^i:C^i\longrightarrow\mathrm{B}^{i+1}(Q^\bullet),
\]
which vanishes on $\mathrm{Z}^i(A^\bullet)$.
Since $\mathrm{B}^{i+1}(Q^\bullet)$ is injective, there is a morphism
\[
  g''^i:A^i\longrightarrow\mathrm{B}^{i+1}(Q^\bullet)
\]
that coincides with $f'^{i+1}\circ d^i$ on $C^i$ and with zero on $\mathrm{Z}^i(A^\bullet)$.
Taking $g^i=g'^i+g''^i$ allows the induction to continue.

\smallskip
\noindent{\rm(ii)}
To embed $A^\bullet=(A^i)$ in a complex belonging to $\mathfrak{J}$, choose for every $i\geq1$ an injective object $Q'^i$ of $\cat{C}$ and an injection $f'^i:A^i\to Q'^i$.
Set
\[
  Q^i=0\quad(i<0),\qquad Q^0=Q'^1,\qquad
  Q^i=Q'^i\oplus Q'^{i+1}\quad(i\geq1),
\]
with the evident differential.
For every $i$, set
\[
  f^i=f'^i+(f'^{i+1}\circ d^i),
\]
where $f'^i=0$ for $i\leq0$.
It is immediate that every $f^i$ is injective and that the $f^i$ commute with the differentials.
\end{proof}

\begin{corollary}[11.5.2.2]
\label{0.11.5.2.2}
Every object $K^\bullet$ of $\cat{K}$ admits a right resolution formed of objects of $\mathfrak{J}$.
If $L^{\bullet\bullet}$ is such a resolution, then for every resolution $L'^{\bullet\bullet}$ of $K^\bullet$ formed of objects of $\cat{K}$, there is a morphism of bicomplexes
\[
  F:L'^{\bullet\bullet}\longrightarrow L^{\bullet\bullet}
\]
compatible with the augmentations, and any two such morphisms $F,F'$ are homotopic.
\end{corollary}

\begin{proof}
This is precisely (M, V, 1.1~a)) applied to the abelian category $\cat{K}$.
\end{proof}

\begin{env}[11.5.2.3]
\label{0.11.5.2.3}
Having established these preliminaries, let $L'^{\bullet\bullet}$ be an injective Cartan--Eilenberg resolution of $K^\bullet$ and let $L^{\bullet\bullet}$ be a resolution of $K^\bullet$ formed of objects of $\mathfrak{J}$.
We show that there is an isomorphism
\[
  \HH^\bullet(T(L'^{\bullet\bullet}))
  \simto
  \HH^\bullet(T(L^{\bullet\bullet})).
\]
Corollary~\sref{0.11.5.2.2} gives a morphism of bicomplexes
\[
  T(L'^{\bullet\bullet})\longrightarrow T(L^{\bullet\bullet}),
\]
and hence a morphism
\[
  'E(T(L'^{\bullet\bullet}))\longrightarrow{}'E(T(L^{\bullet\bullet}))
\]
between their first spectral sequences.
By \sref{0.11.3.3}, these spectral sequences are regular, so by \sref{0.11.1.5} it is enough to check that this morphism is an isomorphism on the $E_2$ terms.
Equivalently, one must show that
\[
  \HH_\mathrm{II}^q(T(L^{i,\bullet}))=\RR^qT(K^i).
\]
Since $L^{i,\bullet}$ is a right resolution of $K^i$, this reduces the question to the following lemma.
\end{env}

\oldpage[0\textsubscript{III}]{38}
\begin{lemma}[11.5.2.4]
\label{0.11.5.2.4}
Let $L^\bullet=(L^i)_{i\in\bb{Z}}$ be a right resolution of an object $A$ of $\cat{C}$ such that $\RR^nT(L^i)=0$ for every $i\in\bb{Z}$ and every $n>0$.
Then
\[
  \HH^\bullet(T(L^\bullet))=\RR^\bullet T(A).
\]
\end{lemma}

\begin{proof}
This is a special case of (T, 2.5.2).
\end{proof}

\begin{env}[11.5.2.5]
\label{0.11.5.2.5}
The proof of \sref{0.11.5.2} now follows immediately, since $L'^{\bullet\bullet}$ is an injective resolution of $K^\bullet$ in the abelian category $\cat{K}$.
In other words, the functor
\[
  K^\bullet\longmapsto\HH^\bullet(T(L'^{\bullet\bullet}))
\]
is precisely the right-derived cohomological functor of $T$ in the category $\cat{K}$ (T, 2.3).
\end{env}

\begin{proposition}[11.5.3]
\label{0.11.5.3}
Under the hypotheses of \sref{0.11.5.1} concerning $\cat{C}$, $\cat{C}'$, and $T$, let $L^{\bullet\bullet}=(L^{ij})$ be a bicomplex such that $L^{ij}=0$ for $j<0$ and such that, for every $i$, $L^{i,\bullet}$ is a resolution of $K^i$.
Suppose further that $\RR^nT(L^{ij})=0$ for every pair $(i,j)$ and every $n>0$.
Then there is a functorial isomorphism
\[
\label{0.11.5.3.1}
\tag{11.5.3.1}
  \RR^\bullet T(K^\bullet)
  \simto
  \HH^\bullet(T(L^{\bullet\bullet})).
\]
\end{proposition}

\begin{proof}
Let $L^{\bullet\bullet}_{(r)}=(L^{ij}_{(r)})$ be the bicomplex defined by
\[
  L^{ij}_{(r)}=0\quad(i<r),\qquad
  L^{ij}_{(r)}=L^{ij}\quad(i\geq r).
\]
It is immediate that $L^{\bullet\bullet}$ is the inductive limit of the $L^{\bullet\bullet}_{(r)}$ as $r$ tends to $-\infty$.
By the hypothesis on $T$ and by \sref{0.11.5.1}, it is therefore enough to prove the proposition when $K^\bullet$ is bounded below, say when $K^i=0$ for $i<0$, with $L^{ij}=0$ for $i<0$.
Let $L'^{\bullet\bullet}=(L'^{ij})$ be a right resolution of $K^\bullet$ formed of objects of $\mathfrak{J}$, as in \sref{0.11.5.2.2}.
There is a morphism of bicomplexes
\[
  L^{\bullet\bullet}\longrightarrow L'^{\bullet\bullet}
\]
compatible with the augmentations, hence a morphism
\[
  'E(T(L^{\bullet\bullet}))\longrightarrow{}'E(T(L'^{\bullet\bullet}))
\]
between the first spectral sequences.
Lemma~\sref{0.11.5.2.4} shows, exactly as in \sref{0.11.5.2.3}, that this morphism is an isomorphism, proving the assertion.
\end{proof}

\begin{remark}[11.5.4]
\label{0.11.5.4}
The preceding arguments show that the conclusions of \sref{0.11.5.2} and \sref{0.11.5.3} hold in the category of complexes $K^\bullet$ bounded below without assuming that $T$ commutes with filtered inductive limits.
Moreover, if one considers only the category $\cat{K}$ of complexes $K^\bullet$ such that $K^i=0$ for $i<0$, the characterization of $\RR^\bullet T(K^\bullet)$ as the system of right-derived functors of $T$ in $\cat{K}$ shows that this cohomological functor is \emph{universal} (T, 2.3).

Another case in which \sref{0.11.5.2} holds without any additional hypothesis on $T$ is when there is an integer $m>0$ such that every object of $\cat{C}$ admits an injective resolution of length at most $m$.
Indeed, in the proof of \sref{0.11.5.1}, all the injective resolutions of objects of $\cat{C}$ that occur may be chosen of length at most $m$.
It follows at once that, for a fixed total degree $n$ and for $-r$ sufficiently large, the terms of total degree $n$ in $T(L^{\bullet\bullet}_{(r)})$ are equal to those in $T(L^{\bullet\bullet})$ and are finite in number.
Consequently,
\[
  \HH^n(T(L^{\bullet\bullet}))
  =\HH^n(T(L^{\bullet\bullet}_{(r)}))
\]
for $-r$ sufficiently large.
With the notation of \sref{0.11.5.2}, one likewise has
\[
  \RR^nT(K^\bullet_{(r)})=\RR^nT(K^\bullet)
\]
for $-r$ sufficiently large (depending on $n$), and similarly for $K'^\bullet$ and $K''^\bullet$, which proves the conclusion.
In the same way, \sref{0.11.5.3} holds without an additional condition on $T$ when $\cat{C}$ satisfies the preceding hypothesis and the resolutions $L^{i,\bullet}$ are taken to have length at most $m$.
\end{remark}

\begin{env}[11.5.5]
\label{0.11.5.5}
The result of \sref{0.11.5.2} extends to covariant multifunctors.
Consider, for example, the situation of \sref{0.11.4.6}, in which filtered inductive limits are assumed to exist and to be exact in $\cat{C}$, $\cat{C}'$, and $\cat{C}''$,
\oldpage[0\textsubscript{III}]{39}
and $T$ is assumed to commute with inductive limits.
Then $\RR^\bullet T(K^\bullet,K'^\bullet)$ is a bifunctor that is cohomological in each of the complexes $K^\bullet,K'^\bullet$.
To see this, one reduces, as in \sref{0.11.5.2}, to the case in which $K^\bullet$ and $K'^\bullet$ are bounded below.
Choosing injective resolutions of $K^\bullet$ and $K'^\bullet$ of the type described in \sref{0.11.5.2.2}, one is reduced to the general property proved in (M, V, 4.1).
\end{env}

\begin{env}[11.5.6]
\label{0.11.5.6}
Similarly, the results of \sref{0.11.4.7} and \sref{0.11.5.3} extend as follows.
Under the hypotheses of \sref{0.11.5.5}, suppose that we are given two bicomplexes
\[
  L^{\bullet\bullet}=(L^{ij}),\qquad
  L'^{\bullet\bullet}=(L'^{ij})
\]
such that $L^{ij}=L'^{ij}=0$ for $j<0$, such that, for every $i$, $L^{i,\bullet}$ is a resolution of $K^i$ and $L'^{i,\bullet}$ is a resolution of $K'^i$, and such that
\[
  \RR^nT(L^{ij},L'^{hk})=0
\]
for $n>0$ and every system of indices $(i,j,h,k)$.
Then there is an isomorphism, functorial in $K^\bullet$ and $K'^\bullet$,
\[
\label{0.11.5.6.1}
\tag{11.5.6.1}
  \RR^\bullet T(K^\bullet,K'^\bullet)
  \simto
  \HH^\bullet(T(L^{\bullet\bullet},L'^{\bullet\bullet})).
\]
This is proved as in \sref{0.11.5.3}, by reducing to the case in which $K^\bullet$ and $K'^\bullet$ are bounded below.

Suppose in addition that, for every pair $(i,j)$ and every pair $(h,k)$, the functors
\[
  A\longmapsto T(A,L'^{hk}),
  \qquad
  A'\longmapsto T(L^{ij},A')
\]
are exact in $\cat{C}$ and $\cat{C}'$, respectively.
Then there are also functorial isomorphisms
\[
\label{0.11.5.6.2}
\tag{11.5.6.2}
  \RR^\bullet T(K^\bullet,K'^\bullet)
  \simto
  \HH^\bullet(T(L^{\bullet\bullet},K'^\bullet))
  \simto
  \HH^\bullet(T(K^\bullet,L'^{\bullet\bullet})).
\]
The proof is analogous to that of \sref{0.11.4.7}.
\end{env}

\begin{env}[11.5.7]
\label{0.11.5.7}
The results of \sref{0.11.5.5} and \sref{0.11.5.6} also hold without assuming that $T$ commutes with inductive limits, provided that one restricts to complexes $K^\bullet,K'^\bullet$ bounded below, or that one assumes that every object of $\cat{C}$ (respectively, of $\cat{C}'$) admits an injective resolution of bounded length and, in \sref{0.11.5.6}, that the second degrees of $L^{\bullet\bullet}$ and $L'^{\bullet\bullet}$ are bounded above.
\end{env}

\subsection{Hyperhomology of a functor with respect to a complex $K_{\bullet}$}
\label{subsection:0.11.6}

\begin{env}[11.6.1]
\label{0.11.6.1}
Let $\cat{C}$ be an abelian category and let $K_\bullet=(K_i)_{i\in\bb{Z}}$ be a complex of objects of $\cat{C}$ whose differential has degree $-1$.
A \emph{left Cartan--Eilenberg resolution} of $K_\bullet$ consists of a bicomplex $L_{\bullet\bullet}=(L_{ij})$, whose differentials have degree $-1$ and for which $L_{ij}=0$ when $j<0$, together with a morphism of simple complexes
\[
  \varepsilon:L_{\bullet,0}\longrightarrow K_\bullet,
\]
such that the conditions obtained from those of \sref{0.11.4.2} by ``reversing the arrows'' are satisfied.
If every object of $\cat{C}$ is a quotient of a projective object, then every complex $K_\bullet$ of $\cat{C}$ admits a \emph{projective Cartan--Eilenberg resolution}, that is, one formed of projective objects $L_{ij}$; it has functorial properties similar to those of \sref{0.11.4.2}.
Moreover, if $K_\bullet$ is bounded below (respectively, above), one may require $L_{ij}=0$ for $i<i_0$ (respectively, $i>i_0$) whenever $K_i=0$ for $i<i_0$ (respectively, $i>i_0$).
If every object of $\cat{C}$ admits a projective resolution of length at most $n$, one may require $L_{ij}=0$ for $j>n$.
\end{env}

\begin{env}[11.6.2]
\label{0.11.6.2}
\oldpage[0\textsubscript{III}]{40}
Suppose that $T$ is an additive covariant functor from $\cat{C}$ to an abelian category $\cat{C}'$.
The \emph{hyperhomology} $\LL_\bullet T(K_\bullet)$ and the \emph{hyperhomology spectral sequences} of $T$ with respect to a complex $K_\bullet$ of $\cat{C}$ (when they exist) are defined from \sref{0.11.4.3}, again by ``reversing the arrows''.
The $E^2$ terms of the two spectral sequences so obtained are
\[
\label{0.11.6.2.1}
\tag{11.6.2.1}
  {}'E^2_{pq}
  =\HH_p(\LL_qT(K_\bullet))
\]
and
\[
\label{0.11.6.2.2}
\tag{11.6.2.2}
  {}''E^2_{pq}
  =\LL_pT(\HH_q(K_\bullet)),
\]
where $\LL_pT$ denotes the $p$th left-derived functor of $T$ for $p\geq0$, and zero for $p<0$, while $\LL_qT(K_\bullet)$ denotes the complex $(\LL_qT(K_i))_{i\in\bb{Z}}$.

The properties of hyperhomology do not all follow from those of hypercohomology merely by ``reversing the arrows'' (unless additional hypotheses of type (AB~5*) from (T, 1.5) are imposed on $\cat{C}'$).
The reason is the regularity conditions on the two preceding spectral sequences, to which one must now apply the criteria of \sref{0.11.3.4}.
Those criteria show that, if filtered inductive limits exist and are exact in $\cat{C}'$, then the complex defined by the bicomplex $T(L_{\bullet\bullet})$ exists and the \emph{second} spectral sequence ${}''E(T(L_{\bullet\bullet}))$ is regular.
If either $K_\bullet$ is bounded below or there is an integer $n$ such that every object of $\cat{C}$ admits a projective resolution of length at most $n$, then the \emph{two} hyperhomology spectral sequences exist (without any hypothesis on $\cat{C}'$) and are biregular.
\end{env}

\begin{proposition}[11.6.3]
\label{0.11.6.3}
Let $\cat{C}$ and $\cat{C}'$ be two abelian categories, and let $T$ be an additive covariant functor from $\cat{C}$ to $\cat{C}'$.
Then:
\begin{enumerate}
  \item[{\rm(i)}] \emph{The hyperhomology $\LL_\bullet T(K_\bullet)$ is a homological functor on the abelian category of complexes of $\cat{C}$ that are bounded below.}

  \item[{\rm(ii)}] \emph{Let $K_\bullet$ be a complex of $\cat{C}$ bounded below.
  If $\LL_nT(K_i)=0$ for every $n>0$ and every $i\in\bb{Z}$, there are functorial isomorphisms}
  \[
  \label{0.11.6.3.1}
  \tag{11.6.3.1}
    \LL_iT(K_\bullet)
    \simto
    \HH_i(T(K_\bullet))
    \qquad(i\in\bb{Z}).
  \]

  \item[{\rm(iii)}] \emph{Let $K_\bullet$ be a complex of $\cat{C}$ bounded below.
  Let $L_{\bullet\bullet}=(L_{ij})$ be a bicomplex such that $L_{ij}=0$ for $j<0$ and, for every $i$, $L_{i,\bullet}$ is a resolution of $K_i$.
  Suppose in addition that $\LL_nT(L_{ij})=0$ for every pair $(i,j)$ and every $n>0$.
  Then there is a functorial isomorphism}
  \[
  \label{0.11.6.3.2}
  \tag{11.6.3.2}
    \LL_\bullet T(K_\bullet)
    \simto
    \HH_\bullet(T(L_{\bullet\bullet})).
  \]
\end{enumerate}

The proofs proceed as those of \sref{0.11.5.2}, \sref{0.11.4.5}, and \sref{0.11.5.3} in the case of complexes bounded below.
We leave the details of these arguments to the reader.
\end{proposition}

\begin{env}[11.6.4]
\label{0.11.6.4}
Entirely analogous results hold for covariant multifunctors that are additive in each argument.
For example, for a bifunctor $T$, the two hyperhomology spectral sequences have $E^2$ terms
\oldpage[0\textsubscript{III}]{41}
\[
\label{0.11.6.4.1}
\tag{11.6.4.1}
  {}'E^2_{pq}
  =\HH_p(\LL_qT(K_\bullet,K'_\bullet))
\]
and
\[
\label{0.11.6.4.2}
\tag{11.6.4.2}
  {}''E^2_{pq}
  =
  \sum_{q'+q''=q}
  \LL_pT(\HH_{q'}(K_\bullet),\HH_{q''}(K'_\bullet)).
\]
Here again it is the \emph{second} spectral sequence that is regular.
The two sequences are biregular when $K_\bullet$ and $K'_\bullet$ are bounded below, or when the objects of the abelian categories under consideration admit projective resolutions of a fixed length.

Moreover:
\end{env}

\begin{proposition}[11.6.5]
\label{0.11.6.5}
Let $\cat{C}$, $\cat{C}'$, and $\cat{C}''$ be three abelian categories, and let $T$ be a covariant biadditive bifunctor from $\cat{C}\times\cat{C}'$ to $\cat{C}''$.
\begin{enumerate}
  \item[{\rm(i)}] \emph{$\LL_\bullet T(K_\bullet,K'_\bullet)$ is a bifunctor that is homological in each of the bounded-below complexes $K_\bullet,K'_\bullet$, formed respectively of objects of $\cat{C}$ and of $\cat{C}'$.}

  \item[{\rm(ii)}] \emph{Suppose that $K_\bullet$ and $K'_\bullet$ are bounded below.
  Let $L_{\bullet\bullet}=(L_{ij})$ and $L'_{\bullet\bullet}=(L'_{ij})$ be two bicomplexes such that $L_{ij}=L'_{ij}=0$ for $j<0$ and such that, for every $i$, $L_{i,\bullet}$ is a resolution of $K_i$ and $L'_{i,\bullet}$ is a resolution of $K'_i$.
  Suppose finally that $\LL_nT(L_{ij},L'_{hk})=0$ for every $n>0$ and every system of indices $(i,j,h,k)$.
  Then there is a functorial isomorphism}
  \[
  \label{0.11.6.5.1}
  \tag{11.6.5.1}
    \LL_\bullet T(K_\bullet,K'_\bullet)
    \simto
    \HH_\bullet(T(L_{\bullet\bullet},L'_{\bullet\bullet})).
  \]

  \item[{\rm(iii)}] \emph{Suppose in addition that, for every pair $(i,j)$ and every pair $(h,k)$, the functors
  \[
    A\longmapsto T(A,L'_{hk}),
    \qquad
    A'\longmapsto T(L_{ij},A')
  \]
  are exact in $\cat{C}$ and $\cat{C}'$, respectively.
  Then there are functorial isomorphisms}
  \[
  \label{0.11.6.5.2}
  \tag{11.6.5.2}
    \LL_\bullet T(K_\bullet,K'_\bullet)
    \simto
    \HH_\bullet(T(L_{\bullet\bullet},K'_\bullet))
    \simto
    \HH_\bullet(T(K_\bullet,L'_{\bullet\bullet})).
  \]
\end{enumerate}

The proofs are analogous to those of \sref{0.11.5.5} and \sref{0.11.5.6}.
\end{proposition}

\subsection{Hyperhomology of a functor with respect to a bicomplex $K_{\bullet\bullet}$}
\label{subsection:0.11.7}

\begin{env}[11.7.1]
\label{0.11.7.1}
Let $\cat{C}$ be an abelian category in which every object is a quotient of a
projective object.
Consider a bicomplex $K_{\bullet\bullet}=(K_{ij})$ of objects of $\cat{C}$ whose
two degrees are bounded below.
We may always restrict ourselves to the case where $K_{ij}=0$ for $i<0$ or
$j<0$, and we shall do so from now on.
We may regard $K_{\bullet\bullet}$ as an ordinary complex consisting of the
objects
\[
  K_{i,\bullet}=(K_{ij})_{j\geq 0}
\]
of the abelian category $\cat{K}$ of nonnegative complexes of objects of
$\cat{C}$.
It follows from Lemma~\sref{0.11.5.2.1} (or, more precisely, from the
``dual'' lemma obtained by ``reversing the arrows'') and from (M,~V,~2.2)
that $K_{\bullet\bullet}$ admits a \emph{projective Cartan--Eilenberg
resolution} in the category $\cat{K}$.
Such a resolution is a tricomplex
$M_{\bullet\bullet\bullet}=(M_{ijk})$ in $\cat{C}$, all of whose degrees are
nonnegative and whose objects are projective, such that, for every $i$,
\[
  M_{i,\bullet\bullet},\qquad
  \mathrm{B}^{\mathrm{I}}_i(M_{\bullet\bullet\bullet}),\qquad
  \mathrm{Z}^{\mathrm{I}}_i(M_{\bullet\bullet\bullet}),\qquad
  \HH^{\mathrm{I}}_i(M_{\bullet\bullet\bullet})
\]
are projective resolutions in $\cat{K}$ of, respectively,
\[
  K_{i,\bullet},\qquad
  \mathrm{B}^{\mathrm{I}}_i(K_{\bullet\bullet}),\qquad
  \mathrm{Z}^{\mathrm{I}}_i(K_{\bullet\bullet}),\qquad
  \HH^{\mathrm{I}}_i(K_{\bullet\bullet}).
\]
In particular, for every pair $(i,j)$, $M_{ij,\bullet}$ is a projective
resolution of $K_{ij}$ in $\cat{C}$.
\end{env}

\begin{proposition}[11.7.2]
\label{0.11.7.2}
Let $T$ be an additive covariant functor from $\cat{C}$ to an abelian
category $\cat{C}'$.
With the notation of \sref{0.11.7.1}, the homology
\[
  \HH_\bullet\bigl(T(M_{\bullet\bullet\bullet})\bigr)
\]
of the simple complex associated with the tricomplex
$T(M_{\bullet\bullet\bullet})$ is canonically isomorphic to the
hyperhomology $\LL_\bullet T(K_{\bullet\bullet})$ of
\oldpage[0\textsubscript{III}]{42}
the simple complex associated with $K_{\bullet\bullet}$
\sref{0.11.6.2}.
It is also the common abutment of six biregular spectral sequences denoted
${}^{(t)}\mathrm{E}$, where $t=a,b,a',b',c$, or $d$, whose
$\mathrm{E}^2$ terms are
\[
\begin{aligned}
  {}^{(a)}\mathrm{E}^2_{pq}
    &= \LL_pT\bigl(\HH_q(K_{\bullet\bullet})\bigr),\\
  {}^{(b)}\mathrm{E}^2_{pq}
    &= \HH_p\bigl(\LL^{\mathrm{II}}_qT(K_{\bullet\bullet})\bigr),\\
  {}^{(a')}\mathrm{E}^2_{pq}
    &= \LL_pT\bigl(\HH^{\mathrm{I}}_q(K_{\bullet\bullet})\bigr),\\
  {}^{(b')}\mathrm{E}^2_{pq}
    &= \HH_p\bigl(\LL_qT(K_{\bullet\bullet})\bigr),\\
  {}^{(c)}\mathrm{E}^2_{pq}
    &= \LL_pT\bigl(\HH^{\mathrm{II}}_q(K_{\bullet\bullet})\bigr),\\
  {}^{(d)}\mathrm{E}^2_{pq}
    &= \HH_p\bigl(\LL^{\mathrm{I}}_qT(K_{\bullet\bullet})\bigr).
\end{aligned}
\]
\end{proposition}

Recall that, for a complex $A_\bullet=(A_i)$, we use the notation
$F(A_\bullet)$ for the complex whose objects are the $F(A_i)$.
For example, $\LL^{\mathrm{II}}_qT(K_{\bullet\bullet})$ denotes the complex
\[
  \bigl(\LL^{\mathrm{II}}_qT(K_{i,\bullet})\bigr)_{i\geq 0},
\]
where $\LL^{\mathrm{II}}_qT(K_{i,\bullet})$ is the hyperhomology of index
$q$ of the functor $T$ relative to the simple complex $K_{i,\bullet}$.

\begin{proof}
Let $L_\bullet$ denote the simple complex associated with
$K_{\bullet\bullet}$, so that
\[
  L_i=\bigoplus_{r+s=i}K_{rs},
\]
and set
\[
  N_{ij}=\bigoplus_{r+s=i}M_{rsj}.
\]
It is clear that, for every $i$, $N_{i,\bullet}$ is a projective resolution
of $L_i$ in $\cat{C}$.
It follows from \sref{0.11.6.3} and \sref{0.11.6.4} that there is a
functorial isomorphism
\[
  \LL_\bullet T(L_\bullet)
  \isoto
  \HH_\bullet\bigl(T(N_{\bullet\bullet})\bigr).
\]
Since the simple complex associated with the bicomplex
$T(N_{\bullet\bullet})$ is also the simple complex associated with the
tricomplex $T(M_{\bullet\bullet\bullet})$, this proves the first assertion.

Moreover, $\LL_\bullet T(L_\bullet)$ is the abutment of the two
hyperhomology spectral sequences \sref{0.11.6.2} of $T$ relative to the
simple complex $L_\bullet$.
These are precisely ${}^{(b')}\mathrm{E}$ and ${}^{(a)}\mathrm{E}$,
respectively.

Now regard $M_{\bullet\bullet\bullet}$ as a bicomplex
$U_{\bullet\bullet}$, where
\[
  U_{ij}=\bigoplus_{r+s=j}M_{irs}.
\]
The homology $\HH_\bullet(T(M_{\bullet\bullet\bullet}))$ is also the
abutment of the two spectral sequences of the bicomplex
$T(U_{\bullet\bullet})$.
For every $i$, the bicomplex $M_{i,\bullet\bullet}$ satisfies the
hypotheses of \sref{0.11.6.3} relative to the simple complex
$K_{i,\bullet}$.
Consequently,
\[
  \HH^{\mathrm{II}}_q\bigl(T(U_{i,\bullet})\bigr)
  =
  \LL^{\mathrm{II}}_qT(K_{i,\bullet}),
\]
and the first spectral sequence of $T(U_{\bullet\bullet})$ is precisely
${}^{(b)}\mathrm{E}$.

On the other hand, for every $r$, $M_{\bullet,r,\bullet}$ is a
Cartan--Eilenberg resolution of the simple complex $K_{\bullet,r}$.
The calculation made in (M,~XV,~2) shows that
\[
  \HH^{\mathrm{I}}_q\bigl(T(M_{\bullet,rs})\bigr)
  =
  T\bigl(\HH^{\mathrm{I}}_q(M_{\bullet,rs})\bigr),
\]
and hence
\[
  \HH^{\mathrm{I}}_q\bigl(T(U_{\bullet,j})\bigr)
  =
  \bigoplus_{r+s=j}
  T\bigl(\HH^{\mathrm{I}}_q(M_{\bullet,rs})\bigr).
\]
In other words, the simple complex
$\HH^{\mathrm{I}}_q(T(U_{\bullet\bullet}))$ is the simple complex
associated with the bicomplex
$T(\HH^{\mathrm{I}}_q(M_{\bullet\bullet\bullet}))$.
For every $q$, $\HH^{\mathrm{I}}_q(M_{\bullet\bullet\bullet})$ is a
projective resolution in $\cat{K}$ of the simple complex
$\HH^{\mathrm{I}}_q(K_{\bullet\bullet})$.
Applying \sref{0.11.6.3}, we obtain
\[
  \HH^{\mathrm{II}}_p
  \bigl(T(\HH^{\mathrm{I}}_q(M_{\bullet\bullet\bullet}))\bigr)
  =
  \LL_pT\bigl(\HH^{\mathrm{I}}_q(K_{\bullet\bullet})\bigr),
\]
and hence the spectral sequence ${}^{(a')}\mathrm{E}$.
Finally, the spectral sequences ${}^{(c)}\mathrm{E}$ and
${}^{(d)}\mathrm{E}$ are obtained by interchanging the roles of the
indices $i$ and $j$ in the definition of the tricomplex
$M_{\bullet\bullet\bullet}$ and applying the preceding arguments to the
resulting tricomplex.
\end{proof}

We call $\LL_\bullet T(K_{\bullet\bullet})$ the \emph{hyperhomology} of
$T$ relative to the bicomplex $K_{\bullet\bullet}$.

\begin{env}[11.7.3]
\label{0.11.7.3}
\emph{Remarks}.
\begin{enumerate}
  \item[{\rm(i)}]
  It follows from \sref{0.11.6.3} that
  $\LL_\bullet T(K_{\bullet\bullet})$ is a homological functor on the
  category of bicomplexes $K_{\bullet\bullet}$ in $\cat{C}$ that are
  bounded below in each degree.

\oldpage[0\textsubscript{III}]{43}
  \item[{\rm(ii)}]
  Let $M_{\bullet\bullet\bullet}$ be a tricomplex in $\cat{C}$ such that,
  for every pair $(r,s)$, $M_{rs,\bullet}$ is a resolution of $K_{rs}$,
  and such that
  \[
    \LL_nT(M_{ijk})=0
  \]
  for every triple $(i,j,k)$ and every $n>0$.
  Then there is an isomorphism
  \[
    \LL_\bullet T(K_{\bullet\bullet})
    \isoto
    \HH_\bullet\bigl(T(M_{\bullet\bullet\bullet})\bigr).
  \]
  Indeed, with the notation used in the proof of \sref{0.11.7.2},
  $N_{i,\bullet}$ is a resolution of $L_i$ such that
  $\LL_nT(N_{ij})=0$ for every $n>0$ and every pair $(i,j)$, so it is
  enough to apply \sref{0.11.6.3},~(iii).

  \item[{\rm(iii)}]
  The results of \sref{0.11.7.2} immediately generalize to covariant
  multifunctors.
  For example, let $\cat{C}'$ be a second abelian category in which every
  object is a quotient of a projective object, let
  $K'_{\bullet\bullet}$ be a bicomplex in $\cat{C}'$ whose two degrees are
  bounded below, and let $T$ be an additive covariant bifunctor from
  $\cat{C}\times\cat{C}'$ to an abelian category $\cat{C}''$.
  If $L_\bullet$ and $L'_\bullet$ are the simple complexes associated with
  $K_{\bullet\bullet}$ and $K'_{\bullet\bullet}$, respectively, we define
  the hyperhomology of $T$ relative to the two bicomplexes
  $K_{\bullet\bullet}$ and $K'_{\bullet\bullet}$ to be
  \[
    \LL_\bullet T(L_\bullet,L'_\bullet).
  \]
  Applying \sref{0.11.6.4} and \sref{0.11.6.5}, we obtain, just as in
  \sref{0.11.7.2}, six spectral sequences abutting to this hyperhomology;
  we leave their explicit formulation to the reader.
\end{enumerate}
\end{env}

\subsection{Supplement on the cohomology of simplicial complexes}
\label{subsection:0.11.8}

\begin{env}[11.8.1]
\label{0.11.8.1}
Let $A$ be a finite set, and let $\Sigma(A)$ be the set of finite
sequences
\[
  \sigma=(\alpha_0,\ldots,\alpha_h)
\]
of elements of $A$ (the ``simplices'' of $A$); set
$|\sigma|=\{\alpha_0,\ldots,\alpha_h\}$.
Recall that the chain complex $\mathrm{C}_\bullet(A)$ is the graded free
abelian group generated by the elements of $\Sigma(A)$, with
$(\alpha_0,\ldots,\alpha_h)$ in degree $h$, and with differential
\[
  d(\alpha_0,\ldots,\alpha_h)
  =
  \sum_{j=0}^{h}(-1)^j
  (\alpha_0,\ldots,\widehat{\alpha_j},\ldots,\alpha_h).
\]
The subgroup $\mathrm{D}_\bullet(A)$ of $\mathrm{C}_\bullet(A)$ generated
by those chains $\sigma=(\alpha_0,\ldots,\alpha_h)$ for which two of the
$\alpha_i$ are equal, and by the chains
\[
  \pi(\sigma)-\varepsilon_\pi\sigma,
\]
where, for every permutation $\pi$,
\[
  \pi(\sigma)
  =
  (\alpha_{\pi(0)},\ldots,\alpha_{\pi(h)})
\]
and $\varepsilon_\pi$ is the sign of $\pi$, is a subcomplex of
$\mathrm{C}_\bullet(A)$.
Its elements are called \emph{degenerate chains}.
Set
\[
  \mathrm{L}_\bullet(A)
  =
  \mathrm{C}_\bullet(A)/\mathrm{D}_\bullet(A);
\]
there is a natural homomorphism of complexes
\[
  p:\mathrm{C}_\bullet(A)\longrightarrow\mathrm{L}_\bullet(A).
\]

Conversely, define a homomorphism of complexes
\[
  j:\mathrm{L}_\bullet(A)\longrightarrow\mathrm{C}_\bullet(A)
\]
as follows.
Choose a total ordering of $A$.
To the class modulo $\mathrm{D}_\bullet(A)$ of a simplex
$\sigma=(\alpha_0,\ldots,\alpha_h)$, associate $0$ if two of the
$\alpha_i$ are equal, and otherwise associate $\varepsilon_\pi(\beta_0,\ldots,\beta_h)$,
where $(\beta_0,\ldots,\beta_h)=\pi(\sigma)$ is the sequence of the $\alpha_i$ in
increasing order.
It is clear that $p\circ j$ is the identity of
$\mathrm{L}_\bullet(A)$.
\end{env}

\begin{env}[11.8.2]
\label{0.11.8.2}
Let $B$ be a second finite set.
If $d'$ and $d''$ are the differentials of
$\mathrm{C}_\bullet(A)$ and $\mathrm{C}_\bullet(B)$, the tensor-product
complex
\[
  \mathrm{C}_\bullet(A)\otimes\mathrm{C}_\bullet(B)
\]
may be regarded as the free abelian group generated by the elements of
$\Sigma(A)\times\Sigma(B)$, with differential
\[
  d(\sigma,\tau)
  =
  (d'\sigma,\tau)+(-1)^{h+1}(\sigma,d''\tau)
  \qquad\text{if }\deg\sigma=h.
\]

The natural homomorphisms
$\mathrm{C}_\bullet(A)\to\mathrm{L}_\bullet(A)$ and
$\mathrm{C}_\bullet(B)\to\mathrm{L}_\bullet(B)$ define a homomorphism
\[
  p:
  \mathrm{C}_\bullet(A)\otimes\mathrm{C}_\bullet(B)
  \longrightarrow
  \mathrm{L}_\bullet(A)\otimes\mathrm{L}_\bullet(B),
\]
the latter tensor product being isomorphic to
\[
  \frac{\mathrm{C}_\bullet(A)\otimes\mathrm{C}_\bullet(B)}
  {\mathrm{C}_\bullet(A)\otimes\mathrm{D}_\bullet(B)
   +\mathrm{D}_\bullet(A)\otimes\mathrm{C}_\bullet(B)}.
\]
Likewise, using the homomorphisms
$\mathrm{L}_\bullet(A)\to\mathrm{C}_\bullet(A)$ and
$\mathrm{L}_\bullet(B)\to\mathrm{C}_\bullet(B)$ defined in
\sref{0.11.8.1}, with total orderings on $A$ and $B$, one obtains a
homomorphism
\[
  j:
  \mathrm{L}_\bullet(A)\otimes\mathrm{L}_\bullet(B)
  \longrightarrow
  \mathrm{C}_\bullet(A)\otimes\mathrm{C}_\bullet(B)
\]
such that $p\circ j$ is the identity.
\end{env}

\oldpage[0\textsubscript{III}]{44}
With this notation:
\begin{proposition}[11.8.3]
\label{0.11.8.3}
There exists a homotopy
\[
  h:
  \mathrm{C}_\bullet(A)\otimes\mathrm{C}_\bullet(B)
  \longrightarrow
  \mathrm{C}_\bullet(A)\otimes\mathrm{C}_\bullet(B)
\]
such that $h(\sigma,\tau)$ is a linear combination of pairs of simplices
$(\sigma_i,\tau_i)$ satisfying
$|\sigma_i|\subset|\sigma|$ and $|\tau_i|\subset|\tau|$, and such that,
if $f=j\circ p$, then
\[
  f-1=h\circ d+d\circ h.
  \tag{11.8.3.1}\label{0.11.8.3.1}
\]
\end{proposition}

\begin{proof}
It suffices to define $h$ on each pair $(\sigma,\tau)$ of simplices, by
induction on the sum of the degrees of $\sigma$ and $\tau$; when that sum
is $0$, one may take $h=0$.
Set
\[
  \omega
  =
  f(\sigma,\tau)-(\sigma,\tau)-h\bigl(d(\sigma,\tau)\bigr).
\]
By the induction hypothesis and the definition of $d$, one has
\[
  \omega
  \in
  \mathrm{C}_\bullet(|\sigma|)
  \otimes
  \mathrm{C}_\bullet(|\tau|).
\]
Moreover,
\[
  d\omega
  =
  f\bigl(d(\sigma,\tau)\bigr)
  -d(\sigma,\tau)
  -d\bigl(h(d(\sigma,\tau))\bigr)
  =
  h\bigl(d(d(\sigma,\tau))\bigr)
  =
  0
\]
by \eqref{0.11.8.3.1} and the induction hypothesis.
Now
\[
  \HH_q\bigl(\mathrm{C}_\bullet(A)\bigr)=0
  \qquad(q>0)
\]
(G, I, 3.7.4), and consequently
\[
  \HH_q\bigl(
    \mathrm{C}_\bullet(A)\otimes\mathrm{C}_\bullet(B)
  \bigr)=0
  \qquad(q>0)
\]
by the K\"unneth formula (G, I, 5.5.2).
Applying this observation with $A$ replaced by $|\sigma|$ and $B$ by
$|\tau|$, one finds
\[
  \omega'
  \in
  \mathrm{C}_\bullet(|\sigma|)
  \otimes
  \mathrm{C}_\bullet(|\tau|)
\]
such that $\omega=d\omega'$.
Taking $h(\sigma,\tau)=\omega'$ verifies \eqref{0.11.8.3.1} for the pair
$(\sigma,\tau)$, and the induction continues.
\end{proof}

\begin{env}[11.8.4]
\label{0.11.8.4}
With the notation of \sref{0.11.8.2}, write
\[
  (\sigma,\tau)\leq(\sigma',\tau')
\]
if $|\sigma|\subset|\sigma'|$ and $|\tau|\subset|\tau'|$.
A \emph{coefficient system} $\sh{S}$ on
$\Sigma(A)\times\Sigma(B)$ consists of a family
$(\Gamma_{\sigma,\tau})$ of abelian groups, where
$\Gamma_{\sigma,\tau}$ depends only on the sets $|\sigma|$ and
$|\tau|$, together with homomorphisms
\[
  \Gamma_{\sigma,\tau}
  \longrightarrow
  \Gamma_{\sigma',\tau'}
  \qquad
  \bigl((\sigma,\tau)\leq(\sigma',\tau')\bigr)
\]
forming an inductive system for this preorder.

Define a cochain complex $\mathrm{C}^\bullet(A,B;\sh{S})$ whose elements
are the families
\[
  \lambda=\bigl(\lambda(\sigma,\tau)\bigr),
  \qquad
  \lambda(\sigma,\tau)\in\Gamma_{\sigma,\tau},
\]
as $(\sigma,\tau)$ ranges over $\Sigma(A)\times\Sigma(B)$.
Its differential is defined as follows.
If
\[
  d(\sigma,\tau)=\sum_i\pm(\sigma_i,\tau_i),
\]
then $|\sigma_i|\subset|\sigma|$ and $|\tau_i|\subset|\tau|$ for every
$i$, and one sets
\[
  d\lambda(\sigma,\tau)
  =
  \sum_i\pm\lambda_i(\sigma_i,\tau_i),
\]
where $\lambda_i(\sigma_i,\tau_i)$ denotes the canonical image of
$\lambda(\sigma_i,\tau_i)$ in $\Gamma_{\sigma,\tau}$.

A cochain $\lambda\in\mathrm{C}^\bullet(A,B;\sh{S})$ is called
\emph{bi-alternating} if $\lambda(\sigma,\tau)=0$ whenever either
$\sigma$ or $\tau$ has two equal terms, and if
\[
  \lambda(\pi(\sigma),\tau)
  =
  \varepsilon_\pi\lambda(\sigma,\tau),
  \qquad
  \lambda(\sigma,\pi'(\tau))
  =
  \varepsilon_{\pi'}\lambda(\sigma,\tau)
\]
for arbitrary permutations $\pi$ and $\pi'$ of the indices.
These cochains plainly generate a subcomplex
$\mathrm{L}^\bullet(A,B;\sh{S})$ of
$\mathrm{C}^\bullet(A,B;\sh{S})$.
\end{env}

\begin{proposition}[11.8.5]
\label{0.11.8.5}
The canonical injection
\[
  \mathrm{L}^\bullet(A,B;\sh{S})
  \longrightarrow
  \mathrm{C}^\bullet(A,B;\sh{S})
\]
induces an isomorphism on cohomology.
\end{proposition}

\begin{proof}
If $p$ and $j$ have the meanings assigned in \sref{0.11.8.2}, then
the maps
\[
  {}'\!p:\lambda\longmapsto\lambda\circ p,
  \qquad
  {}'\!j:\lambda\longmapsto\lambda\circ j
\]
are defined on $\mathrm{L}^\bullet(A,B;\sh{S})$ and
$\mathrm{C}^\bullet(A,B;\sh{S})$, respectively; the first is precisely
the canonical injection.
Since ${}'\!j\circ{}'\!p$ is the identity, it remains only to prove
that ${}'\!p\circ{}'\!j$ is homotopic to the identity.
By \sref{0.11.8.3}, the map
\[
  {}'\!h:\lambda\longmapsto\lambda\circ h
\]
is defined on $\mathrm{C}^\bullet(A,B;\sh{S})$, so transposing
\eqref{0.11.8.3.1} gives the desired result.
\end{proof}

\begin{env}[11.8.6]
\label{0.11.8.6}
\oldpage[0\textsubscript{III}]{45}
Proposition~\sref{0.11.8.5} reduces the computation of the cohomology of
$\mathrm{L}^\bullet(A,B;\sh{S})$ to that of
$\mathrm{C}^\bullet(A,B;\sh{S})$.
On the other hand, by the Eilenberg--Zilber theorem
(G, I, 3.10.2), the latter cohomology is canonically isomorphic to the
cohomology obtained from the chain complex
$\mathrm{P}_\bullet(A,B)$ of linear combinations of pairs
$(\sigma,\tau)\in\Sigma(A)\times\Sigma(B)$ for which $\sigma$ and
$\tau$ have the same degree.
The differential of this complex is
\[
  d(\sigma,\tau)
  =
  \sum_{j,k}(-1)^{j+k}(\sigma_j,\tau_k)
\]
when
\[
  d\sigma=\sum_j(-1)^j\sigma_j,
  \qquad
  d\tau=\sum_k(-1)^k\tau_k.
\]
There are two canonical homomorphisms of complexes
\[
  f:\mathrm{P}_\bullet(A,B)
  \longrightarrow
  \mathrm{C}_\bullet(A)\otimes\mathrm{C}_\bullet(B),
  \qquad
  g:\mathrm{C}_\bullet(A)\otimes\mathrm{C}_\bullet(B)
  \longrightarrow
  \mathrm{P}_\bullet(A,B),
\]
and it is proved there that homotopies $h$ and $h'$ exist such that
\[
  f\circ g-1=d\circ h+h\circ d,
  \qquad
  g\circ f-1=d\circ h'+h'\circ d.
\]
Moreover,
\[
  f(\sigma,\tau)
  \in
  \mathrm{C}_\bullet(|\sigma|)
  \otimes
  \mathrm{C}_\bullet(|\tau|),
  \qquad
  g(\sigma,\tau)
  \in
  \mathrm{P}_\bullet(|\sigma|,|\tau|),
\]
and the homotopies may be chosen so that
\[
  h(\sigma,\tau)
  \in
  \mathrm{C}_\bullet(|\sigma|)
  \otimes
  \mathrm{C}_\bullet(|\tau|),
  \qquad
  h'(\sigma,\tau)
  \in
  \mathrm{P}_\bullet(|\sigma|,|\tau|).
\]
Indeed, $h(\sigma,\tau)$ and $h'(\sigma,\tau)$ may be defined by
induction on the sum of the degrees of $\sigma$ and $\tau$, while
\[
  \HH_q\bigl(
    \mathrm{C}_\bullet(|\sigma|)
    \otimes
    \mathrm{C}_\bullet(|\tau|)
  \bigr)
  =
  \HH_q\bigl(\mathrm{P}_\bullet(|\sigma|,|\tau|)\bigr)
  =
  0
  \qquad(q>0)
\]
(\emph{loc.\ cit.}); the argument is the same as in
\sref{0.11.8.3}.

Now define $\mathrm{P}^\bullet(A,B;\sh{S})$ to be the set of families
\[
  \lambda=\bigl(\lambda(\sigma,\tau)\bigr),
  \qquad
  \lambda(\sigma,\tau)\in\Gamma_{\sigma,\tau},
\]
where $(\sigma,\tau)$ ranges over pairs of simplices of the same degree.
When
\[
  d\sigma=\sum_j(-1)^j\sigma_j,
  \qquad
  d\tau=\sum_k(-1)^k\tau_k,
\]
the formula
\[
  d\lambda(\sigma,\tau)
  =
  \sum_{j,k}(-1)^{j+k}\lambda(\sigma_j,\tau_k)
\]
is defined and gives the differential of
$\mathrm{P}^\bullet(A,B;\sh{S})$.
The maps
\[
  {}'\!f:\lambda\longmapsto\lambda\circ f,\quad
  {}'\!g:\lambda\longmapsto\lambda\circ g,\quad
  {}'\!h:\lambda\longmapsto\lambda\circ h,\quad
  {}'\!h':\lambda\longmapsto\lambda\circ h'
\]
are all defined by the preceding support properties.
Thus ${}'\!f\circ{}'\!g$ and ${}'\!g\circ{}'\!f$ are homotopic to the
identity, giving the required isomorphism between the cohomology of
$\mathrm{C}^\bullet(A,B;\sh{S})$ and that of
$\mathrm{P}^\bullet(A,B;\sh{S})$.
\end{env}

\begin{env}[11.8.7]
\label{0.11.8.7}
\emph{Remark}.
The argument of \sref{0.11.8.3}, applied to
$\mathrm{C}_\bullet(A)$ and $\mathrm{L}_\bullet(A)$, shows that if
$j$ and $p$ are defined as in \sref{0.11.8.1}, then
$f=j\circ p$ again satisfies \eqref{0.11.8.3.1}, with
$|h(\sigma)|\subset|\sigma|$.
As in \sref{0.11.8.5}, one obtains an isomorphism from the cohomology of
$\mathrm{L}^\bullet(A;\sh{S})$ to that of
$\mathrm{C}^\bullet(A;\sh{S})$, the two complexes being defined in the
obvious way.
This is the result whose proof is sketched in (G, I, 3.8.1).
\end{env}

\begin{env}[11.8.8]
\label{0.11.8.8}
Return now to the notation and hypotheses of \sref{0.11.8.2}, and let
\[
  \sh{S}^\bullet=(\sh{S}^k)
\]
be a complex of coefficient systems on
$\Sigma(A)\times\Sigma(B)$.
For each $(\sigma,\tau)$, the groups $\Gamma^k_{\sigma,\tau}$
therefore form a complex of abelian groups $(k\in\bb{Z})$, and for
$(\sigma,\tau)\leq(\sigma',\tau')$ there are commutative diagrams
\[
  \xymatrix{
    \Gamma^k_{\sigma,\tau}\ar[r]\ar[d] &
    \Gamma^{k+1}_{\sigma,\tau}\ar[d]\\
    \Gamma^k_{\sigma',\tau'}\ar[r] &
    \Gamma^{k+1}_{\sigma',\tau'}.
  }
\]
\oldpage[0\textsubscript{III}]{46}
It follows immediately that
\[
  \mathrm{C}^\bullet(A,B;\sh{S}^\bullet)
  =
  \bigl(
    \mathrm{C}^h(A,B;\sh{S}^k)
  \bigr)_{(h,k)\in\bb{Z}\times\bb{Z}}
\]
is a bicomplex of abelian groups, and that
\[
  \mathrm{L}^\bullet(A,B;\sh{S}^\bullet)
  =
  \bigl(
    \mathrm{L}^h(A,B;\sh{S}^k)
  \bigr)
\]
is a sub-bicomplex.
\end{env}

\begin{proposition}[11.8.9]
\label{0.11.8.9}
The canonical injection
\[
  \mathrm{L}^\bullet(A,B;\sh{S}^\bullet)
  \longrightarrow
  \mathrm{C}^\bullet(A,B;\sh{S}^\bullet)
\]
induces an isomorphism on the cohomology of these two bicomplexes.
\end{proposition}

\begin{proof}
For brevity set
\[
  \mathrm{C}^{\bullet\bullet}
  =
  \mathrm{C}^\bullet(A,B;\sh{S}^\bullet),
  \qquad
  \mathrm{L}^{\bullet\bullet}
  =
  \mathrm{L}^\bullet(A,B;\sh{S}^\bullet).
\]
Since $\mathrm{C}^{hk}=\mathrm{L}^{hk}=0$ for $h<0$, the second
spectral sequences of these bicomplexes are regular
\sref{0.11.3.3}.
The homomorphism
$\mathrm{L}^{\bullet\bullet}\to\mathrm{C}^{\bullet\bullet}$ therefore
gives a morphism of spectral sequences
\[
  {}''\mathrm{E}(\mathrm{L}^{\bullet\bullet})
  \longrightarrow
  {}''\mathrm{E}(\mathrm{C}^{\bullet\bullet})
\]
which, on the $\mathrm{E}_2$ terms, is
\[
  \HH^p_{\mathrm{II}}
  \bigl(\HH^q_{\mathrm{I}}(\mathrm{L}^{\bullet\bullet})\bigr)
  \longrightarrow
  \HH^p_{\mathrm{II}}
  \bigl(\HH^q_{\mathrm{I}}(\mathrm{C}^{\bullet\bullet})\bigr).
  \tag{11.8.9.1}\label{0.11.8.9.1}
\]
For every $k\in\bb{Z}$, however, \sref{0.11.8.5} shows that
\[
  \HH^q_{\mathrm{I}}(\mathrm{L}^{\bullet,k})
  \longrightarrow
  \HH^q_{\mathrm{I}}(\mathrm{C}^{\bullet,k})
\]
is bijective.
The conclusion follows from \sref{0.11.1.5}.
\end{proof}

\begin{env}[11.8.10]
\label{0.11.8.10}
Similarly, with the notation of \sref{0.11.8.6}, there is a canonical
homomorphism of bicomplexes
\[
  \mathrm{C}^\bullet(A,B;\sh{S}^\bullet)
  \longrightarrow
  \mathrm{P}^\bullet(A,B;\sh{S}^\bullet),
\]
with the evident notation.
The argument of \sref{0.11.8.9}, now using \sref{0.11.8.6}, shows that
this homomorphism likewise induces an isomorphism on cohomology.
\end{env}

\subsection{A lemma on the complexes of finite type}
\label{subsection:0.11.9}

\begin{proposition}[11.9.1]
\label{0.11.9.1}
Let $\cat{C}$ be an abelian category, and let $K'$ and $K''$ be subsets
of the set of objects of $\cat{C}$, with $K''\subset K'$, satisfying the
following conditions:
\begin{enumerate}
  \item[{\rm(i)}]
  \emph{For every $A'\in K'$ and every epimorphism $u:A\to A'$ in
  $\cat{C}$, there exist an object $B\in K''$ and a morphism $v:B\to A$
  such that $u\circ v$ is an epimorphism.}

  \item[{\rm(ii)}]
  \emph{For every pair $A,B\in K'$ and every epimorphism $u:A\to B$,
  the object $\Ker(u)$ belongs to $K'$.}

  \item[{\rm(iii)}]
  \emph{The product of two objects of $K''$ belongs to $K''$.}
\end{enumerate}

\emph{Let $P_\bullet=(P_i)_{i\in\bb{Z}}$ be a complex in $\cat{C}$ such
that $\HH_i(P_\bullet)\in K'$ for every $i$, and suppose that there is an
integer $d$ for which $\HH_i(P_\bullet)=0$ when $i<d$.
Then there exist a complex $Q_\bullet=(Q_i)_{i\in\bb{Z}}$ in $\cat{C}$
such that $Q_i\in K''$ for every $i$ and $Q_i=0$ for $i<d$, and a
homomorphism of complexes
\[
  u:Q_\bullet\longrightarrow P_\bullet
\]
whose induced morphism
\[
  \HH_\bullet(Q_\bullet)\longrightarrow\HH_\bullet(P_\bullet)
\]
is an isomorphism.}
\end{proposition}

\begin{proof}
We first prove the following consequence of condition~(i).

\medskip\noindent
\emph{(i bis)}
\emph{Let $u:C\to B$ be an epimorphism in $\cat{C}$, let $A$ be an
object of $K'$, and let $v:A\to B$ be a morphism in $\cat{C}$.
Then there exist an object $D\in K''$, an epimorphism $u':D\to A$, and a
morphism $v':D\to C$ such that the diagram}
\[
  \xymatrix{
    D\ar[r]^{u'}\ar[d]_{v'} & A\ar[d]^{v}\\
    C\ar[r]_{u} & B
  }
  \tag{11.9.1.1}\label{0.11.9.1.1}
\]
\emph{is commutative.}

Form the fibre product $C\times_B A$ in $\cat{C}$ and let
\[
  p:C\times_B A\longrightarrow C,
  \qquad
  q:C\times_B A\longrightarrow A
\]
be the canonical projections, so that
\oldpage[0\textsubscript{III}]{47}
\[
  \xymatrix{
    C\times_B A\ar[r]^{q}\ar[d]_{p} & A\ar[d]^{v}\\
    C\ar[r]_{u} & B
  }
  \tag{11.9.1.2}\label{0.11.9.1.2}
\]
is commutative.
The cokernel of $q$ is the quotient of $A$ by $v^{-1}(u(C))$
([27, p.~1-12]).
Since $u$ is an epimorphism, $u(C)=B$ and
$v^{-1}(u(C))=A$; hence $q$ is an epimorphism.
Applying~(i) to $q:C\times_B A\to A$, one obtains an object
$D\in K''$ and a morphism $w:D\to C\times_B A$ such that $q\circ w$ is
an epimorphism.
It is enough to take $u'=q\circ w$ and $v'=p\circ w$.

\medskip
We now prove the proposition by induction.
Suppose, for some $i\geq d-1$, that for every $j\leq i$ we have
constructed objects $Q_j$, differentials
\[
  d^Q_j:Q_j\longrightarrow Q_{j-1}
\]
with $Q_j=0$ for $j<d$ and
$d^Q_{j-1}\circ d^Q_j=0$, and morphisms
\[
  u_j:Q_j\longrightarrow P_j
\]
such that
\[
  d^P_j\circ u_j=u_{j-1}\circ d^Q_j
  \qquad(j\leq i).
\]
Assume moreover that the following conditions hold:
\begin{enumerate}
  \item[$(\mathrm{I}_i)$]
  $Q_j\in K''$ for $j\leq i$, and
  $\mathrm{B}_j(Q_\bullet)\in K'$ for $j<i$.

  \item[$(\mathrm{II}_i)$]
  For $j<i$, the homomorphism
  \[
    \HH_j(Q_\bullet)\longrightarrow\HH_j(P_\bullet)
  \]
  induced by $(u_k)_{k\leq i}$ is an isomorphism.

  \item[$(\mathrm{III}_i)$]
  The composite
  \[
    v_i:
    \mathrm{Z}_i(Q_\bullet)
    \longrightarrow
    \mathrm{Z}_i(P_\bullet)
    \longrightarrow
    \HH_i(P_\bullet),
  \]
  where the first arrow is the restriction of $u_i$ and the second is
  canonical, is an epimorphism.
\end{enumerate}

By~(ii), $\mathrm{Z}_i(Q_\bullet)$ belongs to $K'$: it is the kernel of
the epimorphism
\[
  Q_i\longrightarrow\mathrm{B}_{i-1}(Q_\bullet),
\]
and this uses $(\mathrm{I}_i)$.
Condition~(ii) and $(\mathrm{III}_i)$ also show that
\[
  N_i=\Ker(v_i)
\]
belongs to $K'$.
By (i bis), there exist an object $Q'_{i+1}\in K''$, an epimorphism
\[
  d'_{i+1}:Q'_{i+1}\longrightarrow N_i,
\]
and a morphism $u'_{i+1}:Q'_{i+1}\to P_{i+1}$ such that
\[
  \xymatrix{
    Q'_{i+1}\ar[r]^{d'_{i+1}}\ar[d]_{u'_{i+1}} &
    N_i\ar[d]^{u_i}\\
    P_{i+1}\ar[r]_{d^P_{i+1}} &
    \mathrm{B}_i(P_\bullet)
  }
  \tag{11.9.1.3}\label{0.11.9.1.3}
\]
is commutative.

The canonical morphism
\[
  \mathrm{Z}_{i+1}(P_\bullet)
  \longrightarrow
  \HH_{i+1}(P_\bullet)
\]
is an epimorphism, and $\HH_{i+1}(P_\bullet)\in K'$.
Condition~(i) therefore supplies an object $Q''_{i+1}\in K''$ and a
morphism
\[
  u''_{i+1}:Q''_{i+1}\longrightarrow
  \mathrm{Z}_{i+1}(P_\bullet)
\]
such that the composite
\[
  Q''_{i+1}
  \longrightarrow
  \mathrm{Z}_{i+1}(P_\bullet)
  \longrightarrow
  \HH_{i+1}(P_\bullet)
\]
is an epimorphism.
If $d''_{i+1}:Q''_{i+1}\to N_i$ is the zero morphism, then
\[
  \xymatrix{
    Q''_{i+1}\ar[r]^{d''_{i+1}}\ar[d]_{u''_{i+1}} &
    N_i\ar[d]^{u_i}\\
    \mathrm{Z}_{i+1}(P_\bullet)\ar[r]_{d^P_{i+1}} &
    P_i
  }
  \tag{11.9.1.4}\label{0.11.9.1.4}
\]
is commutative; the lower horizontal arrow is zero.

\oldpage[0\textsubscript{III}]{48}
Set
\[
  Q_{i+1}=Q'_{i+1}\times Q''_{i+1},
  \qquad
  d^Q_{i+1}=d'_{i+1}+d''_{i+1},
  \qquad
  u_{i+1}=u'_{i+1}+u''_{i+1}.
\]
By~(iii), $Q_{i+1}$ belongs to $K''$.
Since
\[
  d^Q_{i+1}(Q_{i+1})
  =
  d'_{i+1}(Q'_{i+1})
  =
  N_i
  \subset
  \mathrm{Z}_i(Q_\bullet),
\]
one has
$d^Q_i\circ d^Q_{i+1}=0$ and
$\mathrm{B}_i(Q_\bullet)=N_i$; this proves
$(\mathrm{I}_{i+1})$.
The commutativity of \eqref{0.11.9.1.3} and
\eqref{0.11.9.1.4} gives
\[
  d^P_{i+1}\circ u_{i+1}
  =
  u_i\circ d^Q_{i+1}.
\]
By the definition of $N_i$, the morphism
\[
  \HH_i(Q_\bullet)
  =
  \mathrm{Z}_i(Q_\bullet)/N_i
  \longrightarrow
  \HH_i(P_\bullet)
\]
induced by the $u_k$ for $k\leq i+1$ is an isomorphism, since $v_i$ is
an epimorphism; hence $(\mathrm{II}_{i+1})$ holds.
Finally, the factor $Q''_{i+1}$ lies in
$\mathrm{Z}_{i+1}(Q_\bullet)$.
The choice of $u''_{i+1}$ shows that
\[
  v_{i+1}:
  \mathrm{Z}_{i+1}(Q_\bullet)
  \longrightarrow
  \mathrm{Z}_{i+1}(P_\bullet)
  \longrightarrow
  \HH_{i+1}(P_\bullet)
\]
is an epimorphism, because its restriction to $Q''_{i+1}$ already is
one.
Thus $(\mathrm{III}_{i+1})$ holds, the induction continues, and the
proposition follows.
\end{proof}

\begin{corollary}[11.9.2]
\label{0.11.9.2}
Let $A$ be a Noetherian ring, not necessarily commutative, and let
$P_\bullet=(P_i)_{i\in\bb{Z}}$ be a complex of right $A$-modules.
Suppose that every $\HH_i(P_\bullet)$ is a finitely generated
$A$-module and that there is an integer $d$ such that
$\HH_i(P_\bullet)=0$ for $i<d$.
Then there exist a complex $Q_\bullet=(Q_i)_{i\in\bb{Z}}$ of free
$A$-modules of finite rank, with $Q_i=0$ for $i<d$, and a homomorphism
of complexes
\[
  u:Q_\bullet\longrightarrow P_\bullet
\]
whose induced morphism
\[
  \HH_\bullet(Q_\bullet)\longrightarrow\HH_\bullet(P_\bullet)
\]
is an isomorphism.
\end{corollary}

\begin{proof}
Apply \sref{0.11.9.1} with $\cat{C}$ the category of right
$A$-modules, with $K'$ the set of finitely generated $A$-modules, and
with $K''$ the set of free $A$-modules of finite rank.
Conditions~(i), (ii), and~(iii) of \sref{0.11.9.1} follow immediately
from the hypothesis that $A$ is Noetherian.
\end{proof}

\begin{env}[11.9.3]
\label{0.11.9.3}
\emph{Remarks}.
\begin{enumerate}
  \item[{\rm(i)}]
  Under the hypotheses of \sref{0.11.9.2}, suppose in addition that the
  $P_i$ are flat right $A$-modules.
  Then, for every left $A$-module $M$, the homomorphism of complexes
  \[
    u\otimes 1:
    Q_\bullet\otimes_A M
    \longrightarrow
    P_\bullet\otimes_A M
  \]
  again induces an isomorphism
  \[
    \HH_\bullet(Q_\bullet\otimes_A M)
    \isoto
    \HH_\bullet(P_\bullet\otimes_A M),
  \]
  as will be seen in Chapter~III.

  \item[{\rm(ii)}]
  The conclusion of \sref{0.11.9.2} need not hold when $A$ is not
  Noetherian.
  Indeed, applying it to a complex having only one nonzero term would
  imply that every finitely generated right $A$-module admits a resolution by
  finitely generated free modules, which is false in general
  (Bourbaki, \emph{Alg\`ebre commutative}, chap.~I, \S2, exerc.~6).

  It is possible, however, to replace the hypothesis that $A$ is
  Noetherian by the assumption that the $\HH_i(P_\bullet)$ admit finite
  $\infty$-presentations (cf.\ Chapter~IV).
\end{enumerate}
\end{env}

\subsection{Euler--Poincar\'e characteristic of a complex of finite length modules}
\label{subsection:0.11.10}

\begin{env}[11.10.1]
\label{0.11.10.1}
Let $A$ be a ring, not necessarily commutative, and let
\[
  M^\bullet:
  0\longrightarrow M^0\longrightarrow M^1
  \longrightarrow\cdots\longrightarrow M^n\longrightarrow 0
  \tag{11.10.1.1}\label{0.11.10.1.1}
\]
be a complex of left $A$-modules of finite length.
The \emph{Euler--Poincar\'e characteristic} of this complex is the
integer
\oldpage[0\textsubscript{III}]{49}
\[
  \chi(M^\bullet)
  =
  \sum_{i=0}^{n}(-1)^i\operatorname{length}(M^i).
  \tag{11.10.1.2}\label{0.11.10.1.2}
\]
\end{env}

\begin{proposition}[11.10.2]
\label{0.11.10.2}
For every finite complex $M^\bullet$ of left $A$-modules of finite
length,
\[
  \chi(M^\bullet)=\chi\bigl(\HH^\bullet(M^\bullet)\bigr),
\]
where $\HH^\bullet(M^\bullet)$ is regarded as a complex with zero
differential.
In particular, if the sequence \eqref{0.11.10.1.1} is exact, then
$\chi(M^\bullet)=0$.
\end{proposition}

\begin{proof}
For brevity write
\[
  \mathrm{B}^i=\mathrm{B}^i(M^\bullet),
  \qquad
  \mathrm{Z}^i=\mathrm{Z}^i(M^\bullet),
  \qquad
  \HH^i=\HH^i(M^\bullet)=\mathrm{Z}^i/\mathrm{B}^i.
\]
The modules $\mathrm{B}^i$, $\mathrm{Z}^i$, and $\HH^i$ all have
finite length.
The exact sequences
\[
  0\longrightarrow\mathrm{B}^i
  \longrightarrow\mathrm{Z}^i
  \longrightarrow\HH^i
  \longrightarrow0
\]
and
\[
  0\longrightarrow\mathrm{Z}^i
  \longrightarrow M^i
  \longrightarrow\mathrm{B}^{i+1}
  \longrightarrow0
\]
give
\[
  \operatorname{length}(\mathrm{Z}^i)
  =
  \operatorname{length}(\HH^i)
  +
  \operatorname{length}(\mathrm{B}^i)
\]
and
\[
  \operatorname{length}(M^i)
  =
  \operatorname{length}(\mathrm{Z}^i)
  +
  \operatorname{length}(\mathrm{B}^{i+1}).
\]
Consequently,
\[
  \operatorname{length}(M^i)-\operatorname{length}(\HH^i)
  =
  \operatorname{length}(\mathrm{B}^{i+1})
  +
  \operatorname{length}(\mathrm{B}^i).
\]
Multiply this equality by $(-1)^i$ and sum for $0\leq i\leq n$.
Since $\mathrm{B}^0=\mathrm{B}^{n+1}=0$, the desired equality follows.
\end{proof}

\begin{corollary}[11.10.3]
\label{0.11.10.3}
Let $\mathrm{E}=(\mathrm{E}_r^{pq})$ be a spectral sequence in the
category of modules over a ring $A$.
Suppose that the $\mathrm{E}_2^{pq}$ are $A$-modules of finite length and
that $\mathrm{E}_2^{pq}\neq0$ for only finitely many pairs $(p,q)$.
Then the Euler--Poincar\'e characteristics
\[
  \chi\bigl(\mathrm{E}_r^{(\bullet)}\bigr)
\]
of all the complexes
\[
  \mathrm{E}_r^{(\bullet)}
  =
  \bigl(\mathrm{E}_r^{(n)}\bigr)_{n\in\bb{Z}}
\]
of \sref{0.11.1.1} are equal.
If, moreover, $\mathrm{E}$ is weakly convergent and one sets
\[
  \mathrm{E}_\infty^{(n)}
  =
  \bigoplus_{p+q=n}\mathrm{E}_\infty^{pq}
  \qquad(n\in\bb{Z}),
\]
then
\[
  \chi\bigl(\mathrm{E}_\infty^{(\bullet)}\bigr)
  =
  \chi\bigl(\mathrm{E}_2^{(\bullet)}\bigr),
\]
where
\[
  \mathrm{E}_\infty^{(\bullet)}
  =
  \bigl(\mathrm{E}_\infty^{(n)}\bigr)_{n\in\bb{Z}}
\]
is regarded as a complex with zero differential.
\end{corollary}

\begin{proof}
First note that if $\mathrm{E}_2^{pq}=0$, then
$\mathrm{E}_r^{pq}=0$ for $2\leq r\leq+\infty$.
Thus every complex $\mathrm{E}_r^{(\bullet)}$ is finite and consists of
$A$-modules of finite length.
For every finite $r$, the equality
\[
  \chi\bigl(\mathrm{E}_r^{(\bullet)}\bigr)
  =
  \chi\bigl(\mathrm{E}_{r+1}^{(\bullet)}\bigr)
\]
therefore follows from \sref{0.11.10.2} and the isomorphism
\[
  \HH^\bullet\bigl(\mathrm{E}_r^{(\bullet)}\bigr)
  \simeq
  \mathrm{E}_{r+1}^{(\bullet)}
\]
of complexes with zero differential.

Because the $\mathrm{E}_r^{pq}$ have finite length, for each pair
$(p,q)$ the sequences
\[
  \bigl(\mathrm{B}_k(\mathrm{E}_2^{pq})\bigr)_{k\geq2},
  \qquad
  \bigl(\mathrm{Z}_k(\mathrm{E}_2^{pq})\bigr)_{k\geq2}
\]
are stationary.
Weak convergence, together with the fact that
$\mathrm{E}_2^{pq}=0$ for all but finitely many pairs $(p,q)$, therefore
implies that some integer $r\geq2$ satisfies
\[
  \mathrm{E}_\infty^{pq}=\mathrm{E}_r^{pq}
\]
for every pair $(p,q)$.
The assertion concerning
$\chi(\mathrm{E}_\infty^{(\bullet)})$ follows.
\end{proof}

\section{Supplement on sheaf cohomology}
\label{section:0.12}

\subsection{Cohomology of sheaves of modules on ringed spaces}
\label{subsection:0.12.1}

\begin{env}[12.1.1]
\label{0.12.1.1}
\oldpage[0\textsubscript{III}]{49}
Let $(X,\sh{O}_X)$ be a ringed space.
Recall that, for every $\sh{O}_X$-Module $\sh{F}$, one defines the
cohomology $\HH^\bullet(X,\sh{F})$, which is a universal cohomological
functor (T, 2.2) from the category $\cat{C}(X)$ of
$\sh{O}_X$-Modules to the category of abelian groups; it is the derived
functor of the left exact functor
\[
  \sh{F}\longmapsto\Gamma(X,\sh{F}).
\]
The functor $\HH^\bullet(X,\sh{F})$ is isomorphic
\oldpage[0\textsubscript{III}]{50}
to the restriction to the category $\cat{C}(X)$ of the cohomological
functor defined in the same way on the category of sheaves of abelian
groups on $X$ (G, II, 7.2.1).
\end{env}

\begin{env}[12.1.2]
\label{0.12.1.2}
Set $A=\Gamma(X,\sh{O}_X)$.
Since every element of $A$ defines an endomorphism of the abelian group
$\Gamma(X,\sh{F})$, it defines by functoriality an endomorphism of the
$\delta$-functor $\HH^\bullet(X,\sh{F})$.
These endomorphisms endow each $\HH^p(X,\sh{F})$ with an
$A$-module structure, and the boundary operator $\partial$ is
$A$-linear.
Moreover, for any two nonnegative integers $p,q$ and any two
$\sh{O}_X$-Modules $\sh{F},\sh{G}$, there is a homomorphism of
$A$-modules, called the \emph{cup product},
\[
\label{0.12.1.2.1}
\tag{12.1.2.1}
  \HH^p(X,\sh{F})\otimes_A\HH^q(X,\sh{G})
  \longrightarrow
  \HH^{p+q}\bigl(X,\sh{F}\otimes_{\sh{O}_X}\sh{G}\bigr)
\]
(G, II, 6.6).
These homomorphisms make the direct sum $S$ of the
$\HH^p(X,\sh{O}_X)$, for $p\geq0$, an anticommutative graded
$A$-algebra, and the direct sum of the $\HH^p(X,\sh{F})$ a graded
$S$-module.

For every open covering $\mathfrak{U}$ of $X$, we shall always denote
by $\check{\mathrm{C}}^\bullet(\mathfrak{U},\sh{F})$ (contrary to
(G, II, 5.1)) the complex of alternating cochains of the nerve of
$\mathfrak{U}$ with values in the coefficient system
$\Gamma(U_\alpha,\sh{F})$.
It is clear that
$\check{\mathrm{C}}^\bullet(\mathfrak{U},\sh{F})$ is a graded
$A$-module, so that the cohomology groups
$\check{\HH}^p(\mathfrak{U},\sh{F})$ of this complex are endowed with
an $A$-module structure.
Moreover, the canonical maps
\[
  \check{\HH}^p(\mathfrak{U},\sh{F})
  \longrightarrow
  \HH^p(X,\sh{F})
\]
(G, II, 5.4) are $A$-linear.
\end{env}

\begin{env}[12.1.3]
\label{0.12.1.3}
Let $(X',\sh{O}_{X'})$ be a second ringed space, and let
$f=(\psi,\theta)$ be a morphism from $X'$ to $X$.
Set $A'=\Gamma(X',\sh{O}_{X'})$; the $\psi$-morphism $\theta$
canonically defines a ring homomorphism $A\to A'$.
Let $\sh{F}$ be an $\sh{O}_X$-Module and $\sh{F}'$ an
$\sh{O}_{X'}$-Module.
For every $f$-morphism $u:\sh{F}\to\sh{F}'$
\sref[0\textsubscript{I}]{0.4.4.1}, we shall see that one can define,
for every $p\geq0$, a di-homomorphism
\[
\label{0.12.1.3.1}
\tag{12.1.3.1}
  u_p:\HH^p(X,\sh{F})\longrightarrow\HH^p(X',\sh{F}').
\]

Indeed, since $\psi^*$ is exact in the category of sheaves of abelian
groups on $X$, the functor
\[
  \sh{F}\longmapsto\HH^\bullet(X',\psi^*(\sh{F}))
\]
is a $\delta$-functor on that category, and it is known that there is a
canonical homomorphism of $\delta$-functors
\[
\label{0.12.1.3.2}
\tag{12.1.3.2}
  \HH^\bullet(X,\sh{F})
  \longrightarrow
  \HH^\bullet(X',\psi^*(\sh{F}))
\]
uniquely determined by the condition that, in degree $0$, it reduces
to the canonical homomorphism
\[
  \Gamma(X,\sh{F})
  \longrightarrow
  \Gamma(X',\psi^*(\sh{F}))
\]
(T, 3.2.2).
Moreover, every element of $A$ determines an endomorphism $\mu$ of
$\Gamma(X,\sh{F})$ and an endomorphism $\mu'$ of
$\Gamma(X',\psi^*(\sh{F}))$ such that the diagram
\[
\label{0.12.1.3.3}
\tag{12.1.3.3}
  \xymatrix{
    \Gamma(X,\sh{F})\ar[r]\ar[d]_\mu
      & \Gamma(X',\psi^*(\sh{F}))\ar[d]^{\mu'}\\
    \Gamma(X,\sh{F})\ar[r]
      & \Gamma(X',\psi^*(\sh{F}))
  }
\]
is commutative.
\oldpage[0\textsubscript{III}]{51}
By the uniqueness property for extending morphisms of universal
cohomological functors (T, 2.2), one obtains unique extensions of
$\mu$ and $\mu'$ to cohomology which make the diagrams analogous to
\sref{0.12.1.3.3} commutative.
Thus \sref{0.12.1.3.2} is a homomorphism of $A$-modules.

Now note that
\[
  f^*(\sh{F})
  =
  \psi^*(\sh{F})
  \otimes_{\psi^*(\sh{O}_X)}
  \sh{O}_{X'},
\]
and hence that there is a canonical di-homomorphism
\[
  \psi^*(\sh{F})\longrightarrow f^*(\sh{F})
\]
from the $\psi^*(\sh{O}_X)$-Module $\psi^*(\sh{F})$ to the
$\sh{O}_{X'}$-Module $f^*(\sh{F})$.
By functoriality, one therefore obtains a functorial di-homomorphism
\[
\label{0.12.1.3.4}
\tag{12.1.3.4}
  \HH^p(X',\psi^*(\sh{F}))
  \longrightarrow
  \HH^p(X',f^*(\sh{F})),
\]
the corresponding rings being $A$ and $A'$.
Composing this di-homomorphism with \sref{0.12.1.3.2}, one obtains a
canonical di-homomorphism, functorial in $\sh{F}$,
\[
\label{0.12.1.3.5}
\tag{12.1.3.5}
  \theta_p:\HH^p(X,\sh{F})
  \longrightarrow
  \HH^p(X',f^*(\sh{F})).
\]

Finally, by functoriality, the homomorphism
$u^\sharp:f^*(\sh{F})\to\sh{F}'$ gives a homomorphism of
$A'$-modules
\[
  \HH^p(X',f^*(\sh{F}))
  \longrightarrow
  \HH^p(X',\sh{F}'),
\]
which, composed with \sref{0.12.1.3.5}, gives
\sref{0.12.1.3.1}.

Let $f'=(\psi',\theta'):X''\to X'$ be a second morphism of ringed
spaces, and let $f''=f\circ f'$ be the composite morphism.
Taking account of the compatibility of the functor $\psi'^*$ with
tensor products \sref[0\textsubscript{I}]{0.4.3.3}, one immediately
checks that the composite of the di-homomorphism
\[
  \HH^p(X',f^*(\sh{F}))
  \longrightarrow
  \HH^p\bigl(X'',f'^*(f^*(\sh{F}))\bigr)
\]
with \sref{0.12.1.3.5} is the corresponding di-homomorphism
\[
  \HH^p(X,\sh{F})
  \longrightarrow
  \HH^p(X'',f''^*(\sh{F})).
\]
\end{env}

\begin{env}[12.1.4]
\label{0.12.1.4}
A direct definition of the homomorphism \sref{0.12.1.3.2} can be
given as follows.
Take an injective resolution
$\sh{L}^\bullet=(\sh{L}^i)$ of $\sh{F}$ consisting of sheaves of
abelian groups on $X$.
Since the functor $\psi^*$ is exact,
$\psi^*(\sh{L}^\bullet)$ is a resolution of $\psi^*(\sh{F})$
consisting of sheaves on $X'$.
If $\sh{L}'^\bullet=(\sh{L}'^i)$ is an injective resolution of
$\psi^*(\sh{F})$ in the category of sheaves of abelian groups on
$X'$, there is a morphism
\[
  \psi^*(\sh{L}^\bullet)\longrightarrow\sh{L}'^\bullet
\]
of complexes of sheaves of abelian groups, compatible with the
augmentations (M, V, 1.1.a), and determined up to homotopy.
One thereby obtains homomorphisms
\[
  \Gamma(X,\sh{L}^\bullet)
  \longrightarrow
  \Gamma(X',\psi^*(\sh{L}^\bullet))
  \longrightarrow
  \Gamma(X',\sh{L}'^\bullet)
\]
of complexes of abelian groups.
Their composite, on passing to cohomology, gives a morphism of
$\delta$-functors
\[
  \HH^\bullet(X,\sh{F})
  \longrightarrow
  \HH^\bullet(X',\psi^*(\sh{F}));
\]
since it coincides with \sref{0.12.1.3.2} in degree $0$, it is
identical to it (T, 2.2).

Now consider an open covering
$\mathfrak{U}=(U_\alpha)$ of $X$, and let
$\mathfrak{U}'=(f^{-1}(U_\alpha))$ be its inverse-image open covering
of $X'$.
For every open subset $V$ of $X$, the canonical homomorphisms
\[
  \Gamma(V,\sh{F})
  \longrightarrow
  \Gamma(f^{-1}(V),f^*(\sh{F}))
\]
immediately define (cf.\ G, II, 5.1) a homomorphism of complexes
\[
  \check{\mathrm{C}}^\bullet(\mathfrak{U},\sh{F})
  \longrightarrow
  \check{\mathrm{C}}^\bullet(\mathfrak{U}',f^*(\sh{F})),
\]
and hence canonical homomorphisms
\[
\label{0.12.1.4.1}
\tag{12.1.4.1}
  \theta_p:
  \check{\HH}^p(\mathfrak{U},\sh{F})
  \longrightarrow
  \check{\HH}^p(\mathfrak{U}',f^*(\sh{F})).
\]

\oldpage[0\textsubscript{III}]{52}
Moreover, there are commutative diagrams
\[
\label{0.12.1.4.2}
\tag{12.1.4.2}
  \xymatrix{
    \check{\HH}^p(\mathfrak{U},\sh{F})
      \ar[r]^-{\theta_p}\ar[d]
      & \check{\HH}^p(\mathfrak{U}',f^*(\sh{F}))\ar[d]\\
    \HH^p(X,\sh{F})
      \ar[r]_-{\theta_p}
      & \HH^p(X',f^*(\sh{F})),
  }
\]
where the vertical arrows are the canonical homomorphisms of
(G, II, 5.2).
To prove the commutativity of \sref{0.12.1.4.2}, consider the complex
of sheaves of alternating cochains of $\sh{F}$ relative to
$\mathfrak{U}$,
$\sh{C}^\bullet(\mathfrak{U},\sh{F})$, for which
\[
  \Gamma\bigl(X,\sh{C}^\bullet(\mathfrak{U},\sh{F})\bigr)
  =
  \check{\mathrm{C}}^\bullet(\mathfrak{U},\sh{F})
\]
(G, II, 5.2).
The canonical homomorphisms
\[
  \Gamma(V,\sh{F})
  \longrightarrow
  \Gamma(f^{-1}(V),\psi^*(\sh{F}))
\]
then define a $\psi$-morphism
\[
  \sh{C}^\bullet(\mathfrak{U},\sh{F})
  \longrightarrow
  \sh{C}^\bullet(\mathfrak{U}',\psi^*(\sh{F})),
\]
and, with the notation above, there is a commutative diagram
\[
  \xymatrix{
    \Gamma\bigl(X,\sh{C}^\bullet(\mathfrak{U},\sh{F})\bigr)
      \ar[r]\ar[d]
      & \Gamma\bigl(X',
          \sh{C}^\bullet(\mathfrak{U}',\psi^*(\sh{F}))\bigr)\ar[d]\\
    \Gamma(X,\sh{L}^\bullet)\ar[r]
      & \Gamma(X',\psi^*(\sh{L}^\bullet))\ar[r]
      & \Gamma(X',\sh{L}'^\bullet).
  }
\]
Passing to cohomology gives commutative diagrams
\[
  \xymatrix{
    \check{\HH}^p(\mathfrak{U},\sh{F})\ar[r]\ar[d]
      & \check{\HH}^p(\mathfrak{U}',\psi^*(\sh{F}))\ar[d]\\
    \HH^p(X,\sh{F})\ar[r]
      & \HH^p(X',\psi^*(\sh{F})),
  }
\]
where the vertical arrows are the canonical homomorphisms of
(G, II, 5.2).
It remains only to combine these diagrams with the commutative diagrams
\[
  \xymatrix{
    \check{\HH}^p(\mathfrak{U}',\psi^*(\sh{F}))\ar[r]\ar[d]
      & \check{\HH}^p(\mathfrak{U}',f^*(\sh{F}))\ar[d]\\
    \HH^p(X',\psi^*(\sh{F}))\ar[r]
      & \HH^p(X',f^*(\sh{F})),
  }
\]
\oldpage[0\textsubscript{III}]{53}
which arise from the homomorphism
$\psi^*(\sh{F})\to f^*(\sh{F})$ and the functoriality of the
canonical homomorphisms of (G, II, 5.2), to obtain the commutativity
of \sref{0.12.1.4.2}.

Note that if
$A=\Gamma(X,\sh{O}_X)$ and $A'=\Gamma(X',\sh{O}_{X'})$, then
\sref{0.12.1.4.1} is a di-homomorphism of modules corresponding to
the rings $A$ and $A'$.
There is a transitivity property for \sref{0.12.1.4.1} with respect
to a composite of two morphisms, analogous to that of
\sref{0.12.1.3.5}.
Finally, in the preceding definitions, instead of an injective
resolution $\sh{L}^\bullet$ of $\sh{F}$, one could equally well have
started with a resolution such that
\[
  \HH^p(X,\sh{L}^i)=0
  \qquad\text{for every $i$ and every $p>0$}
\]
(G, II, 4.7.1).
\end{env}

\begin{env}[12.1.5]
\label{0.12.1.5}
Let $\sh{F},\sh{G},\sh{H}$ be three $\sh{O}_X$-Modules, and consider
an $\sh{O}_X$-homomorphism
\[
  u:\sh{F}\otimes_{\sh{O}_X}\sh{G}\longrightarrow\sh{H},
\]
which gives, in cohomology, the homomorphisms
\[
\label{0.12.1.5.1}
\tag{12.1.5.1}
  \HH^p(X,\sh{F})\otimes_A\HH^q(X,\sh{G})
  \longrightarrow
  \HH^{p+q}(X,\sh{H})
\]
deduced from the cup product \sref{0.12.1.2.1}.
We show that, under the hypotheses and with the notation of
\sref{0.12.1.3}, the diagrams
\[
\label{0.12.1.5.2}
\tag{12.1.5.2}
  \xymatrix@C=10pt{
    \HH^p(X,\sh{F})\otimes_A\HH^q(X,\sh{G})
      \ar[r]\ar[d]
      & \HH^{p+q}(X,\sh{H})\ar[d]\\
    \HH^p(X',f^*(\sh{F}))\otimes_{A'}
      \HH^q(X',f^*(\sh{G}))
      \ar[r]
      & \HH^{p+q}(X',f^*(\sh{H}))
  }
\]
are commutative, where the vertical arrows come from the canonical
homomorphisms \sref{0.12.1.3.5}.

Recall for this purpose that \sref{0.12.1.5.1} can be obtained from
the canonical resolutions (G, II, 4.3)
$\sh{L}^\bullet,\sh{M}^\bullet,\sh{N}^\bullet$ of
$\sh{F},\sh{G},\sh{H}$ respectively, which consist of
$\sh{O}_X$-Modules, and from the linear map
\[
  \sh{L}^\bullet\otimes_{\sh{O}_X}\sh{M}^\bullet
  \longrightarrow
  \sh{N}^\bullet
\]
of complexes of $\sh{O}_X$-Modules corresponding to $u$.
This map provides a homomorphism of complexes of $A$-modules
\[
  \Gamma(X,\sh{L}^\bullet)\otimes_A\Gamma(X,\sh{M}^\bullet)
  \longrightarrow
  \Gamma(X,\sh{N}^\bullet),
\]
and hence, on passing to cohomology, homomorphisms
\[
  \HH^p\bigl(\Gamma(X,\sh{L}^\bullet)\bigr)
  \otimes_A
  \HH^q\bigl(\Gamma(X,\sh{M}^\bullet)\bigr)
  \longrightarrow
  \HH^{p+q}\bigl(\Gamma(X,\sh{N}^\bullet)\bigr)
\]
(G, II, 6.6).
There is clearly a commutative diagram
\[
\label{0.12.1.5.3}
\tag{12.1.5.3}
  \xymatrix@C=8pt{
    \Gamma(X,\sh{L}^\bullet)\otimes_A
      \Gamma(X,\sh{M}^\bullet)
      \ar[r]\ar[d]
      & \Gamma(X,\sh{N}^\bullet)\ar[d]\\
    \Gamma(X',\psi^*(\sh{L}^\bullet))
      \otimes_{\Gamma(X',\psi^*(\sh{O}_X))}
      \Gamma(X',\psi^*(\sh{M}^\bullet))
      \ar[r]
      & \Gamma(X',\psi^*(\sh{N}^\bullet)).
  }
\]
\oldpage[0\textsubscript{III}]{54}
Passing to cohomology gives commutative diagrams
\[
\label{0.12.1.5.4}
\tag{12.1.5.4}
  \xymatrix@C=7pt{
    \HH^p(X,\sh{F})\otimes_A\HH^q(X,\sh{G})
      \ar[r]\ar[d]
      & \HH^{p+q}(X,\sh{H})\ar[d]\\
    \HH^p\bigl(\Gamma(X',\psi^*(\sh{L}^\bullet))\bigr)
      \otimes_{\Gamma(X',\psi^*(\sh{O}_X))}
      \HH^q\bigl(\Gamma(X',\psi^*(\sh{M}^\bullet))\bigr)
      \ar[r]
      & \HH^{p+q}\bigl(\Gamma(X',\psi^*(\sh{N}^\bullet))\bigr).
  }
\]
But since
$\psi^*(\sh{L}^\bullet)$,
$\psi^*(\sh{M}^\bullet)$, and
$\psi^*(\sh{N}^\bullet)$
are resolutions of
$\psi^*(\sh{F})$,
$\psi^*(\sh{G})$, and
$\psi^*(\sh{H})$ respectively, there is a commutative diagram
(G, II, 6.6.1)
\[
\label{0.12.1.5.5}
\tag{12.1.5.5}
  \xymatrix@C=7pt{
    \HH^p\bigl(\Gamma(X',\psi^*(\sh{L}^\bullet))\bigr)
      \otimes_{\Gamma(X',\psi^*(\sh{O}_X))}
      \HH^q\bigl(\Gamma(X',\psi^*(\sh{M}^\bullet))\bigr)
      \ar[r]\ar[d]
      & \HH^{p+q}\bigl(\Gamma(X',\psi^*(\sh{N}^\bullet))\bigr)
        \ar[d]\\
    \HH^p(X',\psi^*(\sh{F}))
      \otimes_{\Gamma(X',\psi^*(\sh{O}_X))}
      \HH^q(X',\psi^*(\sh{G}))
      \ar[r]
      & \HH^{p+q}(X',\psi^*(\sh{H})).
  }
\]
Finally, by functoriality, there is a commutative diagram
\[
\label{0.12.1.5.6}
\tag{12.1.5.6}
  \xymatrix@C=7pt{
    \HH^p(X',\psi^*(\sh{F}))
      \otimes_{\Gamma(X',\psi^*(\sh{O}_X))}
      \HH^q(X',\psi^*(\sh{G}))
      \ar[r]\ar[d]
      & \HH^{p+q}(X',\psi^*(\sh{H}))\ar[d]\\
    \HH^p(X',f^*(\sh{F}))\otimes_{A'}
      \HH^q(X',f^*(\sh{G}))
      \ar[r]
      & \HH^{p+q}(X',f^*(\sh{H})).
  }
\]
Combining the three diagrams \sref{0.12.1.5.4},
\sref{0.12.1.5.5}, and \sref{0.12.1.5.6} yields the desired
commutative diagram \sref{0.12.1.5.2}.
\end{env}

\begin{remark}[12.1.6]
\label{0.12.1.6}
\oldpage[0\textsubscript{III}]{55}
With the notation of \sref{0.12.1.3}, suppose that there is a
commutative diagram
\[
\label{0.12.1.6.1}
\tag{12.1.6.1}
  \xymatrix{
    0\ar[r] & \sh{F}\ar[r]^r\ar[d]_u
      & \sh{G}\ar[r]^s\ar[d]_v
      & \sh{H}\ar[r]\ar[d]_w & 0\\
    0\ar[r] & \sh{F}'\ar[r]_{r'}
      & \sh{G}'\ar[r]_{s'}
      & \sh{H}'\ar[r] & 0,
  }
\]
where $r,s$ are homomorphisms of $\sh{O}_X$-Modules,
$r',s'$ are homomorphisms of $\sh{O}_{X'}$-Modules,
$u,v,w$ are $f$-morphisms, and the rows are exact.
It follows that the diagram
\[
\label{0.12.1.6.2}
\tag{12.1.6.2}
  \xymatrix@C=6pt{
    \cdots\ar[r]
      & \HH^p(X,\sh{F})\ar[r]\ar[d]_{u_p}
      & \HH^p(X,\sh{G})\ar[r]\ar[d]_{v_p}
      & \HH^p(X,\sh{H})\ar[r]^-\partial\ar[d]_{w_p}
      & \HH^{p+1}(X,\sh{F})\ar[r]\ar[d]_{u_{p+1}}
      & \cdots\\
    \cdots\ar[r]
      & \HH^p(X',\sh{F}')\ar[r]
      & \HH^p(X',\sh{G}')\ar[r]
      & \HH^p(X',\sh{H}')\ar[r]_-\partial
      & \HH^{p+1}(X',\sh{F}')\ar[r]
      & \cdots
  }
\]
is commutative.
Indeed, \sref{0.12.1.6.1} factors as
\[
  \xymatrix@R=14pt{
    0\ar[r] & \sh{F}\ar[r]\ar[d]
      & \sh{G}\ar[r]\ar[d]
      & \sh{H}\ar[r]\ar[d] & 0\\
    0\ar[r] & \psi^*(\sh{F})\ar[r]\ar[d]
      & \psi^*(\sh{G})\ar[r]\ar[d]
      & \psi^*(\sh{H})\ar[r]\ar[d] & 0\\
    0\ar[r] & \sh{F}'\ar[r]
      & \sh{G}'\ar[r]
      & \sh{H}'\ar[r] & 0,
  }
\]
where the middle row is exact
\sref[0\textsubscript{I}]{0.3.7.2}.
It is then enough to use the fact that \sref{0.12.1.3.2} is a
homomorphism of $\delta$-functors and that the
$\HH^p(X',\sh{F}')$ form a $\delta$-functor in $\sh{F}'$.
\end{remark}

\begin{env}[12.1.7]
\label{0.12.1.7}
With the hypotheses and notation of \sref{0.12.1.3}, consider now the
case
\[
  \sh{F}=f_*(\sh{F}')=\psi_*(\sh{F}').
\]
We shall see that the di-homomorphism defined in \sref{0.12.1.3},
\[
  \HH^p(X,f_*(\sh{F}'))\longrightarrow\HH^p(X',\sh{F}'),
  \tag{12.1.7.1}
  \label{0.12.1.7.1}
\]
can be obtained (up to an automorphism of $\HH^p(X',\sh{F}')$) as an
edge homomorphism of a spectral sequence for the composite functor
\[
  \sh{F}'\leadsto\Gamma(X,\psi_*(\sh{F}'))
\]
(T, 2.4).
The description of \sref{0.12.1.7.1} given in \sref{0.12.1.4} shows
here that this homomorphism can be obtained as follows.
Take injective resolutions $\sh{L}^\bullet$ and
${\sh{L}'}^\bullet$ of $\psi_*(\sh{F}')$ and $\sh{F}'$,
respectively, and then a homomorphism of complexes
\[
  v:\psi^*(\sh{L}^\bullet)\longrightarrow{\sh{L}'}^\bullet
\]
``above'' the canonical homomorphism
$\psi^*(\psi_*(\sh{F}'))\to\sh{F}'$.
\oldpage[0\textsubscript{III}]{56}
We next note that
\[
  \Gamma(X',{\sh{L}'}^\bullet)
  =\Gamma(X,\psi_*({\sh{L}'}^\bullet))
\]
and that the composite homomorphism
\[
  \Gamma(X,\sh{L}^\bullet)
  \longrightarrow\Gamma(X',\psi^*(\sh{L}^\bullet))
  \xrightarrow{\Gamma(v)}\Gamma(X',{\sh{L}'}^\bullet)
\]
is none other than
\[
  \Gamma(v^\flat):
  \Gamma(X,\sh{L}^\bullet)
  \longrightarrow\Gamma(X,\psi_*({\sh{L}'}^\bullet))
  \tag{12.1.7.2}
  \label{0.12.1.7.2}
\]
\sref[0\textsubscript{I}]{0.3.7.1}; and \sref{0.12.1.7.1} is
obtained by passing to cohomology in \sref{0.12.1.7.2}.

On the other hand, the spectral sequences for the composite functor
\[
  \sh{F}'\leadsto\Gamma(X,\psi_*(\sh{F}'))
\]
are obtained by taking an injective Cartan--Eilenberg resolution
\[
  \sh{M}^{\bullet\bullet}=(\sh{M}^{ij})
\]
of the complex $\psi_*({\sh{L}'}^\bullet)$ in the category of
sheaves of abelian groups on $X$.
The spectral sequences in question are those of the bicomplex
$\Gamma(X,\sh{M}^{\bullet\bullet})$; they are biregular, since
$\sh{M}^{ij}=0$ for $i<0$ or $j<0$.
Now the first spectral sequence of this bicomplex degenerates, since
the sheaves $\psi_*({\sh{L}'}^i)$ are flasque (G, II, 3.1.1), and
therefore
\[
  \HH_\mathrm{II}^q
    \bigl(\Gamma(X,\sh{M}^{i,\bullet})\bigr)
  =\HH^q\bigl(\psi_*({\sh{L}'}^i)\bigr)=0
  \qquad(q>0)
\]
(G, II, 4.4.3).
Consequently, there are bijective edge homomorphisms
\sref{0.11.1.6}
\[
  {}'E_2^{i0}
  =\HH^i\!\left(
      \HH_\mathrm{II}^0
        \bigl(\Gamma(X,\sh{M}^{\bullet\bullet})\bigr)
    \right)
  \longrightarrow
  \HH^i\bigl(\Gamma(X,\sh{M}^{\bullet\bullet})\bigr),
  \tag{12.1.7.3}
  \label{0.12.1.7.3}
\]
and we know \sref{0.11.3.4} that this homomorphism comes, on passing
to cohomology, from the augmentation
\[
  \Gamma(X,\psi_*({\sh{L}'}^\bullet))
  \longrightarrow\Gamma(X,\sh{M}^{\bullet\bullet}),
  \tag{12.1.7.4}
  \label{0.12.1.7.4}
\]
which itself comes from the augmentation
\[
  \eta:\psi_*({\sh{L}'}^\bullet)
  \longrightarrow\sh{M}^{\bullet\bullet}.
\]

For the second spectral sequence, on the other hand, there are edge
homomorphisms
\[
  {}''E_2^{i0}
  =\HH^i\!\left(
      \HH_\mathrm{I}^0
        \bigl(\Gamma(X,\sh{M}^{\bullet\bullet})\bigr)
    \right)
  \longrightarrow
  \HH^i\bigl(\Gamma(X,\sh{M}^{\bullet\bullet})\bigr),
  \tag{12.1.7.5}
  \label{0.12.1.7.5}
\]
which come \sref{0.11.3.4}, on passing to cohomology, from the
homomorphism of complexes
\[
  Z_\mathrm{I}^0
    \bigl(\Gamma(X,\sh{M}^{\bullet\bullet})\bigr)
  \longrightarrow\Gamma(X,\sh{M}^{\bullet\bullet}).
\]
Since $\psi_*$ is left exact, the sequence
\[
  0\longrightarrow\psi_*(\sh{F}')
  \longrightarrow\psi_*({\sh{L}'}^0)
  \longrightarrow\psi_*({\sh{L}'}^1)
\]
is exact.
By the definition of a Cartan--Eilenberg resolution
\sref{0.11.4.2}, we may therefore take
\[
  B_\mathrm{I}^0(\sh{M}^{\bullet\bullet})=0,
  \qquad
  Z_\mathrm{I}^0(\sh{M}^{\bullet\bullet})=\sh{L}^\bullet.
\]
Since the diagram
\[
  \xymatrix@C=50pt@R=24pt{
    \sh{L}^0\ar[r]^{i^0}
      & \sh{M}^{00}\\
    \psi_*(\sh{F}')\ar[u]^\varepsilon
      \ar[r]_{\varepsilon'}\ar[ur]^{\varepsilon''}
      & \psi_*({\sh{L}'}^0)\ar[u]_{\eta^0}
  }
\]
is commutative, the injection of complexes
$i:\sh{L}^\bullet\to\sh{M}^{\bullet\bullet}$ is compatible with the
augmentations $\varepsilon$ and $\varepsilon''$.
Thus there are two homomorphisms of complexes from $\sh{L}^\bullet$
to $\sh{M}^{\bullet\bullet}$,
\[
  \xymatrix@C=52pt@R=22pt{
    \sh{L}^\bullet\ar[r]^i\ar[dr]_{v^\flat}
      & \sh{M}^{\bullet\bullet}\\
    & \psi_*({\sh{L}'}^\bullet)\ar[u]_\eta
  }
\]
\oldpage[0\textsubscript{III}]{57}
compatible with the augmentations $\varepsilon$ and
$\varepsilon''$.
Since $\sh{L}^\bullet$ is an injective resolution and
$\sh{M}^{\bullet\bullet}$ consists of injective sheaves, it follows
from (M, V, 1.1(a)) that these two homomorphisms are homotopic.
The same is therefore true of the two corresponding homomorphisms
\[
  \Gamma(X,\sh{L}^\bullet)
  \longrightarrow\Gamma(X,\sh{M}^{\bullet\bullet}),
\]
and, on passing to cohomology, they consequently give the same
homomorphism.
In other words, we have shown that the edge homomorphism
\sref{0.12.1.7.5}, written
\[
  \HH^p(X,\psi_*(\sh{F}'))
  \longrightarrow\HH^p\bigl(\Gamma(X,\sh{M}^{\bullet\bullet})\bigr),
\]
is the composite of \sref{0.12.1.7.1} and
\sref{0.12.1.7.3}, the latter being written
\[
  \HH^p(X',\sh{F}')
  \longrightarrow\HH^p\bigl(\Gamma(X,\sh{M}^{\bullet\bullet})\bigr)
\]
and having been seen to be an isomorphism.
This proves the assertion.
\end{env}

\subsection{Higher direct images}
\label{subsection:0.12.2}

\begin{env}[12.2.1]
\label{0.12.2.1}
Let $(X,\sh{O}_X)$ and $(Y,\sh{O}_Y)$ be two ringed spaces, and let
$f=(\psi,\theta)$ be a morphism from $X$ to $Y$.  It defines the direct
image functor
\[
  f_*:\cat{C}(X)\longrightarrow\cat{C}(Y),
\]
which is moreover identical to the restriction to $\cat{C}(X)$ of the
functor $\psi_*$ defined on the category of sheaves of abelian groups
on $X$.
The latter functor is additive and left exact; and since every sheaf of
abelian groups on $X$ is isomorphic to a subsheaf of an injective sheaf
of abelian groups, one defines the right derived functors
\[
  \sh{F}\longmapsto\RR^p\psi_*(\sh{F})
\]
of $\psi_*$.  The $\RR^p\psi_*(\sh{F})$ are sheaves of abelian groups
on $Y$, and the $\RR^p\psi_*$ form a universal cohomological functor
(T, 2.3).

Moreover, $\RR^p\psi_*(\sh{F})$ is the sheaf associated with the
presheaf
\[
  V\longmapsto\HH^p(f^{-1}(V),\sh{F})
\]
(T, 3.7.2).
Suppose now that $\sh{F}$ is an $\sh{O}_X$-module.
Then $\HH^p(f^{-1}(V),\sh{F})$ is naturally endowed with the structure
of a $\Gamma(f^{-1}(V),\sh{O}_X)$-module, hence with that of a
$\Gamma(V,\psi_*(\sh{O}_X))$-module; the homomorphism
\[
  \theta:\sh{O}_Y\longrightarrow\psi_*(\sh{O}_X)
\]
then endows it with the structure of a
$\Gamma(V,\sh{O}_Y)$-module.
For the structures so defined, restriction from an open subset $V$ to
an open subset $V'\subset V$ defines a di-homomorphism.  This therefore
defines on every $\RR^p\psi_*(\sh{F})$ the structure of an
$\sh{O}_Y$-module.  We denote this $\sh{O}_Y$-module by
$\RR^p f_*(\sh{F})$; thus $\RR^p f_*$ is defined as an additive
functor from $\cat{C}(X)$ to $\cat{C}(Y)$.

Moreover, the $\RR^p f_*$ form a $\delta$-functor.  Indeed, if
\[
  0\longrightarrow\sh{F}'\longrightarrow\sh{F}
  \longrightarrow\sh{F}''\longrightarrow0
\]
is an exact sequence of $\sh{O}_X$-modules, the preceding description
of the $\RR^p\psi_*$ and of the $\sh{O}_Y$-module structure on
$\RR^p\psi_*(\sh{F})$ immediately shows that the boundary homomorphism
\[
  \partial:\RR^p\psi_*(\sh{F}'')
  \longrightarrow\RR^{p+1}\psi_*(\sh{F}')
\]
is an $\sh{O}_Y$-module homomorphism.
Finally, the $\RR^p f_*$ are identified with the right derived functors
of $f_*$.  Every $\sh{O}_X$-module admits an injective resolution by
$\sh{O}_X$-modules, and such a resolution consists of flasque sheaves
of abelian groups (G, II, 7.1); it can therefore be used to compute the
$\RR^p\psi_*(\sh{F})$, since $\RR^n\psi_*(\sh{G})=0$ for $n\geq1$
and every flasque sheaf $\sh{G}$ (T, 2.4.1, Remark~3, and the
corollary of Proposition~\hyperref[0.3.3.2]{3.3.2}).
It follows that the $\RR^p f_*$ form a universal cohomological functor
from $\cat{C}(X)$ to $\cat{C}(Y)$ (T, 2.3).
\end{env}

\begin{env}[12.2.2]
\label{0.12.2.2}
Let $\sh{F}$ and $\sh{G}$ be two $\sh{O}_X$-modules.
With the notation of \sref{0.12.2.1}, for every open subset $V$ of
$Y$ there is the cup-product homomorphism \sref{0.12.1.2.1}
\[
  \begin{aligned}
  &\HH^p(f^{-1}(V),\sh{F})
    \otimes_{\Gamma(f^{-1}(V),\sh{O}_X)}
    \HH^q(f^{-1}(V),\sh{G})\\
  &\hspace{55mm}\longrightarrow
    \HH^{p+q}\bigl(f^{-1}(V),
      \sh{F}\otimes_{\sh{O}_X}\sh{G}\bigr).
  \end{aligned}
\]
\oldpage[0\textsubscript{III}]{58}
It follows immediately from the definition of the cup product
(G, II, 6.6) that these homomorphisms commute with restriction from
$V$ to an open subspace $V'$ of $V$.
On the other hand, $\theta$ gives a ring homomorphism
\[
  \Gamma(V,\sh{O}_Y)
  \longrightarrow\Gamma(V,\psi_*(\sh{O}_X))
  =\Gamma(f^{-1}(V),\sh{O}_X),
\]
and hence a canonical homomorphism of tensor products
\[
  \begin{aligned}
  &\HH^p(f^{-1}(V),\sh{F})
    \otimes_{\Gamma(V,\sh{O}_Y)}
    \HH^q(f^{-1}(V),\sh{G})\\
  &\quad\longrightarrow
    \HH^p(f^{-1}(V),\sh{F})
    \otimes_{\Gamma(f^{-1}(V),\sh{O}_X)}
    \HH^q(f^{-1}(V),\sh{G}),
  \end{aligned}
\]
which is likewise compatible with restriction from $V$ to $V'$.
By composition, one thus obtains a homomorphism of
$\Gamma(V,\sh{O}_Y)$-modules.  For the sheaves associated with the
presheaves in question, it defines a canonical homomorphism,
functorial in $\sh{F}$ and $\sh{G}$,
\[
  \RR^p f_*(\sh{F})\otimes_{\sh{O}_Y}\RR^q f_*(\sh{G})
  \longrightarrow
  \RR^{p+q}f_*
    \bigl(\sh{F}\otimes_{\sh{O}_X}\sh{G}\bigr).
  \tag{12.2.2.1}
  \label{0.12.2.2.1}
\]
For $p=q=0$, this homomorphism reduces to
\sref[0\textsubscript{I}]{0.4.2.2.1}.
\end{env}

\begin{proposition}[12.2.3]
\label{0.12.2.3}
For every $\sh{O}_X$-module $\sh{F}$ and every locally free
$\sh{O}_Y$-module $\sh{L}$ of finite rank, there are canonical
functorial isomorphisms
\[
  \RR^p f_*(\sh{F})\otimes_{\sh{O}_Y}\sh{L}
  \simto
  \RR^p f_*
    \bigl(\sh{F}\otimes_{\sh{O}_X}f^*(\sh{L})\bigr).
  \tag{12.2.3.1}
  \label{0.12.2.3.1}
\]
\end{proposition}

\begin{proof}
The homomorphism \sref{0.12.2.3.1} is obtained by composing the
following particular case of \sref{0.12.2.2.1},
\[
  \RR^p f_*(\sh{F})
    \otimes_{\sh{O}_Y}f_*\bigl(f^*(\sh{L})\bigr)
  \longrightarrow
  \RR^p f_*
    \bigl(\sh{F}\otimes_{\sh{O}_X}f^*(\sh{L})\bigr),
  \tag{12.2.3.2}
  \label{0.12.2.3.2}
\]
with the homomorphism from the first member of
\sref{0.12.2.3.1} to that of \sref{0.12.2.3.2} induced by the
canonical homomorphism \sref[0\textsubscript{I}]{0.4.4.3.2}.
To verify that \sref{0.12.2.3.1} is an isomorphism when $\sh{L}$ is
locally free, one may immediately reduce to the case
$\sh{L}=\sh{O}_Y$, since the question is local on $Y$ and the
functors in question are additive in $\sh{L}$.
In that case, in view of the definition of \sref{0.12.2.2.1}, the
proposition reduces to verifying that the corresponding homomorphism
of presheaves is bijective; this follows immediately from
$f^*(\sh{O}_Y)=\sh{O}_X$.
\end{proof}

\begin{env}[12.2.4]
\label{0.12.2.4}
Let $(Z,\sh{O}_Z)$ be a third ringed space, and let
$g:Y\to Z$ be a morphism of ringed spaces.
We know (G, II, 7.1 and 3.1.1) that, for every injective
$\sh{O}_X$-module $\sh{G}$, $f_*(\sh{G})$ is a flasque sheaf of
abelian groups.  Consequently, by \sref{0.12.2.1},
\[
  \RR^p g_*\bigl(f_*(\sh{G})\bigr)=0
  \qquad(p>0).
\]
It follows (T, 2.4.1) that the Leray spectral sequence for composite
functors applies to the composite $g_*\circ f_*$.  There is a
biregular spectral sequence whose abutment is the functor
$\RR^\bullet h_*$, where $h=g\circ f$, and whose $E_2$-term is
\[
  E_2^{pq}=\RR^p g_*\bigl(\RR^q f_*(\sh{F})\bigr).
  \tag{12.2.4.1}
  \label{0.12.2.4.1}
\]
\end{env}

\begin{env}[12.2.5]
\label{0.12.2.5}
Under the conditions of \sref{0.12.2.4}, we shall directly define
canonical homomorphisms of $\sh{O}_Z$-modules
\[
  \RR^n g_*\bigl(f_*(\sh{F})\bigr)
  \longrightarrow\RR^n h_*(\sh{F})
  \tag{12.2.5.1}
  \label{0.12.2.5.1}
\]
and
\[
  \RR^n h_*(\sh{F})
  \longrightarrow g_*\bigl(\RR^n f_*(\sh{F})\bigr),
  \tag{12.2.5.2}
  \label{0.12.2.5.2}
\]
\oldpage[0\textsubscript{III}]{59}
which may be identified with the ``edge homomorphisms'' of the Leray
spectral sequence (cf.\ \sref{0.12.1.7}).
It is enough to work with the presheaves with which the higher direct
image sheaves are associated \sref{0.12.2.1}.
For this purpose, let $W$ be any open subset of $Z$, and consider its
inverse image $g^{-1}(W)$ in $Y$.  There is a canonical
di-homomorphism
\[
  \begin{aligned}
  \HH^n\bigl(g^{-1}(W),f_*(\sh{F})\bigr)
  &\longrightarrow
  \HH^n\bigl(f^{-1}(g^{-1}(W)),
    f^*(f_*(\sh{F}))\bigr),
  \end{aligned}
  \tag{12.2.5.3}
  \label{0.12.2.5.3}
\]
the corresponding rings being
$\Gamma(g^{-1}(W),\sh{O}_Y)$ and
$\Gamma(h^{-1}(W),\sh{O}_X)$.
On the other hand, the canonical homomorphism
\sref[0\textsubscript{I}]{0.4.4.3.3} gives, by functoriality,
canonical homomorphisms
\[
  \HH^n\bigl(h^{-1}(W),f^*(f_*(\sh{F}))\bigr)
  \longrightarrow
  \HH^n\bigl(h^{-1}(W),\sh{F}\bigr),
  \tag{12.2.5.4}
  \label{0.12.2.5.4}
\]
which are homomorphisms of
$\Gamma(h^{-1}(W),\sh{O}_X)$-modules.
Taking into account the ring homomorphism
\[
  \Gamma(W,\sh{O}_Z)
  \longrightarrow\Gamma(g^{-1}(W),\sh{O}_Y),
\]
we see that composing \sref{0.12.2.5.3} with
\sref{0.12.2.5.4} gives a homomorphism of presheaves and hence the
homomorphism of sheaves \sref{0.12.2.5.1}.

The definition of \sref{0.12.2.5.2} is still simpler.
By definition, $\RR^n h_*(\sh{F})$ is associated with the presheaf
\[
  W\longmapsto
  \HH^n\bigl(f^{-1}(g^{-1}(W)),\sh{F}\bigr),
\]
while $\RR^n f_*(\sh{F})$ is associated with the presheaf
\[
  V\longmapsto\HH^n(f^{-1}(V),\sh{F}).
\]
There is therefore a canonical homomorphism
\[
  \HH^n\bigl(f^{-1}(g^{-1}(W)),\sh{F}\bigr)
  \longrightarrow
  \Gamma\bigl(g^{-1}(W),\RR^n f_*(\sh{F})\bigr),
\]
and it is immediate that these homomorphisms define a homomorphism of
presheaves, which in turn defines \sref{0.12.2.5.2}.
\end{env}

\begin{env}[12.2.6]
\label{0.12.2.6}
Under the hypotheses of \sref{0.12.2.4}, let $\sh{F}$, $\sh{G}$,
and $\sh{K}$ be three $\sh{O}_X$-modules, and let
\[
  u:\sh{F}\otimes_{\sh{O}_X}\sh{G}\longrightarrow\sh{K}
\]
be an $\sh{O}_X$-homomorphism.
Then the following diagrams are commutative:
\[
  \xymatrix@C=10pt@R=24pt{
    \RR^p g_*(f_*(\sh{F}))
      \otimes_{\sh{O}_Z}\RR^q g_*(f_*(\sh{G}))
      \ar[r]\ar[d]
      & \RR^{p+q}g_*(f_*(\sh{K}))\ar[d]\\
    \RR^p h_*(\sh{F})
      \otimes_{\sh{O}_Z}\RR^q h_*(\sh{G})
      \ar[r]
      & \RR^{p+q}h_*(\sh{K}),
  }
  \tag{12.2.6.1}
  \label{0.12.2.6.1}
\]
and
\[
  \xymatrix@C=10pt@R=24pt{
    \RR^p h_*(\sh{F})
      \otimes_{\sh{O}_Z}\RR^q h_*(\sh{G})
      \ar[r]\ar[d]
      & \RR^{p+q}h_*(\sh{K})\ar[d]\\
    g_*(\RR^p f_*(\sh{F}))
      \otimes_{\sh{O}_Z}g_*(\RR^q f_*(\sh{G}))
      \ar[r]
      & g_*(\RR^{p+q}f_*(\sh{K})).
  }
  \tag{12.2.6.2}
  \label{0.12.2.6.2}
\]
\oldpage[0\textsubscript{III}]{60}
Here the horizontal arrows come from \sref{0.12.2.2.1} (the last one
combined with \sref[0\textsubscript{I}]{0.4.2.2.1}), and the vertical
arrows from \sref{0.12.2.5.1} and \sref{0.12.2.5.2}, respectively.

Indeed, it is enough to verify the assertion for the corresponding
homomorphisms of presheaves.  Returning to the definitions in
\sref{0.12.2.2} and \sref{0.12.2.5}, one is immediately reduced, for
\sref{0.12.2.6.1}, to the commutative diagrams
\sref{0.12.1.5.2}; the verification is still simpler for
\sref{0.12.2.6.2}.
\end{env}

\subsection{Supplements on the Ext functors of sheaves}
\label{subsection:0.12.3}

\begin{env}[12.3.1]
\label{0.12.3.1}
Let $(X,\sh{O}_X)$ be a ringed space.
We shall not recall the definition and principal properties of the
bifunctors
\[
  \Ext^p_{\sh{O}_X}(X;\sh{F},\sh{G}),
  \qquad
  \shExt^p_{\sh{O}_X}(\sh{F},\sh{G}),
\]
the first taking values in the category of
$\Gamma(X,\sh{O}_X)$-modules and the second in the category of
$\sh{O}_X$-modules, nor the biregular spectral sequence
$E(\sh{F},\sh{G})$ relating them (T, 4.2 and G, II, 7.3).
\end{env}

\begin{env}[12.3.2]
\label{0.12.3.2}
Just as in (M, XIV, 1), one defines the notion of an
\emph{extension} of an $\sh{O}_X$-module $\sh{F}$ by an
$\sh{O}_X$-module $\sh{G}$ and the composition law on classes of
equivalent extensions: the arguments used for modules plainly carry
over to an arbitrary abelian category.
The second proof of (M, XIV, 1.1), which uses only the existence of
embeddings into injective objects, is therefore still valid in the
category of $\sh{O}_X$-modules.  It shows that
\[
  \Ext^1_{\sh{O}_X}(X;\sh{F},\sh{G})
\]
is canonically identified with the abelian group of classes of
extensions of $\sh{F}$ by $\sh{G}$.
\end{env}

\begin{proposition}[12.3.3]
\label{0.12.3.3}
Let $(X,\sh{O}_X)$ be a ringed space whose sheaf of rings
$\sh{O}_X$ is coherent.
Then, for every pair of coherent $\sh{O}_X$-modules $\sh{F}$ and
$\sh{G}$ and every $p\geq0$,
\[
  \shExt^p_{\sh{O}_X}(\sh{F},\sh{G})
\]
is a coherent $\sh{O}_X$-module.
\end{proposition}

\begin{proof}
The $\shExt^p_{\sh{O}_X}(\sh{F},\sh{G})$ form a contravariant
cohomological functor in $\sh{F}$.
Since $\sh{F}$ is coherent, for every $p$ and every point $x\in X$
there are an open neighbourhood $U$ of $x$ in $X$ and an exact
sequence of $(\sh{O}_X|U)$-modules
\[
  0\longrightarrow\sh{R}\longrightarrow\sh{L}_{p-1}
  \longrightarrow\cdots\longrightarrow\sh{L}_0
  \longrightarrow\sh{F}|U\longrightarrow0,
\]
where each $\sh{L}_i$ ($0\leq i\leq p-1$) is isomorphic to
$\sh{O}_X^{n_i}|U$, and $\sh{R}$ is coherent.
This follows by induction on $p$ from
\sref[0\textsubscript{I}]{0.5.3.2} and
\sref[0\textsubscript{I}]{0.5.3.4}, since $\sh{O}_X$ is coherent.

For $p\geq1$, moreover,
\[
  \shExt^p_{\sh{O}_X|U}(\sh{L}|U,\sh{G}|U)=0
\]
for every $\sh{O}_X$-module $\sh{L}$ for which $\sh{L}|U$ is
isomorphic to $\sh{O}_X^n|U$ (T, 4.2.3).
The argument of (M, V, 7.2) therefore applies to the contravariant
cohomological functor
\[
  \sh{F}\longmapsto
  \shHom_{\sh{O}_X|U}(\sh{F}|U,\sh{G}|U)
\]
and gives an exact sequence
\[
  \begin{aligned}
  \shHom_{\sh{O}_X|U}(\sh{L}_{p-1},\sh{G}|U)
  &\longrightarrow
  \shHom_{\sh{O}_X|U}(\sh{R},\sh{G}|U)\\
  &\longrightarrow
  \shExt^p_{\sh{O}_X|U}(\sh{F}|U,\sh{G}|U)
  \longrightarrow0.
  \end{aligned}
\]
The first two terms are coherent $(\sh{O}_X|U)$-modules
\sref[0\textsubscript{I}]{0.5.3.5}; hence so is the third
\sref[0\textsubscript{I}]{0.5.3.4}.
\end{proof}

\oldpage[0\textsubscript{III}]{61}
\begin{proposition}[12.3.4]
\label{0.12.3.4}
Let $f:X\to Y$ be a flat morphism of ringed spaces, and let
$\sh{F}$ and $\sh{G}$ be two $\sh{O}_Y$-modules.
\begin{enumerate}
  \item[\rm(i)] There is a homomorphism of cohomological bifunctors
  \[
    f^*\bigl(\shExt^\bullet_{\sh{O}_Y}(\sh{F},\sh{G})\bigr)
    \longrightarrow
    \shExt^\bullet_{\sh{O}_X}
      \bigl(f^*(\sh{F}),f^*(\sh{G})\bigr),
    \tag{12.3.4.1}
    \label{0.12.3.4.1}
  \]
  which in degree $0$ reduces to the canonical homomorphism
  \sref[0\textsubscript{I}]{0.4.4.6}.

  \item[\rm(ii)] There is a canonical morphism of spectral sequences
  \[
    E(\sh{F},\sh{G})
    \longrightarrow E\bigl(f^*(\sh{F}),f^*(\sh{G})\bigr),
    \tag{12.3.4.2}
    \label{0.12.3.4.2}
  \]
  whose map on $E_2$-terms is given by the homomorphisms
  \[
    \begin{aligned}
    \HH^p\bigl(Y,\shExt^q_{\sh{O}_Y}(\sh{F},\sh{G})\bigr)
    \longrightarrow
    \HH^p\bigl(X,\shExt^q_{\sh{O}_X}
      (f^*(\sh{F}),f^*(\sh{G}))\bigr),
    \end{aligned}
    \tag{12.3.4.3}
    \label{0.12.3.4.3}
  \]
  deduced from \sref{0.12.3.4.1} and \sref{0.12.1.3.1}.
\end{enumerate}
\end{proposition}

\begin{proof}
\begin{enumerate}
  \item[\rm(i)]
  Since $f^*$ is exact on the category of $\sh{O}_Y$-modules
  \sref[0\textsubscript{I}]{0.6.7.2}, the functors
  \[
    \sh{G}\longmapsto
      f^*\bigl(\shHom_{\sh{O}_Y}(\sh{F},\sh{G})\bigr),
    \qquad
    \sh{G}\longmapsto
      \shHom_{\sh{O}_X}\bigl(f^*(\sh{F}),f^*(\sh{G})\bigr)
  \]
  are left exact.  The canonical homomorphism
  \sref[0\textsubscript{I}]{0.4.4.6} therefore canonically induces a
  homomorphism between their derived functors.

  To compute these derived functors, take an injective resolution
  $\sh{L}^\bullet=(\sh{L}^i)$ of $\sh{G}$.  One then has
  homomorphisms
  \[
    \begin{aligned}
    \sh{H}^p\!\left(
      f^*\bigl(\shHom_{\sh{O}_Y}
        (\sh{F},\sh{L}^\bullet)\bigr)\right)
    \longrightarrow
    \sh{H}^p\!\left(
      \shHom_{\sh{O}_X}
        (f^*(\sh{F}),f^*(\sh{L}^\bullet))\right)
    \end{aligned}
  \]
  between the cohomology sheaves of complexes.
  By the exactness of $f^*$, the first member is
  \[
    f^*\bigl(\shExt^p_{\sh{O}_Y}(\sh{F},\sh{G})\bigr).
  \]
  The exactness of $f^*$ also implies that
  $f^*(\sh{L}^\bullet)$ is a resolution of $f^*(\sh{G})$.
  If ${\sh{L}'}^\bullet=({\sh{L}'}^i)$ is an injective resolution
  of $f^*(\sh{G})$ in the category of $\sh{O}_X$-modules, there is a
  homomorphism of complexes
  \[
    f^*(\sh{L}^\bullet)\longrightarrow{\sh{L}'}^\bullet,
  \]
  determined up to homotopy, and hence a well-defined homomorphism on
  cohomology.  Its composite with the preceding homomorphism is
  \sref{0.12.3.4.1}.

  \item[\rm(ii)]
  With the preceding notation, there is a homomorphism of complexes of
  $\sh{O}_X$-modules
  \[
    f^*\bigl(\shHom_{\sh{O}_Y}(\sh{F},\sh{L}^\bullet)\bigr)
    \longrightarrow
    \shHom_{\sh{O}_X}
      \bigl(f^*(\sh{F}),{\sh{L}'}^\bullet\bigr).
  \]
  Let $\sh{M}^{\bullet\bullet}$ be a Cartan--Eilenberg injective
  resolution of
  $\shHom_{\sh{O}_Y}(\sh{F},\sh{L}^\bullet)$ in the category of
  $\sh{O}_Y$-modules.  Since $f^*$ is exact,
  $f^*(\sh{M}^{\bullet\bullet})$ is a Cartan--Eilenberg resolution of
  $f^*(\shHom_{\sh{O}_Y}(\sh{F},\sh{L}^\bullet))$.
  If ${\sh{M}'}^{\bullet\bullet}$ is a Cartan--Eilenberg injective
  resolution of
  $\shHom_{\sh{O}_X}(f^*(\sh{F}),{\sh{L}'}^\bullet)$, there is
  \sref{0.11.4.2} a homomorphism, determined up to homotopy,
  \[
    f^*(\sh{M}^{\bullet\bullet})
    \longrightarrow{\sh{M}'}^{\bullet\bullet},
  \]
  compatible with the preceding homomorphism; in other words, there
  is an $f$-morphism
  $\sh{M}^{\bullet\bullet}\to{\sh{M}'}^{\bullet\bullet}$ of
  bicomplexes of sheaves.
  It induces a di-homomorphism
  \[
    \Gamma(Y,\sh{M}^{\bullet\bullet})
    \longrightarrow\Gamma(X,{\sh{M}'}^{\bullet\bullet})
  \]
  of bicomplexes of modules, determined up to homotopy, and therefore
  a well-defined morphism of spectral sequences
  \sref{0.11.3.2}.  This is precisely the desired morphism
  \sref{0.12.3.4.2}, and the description
  \sref{0.12.3.4.3} follows immediately from the definitions.
\end{enumerate}
\end{proof}

\begin{proposition}[12.3.5]
\label{0.12.3.5}
Under the hypotheses of \sref{0.12.3.4}, suppose in addition that the
sheaf of rings $\sh{O}_Y$ is coherent.
Then, for every coherent $\sh{O}_Y$-module $\sh{F}$, the canonical
homomorphisms \sref{0.12.3.4.1} are bijective.
\end{proposition}

\begin{proof}
\oldpage[0\textsubscript{III}]{62}
Since the question is local on $Y$, we may suppose that there is an
exact sequence
\[
  0\longrightarrow\sh{R}\longrightarrow\sh{O}_Y^n
  \longrightarrow\sh{F}\longrightarrow0.
\]
Then $\sh{R}$ is also a coherent $\sh{O}_Y$-module
\sref[0\textsubscript{I}]{0.5.3.4}.
To prove that the homomorphisms
\[
  f^*\bigl(\shExt^p_{\sh{O}_Y}(\sh{F},\sh{G})\bigr)
  \longrightarrow
  \shExt^p_{\sh{O}_X}
    \bigl(f^*(\sh{F}),f^*(\sh{G})\bigr)
\]
are bijective, we argue by induction on $p$.
For $p=0$, the assertion is
\sref[0\textsubscript{I}]{0.6.7.6.1}.
There is a commutative diagram with exact rows
\[
{\tiny
  \xymatrix@C=2pt@R=24pt{
  f^*\!\left(\shExt^{p-1}_{\sh{O}_Y}(\sh{O}_Y^n,\sh{G})\right)
    \ar[r]\ar[d]
  & f^*\!\left(\shExt^{p-1}_{\sh{O}_Y}(\sh{R},\sh{G})\right)
    \ar[r]^\partial\ar[d]
  & f^*\!\left(\shExt^p_{\sh{O}_Y}(\sh{F},\sh{G})\right)
    \ar[r]\ar[d]
  & f^*\!\left(\shExt^p_{\sh{O}_Y}(\sh{O}_Y^n,\sh{G})\right)
    \ar[d]\\
  \shExt^{p-1}_{\sh{O}_X}(\sh{O}_X^n,f^*(\sh{G}))
    \ar[r]
  & \shExt^{p-1}_{\sh{O}_X}(f^*(\sh{R}),f^*(\sh{G}))
    \ar[r]^\partial
  & \shExt^p_{\sh{O}_X}(f^*(\sh{F}),f^*(\sh{G}))
    \ar[r]
  & \shExt^p_{\sh{O}_X}(\sh{O}_X^n,f^*(\sh{G})).
  }}
\]
Here we have used $f^*(\sh{O}_Y)=\sh{O}_X$; since $f^*$ is exact,
both rows are exact.
Moreover,
\[
  \shExt^p_{\sh{O}_Y}(\sh{O}_Y^n,\sh{G})=0,
  \qquad
  \shExt^p_{\sh{O}_X}(\sh{O}_X^n,f^*(\sh{G}))=0
  \qquad(p>0)
\]
(T, 4.2.3).
By the induction hypothesis, the first two vertical arrows in the
preceding diagram are isomorphisms, while the right-hand terms are
zero.  It follows that
\[
  f^*\bigl(\shExt^p_{\sh{O}_Y}(\sh{F},\sh{G})\bigr)
  \longrightarrow
  \shExt^p_{\sh{O}_X}
    \bigl(f^*(\sh{F}),f^*(\sh{G})\bigr)
\]
is an isomorphism.
\end{proof}

\subsection{Hypercohomology of the direct image functor}
\label{subsection:0.12.4}

\begin{env}[12.4.1]
\label{0.12.4.1}
Let $(X,\sh{O}_X)$ and $(Y,\sh{O}_Y)$ be two ringed spaces, and let
$f:X\longrightarrow Y$ be a morphism of ringed spaces.
One can form the hypercohomology of $f_*$ with respect to an arbitrary
complex of $\sh{O}_X$-modules
$\sh{K}^\bullet=(\sh{K}^i)_{i\in\bb{Z}}$
\sref{0.11.4.4}, since filtered inductive limits exist and are exact in
the abelian category of $\sh{O}_Y$-modules (T, 3.1.1).
The $\sh{O}_Y$-modules of hypercohomology
$\RR^p f_*(\sh{K}^\bullet)$ will also be denoted
$\sh{H}^p(f,\sh{K}^\bullet)$, or simply
$\sh{H}^p_f(\sh{K}^\bullet)$.
Recall that $\sh{H}^\bullet(f,\sh{K}^\bullet)$ is the cohomology of
the bicomplex of $\sh{O}_Y$-modules $f_*(\sh{L}^{\bullet\bullet})$,
where $\sh{L}^{\bullet\bullet}$ is a Cartan--Eilenberg injective
resolution of $\sh{K}^\bullet$ in the category of
$\sh{O}_X$-modules.
It is the abutment of two spectral sequences whose $E_2$-terms are
given by
\begin{equation}
\label{0.12.4.1.1}
\tag{12.4.1.1}
  {}'E_2^{pq}
  =\sh{H}^p\bigl(\sh{H}^q(f,\sh{K}^\bullet)\bigr)
\end{equation}
and
\begin{equation}
\label{0.12.4.1.2}
\tag{12.4.1.2}
  {}''E_2^{pq}
  =\sh{H}^p\bigl(f,\sh{H}^q(\sh{K}^\bullet)\bigr)
  \quad
  \bigl(=\RR^p f_*(\sh{H}^q(\sh{K}^\bullet))\bigr).
\end{equation}
In these formulas we have adopted the general notation $T(A^\bullet)$
for the transform of a complex by a functor \sref{0.11.2.1}, and, for
an $\sh{O}_X$-module $\sh{F}$, we write
$\sh{H}^p(f,\sh{F})$ in place of $\RR^p f_*(\sh{F})$.
Recall also that the first spectral sequence is always regular.
Both spectral sequences are biregular when $\sh{K}^\bullet$ is
bounded below,
\oldpage[0\textsubscript{III}]{63}
or when there is an integer $m$ such that every $\sh{O}_X$-module
admits a flasque resolution of length $\leq m$ \sref{0.11.4.4}.
\end{env}

\begin{env}[12.4.2]
\label{0.12.4.2}
We shall similarly denote by $\bb{H}^\bullet(X,\sh{K}^\bullet)$ the
hypercohomology of the functor $\Gamma$ with respect to a complex
$\sh{K}^\bullet$ of $\sh{O}_X$-modules.
Thus the $\bb{H}^p(X,\sh{K}^\bullet)$ are
$\Gamma(X,\sh{O}_X)$-modules.
Moreover, $\bb{H}^\bullet(X,\sh{K}^\bullet)$ may be regarded as a
special case of $\sh{H}^\bullet(f,\sh{K}^\bullet)$, where $f$ is a
morphism from $(X,\sh{O}_X)$ to a ringed space consisting of one point
equipped with the ring $\Gamma(X,\sh{O}_X)$.

For every open subset $V$ of $X$, we shall write
$\bb{H}^\bullet(V,\sh{K}^\bullet)$ in place of
$\bb{H}^\bullet(V,\sh{K}^\bullet|V)$.
\end{env}

\begin{proposition}[12.4.3]
\label{0.12.4.3}
For every integer $p\in\bb{Z}$, the $\sh{O}_Y$-module
$\sh{H}^p(f,\sh{K}^\bullet)$ is canonically isomorphic to the sheaf
associated with the presheaf
\[
  U\longmapsto
  \bb{H}^p\bigl(f^{-1}(U),\sh{K}^\bullet\bigr)
\]
on $Y$.
\end{proposition}

\begin{proof}
With the notation of \sref{0.12.4.1}, the cohomology sheaf
$\sh{H}^p(f_*(\sh{L}^{\bullet\bullet}))$ is associated with the
presheaf
\[
  U\longmapsto
  \HH^p\bigl(\Gamma(U,f_*(\sh{L}^{\bullet\bullet}))\bigr)
  =
  \HH^p\bigl(\Gamma(f^{-1}(U),\sh{L}^{\bullet\bullet})\bigr).
\]
But it is clear that
$\sh{L}^{\bullet\bullet}|f^{-1}(U)$ is a Cartan--Eilenberg injective
resolution of $\sh{K}^\bullet|f^{-1}(U)$ (T, 3.1.3).
Consequently,
\[
  \HH^p\bigl(\Gamma(f^{-1}(U),\sh{L}^{\bullet\bullet})\bigr)
  =
  \bb{H}^p\bigl(f^{-1}(U),\sh{K}^\bullet\bigr)
\]
by definition.
\end{proof}

\begin{proposition}[12.4.4]
\label{0.12.4.4}
The hypercohomology $\sh{H}^\bullet(f,\sh{K}^\bullet)$ is a
cohomological functor of $\sh{K}^\bullet$ in each of the following
cases:
\begin{enumerate}
  \item[\rm{a)}] $\sh{K}^\bullet$ varies in the category of bounded
    below complexes;
  \item[\rm{b)}] there is an integer $m$ such that every
    $\sh{O}_X$-module admits a flasque resolution of length $\leq m$;
  \item[\rm{c)}] $X$ is a Noetherian space.
\end{enumerate}
\end{proposition}

\begin{proof}
Cases~a) and~b) are special cases of \sref{0.11.5.4}.
Case~c), on the other hand, follows from \sref{0.11.5.2}, since in
this case the functor $f_*$ commutes with inductive limits
(G, II, 3.10.1).
\end{proof}

\begin{env}[12.4.5]
\label{0.12.4.5}
Let $\mathfrak{U}=(U_\alpha)$ now be an open covering of $X$.
For every complex of \emph{presheaves}
$\sh{K}^\bullet=(\sh{K}^j)$ on $X$, consider the bicomplex
$\check{\mathrm{C}}^\bullet(\mathfrak{U},\sh{K}^\bullet)$ whose
component of bidegree $(i,j)$ is
$\check{\mathrm{C}}^i(\mathfrak{U},\sh{K}^j)$, the group of
alternating $i$-cochains of the nerve of $\mathfrak{U}$ with values
in $\sh{K}^j$ (G, II, 5.1).
We call the cohomology of this bicomplex the
\emph{hypercohomology of the covering $\mathfrak{U}$ with
coefficients in $\sh{K}^\bullet$}, and denote it by
\[
  \bb{H}^\bullet(\mathfrak{U},\sh{K}^\bullet)
  =
  \HH^\bullet
  \bigl(\check{\mathrm{C}}^\bullet
    (\mathfrak{U},\sh{K}^\bullet)\bigr).
\]
The Leray spectral sequence of a covering
(T, 3.8.1 and G, II, 5.9.1) generalizes to hypercohomology as follows.
\end{env}

\begin{proposition}[12.4.6]
\label{0.12.4.6}
Let $\sh{K}^\bullet=(\sh{K}^i)$ be a complex of
$\sh{O}_X$-modules.
There is a regular spectral functor of $\sh{K}^\bullet$ with
abutment $\bb{H}^\bullet(X,\sh{K}^\bullet)$ and $E_2$-term
\begin{equation}
\label{0.12.4.6.1}
\tag{12.4.6.1}
  E_2^{pq}
  =
  \CHH^p\bigl(\mathfrak{U},h^q(\sh{K}^\bullet)\bigr),
\end{equation}
where $h^q(\sh{K}^\bullet)$ denotes the presheaf complex
\[
  V\longmapsto\HH^q(V,\sh{K}^\bullet)
\]
on $X$.
This spectral sequence is biregular if $\sh{K}^\bullet$ is bounded
below.
\end{proposition}

\begin{proof}
Let $\sh{L}^{\bullet\bullet}$ be a Cartan--Eilenberg injective
resolution of $\sh{K}^\bullet$, and consider the tricomplex
\[
  \check{\mathrm{C}}^\bullet
  (\mathfrak{U},\sh{L}^{\bullet\bullet})
  =
  \bigl(\check{\mathrm{C}}^i
    (\mathfrak{U},\sh{L}^{jk})\bigr).
\]
First regard this tricomplex as a bicomplex in the degrees $i$ and
$j+k$.
Since $i$ takes only nonnegative values, the second spectral sequence
of this bicomplex is regular \sref{0.11.3.3} and degenerate, because
\[
  \CHH^q(\mathfrak{U},\sh{L}^{jk})=0
  \qquad(q>0),
\]
the $\sh{O}_X$-modules $\sh{L}^{jk}$ being flasque sheaves
\oldpage[0\textsubscript{III}]{64}
(G, II, 5.2.3).
It follows that there is a canonical isomorphism
\[
  \HH^n\bigl(
    \check{\mathrm{C}}^\bullet
      (\mathfrak{U},\sh{L}^{\bullet\bullet})
  \bigr)
  \isoto
  \HH^n\bigl(\Gamma(X,\sh{L}^{\bullet\bullet})\bigr)
\]
\sref{0.11.1.6}, by virtue of (G, II, 5.2.2), and therefore, by
definition \sref{0.12.4.2}, an isomorphism
\[
  \HH^n\bigl(
    \check{\mathrm{C}}^\bullet
      (\mathfrak{U},\sh{L}^{\bullet\bullet})
  \bigr)
  \isoto
  \bb{H}^n(X,\sh{K}^\bullet).
\]

Now regard
$\check{\mathrm{C}}^\bullet
(\mathfrak{U},\sh{L}^{\bullet\bullet})$
as a bicomplex in the degrees $i+j$ and $k$.
Since $k$ takes only nonnegative values, the first spectral sequence
of this bicomplex is always regular.
It is biregular if $\sh{L}^{jk}=0$ for $j<j_0$, that is, when
$\sh{K}^\bullet$ is bounded below \sref{0.11.3.3}.
This is the desired spectral sequence.
Indeed, for every $j$, $\sh{L}^{j,\bullet}$ is an injective
resolution of $\sh{K}^j$; consequently, as $i$ varies,
\[
  \HH^q\bigl(
    \check{\mathrm{C}}^i
      (\mathfrak{U},\sh{L}^{j,\bullet})
  \bigr)
\]
is precisely the cochain complex
$\check{\mathrm{C}}^i(\mathfrak{U},h^q(\sh{K}^j))$.
This proves the proposition.
\end{proof}

\begin{corollary}[12.4.7]
\label{0.12.4.7}
If, for every simplex $\sigma$ of the nerve of $\mathfrak{U}$ and
every integer $i$, one has
\[
  \HH^q(U_\sigma,\sh{K}^i)=0
  \qquad(q>0),
\]
then there is a canonical isomorphism
\begin{equation}
\label{0.12.4.7.1}
\tag{12.4.7.1}
  \bb{H}^\bullet(\mathfrak{U},\sh{K}^\bullet)
  \isoto
  \bb{H}^\bullet(X,\sh{K}^\bullet).
\end{equation}
\end{corollary}

\begin{proof}
The hypothesis implies that
$\check{\mathrm{C}}^\bullet
(\mathfrak{U},h^q(\sh{K}^j))=0$ for $q>0$, and hence that
$E_2^{pq}=0$ for $q>0$.
Since the spectral sequence \sref{0.12.4.6.1} is degenerate and
regular, the conclusion follows from the definition
\sref{0.12.4.5} of
$\bb{H}^\bullet(\mathfrak{U},\sh{K}^\bullet)$ and from
\sref{0.11.1.6}.
\end{proof}

\begin{env}[12.4.8]
\label{0.12.4.8}
Let $(X',\sh{O}_{X'})$ be a second ringed space, and let
$f=(\psi,\theta)$ be a morphism from $X'$ to $X$.
By the same method as in \sref{0.12.1.3} and \sref{0.12.1.4}, one
defines a di-homomorphism for the hypercohomology of a complex
$\sh{K}^\bullet$ of $\sh{O}_X$-modules:
\begin{equation}
\label{0.12.4.8.1}
\tag{12.4.8.1}
  \bb{H}^p(X,\sh{K}^\bullet)
  \longrightarrow
  \bb{H}^p(X',f^*(\sh{K}^\bullet)).
\end{equation}
Starting with a Cartan--Eilenberg injective resolution
$\sh{L}^{\bullet\bullet}$ of $\sh{K}^\bullet$, observe that, since
$\psi^*$ is exact, $\psi^*(\sh{L}^{\bullet\bullet})$ is a
Cartan--Eilenberg resolution of $\psi^*(\sh{K}^\bullet)$ in the
category of $\psi^*(\sh{O}_X)$-modules.
There is therefore a morphism
\[
  \psi^*(\sh{L}^{\bullet\bullet})
  \longrightarrow
  {\sh{L}'}^{\bullet\bullet},
\]
where ${\sh{L}'}^{\bullet\bullet}$ is a Cartan--Eilenberg injective
resolution of $\psi^*(\sh{K}^\bullet)$.
It induces a morphism in cohomology
\[
  \bb{H}^\bullet(X,\sh{K}^\bullet)
  \longrightarrow
  \bb{H}^\bullet
  (X',\psi^*(\sh{K}^\bullet)).
\]
Composing it with the morphism induced by functoriality from
$\psi^*(\sh{K}^\bullet)\longrightarrow f^*(\sh{K}^\bullet)$ gives
the desired morphism \sref{0.12.4.8.1}.

Starting with \sref{0.12.4.8.1} and \sref{0.12.4.3}, and arguing as
in \sref{0.12.2.5}, one can then define, for two morphisms of ringed
spaces
\[
  f:X\longrightarrow Y,
  \qquad
  g:Y\longrightarrow Z,
\]
with $h=g\circ f$, homomorphisms for the hypercohomology of a complex
$\sh{K}^\bullet$ of $\sh{O}_X$-modules:
\begin{equation}
\label{0.12.4.8.2}
\tag{12.4.8.2}
  \sh{H}^n(g,f_*(\sh{K}^\bullet))
  \longrightarrow
  \sh{H}^n(h,\sh{K}^\bullet),
\end{equation}
and
\begin{equation}
\label{0.12.4.8.3}
\tag{12.4.8.3}
  \sh{H}^n(h,\sh{K}^\bullet)
  \longrightarrow
  g_*\bigl(\sh{H}^n(f,\sh{K}^\bullet)\bigr).
\end{equation}
We leave the details of the definitions to the reader.
\end{env}

\section{Projective limits in homological algebra}
\label{section:0.13}

\subsection{The Mittag-Leffler condition}
\label{subsection:0.13.1}

\begin{env}[13.1.1]
\label{0.13.1.1}
Let $\cat{C}$ be an abelian category in which infinite products exist
(axiom AB~3* of T, 1.5).
Then the infimum of a family of subobjects of an object of $\cat{C}$
exists, and every projective system of objects of $\cat{C}$ has a
projective limit; the latter is a left exact functor of the projective
system in question (T, 1.8).
\oldpage[0\textsubscript{III}]{65}
Let $(A_\alpha,f_{\alpha\beta})$ be a projective system of objects of
$\cat{C}$ whose index set $\mathrm{I}$ is filtered to the right.
Put
\[
  A=\varprojlim A_\alpha,
\]
and, for each $\alpha\in\mathrm{I}$, let
$f_\alpha:A\longrightarrow A_\alpha$ be the canonical morphism.
For every $\alpha\in\mathrm{I}$, the subobjects
$f_{\alpha\beta}(A_\beta)$, for $\alpha\leq\beta$, form a decreasing
filtered family of subobjects of $A_\alpha$.
The subobject
\[
  A'_\alpha
  =
  \inf_{\beta\geq\alpha}f_{\alpha\beta}(A_\beta)
\]
is called the subobject of \emph{universal images} in $A_\alpha$.
It is clear that
\[
  f_\alpha(A)\subset A'_\alpha,
  \qquad
  f_{\alpha\beta}(A'_\beta)\subset A'_\alpha
  \quad(\alpha\leq\beta).
\]
Thus
$(A'_\alpha,f_{\alpha\beta}|A'_\beta)$ is a projective system and
\[
  A=\varprojlim A'_\alpha.
\]
\end{env}

\begin{env}[13.1.2]
\label{0.13.1.2}
For a projective system $(A_\alpha,f_{\alpha\beta})$ in $\cat{C}$,
the following is called the \emph{Mittag-Leffler condition}:
\begin{itemize}
  \item[\textbf{(ML)}]
  \emph{For every index $\alpha$, there is a $\beta\geq\alpha$ such
  that, for every $\gamma\geq\beta$,}
  \[
    f_{\alpha\gamma}(A_\gamma)
    =
    f_{\alpha\beta}(A_\beta).
  \]
\end{itemize}

It is clear that condition (ML) holds if all the
$f_{\alpha\beta}$ are epimorphisms.
Conversely, suppose that (ML) holds, and for each
$\alpha\in\mathrm{I}$ let $A'_\alpha$ be the subobject of universal
images in $A_\alpha$.
Then, for $\alpha\leq\beta$, the restriction of
$f_{\alpha\beta}$ to $A'_\beta$ is an \emph{epimorphism}
\[
  A'_\beta\longrightarrow A'_\alpha.
\]
Indeed, choose $\gamma\geq\beta$ such that
\[
  f_{\beta\delta}(A_\delta)
  =
  f_{\beta\gamma}(A_\gamma)
  \qquad(\delta\geq\gamma).
\]
Then $A'_\beta=f_{\beta\gamma}(A_\gamma)$, and therefore
\[
  f_{\alpha\beta}(A'_\beta)
  =
  f_{\alpha\gamma}(A_\gamma),
  \qquad
  A'_\alpha
  =
  f_{\alpha\gamma}(A_\gamma)
  =
  f_{\alpha\beta}(A'_\beta).
\]

Condition (ML) also holds when the objects $A_\alpha$ are
\emph{Artinian} in $\cat{C}$, that is, when every family of
subobjects of $A_\alpha$ has a minimal member.
Indeed, a minimal member of the decreasing filtered family
$\bigl(f_{\alpha\beta}(A_\beta)\bigr)$ of subobjects of $A_\alpha$
is necessarily the smallest of those subobjects.
\end{env}

\begin{env}[13.1.3]
\label{0.13.1.3}
\emph{Remark} (13.1.3).
Condition (ML) can also be formulated when $\cat{C}$ is, for
example, the category of sets.
One can again define the subset of universal images in $A_\alpha$,
and the observations made in \sref{0.13.1.1} and
\sref{0.13.1.2} remain valid.
\end{env}

\subsection{The Mittag-Leffler condition for abelian groups}
\label{subsection:0.13.2}

\begin{proposition}[13.2.1]
\label{0.13.2.1}
Let
\[
  0\longrightarrow A_\alpha
  \xrightarrow{u_\alpha}B_\alpha
  \xrightarrow{v_\alpha}C_\alpha
  \longrightarrow0
\]
be an exact sequence of projective systems of abelian groups, relative
to the same filtered index set $\mathrm{I}$.
\begin{enumerate}
  \item[{\rm(i)}]
  If $(B_\alpha)$ satisfies {\rm(ML)}, then so does $(C_\alpha)$.
  \item[{\rm(ii)}]
  If $(A_\alpha)$ and $(C_\alpha)$ satisfy {\rm(ML)}, then so does
  $(B_\alpha)$.
\end{enumerate}
\end{proposition}

\begin{proof}
Let $(f_{\alpha\beta})$, $(g_{\alpha\beta})$, and
$(h_{\alpha\beta})$ be the systems of homomorphisms defining
$(A_\alpha)$, $(B_\alpha)$, and $(C_\alpha)$, respectively.

(i) Suppose that
\[
  g_{\alpha\beta}(B_\beta)
  =
  g_{\alpha\lambda}(B_\lambda)
  \qquad(\lambda\geq\beta).
\]
Since $v_\beta$ and $v_\lambda$ are surjective,
\[
  h_{\alpha\beta}(C_\beta)
  =
  v_\alpha\bigl(g_{\alpha\beta}(B_\beta)\bigr)
  =
  v_\alpha\bigl(g_{\alpha\lambda}(B_\lambda)\bigr)
  =
  h_{\alpha\lambda}(C_\lambda)
\]
for $\lambda\geq\beta$.

(ii) Let $\alpha\in\mathrm{I}$.
Choose $\beta\geq\alpha$ such that
\[
  f_{\alpha\beta}(A_\beta)
  =
  f_{\alpha\lambda}(A_\lambda)
  \qquad(\lambda\geq\beta),
\]
and then choose $\gamma\geq\beta$ such that
\[
  h_{\beta\gamma}(C_\gamma)
  =
  h_{\beta\lambda}(C_\lambda)
  \qquad(\lambda\geq\gamma).
\]
Let $y_\alpha\in g_{\alpha\gamma}(B_\gamma)$, write
$y_\alpha=g_{\alpha\gamma}(y_\gamma)$ with $y_\gamma\in B_\gamma$,
and put $y_\beta=g_{\beta\gamma}(y_\gamma)$.
Then
\[
  v_\beta(y_\beta)
  =
  h_{\beta\gamma}\bigl(v_\gamma(y_\gamma)\bigr).
\]
For every $\lambda\geq\gamma$, there is, by hypothesis, a
$y_\lambda\in B_\lambda$ such that
\[
  h_{\beta\gamma}\bigl(v_\gamma(y_\gamma)\bigr)
  =
  h_{\beta\lambda}\bigl(v_\lambda(y_\lambda)\bigr).
\]
Hence
$v_\beta(y_\beta-g_{\beta\lambda}(y_\lambda))=0$, and therefore
there is an $x_\beta\in A_\beta$
\oldpage[0\textsubscript{III}]{66}
such that
\[
  y_\beta
  =
  g_{\beta\lambda}(y_\lambda)+u_\beta(x_\beta).
\]
It follows that
\[
  y_\alpha
  =
  g_{\alpha\lambda}(y_\lambda)
  +
  u_\alpha\bigl(f_{\alpha\beta}(x_\beta)\bigr).
\]
Since $\lambda\geq\beta$, there is an $x_\lambda\in A_\lambda$ such
that
\[
  f_{\alpha\beta}(x_\beta)=f_{\alpha\lambda}(x_\lambda).
\]
Consequently,
\[
  y_\alpha
  =
  g_{\alpha\lambda}
  \bigl(y_\lambda+u_\lambda(x_\lambda)\bigr)
  \in g_{\alpha\lambda}(B_\lambda),
\]
which proves the assertion.
\end{proof}

\begin{proposition}[13.2.2]
\label{0.13.2.2}
Let $\mathrm{I}$ be a filtered ordered set having a countable
cofinal subset, and let
\[
  0\longrightarrow A_\alpha
  \xrightarrow{u_\alpha}B_\alpha
  \xrightarrow{v_\alpha}C_\alpha
  \longrightarrow0
\]
be an exact sequence of projective systems of abelian groups indexed
by $\mathrm{I}$.
If $(A_\alpha)$ satisfies {\rm(ML)}, then the sequence
\[
  0\longrightarrow\varprojlim A_\alpha
  \longrightarrow\varprojlim B_\alpha
  \longrightarrow\varprojlim C_\alpha
  \longrightarrow0
\]
is exact.
\end{proposition}

\begin{proof}
It remains only to prove that the homomorphism
\[
  v=\varprojlim v_\alpha:
  \varprojlim B_\alpha\longrightarrow\varprojlim C_\alpha
\]
is surjective.
Let $z=(z_\alpha)\in\varprojlim C_\alpha$, and put
\[
  E_\alpha=v_\alpha^{-1}(z_\alpha).
\]
The $E_\alpha$ plainly form a projective system of nonempty sets under
the restrictions of the homomorphisms
$g_{\alpha\beta}:B_\beta\longrightarrow B_\alpha$.

We show that this projective system satisfies (ML).
Identifying $A_\alpha$ with a subgroup of $B_\alpha$ by means of
$u_\alpha$, choose, for every $\alpha\in\mathrm{I}$, an index
$\beta\geq\alpha$ such that
\[
  g_{\alpha\beta}(A_\beta)
  =
  g_{\alpha\lambda}(A_\lambda)
  \qquad(\lambda\geq\beta).
\]
We claim that also
\[
  g_{\alpha\beta}(E_\beta)
  =
  g_{\alpha\lambda}(E_\lambda)
  \qquad(\lambda\geq\beta).
\]
Indeed, take $y_\lambda\in E_\lambda$ and set
\[
  y_\beta=g_{\beta\lambda}(y_\lambda),
  \qquad
  y_\alpha=g_{\alpha\lambda}(y_\lambda).
\]
Let $y'_\alpha=g_{\alpha\beta}(y'_\beta)$ for some
$y'_\beta\in E_\beta$.
Then $y'_\beta-y_\beta=x_\beta\in A_\beta$, and by hypothesis there
is an $x_\lambda\in A_\lambda$ such that
\[
  g_{\alpha\beta}(x_\beta)
  =
  g_{\alpha\lambda}(x_\lambda).
\]
Thus
\[
  y'_\alpha
  =
  g_{\alpha\beta}(y_\beta)+g_{\alpha\beta}(x_\beta)
  =
  g_{\alpha\lambda}(y_\lambda+x_\lambda)
  \in g_{\alpha\lambda}(E_\lambda),
\]
which proves the claim.

It is known (Bourbaki, \emph{Top.\ g\'{e}n.}, chap.~II,
\textsection3, th.~1) that, under the hypotheses imposed on
$\mathrm{I}$, a projective system of nonempty sets satisfying (ML)
has a nonempty projective limit.
There is therefore a point
$y=(y_\alpha)\in\varprojlim E_\alpha$.
Since $v_\alpha(y_\alpha)=z_\alpha$ for every $\alpha$, one has
$z=v(y)$.
\end{proof}

\begin{proposition}[13.2.3]
\label{0.13.2.3}
Under the hypotheses on $\mathrm{I}$ of \sref{0.13.2.2}, let
$(\mathrm{K}^\bullet_\alpha)_{\alpha\in\mathrm{I}}$ be a projective
system of complexes of abelian groups
\[
  \mathrm{K}^\bullet_\alpha
  =
  (\mathrm{K}^n_\alpha)_{n\in\bb{Z}},
\]
whose differentials have degree $+1$.
For every $n$, there is a canonical functorial homomorphism
\begin{equation}
\label{0.13.2.3.1}
\tag{13.2.3.1}
  h_n:
  \HH^n\bigl(\varprojlim\mathrm{K}^\bullet_\alpha\bigr)
  \longrightarrow
  \varprojlim\HH^n(\mathrm{K}^\bullet_\alpha).
\end{equation}
If, for every degree $n$, the projective system of abelian groups
$(\mathrm{K}^n_\alpha)_{\alpha\in\mathrm{I}}$ satisfies {\rm(ML)},
then every $h_n$ is surjective.
If, in addition, the projective system
$(\HH^{n-1}(\mathrm{K}^\bullet_\alpha))_{\alpha\in\mathrm{I}}$
satisfies {\rm(ML)}, then $h_n$ is bijective.
\end{proposition}

\begin{proof}
Put
\[
  \mathrm{K}^n=\varprojlim\mathrm{K}^n_\alpha
  \qquad(n\in\bb{Z}).
\]
The homomorphisms $h_n$ are defined by the commutative diagrams
\[
  \xymatrix{
    \cdots\ar[r]
      & \mathrm{K}^{n-1}\ar[r]\ar[d]
      & \mathrm{K}^n\ar[r]\ar[d]
      & \mathrm{K}^{n+1}\ar[r]\ar[d]
      & \cdots\\
    \cdots\ar[r]
      & \mathrm{K}^{n-1}_\alpha\ar[r]
      & \mathrm{K}^n_\alpha\ar[r]
      & \mathrm{K}^{n+1}_\alpha\ar[r]
      & \cdots,
  }
\]
\oldpage[0\textsubscript{III}]{67}
the differentials in $\mathrm{K}^\bullet$ being the projective limits
of the corresponding differentials in the
$\mathrm{K}^\bullet_\alpha$.
Consider the exact sequences
\[
  (*_n)\qquad
  0\longrightarrow
  \mathrm{B}^n(\mathrm{K}^\bullet_\alpha)
  \longrightarrow
  \mathrm{Z}^n(\mathrm{K}^\bullet_\alpha)
  \longrightarrow
  \HH^n(\mathrm{K}^\bullet_\alpha)
  \longrightarrow0
\]
and
\[
  (**_n)\qquad
  0\longrightarrow
  \mathrm{Z}^{n-1}(\mathrm{K}^\bullet_\alpha)
  \longrightarrow
  \mathrm{K}^{n-1}_\alpha
  \longrightarrow
  \mathrm{B}^n(\mathrm{K}^\bullet_\alpha)
  \longrightarrow0.
\]
The hypothesis and Proposition~\sref{0.13.2.1}~(i) show that the
projective system
$(\mathrm{B}^n(\mathrm{K}^\bullet_\alpha))_{\alpha\in\mathrm{I}}$
satisfies (ML) for every $n$.
It follows from \sref{0.13.2.2} that
\[
  (***_n)\qquad
  0\longrightarrow
  \varprojlim_\alpha\mathrm{B}^n(\mathrm{K}^\bullet_\alpha)
  \longrightarrow
  \varprojlim_\alpha\mathrm{Z}^n(\mathrm{K}^\bullet_\alpha)
  \longrightarrow
  \varprojlim_\alpha\HH^n(\mathrm{K}^\bullet_\alpha)
  \longrightarrow0
\]
is exact.
Now
$\varprojlim_\alpha\mathrm{B}^n(\mathrm{K}^\bullet_\alpha)$
identifies with a subgroup of $\mathrm{K}^{n}$ containing
$\mathrm{B}^n(\mathrm{K}^\bullet)$, while
$\varprojlim_\alpha\mathrm{Z}^n(\mathrm{K}^\bullet_\alpha)$
identifies with a subgroup of $\mathrm{Z}^n(\mathrm{K}^\bullet)$.
Consequently, $h_n$ is surjective.

Suppose now, in addition, that
$(\HH^{n-1}(\mathrm{K}^\bullet_\alpha))_{\alpha\in\mathrm{I}}$
satisfies (ML).
The exact sequences $(*_{n-1})$ and
Proposition~\sref{0.13.2.1}~(ii) show that
$(\mathrm{Z}^{n-1}(\mathrm{K}^\bullet_\alpha))_{\alpha\in\mathrm{I}}$
satisfies (ML).
Applying \sref{0.13.2.2} to the exact sequences $(**_n)$ gives an
exact sequence
\[
  0\longrightarrow
  \varprojlim_\alpha
    \mathrm{Z}^{n-1}(\mathrm{K}^\bullet_\alpha)
  \longrightarrow
  \mathrm{K}^{n-1}
  \xrightarrow{u}
  \varprojlim_\alpha
    \mathrm{B}^n(\mathrm{K}^\bullet_\alpha)
  \longrightarrow0.
\]
Since
$\varprojlim_\alpha\mathrm{B}^n(\mathrm{K}^\bullet_\alpha)
\supset\mathrm{B}^n(\mathrm{K}^\bullet)$
and the composite of $u$ with the injection
\[
  \varprojlim_\alpha\mathrm{B}^n(\mathrm{K}^\bullet_\alpha)
  \longrightarrow\mathrm{K}^n
\]
is the differential
$\mathrm{K}^{n-1}\longrightarrow\mathrm{K}^n$, the surjectivity of
$u$ implies
\[
  \varprojlim_\alpha\mathrm{B}^n(\mathrm{K}^\bullet_\alpha)
  =
  \mathrm{B}^n(\mathrm{K}^\bullet).
\]
Thus $h_n$ is injective.
\end{proof}

\begin{env}[13.2.4]
\label{0.13.2.4}
\emph{Remarks} (13.2.4).
\begin{enumerate}
  \item[{\rm(i)}]
  The argument of \sref{0.13.2.2}
  (cf.\ Bourbaki, \emph{loc.\ cit.}) shows that the conclusion of
  that proposition remains valid if one assumes only that the
  $A_\alpha$ can be equipped with structures of complete metrizable
  spaces in which translations are homeomorphisms, that the maps
  $f_{\alpha\beta}:A_\beta\longrightarrow A_\alpha$ defining the
  projective system are uniformly continuous for the metrics in
  question, and that $(A_\alpha)$ satisfies the condition
  \begin{itemize}
    \item[\textbf{(ML')}]
    \emph{For every index $\alpha$, there is a $\beta\geq\alpha$ such
    that, for every $\gamma\geq\beta$,
    $f_{\alpha\gamma}(A_\gamma)$ is dense in
    $f_{\alpha\beta}(A_\beta)$.}
  \end{itemize}

  This gives an analogous supplement to \sref{0.13.2.3}.
  Suppose that $\mathrm{K}^n_\alpha=0$ for $n<0$ and every $\alpha$;
  that
  $(\mathrm{K}^n_\alpha)_{\alpha\in\mathrm{I}}$ satisfies (ML) for
  $n\geq0$; and that
  $A_\alpha=\HH^0(\mathrm{K}^\bullet_\alpha)$ can be equipped with
  metric-space structures having the properties just stated.
  Then the conclusions of \sref{0.13.2.3} are unchanged for
  $n\geq2$, and $h_1$ is moreover \emph{bijective}.
  Indeed, the argument of \sref{0.13.2.2} also shows that
  $(\mathrm{B}^1(\mathrm{K}^\bullet_\alpha))_{\alpha\in\mathrm{I}}$
  satisfies (ML), that the sequence $(***_1)$ is exact, and finally
  that
  \[
    \varprojlim_\alpha
    \mathrm{B}^1(\mathrm{K}^\bullet_\alpha)
    =
    \mathrm{B}^1(\mathrm{K}^\bullet).
  \]
  This proves, among other things, the assertions of (T, 3.10.2).

  \item[{\rm(ii)}]
  One can introduce the right derived functors
  $\varprojlim^{(n)}$ of the functor $\varprojlim$, and obtain more
  complete statements than those above~\cite{III-28}.
\end{enumerate}
\end{env}

\subsection{Application: cohomology of a projective limit of sheaves}
\label{subsection:0.13.3}

\begin{proposition}[13.3.1]
\label{0.13.3.1}
\oldpage[0\textsubscript{III}]{68}
Let $X$ be a topological space, let
$(\sh{F}_k)_{k\in\bb{N}}$ be a projective system of sheaves of
abelian groups on $X$, and put
$\sh{F}=\varprojlim_k\sh{F}_k$.
Suppose that the following conditions hold:
\begin{enumerate}
  \item[{\rm(i)}]
  There is a basis $\mathfrak{B}$ for the topology of $X$ such that,
  for every $U\in\mathfrak{B}$ and every $i\geq0$, the projective
  system
  $(\HH^i(U,\sh{F}_k))_{k\in\bb{N}}$ satisfies (ML).

  \item[{\rm(ii)}]
  For every $x\in X$ and every $i>0$, one has
  \[
    \varinjlim_U\bigl(\varprojlim_k
    \HH^i(U,\sh{F}_k)\bigr)=0,
  \]
  where $U$ ranges over the neighbourhoods of $x$ belonging to
  $\mathfrak{B}$.

  \item[{\rm(iii)}]
  The homomorphisms
  $u_{hk}:\sh{F}_k\longrightarrow\sh{F}_h$ $(h\leq k)$ defining the
  projective system $(\sh{F}_k)$ are surjective.
\end{enumerate}
Under these conditions, for every $i>0$, the canonical homomorphism
\[
  h_i:\HH^i(X,\sh{F})
  \longrightarrow
  \varprojlim_k\HH^i(X,\sh{F}_k)
\]
is surjective. If, moreover, for a given value of $i$ the projective
system
$(\HH^{i-1}(X,\sh{F}_k))_{k\in\bb{N}}$ satisfies (ML), then $h_i$ is
bijective.
\end{proposition}

\begin{proof}
\emph{a)}
We first suppose that the $\sh{F}_k$, as well as the kernels
$\sh{N}_{hk}$ of the $u_{hk}$, are flasque; we shall then show that
condition~{\rm(iii)} in the statement implies
$\HH^i(X,\sh{F})=0$ for $i>0$.
It suffices to prove that, for every open subset $U$ of $X$ and every
covering $\mathfrak{U}$ of $U$ by open subsets of $U$, one has
\[
  \CHH^i(\mathfrak{U},\sh{F})=0
  \qquad (i>0).
\]
Indeed, it then follows first that, for \v{C}ech cohomology,
$\check{\HH}^{\,i}(U,\sh{F})=0$ for every $i>0$, and hence, by
(G, II, 5.9.2) applied to the set of all open subsets of $X$, that
$\HH^i(X,\sh{F})=0$ for every $i>0$.

Since the $\sh{F}_k$ are flasque, one has
$\CHH^i(\mathfrak{U},\sh{F}_k)=0$ for every $i>0$
(G, II, 5.2.3).
For each $k$, consider the complex
$\check{\mathrm{C}}^\bullet(\mathfrak{U},\sh{F}_k)$ of alternating
cochains on the nerve of the covering $\mathfrak{U}$ (G, II, 5.1);
these complexes plainly form a projective system of complexes of
abelian groups.
Let us show that all the maps
\[
  \check{\mathrm{C}}^\bullet(\mathfrak{U},\sh{F}_k)
  \longrightarrow
  \check{\mathrm{C}}^\bullet(\mathfrak{U},\sh{F}_h)
  \qquad(h\leq k)
\]
are surjective.
By definition, it is enough to show that, for every open subset $V$
of $X$, the map
$\Gamma(V,\sh{F}_k)\longrightarrow\Gamma(V,\sh{F}_h)$ is
surjective. But the sequence
\[
  0\longrightarrow\sh{N}_{hk}\longrightarrow\sh{F}_k
  \longrightarrow\sh{F}_h\longrightarrow0
\]
is exact by hypothesis and therefore gives the exact cohomology
sequence
\[
  \Gamma(V,\sh{F}_k)\longrightarrow\Gamma(V,\sh{F}_h)
  \longrightarrow\HH^1(V,\sh{N}_{hk})=0,
\]
since $\sh{N}_{hk}$ is flasque.

The projective system
$(\check{\mathrm{C}}^\bullet(\mathfrak{U},\sh{F}_k))_{k\in\bb{N}}$
therefore satisfies (ML), and the same is true of
$(\CHH^i(\mathfrak{U},\sh{F}_k))_{k\in\bb{N}}$ for every $i\geq0$:
for $i>0$ this is trivial, while for $i=0$ it follows from what
precedes, since
$\CHH^0(\mathfrak{U},\sh{F}_k)=\Gamma(U,\sh{F}_k)$
(G, II, 5.2.2).
We may therefore apply \sref{0.13.2.3}, which gives
\[
  \CHH^i(\mathfrak{U},\sh{F})
  =
  \varprojlim_k\CHH^i(\mathfrak{U},\sh{F}_k)
  =0
  \qquad(i>0).
\]

\emph{b)}
We now pass to the general case. For each $k\in\bb{N}$, consider the
\emph{canonical resolution}
\[
  \sh{C}^\bullet(X,\sh{F}_k)
  =
  \bigl(\sh{C}^i(X,\sh{F}_k)\bigr)_{i\geq0}
\]
of $\sh{F}_k$ by flasque sheaves (G, II, 4.3).
For every $i\geq0$, it is clear that
$(\sh{C}^i(X,\sh{F}_k))_{k\in\bb{N}}$ is a projective system of
flasque sheaves; let us show that it satisfies the conditions of
part~\emph{a)}.
\oldpage[0\textsubscript{III}]{69}
Indeed, if $\sh{N}_{hk}$ denotes the kernel of $u_{hk}$ for $h\leq k$,
then the sequence
\[
  0\longrightarrow\sh{N}_{hk}\longrightarrow\sh{F}_k
  \longrightarrow\sh{F}_h\longrightarrow0
\]
is exact by~{\rm(iii)}, and our assertion follows from the exactness
of the functor
$\sh{A}\longmapsto\sh{C}^i(X,\sh{A})$ (G, II, 4.3).

Put
\[
  \sh{G}^i=\varprojlim_k\sh{C}^i(X,\sh{F}_k).
\]
By part~\emph{a)}, one has
$\HH^j(X,\sh{G}^i)=0$ for $j>0$ and $i\geq0$.
We shall show that
$\sh{G}^\bullet=(\sh{G}^i)_{i\geq0}$ is a resolution of the sheaf
$\sh{F}$. Since $\HH^j(X,\sh{G}^i)=0$ for $j>0$, the cohomology
$\HH^\bullet(X,\sh{F})$ will then be equal to
$\HH^\bullet(\Gamma(X,\sh{G}^\bullet))$ (G, II, 4.7.1).

Taking the projective limit of the exact sequences
\[
  0\longrightarrow\sh{F}_k
  \longrightarrow\sh{C}^0(X,\sh{F}_k)
  \longrightarrow\sh{C}^1(X,\sh{F}_k)
  \longrightarrow\cdots
\]
plainly gives a complex of sheaves of abelian groups
\[
  0\longrightarrow\sh{F}
  \longrightarrow\sh{G}^0
  \longrightarrow\sh{G}^1
  \longrightarrow\cdots.
\]
To prove our assertion, it remains to establish that
$\sh{H}^i(\sh{G}^\bullet)=0$ for $i>0$.
This sheaf is generated by the presheaf
\[
  U\longmapsto\HH^i\bigl(\Gamma(U,\sh{G}^\bullet)\bigr)
\]
(G, II, 4.1).
Now the complex $\Gamma(U,\sh{G}^\bullet)$ is the projective limit of
the projective system of complexes of abelian groups
\[
  \bigl(\Gamma(U,\sh{C}^\bullet(X,\sh{F}_k))\bigr)_{k\in\bb{N}}
\]
(0\textsubscript{I}, \hyperref[0.3.2.6]{3.2.6}).
We saw in part~\emph{a)} that, for every $i\geq0$, the maps
\[
  \Gamma(U,\sh{C}^i(X,\sh{F}_k))
  \longrightarrow
  \Gamma(U,\sh{C}^i(X,\sh{F}_h))
  \qquad(h\leq k)
\]
are surjective. On the other hand,
\[
  \HH^i(U,\sh{F}_k)
  =
  \HH^i\bigl(\Gamma(U,\sh{C}^\bullet(X,\sh{F}_k))\bigr),
\]
since the canonical resolution $\sh{C}^\bullet(U,\sh{F}_k|U)$ is
induced on $U$ by $\sh{C}^\bullet(X,\sh{F}_k)$.
By hypothesis~{\rm(i)}, for every $U\in\mathfrak{B}$ we may apply
\sref{0.13.2.3} to this projective system of complexes, and hence
\[
  \HH^i\bigl(\Gamma(U,\sh{G}^\bullet)\bigr)
  =
  \varprojlim_k\HH^i(U,\sh{F}_k)
  \qquad(i\geq0).
\]
Hypothesis~{\rm(ii)} then shows, by definition, that the sheaves
$\sh{H}^i(\sh{G}^\bullet)$ vanish for $i>0$.

Thus, for every $i\geq0$,
\[
  \HH^i(X,\sh{F})
  =
  \HH^i\bigl(\Gamma(X,\sh{G}^\bullet)\bigr),
  \qquad
  \Gamma(X,\sh{G}^\bullet)
  =
  \varprojlim_k\Gamma(X,\sh{C}^\bullet(X,\sh{F}_k)).
\]
We have just observed that all the maps
\[
  \Gamma(X,\sh{C}^i(X,\sh{F}_k))
  \longrightarrow
  \Gamma(X,\sh{C}^i(X,\sh{F}_h))
  \qquad(h\leq k)
\]
are surjective; the conclusion therefore follows once more from
\sref{0.13.2.3}.
\end{proof}

\begin{env}[13.3.2]
\label{0.13.3.2}
\emph{Remarks} (13.3.2).
\begin{enumerate}
  \item[{\rm(i)}]
  The statement of \sref{0.13.3.1} is of interest only for $i>0$,
  since, for $i=0$, $h_i$ is always an isomorphism without any
  hypothesis (0\textsubscript{I}, \hyperref[0.3.2.6]{3.2.6}).

  \item[{\rm(ii)}]
  Conditions~{\rm(i)} and~{\rm(ii)} of \sref{0.13.3.1} are satisfied,
  in particular, if
  $\HH^i(U,\sh{F}_k)=0$ for every $k$, every $i>0$, and every
  $U\in\mathfrak{B}$, and if, for $U\in\mathfrak{B}$, the maps
  \[
    \Gamma(U,\sh{F}_k)\longrightarrow\Gamma(U,\sh{F}_h)
  \]
  are surjective. This is the most frequent case in which
  \sref{0.13.3.1} is applied.
\end{enumerate}
\end{env}

\subsection{Mittag-Leffler condition and graded objects associated with projective systems}
\label{subsection:0.13.4}
\label{0.13.4}

\begin{env}[13.4.1]
\label{0.13.4.1}
Let
$\bb{A}=(A_k,u_{kh})_{k\in\bb{Z}}$ be a projective system in an
abelian category $\cat{C}$.
We say that it is \emph{bounded below} if there is a $k_0$ such that
$A_k=0$ for $k<k_0$.
On each $A_k$ we define a \emph{filtration}
$(\mathrm{F}^p(A_k))_{p\in\bb{Z}}$ by
\begin{equation}
\label{0.13.4.1.1}
\tag{13.4.1.1}
  \mathrm{F}^p(A_k)=
  \begin{cases}
    \Ker(A_k\longrightarrow A_{p-1}) & \text{for }p\leq k+1,\\
    0 & \text{for }p\geq k+1.
  \end{cases}
\end{equation}

\oldpage[0\textsubscript{III}]{70}
By hypothesis,
$\mathrm{F}^{k_0}(A_k)=A_k$ and
$\mathrm{F}^{k+1}(A_k)=0$; in other words, the filtration under
consideration is finite (\sref{0.11.1.3}).
The graded objects associated with this filtration are therefore
\[
  \mathrm{gr}^p(A_k)
  =
  \Ker(A_k\longrightarrow A_{p-1})
  \big/
  \Ker(A_k\longrightarrow A_p).
\]
Consequently, $\mathrm{gr}^p(A_k)$ is isomorphic to the image under
$A_k\longrightarrow A_p$ of
$\Ker(A_k\longrightarrow A_{p-1})$.
By the transitivity of the morphisms defining a projective system,
one therefore has
\begin{equation}
\label{0.13.4.1.2}
\tag{13.4.1.2}
  \mathrm{gr}^p(A_k)
  =
  \Ker(A_p\longrightarrow A_{p-1})
  \cap
  \Im(A_k\longrightarrow A_p).
\end{equation}
But \sref{0.13.4.1.1} gives
$\Ker(A_p\longrightarrow A_{p-1})=\mathrm{gr}^p(A_p)$, and hence
\begin{equation}
\label{0.13.4.1.3}
\tag{13.4.1.3}
  \mathrm{gr}^p(A_k)
  =
  \mathrm{gr}^p(A_p)
  \cap
  \Im(A_k\longrightarrow A_p).
\end{equation}

The preceding definitions also show that, for $k\leq h$,
\[
  u_{kh}\bigl(\mathrm{F}^p(A_h)\bigr)
  \subset
  \mathrm{F}^p(A_k),
\]
and consequently, for every $p\in\bb{Z}$, the
$\mathrm{gr}^p(u_{kh})$ define a projective system
$(\mathrm{gr}^p(A_k))_{k\in\bb{Z}}$.
\end{env}

\begin{env}[13.4.2]
\label{0.13.4.2}
We say that the projective system $\bb{A}$ is
\emph{essentially constant} if the morphisms
$A_{k+1}\longrightarrow A_k$ are isomorphisms for all sufficiently
large $k$.
We say that $\bb{A}$ is \emph{strict} if the morphisms
$A_i\longrightarrow A_j$ $(j\leq i)$ are epimorphisms.

When $\bb{A}$ is strict, it follows from \sref{0.13.4.1.3} that, for
$p\leq k\leq h$, the canonical morphism
\[
  \mathrm{gr}^p(A_h)\longrightarrow\mathrm{gr}^p(A_k)
\]
is an isomorphism; in other words, the projective system
$(\mathrm{gr}^p(A_k))_{k\in\bb{Z}}$ is essentially constant.
The sequence of objects $\mathrm{gr}^p(A_p)$, identified for every
$p$ with $\varprojlim_k\mathrm{gr}^p(A_k)$, is then denoted
$\mathrm{gr}^\bullet(\bb{A})$ and is called the
\emph{graded object associated with the strict projective system}
$\bb{A}=(A_k)$.

Suppose now that the bounded-below projective system $\bb{A}$
satisfies (ML).
By \sref{0.13.1.2}, the projective system
$\bb{A}'=(A'_k)$ of ``universal image'' objects is strict; it is also
bounded below.
The graded object $\mathrm{gr}^\bullet(\bb{A}')$ associated with
$\bb{A}'$ is then likewise called the graded object
\emph{associated with} $\bb{A}$ and is denoted
$\mathrm{gr}^\bullet(\bb{A})$.
\end{env}

\begin{proposition}[13.4.3]
\label{0.13.4.3}
Let $\bb{A}=(A_k)_{k\in\bb{Z}}$ be a bounded-below projective system
in an abelian category $\cat{C}$.
The following conditions are equivalent:
\begin{enumerate}
  \item[{\rm a)}] $\bb{A}$ satisfies (ML).
  \item[{\rm b)}] For every $p\in\bb{Z}$, the projective system
  $(\mathrm{gr}^p(A_k))_{k\in\bb{Z}}$ is essentially constant.
\end{enumerate}
Moreover, when these conditions hold, there is, for every
$p\in\bb{Z}$, a canonical isomorphism
\begin{equation}
\label{0.13.4.3.1}
\tag{13.4.3.1}
  \mathrm{gr}^p(\bb{A})
  \isoto
  \varprojlim_k\mathrm{gr}^p(A_k).
\end{equation}
\end{proposition}

\begin{proof}
It follows at once from \sref{0.13.4.1.2} that \emph{a)} implies
\emph{b)}. The same formula, applied to the projective system
$\bb{A}'$ in the notation of \sref{0.13.4.2}, gives
\sref{0.13.4.3.1} by definition.

For $k\leq h$, put
\[
  A_{kh}=\Im(A_h\longrightarrow A_k).
\]
If $k\leq h\leq j$, then
$A_{kj}\subset A_{kh}\subset A_k$.
Equip $A_{kh}$ with the filtration induced by
$(\mathrm{F}^p(A_k))$.
The transitivity of the morphisms defining $\bb{A}$ shows
immediately that this is also the quotient filtration of
$(\mathrm{F}^p(A_h))$; consequently,
\begin{equation}
\label{0.13.4.3.2}
\tag{13.4.3.2}
  \mathrm{gr}^p(A_{kh})
  =
  \Im\bigl(
    \mathrm{gr}^p(A_h)\longrightarrow\mathrm{gr}^p(A_k)
  \bigr).
\end{equation}

\oldpage[0\textsubscript{III}]{71}
Suppose now that \emph{b)} holds.
For every $p\in\bb{Z}$ and every $k\geq p$, there is an integer
$\mathrm{L}(p,k)$ such that the right-hand side of
\sref{0.13.4.3.2} is constant for $h\geq\mathrm{L}(p,k)$.
Since $\mathrm{gr}^p(A_k)=0$ for $p<k_0$, for fixed $k$ only
finitely many of the integers $\mathrm{L}(p,k)$ need be nonzero as
$p$ ranges over the integers $\leq k$.
Let $\mathrm{L}(k)=m$ be the largest of these integers.
For every $h\geq m$, one has
$A_{kh}\subset A_{km}$, and, by the definition of $m$, the canonical
injection
$A_{kh}\longrightarrow A_{km}$ induces an isomorphism
\[
  \mathrm{gr}^\bullet(A_{kh})
  \simto
  \mathrm{gr}^\bullet(A_{km}).
\]
Since the filtrations are finite, this injection is itself bijective
(Bourbaki, \emph{Alg.\ comm.}, chap.~III, \textsection2,
n\textsuperscript{o}~8, th.~1).
This proves that $\bb{A}$ satisfies (ML).
\end{proof}

\begin{env}[13.4.4]
\label{0.13.4.4}
Suppose that the projective limit
$A=\varprojlim_k A_k$ exists in $\cat{C}$.
In the definitions of \sref{0.13.4.1}, one may then replace $A_k$ by
$A$, and the filtration thereby defined on $A$ still satisfies
\begin{equation}
\label{0.13.4.4.1}
\tag{13.4.4.1}
  \mathrm{gr}^p(A)
  =
  \mathrm{gr}^p(A_p)
  \cap
  \Im(A\longrightarrow A_p).
\end{equation}
\end{env}

\begin{corollary}[13.4.5]
\label{0.13.4.5}
Suppose that $\cat{C}$ is the category of abelian groups.
If the projective system $\bb{A}$ satisfies (ML) and
$A=\varprojlim_k A_k$, then for every $p\in\bb{Z}$ there is a
canonical isomorphism
\begin{equation}
\label{0.13.4.5.1}
\tag{13.4.5.1}
  \mathrm{gr}^p(A)
  \isoto
  \varprojlim_k\mathrm{gr}^p(A_k).
\end{equation}
\end{corollary}

\begin{proof}
Indeed,
\[
  \Im(A_k\longrightarrow A_p)
  =
  \Im(A\longrightarrow A_p)
\]
for all sufficiently large $k$
(Bourbaki, \emph{Top.\ g\'{e}n.}, chap.~II,
3\textsuperscript{e}~\'{e}d., \textsection3,
n\textsuperscript{o}~5, th.~1), and the conclusion follows from
\sref{0.13.4.1.3} and \sref{0.13.4.4.1}.
\end{proof}

\subsection{Projective limits of spectral sequences of filtered complexes}
\label{subsection:0.13.5}

\begin{env}[13.5.1]
\label{0.13.5.1}
Let $\cat{C}$ be an abelian category and let $X^\bullet$ be a complex
of objects of $\cat{C}$ equipped with a filtration
$(\mathrm{F}^p(X^\bullet))_{p\in\bb{Z}}$ such that
$\mathrm{F}^{p_0}(X^\bullet)=X^\bullet$ for some index $p_0$.
For every $k\in\bb{Z}$, consider the complex
\[
  X^\bullet_k=X^\bullet/\mathrm{F}^{k+1}(X^\bullet);
\]
it is canonically equipped with the filtration consisting of
\[
  \mathrm{F}^p(X^\bullet_k)
  =
  \mathrm{F}^p(X^\bullet)/\mathrm{F}^{k+1}(X^\bullet)
  \quad\text{for }p\leq k,
\]
and $\mathrm{F}^p(X^\bullet_k)=0$ for $p\geq k+1$.
Moreover, there are canonical morphisms
$X^\bullet_{k+1}\longrightarrow X^\bullet_k$, which make
$\bb{X}^\bullet=(X^\bullet_k)_{k\in\bb{Z}}$ a \emph{projective system
of filtered complexes of objects of $\cat{C}$}.
We note that this projective system is \emph{strict} and that
$X^\bullet_k=0$ for $k<p_0$.
\end{env}

\begin{env}[13.5.2]
\label{0.13.5.2}
More generally, consider a bounded-below strict projective system
$\bb{X}^\bullet=(X^\bullet_k)_{k\in\bb{Z}}$ of complexes of objects of
$\cat{C}$, and equip each $X^\bullet_k$ with the filtration defined in
\sref{0.13.4.1} (applied in the abelian category of bounded-below
complexes of objects of $\cat{C}$).
The morphisms
$X^\bullet_k\longrightarrow X^\bullet_p$ ($p\leq k$) then become
morphisms of \emph{filtered} complexes with finite filtrations.
The functoriality of the spectral sequences of filtered complexes
(\sref{0.11.2.3}) shows that the transition morphisms of
$\bb{X}^\bullet$ induce morphisms making
\[
  \mathrm{E}(\bb{X}^\bullet)
  =
  \bigl(\mathrm{E}(X^\bullet_k)\bigr)_{k\in\bb{Z}}
\]
a \emph{projective system of spectral sequences}.
\end{env}

\begin{lemma}[13.5.3]
\label{0.13.5.3}
Suppose that the projective system
$\bb{X}^\bullet=(X^\bullet_k)_{k\in\bb{Z}}$ of filtered complexes is
obtained as in \sref{0.13.5.2}. Then:
\begin{enumerate}
  \item[{\rm a)}] For $r\geq p-p_0$, one has
  \[
    \mathrm{B}_r^{pq}(X^\bullet_k)
    =
    \mathrm{B}_\infty^{pq}(X^\bullet_k)
  \]
  for every $k\in\bb{Z}$.
  \item[{\rm b)}] For $k+1\leq p+r$, one has
  \[
    \mathrm{Z}_r^{pq}(X^\bullet_k)
    =
    \mathrm{Z}_\infty^{pq}(X^\bullet_k).
  \]
  \item[{\rm c)}] For $k+1\geq p+r$, the morphisms
  \[
    \mathrm{Z}_r^{pq}(X^\bullet_h)
    \longrightarrow
    \mathrm{Z}_r^{pq}(X^\bullet_k)
    \quad\text{and}\quad
    \mathrm{B}_r^{pq}(X^\bullet_h)
    \longrightarrow
    \mathrm{B}_r^{pq}(X^\bullet_k)
  \]
  are isomorphisms for every $h\geq k$.
\end{enumerate}
\end{lemma}

\oldpage[0\textsubscript{III}]{72}
These three properties follow immediately from the definitions in
\sref{0.11.2.2}, taking into account that
$\mathrm{F}^{p-r+1}(X^\bullet_k)=X^\bullet_k$ when $p-r<p_0$.

\begin{env}[13.5.4]
\label{0.13.5.4}
Suppose that the hypotheses of \sref{0.13.5.3} hold.
Then, for fixed $p$, $q$, and $r$ (with $r$ finite), the projective
systems
\[
  \bigl(\mathrm{Z}_r^{pq}(X^\bullet_k)\bigr)_{k\in\bb{Z}},
  \qquad
  \bigl(\mathrm{B}_r^{pq}(X^\bullet_k)\bigr)_{k\in\bb{Z}},
  \qquad
  \bigl(\mathrm{E}_r^{pq}(X^\bullet_k)\bigr)_{k\in\bb{Z}}
\]
are \emph{essentially constant}; their respective projective limits
will be denoted by
$\mathrm{Z}_r^{pq}(\bb{X}^\bullet)$,
$\mathrm{B}_r^{pq}(\bb{X}^\bullet)$, and
\[
  \mathrm{E}_r^{pq}(\bb{X}^\bullet)
  =
  \mathrm{Z}_r^{pq}(\bb{X}^\bullet)/
  \mathrm{B}_r^{pq}(\bb{X}^\bullet).
\]
The objects $\mathrm{Z}_r^{pq}(\bb{X}^\bullet)$ and
$\mathrm{B}_r^{pq}(\bb{X}^\bullet)$ identify canonically with
subobjects of $\mathrm{E}_1^{pq}(\bb{X}^\bullet)$.
The definition of the $d_r^{pq}$ (M, XV, 1) shows that these
morphisms, taken relative to the $X^\bullet_k$, are also essentially
constant, and hence define morphisms
\begin{equation}
\label{0.13.5.4.1}
\tag{13.5.4.1}
  d_r^{pq}:
  \mathrm{E}_r^{pq}(\bb{X}^\bullet)
  \longrightarrow
  \mathrm{E}_r^{p+r,q-r+1}(\bb{X}^\bullet)
\end{equation}
such that
$d_r^{p+r,q-r+1}\circ d_r^{pq}=0$.
Moreover, there are canonical isomorphisms from
$\Ker(d_r^{pq})$ onto
\[
  \mathrm{Z}_{r+1}^{pq}(\bb{X}^\bullet)/
  \mathrm{B}_r^{pq}(\bb{X}^\bullet)
\]
and from $\Im(d_r^{pq})$ onto
\[
  \mathrm{B}_{r+1}^{p+r,q-r+1}(\bb{X}^\bullet)/
  \mathrm{B}_r^{p+r,q-r+1}(\bb{X}^\bullet).
\]
\end{env}

\begin{lemma}[13.5.5]
\label{0.13.5.5}
Under the hypotheses of \sref{0.13.5.3}, for
$s\geq r>p-p_0$ there is a canonical monomorphism
\begin{equation}
\label{0.13.5.5.1}
\tag{13.5.5.1}
  i:
  \mathrm{E}_s^{pq}(\bb{X}^\bullet)
  \longrightarrow
  \mathrm{E}_r^{pq}(\bb{X}^\bullet)
\end{equation}
and a canonical isomorphism
\begin{equation}
\label{0.13.5.5.2}
\tag{13.5.5.2}
  j_r:
  \mathrm{E}_r^{pq}(\bb{X}^\bullet)
  \simto
  \mathrm{E}_\infty^{pq}(X^\bullet_{p+r-1})
\end{equation}
such that the diagram
\begin{equation}
\label{0.13.5.5.3}
\tag{13.5.5.3}
  \xymatrix{
    \mathrm{E}_s^{pq}(\bb{X}^\bullet)
      \ar[r]^{j_s}\ar[d]_i
    &
    \mathrm{E}_\infty^{pq}(X^\bullet_{p+s-1})
      \ar[d]
    \\
    \mathrm{E}_r^{pq}(\bb{X}^\bullet)
      \ar[r]^{j_r}
    &
    \mathrm{E}_\infty^{pq}(X^\bullet_{p+r-1})
  }
\end{equation}
is commutative, where the right vertical arrow is induced by the
morphism
$X^\bullet_{p+s-1}\longrightarrow X^\bullet_{p+r-1}$.
\end{lemma}

\begin{proof}
The existence of $i$ follows from
\[
  \mathrm{B}_r^{pq}(X^\bullet_k)
  =
  \mathrm{B}_\infty^{pq}(X^\bullet_k)
  \qquad (r>p-p_0)
\]
(\sref{0.13.5.3},~a)).
For $k+1\leq p+r$, one has
\[
  \mathrm{Z}_r^{pq}(X^\bullet_k)
  =
  \mathrm{Z}_\infty^{pq}(X^\bullet_k)
\]
(\sref{0.13.5.3},~b)); in particular,
\[
  \mathrm{Z}_r^{pq}(X^\bullet_{p+r-1})
  =
  \mathrm{Z}_\infty^{pq}(X^\bullet_{p+r-1}).
\]
On the other hand,
$\mathrm{Z}_r^{pq}(X^\bullet_{p+r-1})$ and
$\mathrm{B}_r^{pq}(X^\bullet_{p+r-1})$ identify canonically with
$\mathrm{Z}_r^{pq}(\bb{X}^\bullet)$ and
$\mathrm{B}_r^{pq}(\bb{X}^\bullet)$ by
\sref{0.13.5.3},~c).
This proves the existence of $j_r$ and the commutativity of
\sref{0.13.5.5.3}.
\end{proof}

\begin{corollary}[13.5.6]
\label{0.13.5.6}
Under the hypotheses of \sref{0.13.5.3}, if one of the projective
limits
\[
  \varprojlim_r\mathrm{E}_r^{pq}(\bb{X}^\bullet),
  \qquad
  \varprojlim_k\mathrm{E}_\infty^{pq}(X^\bullet_k)
\]
exists, then so does the other, and there is a canonical isomorphism
\begin{equation}
\label{0.13.5.6.1}
\tag{13.5.6.1}
  j_\infty:
  \varprojlim_r\mathrm{E}_r^{pq}(\bb{X}^\bullet)
  \simto
  \varprojlim_k\mathrm{E}_\infty^{pq}(X^\bullet_k).
\end{equation}
Moreover, the projective system
$(\mathrm{E}_r^{pq}(\bb{X}^\bullet))_{r\in\bb{Z}}$ is essentially
constant (\sref{0.13.4.2}) if and only if the projective system
$(\mathrm{E}_\infty^{pq}(X^\bullet_k))_{k\in\bb{Z}}$ is essentially
constant.
\end{corollary}

\begin{env}[13.5.7]
\label{0.13.5.7}
We denote by
$\mathrm{B}_\infty^{pq}(\bb{X}^\bullet)$ and
$\mathrm{Z}_\infty^{pq}(\bb{X}^\bullet)$ the subobjects of
$\mathrm{E}_1^{pq}(\bb{X}^\bullet)$ equal respectively to
$\mathrm{B}_r^{pq}(\bb{X}^\bullet)$ for $r>p-p_0$ and to
$\inf_r\mathrm{Z}_r^{pq}(\bb{X}^\bullet)$, when the latter exists.
Thus
\[
  \varprojlim_r\mathrm{E}_r^{pq}(\bb{X}^\bullet)
\]
\oldpage[0\textsubscript{III}]{73}
identifies canonically with
\[
  \mathrm{E}_\infty^{pq}(\bb{X}^\bullet)
  =
  \mathrm{Z}_\infty^{pq}(\bb{X}^\bullet)/
  \mathrm{B}_\infty^{pq}(\bb{X}^\bullet).
\]
We note that the objects
$\mathrm{Z}_r^{pq}(\bb{X}^\bullet)$,
$\mathrm{B}_r^{pq}(\bb{X}^\bullet)$,
$\mathrm{E}_r^{pq}(\bb{X}^\bullet)$
($1\leq r\leq+\infty$), and the morphisms $d_r^{pq}$ depend
\emph{functorially} on the projective system $\bb{X}^\bullet$, subject
to the restrictions in \sref{0.13.5.5}, and the morphisms defined in
\sref{0.13.5.5} and \sref{0.13.5.6} are functorial.
\end{env}

\subsection{Spectral sequence of a functor relative to an object equipped with a finite filtration}
\label{subsection:0.13.6}

\begin{env}[13.6.1]
\label{0.13.6.1}
Let $\cat{C}$ and $\cat{C}'$ be two abelian categories, and let
$\mathrm{T}:\cat{C}\longrightarrow\cat{C}'$ be an additive covariant
functor.
Suppose that every object of $\cat{C}$ is isomorphic to a subobject of
an injective object, so that the right derived functors
$\RR^p\mathrm{T}$ ($p\geq0$) exist.
\end{env}

\begin{lemma}[13.6.2]
\label{0.13.6.2}
Let $A$ be an object of $\cat{C}$ equipped with a finite filtration
$(\mathrm{F}^i(A))_{i\in\bb{Z}}$.
There exists an injective resolution
$X^\bullet=(X^j)_{j\geq0}$ of $A$, equipped with a finite filtration
$(\mathrm{F}^i(X^\bullet))_{i\in\bb{Z}}$, such that
\[
  \mathrm{F}^i(A)=A
  \quad\text{(respectively, }\mathrm{F}^i(A)=0\text{)}
\]
implies
\[
  \mathrm{F}^i(X^\bullet)=X^\bullet
  \quad\text{(respectively, }\mathrm{F}^i(X^\bullet)=0\text{)},
\]
and such that, for every $i\in\bb{Z}$,
$\mathrm{F}^i(X^\bullet)$ is an injective resolution of
$\mathrm{F}^i(A)$.
\end{lemma}

\begin{proof}
Let $p$ (respectively, $q>p$) be the largest index such that
$\mathrm{F}^p(A)=A$ (respectively, the smallest index such that
$\mathrm{F}^q(A)=0$).
We argue by induction on $q-p$, the lemma being immediate when
$q-p=1$.
Suppose that an injective resolution ${X''}^\bullet$ of
$A/\mathrm{F}^{q-1}(A)$ having the required properties has been
constructed.
Consider the exact sequence
\[
  0\longrightarrow\mathrm{F}^{q-1}(A)
  \longrightarrow A
  \longrightarrow A/\mathrm{F}^{q-1}(A)
  \longrightarrow0,
\]
take an injective resolution ${X'}^\bullet$ of
$\mathrm{F}^{q-1}(A)$, and then choose an injective resolution
$X^\bullet$ of $A$ so as to obtain an exact sequence
\[
  0\longrightarrow {X'}^\bullet
  \longrightarrow X^\bullet
  \longrightarrow {X''}^\bullet
  \longrightarrow0
\]
compatible with the preceding sequence (M, V, 2.2).
It is clear that $X^\bullet$ has the required properties.
\end{proof}

\begin{corollary}[13.6.3]
\label{0.13.6.3}
Let $B$ be a second object of $\cat{C}$, equipped with a finite
filtration $(\mathrm{F}^i(B))_{i\in\bb{Z}}$, let $s$ be an integer,
and let $u:A\longrightarrow B$ be a morphism such that
\[
  u(\mathrm{F}^i(A))\subset\mathrm{F}^{i+s}(B)
\]
for every $i\in\bb{Z}$.
If $Y^\bullet=(Y^j)_{j\geq0}$ is an injective resolution of $B$
equipped with a filtration
$(\mathrm{F}^i(Y^\bullet))_{i\in\bb{Z}}$ having the properties stated
in \sref{0.13.6.2}, then there exists a morphism
$v:X^\bullet\longrightarrow Y^\bullet$, compatible with $u$, such
that
\[
  v(\mathrm{F}^i(X^\bullet))
  \subset
  \mathrm{F}^{i+s}(Y^\bullet)
\]
for every $i\in\bb{Z}$.
Moreover, any two such morphisms $v$ and $v'$ are homotopic.
\end{corollary}

This follows immediately, by induction on $q-p$, from the preceding
construction and from (M, V, 2.3).

\begin{env}[13.6.4]
\label{0.13.6.4}
Under the hypotheses of \sref{0.13.6.2}, consider the complex
$\mathrm{T}(X^\bullet)$ in $\cat{C}'$.
It is filtered by the complexes
$\mathrm{T}(\mathrm{F}^i(X^\bullet))$, since
$\mathrm{F}^i(X^\bullet)$ is a direct summand of $X^\bullet$.
It follows from \sref{0.13.6.3} that the \emph{spectral sequence} of
this filtered complex depends, up to isomorphism, only on the
filtered object $A$.
Its abutment is the cohomology $\RR^\bullet\mathrm{T}(A)$, equipped
with the filtration
\begin{equation}
\label{0.13.6.4.1}
\tag{13.6.4.1}
  \begin{split}
  \mathrm{F}^p(\RR^n\mathrm{T}(A))
  &=
  \Im\bigl(
    \RR^n\mathrm{T}(\mathrm{F}^p(A))
    \longrightarrow
    \RR^n\mathrm{T}(A)
  \bigr)
  \\
  &=
  \Ker\bigl(
    \RR^n\mathrm{T}(A)
    \longrightarrow
    \RR^n\mathrm{T}(A/\mathrm{F}^p(A))
  \bigr),
  \end{split}
\end{equation}
by \sref{0.11.2.2}, and its $\mathrm{E}_1$-term is
\begin{equation}
\label{0.13.6.4.2}
\tag{13.6.4.2}
  \mathrm{E}_1^{pq}
  =
  \RR^{p+q}\mathrm{T}(\mathrm{gr}^p(A)),
\end{equation}
where, as usual,
$\mathrm{gr}^p(A)=\mathrm{F}^p(A)/\mathrm{F}^{p+1}(A)$.
It is clear from \sref{0.11.2.2} that the filtration of the abutment
is finite and that, for fixed $p$ and $q$, the sequences
\oldpage[0\textsubscript{III}]{74}
\[
  \mathrm{B}_r^{pq}(A)
  =
  \mathrm{B}_r^{pq}(\mathrm{T}(X^\bullet))
  \quad\text{and}\quad
  \mathrm{Z}_r^{pq}(A)
  =
  \mathrm{Z}_r^{pq}(\mathrm{T}(X^\bullet))
\]
are stationary.
The preceding spectral sequence is therefore \emph{biregular}
(\sref{0.11.1.3}).
We denote it by
$\mathrm{E}(A)=(\mathrm{E}_r^{pq}(A))$ and call it the
\emph{spectral sequence of the functor $\mathrm{T}$ relative to the
filtered object $A$}.
\end{env}

\begin{env}[13.6.5]
\label{0.13.6.5}
Now suppose that the hypotheses of \sref{0.13.6.3} hold, and retain
its notation.
Since $\mathrm{F}^i(X^\bullet)$ (respectively,
$\mathrm{F}^i(Y^\bullet)$) is a direct summand of $X^\bullet$
(respectively, $Y^\bullet$), one has
\[
  (\mathrm{T}v)\bigl(\mathrm{T}(\mathrm{F}^i(X^\bullet))\bigr)
  \subset
  \mathrm{T}(\mathrm{F}^{i+s}(Y^\bullet))
\]
for every $i\in\bb{Z}$.
The definitions in \sref{0.11.2.2} then show that, for
$1\leq r\leq+\infty$, $\mathrm{T}v$ defines morphisms
\[
  \mathrm{B}_r^{pq}(\mathrm{T}(X^\bullet))
  \longrightarrow
  \mathrm{B}_r^{p+s,q-s}(\mathrm{T}(Y^\bullet))
\]
and
\[
  \mathrm{Z}_r^{pq}(\mathrm{T}(X^\bullet))
  \longrightarrow
  \mathrm{Z}_r^{p+s,q-s}(\mathrm{T}(Y^\bullet)),
\]
and hence a morphism
\[
  w_r:
  \mathrm{E}_r^{pq}(A)
  \longrightarrow
  \mathrm{E}_r^{p+s,q-s}(B).
\]
Similarly, on the abutments one has morphisms
\[
  u_n:
  \RR^n\mathrm{T}(A)
  \longrightarrow
  \RR^n\mathrm{T}(B)
\]
such that
\[
  u_n\bigl(\mathrm{F}^p(\RR^n\mathrm{T}(A))\bigr)
  \subset
  \mathrm{F}^{p+s}(\RR^n\mathrm{T}(B)).
\]
The definition of the $d_r^{pq}$ (M, XV, 1) also shows that the
diagrams
\[
  \xymatrix{
    \mathrm{E}_r^{pq}(A)
      \ar[r]^{d_r^{pq}}\ar[d]_{w_r}
    &
    \mathrm{E}_r^{p+r,q-r+1}(A)
      \ar[d]^{w_r}
    \\
    \mathrm{E}_r^{p+s,q-s}(B)
      \ar[r]_{d_r^{p+s,q-s}}
    &
    \mathrm{E}_r^{p+r+s,q-r-s+1}(B)
  }
\]
are commutative.
It follows that there is an analogous commutative diagram for the
isomorphisms $\alpha_r^{pq}$, which we leave to the reader to make
explicit.
Finally (\emph{loc.\ cit.}), one also has commutative diagrams on the
abutments:
\[
  \xymatrix{
    \mathrm{E}_\infty^{pq}(A)
      \ar[r]^{\beta^{pq}}\ar[d]_{w_\infty}
    &
    \mathrm{gr}^p(\RR^{p+q}\mathrm{T}(A))
      \ar[d]^{u_{p+q}}
    \\
    \mathrm{E}_\infty^{p+s,q-s}(B)
      \ar[r]_{\beta^{p+s,q-s}}
    &
    \mathrm{gr}^{p+s}(\RR^{p+q}\mathrm{T}(B)).
  }
\]
\end{env}

\begin{env}[13.6.6]
\label{0.13.6.6}
Suppose in particular that there is a ring $S$, equipped with a
\emph{filtration} $(\mathrm{F}^i(S))_{i\in\bb{Z}}$, and a ring
homomorphism
\begin{equation}
\label{0.13.6.6.1}
\tag{13.6.6.1}
  h:S\longrightarrow\Hom_{\cat{C}}(A,A)
\end{equation}
\oldpage[0\textsubscript{III}]{75}
such that, for every $t\in\mathrm{F}^i(S)$,
\[
  h_t(\mathrm{F}^j(A))
  \subset
  \mathrm{F}^{i+j}(A)
\]
for every pair $i,j$.
For brevity, we then say that $A$ is equipped with the structure of a
\emph{filtered $S$--$\cat{C}$-module} over the filtered ring $S$.

Passing to the associated graded objects, every $h_t$, for
$t\in\mathrm{F}^i(S)$, defines a graded endomorphism
$\overline{h}_t$ of $\mathrm{gr}^\bullet(A)$, homogeneous of degree
$i$.
Moreover, this endomorphism depends only on the class of $t$ in
$\mathrm{gr}^i(S)$, and one thereby obtains a homomorphism of
\emph{graded rings}
\[
  \overline{h}:
  \mathrm{gr}^\bullet(S)
  \longrightarrow
  \Hom_{\cat{C}}^\bullet
  \bigl(\mathrm{gr}^\bullet(A),\mathrm{gr}^\bullet(A)\bigr),
\]
where the right-hand side is the ring of graded endomorphisms of
$\mathrm{gr}^\bullet(A)$.
We say that $\mathrm{gr}^\bullet(A)$ is equipped with the structure
of a \emph{graded $\mathrm{gr}^\bullet(S)$--$\cat{C}$-module}.

It follows from \sref{0.13.6.5} that, for
$1\leq r\leq+\infty$, every
$\overline{t}\in\mathrm{gr}^i(S)$ canonically defines, on the
bigraded objects
\[
  \bigl(\mathrm{B}_r^{pq}(A)\bigr)_{(p,q)\in\bb{Z}\times\bb{Z}},
  \qquad
  \bigl(\mathrm{Z}_r^{pq}(A)\bigr)_{(p,q)\in\bb{Z}\times\bb{Z}},
\]
and
\[
  \mathrm{E}_r(A)
  =
  \bigl(\mathrm{E}_r^{pq}(A)\bigr)_{(p,q)\in\bb{Z}\times\bb{Z}},
\]
bigraded endomorphisms of degree $(i,-i)$.
On $\mathrm{E}_r(A)$, for finite $r$, this endomorphism commutes with
the bigraded endomorphism defined by the $d_r^{pq}$.
These endomorphisms satisfy the usual associativity and
distributivity conditions with respect to addition in
$\mathrm{gr}^\bullet(S)$ and in the bigraded objects in question.
For brevity, we say that the latter are
\emph{bigraded $\mathrm{gr}^\bullet(S)$--$\cat{C}'$-modules};
it is immediate that the $\alpha_r^{pq}$ define an isomorphism for
this type of structure.

For every integer $n$, let $\mathrm{B}_r^{(n)}(A)$
(respectively, $\mathrm{Z}_r^{(n)}(A)$ and
$\mathrm{E}_r^{(n)}(A)$) denote the graded subobject of
$\mathrm{B}_r^{\bullet,\bullet}(A)$
(respectively, $\mathrm{Z}_r^{\bullet,\bullet}(A)$ and
$\mathrm{E}_r^{\bullet,\bullet}(A)$) consisting of the terms for
which $p+q=n$ ($1\leq r\leq+\infty$).
It is immediate that these are
\emph{graded $\mathrm{gr}^\bullet(S)$--$\cat{C}'$-modules}.

Finally, every $\overline{t}\in\mathrm{gr}^i(S)$ defines, for every
$n$, a graded endomorphism of degree $i$ on the graded object
$\mathrm{gr}^\bullet(\RR^n\mathrm{T}(A))$, which thereby acquires the
structure of a graded
$\mathrm{gr}^\bullet(S)$--$\cat{C}'$-module.
For $p+q=n$, the $\beta^{pq}$ define an isomorphism
\[
  \mathrm{E}^{(n)}(A)
  \isoto
  \mathrm{gr}^\bullet(\RR^n\mathrm{T}(A))
\]
for this type of structure.

We note that, when $\cat{C}'$ is the category of abelian groups, the
structures of an $S$--$\cat{C}'$-module (respectively, of a graded or
bigraded $\mathrm{gr}^\bullet(S)$--$\cat{C}'$-module) are simply the
usual structures of an $S$-module (respectively, a graded or
bigraded $\mathrm{gr}^\bullet(S)$-module).
\end{env}

\subsection{Derived functors of a projective limit of arguments}
\label{subsection:0.13.7}
\label{0.13.7}

\begin{env}[13.7.1]
\label{0.13.7.1}
Let $\cat{C}$ and $\cat{C}'$ be two abelian categories, and suppose
that every object of $\cat{C}$ is isomorphic to a subobject of an
injective object.
Let
\[
  \mathrm{T}:\cat{C}\longrightarrow\cat{C}'
\]
be an additive covariant functor.
Consider a strict projective system
$\bb{A}=(A_k)_{k\in\bb{Z}}$ in $\cat{C}$, bounded below; more
precisely, suppose that $A_k=0$ for $k<k_0$.
To this system we canonically associate, on each $A_k$, the
filtration $(\mathrm{F}^p(A_k))_{p\in\bb{Z}}$ given by
\sref{0.13.4.1.1}.
Since the projective system is strict, the canonical morphisms
\begin{equation}
\label{0.13.7.1.1}
\tag{13.7.1.1}
  \mathrm{F}^i(A_h)/\mathrm{F}^j(A_h)
  \longrightarrow
  \mathrm{F}^i(A_k)/\mathrm{F}^j(A_k)
  \qquad(h\geq k)
\end{equation}
are isomorphisms for $i\leq j\leq k+1$.
Recall also that
$\mathrm{F}^{k_0}(A_k)=A_k$ and
$\mathrm{F}^{k+1}(A_k)=0$ for every $k$.
\end{env}

\begin{env}[13.7.2]
\label{0.13.7.2}
For each $k$, construct an injective resolution
$\mathrm{X}_k^\bullet=(\mathrm{X}_k^j)_{j\geq0}$ of $A_k$ having the
properties of \sref{0.13.6.2}.
The canonical morphisms $A_{k+1}\longrightarrow A_k$ allow us, by
\sref{0.13.6.3}, to define for each $k$ a morphism of complexes
\[
  \mathrm{X}_{k+1}^\bullet\longrightarrow\mathrm{X}_k^\bullet
\]
compatible with the filtrations and making
$\bb{X}^\bullet=(\mathrm{X}_k^\bullet)_{k\in\bb{Z}}$ a projective
system of complexes.
\oldpage[0\textsubscript{III}]{76}
We may moreover suppose that this projective system is strict.
Indeed, by \sref{0.13.7.1.1},
\[
  A_k
  \simeq
  A_{k+1}/\mathrm{F}^{k+1}(A_{k+1}).
\]
Thus, in constructing $\mathrm{X}_{k+1}^\bullet$, we may take the
injective resolution of
$A_{k+1}/\mathrm{F}^{k+1}(A_{k+1})$ to be
\emph{equal} to $\mathrm{X}_k^\bullet$.
It then follows from (M, V, 2.3) that the morphism of complexes
$\mathrm{X}_{k+1}^\bullet\longrightarrow\mathrm{X}_k^\bullet$ may be
constructed so that, on passing to quotients, it induces an
isomorphism
\[
  \mathrm{X}_{k+1}^\bullet/
  \mathrm{F}^{k+1}(\mathrm{X}_{k+1}^\bullet)
  \simto
  \mathrm{X}_k^\bullet
\]
respecting the filtrations; this is precisely the condition of
\sref{0.13.5.1}.
\end{env}

\begin{env}[13.7.3]
\label{0.13.7.3}
By construction, the projective system
$\mathrm{T}(\bb{X}^\bullet)$ of complexes
$\mathrm{T}(\mathrm{X}_k^\bullet)$ satisfies the hypotheses of
\sref{0.13.5.3}.
The results of \sref{0.13.5.4}, \sref{0.13.5.5}, and
\sref{0.13.5.6} therefore apply to the spectral sequences
\[
  \mathrm{E}\bigl(\mathrm{T}(\mathrm{X}_k^\bullet)\bigr)
  =
  \mathrm{E}(A_k).
\]
For $1\leq r\leq+\infty$, we write
$\mathrm{E}_r^{pq}(\bb{A})$ in place of
$\mathrm{E}_r^{pq}(\mathrm{T}(\bb{X}^\bullet))$
(cf.\ \sref{0.13.5.7} when $r=+\infty$), and use the same convention
for the analogous notation.
In particular,
\begin{equation}
\label{0.13.7.3.1}
\tag{13.7.3.1}
  \mathrm{E}_1^{pq}(\bb{A})
  =
  \RR^{p+q}\mathrm{T}\bigl(\mathrm{gr}^p(\bb{A})\bigr),
\end{equation}
by \sref{0.13.6.4.2} and the fact that the projective system
$(\mathrm{gr}^p(A_k))_{k\in\bb{Z}}$ is essentially constant.
These results, together with \sref{0.13.4.3}, give the following
proposition, first proved by Shih Weishu by a different method
(unpublished).
\end{env}

\begin{proposition}[13.7.4]
\label{0.13.7.4}
\emph{(Shih).}
Let $n$ be an integer.
The following two conditions are equivalent:
\begin{enumerate}
  \item[{\rm a)}] For every pair $(p,q)$ such that $p+q=n$, the
  projective system
  $(\mathrm{E}_r^{pq}(\bb{A}))_{r\geq2}$ is essentially constant.
  \item[{\rm b)}] The projective system
  \[
    \RR^n\mathrm{T}(\bb{A})
    =
    \bigl(\RR^n\mathrm{T}(A_k)\bigr)_{k\in\bb{Z}}
  \]
  satisfies {\rm(ML)}.
\end{enumerate}
Moreover, when these conditions hold, there is a canonical
isomorphism
\begin{equation}
\label{0.13.7.4.1}
\tag{13.7.4.1}
  \mathrm{gr}^p\bigl(\RR^n\mathrm{T}(\bb{A})\bigr)
  \isoto
  \mathrm{E}_\infty^{p,n-p}(\bb{A})
  \qquad(p\in\bb{Z}).
\end{equation}
\end{proposition}

\begin{proof}
By \sref{0.13.5.6}, condition {\rm a)} says precisely that, for
$p+q=n$, the projective system
$(\mathrm{E}_\infty^{pq}(A_k))_{k\in\bb{Z}}$ is essentially
constant.
On the other hand,
$\mathrm{gr}^p(\RR^n\mathrm{T}(A_k))$ is canonically isomorphic to
$\mathrm{E}_\infty^{p,n-p}(A_k)$.
It follows from \sref{0.13.4.3} that {\rm a)} and {\rm b)} are
equivalent, and \sref{0.13.7.4.1} is just
\sref{0.13.5.6.1} in the case under consideration.
\end{proof}

\begin{corollary}[13.7.5]
\label{0.13.7.5}
Let $\sh{F}=(\sh{F}_k)_{k\in\bb{N}}$ be a projective system of
sheaves of abelian groups satisfying conditions {\rm(i)}, {\rm(ii)},
and {\rm(iii)} of \sref{0.13.3.1}, and put
$\sh{F}=\varprojlim_k\sh{F}_k$.
Suppose that, for the functor
$\sh{G}\longmapsto\Gamma(X,\sh{G})$, the projective system
$(\mathrm{E}_r^{pq}(\sh{F}))_r$ is essentially constant whenever
$p+q=n$ or $p+q=n+1$.
On $\HH^{n+1}(X,\sh{F})$, consider the filtration
\[
  \mathrm{F}^p\bigl(\HH^{n+1}(X,\sh{F})\bigr)
  =
  \Ker\bigl(
    \HH^{n+1}(X,\sh{F})
    \longrightarrow
    \HH^{n+1}(X,\sh{F}_{p-1})
  \bigr).
\]
Then there is a canonical isomorphism
\begin{equation}
\label{0.13.7.5.1}
\tag{13.7.5.1}
  \mathrm{gr}^p\bigl(\HH^{n+1}(X,\sh{F})\bigr)
  \isoto
  \mathrm{E}_\infty^{p,n-p+1}(\sh{F})
  \qquad(p\in\bb{Z}).
\end{equation}
\end{corollary}

\begin{proof}
\oldpage[0\textsubscript{III}]{77}
Apply \sref{0.13.7.4} with $\cat{C}$ the category of sheaves of
abelian groups on $X$, $\cat{C}'$ the category of abelian groups,
and $\mathrm{T}=\Gamma$.
It gives, for every $p\in\bb{Z}$, a canonical isomorphism
\[
  \mathrm{gr}^p\bigl(\RR^{n+1}\Gamma(\sh{F})\bigr)
  \isoto
  \mathrm{E}_\infty^{p,n-p+1}(\sh{F}).
\]
Moreover, by \sref{0.13.7.4}, the projective system
$(\HH^n(X,\sh{F}_k))_{k\in\bb{Z}}$ satisfies {\rm(ML)}; hence
\sref{0.13.3.1} gives a canonical isomorphism
\[
\label{0.13.7.5.1b}
  \HH^{n+1}(X,\sh{F})
  \isoto
  \varprojlim_k\HH^{n+1}(X,\sh{F}_k).
\]

Since the projective system
$\RR^{n+1}\Gamma(\sh{F})$ satisfies {\rm(ML)} by
\sref{0.13.7.4}, there is a canonical isomorphism
\[
  \mathrm{gr}^p\bigl(\RR^{n+1}\Gamma(\sh{F})\bigr)
  \isoto
  \varprojlim_k
  \mathrm{gr}^p\bigl(\HH^{n+1}(X,\sh{F}_k)\bigr)
\]
by \sref{0.13.4.3}, and a canonical isomorphism
\[
  \varprojlim_k
  \mathrm{gr}^p\bigl(\HH^{n+1}(X,\sh{F}_k)\bigr)
  \isoto
  \mathrm{gr}^p\left(
    \varprojlim_k\HH^{n+1}(X,\sh{F}_k)
  \right)
\]
by \sref{0.13.4.5}.
It remains only to check that the preceding isomorphism is
compatible with the filtrations on its two sides.
This follows immediately from the definitions and, for every $p$,
from the commutativity of
\[
  \xymatrix{
    \HH^{n+1}(X,\sh{F})
      \ar[r]^-{\sim}\ar[dr]
    &
    \varprojlim_k\HH^{n+1}(X,\sh{F}_k)
      \ar[d]
    \\
    &
    \HH^{n+1}(X,\sh{F}_{p-1}).
  }
\]
\end{proof}

\begin{env}[13.7.6]
\label{0.13.7.6}
Let $S$ be a ring equipped with a filtration
$(\mathrm{F}^i(S))_{i\in\bb{Z}}$ such that
$\mathrm{F}^0(S)=S$ (and hence
$\mathrm{gr}^i(S)=0$ for $i<0$).
Suppose that each $A_k$, with the filtration defined in
\sref{0.13.7.1}, is a filtered
$S$--$\cat{C}$-module in the sense of \sref{0.13.6.6}, and that the
morphisms $A_h\longrightarrow A_k$ for $k\leq h$ are morphisms of
filtered $S$--$\cat{C}$-modules.
For brevity, we then say that $\bb{A}$ is a
\emph{projective system of filtered $S$--$\cat{C}$-modules}.

It is immediate that, for $k\leq h$ and
$1\leq r\leq+\infty$, the morphisms
\[
  \mathrm{B}_r^{pq}(A_h)\longrightarrow\mathrm{B}_r^{pq}(A_k),
  \qquad
  \mathrm{Z}_r^{pq}(A_h)\longrightarrow\mathrm{Z}_r^{pq}(A_k)
\]
are morphisms of bigraded
$\mathrm{gr}^\bullet(S)$--$\cat{C}'$-modules
(\sref{0.13.6.5}), and that, for finite $r$, the families
\[
  \bigl(\mathrm{Z}_r^{pq}(\bb{A})\bigr),\qquad
  \bigl(\mathrm{B}_r^{pq}(\bb{A})\bigr),\qquad
  \bigl(\mathrm{E}_r^{pq}(\bb{A})\bigr)
\]
are bigraded
$\mathrm{gr}^\bullet(S)$--$\cat{C}'$-modules, the first two being
submodules of $(\mathrm{E}_1^{pq}(\bb{A}))$.
We again write
$\mathrm{Z}_r^{(n)}(\bb{A})$,
$\mathrm{B}_r^{(n)}(\bb{A})$, and
$\mathrm{E}_r^{(n)}(\bb{A})$ for the respective graded subobjects
obtained by retaining only the terms with $p+q=n$; these are graded
$\mathrm{gr}^\bullet(S)$--$\cat{C}'$-modules.

When the projective system
$(\mathrm{E}_r^{pq}(\bb{A}))_r$ is essentially constant,
$\mathrm{E}_\infty^{pq}(\bb{A})$ is likewise a bigraded
$\mathrm{gr}^\bullet(S)$--$\cat{C}'$-module, and each
$\mathrm{E}_\infty^{(n)}(\bb{A})$ is a graded
$\mathrm{gr}^\bullet(S)$--$\cat{C}'$-module.
Moreover, for each $k$, the isomorphisms
\[
  \beta^{p,n-p}:
  \mathrm{E}_\infty^{p,n-p}(A_k)
  \isoto
  \mathrm{gr}^p\bigl(\RR^n\mathrm{T}(A_k)\bigr)
\]
form an isomorphism of graded
$\mathrm{gr}^\bullet(S)$--$\cat{C}'$-modules from
$\mathrm{E}_\infty^{(n)}(A_k)$ onto
$\mathrm{gr}^\bullet(\RR^n\mathrm{T}(A_k))$.
Under the preceding hypotheses, the isomorphism
\[
  \beta^{p,n-p}:
  \mathrm{E}_\infty^{(n)}(\bb{A})
  \isoto
  \varprojlim_k
  \mathrm{gr}^\bullet\bigl(\RR^n\mathrm{T}(A_k)\bigr)
\]
therefore also respects these structures.
The same is plainly true of the canonical isomorphism
\[
  \mathrm{gr}^\bullet\bigl(\RR^n\mathrm{T}(\bb{A})\bigr)
  \isoto
  \varprojlim_k
  \mathrm{gr}^\bullet\bigl(\RR^n\mathrm{T}(A_k)\bigr);
\]
consequently, the isomorphisms \sref{0.13.7.4.1} form an
isomorphism of graded
$\mathrm{gr}^\bullet(S)$--$\cat{C}'$-modules.
\end{env}

\begin{proposition}[13.7.7]
\label{0.13.7.7}
Let $S$ be a Noetherian $\mathfrak{J}$-adic ring.
Suppose that $\cat{C}$ is an abelian category in which every object
is isomorphic to a subobject of an injective object, and let
$\mathrm{T}$ be an additive covariant functor from $\cat{C}$ to the
category of abelian groups.
Let $\bb{A}=(A_k)_{k\in\bb{Z}}$ be a
\oldpage[0\textsubscript{III}]{78}
strict projective system of filtered $S$--$\cat{C}$-modules, for the $\mathfrak{J}$-adic
filtration on $S$, bounded below.
For a fixed integer $n$, suppose that the following condition holds:
\begin{itemize}
  \item[$(\mathrm{F}_n)$]
  \emph{The graded $\mathrm{gr}^\bullet(S)$-module}
  \[
    \mathrm{E}_1^{(m)}(\bb{A})
    =
    \bigl(
      \RR^m\mathrm{T}(\mathrm{gr}^p(\bb{A}))
    \bigr)_{p\in\bb{Z}}
  \]
  \emph{of \sref{0.13.7.3.1} is of finite type for
  $m=n$ and $m=n+1$.}
\end{itemize}
Then:
\begin{enumerate}
  \item[{\rm(i)}] The projective systems
  $(\RR^n\mathrm{T}(A_k))_{k\in\bb{Z}}$ and
  $(\RR^{n+1}\mathrm{T}(A_k))_{k\in\bb{Z}}$ satisfy {\rm(ML)}.
  \item[{\rm(ii)}] If
  \[
    \RR'^{\,n}\mathrm{T}(\bb{A})
    =
    \varprojlim_k\RR^n\mathrm{T}(A_k),
  \]
  then $\RR'^{\,n}\mathrm{T}(\bb{A})$ is an $S$-module of finite
  type.
  \item[{\rm(iii)}] The filtration on
  $\RR'^{\,n}\mathrm{T}(\bb{A})$ defined by
  \[
    \mathrm{F}^p\bigl(\RR'^{\,n}\mathrm{T}(\bb{A})\bigr)
    =
    \Ker\bigl(
      \RR'^{\,n}\mathrm{T}(\bb{A})
      \longrightarrow
      \RR^n\mathrm{T}(A_{p-1})
    \bigr)
    \qquad(p\in\bb{Z})
  \]
  is $\mathfrak{J}$-good; that is,
  \[
    \mathfrak{J}\,
    \mathrm{F}^p\bigl(\RR'^{\,n}\mathrm{T}(\bb{A})\bigr)
    \subset
    \mathrm{F}^{p+1}\bigl(\RR'^{\,n}\mathrm{T}(\bb{A})\bigr)
  \]
  for every $p$, with equality for all sufficiently large $p$.
  In particular, the topology on
  $\RR'^{\,n}\mathrm{T}(\bb{A})$ defined by this filtration is the
  $\mathfrak{J}$-adic topology.
  \item[{\rm(iv)}] The projective system
  $(\mathrm{E}_r^{pq}(\bb{A}))_r$ is essentially constant whenever
  $p+q=n$ or $p+q=n+1$.
  Thus $\mathrm{E}_\infty^{pq}(\bb{A})$ is defined
  (\sref{0.13.5.7}), and the isomorphisms
  \begin{equation}
  \label{0.13.7.7.1}
  \tag{13.7.7.1}
    \mathrm{gr}^p\bigl(\RR'^{\,n}\mathrm{T}(\bb{A})\bigr)
    \isoto
    \mathrm{E}_\infty^{p,n-p}(\bb{A})
    \qquad(p\in\bb{Z})
  \end{equation}
  form an isomorphism of graded
  $\mathrm{gr}^\bullet(S)$-modules.
\end{enumerate}
\end{proposition}

\begin{proof}
In view of \sref{0.13.7.7.1}, and taking account of the
isomorphisms \sref{0.13.7.4.1}, we shall by abuse of notation also
write $\RR^n\mathrm{T}(\bb{A})$ for the projective limit
$\RR'^{\,n}\mathrm{T}(\bb{A})$ of the projective system
$(\RR^n\mathrm{T}(A_k))$.

The graded ring $\mathrm{gr}^\bullet(S)$ is Noetherian
(Bourbaki, \emph{Alg.\ comm.}, chap.~III, \textsection2,
n\textsuperscript{o}~9, cor.~5 of th.~2).
Hence the increasing sequence of graded
$\mathrm{gr}^\bullet(S)$-submodules
$\mathrm{B}_r^{(m)}(\bb{A})$ of
$\mathrm{E}_1^{(m)}(\bb{A})$ is stationary for $m=n$ and
$m=n+1$.
Condition {\rm(b)} of \sref{0.11.1.10} is therefore satisfied.
It follows that condition {\rm(a)} of \sref{0.13.7.4} holds for
$n$ and for $n+1$, proving {\rm(i)}.

Furthermore, the isomorphisms \sref{0.13.7.4.1}, together with the
remarks of \sref{0.13.7.6}, show that
$\mathrm{gr}^\bullet(\RR^n\mathrm{T}(\bb{A}))$ is a graded
$\mathrm{gr}^\bullet(S)$-module isomorphic to
\[
  \mathrm{E}_\infty^{(n)}(\bb{A})
  =
  \mathrm{Z}_\infty^{(n)}(\bb{A})
  /
  \mathrm{B}_\infty^{(n)}(\bb{A}).
\]
Since $\mathrm{Z}_\infty^{(n)}(\bb{A})$ is a submodule of
$\mathrm{E}_1^{(n)}(\bb{A})$, it is of finite type, and so is
$\mathrm{E}_\infty^{(n)}(\bb{A})$.
Finally, for the filtration
$(\mathrm{F}^p(\RR'^{\,n}\mathrm{T}(\bb{A})))$, it follows from
\sref{0.13.4.5} that
\[
  \mathrm{gr}^\bullet\bigl(\RR^n\mathrm{T}(\bb{A})\bigr)
  \quad\hbox{and}\quad
  \mathrm{gr}^\bullet\bigl(\RR'^{\,n}\mathrm{T}(\bb{A})\bigr)
\]
are isomorphic graded $\mathrm{gr}^\bullet(S)$-modules.
This proves {\rm(iv)}.
Assertions {\rm(ii)} and {\rm(iii)} will follow from these results
and the following lemma.
\end{proof}

\begin{lemma}[13.7.7.2]
\label{0.13.7.7.2}
Let $S$ be a Noetherian $\mathfrak{J}$-adic ring, and let $M$
be an $S$-module equipped with a codiscrete filtration
$(\mathrm{F}^p(M))_{p\in\bb{Z}}$ such that
\[
  \mathfrak{J}\mathrm{F}^p(M)
  \subset
  \mathrm{F}^{p+1}(M).
\]
Thus $M$ is a filtered module over $S$ endowed with its
$\mathfrak{J}$-adic filtration.
Suppose moreover that $M$ is separated for the topology defined by
$(\mathrm{F}^p(M))$.
Then the following conditions are equivalent:
\begin{enumerate}
  \item[{\rm a)}] $M$ is an $S$-module of finite type and
  $(\mathrm{F}^p(M))$ is a $\mathfrak{J}$-good filtration.
  \item[{\rm b)}] $\mathrm{gr}^\bullet(M)$ is a
  $\mathrm{gr}^\bullet(S)$-module of finite type.
  \item[{\rm c)}] The $\mathrm{gr}^p(M)$ are $S$-modules of
  finite type and, for all sufficiently large $p$, the canonical
  homomorphisms
  \begin{equation}
  \label{0.13.7.7.3}
  \tag{13.7.7.3}
    \mathfrak{J}\otimes_S\mathrm{gr}^p(M)
    \longrightarrow
    \mathrm{gr}^{p+1}(M)
  \end{equation}
  \oldpage[0\textsubscript{III}]{79}
  are surjective.
  They are induced by
  $\mathfrak{J}\otimes_S\mathrm{F}^p(M)
  \longrightarrow\mathrm{F}^{p+1}(M)$, taking account of the fact
  that the image of the composite homomorphism
  \[
    \mathfrak{J}\otimes_S\mathrm{F}^{p+1}(M)
    \longrightarrow
    \mathfrak{J}\otimes_S\mathrm{F}^p(M)
    \longrightarrow
    \mathrm{F}^{p+1}(M)
  \]
  is
  $\mathfrak{J}\mathrm{F}^{p+1}(M)
  \subset\mathrm{F}^{p+2}(M)$.
\end{enumerate}

For the proof, see Bourbaki, \emph{Alg.\ comm.}, chap.~III,
\textsection3, n\textsuperscript{o}~1, prop.~3.
\end{lemma}

\begin{env}[13.7.7.4]
\label{0.13.7.7.4}
To apply the lemma \sref{0.13.7.7.2}, it remains to observe that the
topology on $\RR'^{\,n}\mathrm{T}(\bb{A})$ defined by the filtration
under consideration makes it a separated and complete $S$-module:
this is the projective-limit topology of the discrete groups
$\RR^n\mathrm{T}(A_k)$.
On the other hand, if $A_k=0$ for $k<k_0$, then also
$\RR^n\mathrm{T}(A_k)=0$ for $k<k_0$; hence
\[
  \mathrm{F}^{k_0}\bigl(\RR'^{\,n}\mathrm{T}(\bb{A})\bigr)
  =
  \RR'^{\,n}\mathrm{T}(\bb{A}).
\]
Thus all the hypotheses of the lemma are indeed satisfied.
\end{env}

\begin{corollary}[13.7.8]
\label{0.13.7.8}
If hypothesis $(\mathrm{F}_n)$ holds, then for every $f\in S$ there
is a canonical isomorphism
\begin{equation}
\label{0.13.7.8.1}
\tag{13.7.8.1}
  \varprojlim_k\bigl(\RR^n\mathrm{T}(A_k)\bigr)_f
  \isoto
  \RR^n\mathrm{T}(\bb{A})\otimes_S S_{\{f\}}.
\end{equation}
\end{corollary}

\begin{proof}
Indeed, $\RR^n\mathrm{T}(\bb{A})$ is an $S$-module of finite type,
and $S_{\{f\}}$ is a Noetherian adic $S$-algebra
(\sref[0\textsubscript{I}]{0.7.6.11}), namely the separated
completion of $S_f$ for the $\mathfrak{J}$-preadic topology
(\sref[0\textsubscript{I}]{0.7.6.2}).
It follows from \sref[0\textsubscript{I}]{0.7.7.8} and
\sref[0\textsubscript{I}]{0.7.7.1} that
\[
  \RR^n\mathrm{T}(\bb{A})\otimes_S S_{\{f\}}
\]
is isomorphic to the separated completion of
$\RR^n\mathrm{T}(\bb{A})\otimes_S S_f$ for the
$\mathfrak{J}$-preadic topology.
A fundamental system of neighbourhoods of zero for this topology
is
\[
  \bigl(
    \mathfrak{J}^p\RR^n\mathrm{T}(\bb{A})
  \bigr)\otimes_S S_f;
\]
the groups
\[
  \mathrm{F}^p\bigl(\RR^n\mathrm{T}(\bb{A})\bigr)
  \otimes_S S_f
\]
therefore also form such a system.
The latter group is the kernel of the canonical map
\[
  \bigl(\RR^n\mathrm{T}(\bb{A})\bigr)_f
  \longrightarrow
  \bigl(\RR^n\mathrm{T}(A_{p-1})\bigr)_f.
\]
Consequently, the separated group associated with
$\RR^n\mathrm{T}(\bb{A})\otimes_S S_f$ identifies with a subgroup
$G$ of
\[
  \varprojlim_k\bigl(\RR^n\mathrm{T}(A_k)\bigr)_f.
\]
But the projective system
$((\RR^n\mathrm{T}(A_k))_f)_k$ plainly satisfies {\rm(ML)}, and the
image of $(\RR^n\mathrm{T}(\bb{A}))_f$ in each
$(\RR^n\mathrm{T}(A_k))_f$ equals the common image of the
$(\RR^n\mathrm{T}(A_h))_f$ for all sufficiently large $h\geq k$.
It follows at once that $G$ is \emph{everywhere dense} in
$\varprojlim_k(\RR^n\mathrm{T}(A_k))_f$.
Since the latter group is separated and complete, the corollary is
proved.
\end{proof}


\clearpage
\phantomsection
\label{section:ega3}
\pdfbookmark[0]{Chapter III}{chapter.III}
\addcontentsline{toc}{part}{Chapter III: Cohomological study of coherent sheaves}
\section*{CHAPTER III}
\setcounter{section}{0}
\setcounter{subsection}{0}
\setcounter{equation}{0}
\begin{center}
\Large\bfseries COHOMOLOGICAL STUDY OF COHERENT SHEAVES
\end{center}
\oldpage[III]{81}

\section*{Summary}

\begin{longtable}{@{}p{0.08\textwidth}p{0.86\textwidth}@{}}
  \hyperref[section:III.1]{\textsection1}. & Cohomology of affine schemes.\\
  \hyperref[section:III.2]{\textsection2}. & Cohomological study of projective morphisms.\\
  \hyperref[section:III.3]{\textsection3}. & Finiteness theorem for proper morphisms.\\
  \hyperref[section:III.4]{\textsection4}. & The fundamental theorem of proper morphisms. Applications.\\
  \hyperref[section:III.5]{\textsection5}. & An existence theorem for coherent algebraic sheaves.\\
  \textsection8. & The duality theorem for projective bundles\\
  \textsection9. & Relative cohomology and local cohomology; local duality\\
  \textsection10. & Relations between projective cohomology and local cohomology. Formal completion technique along a divisor\\
  \textsection11. & Global and local Picard groups\footnotemark
\end{longtable}
\footnotetext{EGA IV does not depend on \textsection\textsection8--11, and will probably be published before these chapters. \emph{[Trans.] These last four chapters were never published.}}
\bigskip

This chapter gives the fundamental theorems concerning the cohomology of coherent algebraic sheaves, with the exception of theorems explaining the theory of residues (duality theorems), which will be the subject of a later chapter.
These results will prove to be essential tools in all that follows, even in questions which are not truly cohomological in their nature;
the reader will see the first such examples starting from \hyperref[section:III.4]{\textsection4}.
Finally, in \textsection\textsection8--11, we will develop certain results, related to the duality of coherent sheaves, that are particularly important for applications, and which can be explained even before the introduction of the full general theory of residues.

\oldpage[III]{82}
The content of \textsection\hyperref[section:III.1]{\textsection1} and \hyperref[section:III.2]{2} is due to J.-P.~Serre, and the reader will observe that we have had only to follow (FAC).
\textsection8 and 9 are equally inspired by (FAC) (the changes necessitated by the different contexts, however, being less evident).
Finally, as we said in the Introduction, \hyperref[section:III.4]{\textsection4} should be considered as the formalisation, in modern language, of the fundamental ``invariance theorem'' of Zariski's ``theory of holomorphic functions''.

We draw attention to the fact that the results of n\textsuperscript{o}\hyperref[subsection:III.3.4]{3.4} (and the preliminary propositions of \sref[0]{0.13.4} to \sref[0]{0.13.7}) will not be used in what follows \hyperref[section:ega3]{Chapter~III}, and can thus be skipped in a first reading.
\bigskip

\section{Cohomology of affine schemes}
\label{section:III.1}

\subsection{Review of the exterior algebra complex}
\label{subsection:III.1.1}

\begin{env}[1.1.1]
\label{III.1.1.1}
Let $A$ be a ring, $\mathbf{f}=(f_i)_{1\leq i\leq r}$ a system of $r$ elements of $A$.
The \emph{exterior algebra complex $K_\bullet(\mathbf{f})$} corresponding to $\mathbf{f}$ is a chain complex (G, I, 2.2) defined in the following way: the graded $A$-module $K_\bullet(\mathbf{f})$ is equal to the \emph{exterior algebra $\wedge(A^r)$}, graded in the usual way, and the boundary map is the \emph{interior multiplication $i_\mathbf{f}$} by $\mathbf{f}$ considered as an element of the dual $\dual{(A^r)}$; we recall that $i_\mathbf{f}$ is an \emph{antiderivation} of degree $-1$ of $\wedge(A^r)$, and if $(\mathbf{e}_i)_{1\leq i\leq r}$ is the canonical basis of $A^r$, then we have $i_\mathbf{f}(\mathbf{e}_i)=f_i$; the verification of the condition $i_\mathbf{f}\circ i_\mathbf{f}=0$ is immediate.

An equivalent definition is the following: for each $i$, we consider a chain complex $K_\bullet(f_i)$ defined as follows: $K_0(f_i)=K_1(f_i)=A$, $K_n(f_i)=0$ for $n\neq 0,1$: the boundary map is defined by the condition that $d_1:A\to A$ is \emph{multiplication by $f_i$}.
We then take $K_\bullet(\mathbf{f})$ to be the \emph{tensor product $K_\bullet(f_1)\otimes K_\bullet(f_2)\otimes\cdots\otimes K_\bullet(f_r)$} (G, I, 2.7) with its total degree; the verification of the isomorphism from this complex to the complex defined above is immediate.
\end{env}

\begin{env}[1.1.2]
\label{III.1.1.2}
For every $A$-module $M$, we define the \emph{chain complex}
\[
\label{III.1.1.2.1}
  K_\bullet(\mathbf{f},M)=K_\bullet(\mathbf{f})\otimes_A M
  \tag{1.1.2.1}
\]
and the \emph{cochain complex} (G, I, 2.2)
\[
\label{III.1.1.2.2}
  K^\bullet(\mathbf{f},M)=\Hom_A(K_\bullet(\mathbf{f}),M).
  \tag{1.1.2.2}
\]

If $g$ is a $k$-cochain of this latter complex, and if we set
\[
  g(i_1,\dots,i_k)=g(\mathbf{e}_{i_1}\wedge\cdots\wedge\mathbf{e}_{i_k}),
\]
then $g$ identifies with an \emph{alternating} map from $[1,r]^k$ to $M$, and it follows from the above definitions that we have
\[
\label{III.1.1.2.3}
  d^k g(i_1,i_2,\dots,i_{k+1})=\sum_{h=1}^{k+1}(-1)^{h-1}f_{i_h}g(i_1,\dots,\widehat{i_h},\dots,i_{k+1}).
  \tag{1.1.2.3}
\]
\end{env}

\begin{env}[1.1.3]
\label{III.1.1.3}
\oldpage[III]{83}
From the above complexes, we deduce as usual the \emph{homology and cohomology $A$-modules} (G, I, 2.2)
\[
  \HH_\bullet(\mathbf{f},M)=\HH_\bullet(K_\bullet(\mathbf{f},M)),
  \tag{1.1.3.1}
\]
\[
  \HH^\bullet(\mathbf{f},M)=\HH^\bullet(K^\bullet(\mathbf{f},M)).
  \tag{1.1.3.2}
\]

We define an \emph{$A$-isomorphism $K_i(\mathbf{f},M)\isoto K^{r-i}(\mathbf{f},M)$} by sending each chain $z=\sum(\mathbf{e}_{i_1}\wedge\cdots\wedge\mathbf{e}_{i_k})\otimes z_{i_1,\dots,i_k}$ to the cochain $g_z$ such that $g_z(j_1,\dots,j_{r-k})=\varepsilon z_{i_1,\dots,i_k}$, where $(j_h)_{1\leq h\leq r-k}$ is the strictly increasing sequence complementary to the strictly increasing sequence $(i_h)_{1\leq h\leq k}$ in $[1,r]$ and $\varepsilon=(-1)^\nu$, where $\nu$ is the number of inversions of the permutation $i_1,\dots,i_k,j_1,\dots,j_{r-k}$ of $[1,r]$.
We verify that $g_{dz}=d(g_z)$, which gives an isomorphism
\[
  \HH^i(\mathbf{f},M)\isoto\HH_{r-i}(\mathbf{f},M)\text{ for }0\leq i\leq r.
  \tag{1.1.3.3}\label{III.1.1.3.3}
\]

In this chapter, we will especially consider the cohomology modules $\HH^\bullet(\mathbf{f},M)$.

For a given $\mathbf{f}$, it is immediate (G, I, 2.1) that $M\mapsto\HH^\bullet(\mathbf{f},M)$ is a \emph{cohomological functor} (T, II, 2.1) from the category of $A$-modules to the category of graded $A$-modules, zero in degrees $<0$ and $>r$.
In addition, we have
\[
  \HH^0(\mathbf{f},M)=\Hom_A(A/(\mathbf{f}),M),
  \tag{1.1.3.4}
\]
denoting by $(\mathbf{f})$ the ideal of $A$ generated by $f_1,\dots,f_r$; this follows immediately from (\hyperref[III.1.1.2.3]{1.1.2.3}), and it is clear that $\HH^0(\mathbf{f},M)$ identifies with the submodule of $M$ \emph{killed by $(\mathbf{f})$}.
Similarly, we have by (\hyperref[III.1.1.2.3]{1.1.2.3}) that
\[
  \HH^r(\mathbf{f},M)=M/\bigg(\sum_{i=1}^r f_i M\bigg)=(A/(\mathbf{f}))\otimes_A M.
  \tag{1.1.3.5}\label{III.1.1.3.5}
\]

We will use the following known result, which we will recall a proof of to be complete:
\end{env}

\begin{proposition}[1.1.4]
\label{III.1.1.4}
Let $A$ be a ring, $\mathbf{f}=(f_i)_{1\leq i\leq r}$ a finite family of elements of $A$, and $M$ an $A$-module.
If, for $1\leq i\leq r$, the scaling $z\mapsto f_i\cdot z$ on $M_{i-1}=M/(f_1 M+\cdots+f_{i-1}M)$ is injective, then we have $\HH^i(\mathbf{f},M)=0$ for $i\neq r$.
\end{proposition}

\begin{proof}
It suffices to prove that $\HH_i(\mathbf{f},M)=0$ for all $i>0$ according to (\hyperref[III.1.1.3.3]{1.1.3.3}).
We argue by induction on $r$, the case $r=0$ being trivial.
Set $\mathbf{f}'=(f_i)_{1\leq i\leq r-1}$; this family satisfies the conditions in the statement, so if we set $L_\bullet=K_\bullet(\mathbf{f}',M)$, then we have $\HH_i(L_\bullet)=0$ for $i>0$ by hypothesis, and $\HH_0(L_\bullet)=M_{r-1}$ by virtue of (\hyperref[III.1.1.3.3]{1.1.3.3}) and (\hyperref[III.1.1.3.5]{1.1.3.5}).
To abbreviate, set $K_\bullet=K_\bullet(f_r)=K_0\oplus K_1$, with $K_0=K_1=A$, $d_1:K_1\to K_0$ multiplication by $f_r$; we have by definition \sref{III.1.1.1} that $K_\bullet(\mathbf{f},M)=K_\bullet\otimes_A L_\bullet$.
We have the following lemma:

\begin{lemma}[1.1.4.1]
\label{III.1.1.4.1}
Let $K_\bullet$ be a chain complex of free $A$-modules, zero except in dimensions $0$ and $1$.
For every chain complex $L_\bullet$ of $A$-modules, we have an exact sequence
\[
  0\to\HH_0(K_\bullet\otimes\HH_p(L_\bullet))\to\HH_p(K_\bullet\otimes L_\bullet)\to\HH_1(K_\bullet\otimes\HH_{p-1}(L_\bullet))\to 0
\]
for every index $p$.
\end{lemma}

\oldpage[III]{84}
This is a particular case of an exact sequence of low-order terms of the K\"unneth spectral sequence (M, XVII, 5.2 (a) and G, I, 5.5.2); it can be proved directly as follows.
Consider $K_0$ and $K_1$ as chain complexes (zero in dimensions $\neq 0$ and $\neq 1$ respectively); we then have an exact sequence of complexes
\[
  0\to K_0\otimes L_\bullet\to K_\bullet\otimes L_\bullet\to K_1\otimes L_\bullet\to 0,
\]
to which we can apply the exact sequence in homology
\[
  \begin{aligned}
  \cdots&\to\HH_{p+1}(K_1\otimes L_\bullet)
  \xrightarrow{\partial}\HH_p(K_0\otimes L_\bullet)
  \to\HH_p(K_\bullet\otimes L_\bullet)\\
  &\to\HH_p(K_1\otimes L_\bullet)
  \xrightarrow{\partial}\HH_{p-1}(K_0\otimes L_\bullet)\to\cdots.
  \end{aligned}
\]
But it is evident that $\HH_p(K_0\otimes L_\bullet)=K_0\otimes\HH_p(L_\bullet)$ and $\HH_p(K_1\otimes L_\bullet)=K_1\otimes\HH_{p-1}(L_\bullet)$ for all $p$; in addition, we verify immediately that the operator $\partial:K_1\otimes\HH_p(L_\bullet)\to K_0\otimes\HH_p(L_\bullet)$ is none other than $d_1\otimes 1$; the lemma thus follows from the above exact sequence and the definition of $\HH_0(K_\bullet\otimes\HH_p(L_\bullet))$ and $\HH_1(K_\bullet\otimes\HH_{p-1}(L_\bullet))$.

The lemma having been established, the end of the proof of Proposition \sref{III.1.1.4} is immediate: the induction hypothesis and Lemma \sref{III.1.1.4.1} give $\HH_p(K_\bullet\otimes L_\bullet)=0$ for $p\geq 2$; in addition if we show that $\HH_1(K_\bullet,\HH_0(L_\bullet))=0$, then we also deduce from Lemma \sref{III.1.1.4.1} that $\HH_1(K_\bullet\otimes L_\bullet)=0$; but by definition, $\HH_1(K_\bullet,\HH_0(L_\bullet))$ is none other than the kernel of the scaling $z\mapsto f_r\cdot z$ on $M_{r-1}$, and as by hypothesis this kernel is zero, this finishes the proof.
\end{proof}

\begin{env}[1.1.5]
\label{III.1.1.5}
Let $\mathbf{g}=(g_i)_{1\leq i\leq r}$ be a second sequence of $r$ elements of $A$, and set $\mathbf{f}\mathbf{g}=(f_i g_i)_{1\leq i\leq r}$.
We can define a canonical homomorphism of complexes
\[
  \vphi_\mathbf{g}:K_\bullet(\mathbf{f}\mathbf{g})\to K_\bullet(\mathbf{f})
  \tag{1.1.5.1}
\]
as the canonical extension to the exterior algebra $\wedge(A^r)$ of the $A$-linear map $(x_1,\dots,x_r)\mapsto(g_1 x_1,\dots,g_r x_r)$ from $A^r$ to itself.
To see that we have a homomorphism of complexes, it suffices to note, in general, that if $u:E\to F$ is an $A$-linear map, and if $\mathbf{x}\in\dual{F}$ and $\mathbf{y}={}^t u(\mathbf{x})\in\dual{E}$, then we have the formula
\[
  (\wedge u)\circ i_\mathbf{y}=i_\mathbf{x}\circ(\wedge u);
  \tag{1.1.5.2}
\]
indeed, both sides are antiderivations of $\wedge F$, and it suffices to check that they coincide on $F$, which follows immediately from the definitions.

When we identify $K_\bullet(\mathbf{f})$ with the tensor product of the $K_\bullet(f_i)$ \sref{III.1.1.1}, $\vphi_\mathbf{g}$ is the tensor product of the $\vphi_{g_i}$, where $\vphi_{g_i}$ is the identity in degree $0$ and multiplication by $g_i$ in degree $1$.
\end{env}

\begin{env}[1.1.6]
\label{III.1.1.6}
In particular, for every pair of integers $m$ and $n$ such that $0\leq n\leq m$, we have homomorphisms of complexes
\[
  \vphi_{\mathbf{f}^{m-n}}:K_\bullet(\mathbf{f}^m)\to K_\bullet(\mathbf{f}^n)
  \tag{1.1.6.1}
\]
and as a result, homomorphisms
\[
  \vphi_{\mathbf{f}^{m-n}}:K^\bullet(\mathbf{f}^n,M)\to K^\bullet(\mathbf{f}^m,M),
  \tag{1.1.6.2}
\]
\[
  \vphi_{\mathbf{f}^{m-n}}:\HH^\bullet(\mathbf{f}^n,M)\to\HH^\bullet(\mathbf{f}^m,M).
  \tag{1.1.6.3}
\]

\oldpage[III]{85}
The latter homomorphisms evidently satisfy the transitivity condition $\vphi_{\mathbf{f}^{m-p}}=\vphi_{\mathbf{f}^{m-n}}\circ\vphi_{\mathbf{f}^{n-p}}$ for $p\leq n\leq m$; they therefore define two \emph{inductive systems} of $A$-modules; we set
\[
  C^\bullet((\mathbf{f}),M)=\varinjlim_n K^\bullet(\mathbf{f}^n,M),
  \tag{1.1.6.4}
\]
\[
  \HH^\bullet((\mathbf{f}),M)=\HH^\bullet(C^\bullet((\mathbf{f}),M))=\varinjlim_n\HH^\bullet(\mathbf{f}^n,M),
  \tag{1.1.6.5}
\]
the last equality following from the fact that passing to the inductive limit commutes with the functor $\HH^\bullet$ (G, I, 2.1).
We will later see \sref{III.1.4.3} that $\HH^\bullet((\mathbf{f}),M)$ depends in fact only on the \emph{ideal $(\mathbf{f})$} of $A$ (and even only on the $(\mathbf{f})$-pre-adic topology on $A$), which justifies the notations.

It is clear that $M\mapsto C^\bullet((\mathbf{f}),M)$ is an exact $A$-linear functor, and $M\mapsto\HH^\bullet((\mathbf{f}),M)$ is a cohomological functor.
\end{env}

\begin{env}[1.1.7]
\label{III.1.1.7}
Set $\mathbf{f}=(f_i)\in A^r$ and $\mathbf{g}=(g_i)\in A^r$; denote by $e_\mathbf{g}$ the left multiplication by the vector $\mathbf{g}\in A^r$ on the exterior algebra $\wedge(A^r)$; we know that we have the \emph{homotopy formula}
\[
  i_\mathbf{f}e_\mathbf{g}+e_\mathbf{g}i_\mathbf{f}=\langle\mathbf{g},\mathbf{f}\rangle 1
  \tag{1.1.7.1}
\]
in the $A$-module $\wedge(A^r)$ ($1$ denotes the identity automorphism of $\wedge(A^r)$); this relation also implies that \emph{in the complex $K_\bullet(\mathbf{f})$} we have
\[
  de_\mathbf{g}+e_\mathbf{g}d=\langle\mathbf{g},\mathbf{f}\rangle 1.
  \tag{1.1.7.2}
\]

If the ideal $(\mathbf{f})$ is equal to $A$, then there exists a $\mathbf{g}\in A^r$ such that $\langle\mathbf{g},\mathbf{f}\rangle=\sum_{i=1}^r g_i f_i=1$.
As a result (G, I, 2.4):
\end{env}

\begin{proposition}[1.1.8]
\label{III.1.1.8}
Suppose that the ideal $(\mathbf{f})$ generated by the $f_i$ is equal to $A$.
Then the complex $K_\bullet(\mathbf{f})$ is homotopically trivial, and so are the complexes $K_\bullet(\mathbf{f},M)$ and $K^\bullet(\mathbf{f},M)$ for every $A$-module $M$.
\end{proposition}

\begin{corollary}[1.1.9]
\label{III.1.1.9}
If $(\mathbf{f})=A$, then we have $\HH^\bullet(\mathbf{f},M)=0$ and $\HH^\bullet((\mathbf{f}),M)=0$ for every $A$-module $M$.
\end{corollary}

\begin{proof}
Indeed, we then have $(\mathbf{f}^n)=A$ for all $n$.
\end{proof}

\begin{remark}[1.1.10]
\label{III.1.1.10}
With the same notations as above, set $X=\Spec(A)$ and $Y$ the closed subprescheme of $X$ defined by the ideal $(\mathbf{f})$.
We will prove in \textsection9 that $\HH^\bullet((\mathbf{f}),M)$ is isomorphic to the cohomology $\HH_Y^\bullet(X,\widetilde{M})$ corresponding to the antifilter $\Phi$ of closed subsets of $Y$ (T, 3.2).
We will also show that Proposition \sref{III.1.2.3} applied to $X$ and to $\sh{F}=\widetilde{M}$ is a particular case of an exact sequence in cohomology
\[
  \cdots\to\HH_Y^p(X,\sh{F})\to\HH^p(X,\sh{F})\to\HH^p(X\setmin Y,\sh{F})\to\HH_Y^{p+1}(X,\sh{F})\to\cdots.
\]
\end{remark}

\subsection{\v Cech cohomology of an open cover}
\label{subsection:III.1.2}

\begin{notation}[1.2.1]
\label{III.1.2.1}
In this section, we denote:
\begin{enumerate}
  \item $X$ a prescheme;
  \item $\sh{F}$ a quasi-coherent $\sh{O}_X$-module;
\oldpage[III]{86}
  \item $A=\Gamma(X,\sh{O}_X)$, $M=\Gamma(X,\sh{F})$;
  \item $\mathbf{f}=(f_i)_{1\leq i\leq r}$ a finite system of elements of $A$;
  \item $U_i=X_{f_i}$, the open set \sref[0]{0.5.5.2} of the $x\in X$ such that $f_i(x)\neq 0$;
  \item $U=\bigcup_{i=1}^r U_i$;
  \item $\mathfrak{U}$ the cover $(U_i)_{1\leq i\leq r}$ of $U$.
\end{enumerate}
\end{notation}

\begin{env}[1.2.2]
\label{III.1.2.2}
Suppose that $X$ is either a prescheme whose underlying space is \emph{Noetherian} or a \emph{scheme} whose underlying space is \emph{quasi-compact}.
We then know \sref[I]{I.9.3.3} that we have $\Gamma(U_i,\sh{F})=M_{f_i}$.
We set
\[
  U_{i_0 i_1\cdots i_p}=\bigcap_{k=0}^p U_{i_k}=X_{f_{i_0}f_{i_1}\cdots f_{i_p}}
\]
\sref[0]{0.5.5.3}; so we also have
\[
  \Gamma(U_{i_0 i_1\cdots i_p},\sh{F})=M_{f_{i_0}f_{i_1}\cdots f_{i_p}}.
  \tag{1.2.2.1}
\]

We have \sref[0]{0.1.6.1} that $M_{f_{i_0}f_{i_1}\cdots f_{i_p}}$ identifies with the inductive limit $\varinjlim_n M_{i_0 i_1\cdots i_p}^{(n)}$, where the inductive system is formed by the $M_{i_0 i_1\cdots i_p}^{(n)}=M$, the homomorphisms $\vphi_{nm}:M_{i_0 i_1\cdots i_p}^{(m)}\to M_{i_0 i_1\cdots i_p}^{(n)}$ being multiplication by $(f_{i_0}f_{i_1}\cdots f_{i_p})^{n-m}$ for $m\leq n$.
We denote by $C_n^p(M)$ the set of \emph{alternating} maps from $[1,r]^{p+1}$ to $M$ (for all $n$); these $A$-modules also form an inductive system with respect to the $\vphi_{nm}$.
If $C^p(\mathfrak{U},\sh{F})$ is the group of \emph{alternating} \v Cech $p$-cochains relative to the cover $\mathfrak{U}$, with coefficients in $\sh{F}$ (G, II, 5.1), then it follows from the above that we can write
\[
  C^p(\mathfrak{U},\sh{F})=\varinjlim_n C_n^p(M).
  \tag{1.2.2.2}
\]

With the notations of \sref{III.1.1.2}, $C_n^p(M)$ identifies with $K^{p+1}(\mathbf{f}^n,M)$, and the map $\vphi_{nm}$ identifies with the map $\vphi_{\mathbf{f}^{n-m}}$ defined in \sref{III.1.1.6}.
We thus have, for every $p\geq 0$, a canonical functorial isomorphism
\[
  C^p(\mathfrak{U},\sh{F})\isoto C^{p+1}((\mathbf{f}),M).
  \tag{1.2.2.3}\label{III.1.2.2.3}
\]

In addition, the formula (\hyperref[III.1.1.2.3]{1.1.2.3}) and the definition of the cohomology of a cover (G, II, 5.1) show that the isomorphisms (\hyperref[III.1.2.2.3]{1.2.2.3}) are compatible with the coboundary maps.
\end{env}

\begin{proposition}[1.2.3]
\label{III.1.2.3}
If $X$ is a prescheme whose underlying space is Noetherian or a scheme whose underlying space is quasi-compact, then there exists a canonical functorial isomorphism in $\sh{F}$
\[
  \HH^p(\mathfrak{U},\sh{F})\isoto\HH^{p+1}((\mathbf{f}),M)\text{ for }p\geq 1.
  \tag{1.2.3.1}\label{III.1.2.3.1}
\]

In addition, we have a functorial exact sequence in $\sh{F}$
\[
  0\to\HH^0((\mathbf{f}),M)\to M\to\HH^0(\mathfrak{U},\sh{F})\to\HH^1((\mathbf{f}),M)\to 0.
  \tag{1.2.3.2}\label{III.1.2.3.2}
\]
\end{proposition}

\oldpage[III]{87}
\begin{proof}
The isomorphisms (\hyperref[III.1.2.3.1]{1.2.3.1}) are immediate consequences of what we saw in \sref{III.1.2.2}.
On the other hand, we have $C^0(\mathfrak{U},\sh{F})=C^1((\mathbf{f}),M)$; as a result, $\HH^0(\mathfrak{U},\sh{F})$ identifies with the subgroup of $1$-cocycles of $C^1((\mathbf{f}),M)$; as $M=C^0((\mathbf{f}),M)$, the exact sequence (\hyperref[III.1.2.3.2]{1.2.3.2}) is none other than the one given by the definition of the cohomology groups $\HH^0((\mathbf{f}),M)$ and $\HH^1((\mathbf{f}),M)$.
\end{proof}

\begin{corollary}[1.2.4]
\label{III.1.2.4}
Suppose that the $X_{f_i}$ are quasi-compact and that there exists $g_i\in\Gamma(U,\sh{O}_X)$ such that $\sum_i g_i(f_i|U)=1|U$.
Then for every quasi-coherent $(\sh{O}_X|U)$-module $\sh{G}$, we have $\HH^p(\mathfrak{U},\sh{G})=0$ for $p>0$; if in addition $U=X$, then the canonical homomorphism (\hyperref[III.1.2.3.2]{1.2.3.2}) $\Gamma(X,\sh{G})\to\HH^0(\mathfrak{U},\sh{G})$ is bijective.
\end{corollary}

\begin{proof}
As by hypothesis the $U_i=X_{f_i}$ are quasi-compact, so is $U$, and we can reduce to the case where $U=X$; setting $N=\Gamma(X,\sh{G})$, the hypothesis then implies that $\HH^p((\mathbf{f}),N)=0$ for all $p\geq 0$ \sref{III.1.1.9}.
The corollary then follows immediately from (\hyperref[III.1.2.3.1]{1.2.3.1}) and (\hyperref[III.1.2.3.2]{1.2.3.2}).
\end{proof}

We note that since $\HH^0(\mathfrak{U},\sh{F})=\HH^0(U,\sh{F})$ (G, II, 5.2.2), we have again proved \sref[I]{I.1.3.7} as a special case.

\begin{remark}[1.2.5]
\label{III.1.2.5}
Suppose that $X$ is an \emph{affine scheme}; then the $U_i=X_{f_i}=D(f_i)$ are affine open sets, as well as the $U_{i_0 i_1\cdots i_p}$ (but $U$ is not necessarily affine).
In this case, the functors $\Gamma(X,\sh{F})$ and $\Gamma(U_{i_0 i_1\cdots i_p},\sh{F})$ are exact in $\sh{F}$ \sref[I]{I.1.3.11}.
If we have an exact sequence $0\to\sh{F}'\to\sh{F}\to\sh{F}''\to 0$ of quasi-coherent $\sh{O}_X$-modules, then the sequence of complexes
\[
  0\to C^\bullet(\mathfrak{U},\sh{F}')\to C^\bullet(\mathfrak{U},\sh{F})\to C^\bullet(\mathfrak{U},\sh{F}'')\to 0
\]
is exact, and thus gives an exact sequence in cohomology
\[
  \cdots\to\HH^p(\mathfrak{U},\sh{F}')\to\HH^p(\mathfrak{U},\sh{F})\to\HH^p(\mathfrak{U},\sh{F}'')\xrightarrow{\partial}\HH^{p+1}(\mathfrak{U},\sh{F}')\to\cdots.
\]
On the other hand, if we set $M'=\Gamma(X,\sh{F}')$ and $M''=\Gamma(X,\sh{F}'')$, then the sequence $0\to M'\to M\to M''\to 0$ is exact; as $C^\bullet((\mathbf{f}),M)$ is an exact functor in $M$, we also have the exact sequence in cohomology
\[
  \cdots\to\HH^p((\mathbf{f}),M')\to\HH^p((\mathbf{f}),M)\to\HH^p((\mathbf{f}),M'')\xrightarrow{\partial}\HH^{p+1}((\mathbf{f}),M')\to\cdots.
\]

This being so, as the diagram
\[
  \xymatrix{
    0\ar[r] &
    C^\bullet(\mathfrak{U},\sh{F}')\ar[r]\ar[d] &
    C^\bullet(\mathfrak{U},\sh{F})\ar[r]\ar[d] &
    C^\bullet(\mathfrak{U},\sh{F}'')\ar[r]\ar[d] &
    0\\
    0\ar[r] &
    C^\bullet((\mathbf{f}),M')\ar[r] &
    C^\bullet((\mathbf{f}),M)\ar[r] &
    C^\bullet((\mathbf{f}),M'')\ar[r] &
    0
  }
\]
\oldpage[III]{88}
is commutative, we conclude that the diagrams
\[
  \xymatrix{
    \HH^p(\mathfrak{U},\sh{F}'')\ar[r]^\partial\ar[d] &
    \HH^{p+1}(\mathfrak{U},\sh{F}')\ar[d]\\
    \HH^{p+1}((\mathbf{f}),M'')\ar[r]^\partial &
    \HH^{p+2}((\mathbf{f}),M')
  }
  \tag{1.2.5.1}
\]
are commutative for all $p$~(G, I, 2.1.1).
\end{remark}

\subsection{Cohomology of an affine scheme}
\label{subsection:III.1.3}

\begin{theorem}[1.3.1]
\label{III.1.3.1}
Let $X$ be an affine scheme.
For every quasi-coherent $\sh{O}_X$-module $\sh{F}$, we have $\HH^p(X,\sh{F})=0$ for all $p>0$.
\end{theorem}

\begin{proof}
Let $\mathfrak{U}$ be a finite cover of $X$ by the affine open sets $X_{f_i}=D(f_i)$ ($1\leq i\leq r$); we then know that the ideal of $A=\Gamma(X,\sh{O}_X)$ generated by the $f_i$ is equal to $A$.
We thus conclude from Corollary~\sref{III.1.2.4} that we have $\HH^p(\mathfrak{U},\sh{F})=0$ for $p>0$.
As there are finite covers of $X$ by affine open sets which are arbitrarily fine \sref[I]{I.1.1.10}, the definition of \v Cech cohomology (G, II, 5.8) shows that we also have $\CHH^p(X,\sh{F})=0$ for $p>0$.
But this also applies to every prescheme $X_f$ for $f\in A$ \sref[I]{I.1.3.6}, hence $\CHH^p(X_f,\sh{F})=0$ for $p>0$.
As we have $X_f\cap X_g=X_{fg}$, we deduce that we also have $\HH^p(X,\sh{F})=0$ for all $p>0$, by virtue of (G, II, 5.9.2).
\end{proof}

\begin{corollary}[1.3.2]
\label{III.1.3.2}
Let $Y$ be a prescheme, $f:X\to Y$ an affine morphism \sref[II]{II.1.6.1}.
For every quasi-coherent $\sh{O}_X$-module $\sh{F}$, we have $\RR^q f_*(\sh{F})=0$ for $q>0$.
\end{corollary}

\begin{proof}
By definition $\RR^q f_*(\sh{F})$ is the $\sh{O}_Y$-module associated to the presheaf $U\mapsto\HH^q(f^{-1}(U),\sh{F})$, where $U$ varies over the open subsets of $Y$.
But the affine open sets form a basis for $Y$, and for such an open set $U$, $f^{-1}(U)$ is affine \sref[II]{II.1.3.2}, hence $\HH^q(f^{-1}(U),\sh{F})=0$ by Theorem~\sref{III.1.3.1}, which proves the corollary.
\end{proof}

\begin{corollary}[1.3.3]
\label{III.1.3.3}
Let $Y$ be a prescheme, $f:X\to Y$ an affine morphism.
For every quasi-coherent $\sh{O}_X$-module $\sh{F}$, the canonical homomorphism $\HH^p(Y,f_*(\sh{F}))\to\HH^p(X,\sh{F})$ (\textbf{0},~12.1.3.1) is bijective for all $p$.
\end{corollary}

\begin{proof}
It suffices (by \sref[0]{0.12.1.7}) to show that the edge homomorphisms $''E_2^{p0}=\HH^p(Y,f_*(\sh{F}))\to\HH^p(X,\sh{F})$ of the second spectral sequence of the composite functor $\Gamma f_*$ are bijective.
But the $E_2$-term of this spectral sequence is given by $''E_2^{pq}=\HH^p(Y,\RR^q f_*(\sh{F}))$ (G, II, 4.17.1), so it follows from Corollary~\sref{III.1.3.2} that $''E_2^{pq}=0$ for $q>0$, and the spectral sequence degenerates; hence our assertion \sref[0]{0.11.1.6}.
\end{proof}

\begin{corollary}[1.3.4]
\label{III.1.3.4}
Let $f:X\to Y$ be an affine morphism, $g:Y\to Z$ a morphism.
\oldpage[III]{89}
For every quasi-coherent $\sh{O}_X$-module $\sh{F}$, the canonical homomorphism $\RR^p g_*(f_*(\sh{F}))\to\RR^p(g\circ f)_*(\sh{F})$ (\textbf{0},~12.2.5.1) is bijective for all $p$.
\end{corollary}

\begin{proof}
It suffices to note that, according to Corollary~\sref{III.1.3.3}, for every affine open subset $W$ of $Z$, the canonical homomorphism $\HH^p(g^{-1}(W),f_*(\sh{F}))\to\HH^p(f^{-1}(g^{-1}(W)),\sh{F})$ is bijective; this proves that the homomorphism of presheaves defining the canonical homomorphism $\RR^p g_*(f_*(\sh{F}))\to\RR^p(g\circ f)_*(\sh{F})$ is bijective \sref[0]{0.12.2.5}.
\end{proof}

\subsection{Application to the cohomology of arbitrary preschemes}
\label{subsection:III.1.4}

\begin{proposition}[1.4.1]
\label{III.1.4.1}
Let $X$ be a scheme, $\mathfrak{U}=(U_\alpha)$ be a cover of $X$ by affine open sets.
For every quasi-coherent $\sh{O}_X$-module $\sh{F}$, the cohomology modules $\HH^\bullet(X,\sh{F})$ and $\HH^\bullet(\mathfrak{U},\sh{F})$ (over $\Gamma(X,\sh{O}_X)$) are canonically isomorphic.
\end{proposition}

\begin{proof}
  As $X$ is a scheme, every finite intersection $V$ of open sets in the cover $\mathfrak{U}$ is affine \sref[I]{I.5.5.6}, so $\HH^q(V,\sh{F})=0$ for $q\geq 1$ by Theorem~\sref{III.1.3.1}.
The proposition then follows from a theorem of Leray (G, II, 5.4.1).
\end{proof}

\begin{remark}[1.4.2]
\label{III.1.4.2}
We note that the result of Proposition~\sref{III.1.4.1} is still true when the finite intersections of the sets $U_\alpha$ are affine, even when we do not necessarily assume that $X$ is a scheme.
\end{remark}

\begin{corollary}[1.4.3]
\label{III.1.4.3}
Let $X$ be a scheme with quasi-compact underlying space, $A=\Gamma(X,\sh{O}_X)$, and $\mathbf{f}=(f_i)_{1\leq i\leq r}$ a finite sequence of elements of $A$ such that the $X_{f_i}$ (notation of \sref{III.1.2.1}) are affine.
Then (with the notations of \sref{III.1.2.1}), for every quasi-coherent $\sh{O}_X$-module $\sh{F}$, we have a canonical isomorphism which is functorial in $\sh{F}$
\[
  \HH^q(U,\sh{F})\isoto\HH^{q+1}((\mathbf{f}),M)\text{ for }q\geq 1,
  \tag{1.4.3.1}\label{III.1.4.3.1}
\]
and an exact sequence which is functorial in $\sh{F}$
\[
  0\to\HH^0((\mathbf{f}),M)\to M\to\HH^0(U,\sh{F})\to\HH^1((\mathbf{f}),M)\to 0.
  \tag{1.4.3.2}
\]
\end{corollary}

\begin{proof}
This follows immediately from Propositions~\sref{III.1.4.1} and \sref{III.1.2.3}.
\end{proof}

\begin{env}[1.4.4]
\label{III.1.4.4}
If $X$ is an \emph{affine scheme}, then it follows from Remark~\sref{III.1.2.5} and Proposition~\sref{III.1.4.1} that for all $q\geq 0$, the diagrams
\[
  \xymatrix{
    \HH^q(U,\sh{F}'')\ar[r]^\partial\ar[d] &
    \HH^{q+1}(U,\sh{F}')\ar[d]\\
    \HH^{q+1}((\mathbf{f}),M'')\ar[r]^\partial &
    \HH^{q+2}((\mathbf{f}),M')
  }
  \tag{1.4.4.1}
\]
corresponding to an exact sequence $0\to\sh{F}'\to\sh{F}\to\sh{F}''\to 0$ of quasi-coherent $\sh{O}_X$-modules (with the notations of Remark~\sref{III.1.2.5}) are commutative.
\end{env}

\begin{proposition}[1.4.5]
\label{III.1.4.5}
\oldpage[III]{90}
Let $X$ be a quasi-compact scheme, $\sh{L}$ an invertible $\sh{O}_X$-module, and consider the graded ring $A_\bullet=\Gamma_\bullet(\sh{L})$ \sref[0]{0.5.4.6}; then $\HH^\bullet(\sh{F},\sh{L})=\bigoplus_{n\in\bb{Z}}\HH^\bullet(X,\sh{F}\otimes\sh{L}^{\otimes n})$ is a graded $A_\bullet$-module, and for all $f\in A_n$, we have a canonical isomorphism
\[
  \HH^\bullet(X_f,\sh{F})\isoto(\HH^\bullet(\sh{F},\sh{L}))_{(f)}
  \tag{1.4.5.1}\label{III.1.4.5.1}
\]
of $(A_\bullet)_{(f)}$-modules.
\end{proposition}

\begin{proof}
As $X$ is a quasi-compact scheme, we can calculate the cohomology of all the $\sh{O}_X$-modules $\sh{F}\otimes\sh{L}^{\otimes n}$ using the same finite cover $\mathfrak{U}=(U_i)$ consisting of the affine open sets such that the restriction $\sh{L}|U_i$ is isomorphic to $\sh{O}_X|U_i$ for each $i$ \sref{III.1.4.1}.
It is then immediate that the $U_i\cap X_f$ are affine open sets \sref[I]{I.1.3.6}, and we can thus calculate the cohomology $\HH^\bullet(X_f,\sh{F}\otimes\sh{L}^{\otimes n})$ using the cover $\mathfrak{U}|X_f=(U_i\cap X_f)$ \sref{III.1.4.1}.
It is immediate that for all $f\in A_n$, multiplication by $f$ defines a homomorphism $C^\bullet(\mathfrak{U},\sh{F}\otimes\sh{L}^m)\to C^\bullet(\mathfrak{U},\sh{F}\otimes\sh{L}^{\otimes(m+n)})$, hence a homomorphism $\HH^\bullet(\mathfrak{U},\sh{F}\otimes\sh{L}^{\otimes m})\to\HH^\bullet(\mathfrak{U},\sh{F}\otimes\sh{L}^{\otimes(m+n)})$, which establishes the first assertion.
On the other hand, for a given $f\in A_n$, it follows from \sref[I]{I.9.3.2} that we have an isomorphism of complexes of $(A_\bullet)_{(f)}$-modules
\[
  C^\bullet(\mathfrak{U}|X_f,\sh{F})\isoto\bigg(\!\!C^\bullet\bigg(\mathfrak{U},\bigoplus_{n\in\bb{Z}}\sh{F}\otimes\sh{L}^{\otimes n}\bigg)\!\!\bigg)_{(f)},
\]
taking into account \sref[I]{I.1.3.9}[ii].
Passing to the cohomology of these complexes, we induce the isomorphism (\hyperref[III.1.4.5.1]{1.4.5.1}), recalling that the functor $M\mapsto M_{(f)}$ is exact on the category of graded $A_\bullet$-modules.
\end{proof}

\begin{corollary}[1.4.6]
\label{III.1.4.6}
Suppose that the hypotheses of Proposition~\sref{III.1.4.5} are satisfied, and in addition suppose that $\sh{L}=\sh{O}_X$.
If we set $A=\Gamma(X,\sh{O}_X)$, then for all $f\in A$, we have a canonical isomorphism $\HH^\bullet(X_f,\sh{F})\isoto(\HH^\bullet(X,\sh{F}))_f$ of $A_f$-modules.
\end{corollary}

\begin{corollary}[1.4.7]
\label{III.1.4.7}
Let $X$ be a quasi-compact scheme, $f$ an element of $\Gamma(X,\sh{O}_X)$.
\begin{enumerate}
  \item[{\rm(i)}] Suppose that the open set $X_f$ is affine.
    Then for every quasi-coherent $\sh{O}_X$-module $\sh{F}$, every $i>0$, and every $\xi\in\HH^i(X,\sh{F})$, there exists an integer $n>0$ such that $f^n\xi=0$.
  \item[{\rm(ii)}] Conversely, suppose that $X_f$ is quasi-compact and that for every quasi-coherent sheaf of ideals $\sh{J}$ of $\sh{O}_X$ and every $\zeta\in\HH^1(X,\sh{J})$, there exists an $n>0$ such that $f^n\zeta=0$.
    Then $X_f$ is affine.
\end{enumerate}
\end{corollary}

\begin{proof}
\medskip\noindent
\begin{enumerate}
  \item[(i)] If $X_f$ is affine, then we have $\HH^i(X_f,\sh{F})=0$ for all $i>0$ \sref{III.1.3.1}, so the assertion follows directly from Corollary~\sref{III.1.4.6}.
  \item[(ii)] By virtue of Serre's criterion \sref[II]{II.5.2.1}, it suffices to prove that for every quasi-coherent sheaf of ideals $\sh{K}$ of $\sh{O}_X|X_f$, we have $\HH^1(X_f,\sh{K})=0$.
    As $X_f$ is a quasi-compact open set in a quasi-compact scheme $X$, there exists a quasi-coherent sheaf of ideals $\sh{J}$ of $\sh{O}_X$ such that $\sh{K}=\sh{J}|X_f$ \sref[I]{I.9.4.2}.
    According to Corollary~\sref{III.1.4.6}, we have $\HH^1(X_f,\sh{K})=(\HH^1(X,\sh{J}))_f$, and the hypothesis implies that the right hand side is zero, hence the assertion.
\end{enumerate}
\end{proof}

\begin{remark}[1.4.8]
\label{III.1.4.8}
We note that Corollary~\sref{III.1.4.7}[i] gives a simpler proof of the relation (\textbf{II},~4.5.13.2).
\end{remark}

\begin{lemma}[1.4.9]
\label{III.1.4.9}
\oldpage[III]{91}
Let $X$ be a quasi-compact scheme, $\mathfrak{U}=(U_i)_{1\leq i\leq n}$ a finite cover of $X$ by affine open sets, and ${\sh{F}}$ a quasi-coherent $\sh{O}_X$-module.
The complex of sheaves $\sh{C}^\bullet(\mathfrak{U},\sh{F})$ defined by the cover $\mathfrak{U}$ (G, II, 5.2) is then a quasi-coherent $\sh{O}_X$-module.
\end{lemma}

\begin{proof}
It follows from the definitions (G, II, 5.2) that $\sh{C}^p(\mathfrak{U},\sh{F})$ is the direct sum of the direct image sheaves of the $\sh{F}|U_{i_0\cdots i_p}$ under the canonical injection $U_{i_0\cdots i_p}\to X$.
The hypothesis that $X$ is a scheme implies that these injections are affine morphisms \sref[I]{I.5.5.6}, hence the $\sh{C}^p(\mathfrak{U},\sh{F})$ are quasi-coherent \sref[II]{II.1.2.6}.
\end{proof}

\begin{proposition}[1.4.10]
\label{III.1.4.10}
Let $u:X\to Y$ be a separated and quasi-compact morphism.
For every quasi-coherent $\sh{O}_X$-module $\sh{F}$, the $\RR^q u_*(\sh{F})$ are quasi-coherent $\sh{O}_Y$-modules.
\end{proposition}

\begin{proof}
The question is local on $Y$, so we can suppose that $Y$ is affine.
Then $X$ is a finite union of affine open sets $U_i$ ($1\leq i\leq n$); let $\mathfrak{U}$ be the cover $(U_i)$.
In addition, as $Y$ is a scheme, it follows from \sref[I]{I.5.5.10} that for every affine open $V\subset Y$, the canonical injection $u^{-1}(V)\to X$ is an affine morphism; we conclude (Proposition~\sref{III.1.4.1} and (G, II, 5.2)) that we have a canonical isomorphism
\[
  \HH^\bullet(u^{-1}(V),\sh{F})\isoto\HH^\bullet(\Gamma(V,\sh{K}^\bullet)),
  \tag{1.4.10.1}\label{III.1.4.10.1}
\]
where we set $\sh{K}^\bullet=u_*(\sh{C}^\bullet(\mathfrak{U},\sh{F}))$.
  According to Lemma~\sref{III.1.4.9} and \sref[I]{I.9.2.2}, $\sh{K}^\bullet$ is a quasi-coherent $\sh{O}_Y$-module; moreover, it constitutes a \emph{complex of sheaves} since so is $\sh{C}^\bullet(\mathfrak{U},\sh{F})$.
It then follows from the definition of the cohomology $\sh{H}^\bullet(\sh{K}^\bullet)$ (G, II, 4.1) that the latter consists of quasi-coherent $\sh{O}_Y$-modules \sref[I]{I.4.1.1}.
As (for $V$ affine in $Y$) the functor $\Gamma(V,\sh{G})$ is exact in $\sh{G}$ on the category of quasi-coherent $\sh{O}_Y$-modules, we have (G, II, 4.1)
\[
  \HH^\bullet(\Gamma(V,\sh{K}^\bullet))=\Gamma(V,\sh{H}^\bullet(\sh{K}^\bullet)).
  \tag{1.4.10.2}\label{III.1.4.10.2}
\]

Finally, we note that it follows from the definition of the canonical homomorphism
\[
  \HH^\bullet(\mathfrak{U},\sh{F})\to\HH^\bullet(X,\sh{F}),
\]
given in (G, II, 5.2), that if $V'\subset V$ is a second affine open subset of $Y$, then the
diagram
\[
  \xymatrix{
    \HH^\bullet(u^{-1}(V),\sh{F})\ar[r]^\sim\ar[d] &
    \HH^\bullet(\Gamma(V,\sh{K}^\bullet))\ar[d]\\
    \HH^\bullet(u^{-1}(V'),\sh{F})\ar[r]^\sim &
    \HH^\bullet(\Gamma(V',\sh{K}^\bullet))
  }
\]
is commutative.
We thus conclude from the above that the isomorphisms (\hyperref[III.1.4.10.1]{1.4.10.1}) define an isomorphism of $\sh{O}_Y$-modules
\[
  \RR^\bullet u_*(\sh{F})\isoto\sh{H}^\bullet(\sh{K}^\bullet),
  \tag{1.4.10.3}\label{III.1.4.10.3}
\]
\oldpage[III]{92}
and as a result, $\RR^\bullet u_*(\sh{F})$ is quasi-coherent.
\end{proof}

In addition, it follows from (\hyperref[III.1.4.10.3]{1.4.10.3}), (\hyperref[III.1.4.10.2]{1.4.10.2}), and (\hyperref[III.1.4.10.1]{1.4.10.1}) that:
\begin{corollary}[1.4.11]
\label{III.1.4.11}
Under the hypotheses of Proposition~\sref{III.1.4.10}, for every affine open set $V$ of $Y$, the canonical homomorphism
\[
  \HH^q(u^{-1}(V),\sh{F})\to\Gamma(V,\RR^q u_*(\sh{F}))
  \tag{1.4.11.1}\label{III.1.4.11.1}
\]
is an isomorphism for all $q\geq 0$.
\end{corollary}

\begin{corollary}[1.4.12]
\label{III.1.4.12}
Suppose that the hypotheses of Proposition~\sref{III.1.4.10} are satisfied, and in addition suppose that $Y$ is quasi-compact.
Then there exists an integer $r>0$ such that for every quasi-coherent $\sh{O}_X$-module $\sh{F}$ and every integer $q>r$, we have $\RR^q u_*(\sh{F})=0$.
If $Y$ is affine, then we can take for $r$ an integer such that there exists a cover of $X$ consisting of $r$ affine open sets.
\end{corollary}

\begin{proof}
As we can cover $Y$ by a finite number of affine open sets, we can reduce to proving the second assertion, by virtue of Corollary~\sref{III.1.4.11}.
If $\mathfrak{U}$ is a cover of $X$ by $r$ affine open sets, then we have $\HH^q(\mathfrak{U},\sh{F})=0$ for $q>r$, since the cochains of $C^q(\mathfrak{U},\sh{F})$ are alternating; the assertion thus follows from Proposition~\sref{III.1.4.1}.
\end{proof}

\begin{corollary}[1.4.13]
\label{III.1.4.13}
Suppose that the hypotheses of Proposition~\sref{III.1.4.10} are satisfied, and in addition suppose that $Y=\Spec(A)$ is affine.
Then for every quasi-coherent $\sh{O}_X$-module $\sh{F}$ and every $f\in A$, we have
\[
  \Gamma(Y_f,\RR^q u_*(\sh{F}))=(\Gamma(Y,\RR^q u_*(\sh{F})))_f
\]
up to canonical isomorphism.
\end{corollary}

\begin{proof}
This follows from the fact that $\RR^q u_*(\sh{F})$ is a quasi-coherent $\sh{O}_Y$-module \sref[I]{I.1.3.7}.
\end{proof}

\begin{proposition}[1.4.14]
\label{III.1.4.14}
Let $f:X\to Y$ be a separated and quasi-compact morphism, $g:Y\to Z$ an affine morphism.
For every quasi-coherent $\sh{O}_X$-module $\sh{F}$, the canonical homomorphism $\RR^p(g\circ f)_*(\sh{F})\to g_*(\RR^p f_*(\sh{F}))$ (\textbf{0},~12.2.5.2) is bijective for all $p$.
\end{proposition}

\begin{proof}
For every affine open subset $W$ of $Z$, $g^{-1}(W)$ is an affine open subset of $Y$.
The homomorphism of presheaves defining the canonical homomorphism
\[
  \RR^p(g\circ f)_*(\sh{F})\to g_*(\RR^p f_*(\sh{F}))
\]
\sref[0]{0.12.2.5} is thus bijective by Corollary~\sref{III.1.4.11}.
\end{proof}

\begin{proposition}[1.4.15]
\label{III.1.4.15}
Let $u:X\to Y$ be a separated morphism quasi-compact, $v:Y'\to Y$ a flat morphism of preschemes \sref[0]{0.6.7.1}; let $u'=u_{(Y')}$, such that we have the commutative diagram
\[
  \xymatrix{
    X\ar[d]_u &
    X'=X_{(Y')}\ar[l]_{v'}\ar[d]^{u'}\\
    Y &
    Y'\ar[l]_{v}.
  }
  \tag{1.4.15.1}
  \label{III.1.4.15.1}
\]
Then for every quasi-coherent $\sh{O}_X$-module $\sh{F}$, $\RR^q u_*'(\sh{F}')$ is canonically isomorphic to $\RR^q u_*(\sh{F})\otimes_{\sh{O}_Y}\sh{O}_{Y'}=v^*(\RR^q u_*(\sh{F}))$ for all $q\geq 0$, where $\sh{F}'={v'}^*(\sh{F})=\sh{F}\otimes_{\sh{O}_Y}\sh{O}_{Y'}$.
\end{proposition}

\oldpage[III]{93}
\begin{proof}
The canonical homomorphism $\rho:\sh{F}\to v_*'({v'}^*(\sh{F}))$ (\textbf{0},~4.4.3.2) defines by functoriality a homomorphism
\[
  \RR^q u_*(\sh{F})\to\RR^q u_*(v_*'(\sh{F}')).
  \tag{1.4.15.2}\label{III.1.4.15.2}
\]

On the other hand, we have, by setting $w=u\circ v'=v\circ u'$, the canonical homomorphisms (\textbf{0}~,12.2.5.1 and 12.2.5.2)
\[
  \RR^q u_*(v_*'(\sh{F}'))\to\RR^q w_*(\sh{F}')\to v_*(\RR^q u_*'(\sh{F}')).
  \tag{1.4.15.3}\label{III.1.4.15.3}
\]

Composing (\hyperref[III.1.4.15.3]{1.4.15.3}) and (\hyperref[III.1.4.15.2]{1.4.15.2}), we have a homomorphism
\[
  \psi:\RR^q u_*(\sh{F})\to v_*(\RR^q u_*'(\sh{F}')),
\]
and finally we obtain a canonical homomorphism (whose definition does not make \emph{any assumptions} on $v$)
\[
  \psi^\sharp:v^*(\RR^q u_*(\sh{F}))\to\RR^q u_*'(\sh{F}'),
  \tag{1.4.15.4}
\]
and it is necessary to prove that it is an isomorphism when $v$ is \emph{flat}.
It is clear that the question is local on $Y$ and $Y'$, and we can therefore suppose that $Y=\Spec(A)$ and $Y'=\Spec(B)$; we will also use the following lemma:
\begin{lemma}[1.4.15.5]
\label{III.1.4.15.5}
Let $\vphi:A\to B$ be a ring homomorphism, $Y=\Spec(A)$, $X=\Spec(B)$, $f:X\to Y$ the morphism corresponding to $\vphi$, and $M$ a $B$-module.
For the $\sh{O}_X$-module $\widetilde{M}$ to be $f$-flat \sref[0]{0.6.7.1}, it is necessary and sufficient for $M$ to be a flat $A$-module.
In particular, for the morphism $f$ to be flat, it is necessary and sufficient for $B$ to be a flat $A$-module.
\end{lemma}

This follows from the definition \sref[0]{0.6.7.1} and from \sref[0]{0.6.3.3}, taking into account \sref[I]{I.1.3.4}.

This being so, it follows from (\hyperref[III.1.4.11.1]{1.4.11.1}) and the definitions of the homomorphisms (\hyperref[III.1.4.15.3]{1.4.15.3}) (cf.~\sref[0]{0.12.2.5}) that $\psi$ then corresponds to the composite morphism
\[
  \HH^q(X,\sh{F})\xrightarrow{\rho_q}\HH^q(X,v_*'({v'}^*(\sh{F})))\xrightarrow{\theta_q}\HH^q(X',{v'}^*(v_*'({v'}^*(\sh{F}))))\xrightarrow{\sigma_q}\HH^q(X',{v'}^*(\sh{F})),
\]
where $\rho_q$ and $\sigma_q$ are the homomorphisms in cohomology corresponding to the canonical morphisms $\rho$ and $\sigma:{v'}^*(v_*'(\sh{G}'))\to\sh{G}'$, and $\theta_q$ is the $\vphi$-morphism (\textbf{0},~12.1.3.1) relative to the $\sh{O}_X$-module $v_*'({v'}^*(\sh{F}))$.
But by the functoriality of $\theta_q$, we have the commutative diagram
\[
  \xymatrix{
    \HH^q(X,\sh{F})\ar[rr]^{\rho_q}\ar[dd]_{\theta_q} & &
    \HH^q(X,v_*'({v'}^*(\sh{F})))\ar[dd]^{\theta_q}\\\\
    \HH^q(X',{v'}^*(\sh{F}))\ar[rr]^{{v'}^*(\rho_q)} & &
    \HH^q(X',{v'}^*(v_*'({v'}^*(\sh{F})))),
  }
\]
\oldpage[III]{94}
and as by definition \sref[0]{0.4.4.3} ${v'}^*(\rho)$ is the inverse of $\sigma$, we see that the composite morphism considered above is finally none other than $\theta_q$; as a result, $\psi^\sharp$ is the associated $B$-homomorphism $\HH^q(X,\sh{F})\otimes_A B\to\HH^q(X',\sh{F}')$.
As $u$ is quasi-compact, $X$ is a finite union of affine open sets $U_i$ ($1\leq i\leq r$); let $\mathfrak{U}$ be the cover $(U_i)$.
As $v$ is an affine morphism, so is $v'$ \sref[II]{II.1.6.2}[iii], and as a result the $U_i'={v'}^{-1}(U_i)$ form an affine open cover $\mathfrak{U}'$ of $X'$.
We then know \sref[0]{0.12.1.4.2} that the diagram
\[
  \xymatrix{
    \HH^q(\mathfrak{U},\sh{F})\ar[r]^{\theta_q}\ar[d] &
    \HH^q(\mathfrak{U}',\sh{F}')\ar[d]\\
    \HH^q(X,\sh{F})\ar[r]^{\theta_q} &
    \HH^q(X',\sh{F}')
  }
\]
is commutative, and the vertical arrows are isomorphisms since $X$ and $X'$ are schemes \sref{III.1.4.1}.
As a result, it suffices to prove that the canonical $\vphi$-morphism $\theta_q:\HH^q(\mathfrak{U},\sh{F})\to\HH^q(\mathfrak{U}',\sh{F}')$ is such that the associated $B$-homomorphism
\[
  \HH^q(\mathfrak{U},\sh{F})\otimes_A B\to\HH^q(\mathfrak{U}',\sh{F}')
\]
is an isomorphism.
For every sequence $\mathbf{s}=(i_k)_{0\leq k\leq p}$ of $p+1$ indices of $[1,r]$, set $U_\mathbf{s}=\bigcap_{k=0}^p U_{i_k}$, $U_\mathbf{s}'=\bigcap_{k=0}^p U_{i_k}'={v'}^{-1}(U_\mathbf{s})$, $M_\mathbf{s}=\Gamma(U_\mathbf{s},\sh{F})$, and $M_\mathbf{s}'=\Gamma(U_\mathbf{s}',\sh{F}')$.
The canonical map $M_\mathbf{s}\otimes_A B\to M_\mathbf{s}'$ is an isomorphism \sref[I]{I.1.6.5}, hence the canonical map $C^p(\mathfrak{U},\sh{F})\otimes_A B\to C^p(\mathfrak{U}',\sh{F}')$ is an isomorphism, by which $d\otimes 1$ identifies with the coboundary map $C^p(\mathfrak{U}',\sh{F}')\to C^{p+1}(\mathfrak{U}',\sh{F}')$.
As $B$ is a \emph{flat} $A$-module, it follows from the definition of the cohomology modules that the canonical map $\HH^q(\mathfrak{U},\sh{F})\otimes_A B\to\HH^q(\mathfrak{U}',\sh{F}')$ is an isomorphism \sref[0]{0.6.1.1}.
This result will later be generalized in \hyperref[section:III.6]{\textsection6}.
\end{proof}

\begin{corollary}[1.4.16]
\label{III.1.4.16}
Let $A$ be a ring, $X$ an $A$-scheme of finite type, and $B$ an $A$-algebra which is faithfully flat over $A$.
For $X$ to be affine, it is necessary and sufficient for $X\otimes_A B$ to be.
\end{corollary}

\begin{proof}
The condition is evidently necessary \sref[I]{I.3.2.2}; we show that it is sufficient.
As $X$ is separated over $A$ and the morphism $\Spec(B)\to\Spec(A)$ is flat, it follows from Proposition~\sref{III.1.4.15} that we have
\[
\label{III.1.4.16.1}
  \HH^i(X\otimes_A B,\sh{F}\otimes_A B)=\HH^i(X,\sh{F})\otimes_A B
  \tag{1.4.16.1}
\]
for every $i\geq 0$ and every quasi-coherent $\sh{O}_X$-module $\sh{F}$.
If $X\otimes_A B$ is affine, the left hand side of (\hyperref[III.1.4.16.1]{1.4.16.1}) is zero for $i=1$, hence so is $\HH^1(X,\sh{F})$ since $B$ is a faithfully flat $A$-module.
As $X$ is a quasi-compact scheme, we finish the proof by Serre's criterion \sref[II]{II.5.2.1}.
\end{proof}

\begin{proposition}[1.4.17]
\label{III.1.4.17}
Let $X$ be a prescheme, $0\to\sh{F}\xrightarrow{u}\sh{G}\xrightarrow{v}\sh{H}\to 0$ an exact sequence of $\sh{O}_X$-modules.
If $\sh{F}$ and $\sh{H}$ are quasi-coherent, then so is $\sh{G}$.
\end{proposition}

\oldpage[III]{95}
\begin{proof}
The question is local on $X$, so we can suppose that $X=\Spec(A)$ is affine, and it then suffices to prove that $\sh{G}$ satisfies the conditions (d1) and (d2) of \sref[I]{I.1.4.1} (with $V=X$).
The verification of (d2) is immediate, because if $t\in\Gamma(X,\sh{G})$ is zero when restricted to $D(f)$, then so is its image $v(t)\in\Gamma(X,\sh{H})$; therefore there exists an $m>0$ such that $f^m v(t)=v(f^m t)=0$ \sref[I]{I.1.4.1}, and as $\Gamma$ is left exact, $f^m t=u(s)$, where $s\in\Gamma(X,\sh{F})$; as $u$ is injective, the restriction of $s$ to $D(f)$ is zero, hence \sref[I]{I.1.4.1} there exists an integer $n>0$ such that $f^n s=0$; we finally deduce that $f^{m+n}t=u(f^n s)=0$.

We now check (d1); let $t'\in\Gamma(D(f),\sh{G})$; as $\sh{H}$ is quasi-coherent, there exists an integer $m$ such that $f^m v(t')=v(f^m t')$ extends to a section $z\in\Gamma(X,\sh{H})$ \sref[I]{I.1.4.1}.
But in virtue of Theorem~\sref{III.1.3.1} (or (\textbf{I},~5.1.9.2)) applied to the quasi-coherent $\sh{O}_X$-module $\sh{F}$, the sequence $\Gamma(X,\sh{G})\to\Gamma(X,\sh{H})\to 0$ is exact, so there exists $t\in\Gamma(X,\sh{G})$ such that $z=v(t)$; we thus see that $v(f^m t'-t'')=0$, denoting by $t''$ the restriction of $t$ to $D(f)$; thus we have $f^m t'-t''=u(s')$, where $s'\in\Gamma(D(f),\sh{F})$.
But as $\sh{F}$ is quasi-coherent, there exists an integer $n>0$ such that $f^n s'$ extends to a section $s\in\Gamma(X,\sh{F})$; as $f^{m+n}t'-f^n t''=u(f^n s')$, we see that $f^{m+n}t'$ is the restriction to $D(f)$ of a section $f^n t+u(s)\in\Gamma(X,\sh{G})$, which finishes the proof.
\end{proof}

\section{Cohomological study of projective morphisms}
\label{section:III.2}


\subsection{Explicit calculations of certain cohomology groups}
\label{subsection:III.2.1}

\begin{env}[2.1.1]
\label{III.2.1.1}
\oldpage[III]{95}
Let $X$ be a prescheme, $\sh{L}$ an invertible $\sh{O}_X$-module; consider the graded ring \sref[0]{0.5.4.6}
\[
\label{III.2.1.1.1}
  S = \Gamma_*(X,\sh{L}) = \bigoplus_{n\in\bb{Z}} \Gamma(X,\sh{L}^{\otimes n}).
  \tag{2.1.1.1}
\]

Let $(f_i)_{1\leq i\leq r}$ be a finite family of \emph{homogeneous} elements of $S$, say $f_i\in S_{d_i}$; set $U_i = X_{f_i}$, $U = \bigcup_i U_i$, and denote by $\mathfrak{U}$ the covering $(U_i)$ of $U$.
For every quasi-coherent $\sh{O}_X$-module $\sh{F}$, we set
\[
\label{III.2.1.1.2}
  \HH^\bullet(U,\sh{F}(*)) = \bigoplus_{n\in\bb{Z}} \HH^\bullet(U,\sh{F}\otimes\sh{L}^{\otimes n})
  \tag{2.1.1.2}
\]
\[
\label{III.2.1.1.3}
  \HH^\bullet(\mathfrak{U},\sh{F}(*)) = \bigoplus_{n\in\bb{Z}} \HH^\bullet(\mathfrak{U},\sh{F}\otimes\sh{L}^{\otimes n}).
  \tag{2.1.1.3}
\]

The abelian groups (\hyperref[III.2.1.1.2]{2.1.1.2}) and (\hyperref[III.2.1.1.3]{2.1.1.3}) are \emph{bigraded}, by taking
\[
  (\HH^\bullet(U,\sh{F}(*)))_{mn} = \HH^m(U,\sh{F}\otimes\sh{L}^{\otimes n})
\]
and an analogous definition for (\hyperref[III.2.1.1.3]{2.1.1.3}).
For the second grading, it is clear that these groups are graded $S$-modules, as follows, for example, from the fact that $\sh{F}\mapsto\HH^m(U,\sh{F})$ and $\sh{F}\mapsto\HH^m(\mathfrak{U},\sh{F})$ are functors.
\end{env}

\begin{env}[2.1.2]
\label{III.2.1.2}
Consider now the graded $S$-module \sref[0]{0.5.4.6}
\[
\label{III.2.1.2.1}
  M = \Gamma_*(\sh{L},\sh{F}) = \HH^0(X,\sh{F}(*)) = \bigoplus_{n\in\bb{Z}} \Gamma(X,\sh{F}\otimes\sh{L}^{\otimes n}).
  \tag{2.1.2.1}
\]

\oldpage[III]{96}
If $X$ is a prescheme whose underlying space is Noetherian, or a quasi-compact scheme, it follows from \sref[I]{I.9.3.1} that, setting as usual $U_{i_0 i_1\cdots i_p} = \bigcap_{k=0}^p U_{i_k}$, we have, up to canonical isomorphism,
\[
  \Gamma(U_{i_0 i_1\cdots i_p},\sh{F}(*)) = \HH^0(U_{i_0 i_1\cdots i_p},\sh{F}(*)) = M_{f_{i_0}f_{i_1}\cdots f_{i_p}}.
\]

We can further, with the notation of (\hyperref[III.1.2.2]{1.2.2}), identify $M_{f_{i_0}f_{i_1}\cdots f_{i_p}}$ with $\varinjlim_n M_{i_0 i_1\cdots i_p}^{(n)}$.
This identification is an isomorphism of \emph{graded} $S$-modules if we define the degree of a homogeneous element $z\in\varinjlim_n M_{i_0 i_1\cdots i_p}^{(n)}$ as follows: $z$ is the canonical image of a homogeneous element $x\in M_{i_0 i_1\cdots i_p}^{(n)} = M$ of degree $m$, and we take the degree of $z$ to be the number $m-n(d_{i_0}+d_{i_1}+\cdots+d_{i_p})$.
Taking into account the definition of the homomorphisms $\vphi_{kh}:M_{i_0 i_1\cdots i_p}^{(h)}\to M_{i_0 i_1\cdots i_p}^{(k)}$ (\hyperref[III.1.2.2]{1.2.2}), we see immediately that this definition does not depend on the ``representative'' $x$ of $z$ that was chosen.
Denoting, as in (\hyperref[III.1.2.2]{1.2.2}), by $C_n^p(M)$ the set of alternating maps from $[1,r]^{p+1}$ to $M$ (for every $n$), we similarly define a graded $S$-module structure on $\varinjlim_n C_n^p(M)$; as in (\hyperref[III.1.2.2]{1.2.2}), we then have
\[
\label{III.2.1.2.2}
  C^p(\mathfrak{U},\sh{F}(*)) = \varinjlim_n C_n^p(M)
  \tag{2.1.2.2}
\]
with the isomorphism between the two sides \emph{respecting degrees}.
We also have, as in (\hyperref[III.1.2.2]{1.2.2}),
\[
\label{III.2.1.2.3}
  C^p(\mathfrak{U},\sh{F}(*)) = C^{p+1}((\mathbf{f}),M) = \varinjlim_n K^{p+1}(\mathbf{f}^n,M)
  \tag{2.1.2.3}
\]
with the isomorphism \emph{preserving degrees}: the degree of an element of $\varinjlim_n K^{p+1}(\mathbf{f}^n,M)$, which is the canonical image of a cochain $g\in K^{p+1}(\mathbf{f}^n,M)$ whose values $g(i_0,\dots,i_p)$ lie in one and the same homogeneous component $M_m$ of $M$, is $m-n(d_{i_0}+\cdots+d_{i_p})$, independently of the choice of this cochain as a representative of the element in question.

Since the preceding isomorphisms are compatible with the coboundary operators, we conclude, as in (\hyperref[III.1.2.2]{1.2.2}), that:
\end{env}

\begin{proposition}[2.1.3]
\label{III.2.1.3}
Let $X$ be a prescheme whose underlying space is Noetherian, or a quasi-compact scheme.
There is a canonical isomorphism, functorial in $\sh{F}$,
\[
\label{III.2.1.3.1}
  \HH^p(\mathfrak{U},\sh{F}(*)) \simto \HH^{p+1}((\mathbf{f}),M)
  \qquad\text{for every } p\geq 1.
  \tag{2.1.3.1}
\]
Moreover, there is an exact sequence, functorial in $\sh{F}$,
\[
\label{III.2.1.3.2}
  0 \to \HH^0((\mathbf{f}),M) \to M \to \HH^0(\mathfrak{U},\sh{F}(*)) \to \HH^1((\mathbf{f}),M) \to 0.
  \tag{2.1.3.2}
\]
All the homomorphisms introduced are of degree $0$ for the graded $S$-module structures, where $S$ is the ring (\hyperref[III.2.1.1.1]{2.1.1.1}).
\end{proposition}

\begin{corollary}[2.1.4]
\label{III.2.1.4}
\oldpage[III]{97}
If $X$ is a quasi-compact scheme and the $U_i=X_{f_i}$ are affine, there is a canonical isomorphism, functorial in $\sh{F}$ and of degree $0$,
\[
\label{III.2.1.4.1}
  \HH^p(U,\sh{F}(*)) \simto \HH^{p+1}((\mathbf{f}),M)
  \qquad\text{for every } p\geq 1,
  \tag{2.1.4.1}
\]
and an exact sequence, functorial in $\sh{F}$,
\[
\label{III.2.1.4.2}
  0 \to \HH^0((\mathbf{f}),M) \to M \to \HH^0(U,\sh{F}(*)) \to \HH^1((\mathbf{f}),M) \to 0,
  \tag{2.1.4.2}
\]
in which all the homomorphisms are of degree $0$.
\end{corollary}

\begin{proof}
It suffices to apply \sref{III.1.4.1} to the result of \sref{III.2.1.3}.
\end{proof}

The ``local'' proposition analogous to \sref{III.2.1.3} is the following.

\begin{proposition}[2.1.5]
\label{III.2.1.5}
Let $S$ be a positively graded ring, let the $f_i$ ($1\leq i\leq r$) be homogeneous elements of $S_+$ of degrees $d_i$, and let $M$ be a graded $S$-module.
Let $X=\Proj(S)$ be the homogeneous prime spectrum of $S$, and set $U_i=\mathrm{D}_+(f_i)$, $U=\bigcup_i U_i$, and $\HH^\bullet(U,\widetilde{M}(*))=\bigoplus_{n\in\bb{Z}}\HH^\bullet(U,(M(n))^\sim)$.
There are canonical isomorphisms, functorial in $M$ and of degree $0$ for the graded $S$-module structures,
\[
\label{III.2.1.5.1}
  \HH^p(U,\widetilde{M}(*)) \simto \HH^{p+1}((\mathbf{f}),M)
  \qquad\text{for every } p\geq 1,
  \tag{2.1.5.1}
\]
and an exact sequence, functorial in $M$,
\[
\label{III.2.1.5.2}
  0 \to \HH^0((\mathbf{f}),M) \to M \to \HH^0(U,\widetilde{M}(*)) \to \HH^1((\mathbf{f}),M) \to 0,
  \tag{2.1.5.2}
\]
in which all the homomorphisms are of degree $0$.
\end{proposition}

\begin{proof}
Indeed, by definition \sref[II]{II.2.5.2},
\[
  \Gamma(U_{i_0i_1\cdots i_p},(M(n))^\sim)
  =(M_{f_{i_0}\cdots f_{i_p}})_n,
\]
and therefore $\Gamma(U_{i_0i_1\cdots i_p},\widetilde{M}(*))=M_{f_{i_0}\cdots f_{i_p}}$.
The rest of the argument is the same as in the proof of \sref{III.2.1.4}, taking into account that $X$ is a scheme.
\end{proof}

\begin{env}[2.1.6]
\label{III.2.1.6}
\emph{Remarks.} ---
(i) Under the hypotheses of \sref{III.2.1.5}, the functors $\Gamma(U_{i_0i_1\cdots i_p},\widetilde{M}(*))$ are exact in $M$, by \sref[0]{0.1.3.2}.
The same argument as in (\hyperref[III.1.2.5]{1.2.5}) therefore shows that, if $0\to M'\to M\to M''\to0$ is an exact sequence of graded $S$-modules whose homomorphisms have degree $0$, then for every $p\geq0$ there are commutative diagrams
\[
\label{III.2.1.6.1}
  \xymatrix{
    \HH^p(U,\widetilde{M''}(*)) \ar[r]^-\partial \ar[d] &
    \HH^{p+1}(U,\widetilde{M'}(*)) \ar[d] \\
    \HH^{p+1}((\mathbf{f}),M'') \ar[r]^-\partial &
    \HH^{p+2}((\mathbf{f}),M')
  }
  \tag{2.1.6.1}
\]

(ii) Proposition~\sref{III.2.1.5} is especially useful when $S$ is an $A$-algebra generated by finitely many elements of degree $1$, with $A$ Noetherian; in that case every quasi-coherent $\sh{O}_X$-module is of the form $\widetilde{M}$ \sref[II]{II.2.7.7}.
\end{env}

\begin{env}[2.1.7]
\label{III.2.1.7}
\oldpage[III]{98}
We now apply \sref{III.2.1.5} in the case $S=A[\mathrm{T}_0,\dots,\mathrm{T}_r]$, where $A$ is an arbitrary ring and the $\mathrm{T}_i$ are indeterminates, taking $M=S$ and $f_i=\mathrm{T}_i$.
We are thus essentially reduced to computing $\HH^\bullet((\mathbf{T}),S)$, where $\mathbf{T}=(\mathrm{T}_i)_{0\leq i\leq r}$.
\end{env}

\begin{lemma}[2.1.8]
\label{III.2.1.8}
If $S=A[\mathrm{T}_0,\dots,\mathrm{T}_r]$, then, with $\mathbf{T}=(\mathrm{T}_i)_{0\leq i\leq r}$,
\[
\label{III.2.1.8.1}
  \HH^i(\mathbf{T}^n,S)=0
  \qquad\text{if }i\neq r+1,
  \tag{2.1.8.1}
\]
\[
\label{III.2.1.8.2}
  \HH^{r+1}(\mathbf{T}^n,S)=S/(\mathbf{T}^n).
  \tag{2.1.8.2}
\]
The $A$-module $\HH^{r+1}(\mathbf{T}^n,S)$ therefore has an $A$-basis formed by the classes modulo $(\mathbf{T}^n)$ of the monomials
\[
  \mathrm{T}_{\mathbf{p}}=\mathrm{T}_0^{p_0}\cdots\mathrm{T}_r^{p_r},
  \qquad \mathbf{p}=(p_0,\dots,p_r),\quad 0\leq p_i<n\ \text{for every }i.
\]
\end{lemma}

\begin{proof}
This follows immediately from (\hyperref[III.1.1.3.5]{1.1.3.5}) and Proposition~\sref{III.1.1.4}, whose hypotheses are trivially satisfied.
\end{proof}

\begin{env}[2.1.9]
\label{III.2.1.9}
Passing to the inductive limit over $n$, relations (\hyperref[III.2.1.8.1]{2.1.8.1}) give $\HH^i((\mathbf{T}),S)=0$ for $i\neq r+1$.
For $i=r+1$, the inductive system consists of the modules $S/(\mathbf{T}^n)$, the homomorphism
\[
  \vphi_{nm}:S/(\mathbf{T}^n)\longrightarrow S/(\mathbf{T}^m)
  \qquad (0\leq n\leq m)
\]
being multiplication by $(\mathrm{T}_0\cdots\mathrm{T}_r)^{m-n}$.
For $n\geq\sup(p_i)_{0\leq i\leq r}$, denote by $\xi_{\mathbf{p}}^{(n)}=\xi_{p_0,\dots,p_r}^{(n)}$ the class of $\mathrm{T}_0^{n-p_0}\cdots\mathrm{T}_r^{n-p_r}$ modulo $(\mathbf{T}^n)$.
Then $\vphi_{nm}(\xi_{\mathbf{p}}^{(n)})=\xi_{\mathbf{p}}^{(m)}$, so these elements have the same canonical image $\xi_{\mathbf{p}}=\xi_{p_0,\dots,p_r}$ in the inductive limit $\HH^{r+1}((\mathbf{T}),S)$.
By the definition of degree in (\hyperref[III.2.1.2]{2.1.2}), the degree of $\xi_{\mathbf{p}}$ is $-|\mathbf{p}|=-(p_0+p_1+\cdots+p_r)$.
It is clear that the $\xi_{\mathbf{p}}^{(n)}$ with $0<p_i\leq n$ for $0\leq i\leq r$ form a basis of $S/(\mathbf{T}^n)$.
We therefore deduce immediately from \sref{III.2.1.8}:
\end{env}

\begin{corollary}[2.1.10]
\label{III.2.1.10}
With the notation of \sref{III.2.1.8},
\[
\label{III.2.1.10.1}
  \HH^i((\mathbf{T}),S)=0
  \qquad\text{for }i\neq r+1,
  \tag{2.1.10.1}
\]
and $\HH^{r+1}((\mathbf{T}),S)$ is a free $A$-module with basis the elements $\xi_{p_0,\dots,p_r}$ such that $p_i>0$ for $0\leq i\leq r$.
\end{corollary}

\begin{env}[2.1.11]
\label{III.2.1.11}
\emph{Remark.} ---
Let $N$ be any $A$-module and let $M=S\otimes_A N$.
The argument of \sref{III.2.1.8} shows more generally that
\[
\label{III.2.1.11.1}
  \HH^i(\mathbf{T}^n,M)=0
  \qquad\text{if }i\neq r+1,
  \tag{2.1.11.1}
\]
\[
\label{III.2.1.11.2}
  \HH^{r+1}(\mathbf{T}^n,M)=(S/(\mathbf{T}^n))\otimes_A N.
  \tag{2.1.11.2}
\]
Indeed, the latter formula follows directly from (\hyperref[III.1.1.3.5]{1.1.3.5}).
Moreover, it is clear that
\[
  M/(\mathrm{T}_0^nM+\cdots+\mathrm{T}_{i-1}^nM)
  =\bigl(S/(\mathrm{T}_0^nS+\cdots+\mathrm{T}_{i-1}^nS)\bigr)\otimes_A N,
\]
since the ideal $\mathrm{T}_0^nS+\cdots+\mathrm{T}_{i-1}^nS$ is a direct summand of the $A$-module $S$.
This permits us to apply \sref{III.1.1.4} to $M$, and yields (\hyperref[III.2.1.11.1]{2.1.11.1}).
\end{env}

Combining \sref{III.2.1.10} and \sref{III.2.1.5}, we obtain:

\begin{proposition}[2.1.12]
\label{III.2.1.12}
Let $A$ be a ring, let $r$ be an integer $>0$, and let $X=\bb{P}_A^r$ \sref[II]{II.4.1.1}.
Then:
\begin{enumerate}
  \item[{\rm(i)}] $\HH^i(X,\sh{O}_X(*))=0$ for $i\neq0,r$.
  \item[{\rm(ii)}] The canonical homomorphism $\alpha:S\to\HH^0(X,\sh{O}_X(*))$ \sref[II]{II.2.6.2} is bijective.
\oldpage[III]{99}
  \item[{\rm(iii)}] $\HH^r(X,\sh{O}_X(*))$ is a free $A$-module with basis the elements $\xi_{p_0,\dots,p_r}$, where $p_i>0$ for $0\leq i\leq r$; the element $\xi_{p_0,\dots,p_r}$ has degree $-|\mathbf{p}|=-(p_0+\cdots+p_r)$, and
  \[
    \mathrm{T}_i\xi_{p_0,\dots,p_r}
    =\xi_{p_0,\dots,p_{i-1},p_i-1,p_{i+1},\dots,p_r}.
  \]
\end{enumerate}
\end{proposition}

\begin{proof}
In the exact sequence (\hyperref[III.2.1.5.2]{2.1.5.2}) applied to $M=S=A[\mathrm{T}_0,\dots,\mathrm{T}_r]$, we have $\HH^0((\mathbf{T}),S)=0$ and $\HH^1((\mathbf{T}),S)=0$ by \sref{III.2.1.10.1}; moreover, Proposition~\sref{III.2.1.5} applies with $U=X$, since $X$ is the union of the $\mathrm{D}_+(\mathrm{T}_i)$ \sref[II]{II.2.3.14}.
It remains to identify the map $S\to\HH^0(X,\sh{O}_X(*))$ in (\hyperref[III.2.1.5.2]{2.1.5.2}) with the canonical map $\alpha$.
This follows from the canonical identification of $\HH^0(U,\sh{O}_X(*))$ and $\HH^0(\mathfrak{U},\sh{O}_X(*))$.
\end{proof}

\begin{corollary}[2.1.13]
\label{III.2.1.13}
The only pairs $(i,n)$ for which $\HH^i(X,\sh{O}_X(n))\neq0$ can occur are $i=0$, $n\geq0$, and $i=r$, $n\leq-(r+1)$.
\end{corollary}

If $A\neq0$, then $\HH^i(X,\sh{O}_X(n))\neq0$ does indeed hold for all the pairs listed in \sref{III.2.1.13}.
This follows from \sref{III.2.1.12}, since $S_n\neq0$ for every degree $n\geq0$.

In the applications made in this chapter, we will mostly use the following less precise result.

\begin{corollary}[2.1.14]
\label{III.2.1.14}
The $A$-modules $\HH^i(X,\sh{O}_X(n))$ are free of finite type; if $i>0$, they vanish for $n>0$.
\end{corollary}

\begin{proposition}[2.1.15]
\label{III.2.1.15}
Let $Y$ be a prescheme, let $\sh{E}$ be a locally free $\sh{O}_Y$-module of rank $r+1$, let $X=\bb{P}(\sh{E})$ be the projective bundle defined by $\sh{E}$, and let $f:X\to Y$ be the structural morphism.
The only values of $i$ and $n$ for which $\RR^if_*(\sh{O}_X(n))\neq0$ can occur are $i=0$, $n\geq0$, and $i=r$, $n\leq-(r+1)$.
Moreover, the canonical homomorphism \sref[II]{II.3.3.2}
\[
  \alpha:\bb{S}_{\sh{O}_Y}(\sh{E})
  \longrightarrow \Gamma_*(\sh{O}_X)
  =\RR^0f_*(\sh{O}_X(*))
  =\bigoplus_{n\in\bb{Z}}f_*(\sh{O}_X(n))
\]
is an isomorphism.
\end{proposition}

\begin{proof}
The question is local on $Y$, so we may assume that $Y$ is affine with ring $A$ and that $\sh{E}=\widetilde{E}$, where $E=A^{r+1}$.
We are then immediately reduced to \sref{III.2.1.12}, taking \sref{III.1.4.11} into account.
\end{proof}

\begin{env}[2.1.16]
\label{III.2.1.16}
\emph{Remark.} ---
We will later complete the results of \sref{III.2.1.15} by proving the following propositions.
Set
\[
  \boldsymbol{\omega}=f^*\bigl(\bigwedge^{r+1}\sh{E}\bigr)(-r-1),
\]
which is an invertible $\sh{O}_X$-module.
Then:
\begin{enumerate}
  \item[{\rm(i)}] There is a canonical isomorphism
  \[
  \label{III.2.1.16.1}
    \rho:\sh{O}_Y \simto \RR^rf_*(\boldsymbol{\omega}).
    \tag{2.1.16.1}
  \]
  \item[{\rm(ii)}] The cup-product pairing \sref[0]{0.12.2.2}
  \[
  \label{III.2.1.16.2}
    \RR^rf_*(\sh{O}_X(n))
    \times \RR^0f_*(\boldsymbol{\omega}(-n))
    \longrightarrow \RR^rf_*(\boldsymbol{\omega}),
    \tag{2.1.16.2}
  \]
  composed with $\rho^{-1}$, defines an \emph{isomorphism} from $\RR^rf_*(\sh{O}_X(n))$ onto the \emph{dual} of the locally free $\sh{O}_Y$-module
  \[
    \RR^0f_*(\boldsymbol{\omega}(-n))
    =\bigl(\bigwedge^{r+1}\sh{E}\bigr)
      \otimes_{\sh{O}_Y}
      \bigl(\bb{S}_{\sh{O}_Y}(\sh{E})\bigr)_{-n-r-1}.
  \]
\end{enumerate}
\end{env}

\subsection{The fundamental theorem of projective morphisms}
\label{subsection:III.2.2}
\oldpage[III]{100}

\begin{theorem}[2.2.1]
\label{III.2.2.1}
(Serre).
Let $Y$ be a Noetherian prescheme, let $f:X\to Y$ be a proper morphism, and let $\sh{L}$ be an invertible $\sh{O}_X$-module ample relative to $f$.
For every $\sh{O}_X$-module $\sh{F}$, set
\[
  \sh{F}(n)=\sh{F}\otimes_{\sh{O}_X}\sh{L}^{\otimes n}
  \qquad(n\in\bb{Z}).
\]
Then, for every coherent $\sh{O}_X$-module $\sh{F}$:
\begin{enumerate}
  \item[{\rm(i)}] the $\sh{O}_Y$-modules $\RR^qf_*(\sh{F})$ are coherent;
  \item[{\rm(ii)}] there exists an integer $N$ such that, for $n\geq N$,
  \[
    \RR^qf_*(\sh{F}(n))=0
    \qquad\text{for every }q>0;
  \]
  \item[{\rm(iii)}] there exists an integer $N$ such that, for $n\geq N$, the canonical homomorphism
  \[
    f^*\bigl(f_*(\sh{F}(n))\bigr)\longrightarrow\sh{F}(n)
  \]
  is surjective.
\end{enumerate}
\end{theorem}

\begin{proof}
First note that if the theorem is true after replacing $\sh{L}$ by $\sh{L}^{\otimes d}$, where $d>0$, then it is true in its original form.
Indeed, writing
\[
  \sh{F}(n)
  =\bigl(\sh{F}\otimes\sh{L}^{\otimes r}\bigr)
   \otimes\sh{L}^{\otimes hd},
   \qquad h>0,\quad 0\leq r<d,
\]
the hypothesis gives, for each $r$, an integer $N_r$ such that (ii) and (iii) hold for $\sh{F}\otimes\sh{L}^{\otimes r}$ whenever $h\geq N_r$.
Taking $N$ to be the largest of the integers $dN_r$, assertions (ii) and (iii) therefore hold for $n\geq N$.

We may consequently suppose that $\sh{L}$ is very ample relative to $f$ \sref[II]{II.4.6.11}.
There is then a dominant open immersion over $Y$
\[
  i:X\longrightarrow P,
  \qquad P=\Proj(\sh{S}),
\]
where $\sh{S}$ is a quasi-coherent $\sh{O}_Y$-algebra graded in positive degrees, $\sh{S}_1$ is of finite type and generates $\sh{S}$, and $\sh{L}$ is isomorphic to $i^*(\sh{O}_P(1))$ \sref[II]{II.4.4.7}.
Since $f$ is proper, $i$ is proper as well \sref[II]{II.5.4.4}, and hence is an isomorphism $X\simto P$.
We are thus reduced to the case $X=\Proj(\sh{S})$ and $\sh{L}=\sh{O}_X(1)$.
The theorem is then a consequence of the following proposition.
\end{proof}

\begin{proposition}[2.2.2]
\label{III.2.2.2}
Let $A$ be a Noetherian ring and let $S$ be an $A$-algebra graded in positive degrees such that the $A$-module $S_1$ has a system of $r+1$ generators and generates the algebra $S$.
Set $X=\Proj(S)$.
For every coherent $\sh{O}_X$-module $\sh{F}$:
\begin{enumerate}
  \item[{\rm(i)}] the $A$-modules $\HH^q(X,\sh{F})$ are of finite type;
  \item[{\rm(ii)}] $\HH^q(X,\sh{F})=0$ for $q>r$;
  \item[{\rm(iii)}] there exists an integer $N$ such that, for $n\geq N$,
  \[
    \HH^q(X,\sh{F}(n))=0
    \qquad\text{for every }q>0;
  \]
  \item[{\rm(iv)}] there exists an integer $N$ such that, for $n\geq N$, $\sh{F}(n)$ is generated by its sections over $X$.
\end{enumerate}
\end{proposition}

\begin{proof}
Let us first show how \sref{III.2.2.2} implies \sref{III.2.2.1}.
In the special case $X=\Proj(\sh{S})$ to which \sref{III.2.2.1} was reduced above, $Y$ is quasi-compact.
It can therefore be covered by finitely many affine open subsets $U_\alpha$, with Noetherian rings, such that the restriction of $\sh{S}_1$ to each $U_\alpha$ is generated by finitely many sections over $U_\alpha$.
Once \sref{III.2.2.2} is known, it is enough to take in parts (ii) and (iii) of \sref{III.2.2.1} the largest of the corresponding integers for the $U_\alpha$, taking \sref{III.1.4.11} and \sref[II]{II.3.4.7} into account.

To prove \sref{III.2.2.2}, observe that $X$ identifies with a closed subprescheme of $P=\bb{P}_A^r$ \sref[II]{II.3.6.2}.
Moreover, if $j:X\to P$ is the canonical injection, then $j_*(\sh{F})$ is a coherent $\sh{O}_P$-module and
\[
  j_*(\sh{F}(n))=(j_*\sh{F})(n)
\]
by \sref[II]{II.3.4.5} and \sref[II]{II.3.5.2}.
In view of (G, II, Corollary to Theorem~4.9.1), we are therefore reduced to proving \sref{III.2.2.2} in the special case
\[
  X=\bb{P}_A^r,
  \qquad S=A[T_0,\ldots,T_r].
\]

\oldpage[III]{101}

Since $X$ is covered by the $r+1$ affine open subsets $D_+(T_i)$, assertion (ii) follows from \sref{III.1.4.12}.
Assertion (iv) has already been proved in \sref[II]{II.2.7.9}.

We now prove (i) and (iii) simultaneously.
They hold for $\sh{F}=\sh{O}_X(m)$ by \sref{III.2.1.13}, and hence also when $\sh{F}$ is a direct sum of finitely many $\sh{O}_X$-modules of the form $\sh{O}_X(m_j)$.
They are trivially true for $q>r$ by (ii).
We proceed by descending induction on $q$.

The sheaf $\sh{F}$ is isomorphic to a quotient of a direct sum $\sh{E}$ of finitely many sheaves $\sh{O}_X(m_j)$ \sref[II]{II.2.7.10}.
Thus there is an exact sequence
\[
  0\longrightarrow\sh{R}\longrightarrow\sh{E}
   \longrightarrow\sh{F}\longrightarrow0,
\]
where $\sh{R}$ is coherent \sref[0]{0.5.3.3} and $\sh{E}$ satisfies (i) and (iii).
Since twisting is exact, for every $n\in\bb{Z}$ we also have an exact sequence
\[
  0\longrightarrow\sh{R}(n)\longrightarrow\sh{E}(n)
   \longrightarrow\sh{F}(n)\longrightarrow0,
\]
and hence an exact cohomology sequence
\[
  \HH^{q-1}(X,\sh{E}(n))
  \longrightarrow\HH^{q-1}(X,\sh{F}(n))
  \longrightarrow\HH^q(X,\sh{R}(n)).
\]
Because $\sh{E}(n)$ is a direct sum of sheaves $\sh{O}_X(n+m_j)$ \sref[II]{II.2.5.14}, the module $\HH^{q-1}(X,\sh{E}(n))$ is of finite type.
The module $\HH^q(X,\sh{R}(n))$ is likewise of finite type by the induction hypothesis.
Since $A$ is Noetherian, it follows that $\HH^{q-1}(X,\sh{F}(n))$ is of finite type, proving (i).

By the induction hypothesis we may choose $N$ so that $\HH^q(X,\sh{R}(n))=0$ for $n\geq N$.
We may also choose $N$ so that $\HH^{q-1}(X,\sh{E}(n))=0$ for $n\geq N$, since $\sh{E}$ satisfies (iii).
The exact cohomology sequence then gives
\[
  \HH^{q-1}(X,\sh{F}(n))=0
  \qquad(n\geq N),
\]
which completes the proof.
\end{proof}

\begin{corollary}[2.2.3]
\label{III.2.2.3}
Under the hypotheses of \sref{III.2.2.1}, let
\[
  \sh{F}\longrightarrow\sh{G}\longrightarrow\sh{H}
\]
be an exact sequence of coherent $\sh{O}_X$-modules.
There exists an integer $N$ such that, for $n\geq N$, the sequence
\[
  f_*(\sh{F}(n))
  \longrightarrow f_*(\sh{G}(n))
  \longrightarrow f_*(\sh{H}(n))
\]
is exact.
\end{corollary}

\begin{proof}
Let $\sh{F}'$, $\sh{G}'$, and $\sh{G}''$ be respectively the kernel, image, and cokernel of $\sh{F}\to\sh{G}$.
Then $\sh{G}'$ is the kernel and $\sh{G}''$ the image of $\sh{G}\to\sh{H}$; let $\sh{H}''$ be the cokernel of this latter homomorphism.
All these $\sh{O}_X$-modules are coherent \sref[0]{0.5.3.4}.
Since twisting is exact, it is enough to prove that, for $n$ sufficiently large, each of the sequences
\begin{align*}
  0&\longrightarrow f_*(\sh{F}'(n))
    \longrightarrow f_*(\sh{F}(n))
    \longrightarrow f_*(\sh{G}'(n))
    \longrightarrow0,\\
  0&\longrightarrow f_*(\sh{G}'(n))
    \longrightarrow f_*(\sh{G}(n))
    \longrightarrow f_*(\sh{G}''(n))
    \longrightarrow0,\\
  0&\longrightarrow f_*(\sh{G}''(n))
    \longrightarrow f_*(\sh{H}(n))
    \longrightarrow f_*(\sh{H}''(n))
    \longrightarrow0
\end{align*}
is exact.
Consequently, we may suppose that
\[
  0\longrightarrow\sh{F}\longrightarrow\sh{G}
   \longrightarrow\sh{H}\longrightarrow0
\]
is exact.
The cohomology sequence is then
\[
  0\longrightarrow f_*(\sh{F}(n))
  \longrightarrow f_*(\sh{G}(n))
  \longrightarrow f_*(\sh{H}(n))
  \longrightarrow \RR^1f_*(\sh{F}(n))
  \longrightarrow\cdots,
\]
and the conclusion follows from \sref{III.2.2.1}, (ii).
\end{proof}

\oldpage[III]{102}

\begin{corollary}[2.2.4]
\label{III.2.2.4}
Let $Y$ be a Noetherian prescheme, let $f:X\to Y$ be a morphism of finite type, and let $\sh{L}$ be an invertible $\sh{O}_X$-module ample relative to $f$.
For every $\sh{O}_X$-module $\sh{F}$, set
\[
  \sh{F}(n)=\sh{F}\otimes_{\sh{O}_X}\sh{L}^{\otimes n}
  \qquad(n\in\bb{Z}).
\]
Let
\[
  \sh{F}\longrightarrow\sh{G}\longrightarrow\sh{H}
\]
be an exact sequence of coherent $\sh{O}_X$-modules such that the support of $\operatorname{Im}(\sh{F}\to\sh{G})$ is proper over $Y$ \sref[II]{II.5.4.10}.
Then there exists an integer $N$ such that, for $n\geq N$, the sequence
\[
  f_*(\sh{F}(n))
  \longrightarrow f_*(\sh{G}(n))
  \longrightarrow f_*(\sh{H}(n))
\]
is exact.
\end{corollary}

\begin{proof}
The argument at the beginning of \sref{III.2.2.1} shows that if the corollary is true for $\sh{L}^{\otimes d}$, where $d>0$, then it is true for $\sh{L}$.
We may therefore suppose that $\sh{L}$ is very ample relative to $f$ \sref[II]{II.4.6.11}.
We can then identify $X$ with an open subset of a $Y$-prescheme
\[
  Z=\Proj(\sh{S}),
\]
where $\sh{S}$ is a quasi-coherent $\sh{O}_Y$-algebra graded in positive degrees, $\sh{S}_1$ is of finite type and generates $\sh{S}$, and
\[
  \sh{L}=i^*(\sh{O}_Z(1)),
\]
with $i:X\to Z$ the canonical immersion \sref[II]{II.4.4.7}.

Set
$\sh{G}_1=\operatorname{Im}(\sh{F}\to\sh{G})
=\operatorname{Ker}(\sh{G}\to\sh{H})$, which is coherent. Since we have the
exact sequence $0\to\sh{G}_1(n)\to\sh{G}(n)\to\sh{H}(n)$ and the functor
$f_*$ is left exact, the sequence
$0\to f_*(\sh{G}_1(n))\to f_*(\sh{G}(n))\to f_*(\sh{H}(n))$ is exact for every
$n$; it therefore suffices to show that, for sufficiently large $n$, the
sequence $f_*(\sh{F}(n))\to f_*(\sh{G}_1(n))\to0$ is exact. In other words, we
may restrict to the case where $\sh{H}=0$ and the support of $\sh{G}$ is
proper over $Y$, hence \emph{closed in $Z$} (II, 5.4.10). Consequently
$\sh{G}'=i_*(\sh{G})$ is a coherent $\sh{O}_Z$-Module such that
$\sh{G}'|X=\sh{G}$. We know from (I, 9.4.3) that there exists a coherent
$\sh{O}_Z$-Module $\sh{F}'$ such that $\sh{F}'|X=\sh{F}$. Moreover, since $X$
is open in $Z$ and $\operatorname{Supp}(\sh{G})\subset X$ is closed in $Z$, it
is immediate that one defines a surjective homomorphism
$u':\sh{F}'\to\sh{G}'$ of sheaves such that $u'|X$ is the given surjective
homomorphism $u:\sh{F}\to\sh{G}$, by taking $u'|U=0$ for every open subset $U$
of $Z$ not meeting $\operatorname{Supp}(\sh{G})$ and $u'|U=u|U$ for every open
subset $U\subset X$. This being so, as at the beginning of the proof of
(2.2.2), we see that we may reduce to the case where $Y$ is affine; it is
therefore enough to show that, for sufficiently large $n$, the homomorphism
$\Gamma(u):\Gamma(X,\sh{F}(n))\to\Gamma(X,\sh{G}(n))$ is surjective. We have
the commutative diagram
\[
\xymatrix@C=35mm@R=12mm{
\Gamma(Z,\sh{F}'(n)) \ar[r]^{\Gamma(u')} \ar[d]_v &
\Gamma(Z,\sh{G}'(n)) \ar[d]^w \\
\Gamma(X,\sh{F}(n)) \ar[r]_{\Gamma(u)} & \Gamma(X,\sh{G}(n))
}
\]
where $v$ and $w$ are the restriction homomorphisms. The definition of
$\sh{G}'$ moreover shows that $w$ is \emph{bijective}; and by (2.2.3),
$\Gamma(u')$ is \emph{surjective} for sufficiently large $n$, so the same is
true of $\Gamma(u)$.
\end{proof}

\begin{env}[2.2.5]
\label{III.2.2.5}
\emph{Remarks.}

(i) Assertion (i) of \sref{III.2.2.1} remains true if one assumes only that $Y$ is \emph{locally Noetherian}.
Indeed, the property is plainly local on $Y$.
Moreover, the hypotheses of \sref{III.2.2.1} imply that, for every open subset $U\subset Y$, the restriction of $f$ to $f^{-1}(U)$ is a projective morphism
\[
  f^{-1}(U)\longrightarrow U
\]
\sref[II]{II.5.5.5}[iii], and $\sh{L}|_{f^{-1}(U)}$ is ample for this morphism \sref[II]{II.4.6.4}.

(ii) Assertion (iii) of \sref{III.2.2.1} remains valid, as we have seen, if one assumes only that $X$ is a quasi-compact scheme, or a prescheme whose underlying space is Noetherian, and that $f:X\to Y$ is quasi-compact \sref[II]{II.4.6.8}.
It should be noted, however, that even when $Y$ is the \emph{spectrum of a field} $K$ and $f$ is \emph{quasi-projective}, assertion (ii) of \sref{III.2.2.1} need no longer hold.

For example, let
\[
  X'=\Spec(K[T_0,\ldots,T_r])
\]
and let $X$ be the union of the affine open subsets $D(T_i)$ of $X'$, for $0\leq i\leq r$.
Since the immersion $X\to X'$ is quasi-compact, the structural morphism $f:X\to Y$ is quasi-affine \sref[II]{II.5.1.10}, and hence $\sh{O}_X$ is very ample relative to $f$ \sref[II]{II.5.1.6}.
But
\[
  \Gamma(X,\sh{O}_X)
  =\bigcap_{i=0}^r(K[T_0,\ldots,T_r])_{T_i}
  =K[T_0,\ldots,T_r]
\]
by \sref[I]{I.8.2.1.1}.
It follows from \sref{III.1.4.3.1} and \sref{III.1.1.3.5} that
\[
  \HH^r(X,\sh{O}_X^{\otimes n})
  =\HH^r(X,\sh{O}_X)
  \neq0
\]
for every $n$.
(iii) Under the hypotheses of (2.2.4) concerning $X$, $Y$, $f$, and $\sh{L}$,
it is immediate that if $\sh{H}$ is a coherent $\sh{O}_X$-Module whose support
is \emph{proper over $Y$}, an argument analogous to that of (2.2.4) shows
that there exists an integer $N$ such that, for $n\geq N$,
$\RR^if_*(\sh{H}(n))=0$ for every $i>0$. It follows from the exact
cohomology sequence that if $u:\sh{F}\to\sh{G}$ is a homomorphism of coherent
$\sh{O}_X$-Modules such that $\Ker(u)$ and $\Coker(u)$ have supports
\emph{proper over $Y$}, then there exists $N$ such that for $n\geq N$, the
corresponding homomorphism
$\RR^if_*(\sh{F}(n))\to\RR^if_*(\sh{G}(n))$ is \emph{bijective} for every
$i>0$.
\end{env}


\subsection{Application to graded sheaves of algebras and modules}
\label{subsection:III.2.3}

\begin{theorem}[2.3.1]
\label{III.2.3.1}
Let $Y$ be a Noetherian prescheme, let $\sh{S}$ be a quasi-coherent $\sh{O}_Y$-algebra of finite type, graded in positive degrees, let
\[
  X=\Proj(\sh{S}),
\]
and let $q:X\to Y$ be the structural morphism.
Let $\sh{M}$ be a quasi-coherent graded $\sh{S}$-module satisfying condition \textbf{(TF)}.
There then exists an integer $N$ such that, for $n\geq N$, the canonical homomorphism \sref[II]{II.8.14.5.1}
\[
  \alpha_n:\sh{M}_n
  \longrightarrow q_*\bigl(\shProj_0(\sh{M}(n))\bigr)
  =q_*\bigl((\shProj(\sh{M}))_n\bigr)
\]
\oldpage[III]{103}
is bijective.
In other words, the canonical homomorphism
\[
  \alpha:\sh{M}\longrightarrow\bbGamma_*(\shProj(\sh{M}))
\]
is a \textbf{(TN)}-isomorphism.
\end{theorem}

\begin{proof}
We may restrict to the case in which $\sh{M}$ is an $\sh{S}$-module of finite type \sref[II]{II.3.4.2}.
Since $Y$ is quasi-compact, there is an integer $d>0$ such that $\sh{S}^{(d)}$ is generated by the quasi-coherent $\sh{O}_Y$-module $\sh{S}_d$, which is of finite type \sref[II]{II.3.1.10}, and hence coherent because $Y$ is Noetherian.
Now $\sh{M}$ is the direct sum of the $\sh{M}^{(d,k)}$, for $0\leq k<d$, and each $\sh{M}^{(d,k)}$ is a quasi-coherent $\sh{S}^{(d)}$-module of finite type, by \sref[II]{II.2.1.6}[iii].
Since the question is local on $Y$, it plainly suffices to prove that each canonical homomorphism
\[
  \alpha:\sh{M}^{(d,k)}
  \longrightarrow
  \bbGamma_*\bigl((\shProj(\sh{M}))^{(d,k)}\bigr)
\]
is a \textbf{(TN)}-isomorphism.
Taking \sref[II]{II.8.14.13} into account, and in particular diagram \sref[II]{II.8.14.13.4}, we are thus reduced to proving the theorem when $\sh{S}$ is generated by $\sh{S}_1$ and $\sh{S}_1$ is a coherent $\sh{O}_Y$-module.

Since $Y$ is Noetherian, the argument used at the beginning of the proof of \sref{III.2.2.2} shows that we may reduce to the case
\[
  Y=\Spec(A),\qquad
  \sh{S}=\widetilde{S},\qquad
  \sh{M}=\widetilde{M},
\]
where $A$ is a Noetherian ring, $S_1$ is an $A$-module of finite type, and $M$ is a graded $S$-module of finite type.
We shall show that it then suffices to prove the theorem for $M=S$.
Indeed, in the general case there is an exact sequence
\[
  L'\longrightarrow L\longrightarrow M\longrightarrow0,
\]
where $L$ and $L'$ are direct sums of graded modules of the form $S(m)$.
If the result holds for $M=S$, it also holds for $M=S(m)$, and hence for $L$ and $L'$.
Consider the commutative diagram
\[
  \xymatrix{
    \widetilde{L'_n}\ar[r]\ar[d]_{\alpha_n} &
    \widetilde{L_n}\ar[r]\ar[d]^{\alpha_n} &
    \widetilde{M_n}\ar[r]\ar[d]^{\alpha_n} &
    0\\
    q_*(\widetilde{L'}(n))\ar[r] &
    q_*(\widetilde{L}(n))\ar[r] &
    q_*(\widetilde{M}(n))\ar[r] &
    0 .
  }
\]
\oldpage[III]{104}
The second row is exact by \sref{III.2.2.3} once $n$ is sufficiently large.
The same is true of the first row, and since the two left-hand vertical arrows are isomorphisms, the third is as well.

It therefore remains to prove the theorem for $M=S$.
Suppose first that
\[
  S=A[T_0,\ldots,T_r],
\]
where the $T_i$ are indeterminates.
In this case the assertion is precisely \sref{III.2.1.11}[ii].
In the general case, $S$ identifies with a quotient of a ring
\[
  S'=A[T_0,\ldots,T_r]
\]
by a graded ideal, so $X$ is a closed subscheme of
\[
  X'=\bb{P}_A^r
\]
\sref[II]{II.2.9.2}.
If $j:X\to X'$ is the canonical injection, then $j_*(\widetilde{S}(n))$ is the $\sh{O}_{X'}$-module
\[
  (\shProj(\widetilde{S}))(n),
\]
where $S$ is regarded as a graded $S'$-module; this follows at once from \sref[II]{II.2.8.7}.
Since $j_*(\widetilde{S}(n))$ is an $\sh{O}_{X'}$-module satisfying \textbf{(TF)}, the canonical homomorphism
\[
  \alpha_n:S_n\longrightarrow
  \Gamma\bigl(X',j_*(\widetilde{S}(n))\bigr)
\]
is bijective for all sufficiently large $n$, by the polynomial case just proved.
This completes the proof, since
\[
  \Gamma\bigl(X',j_*(\widetilde{S}(n))\bigr)
  =\Gamma(X,\widetilde{S}(n)).
\]
\end{proof}

\begin{corollary}[2.3.2]
\label{III.2.3.2}
Under the hypotheses of \sref{III.2.3.1}, let
\[
  \sh{S}_X=\bigoplus_{n\in\bb{Z}}\sh{O}_X(n),
\]
and let $\sh{F}$ be a quasi-coherent graded $\sh{S}_X$-module of finite type.
Then $\bbGamma_*(\sh{F})$ satisfies condition \textbf{(TF)}.
\end{corollary}

\begin{proof}
In the proof of \sref{III.2.3.1} we saw that $X$, which is isomorphic to $\Proj(\sh{S}^{(d)})$ \sref[II]{II.3.1.8}, is of finite type over $Y$ \sref[II]{II.3.4.1}.
It follows from \sref[II]{II.8.14.9} that $\sh{F}$ is isomorphic to a graded $\sh{S}_X$-module of the form $\shProj(\sh{M})$, where $\sh{M}$ is a quasi-coherent graded $\sh{S}$-module of finite type.
By \sref{III.2.3.1}, $\bbGamma_*(\sh{F})$ is \textbf{(TN)}-isomorphic to $\sh{M}$ and consequently satisfies \textbf{(TF)}.
\end{proof}

\begin{env}[2.3.3]
\label{III.2.3.3}
\emph{Scholium.}
Let $Y$ be a Noetherian prescheme, let $\sh{S}$ be a graded $\sh{O}_Y$-algebra satisfying the hypotheses of \sref{III.2.3.1}, and let $X=\Proj(\sh{S})$.
Let $\cat{K}_{\sh{S}}$ be the abelian category of quasi-coherent graded $\sh{S}$-modules satisfying \textbf{(TF)}, and let $\cat{K}'_{\sh{S}}$ be the subcategory of $\cat{K}_{\sh{S}}$ consisting of the $\sh{S}$-modules satisfying \textbf{(TN)}.
Finally, let $\cat{K}_X$ be the category of quasi-coherent graded $\sh{S}_X$-modules $\sh{F}$ of finite type; since $\sh{S}_X$ is periodic \sref[II]{II.8.14.4} and \sref[II]{II.8.14.12}, this is equivalent to requiring the $\sh{F}_i$ to be coherent $\sh{O}_X$-modules.

By \sref[II]{II.8.14.8}, \sref[II]{II.8.14.10}, and \sref{III.2.3.2}, the functors
\[
  \sh{M}\longmapsto\shProj(\sh{M}),
  \qquad
  \sh{F}\longmapsto\bbGamma_*(\sh{F})
\]
define an equivalence (T, I, 1.2) between the quotient category
\[
  \cat{K}_{\sh{S}}/\cat{K}'_{\sh{S}}
\]
(T, I, 1.11) and the category $\cat{K}_X$.
When $\sh{S}$ is generated by $\sh{S}_1$, the category $\cat{K}_X$ may be replaced by the category of coherent $\sh{O}_X$-modules \sref[II]{II.8.14.12}.
\end{env}

\begin{proposition}[2.3.4]
\label{III.2.3.4}
Let $Y$ be a Noetherian prescheme.

(i) Let $\sh{S}$ be a quasi-coherent $\sh{O}_Y$-algebra of finite type, graded in positive degrees.
Set
\[
  X=\Proj(\sh{S}),\qquad
  \sh{S}_X=\shProj(\sh{S})
  =\bigoplus_{n\in\bb{Z}}\sh{O}_X(n).
\]
Then $\sh{S}_X$ is a periodic quasi-coherent graded $\sh{O}_X$-algebra \sref[II]{II.8.14.12}, its homogeneous components
\[
  (\sh{S}_X)_n=\sh{O}_X(n)
\]
are coherent $\sh{O}_X$-modules, and, if $d>0$ is a period of $\sh{S}_X$, then
\[
  (\sh{S}_X)_d=\sh{O}_X(d)
\]
is an invertible $Y$-ample $\sh{O}_X$-module.
Moreover, the canonical homomorphism
\[
  \alpha:\sh{S}\longrightarrow\bbGamma_*(\sh{S}_X)
\]
is a \textbf{(TN)}-isomorphism.

(ii) Conversely, let $q:X\to Y$ be a projective morphism, and let $\sh{S}'$ be a graded $\sh{O}_X$-algebra whose homogeneous components $\sh{S}'_n$, for $n\in\bb{Z}$, are coherent $\sh{O}_X$-modules.
Suppose that $\sh{S}'$ has a period $d>0$ such that $\sh{S}'_d$ is an invertible $\sh{O}_X$-module ample for $q$.
Then
\[
  \sh{S}=\bigoplus_{n\geq0}q_*(\sh{S}'_n)
\]
is a quasi-coherent $\sh{O}_Y$-algebra of finite type, graded in positive degrees, and there is a $Y$-isomorphism
\[
  r:X\simto\Proj(\sh{S})
\]
such that $r^*(\shProj(\sh{S}))$ is isomorphic to $\sh{S}'$ as a graded $\sh{O}_X$-algebra.
\end{proposition}

\begin{proof}
(i) Almost all the assertions have already been proved, the last being a special case of \sref{III.2.3.1}.
The periodicity of $\sh{S}_X$ was proved in \sref[II]{II.8.14.14}, and the assertion that there is a period $d>0$ for which $\sh{O}_X(d)$ is invertible and $Y$-ample is precisely \sref[II]{II.4.6.18}.
Finally, for $0\leq k<d$, $(\sh{S}_X)^{(d,k)}$ is an $(\sh{S}_X)^{(d)}$-module of finite type \sref[II]{II.8.14.14}.
It follows that each $(\sh{S}_X)_n$ is a quasi-coherent $\sh{O}_X$-module of finite type \sref[II]{II.2.1.6}[ii]; since $\sh{O}_X$ is coherent, so is each $(\sh{S}_X)_n$.

(ii) After replacing the period $d$ by one of its multiples, we may suppose that
\[
  \sh{L}=\sh{S}'_d
\]
is very ample relative to $q$ \sref[II]{II.4.6.11}.
By hypothesis,
\[
  \sh{S}'^{(d)}
  =\bigoplus_{n\in\bb{Z}}\sh{L}^{\otimes n},
\]
and hence
\[
  \sh{S}^{(d)}
  =\bigoplus_{n\geq0}q_*(\sh{L}^{\otimes n}).
\]
By \sref[II]{II.3.1.8} and \sref[II]{II.3.2.9}, there is a $Y$-isomorphism
\[
  s:X'=\Proj(\sh{S})\simto
  X''=\Proj(\sh{S}^{(d)})
\]
such that
\oldpage[III]{105}
\[
  s^*(\sh{O}_{X''}(n))=\sh{O}_{X'}(nd)
\]
for every $n$ \sref[II]{II.3.5.2}.
We shall therefore establish the existence of a $Y$-isomorphism $X\simto X'$ by proving the following assertion.

\begin{env}[2.3.4.1]
\label{III.2.3.4.1}
Let $Y$ be a Noetherian prescheme, let $q:X\to Y$ be a projective morphism, and let $\sh{L}$ be an invertible $\sh{O}_X$-module very ample for $q$.
Then
\[
  \sh{S}=\bigoplus_{n\geq0}q_*(\sh{L}^{\otimes n})
\]
is a quasi-coherent graded $\sh{O}_Y$-algebra of finite type such that
\[
  \sh{S}_n=\sh{S}_1^n
\]
for all sufficiently large $n$, and there is a $Y$-isomorphism
\[
  r:X\simto P=\Proj(\sh{S})
\]
such that
\[
  \sh{L}=r^*(\sh{O}_P(1)).
\]
\end{env}

Since $q$ is projective, it follows from \sref[II]{II.5.4.4} and \sref[II]{II.4.4.7} that there is a $Y$-isomorphism
\[
  r':X\simto P'=\Proj(\sh{T}),
\]
where $\sh{T}$ is a quasi-coherent $\sh{O}_Y$-algebra, $\sh{T}_1$ is of finite type and generates $\sh{T}$, and
\[
  \sh{L}=r'^*(\sh{O}_{P'}(1)).
\]
If $q':P'\to Y$ is the structural morphism, then
\[
  \sh{S}=\bigoplus_{n\geq0}q'_*(\sh{O}_{P'}(n)).
\]
By \sref{III.2.3.1}, the canonical homomorphism
\[
  \alpha_n:\sh{T}_n\longrightarrow
  \sh{S}_n=q'_*(\sh{O}_{P'}(n))
\]
is bijective for all sufficiently large $n$.
Since $\sh{T}_n=\sh{T}_1^n$, it follows \emph{a fortiori} that $\sh{S}_n=\sh{S}_1^n$ for all sufficiently large $n$.

The canonical homomorphism $\alpha:\sh{T}\to\sh{S}$ of graded $\sh{O}_Y$-algebras is a \textbf{(TN)}-isomorphism.
Consequently
\phantomsection\label{III.2.3.5.1}
\[
  \Phi=\Proj(\alpha):\Proj(\sh{S})\longrightarrow\Proj(\sh{T})
\]
is an isomorphism \sref[II]{II.3.6.1}.
Moreover, \sref[II]{II.3.5.2} and the fact that the graded $\sh{T}$-modules obtained from $\sh{S}(n)$ by restriction of scalars are \textbf{(TN)}-isomorphic to $\sh{T}(n)$ give
\[
  \Phi_*(\sh{O}_P(n))=\sh{O}_{P'}(n)
\]
for every $n$ \sref[II]{II.3.4.2}.
It remains only to show that $\sh{S}$ is of finite type over $\sh{O}_Y$.
The modules
\[
  \sh{S}_n=q'_*(\sh{O}_{P'}(n))
\]
are coherent by \sref{III.2.2.1}; and if $\sh{S}_n=\sh{S}_1^n$ for $n\geq n_0$, then $\sh{S}$ is generated by the coherent module
\[
  \bigoplus_{i\leq n_0}\sh{S}_i.
\]
This proves the assertion \sref[I]{I.9.6.2}.

We return to the proof of \sref{III.2.3.4}, retaining its notation.
We have proved that there is a $Y$-isomorphism
\[
  r'':X\simto X''
\]
such that
\[
  r''_*(\sh{L}^{\otimes n})=\sh{O}_{X''}(n)
\]
for every $n\in\bb{Z}$.
Let $q'':X''\to Y$ be the structural morphism.
The algebra $\sh{S}'$ is the direct sum of the graded $\sh{S}'^{(d)}$-modules $\sh{S}'^{(d,k)}$.
Each of these is a quasi-coherent $\sh{S}'^{(d)}$-module of finite type, by the periodicity of $\sh{S}'$ and the hypothesis that the $\sh{S}'_n$ are $\sh{O}_X$-modules of finite type \sref[II]{II.8.14.12}.
Set
\[
  \sh{F}^{(k)}=r''_*(\sh{S}'^{(d,k)}),
\]
so that the $\sh{F}^{(k)}$ are quasi-coherent graded $\sh{S}_{X''}$-modules of finite type.
By \sref[II]{II.8.14.8}, the canonical homomorphism
\[
  \beta:\shProj\bigl(\bbGamma_*(\sh{F}^{(k)})\bigr)
  \longrightarrow\sh{F}^{(k)}
\]
is an isomorphism of $\sh{S}_{X''}$-modules.
On the other hand,
\[
  q''_*\bigl((\sh{F}^{(k)})_n\bigr)
  =q_*\bigl((\sh{S}'^{(d,k)})_n\bigr),
\]
and for $n\geq0$ the last $\sh{O}_Y$-module is, by definition, equal to $(\sh{S}^{(d,k)})_n$.
In other words, the canonical injection
\[
  \sh{S}^{(d,k)}
  \longrightarrow\bbGamma_*(\sh{F}^{(k)})
\]
is a \textbf{(TN)}-isomorphism.
Hence
\[
  \shProj(\sh{S}^{(d,k)})
  =\shProj\bigl(\bbGamma_*(\sh{F}^{(k)})\bigr),
\]
and consequently
\[
  r''{}^*\bigl(\shProj(\sh{S}^{(d,k)})\bigr)
  \simeq\sh{S}'^{(d,k)}.
\]
It remains to observe that, up to canonical isomorphism,
\[
  \shProj(\sh{S}^{(d,k)})
  =s_*\bigl((\shProj(\sh{S}))^{(d,k)}\bigr)
\]
\sref[II]{II.8.14.13.1}; this proves that $r^*(\shProj(\sh{S}))$ is isomorphic to $\sh{S}'$.

Finally, by \sref{III.2.3.2}, each $\bbGamma_*(\sh{F}^{(k)})$ satisfies condition \textbf{(TF)}, and hence so does each $\sh{S}^{(d,k)}$.
Moreover, the $\sh{S}_n=q_*(\sh{S}'_n)$ are coherent by \sref{III.2.2.1}, because the $\sh{S}'_n$ are coherent.
It follows at once that the $\sh{S}^{(d,k)}$ are $\sh{S}^{(d)}$-modules of finite type.
Since \sref{III.2.3.4.1} shows that $\sh{S}^{(d)}$ is an $\sh{O}_Y$-algebra of finite type, $\sh{S}$ is likewise an $\sh{O}_Y$-algebra of finite type.
\end{proof}

\oldpage[III]{106}
\begin{proposition}[2.3.5]
\label{III.2.3.5}
Let $Y$ be an integral Noetherian prescheme, let $X$ be an integral prescheme, and let
\[
  f:X\longrightarrow Y
\]
be a projective birational morphism.
There is then a coherent fractional ideal
\[
  \sh{J}\subset\sh{R}(Y)
\]
\sref[II]{II.8.1.2} such that $X$ is $Y$-isomorphic to the prescheme obtained by blowing up $\sh{J}$ \sref[II]{II.8.1.3}.
Moreover, there is an open subset $U$ of $Y$ such that the restriction of $f$ to $f^{-1}(U)$ is an isomorphism from $f^{-1}(U)$ onto $U$ (cf.\ \textbf{I}, 6.5.5), and such that $\sh{J}|U$ is invertible.
\end{proposition}

\begin{proof}
There is an invertible $\sh{O}_X$-module $\sh{L}$ very ample for $f$ \sref[II]{II.4.4.2} and \sref[II]{II.5.3.2}.
Applying \sref{III.2.3.4.1}, we identify $X$ with $\Proj(\sh{S})$, where
\[
  \sh{S}=\bigoplus_{n\geq0}f_*(\sh{L}^{\otimes n}).
\]
The modules $f_*(\sh{L}^{\otimes n})$ are torsion-free over $\sh{O}_Y$ \sref[I]{I.7.4.5}.
Thus $\sh{S}$ identifies canonically with a sub-$\sh{O}_Y$-module of
\[
  \sh{S}\otimes_{\sh{O}_Y}\sh{R}(Y)
\]
\sref[I]{I.7.4.1}.
The latter is a simple sheaf \sref[I]{I.7.3.6}, and is therefore determined by its restriction to a nonempty open subset.

Choose a nonempty open subset $U'\subset U$ on which $\sh{L}|f^{-1}(U')$ is isomorphic to $\sh{O}_X|f^{-1}(U')$, and hence to $\sh{O}_Y|U'$.
The restrictions $f_*(\sh{L}^{\otimes n})|U'$ are then isomorphic to $\sh{O}_Y|U'$.
It follows that
\[
  \sh{S}\otimes_{\sh{O}_Y}\sh{R}(Y)
  \simeq\sh{R}(Y)[T],
\]
where $T$ is an indeterminate.
By \sref{III.2.3.4.1}, $\sh{S}$ is \textbf{(TN)}-isomorphic to the sub-$\sh{O}_Y$-algebra generated by the canonical image of $f_*(\sh{L})$ in this polynomial algebra.
Under the displayed identification, that image is $\sh{J}\cdot T$, where $\sh{J}$ is a coherent sub-$\sh{O}_Y$-module of $\sh{R}(Y)$ \sref{III.2.2.1}.
Its restriction to $U'$ is isomorphic to $\sh{O}_Y|U'$, and, after shrinking $U$ if necessary, $\sh{J}|U$ is invertible.
Consequently
\[
  \sh{S}\quad\text{is \textbf{(TN)}-isomorphic to}\quad
  \bigoplus_{n\geq0}\sh{J}^n,
\]
which completes the proof.
\end{proof}

\begin{corollary}[2.3.6]
\label{III.2.3.6}
Under the hypotheses of \sref{III.2.3.5}, suppose in addition that for every nonzero coherent sub-$\sh{O}_Y$-module
\[
  \sh{J}\subset\sh{R}(Y)
\]
there is an invertible $\sh{O}_Y$-module $\sh{L}$ such that
\[
  \Gamma\bigl(Y,\sh{L}\otimes_{\sh{O}_Y}
  \shHom(\sh{J},\sh{O}_Y)\bigr)\neq0.
\]
Then, in the statement of \sref{III.2.3.5}, $\sh{J}$ may be taken to be an ideal of $\sh{O}_Y$.
This additional condition is always satisfied if there exists an ample $\sh{O}_Y$-module.
\end{corollary}

\begin{proof}
Indeed, by \sref[0]{0.5.4.2},
\[
  \sh{L}\otimes\shHom(\sh{J},\sh{O}_Y)
  =\shHom\bigl(\sh{L}^{-1},\shHom(\sh{J},\sh{O}_Y)\bigr)
  =\shHom(\sh{J}\otimes\sh{L}^{-1},\sh{O}_Y).
\]
The hypothesis therefore says that there is a nonzero homomorphism
\[
  u:\sh{J}\otimes\sh{L}^{-1}\longrightarrow\sh{O}_Y.
\]
For every $y\in Y$, the module $(\sh{J}\otimes\sh{L}^{-1})_y$ identifies with a sub-$\sh{O}_y$-module of the field of fractions $\sh{R}(Y)_y$ of $\sh{O}_y$ \sref[I]{I.7.1.5}.
Thus $u_y$ is necessarily injective, and $u$ identifies $\sh{J}\otimes\sh{L}^{-1}$ with an ideal $\sh{J}'$ of $\sh{O}_Y$.
But
\[
  \Proj\left(\bigoplus_{n\geq0}\sh{J}^n\right)
  \quad\text{and}\quad
  \Proj\left(\bigoplus_{n\geq0}
  (\sh{J}\otimes\sh{L}^{-1})^n\right)
\]
are $Y$-isomorphic \sref[II]{II.3.1.8}; this proves the first assertion.

For the second, note that
\[
  \sh{F}=\shHom_{\sh{O}_Y}(\sh{J},\sh{O}_Y)
\]
is coherent and nonzero, since there is an open subset $U$ of $Y$ on which $\sh{J}|U$ is invertible.
If $\sh{L}$ is an ample $\sh{O}_Y$-module, there is an integer $n$ such that
\[
  \sh{F}(n)=\sh{F}\otimes\sh{L}^{\otimes n}
\]
is generated by its sections over $Y$ \sref[II]{II.4.5.5}.
\emph{A fortiori},
\[
  \Gamma(Y,\sh{F}(n))\neq0,
\]
which proves the assertion.
\end{proof}

\begin{corollary}[2.3.7]
\label{III.2.3.7}
Let $X$ and $Y$ be integral schemes projective over a field $k$, and let
\[
  f:X\longrightarrow Y
\]
be a birational $k$-morphism.
Then $X$ is $k$-isomorphic to a $Y$-scheme obtained by blowing up a closed subscheme $Y'$ of $Y$, not necessarily reduced.
\end{corollary}

\begin{proof}
\oldpage[III]{107}
Indeed, $f$ is projective \sref[II]{II.5.5.5}[v], and since $Y$ is projective over $k$, the additional condition of \sref{III.2.3.6} is satisfied.
It therefore suffices to take the closed subscheme $Y'$ of $Y$ defined by the coherent ideal $\sh{J}$ of \sref{III.2.3.6}.
\end{proof}

\begin{env}[2.3.8]
\label{III.2.3.8}
\emph{Remark.}
In Chapter~IV, in studying the notion of divisor, we shall see that if, in the statement of \sref{III.2.3.5}, the local rings $\sh{O}_y$, for $y\in Y$, are assumed factorial---as is the case, for example, when $Y$ is nonsingular---then $X$ may be obtained from $Y$ by blowing up a closed subprescheme $Y'$ of $Y$ whose underlying space is contained in $Y-U$.
\end{env}


\subsection{A generalisation of the fundamental theorem}
\label{subsection:III.2.4}

\begin{theorem}[2.4.1]
\label{III.2.4.1}
Let $Y$ be a Noetherian prescheme, let $\sh{S}$ be a quasi-coherent $\sh{O}_Y$-algebra of finite type, let
\[
  f:X\longrightarrow Y
\]
be a projective morphism, set $\sh{S}'=f^*(\sh{S})$, and let $\sh{M}$ be a quasi-coherent $\sh{S}'$-module of finite type.
Then:
\begin{enumerate}
  \item[(i)] For every $p\in\bb Z$, $\RR^p f_*(\sh{M})$ is an $\sh{S}$-module of finite type.

  \item[(ii)] Suppose in addition that $\sh{L}$ is an invertible $\sh{O}_X$-module ample for $f$, and, for every $n\in\bb Z$, set
  \[
    \sh{M}(n)=\sh{M}\otimes\sh{L}^{\otimes n}.
  \]
  There is an integer $N$ such that, for $n\geq N$,
  \[
  \label{III.2.4.1.1}
    \RR^p f_*(\sh{M}(n))=0
    \tag{2.4.1.1}
  \]
  for every $p>0$, and the canonical homomorphism
  \[
    f^*\bigl(f_*(\sh{M}(n))\bigr)\longrightarrow\sh{M}(n)
  \]
  \sref[0]{0.4.4.3} is surjective.
\end{enumerate}
\end{theorem}

\begin{proof}
Set
\[
  Y'=\Spec(\sh{S}),\qquad X'=\Spec(\sh{S}'),
\]
so that $X'=X\times_Y Y'$ \sref[II]{II.1.5.5}.
Let
\[
  g:Y'\longrightarrow Y,\qquad g':X'\longrightarrow X
\]
be the structural morphisms, which are affine by definition, and let
\[
  f'=f_{(Y')}:X'\longrightarrow Y'.
\]
We thus have a commutative diagram
\[
  \xymatrix{
    X\ar[d]_f &
    X'\ar[l]_{g'}\ar[d]^{f'}\ar[dl]^{h}\\
    Y &
    Y'\ar[l]^{g},
  }
\]
where $f'$ is projective \sref[II]{II.5.5.5}[iii] and
\[
  h=f\circ g'=g\circ f'.
\]

\emph{(i)} Let $\widetilde{\sh{M}}$ be the $\sh{O}_{X'}$-module associated with the quasi-coherent $\sh{S}'$-module $\sh{M}$ when $X'$ is regarded as an affine $X$-scheme \sref[II]{II.1.4.3}; thus
\[
  \sh{M}=g'_*(\widetilde{\sh{M}}).
\]
Since $\sh{M}$ is an $\sh{S}'$-module of finite type, $\widetilde{\sh{M}}$ is an $\sh{O}_{X'}$-module of finite type \sref[II]{II.1.4.5}.
Moreover, $h$ is of finite type, since $g$ and $f'$ are of finite type \sref[II]{II.1.3.7} and \sref[I]{I.6.3.4}[ii].
Hence $X'$ is Noetherian \sref[I]{I.6.3.7}, and $\widetilde{\sh{M}}$ is therefore coherent.

Since $g'$ is affine, the canonical homomorphism
\[
  \RR^p f_*(\sh{M})
  \longrightarrow
  \RR^p h_*(\widetilde{\sh{M}})
\]
is bijective \sref{III.1.3.4}.
It is moreover a homomorphism of $\sh{S}$-modules.
Indeed, the canonical homomorphism
\[
\label{III.2.4.1.2}
  g'_*(\sh{O}_{X'})
  \otimes_{\sh{O}_X}
  g'_*(\widetilde{\sh{M}})
  \longrightarrow
  g'_*(\widetilde{\sh{M}})
  \tag{2.4.1.2}
\]
defines the $\sh{S}'$-module structure on $\sh{M}$, since
\[
  \sh{S}'=g'_*(\sh{O}_{X'}).
\]
It canonically gives rise to a homomorphism
\[
  f_*\bigl(g'_*(\sh{O}_{X'})\bigr)
  \otimes
  \RR^p f_*\bigl(g'_*(\widetilde{\sh{M}})\bigr)
  \longrightarrow
  \RR^p f_*\bigl(g'_*(\widetilde{\sh{M}})\bigr).
\]
\oldpage[III]{108}
By \sref[0]{0.12.2.2}, and because \eqref{III.2.4.1.2} itself is obtained, by applying \sref[0]{0.4.2.2.1}, from the homomorphism
\[
  \sh{O}_{X'}\otimes\widetilde{\sh{M}}
  \longrightarrow\widetilde{\sh{M}}
\]
that defines the $\sh{O}_{X'}$-module structure on $\widetilde{\sh{M}}$, the diagram
\[
  \xymatrix{
    f_*\bigl(g'_*(\sh{O}_{X'})\bigr)
    \otimes\RR^p f_*\bigl(g'_*(\widetilde{\sh{M}})\bigr)
    \ar[r]\ar[d] &
    \RR^p f_*\bigl(g'_*(\widetilde{\sh{M}})\bigr)\ar[d]\\
    h_*(\sh{O}_{X'})
    \otimes\RR^p h_*(\widetilde{\sh{M}})
    \ar[r] &
    \RR^p h_*(\widetilde{\sh{M}})
  }
\]
is commutative \sref[0]{0.12.2.6}.
Composing the horizontal arrows with the homomorphism induced by the canonical homomorphism
\[
  \sh{S}
  \longrightarrow f_*(f^*(\sh{S}))
  =f_*(\sh{S}')
  =f_*\bigl(g'_*(\sh{O}_{X'})\bigr)
  =h_*(\sh{O}_{X'}),
\]
we obtain the assertion.

On the other hand, since $g$ is affine and $f'$ is separated and quasi-compact, the canonical homomorphism
\[
  \RR^p h_*(\widetilde{\sh{M}})
  \longrightarrow
  g_*\bigl(\RR^p f'_*(\widetilde{\sh{M}})\bigr)
\]
is bijective \sref{III.1.4.14}.
As above, it is an isomorphism of $\sh{S}$-modules, this time by using the commutativity in \sref[0]{0.12.2.6.2}.
Now $f'$ is projective and $\widetilde{\sh{M}}$ is coherent, so $\RR^p f'_*(\widetilde{\sh{M}})$ is a coherent $\sh{O}_{Y'}$-module by \sref{III.2.2.1}.
It follows that
\[
  g_*\bigl(\RR^p f'_*(\widetilde{\sh{M}})\bigr)
\]
is an $\sh{S}$-module of finite type \sref[II]{II.1.4.5}.

\emph{(ii)} Put
\[
  \sh{L}'=g'^*(\sh{L}),
\]
an invertible $\sh{O}_{X'}$-module.
For every $n\in\bb Z$, up to canonical isomorphism,
\[
  g'_*\bigl(\widetilde{\sh{M}}\otimes\sh{L}'^{\otimes n}\bigr)
  =
  g'_*(\widetilde{\sh{M}})\otimes\sh{L}^{\otimes n}
  =
  \sh{M}(n)
\]
\sref[0]{0.5.4.10}.
Applying the argument of \emph{(i)} to
\[
  \widetilde{\sh{M}}\otimes\sh{L}'^{\otimes n}
\]
shows that
\[
  \RR^p f_*
  \bigl(g'_*(\widetilde{\sh{M}}\otimes\sh{L}'^{\otimes n})\bigr)
  \simeq
  g_*\bigl(\RR^p f'_*
  (\widetilde{\sh{M}}\otimes\sh{L}'^{\otimes n})\bigr).
\]
Now $\sh{L}'$ is ample for $f'$ \sref[II]{II.4.6.13}[iii].
It follows from \sref{III.2.2.1} that there is an integer $N$ such that
\[
  \RR^p f'_*
  \bigl(\widetilde{\sh{M}}\otimes\sh{L}'^{\otimes n}\bigr)=0
\]
for every $p>0$ and every $n\geq N$.
This proves \eqref{III.2.4.1.1}.

It follows once more from \sref{III.2.2.1} that $N$ may be chosen so that, for $n\geq N$, the canonical homomorphism
\[
  f'^*\bigl(f'_*(\widetilde{\sh{M}}\otimes\sh{L}'^{\otimes n})\bigr)
  \longrightarrow
  \widetilde{\sh{M}}\otimes\sh{L}'^{\otimes n}
\]
is surjective.
Since $g'_*$ is exact \sref[II]{II.1.4.4}, the corresponding homomorphism
\[
  g'_*\Bigl(
    f'^*\bigl(f'_*(\widetilde{\sh{M}}\otimes\sh{L}'^{\otimes n})\bigr)
  \Bigr)
  \longrightarrow
  g'_*(\widetilde{\sh{M}}\otimes\sh{L}'^{\otimes n})
  =\sh{M}(n)
\]
is surjective.
But \sref[II]{II.1.5.2} gives
\[
  g'_*\Bigl(
    f'^*\bigl(f'_*(\widetilde{\sh{M}}\otimes\sh{L}'^{\otimes n})\bigr)
  \Bigr)
  =
  f^*\Bigl(
    g_*\bigl(f'_*(\widetilde{\sh{M}}\otimes\sh{L}'^{\otimes n})\bigr)
  \Bigr),
\]
and $g_*\circ f'_*=f_*\circ g'_*$.
Consequently,
\[
  g'_*\Bigl(
    f'^*\bigl(f'_*(\widetilde{\sh{M}}\otimes\sh{L}'^{\otimes n})\bigr)
  \Bigr)
  =
  f^*\bigl(f_*(\sh{M}(n))\bigr),
\]
which completes the proof.
\end{proof}

\begin{env}[2.4.2]
\label{III.2.4.2}
We shall apply \sref{III.2.4.1}, in particular, when $\sh{S}$ is an $\sh{O}_Y$-algebra graded in positive degrees and
\[
  \sh{M}=\bigoplus_{k\in\bb Z}\sh{M}_k
\]
is a graded $\sh{S}'$-module.
\oldpage[III]{109}
Under the same finiteness hypotheses on $\sh{S}$ and $\sh{M}$, it follows from \sref{III.2.4.1} that
\[
  \bigoplus_{k\in\bb Z}\RR^p f_*(\sh{M}_k)
\]
is an $\sh{S}$-module of finite type for every $p$.
Under the additional hypotheses of \sref{III.2.4.1}[ii], there is moreover an integer $N$ such that, for $n\geq N$,
\[
  \RR^p f_*(\sh{M}_k(n))=0
\]
for every $p>0$ and every $k\in\bb Z$, and the canonical homomorphism
\[
  f^*\bigl(f_*(\sh{M}_k(n))\bigr)
  \longrightarrow
  \sh{M}_k(n)
\]
is surjective for every $k\in\bb Z$.
\end{env}


\subsection{Euler--Poincar\'e characteristic and the Hilbert polynomial}
\label{subsection:III.2.5}

\begin{env}[2.5.1]
\label{III.2.5.1}
Let $A$ be an Artinian ring and let $X$ be an $A$-scheme projective over
\[
  Y=\Spec(A).
\]
For every coherent $\sh{O}_X$-module $\sh{F}$, the $A$-modules
\[
  \HH^i(X,\sh{F})\qquad(i\geq0)
\]
are of finite type \sref{III.2.2.1}, and hence have finite length, since $A$ is Artinian.
Moreover, \sref{III.2.2.1} shows that $\HH^i(X,\sh{F})=0$ for all but finitely many $i\geq0$.
Thus the integer
\[
\label{III.2.5.1.1}
  \chi_A(\sh{F})
  =
  \sum_{i=0}^{\infty}
  (-1)^i\operatorname{length}_A\bigl(\HH^i(X,\sh{F})\bigr)
  \tag{2.5.1.1}
\]
is defined for every coherent $\sh{O}_X$-module $\sh{F}$.
When $A$ is an Artinian local ring, $\chi_A(\sh{F})$ is called the \emph{Euler--Poincar\'e characteristic} of $\sh{F}$ with respect to $A$.
For $\sh{F}=\sh{O}_X$, the integer $\chi_A(\sh{O}_X)$ is called the \emph{arithmetic genus} of $X$ with respect to $A$.
\end{env}

\begin{proposition}[2.5.2]
\label{III.2.5.2}
Let
\[
  0\longrightarrow\sh{F}'
  \longrightarrow\sh{F}
  \longrightarrow\sh{F}''
  \longrightarrow0
\]
be an exact sequence of coherent $\sh{O}_X$-modules.
Then
\[
\label{III.2.5.2.1}
  \chi_A(\sh{F})
  =
  \chi_A(\sh{F}')+\chi_A(\sh{F}'').
  \tag{2.5.2.1}
\]
\end{proposition}

\begin{proof}
The cohomology modules of $\sh{F}$, $\sh{F}'$, and $\sh{F}''$ vanish in all sufficiently large degrees.
There is therefore an integer $r>0$ for which the long exact cohomology sequence has the form
\[
  0\longrightarrow\HH^0(X,\sh{F}')
  \longrightarrow\HH^0(X,\sh{F})
  \longrightarrow\HH^0(X,\sh{F}'')
  \longrightarrow\HH^1(X,\sh{F}')
  \longrightarrow\cdots
\]
\[
  \cdots\longrightarrow\HH^r(X,\sh{F}')
  \longrightarrow\HH^r(X,\sh{F})
  \longrightarrow\HH^r(X,\sh{F}'')
  \longrightarrow0.
\]
In an exact sequence of $A$-modules of finite length with zero terms at both ends, the alternating sum of the lengths is zero \sref[0]{0.11.10.1}.
Applying this fact gives \eqref{III.2.5.2.1} immediately.
\end{proof}

The result of \sref{III.2.5.2} applies whenever $X$ is a quasi-compact $A$-scheme and the $A$-modules $\HH^i(X,\sh{F})$ are of finite type for every coherent $\sh{O}_X$-module $\sh{F}$ \sref{III.1.4.12}.

\begin{theorem}[2.5.3]
\label{III.2.5.3}
Let $A$ be an Artinian local ring, let $X$ be a scheme projective over
\[
  Y=\Spec(A),
\]
let $\sh{L}$ be an invertible $\sh{O}_X$-module very ample relative to $Y$, and let $\sh{F}$ be a coherent $\sh{O}_X$-module.
For every $n\in\bb Z$, put
\[
  \sh{F}(n)=\sh{F}\otimes_{\sh{O}_X}\sh{L}^{\otimes n}.
\]
\begin{enumerate}
  \item[(i)] There is a unique polynomial $P\in\bb Q[T]$ such that
  \[
    \chi_A(\sh{F}(n))=P(n)
  \]
  for every $n\in\bb Z$.
  The polynomial $P$ is called the \emph{Hilbert polynomial} of $\sh{F}$ with respect to $A$.

  \item[(ii)] For all sufficiently large $n$,
  \[
    \chi_A(\sh{F}(n))
    =
    \operatorname{length}_A\Gamma(X,\sh{F}(n)).
  \]

  \item[(iii)] The coefficient of the highest-degree term of $\chi_A(\sh{F}(n))$ is nonnegative.
\end{enumerate}
\end{theorem}

We shall moreover prove in Chapter~IV, in the section devoted to dimension, that \emph{the degree of $\chi_A(\sh{F}(n))$ equals the dimension of the support of $\sh{F}$}.

\begin{proof}
For all sufficiently large $n$, we have
\[
  \HH^i(X,\sh{F}(n))=0
\]
for every $i>0$ \sref{III.2.2.1}.
\oldpage[III]{110}
Consequently,
\[
  \chi_A(\sh{F}(n))
  =
  \operatorname{length}_A\HH^0(X,\sh{F}(n))
  =
  \operatorname{length}_A\Gamma(X,\sh{F}(n))
\]
for all sufficiently large $n$, which proves \emph{(ii)}.
It follows that $\chi_A(\sh{F}(n))\geq0$ for all sufficiently large $n$, and \emph{(iii)} therefore follows from \emph{(i)}.
The uniqueness assertion in \emph{(i)} is immediate, so it remains to prove the existence of $P$.

We first show that we may assume
\[
  \mathfrak{m}\sh{F}=0,
\]
where $\mathfrak{m}$ is the maximal ideal of $A$.
Indeed, there is an integer $s>0$ such that $\mathfrak{m}^s=0$, and $\sh{F}(n)$ therefore has the finite filtration
\[
  \sh{F}(n)
  \supset\mathfrak{m}\sh{F}(n)
  \supset\cdots
  \supset\mathfrak{m}^{s-1}\sh{F}(n)
  \supset0.
\]
Induction using \eqref{III.2.5.2.1} gives
\[
  \chi_A(\sh{F}(n))
  =
  \sum_{k=1}^{s}
  \chi_A\left(
    \mathfrak{m}^{k-1}\sh{F}(n)/
    \mathfrak{m}^{k}\sh{F}(n)
  \right).
\]
But
\[
  \mathfrak{m}^{k-1}\sh{F}(n)/
  \mathfrak{m}^{k}\sh{F}(n)
  =
  \sh{F}'_k(n),
  \qquad
  \sh{F}'_k=
  \mathfrak{m}^{k-1}\sh{F}/\mathfrak{m}^{k}\sh{F},
\]
which proves the reduction.

Suppose, then, that $\mathfrak{m}\sh{F}=0$.
Let $X'$ be the closed subscheme of $X$ obtained as the inverse image, under the structural morphism
\[
  X\longrightarrow\Spec(A),
\]
of the unique closed point of $\Spec(A)$, and let $j:X'\to X$ be the canonical injection.
Then
\[
  \sh{F}=j_*(\sh{F}')
\]
for a coherent $\sh{O}_{X'}$-module $\sh{F}'$, and $X'$ is projective over $\Spec(K)$, where
\[
  K=A/\mathfrak{m}.
\]
If $\sh{L}'=j^*(\sh{L})$, then $\sh{L}'$ is very ample relative to $\Spec(K)$ \sref[II]{II.4.4.10}, and
\[
  \sh{F}(n)=j_*(\sh{F}'(n)),
  \qquad
  \sh{F}'(n)
  =
  \sh{F}'\otimes_{\sh{O}_{X'}}\sh{L}'^{\otimes n}
\]
\sref[0]{0.5.4.10}.
It follows that
\[
  \chi_A(\sh{F}(n))
  =
  \chi_K(\sh{F}'(n))
\]
(G, II, 4.9.1), and we are reduced to the case in which $A$ is a field.

Now $X$ may be regarded as a closed subscheme of
\[
  P=\mathbf{P}_A^r
\]
for a suitable $r$ \sref[II]{II.5.5.4}[ii].
If $i:X\to P$ is the canonical injection, the same argument gives
\[
  \chi_A(\sh{F}(n))
  =
  \chi_A\bigl(i_*(\sh{F})(n)\bigr).
\]
We may therefore restrict to the case
\[
  X=\mathbf{P}_A^r=\Proj(S),
  \qquad
  S=A[T_0,\ldots,T_r],
\]
where $A$ is a field.

In this case $\sh{F}=\widetilde{M}$ for a graded $S$-module $M$ of finite type \sref[II]{II.2.7.8}.
By the Hilbert syzygy theorem (M, VIII, 6.5), $M$ has a finite resolution by graded free $S$-modules of finite type:
\[
  0\longrightarrow L_q
  \longrightarrow L_{q-1}
  \longrightarrow\cdots
  \longrightarrow L_1
  \longrightarrow M
  \longrightarrow0.
\]
Since the functor $M\mapsto\widetilde{M}$ is exact \sref[II]{II.2.5.4}, we obtain an exact sequence
\[
  0\longrightarrow\widetilde{L}_q
  \longrightarrow\widetilde{L}_{q-1}
  \longrightarrow\cdots
  \longrightarrow\widetilde{L}_1
  \longrightarrow\widetilde{M}
  \longrightarrow0,
\]
and hence, for every $n\in\bb Z$, an exact sequence
\[
  0\longrightarrow\widetilde{L}_q(n)
  \longrightarrow\widetilde{L}_{q-1}(n)
  \longrightarrow\cdots
  \longrightarrow\widetilde{L}_1(n)
  \longrightarrow\widetilde{M}(n)
  \longrightarrow0.
\]
Induction on $q$, using \sref{III.2.5.2}, yields
\[
  \chi_A(\widetilde{M}(n))
  =
  \sum_{j=1}^{q}
  (-1)^{j+1}\chi_A(\widetilde{L}_j(n)).
\]
To prove \emph{(i)}, we are therefore reduced to the case in which $M$ is graded, free, and of finite type, and hence to the case
\[
  M=S(h)
\]
for some $h\in\bb Z$.
Then
\[
  \widetilde{M}(n)
  =
  \widetilde{M(n)}
  =
  \widetilde{S(n+h)}
\]
\sref[II]{II.2.5.15}, so the theorem follows from the lemma below.
\end{proof}

\oldpage[III]{111}
\begin{lemma}[2.5.3.1]
\label{III.2.5.3.1}
Let $A$ be a field, let $r>0$ be an integer, and let
\[
  X=\mathbf{P}_A^r.
\]
Then
\[
  \chi_A(\sh{O}_X(n))
  =
  \binom{n+r}{r}
\]
for every $n\in\bb Z$.
\end{lemma}

\begin{proof}
For $n\geq0$,
\[
  \chi_A(\sh{O}_X(n))
  =
  \operatorname{length}_A\HH^0(X,\sh{O}_X(n)),
\]
which is the number of monomials in the $T_i$ of total degree $n$, namely
\[
  \binom{n+r}{r}
\]
\sref{III.2.1.12}.

For $n\leq-r-1$, we likewise have
\[
  \chi_A(\sh{O}_X(n))
  =
  (-1)^r\operatorname{length}_A\HH^r(X,\sh{O}_X(n)).
\]
Writing $n=-r-h$, the dimension of $\HH^r(X,\sh{O}_X(n))$ over $A$ is the number of sequences
\[
  (p_i)_{0\leq i\leq r}
\]
of positive integers satisfying
\[
  \sum_{i=0}^{r}p_i=r+h
\]
\sref{III.2.1.12}.
Equivalently, it is the number of sequences of nonnegative integers
\[
  (q_i)_{0\leq i\leq r}
\]
such that
\[
  \sum_{i=0}^{r}q_i=h-1.
\]
This number is
\[
  \binom{h+r-1}{r}
  =
  (-1)^r\binom{n+r}{r}.
\]
Finally, if $-r\leq n<0$, then
\[
  \binom{n+r}{r}=0,
\]
and $\HH^i(X,\sh{O}_X(n))=0$ for every $i\geq0$ \sref{III.2.1.12}.
This proves the lemma.
\end{proof}

\begin{corollary}[2.5.4]
\label{III.2.5.4}
Let $A$ be an Artinian local ring, let $S$ be a graded $A$-algebra of finite type generated by $S_1$, let $M$ be a graded $S$-module of finite type, and let
\[
  X=\Proj(S).
\]
Then, for all sufficiently large $n$,
\[
  \chi_A(\widetilde{M}(n))
  =
  \operatorname{length}_A(M_n).
\]
\end{corollary}

\begin{proof}
For all sufficiently large $n$, the modules
\[
  M_n
  \quad\text{and}\quad
  \Gamma(X,\widetilde{M}(n))
\]
are isomorphic \sref{III.2.3.1}.
\end{proof}


\subsection{Application: ampleness criteria}
\label{subsection:III.2.6}

\begin{proposition}[2.6.1]
\label{III.2.6.1}
Let $Y$ be a Noetherian prescheme, let $f:X\to Y$ be a proper morphism,
and let $\sh{L}$ be an invertible $\sh{O}_X$-module.
The following conditions are equivalent:
\begin{enumerate}
  \item[\emph{a)}] $\sh{L}$ is ample for $f$.
  \item[\emph{b)}] For every coherent $\sh{O}_X$-module $\sh{F}$, there
  exists an integer $N$ such that, for $n\geq N$,
  \[
    \RR^qf_*(\sh{F}\otimes\sh{L}^{\otimes n})=0
  \]
  for every $q>0$.
  \item[\emph{c)}] For every coherent ideal $\sh{J}$ of $\sh{O}_X$, there
  exists an integer $N$ such that, for $n\geq N$,
  \[
    \RR^1f_*(\sh{J}\otimes\sh{L}^{\otimes n})=0.
  \]
\end{enumerate}
\end{proposition}

\begin{proof}
We have seen that \emph{a)} implies \emph{b)} \sref{III.2.2.1},~(ii).
It is immediate that \emph{b)} implies \emph{c)}, and it remains to prove
that \emph{c)} implies \emph{a)}.
We may restrict to the case where $Y$ is affine \sref[II]{II.4.6.4}.
It is enough to show that, as $h$ ranges through the sections of
$\sh{L}^{\otimes n}$ over $X$, with $n>0$, the affine open subsets $X_h$
form a covering of $X$ \sref[II]{II.4.5.2}.

For this purpose, let $x$ be a closed point of $X$ and let $U$ be an
affine open neighbourhood of $x$.
We shall find an integer $n$ and a section
\[
  h\in\Gamma(X,\sh{L}^{\otimes n})
\]
such that $x\in X_h\subset U$.
The open subset $X_h$ will then be affine \sref[I]{I.1.3.6}.
The union of all such $X_h$ is an open subset of $X$ containing every
closed point, and therefore is all of $X$, since $X$ is Noetherian
\sref[I]{I.6.3.7} and \sref[0]{0.2.1.3}.

Let $\sh{J}$, respectively $\sh{J}'$, be the quasi-coherent ideal of
$\sh{O}_X$ defining the reduced closed subprescheme whose underlying
space is $X-U$, respectively $(X-U)\cup\{x\}$ \sref[I]{I.5.2.1}.
The ideals $\sh{J}$ and $\sh{J}'$ are coherent \sref[I]{I.6.1.1},
$\sh{J}'\subset\sh{J}$, and
\[
  \sh{J}''=\sh{J}/\sh{J}'
\]
is a coherent $\sh{O}_X$-module \sref[0]{0.5.3.3}, supported by
$\{x\}$, with $\sh{J}''_x=\kres(x)$.
Since $\sh{L}$ is locally free, the sequence
\[
  0\longrightarrow
  \sh{J}'\otimes\sh{L}^{\otimes n}
  \longrightarrow
  \sh{J}\otimes\sh{L}^{\otimes n}
  \longrightarrow
  \sh{J}''\otimes\sh{L}^{\otimes n}
  \longrightarrow0
\]
is exact for every $n$.
By hypothesis, for all sufficiently large $n$,
\[
  \HH^1(X,\sh{J}'\otimes\sh{L}^{\otimes n})=0.
\]

\oldpage[III]{112}

The exact cohomology sequence therefore shows that
\[
  \Gamma(X,\sh{J}\otimes\sh{L}^{\otimes n})
  \longrightarrow
  \Gamma(X,\sh{J}''\otimes\sh{L}^{\otimes n})
\]
is surjective.
A section $g$ of $\sh{J}''\otimes\sh{L}^{\otimes n}$ over $X$ such that
$g(x)\neq0$ is consequently the image of a section
\[
  h\in\Gamma(X,\sh{J}\otimes\sh{L}^{\otimes n})
  \subset\Gamma(X,\sh{L}^{\otimes n});
\]
the inclusion follows from \sref[0]{0.5.4.1}.
By construction, $h(x)\neq0$, while $h(z)=0$ for every $z\notin U$.
Thus $x\in X_h\subset U$, which proves the assertion.
\end{proof}

\begin{proposition}[2.6.2]
\label{III.2.6.2}
Let $Y$ be a Noetherian prescheme, let $f:X\to Y$ be a morphism of
finite type, let $g:X'\to X$ be a finite surjective morphism, let
$\sh{L}$ be an invertible $\sh{O}_X$-module, and set
\[
  \sh{L}'=g^*(\sh{L}).
\]
Suppose that there exists a subset $Z$ of $X$, proper over $Y$
\sref[II]{II.5.4.10}, such that, for every $x\in X-Z$, either $X$ is
normal at $x$, or
\[
  \bigl(g_*(\sh{O}_{X'})\bigr)_x
\]
is a free $\sh{O}_{X,x}$-module.
Then $\sh{L}$ is ample for $f$ if and only if $\sh{L}'$ is ample for
$f\circ g$.
\end{proposition}

\begin{proof}
\begin{env}[2.6.2.1]
\label{III.2.6.2.1}
Since $g$ is affine, the condition is necessary \sref[II]{II.5.1.12}.
To prove that it is sufficient, we may suppose that $Y$ is affine
\sref[II]{II.4.6.4}.
We first show that we may also restrict to the case in which $X$ is
reduced.

Let $j:X_{\mathrm{red}}\to X$ be the canonical immersion, and put
\[
  X_1=X_{\mathrm{red}},
  \qquad
  X'_1=X'\times_X X_1.
\]
We then have the commutative diagram
\[
  \xymatrix{
    X'\ar[d]_g &
    X'_1\ar[l]_{j'}\ar[d]^{g_1}\\
    X &
    X_1\ar[l]^{j}
  }
  \label{III.2.6.2.2}\tag{2.6.2.2}
\]
The morphism $f\circ j$ is of finite type \sref[I]{I.6.3.4}, and
$g_1$ is finite \sref[II]{II.6.1.5},~(iii).
If $\sh{L}'$ is ample for $f\circ g$, then
$j'^*(\sh{L}')$ is ample for $f\circ g\circ j'$, because $j'$ is a
closed immersion \sref[II]{II.5.1.12} and \sref[I]{I.4.3.2}.
If $Z_1=j^{-1}(Z)$, then $Z_1$ is proper over $Y$
\sref[II]{II.5.4.10}.
Moreover, if $X$ is normal at a point $x$, then so is $X_{\mathrm{red}}$
at that point; and if
$\bigl(g_*(\sh{O}_{X'})\bigr)_x$ is a free $\sh{O}_{X,x}$-module,
then \sref[II]{II.1.5.2} shows that
\[
  \bigl((g_1)_*(\sh{O}_{X'_1})\bigr)_x
\]
is a free $\sh{O}_{X_1,x}$-module.
Finally, since $X$ is Noetherian \sref[I]{I.6.3.7}, the ampleness of
$j^*(\sh{L})$ implies that of $\sh{L}$ \sref[II]{II.4.5.14}.
As
\[
  j'^*(\sh{L}')=g_1^*(j^*(\sh{L})),
\]
this proves the announced reduction.
Henceforth we suppose that $Y$ is affine and that $X$ is reduced.
\end{env}

The hypotheses of \sref[II]{II.6.6.11} are now satisfied.
There exist a reduced $Y$-prescheme $X_2$ and a finite birational
$Y$-morphism
\[
  h:X_2\to X
\]
whose restriction to $h^{-1}(X-Z)$ is an isomorphism onto $X-Z$, and
such that $h^*(\sh{L})$ is ample.
Replacing $X'$ by $X_2$, we are reduced to proving the proposition under
the additional assumption that $g$ has the properties just listed for
$h$.
We continue to write $Z$ for a closed subprescheme of $X$, proper over
$Y$ \sref[II]{II.5.4.10}, whose underlying space is the given subset
$Z$.

\begin{env}[2.6.2.3]
\label{III.2.6.2.3}
Let $X_1$ be a closed subprescheme of $X$, let $j:X_1\to X$ be the
canonical immersion, and let
\[
  X'_1=g^{-1}(X_1)=X'\times_X X_1,
\]
with canonical immersion $j':X'_1\to X'$.
The corresponding square is again \eqref{III.2.6.2.2}.
Put
\[
  \sh{L}_1=j^*(\sh{L}),
  \qquad
  \sh{L}'_1=j'^*(\sh{L}')=g_1^*(\sh{L}_1).
\]
Then $\sh{L}'_1$ is ample for $f\circ g\circ j'$
\sref[II]{II.5.1.12}.
If $Z_1=j^{-1}(Z)$, the closed subprescheme $Z_1$ of $X_1$ is proper
over $Y$ \sref[II]{II.5.4.2},~(ii).
In other words, the hypotheses of \sref{III.2.6.2} hold for
$X_1$, $\sh{L}_1$, $g_1$, and $Z_1$.
\end{env}

\oldpage[III]{113}

This permits us to prove \sref{III.2.6.2} by Noetherian induction
\sref[0]{0.2.2.2} when the restriction of $g$ to $g^{-1}(X-Z)$ is an
isomorphism onto $X-Z$; as seen in \sref{III.2.6.2.1}, that case is
sufficient.
It is enough to show that if the conclusion of \sref{III.2.6.2} holds
for $\sh{L}_1$ on every closed subprescheme $X_1$ of $X$ whose
underlying space is not $X$, then it also holds for $\sh{L}$.

\begin{env}[2.6.2.4]
\label{III.2.6.2.4}
Set
\[
  \sh{A}=\sh{O}_X,
  \qquad
  \sh{B}=g_*(\sh{O}_{X'}).
\]
Since $g$ is birational and $X$ is reduced, $\sh{B}$ is a
sub-$\sh{A}$-algebra of the sheaf $\sh{R}(X)$ of total quotient rings,
and it is a coherent $\sh{A}$-module.
Moreover,
\[
  \sh{B}|_{X-Z}=\sh{A}|_{X-Z}.
\]
Let $\sh{K}$ be the conductor of $\sh{B}$ over $\sh{A}$, that is, the
largest quasi-coherent sub-$\sh{A}$-module of $\sh{A}$ such that
\[
  \sh{B}\sh{K}\subset\sh{A}.
\]
Equivalently, $\sh{K}$ is the annihilator of the
$\sh{A}$-module $\sh{B}/\sh{A}$ \sref[0]{0.5.3.7}.
It follows that
\[
  \sh{B}\sh{K}=\sh{K}.
\]
Clearly $\sh{K}_x=\sh{A}_x$ at every point admitting a neighbourhood
$W_x$ on which $g:g^{-1}(W_x)\to W_x$ is an isomorphism.
This holds in particular at every point of $X-Z$, and in a
neighbourhood of every generic point of an irreducible component of
$X$.

Let
\[
  Z_1=\Spec(\sh{A}/\sh{K})
\]
be the closed subprescheme of $X$ defined by $\sh{K}$.
It is proper over $Y$, because its underlying space is closed in $Z$
\sref[II]{II.5.4.10}.
The definition of $\sh{K}$ also gives
\[
  \sh{B}|_{X-Z_1}=\sh{A}|_{X-Z_1}.
\]
We may therefore reduce to the case
\[
  Z=\Spec(\sh{A}/\sh{K}).
\]
Since $X-Z_1$ is a nonempty open subset of $X$, we may in addition
suppose that the underlying space of $Z$ is distinct from $X$.
\end{env}

\begin{env}[2.6.2.5]
\label{III.2.6.2.5}
Regard $X'$ as $\Spec(\sh{B})$, since $g$ is affine, and let
\[
  \sh{K}'=\widetilde{\sh{K}}.
\]
This is a coherent ideal of $\sh{O}_{X'}$ satisfying
$g_*(\sh{K}')=\sh{K}$ \sref[II]{II.1.4.1}.
The closed subprescheme
\[
  Z'=g^{-1}(Z)=Z\times_X X'
\]
of $X'$ is defined by $\sh{K}'$ and is equal to
\[
  \Spec(\sh{B}/\sh{K})
\]
\sref[II]{II.1.4.10}.
The morphism $h:Z'\to Z$ is finite \sref[II]{II.6.1.5},~(iii);
consequently $Z'$ is proper over $Y$ \sref[II]{II.6.1.11} and
\sref[II]{II.5.4.2},~(ii).

We must prove that, for every $x\in X$ and every open neighbourhood $U$
of $x$, there exist an integer $n>0$ and a section $s$ of
$\sh{L}^{\otimes n}$ over $X$ such that
\[
  x\in X_s\subset U
\]
\sref[II]{II.4.5.2}.
We distinguish two cases.

\smallskip
\noindent
$1^\circ$ Suppose first that $x\in X-Z$.
We may clearly suppose that $U\subset X-Z$, so that
$U'=g^{-1}(U)$ does not meet $Z'$.
Since $\sh{L}'$ is ample, there exist $n>0$ and a section $s'$ of
$\sh{L}'^{\otimes n}$ over $X'$ such that
\[
  x'=g^{-1}(x)\in X'_{s'}\subset g^{-1}(U)
\]
\sref[II]{II.4.5.2}.
We may also suppose that
$\sh{K}'\otimes\sh{L}'^{\otimes n}$ is generated by its global sections
\sref[II]{II.4.5.5}.
As $\sh{K}'_{x'}=\sh{O}_{X',x'}$, one of these sections, say $s''$, is
nonzero at $x'$.
Multiplying it by $s'$, and thus replacing $n$ by $2n$, we may arrange
that
\[
  x'\in X'_{s''}\subset g^{-1}(U).
\]
By \sref[0]{0.5.4.10}, there is a canonical isomorphism
\[
  \Gamma(X,\sh{K}\otimes\sh{L}^{\otimes n})
  \xrightarrow{\sim}
  \Gamma(X',\sh{K}'\otimes\sh{L}'^{\otimes n}).
\]
The section $s$ of $\sh{K}\otimes\sh{L}^{\otimes n}$ corresponding to
$s''$ plainly has the required properties.

\smallskip
\noindent
$2^\circ$ Suppose now that $x\in Z$.
Let $\sh{J}$ be the coherent ideal of $\sh{O}_X$ defining the reduced
closed subprescheme whose underlying space is $X-U$, and consider in
$\sh{B}$ the coherent ideals

\oldpage[III]{114}

\[
  \sh{J}\sh{B}
  \qquad\text{and}\qquad
  \sh{J}\sh{K}
  =
  \sh{J}(\sh{K}\sh{B})
  =
  \sh{K}(\sh{J}\sh{B}).
\]
We have the diagram of inclusions
\[
  \xymatrix@R=1.5em{
    \sh{J}\sh{B}\ar@{^{(}->}[r] &
    \sh{B}\\
    \sh{J}\ar@{^{(}->}[r]\ar@{^{(}->}[u] &
    \sh{A}\ar@{^{(}->}[u]\\
    \sh{J}\sh{K}\sh{B}=\sh{J}\sh{K}
      \ar@{^{(}->}[r]\ar@{^{(}->}[u] &
    \sh{K}\ar@{^{(}->}[u]
  }
  \label{III.2.6.2.6}\tag{2.6.2.6}
\]
Let $\sh{J}'=(\sh{J}\sh{B})^\sim$ be the coherent ideal of
$\sh{O}_{X'}$ such that
\[
  g_*(\sh{J}')=\sh{J}\sh{B}.
\]
Then
\[
  \sh{J}'\sh{K}'=(\sh{J}\sh{K}\sh{B})^\sim,
  \qquad
  \sh{J}'/\sh{J}'\sh{K}'
  =
  (\sh{J}\sh{B}/\sh{J}\sh{K}\sh{B})^\sim
\]
\sref[II]{II.1.4.4}.
Since $\sh{J}|_V=\sh{J}\sh{K}|_V$ on every open subset $V$ disjoint
from $Z$, the support of $\sh{J}'/\sh{J}'\sh{K}'$ is contained in
$Z'$.
Because $Z'$ is proper over $Y$, \sref{III.2.2.4} shows that, for all
sufficiently large $n$, the canonical map
\[
  \Gamma(X',\sh{J}'\otimes\sh{L}'^{\otimes n})
  \longrightarrow
  \Gamma\bigl(
    X',
    (\sh{J}'/\sh{J}'\sh{K}')\otimes\sh{L}'^{\otimes n}
  \bigr)
\]
is surjective.
By \sref[0]{0.5.4.10}, it follows that the canonical map
\[
  \Gamma(X,\sh{J}\sh{B}\otimes\sh{L}^{\otimes n})
  \longrightarrow
  \Gamma\bigl(
    X,
    (\sh{J}\sh{B}/\sh{J}\sh{K}\sh{B})
      \otimes\sh{L}^{\otimes n}
  \bigr)
\]
is surjective.

Let $i:Z\to X$ and $i':Z'\to X'$ be the canonical immersions.
They fit into the commutative diagram
\[
  \xymatrix{
    X'\ar[d]_g &
    Z'\ar[l]_{i'}\ar[d]^h\\
    X &
    Z\ar[l]^{i}
  }
\]
Put
\[
  \sh{M}=i^*(\sh{L}),
  \qquad
  \sh{M}'=i'^*(\sh{L}').
\]
Since $\sh{L}'$ is ample, $\sh{M}'$ is ample
\sref[II]{II.5.1.12}, and
\[
  \sh{M}'=h^*(\sh{M}).
\]
The Noetherian induction hypothesis, applied because $Z\neq X$, now
shows that $\sh{M}$ is ample.
Consequently, for all sufficiently large $n$,
\[
  i^*(\sh{J}/\sh{J}\sh{K})\otimes\sh{M}^{\otimes n}
\]
is generated by its sections over $Z$ \sref[II]{II.4.5.5}.
Since
\[
  \sh{J}/\sh{J}\sh{K}
  =
  i_*\bigl(i^*(\sh{J}/\sh{J}\sh{K})\bigr),
\]
\sref[0]{0.5.4.10} gives a section $s$ of
\[
  (\sh{J}/\sh{J}\sh{K})\otimes\sh{L}^{\otimes n}
\]
over $X$, for some sufficiently large $n$, such that $s(x)\neq0$.
Indeed, $\sh{J}_x=\sh{A}_x$ by the definition of $\sh{J}$, whereas
$\sh{K}_x\neq\sh{A}_x$ by hypothesis.

The diagram \eqref{III.2.6.2.6} shows that $s$ is also a section of
\[
  (\sh{J}\sh{B}/\sh{J}\sh{K}\sh{B})
  \otimes\sh{L}^{\otimes n}
\]
over $X$.
It is therefore the canonical image of a section
\[
  t\in\Gamma(X,\sh{J}\sh{B}\otimes\sh{L}^{\otimes n}).
\]
By definition, the image of $t$ modulo
$\sh{J}\sh{K}\sh{B}\otimes\sh{L}^{\otimes n}$ lies in
\[
  \frac{\sh{J}\otimes\sh{L}^{\otimes n}}
       {\sh{J}\sh{K}\sh{B}\otimes\sh{L}^{\otimes n}}.
\]
The diagram \eqref{III.2.6.2.6} therefore implies that $t$ is a section
of $\sh{J}\otimes\sh{L}^{\otimes n}$, hence in particular a section of
$\sh{L}^{\otimes n}$.
We have $t(x)\neq0$, so $x\in X_t$.
By the definition of $\sh{J}$, the section $t$ vanishes on $X-U$, the
support of $\sh{O}_X/\sh{J}$.
Thus $X_t\subset U$, and the proof is complete.
\end{env}
\end{proof}

\oldpage[III]{115}

\begin{env}[2.6.3]
\label{III.2.6.3}
\emph{Remark.}
When $X$ is proper over $Y$, Proposition~\sref{III.2.6.2} can be proved
more simply by arguing as in Chevalley's theorem
\sref[II]{II.6.7.1}, with the aid of Proposition~\sref{III.2.6.1} and
the lemma \sref[II]{II.6.7.1.1}.
\end{env}

\section{Finiteness theorem for proper morphisms}
\label{section:III.3}

\subsection{The d\'evissage lemma}
\label{subsection:III.3.1}

\begin{definition}[3.1.1]
\label{III.3.1.1}
Let $\C$ be an abelian category.
We say that a subset $\C'$ of the set of objects of $\C$ is \emph{exact} if $0\in\C'$ and if, for every exact sequence $0\to A'\to A\to A''\to 0$ in $\C$ such that two of the objects $A$, $A'$, $A''$ are in $\C'$, then the third is also in $\C'$.
\end{definition}

\begin{theorem}[3.1.2]
\label{III.3.1.2}
Let $X$ be a Noetherian prescheme; we denote by $\C$ the abelian category of coherent $\sh{O}_X$-modules.
Let $\C'$ be an exact subset of $\C$, $X'$ a closed subset of the underlying space of $X$.
Suppose that for every closed irreducible subset $Y$ of $X'$, with generic point $y$, there exists an $\sh{O}_X$-module $\sh{G}\in\C'$ of support $Y$ such that $\sh{G}_y$ is a $\kres(y)$-vector space of dimension~$1$.
Then every coherent $\sh{O}_X$-module with support contained in $X'$ is in $\C'$ (and in particular, if $X'=X$, then we have $\C'=\C$).
\end{theorem}

\begin{proof}
Consider the following property $\textbf{P}(Y)$ of a closed subset $Y$ of $X'$: every coherent $\sh{O}_X$-module with support contained in $Y$ is in $\C'$.
By virtue of the principle of Noetherian induction \sref[0]{0.2.2.2}, we see that we can reduce to showing that \emph{if $Y$ is a closed subset of $X'$ such that the property $\textbf{P}(Y')$ is true for every closed subset $Y'$ of $Y$, distinct from $Y$, then $\textbf{P}(Y)$ is true}.

Therefore, let $\sh{F}\in\C$ have support contained in $Y$, and we show that $\sh{F}\in\C'$.
Denote also by $Y$ the closed reduced subprescheme of $X$ having $Y$ for its underlying space \sref[I]{I.5.2.1}; it is defined by a coherent sheaf of ideals $\sh{J}$ of $\sh{O}_X$.
We know \sref[I]{I.9.3.4} that there exists an integer $n>0$ such that $\sh{J}^n\sh{F}=0$; for $1\leq k\leq n$, we thus have an exact sequence
\[
  0\to\sh{J}^{k-1}\sh{F}/\sh{J}^k\sh{F}\to\sh{F}/\sh{J}^k\sh{F}\to\sh{F}/\sh{J}^{k-1}\sh{F}\to 0
\]
 of coherent $\sh{O}_X$-modules (\sref[0]{0.5.3.6} and \sref[0]{0.5.3.3}); as $\C'$ is exact, we see, by induction on $k$, that it suffices to show that each of the $\sh{F}_k=\sh{J}^{k-1}\sh{F}/\sh{J}^k\sh{F}$ is in $\C'$.
We thus reduce to proving that $\sh{F}\in\C'$ under the additional hypothesis that $\sh{J}\sh{F}=0$; it is equivalent to say that $\sh{F}=j_*(j^*(\sh{F}))$, where $j$ is the canonical injection $Y\to X$.
Let us now consider two cases:
\begin{enumerate}
  \item[(a)] $Y$ is \emph{reducible}.
    Let $Y=Y'\cup Y''$, where $Y'$ and $Y''$ are closed subsets of $Y$, distinct from $Y$; denote also by $Y'$ and $Y''$ the closed reduced subpreschemes of $X$ having $Y'$ and $Y''$ for their respective underlying spaces, which are defined respectively by sheaves of ideals $\sh{J}'$ and $\sh{J}''$ of $\sh{O}_X$.
    Set $\sh{F}'=\sh{F}\otimes_{\sh{O}_X}(\sh{O}_X/\sh{J}')$ and $\sh{F}''=\sh{F}\otimes_{\sh{O}_X}(\sh{O}_X/\sh{J}'')$.
    The canonical homomorphisms $\sh{F}\to\sh{F}'$ and $\sh{F}\to\sh{F}''$ thus define a homomorphism $u:\sh{F}\to\sh{F}'\oplus\sh{F}''$.
    We show that for every $z\not\in Y'\cap Y''$, the homomorphism $u_z:\sh{F}_z\to\sh{F}_z'\oplus\sh{F}_z''$ is \emph{bijective}.
    Indeed, we have $\sh{J}'\cap\sh{J}''=\sh{J}$, since the question is local and
\oldpage[III]{116}
    the above equality follows from (\sref[I]{I.5.2.1} and \sref[I]{I.1.1.5}); if $z\not\in Y''$, then we have $\sh{J}_z'=\sh{J}_z$, hence $\sh{F}_z'=\sh{F}_z$ and $\sh{F}_z''=0$, which establishes our assertion in this case; we reason similarly for $z\not\in Y'$.
    As a result, the kernel and cokernel of $u$, which are in $\C$ \sref[0]{0.5.3.4}, have their support in $Y'\cap Y''$, and thus are in $\C'$ by hypothesis; for the same reason, $\sh{F}'$ and $\sh{F}''$ are in $\C'$, hence also $\sh{F}'\oplus\sh{F}''$, as $\C'$ is exact.
    The conclusion then follows from the consideration of the two exact sequences
    \[
      0\to\Im u\to\sh{F}'\oplus\sh{F}''\to\Coker u\to 0,
    \]
    \[
      0\to\Ker u\to\sh{F}\to\Im u\to 0,
    \]
    and the hypothesis that $\C'$ is exact.
  \item[(b)] $Y$ is irreducible, and as a result, the subprescheme $Y$ of $X$ is \emph{integral}.
    If $y$ is its generic point, then we have $(\sh{O}_Y)_y=\kres(y)$, and as $j^*(\sh{F})$ is a coherent $\sh{O}_Y$-module, $\sh{F}_y=(j^*(\sh{F}))_y$ is a $\kres(y)$-vector space of finite dimension~$m$.
    By hypothesis, there is a coherent $\sh{O}_X$-module $\sh{G}\in\C'$ (necessarily of support $Y$) such that $\sh{G}_y$ is a $\kres(y)$-vector space of dimension~$1$.
    As a result, there is a $\kres(y)$-isomorphism $(\sh{G}_y)^m\isoto\sh{F}_y$, which is also an $\sh{O}_Y$-isomorphism, and as $\sh{G}^m$ and $\sh{F}$ are coherent, there exists an open neighbourhood $W$ of $y$ in $X$ and an isomorphism $\sh{G}^m|W\isoto\sh{F}|W$ \sref[0]{0.5.2.7}.
    Let $\sh{H}$ be the graph of this isomorphism, which is a coherent $(\sh{O}_X|W)$-submodule of $(\sh{G}^m\oplus\sh{F})|W$, canonically isomorphic to $\sh{G}^m|W$ and to $\sh{F}|W$; there thus exists a coherent $\sh{O}_X$-submodule $\sh{H}_0$ of $\sh{G}^m\oplus\sh{F}$, inducing $\sh{H}$ on $W$ and $0$ on $X\setmin Y$, since $\sh{G}^m$ and $\sh{F}$ have $Y$ for their support~\sref[I]{I.9.4.7}.
    The restrictions $v:\sh{H}_0\to\sh{G}^m$ and $w:\sh{H}_0\to\sh{F}$ of the canonical projections of $\sh{G}^m\oplus\sh{F}$ are then homomorphisms of coherent $\sh{O}_X$-modules, which, on $W$ and on $X\setmin Y$, reduce to isomorphisms; in other words, the kernels and cokernels of $v$ and $w$ have their support in the closed set $Y\setmin(Y\cap W)$, distinct from $Y$.
    They are in $\C'$; on the other hand, we have $\sh{G}^m\in\C'$ since $\sh{G}\in\C'$ and since $\C'$ is exact.
    We conclude successively, by the exactness of $\C'$, that $\sh{H}_0\in\C'$, then $\sh{F}\in\C'$.
Q.E.D.
\end{enumerate}
\end{proof}

\begin{corollary}[3.1.3]
\label{III.3.1.3}
Suppose that the exact subset $\C'$ of $\C$ has in addition the property that any coherent direct factor of a coherent $\sh{O}_X$-module $\sh{M}\in\C'$ is also in $\C'$.
In this case, the conclusion of Theorem~\sref{III.3.1.2} is still valid when the condition ``$\sh{G}_y$ is a $\kres(y)$-vector space of dimension~$1$'' is replaced by $\sh{G}_y\neq 0$ (this is equivalent to $\Supp(\sh{G})=Y$).
\end{corollary}

\begin{proof}
The reasoning of Theorem~\sref{III.3.1.2} must be modified only in the case~(b); now $\sh{G}_y$ is a $\kres(y)$-vector space of dimension $q>0$, and as a result, we have an $\sh{O}_Y$-isomorphism $(\sh{G}_y)^m\isoto(\sh{F}_y)^q$; the end of the reasoning in Theorem~\sref{III.3.1.2} then proves that $\sh{F}^q\in\C'$, and the additional hypothesis on $\C'$ implies that $\sh{F}\in\C'$.
\end{proof}

\subsection{The finiteness theorem: the case of usual schemes}
\label{subsection:III.3.2}
\label{III.3.2}

\begin{theorem}[3.2.1]
\label{III.3.2.1}
Let $Y$ be a locally Noetherian prescheme, $f:X\to Y$ a proper morphism.
For every coherent $\sh{O}_X$-module $\sh{F}$, the $\sh{O}_Y$-modules $\RR^q f_*(\sh{F})$ are coherent for $q\geq 0$.
\end{theorem}

\begin{proof}
Since the question is local on $Y$, we can suppose $Y$ Noetherian, thus $X$ Noetherian~\sref[I]{I.6.3.7}.
The coherent $\sh{O}_X$-modules $\sh{F}$ for which the conclusion of Theorem~\sref{III.3.2.1} is true form an \emph{exact} subset $\C'$ of the category $\C$ of coherent $\sh{O}_X$-modules.
\oldpage[III]{117}
Indeed, let $0\to\sh{F}'\to\sh{F}\to\sh{F}''\to 0$ be an exact sequence of coherent $\sh{O}_X$-modules; suppose for example that $\sh{F}'$ and $\sh{F}''$ belong to $\C'$; we have the long exact sequence in cohomology
\[
  \RR^{q-1}f_*(\sh{F}'')\xrightarrow{\partial}\RR^q f_*(\sh{F}')\to\RR^q f_*(\sh{F})\to\RR^q f_*(\sh{F}'')\xrightarrow{\partial}\RR^{q+1}f_*(\sh{F}'),
\]
in which by hypothesis the outer four terms are coherent; it is the same for the middle term $\RR^q f_*(\sh{F})$ by (\sref[0]{0.5.3.4}~and~\sref[0]{0.5.3.3}).
We show in the same way that when $\sh{F}$ and $\sh{F}'$ (resp.~$\sh{F}$ and $\sh{F}''$) are in $\C'$, then so is $\sh{F}''$ (resp.~$\sh{F}'$).
In addition, every coherent \emph{direct factor} $\sh{F}'$ of an $\sh{O}_X$-module $\sh{F}\in\C'$ belongs to $\C'$: indeed, $\RR^q f_*(\sh{F}')$ is then a direct factor of $\RR^q f_*(\sh{F})$~(G,~II,~4.4.4), therefore it is of finite type, and as it is quasi-coherent~\sref{III.1.4.10}, it is coherent, as $Y$ is Noetherian.
By virtue of Corollary~\sref{III.3.1.3}, we reduce to proving that when $X$ is \emph{irreducible} with generic point $x$, there exists \emph{one} coherent $\sh{O}_X$-module $\sh{F}$ belonging to $\C'$, such that $\sh{F}_x\neq 0$: indeed, if this point is established, then it can be applied to any irreducible closed subprescheme $Y$ of $X$, since if $j:Y\to X$ is the canonical injection, then $f\circ j$ is proper~\sref[II]{II.5.4.2}, and if $\sh{G}$ is a coherent $\sh{O}_Y$-module with support $Y$, then $j_*(\sh{G})$ is a coherent $\sh{O}_X$-module such that $\RR^q(f\circ j)_*(\sh{G})=\RR^q f_*(j_*(\sh{G}))$~(G,~II,~4.9.1), therefore we can apply Corollary~\sref{III.3.1.3}.

By virtue of Chow's lemma~\sref[II]{II.5.6.2}, there exists an irreducible prescheme $X'$ and a \emph{projective} and surjective morphism $g:X'\to X$ such that $f\circ g:X'\to Y$ is \emph{projective}.
There exists an ample $\sh{O}_{X'}$-module $\sh{L}$ for $g$~\sref[II]{II.5.3.1}; we apply the fundamental theorem of projective morphisms~\sref{III.2.2.1} to $g:X'\to X$ and with $\sh{L}$: there thus exists an integer $n$ such that $\sh{F}=g_*(\sh{O}_{X'}(n))$ is a coherent $\sh{O}_X$-module and $\RR^q g_*(\sh{O}_{X'}(n))=0$ for all $q>0$; in addition, as $g^*(g_*(\sh{O}_{X'}(n)))\to\sh{O}_{X'}(n)$ is surjective for $n$ large enough~\sref{III.2.2.1}, we see that we can suppose, at the generic point $x$ of $X$, that we have $\sh{F}_x\neq 0$~\sref[II]{II.3.4.7}.
On the other hand, as $f\circ g$ is projective and $Y$ is Noetherian, the $\RR^q(f\circ g)_*(\sh{O}_{X'}(n))$ are \emph{coherent}~\sref{III.2.2.1}.
This being so, $\RR^\bullet(f\circ g)_*(\sh{O}_{X'}(n))$ is the abutment of a Leray spectral sequence, whose $E_2$-term is given by $E_2^{pq}=\RR^p f_*(\RR^q g_*(\sh{O}_{X'}(n)))$; the above shows that this spectral sequence degenerates, and we then know~\sref[0]{0.11.1.6} that $E_2^{p0}=\RR^p f_*(\sh{F})$ is isomorphic to $\RR^p(f\circ g)_*(\sh{O}_{X'}(n))$, which finishes the proof.
\end{proof}

\begin{corollary}[3.2.2]
\label{III.3.2.2}
Let $Y$ be a locally Noetherian prescheme.
For every proper morphism $f:X\to Y$, the direct image under $f$ of any coherent $\sh{O}_X$-module is a coherent $\sh{O}_Y$-module.
\end{corollary}

\begin{corollary}[3.2.3]
\label{III.3.2.3}
Let $A$ be a Noetherian ring, $X$ a proper scheme over $A$; for every coherent $\sh{O}_X$-module $\sh{F}$, the $\HH^p(X,\sh{F})$ are $A$-modules of finite type, and there exists an integer $r>0$ such that for every coherent $\sh{O}_X$-module $\sh{F}$ and all $p>r$, $\HH^p(X,\sh{F})=0$.
\end{corollary}

\begin{proof}
The second assertion has already been proved~\sref{III.1.4.12}; the first follows from the finiteness theorem~\sref{III.3.2.1}, taking into account Corollary~\sref{III.1.4.11}.
\end{proof}

In particular, if $X$ is a \emph{proper algebraic scheme} over a field $k$, then, for every coherent $\sh{O}_X$-module $\sh{F}$, the $\HH^p(X,\sh{F})$ are \emph{finite-dimensional} $k$-vector spaces.

\begin{corollary}[3.2.4]
\label{III.3.2.4}
Let $Y$ be a locally Noetherian prescheme, $f:X\to Y$ a morphism of finite type.
For every coherent $\sh{O}_X$-module $\sh{F}$ whose support is proper over $Y$~\sref[II]{II.5.4.10}, the $\sh{O}_Y$-modules $\RR^q f_*(\sh{F})$ are coherent.
\end{corollary}

\oldpage[III]{118}
\begin{proof}
Since the question is local on $Y$, we can suppose $Y$ Noetherian, and it is the same for $X$~\sref[I]{I.6.3.7}.
By hypothesis, every closed subprescheme $Z$ of $X$ whose underlying space is $\Supp(\sh{F})$ is proper over $Y$, in other words, if $j:Z\to X$ is the canonical injection, then $f\circ j:Z\to Y$ is proper.
We can suppose that $Z$ is such that $\sh{F}=j_*(\sh{G})$, where $\sh{G}=j^*(\sh{F})$ is a coherent $\sh{O}_Z$-module~\sref[I]{I.9.3.5}; as we have $\RR^q f_*(\sh{F})=\RR^q(f\circ j)_*(\sh{G})$ by Corollary~\sref{III.1.3.4}, the conclusion follows immediately from Theorem~\sref{III.3.2.1}.
\end{proof}

\subsection{Generalization of the finiteness theorem (usual schemes)}
\label{subsection:III.3.3}

\begin{proposition}[3.3.1]
\label{III.3.3.1}
Let $Y$ be a Noetherian prescheme, $\sh{S}$ a quasi-coherent $\sh{O}_Y$-algebra of finite type, graded in positive degrees, $Y'=\Proj(\sh{S})$, and $g:Y'\to Y$ the structure morphism.
Let $f:X\to Y$ be a proper morphism, $\sh{S}'=f^*(\sh{S})$, $\sh{M}=\bigoplus_{k\in\bb{Z}}\sh{M}_k$ a quasi-coherent graded $\sh{S}'$-module of finite type.
Then the $\RR^p f_*(\sh{M})=\bigoplus_{k\in\bb{Z}}\RR^p f_*(\sh{M}_k)$ are graded $\sh{S}$-modules of finite type for all $p$.
Suppose in addition that $\sh{S}$ is generated by $\sh{S}_1$; then, for every $p\in\bb{Z}$, there exists an integer $k_p$ such that for all $k\geq k_p$ and all $r\geq 0$, we have
\[
\label{III.3.3.1.1}
  \RR^p f_*(\sh{M}_{k+r})=\sh{S}_r\RR^p f_*(\sh{M}_k).
  \tag{3.3.1.1}
\] 
\end{proposition}

\begin{proof}
The first assertion is identical to the statement of Theorem~\sref{III.2.4.1}[i], where we have simply replaced ``projective morphism'' by ``proper morphism''.
In the proof of Theorem~\sref{III.2.4.1}[i], the hypothesis on $f$ was only used to show (with the notation of this proof) that $\RR^p f_*'(\widetilde{\sh{M}})$ is a coherent $\sh{O}_{Y'}$-module.
With the hypothesis of Proposition~\sref{III.3.3.1}, $f'$ is proper~\sref[II]{II.5.4.2}[iii], so we can resume without change in the proof of Theorem~\sref{III.2.4.1}[i], thanks to the finiteness theorem~\sref{III.3.2.1}.

As for the second assertion, it suffices to remark that there is a finite affine open cover $(U_i)$ of $Y$ such that the restrictions to the $U_i$ of the two sides of~(\hyperref[III.3.3.1.1]{3.3.1.1}) are equal for all $k\geq k_{p,i}$~\sref[II]{II.2.1.6}[ii]; it suffices to take for $k_p$ the largest of the $k_{p,i}$.
\end{proof}

\begin{corollary}[3.3.2]
\label{III.3.3.2}
Let $A$ be a Noetherian ring, $\mathfrak{m}$ an ideal of $A$, $X$ a proper $A$-scheme, and $\sh{F}$ a coherent $\sh{O}_X$-module.
Then, for all $p\geq 0$, the direct sum $\bigoplus_{k\geq 0}\HH^p(X,\mathfrak{m}^k\sh{F})$ is a module of finite type over the ring $S=\bigoplus_{k\geq 0}\mathfrak{m}^k$; in particular, there exists an integer $k_p\geq 0$ such that for all $k\geq k_p$ and all $r\geq 0$, we have
\[
\label{III.3.3.2.1}
  \HH^p(X,\mathfrak{m}^{k+r}\sh{F})=\mathfrak{m}^r\HH^p(X,\mathfrak{m}^k\sh{F}).
  \tag{3.3.2.1}
\]
\end{corollary}

\begin{proof}
It suffices to apply Proposition~\sref{III.3.3.1} with $Y=\Spec(A)$, $\sh{S}=\widetilde{S}$, $\sh{M}_k=\mathfrak{m}^k\sh{F}$, taking into account Corollary~\sref{III.1.4.11}.
\end{proof}

It should be remembered that the $S$-module structure on $\bigoplus_{k\geq 0}\HH^p(X,\mathfrak{m}^k\sh{F})$ is obtained by considering, for every $a\in\mathfrak{m}^r$, the map $\HH^p(X,\mathfrak{m}^k\sh{F})\to\HH^p(X,\mathfrak{m}^{k+r}\sh{F})$, which comes from the passage to cohomology of the multiplication map $\mathfrak{m}^k\sh{F}\to\mathfrak{m}^{k+r}\sh{F}$ defined by $a$~\sref{III.2.4.1}.

\subsection{Finiteness theorem: the case of formal schemes}
\label{subsection:III.3.4}

\oldpage[III]{119}
The results of this section (except the definition~\sref{III.3.4.1}) will not be used in the rest of this chapter.

\begin{env}[3.4.1]
\label{III.3.4.1}
Let $\mathfrak{X}$ and $\mathfrak{S}$ be two locally Noetherian formal preschemes~\sref[I]{I.10.4.2}, $f:\mathfrak{X}\to\mathfrak{S}$ a morphism of formal preschemes.
We say that $f$ is a \emph{proper} morphism if it satisfies the following conditions:
\begin{enumerate}
  \item[1st.] \emph{$f$ is a morphism of finite type~\sref[I]{I.10.13.3}}.
  \item[2nd.] \emph{If $\sh{K}$ is a sheaf of ideals of definition for $\mathfrak{S}$ and if we set $\sh{J}=f^*(\sh{K})\sh{O}_\mathfrak{X}$, $X_0=(\mathfrak{X},\sh{O}_\mathfrak{X}/\sh{J})$, $S_0=(\mathfrak{S},\sh{O}_\mathfrak{S}/\sh{K})$, then the morphism $f_0:X_0\to S_0$ induced by $f$~\sref[I]{I.10.5.6} is proper}.
\end{enumerate}
It is immediate that this definition does not depend on the sheaf of ideals of definition $\sh{K}$ for $\mathfrak{S}$ considered; indeed, if $\sh{K}'$ is a second sheaf of ideals of definition such that $\sh{K}'\subset\sh{K}$, and if we set $\sh{J}'=f^*(\sh{K}')\sh{O}_\mathfrak{X}$, $X_0'=(\mathfrak{X},\sh{O}_\mathfrak{X}/\sh{J}')$, $S_0'=(\mathfrak{S},\sh{O}_\mathfrak{S}/\sh{K}')$, then the morphism $f_0':X_0'\to S_0'$ induced by $f$ is such that the diagram
\[
  \xymatrix{
    X_0\ar[r]^{f_0}\ar[d]_i &
    S_0\ar[d]^j\\
    X_0'\ar[r]_{f_0'} &
    S_0'
  }
\]
is commutative, $i$ and $j$ being surjective immersions; it is equivalent to say that $f_0$ or $f_0'$ is proper, by virtue of~\sref[II]{II.5.4.5}.

We note that, for all $n\geq 0$, if we set $X_n=(\mathfrak{X},\sh{O}_\mathfrak{X}/\sh{J}^{n+1})$, $S_n=(\mathfrak{S},\sh{O}_\mathfrak{S}/\sh{K}^{n+1})$, then the morphism $f_n:X_n\to S_n$ induced by $f$~\sref[I]{I.10.5.6} is proper for all $n$ whenever it is for $n=0$~\sref[II]{II.5.4.6}.

If $g:Y\to Z$ is a proper morphism of locally Noetherian usual preschemes, $Z'$ a closed subset of $Z$, $Y'$ a closed subset of $Y$ such that $g(Y')\subset Z'$, then the extension $\widehat{g}:Y_{/Y'}\to Z_{/Z'}$ of $g$ to the completions~\sref[I]{I.10.9.1} is a proper morphism of formal preschemes, as it follows from the definition and from~\sref[II]{II.5.4.5}.

Let $\mathfrak{X}$ and $\mathfrak{S}$ be two locally Noetherian formal preschemes, $f:\mathfrak{X}\to\mathfrak{S}$ a morphism \emph{of finite type}~\sref[I]{I.10.13.3}; the notation being the same as above, we say that a subset $Z$ of the underlying space of $\mathfrak{X}$ is \emph{proper} over $\mathfrak{S}$ (or proper for $f$) if, considered as a subset of $X_0$, $Z$ is \emph{proper over $S_0$}~\sref[II]{II.5.4.10}.
All the properties of proper subsets of usual preschemes stated in~\sref[II]{II.5.4.10} are still true for the proper subsets of formal preschemes, as it follows immediately from the definitions.
\end{env}

\begin{theorem}[3.4.2]
\label{III.3.4.2}
Let $\mathfrak{X}$ and $\mathfrak{Y}$ be locally Noetherian formal preschemes, $f:\mathfrak{X}\to\mathfrak{Y}$ a proper morphism.
For every coherent $\sh{O}_\mathfrak{X}$-module $\sh{F}$, the $\sh{O}_\mathfrak{Y}$-modules $\RR^q f_*(\sh{F})$ are coherent for all $q\geq 0$.
\end{theorem}

Let $\sh{J}$ be a sheaf of ideals of definition for $\mathfrak{Y}$, $\sh{K}=f^*(\sh{J})\sh{O}_\mathfrak{X}$, and consider the $\sh{O}_\mathfrak{X}$-modules
\[
\label{III.3.4.2.1}
  \sh{F}_k=\sh{F}\otimes_{\sh{O}_\mathfrak{Y}}(\sh{O}_\mathfrak{Y}/\sh{J}^{k+1})=\sh{F}/\sh{K}^{k+1}\sh{F}\quad(k\geq 0)
  \tag{3.4.2.1}
\]
which evidently form a \emph{projective system} of topological $\sh{O}_\mathfrak{X}$-modules, such that
\oldpage[III]{120}
\[
  \sh{F}=\varprojlim_k\sh{F}_k
\]
by \sref[I]{I.10.11.3}.
On the other hand, it follows from Theorem~\sref{III.3.4.2} that each of the $\RR^q f_*(\sh{F})$, being coherent, is naturally equipped with a topological $\sh{O}_\mathfrak{Y}$-module structure~\sref[I]{I.10.11.6}, and so are the $\RR^q f_*(\sh{F}_k)$.
The canonical homomorphisms $\sh{F}\to\sh{F}_k=\sh{F}/\sh{K}^{k+1}\sh{F}$ canonically correspond to homomorphisms
\[
  \RR^q f_*(\sh{F})\to\RR^q f_*(\sh{F}_k),
\]
which are necessarily continuous for the topological $\sh{O}_\mathfrak{Y}$-module structures above~\sref[I]{I.10.11.6}, and form a projective system, giving the limit a canonical functorial homomorphism
\[
\label{III.3.4.2.2}
  \RR^q f_*(\sh{F})\to\varprojlim_k\RR^q f_*(\sh{F}_k),
  \tag{3.4.2.2}
\]
which will be a continuous homomorphism of topological $\sh{O}_\mathfrak{Y}$-modules.
We will prove along with Theorem~\sref{III.3.4.2} the
\begin{corollary}[3.4.3]
\label{III.3.4.3}
Each of the homomorphisms~(\hyperref[III.3.4.2.2]{3.4.2.2}) is a topological isomorphism.
In addition, if $\mathfrak{Y}$ is Noetherian, then the projective system $(\RR^q f_*(\sh{F}/\sh{K}^{k+1}\sh{F}))_{k\geq 0}$ satisfies the \emph{(ML)}-condition~\sref[0]{0.13.1.1}.
\end{corollary}
We will begin by establishing Theorem~\sref{III.3.4.2} and Corollary~\sref{III.3.4.3} when $\mathfrak{Y}$ is a Noetherian formal affine scheme~\sref[I]{I.10.4.1}:
\begin{corollary}[3.4.4]
\label{III.3.4.4}
Under the hypotheses of Theorem~\sref{III.3.4.2}, suppose in addition that $\mathfrak{Y}=\Spf(A)$, where $A$ is an adic Noetherian ring.
Let $\mathfrak{J}$ be an ideal of definition for $A$, and set $\sh{F}_k=\sh{F}/\mathfrak{J}^{k+1}\sh{F}$ for $k\geq 0$.
Then the $\HH^n(\mathfrak{X},\sh{F})$ are $A$-modules of finite type; the projective system $(\HH^n(\mathfrak{X},\sh{F}_k))_{k\geq 0}$ satisfies the \emph{(ML)}-condition for all $n$; if we set
\[
\label{III.3.4.4.1}
  N_{n,k}=\Ker\big(\HH^n(\mathfrak{X},\sh{F})\to\HH^n(\mathfrak{X},\sh{F}_k)\big)
  \tag{3.4.4.1}
\]
(also equal to $\Im(\HH^n(\mathfrak{X},\mathfrak{J}^{k+1}\sh{F})\to\HH^n(\mathfrak{X},\sh{F}))$ by the exact sequence in cohomology), then the $N_{n,k}$ define on $\HH^n(\mathfrak{X},\sh{F})$ a $\mathfrak{J}$-good filtration~\sref[0]{0.13.7.7}; finally, the canonical homomorphism
\[
\label{III.3.4.4.2}
  \HH^n(\mathfrak{X},\sh{F})\to\varprojlim_k\HH^n(\mathfrak{X},\sh{F}_k)
  \tag{3.4.4.2}
\]
is a topological isomorphism for all $n$ (the left hand side being equipped with the $\mathfrak{J}$-adic topology, the $\HH^n(\mathfrak{X},\sh{F}_k)$ with the discrete topology).
\end{corollary}

Set
\[
\label{III.3.4.4.3}
  S=\gr(A)=\bigoplus_{k\geq 0}\mathfrak{J}^k/\mathfrak{J}^{k+1},\ \sh{M}=\gr(\sh{F})=\bigoplus_{k\geq 0}\mathfrak{J}^k\sh{F}/\mathfrak{J}^{k+1}\sh{F}. \tag{3.4.4.3}
\]
We know that $\mathfrak{J}^\Delta$ is a sheaf of ideals of definition for $\mathfrak{Y}$~\sref[I]{I.10.3.1}; let $\sh{K}=f^*(\mathfrak{J}^\Delta)\sh{O}_\mathfrak{X}$, $X_0=(\mathfrak{X},\sh{O}_\mathfrak{X}/\sh{K})$, $Y_0=(\mathfrak{Y},\sh{O}_\mathfrak{Y}/\mathfrak{J}^\Delta)=\Spec(A_0)$, with $A_0=A/\mathfrak{J}$.
It is clear that the $\sh{M}_k=\mathfrak{J}^k\sh{F}/\mathfrak{J}^{k+1}\sh{F}$ are coherent $\sh{O}_{X_0}$-modules~\sref[I]{I.10.11.3}.
Consider on the other hand the quasi-coherent graded $\sh{O}_{X_0}$-algebra
\[
\label{III.3.4.4.4}
  \sh{S}=\sh{O}_{X_0}\otimes_{A_0}S.
  \tag{3.4.4.4}
\]
Multiplication defines a surjective homomorphism of graded algebras
\[
  \sh{S}\longrightarrow\gr(\sh{O}_\mathfrak{X})
  =\bigoplus_{k\geq 0}\sh{K}^k/\sh{K}^{k+1},
\]
through which $\sh{M}$ is a graded $\sh{S}$-module.

The hypothesis that $\sh{F}$ is an $\sh{O}_\mathfrak{X}$-module of finite type implies first that $\sh{M}$ is
\oldpage[III]{121}
a graded $\sh{S}$-module \emph{of finite type}.
Indeed, the question is local on $\mathfrak{X}$, and we can thus suppose that $\mathfrak{X}=\Spf(B)$, where $B$ is an adic Noetherian ring, and $\sh{F}=N^\Delta$, where $N$ is a $B$-module of finite type~\sref[I]{I.10.10.5}.
Set $B_0=B/\mathfrak{J}B$ and $N_0=N/\mathfrak{J}N$.
Then $X_0=\Spec(B_0)$, and the quasi-coherent $\sh{O}_{X_0}$-modules $\sh{S}$ and $\sh{M}$ are respectively associated with
\[
  S'=\bigoplus_{k\geq 0}((\mathfrak{J}^k/\mathfrak{J}^{k+1})\otimes_{A_0}B_0),\qquad
  M'=\bigoplus_{k\geq 0}\mathfrak{J}^kN/\mathfrak{J}^{k+1}N.
\]
For every $k\geq 0$, multiplication induces a surjective map
\[
  (\mathfrak{J}^k/\mathfrak{J}^{k+1})\otimes_{A_0}N_0
  \longrightarrow\mathfrak{J}^kN/\mathfrak{J}^{k+1}N,
  \qquad \overline a\otimes\overline n\longmapsto\overline{an}.
\]
These maps define a surjective homomorphism of graded $S'$-modules $S'\otimes_{B_0}N_0\longrightarrow M'$.
Since $N_0$ is a $B_0$-module of finite type, $M'$ is therefore an $S'$-module of finite type, hence our assertion~\sref[I]{I.1.3.13}.

As the morphism $f_0:X_0\to Y_0$ is \emph{proper} by hypothesis, we can apply Proposition~\sref{III.3.3.1} to $\sh{S}$, $\sh{M}$, and the morphism $f_0$: taking into account Corollary~\sref{III.1.4.11}, we conclude that for \emph{all $n\geq 0$}, $\bigoplus_{k\geq 0}\HH^n(X_0,\sh{M}_k)$ is a graded $S$-module \emph{of finite type}.
This proves that the condition ($\text{F}_n$) of~\sref[0]{0.13.7.7} is satisfied for \emph{all $n\geq 0$}, when we consider the strict projective system $(\sh{F}/\mathfrak{J}^k\sh{F})_{k\geq 0}$ of sheaves of abelian groups on $X_0$, each equipped with its natural ``filtered $A$-module'' structure.
We can thus apply~\sref[0]{0.13.7.7}, which proves that:
\begin{enumerate}
  \item[1st.] The projective system $(\HH^n(\mathfrak{X},\sh{F}_k))_{k\geq 0}$ satisfies the (ML)-condition.
  \item[2nd.] If $\HH^{\prime n}=\varprojlim_k\HH^n(\mathfrak{X},\sh{F}_k)$, then $\HH^{\prime n}$ is an $A$-module of finite type.
  \item[3rd.] The filtration defined on $\HH^{\prime n}$ by the kernels of the canonical homomorphisms $\HH^{\prime n}\to\HH^n(\mathfrak{X},\sh{F}_k)$ is $\mathfrak{J}$-good.
\end{enumerate}

Note that on the other hand, if we set $X_k=(\mathfrak{X},\sh{O}_\mathfrak{X}/\sh{K}^{k+1})$, then $\sh{F}_k$ is a coherent $\sh{O}_{X_k}$-module~\sref[I]{I.10.11.3}, and if $U$ is an affine open set in $X_0$, then $U$ is also an affine open set in each of the $X_k$~\sref[I]{I.5.1.9}, so $\HH^n(U,\sh{F}_k)=0$ for all $n>0$ and all $k$~\sref{III.1.3.1} and $\HH^0(U,\sh{F}_k)\to\HH^0(U,\sh{F}_h)$ is surjective for $h\leq k$~\sref[I]{I.1.3.9}.
We are thus in the conditions of~\sref[0]{0.13.3.2} and applying~\sref[0]{0.13.3.1} proves that $\HH^{\prime n}$ canonically identifies with $\HH^n(\mathfrak{X},\varprojlim_k\sh{F}_k)=\HH^n(\mathfrak{X},\sh{F})$; this finishes the proof of Corollary~\sref{III.3.4.4}.

\begin{env}[3.4.5]
\label{III.3.4.5}
We return to the proof of~\sref{III.3.4.2} and~\sref{III.3.4.3}.
We first prove the propositions for the case $\mathfrak{Y}=\Spf(A)$ envisaged in~\sref{III.3.4.4}; for this, for all $g\in A$, apply~\sref{III.3.4.4} to the Noetherian affine formal scheme induced on the open set $\mathfrak{Y}_g=\mathfrak{D}(g)$ of $\mathfrak{Y}$, which is equal to $\Spf(A_{\{g\}})$, and to the formal prescheme induced by $\mathfrak{X}$ on $f^{-1}(\mathfrak{Y}_g)$; note that $\mathfrak{Y}_g$ is also an affine open set in the prescheme $Y_k=(\mathfrak{Y},\sh{O}_\mathfrak{Y}/(\mathfrak{J}^\Delta)^{k+1})$, and as $\sh{F}_k$ is a coherent $\sh{O}_{X_k}$-module, we have
\[
  \HH^n(f^{-1}(\mathfrak{Y}_g),\sh{F}_k)=\Gamma(\mathfrak{Y}_g,\RR^n f_*(\sh{F}_k))
\]
for all $k\geq 0$ by virtue of Corollary~\sref{III.1.4.11}.
The canonical homomorphism
\[
  \HH^n(f^{-1}(\mathfrak{Y}_g),\sh{F})\to\varprojlim_k\Gamma(\mathfrak{Y}_g,\RR^n f_*(\sh{F}_k))
\]
is an isomorphism; but we have~\sref[0]{0.3.2.6}
\[
  \varprojlim_k\Gamma(\mathfrak{Y}_g,\RR^n f_*(\sh{F}_k))=\Gamma(\mathfrak{Y}_g,\varprojlim_k\RR^n f_*(\sh{F}_k)),
\]
\oldpage[III]{122}
and as the sheaf $\RR^n f_*(\sh{F})$ is the sheaf associated to the presheaf $\mathfrak{Y}_g\mapsto\HH^n(f^{-1}(\mathfrak{Y}_g),\sh{F})$ on the $\mathfrak{Y}_g$~\sref[0]{0.3.2.1}, we have shown that the homomorphism~(\hyperref[III.3.4.2.2]{3.4.2.2}) is \emph{bijective}.
Let us now prove that $\RR^n f_*(\sh{F})$ is a coherent $\sh{O}_\mathfrak{Y}$-module, and more precisely that we have
\[
\label{III.3.4.5.1}
  \RR^n f_*(\sh{F})=\big(\HH^n(\mathfrak{X},\sh{F})\big)^\Delta.
  \tag{3.4.5.1}
\]

With the above notation, we have, since $\sh{F}_k$ is a coherent $\sh{O}_{X_k}$-module~\sref{III.1.4.13},
\[
  \Gamma(\mathfrak{Y}_g,\RR^n f_*(\sh{F}_k))=(\Gamma(\mathfrak{Y},\RR^n f_*(\sh{F}_k)))_g=(\HH^n(\mathfrak{X},\sh{F}_k))_g.
\]

Now the $\HH^n(\mathfrak{X},\sh{F}_k)$ form a projective system satisfying (ML), and their projective limit $\HH^n(\mathfrak{X},\sh{F})$ is an $A$-module of finite type.
We conclude~\sref[0]{0.13.7.8} that we have
\[
  \varprojlim_k\big((\HH^n(\mathfrak{X},\sh{F}_k))_g\big)=\HH^n(\mathfrak{X},\sh{F})\otimes_A A_{\{g\}}=\Gamma(\mathfrak{Y}_g,(\HH^n(\mathfrak{X},\sh{F}))^\Delta),
\]
taking into account~\sref[I]{I.10.10.8} applied to $A$ and $A_{\{g\}}$; this proves~(\hyperref[III.3.4.5.1]{3.4.5.1}) since $\Gamma(\mathfrak{Y}_g,\RR^n f_*(\sh{F}))=\varprojlim_k\Gamma(\mathfrak{Y}_g,\RR^n f_*(\sh{F}_k))$.

As (\hyperref[III.3.4.2.2]{3.4.2.2}) is then an isomorphism of coherent $\sh{O}_\mathfrak{Y}$-modules, it is necessarily a \emph{topological} isomorphism~\sref[I]{I.10.11.6}.
Finally, it follows from the relations $\RR^n f_*(\sh{F}_k)=(\HH^n(\mathfrak{X},\sh{F}_k))^\Delta$ that the projective system $(\RR^n f_*(\sh{F}_k))_{k\geq 0}$ satisfies (ML)~\sref[I]{I.10.10.2}.

Once \sref{III.3.4.2} and \sref{III.3.4.3} are proved in the case where the formal prescheme $\mathfrak{Y}$ is affine Noetherian, it is immediate to pass to the general case for~\sref{III.3.4.2} and the first assertion of~\sref{III.3.4.3}, which are local on $\mathfrak{Y}$.
As for the second assertion of~\sref{III.3.4.3}, it suffices, $\mathfrak{Y}$ being Noetherian, to cover it by a finite number of Noetherian affine open sets $U_i$ and to note that the restrictions of the projective system $(\RR^q f_*(\sh{F}_k))$ to each of the $U_i$ satisfies (ML).
\end{env}

Along the way, we have in addition proved:
\begin{corollary}[3.4.6]
\label{III.3.4.6}
Under the hypotheses of Corollary~\sref{III.3.4.4}, the canonical homomorphism
\[
\label{III.3.4.6.1}
  \HH^q(\mathfrak{X},\sh{F})\to\Gamma(\mathfrak{Y},\RR^q f_*(\sh{F}))
  \tag{3.4.6.1}
\]
is bijective.
\end{corollary}

\section{The fundamental theorem of proper morphisms. Applications}
\label{section:III.4}


\subsection{The fundamental theorem}
\label{subsection:III.4.1}

\begin{env}[4.1.1]
\label{III.4.1.1}
Let $X$ and $Y$ be Noetherian usual preschemes, let $f:X\to Y$ be a
proper morphism, let $Y'$ be a closed subset of $Y$, and let
\[
  X'=f^{-1}(Y').
\]
We denote by $\widehat{X}$ and $\widehat{Y}$ the formal preschemes
$X_{/X'}$ and $Y_{/Y'}$, obtained by completing $X$ and $Y$ along
$X'$ and $Y'$, respectively \sref[I]{I.10.8.5}.
We denote by $\widehat{f}:\widehat{X}\to\widehat{Y}$ the extension of
$f$ to these completions \sref[I]{I.10.9.1}.
For every coherent $\sh{O}_X$-module $\sh{F}$, we denote by

\oldpage[III]{123}

\[
  \widehat{\sh{F}}=\sh{F}_{/X'}
\]
its completion along $X'$ \sref[I]{I.10.8.4}; this is a coherent
$\sh{O}_{\widehat{X}}$-module \sref[I]{I.10.8.8}.
\end{env}

\begin{env}[4.1.2]
\label{III.4.1.2}
Let $\sh{J}$ be a coherent ideal of $\sh{O}_Y$ such that
\[
  \Supp(\sh{O}_Y/\sh{J})=Y'
\]
\sref[I]{I.5.2.1}.
We know \sref[I]{I.4.4.5} that
\[
  \sh{K}=f^*(\sh{J})\sh{O}_X
\]
is a coherent ideal of $\sh{O}_X$ satisfying
\[
  \Supp(\sh{O}_X/\sh{K})=X'.
\]
For every $k\geq0$, consider the coherent $\sh{O}_X$-module
\[
  \sh{F}_k
  =
  \sh{F}\otimes_{\sh{O}_Y}(\sh{O}_Y/\sh{J}^{k+1})
  =
  \sh{F}/\sh{K}^{k+1}\sh{F}.
\]
The $\sh{O}_Y$-modules
\[
  \RR^n f_*(\sh{F})
  \quad\text{and}\quad
  \RR^n f_*(\sh{F}_k)
\]
are coherent for every $n$ \sref{III.3.2.1}.
For every $k\geq0$ and every $n$, the canonical homomorphism
$\sh{F}\to\sh{F}_k$ induces, by functoriality, a homomorphism
\[
  \RR^n f_*(\sh{F})
  \longrightarrow
  \RR^n f_*(\sh{F}_k).
  \tag{4.1.2.1}
  \label{III.4.1.2.1}
\]
Since $\sh{F}_k$ is an $\sh{O}_X/\sh{K}^{k+1}$-module,
$\RR^n f_*(\sh{F}_k)$ is an
$\sh{O}_Y/\sh{J}^{k+1}$-module \sref[0]{0.12.2.1}.
We therefore deduce from \eqref{III.4.1.2.1} a homomorphism
\[
  \RR^n f_*(\sh{F})
  \otimes_{\sh{O}_Y}(\sh{O}_Y/\sh{J}^{k+1})
  \longrightarrow
  \RR^n f_*(\sh{F}_k).
  \tag{4.1.2.2}
  \label{III.4.1.2.2}
\]
The two sides of \eqref{III.4.1.2.2} form projective systems.
The projective limit of the left-hand side is the completion
\[
  \bigl(\RR^n f_*(\sh{F})\bigr)_{/Y'},
\]
which we shall denote by
\[
  \widehat{\RR^n f_*(\sh{F})}.
\]
The homomorphisms \eqref{III.4.1.2.2} are compatible with the projective
systems and hence define, on passing to the limit, a canonical
homomorphism
\[
  \vphi_n:
  \widehat{\RR^n f_*(\sh{F})}
  \longrightarrow
  \varprojlim_k\RR^n f_*(\sh{F}_k).
  \tag{4.1.2.3}
  \label{III.4.1.2.3}
\]
Moreover, \eqref{III.4.1.2.2} is a homomorphism of
$\sh{O}_Y/\sh{J}^{k+1}$-modules.
By \sref[I]{I.10.8.3}, it may therefore be regarded as a continuous
homomorphism of topological $\sh{O}_{\widehat{Y}}$-modules endowed with
their pseudo-discrete topologies \sref[0]{0.3.8.1}.
Consequently, $\vphi_n$ is a continuous homomorphism of topological
$\sh{O}_{\widehat{Y}}$-modules.
\end{env}

\begin{env}[4.1.3]
\label{III.4.1.3}
Let $i:\widehat{X}\to X$ be the canonical morphism of ringed spaces
defined in \sref[I]{I.10.8.7}.
We have the commutative diagram
\[
  \xymatrix{
    X_k\ar[r]^{h_k}\ar[dr]_{i_k} &
    \widehat{X}\ar[d]^{i}\\
    & X
  }
  \tag{4.1.3.1}
  \label{III.4.1.3.1}
\]
where $X_k$ is the closed subprescheme of $X$ defined by
$\sh{K}^{k+1}$, $i_k$ is the canonical immersion, and $h_k$ is the
morphism of ringed spaces corresponding to the identity on the
underlying spaces and to the canonical homomorphism
\[
  \sh{O}_{\widehat{X}}
  \longrightarrow
  \sh{O}_X/\sh{K}^{k+1}
\]
\sref[I]{I.10.5.2}.
Moreover,
\[
  \widehat{\sh{F}}=i^*(\sh{F})
\]
up to canonical isomorphism \sref[I]{I.10.8.8}.
We have
\[
  \HH^n(X_k,i_k^*(\sh{F}_k))
  =
  \HH^n(X,\sh{F}_k)
  \tag{4.1.3.2}
  \label{III.4.1.3.2}
\]
up to canonical isomorphism.

\oldpage[III]{124}

Indeed,
\[
  \sh{F}_k=(i_k)_*(i_k^*(\sh{F}_k))
\]
(G,~II,~4.9.1).
The canonical homomorphism
\[
  \HH^n(\widehat{X},\widehat{\sh{F}})
  \longrightarrow
  \HH^n(X_k,h_k^*(\widehat{\sh{F}}))
\]
\sref[0]{0.12.1.3.5} may thus also be written
\[
  \HH^n(\widehat{X},\widehat{\sh{F}})
  \longrightarrow
  \HH^n(X,\sh{F}_k).
  \tag{4.1.3.3}
  \label{III.4.1.3.3}
\]
These homomorphisms form a projective system and therefore define a
canonical homomorphism
\[
  \psi_{n,X}:
  \HH^n(\widehat{X},\widehat{\sh{F}})
  \longrightarrow
  \varprojlim_k\HH^n(X,\sh{F}_k).
  \tag{4.1.3.4}
  \label{III.4.1.3.4}
\]
Replacing $X$ by an open subset of the form $f^{-1}(V)$, where $V$ is
an affine open subset of $Y$, and taking \sref{III.1.4.11} into account,
we obtain homomorphisms
\[
  \psi_{n,V}:
  \HH^n(\widehat{X}\cap f^{-1}(V),\widehat{\sh{F}})
  \longrightarrow
  \varprojlim_k\Gamma(V,\RR^n f_*(\sh{F}_k)).
  \tag{4.1.3.5}
  \label{III.4.1.3.5}
\]
These homomorphisms plainly commute with restriction from $V$ to a
smaller affine open subset.
They therefore define a canonical homomorphism of sheaves
\[
  \psi_n:
  \RR^n\widehat{f}_*(\widehat{\sh{F}})
  \longrightarrow
  \varprojlim_k\RR^n f_*(\sh{F}_k).
  \tag{4.1.3.6}
  \label{III.4.1.3.6}
\]
\end{env}

\begin{env}[4.1.4]
\label{III.4.1.4}
Finally, let $j:\widehat{Y}\to Y$ be the canonical morphism of ringed
spaces \sref[I]{I.10.8.7}.
Since $\RR^n f_*(\sh{F})$ is a coherent $\sh{O}_Y$-module
\sref{III.3.2.1}, we have, up to canonical isomorphism,
\[
  j^*(\RR^n f_*(\sh{F}))
  =
  \widehat{\RR^n f_*(\sh{F})}
\]
\sref[I]{I.10.8.8}.
There is consequently a canonical homomorphism
\[
  \rho_n:
  \widehat{\RR^n f_*(\sh{F})}
  =
  j^*(\RR^n f_*(\sh{F}))
  \longrightarrow
  \RR^n\widehat{f}_*(i^*(\sh{F}))
  =
  \RR^n\widehat{f}_*(\widehat{\sh{F}})
  \tag{4.1.4.1}
  \label{III.4.1.4.1}
\]
defined for ringed spaces in general; see the proof of
\sref{III.1.4.15}.
We shall show that the diagram
\[
  \xymatrix{
    \widehat{\RR^n f_*(\sh{F})}
      \ar[rr]^{\rho_n}\ar[dr]_{\vphi_n}
    &&
    \RR^n\widehat{f}_*(\widehat{\sh{F}})
      \ar[dl]^{\psi_n}\\
    &
    \displaystyle\varprojlim_k\RR^n f_*(\sh{F}_k)
    &
  }
  \tag{4.1.4.2}
  \label{III.4.1.4.2}
\]
is commutative.
It is enough to prove the commutativity of the diagram formed by the
corresponding homomorphisms of presheaves.
We may therefore restrict to the case where $Y$ is affine, and it is
enough to prove that
\[
  \xymatrix{
    \widehat{\HH^n(X,\sh{F})}
      \ar[rr]^{\rho_n}\ar[dr]_{\vphi_n}
    &&
    \HH^n(\widehat{X},\widehat{\sh{F}})
      \ar[dl]^{\psi_{n,X}}\\
    &
    \displaystyle\varprojlim_k\HH^n(X,\sh{F}_k)
    &
  }
  \tag{4.1.4.3}
  \label{III.4.1.4.3}
\]
is commutative.
The commutativity of \eqref{III.4.1.3.1}, together with the relations
between cohomology groups established in \sref{III.4.1.3}, gives the
commutative diagram

\oldpage[III]{125}

\[
  \xymatrix{
    \HH^n(X,\sh{F})
      \ar[r]\ar[dr]
    &
    \HH^n(\widehat{X},\widehat{\sh{F}})
      =
    \HH^n(\widehat{X},i^*(\sh{F}))
      \ar[d]\\
    &
    \HH^n(X,\sh{F}_k)
      =
    \HH^n(X_k,i_k^*(\sh{F})).
  }
\]
The commutativity of \eqref{III.4.1.4.3} follows immediately.
\end{env}

\begin{theorem}[4.1.5]
\label{III.4.1.5}
Let $f:X\to Y$ be a proper morphism of Noetherian preschemes, let
$Y'$ be a closed subset of $Y$, and put $X'=f^{-1}(Y')$.
Then, for every coherent $\sh{O}_X$-module $\sh{F}$,
\[
  \RR^n\widehat{f}_*(\widehat{\sh{F}})
\]
is a coherent $\sh{O}_{\widehat{Y}}$-module, and the homomorphisms
$\vphi_n$, $\psi_n$, and $\rho_n$ in diagram
\eqref{III.4.1.4.2} are topological isomorphisms.
\end{theorem}

\begin{proof}
It is enough to prove that $\vphi_n$ and $\psi_n$ are isomorphisms.
Indeed, $\RR^n f_*(\sh{F})$ is coherent \sref{III.3.2.1}, so its
completion $\widehat{\RR^n f_*(\sh{F})}$ is coherent
\sref[I]{I.10.8.8}; the bicontinuity of $\vphi_n$, $\psi_n$, and
$\rho_n$ is then automatic \sref[I]{I.10.11.6}.
\end{proof}

\begin{env}[4.1.6]
\label{III.4.1.6}
\emph{Remarks.}
\begin{enumerate}
  \item[(i)] Put
  \[
    \widehat{\sh{F}}_k
    =
    \widehat{\sh{F}}/
    \widehat{\sh{K}}^{\,k+1}\widehat{\sh{F}}.
  \]
  Then, immediately,
  \[
    \sh{F}_k=i_*(\widehat{\sh{F}}_k),
  \]
  and the canonical homomorphism \eqref{III.4.1.3.6} is precisely the
  homomorphism already defined in (\hyperref[III.3.4.2.2]{3.4.2.2}):
  \[
    \RR^n\widehat{f}_*(\widehat{\sh{F}})
    \longrightarrow
    \varprojlim_k
    \RR^n\widehat{f}_*(\widehat{\sh{F}}_k).
    \tag{4.1.6.1}
    \label{III.4.1.6.1}
  \]
  Thus the fact that $\psi_n$ is an isomorphism is a special case of
  \sref{III.3.4.3}.
  We shall nevertheless give a direct proof below, avoiding the
  delicate considerations concerning projective limits of spectral
  sequences \sref[0]{0.13.7} on which the general theorem
  \sref{III.3.4.3} rests.

  \item[(ii)] Since the $\psi_n$ are isomorphisms, saying that the
  $\vphi_n$ are isomorphisms is equivalent to saying that
  \[
    \rho_n=\psi_n^{-1}\circ\vphi_n
  \]
  are isomorphisms.
  Theorem~\sref{III.4.1.5} states, in particular, that the formation of
  $\RR^n f_*$ commutes with completion; it may therefore be called the
  first comparison theorem between the ``algebraic'' and the
  ``formal'' theories.
\end{enumerate}
We shall begin by proving the affine form of \sref{III.4.1.5}.
\end{env}

\begin{corollary}[4.1.7]
\label{III.4.1.7}
Under the hypotheses of \sref{III.4.1.5}, suppose in addition that
$Y=\Spec(A)$, where $A$ is Noetherian, and that
$\sh{J}=\widetilde{\mathfrak{J}}$, where $\mathfrak{J}$ is an ideal
of $A$, so that
\[
  \sh{F}_k=\sh{F}/\mathfrak{J}^{k+1}\sh{F}.
\]
The canonical homomorphism
\[
  \vphi_n:
  \widehat{\HH^n(X,\sh{F})}
  \longrightarrow
  \varprojlim_k\HH^n(X,\sh{F}_k)
  \tag{4.1.7.1}
  \label{III.4.1.7.1}
\]
where the first term is the separated completion of
$\HH^n(X,\sh{F})$ for the $\mathfrak{J}$-preadic topology, is an
isomorphism.
For every $n$, the projective system
\[
  \bigl(\HH^n(X,\sh{F}_k)\bigr)_{k\geq0}
\]
satisfies condition \emph{(ML)}, and the canonical homomorphism
\[
  \psi_n:
  \HH^n(\widehat{X},\widehat{\sh{F}})
  \longrightarrow
  \varprojlim_k\HH^n(X,\sh{F}_k)
  \tag{4.1.7.2}
  \label{III.4.1.7.2}
\]

\oldpage[III]{126}

is an isomorphism.
Finally, the filtration on $\HH^n(X,\sh{F})$ defined by the kernels
of the canonical homomorphisms
\[
  \HH^n(X,\sh{F})\longrightarrow\HH^n(X,\sh{F}_k)
\]
is $\mathfrak{J}$-good \sref[0]{0.13.7.7}, and $\vphi_n$ is a
topological isomorphism.\footnote{The following proof, which is
simpler than the original proof, and the supplement concerning the
filtration of $\HH^n(X,\sh{F})$, were communicated to us by
J.-P. Serre.}
\end{corollary}

\begin{proof}
Fix $n\geq0$ throughout this proof and, for simplicity, put
\[
  \mathrm{H}=\HH^n(X,\sh{F}),
  \qquad
  \mathrm{H}_k=\HH^n(X,\sh{F}_k).
  \tag{4.1.7.3}
  \label{III.4.1.7.3}
\]
Put also
\[
  \mathrm{R}_k
  =
  \Ker(\mathrm{H}\longrightarrow\mathrm{H}_k),
  \qquad
  \text{a sub-$A$-module of $\mathrm{H}$.}
  \tag{4.1.7.4}
  \label{III.4.1.7.4}
\]
The cohomology exact sequence
\[
  \HH^n(X,\mathfrak{J}^{k+1}\sh{F})
  \longrightarrow
  \HH^n(X,\sh{F})
  \longrightarrow
  \HH^n(X,\sh{F}_k)
  \xrightarrow{\partial}
  \HH^{n+1}(X,\mathfrak{J}^{k+1}\sh{F})
  \longrightarrow
  \HH^{n+1}(X,\sh{F})
\]
shows that
\[
  \mathrm{R}_k
  =
  \Im\bigl(
    \HH^n(X,\mathfrak{J}^{k+1}\sh{F})
    \longrightarrow
    \HH^n(X,\sh{F})
  \bigr).
\]
We put
\[
  \begin{split}
  \mathrm{Q}_k
  &=
  \Ker\bigl(
    \HH^{n+1}(X,\mathfrak{J}^{k+1}\sh{F})
    \longrightarrow
    \HH^{n+1}(X,\sh{F})
  \bigr)
  \\
  &=
  \Im\bigl(
    \HH^n(X,\sh{F}_k)
    \longrightarrow
    \HH^{n+1}(X,\mathfrak{J}^{k+1}\sh{F})
  \bigr).
  \end{split}
  \tag{4.1.7.5}
  \label{III.4.1.7.5}
\]
We therefore have the exact sequence
\[
  0\longrightarrow\mathrm{R}_k
  \longrightarrow\mathrm{H}
  \longrightarrow\mathrm{H}_k
  \longrightarrow\mathrm{Q}_k
  \longrightarrow0.
  \tag{4.1.7.6}
  \label{III.4.1.7.6}
\]

\begin{env}[4.1.7.7]
\label{III.4.1.7.7}
Let $x\in\mathfrak{J}^m$, where $m\geq0$.
Multiplication by $x$ on $\mathfrak{J}^k\sh{F}$ is a homomorphism
\[
  \mathfrak{J}^k\sh{F}
  \longrightarrow
  \mathfrak{J}^{k+m}\sh{F},
\]
and hence gives a homomorphism
\[
  \mu_{x,m}:
  \HH^n(X,\mathfrak{J}^k\sh{F})
  \longrightarrow
  \HH^n(X,\mathfrak{J}^{k+m}\sh{F}).
  \tag{4.1.7.8}
  \label{III.4.1.7.8}
\]
Let
\[
  S=\bigoplus_{k\geq0}\mathfrak{J}^k
\]
be the graded $A$-algebra.
We know that the multiplications $\mu_{x,m}$ endow
\[
  \mathrm{E}
  =
  \bigoplus_{k\geq0}
  \HH^n(X,\mathfrak{J}^k\sh{F})
\]
with the structure of a graded $S$-module of finite type
\sref{III.3.3.2}; the graded ring $S$ is Noetherian
\sref[II]{II.2.1.5}.
\end{env}

\begin{lemma}[4.1.7.9]
\label{III.4.1.7.9}
The submodules $\mathrm{R}_k$ of $\mathrm{H}$ define a
$\mathfrak{J}$-good filtration on $\mathrm{H}$.
\end{lemma}

\begin{proof}
First let us show that
\[
  \mathfrak{J}^m\mathrm{R}_k
  \subset
  \mathrm{R}_{k+m},
  \tag{4.1.7.10}
  \label{III.4.1.7.10}
\]
where multiplication on
$\mathrm{H}=\HH^n(X,\sh{F})$ by an element
$x\in\mathfrak{J}^m$ is the map $\mu_{x,0}$.
For every $x\in\mathfrak{J}^m$, the diagram
\[
  \xymatrix{
    \mathfrak{J}^{k+1}\sh{F}
      \ar[r]^{x}\ar[d]
    &
    \mathfrak{J}^{k+m+1}\sh{F}
      \ar[d]
    \\
    \sh{F}
      \ar[r]_{x}
    &
    \sh{F}
  }
  \tag{4.1.7.11}
  \label{III.4.1.7.11}
\]

\oldpage[III]{127}

is commutative; its horizontal arrows are multiplication by $x$, and
its vertical arrows are the canonical injections.
Consequently, the corresponding diagram
\[
  \xymatrix{
    \HH^n(X,\mathfrak{J}^{k+1}\sh{F})
      \ar[r]^{\mu_{x,m}}\ar[d]
    &
    \HH^n(X,\mathfrak{J}^{k+m+1}\sh{F})
      \ar[d]
    \\
    \HH^n(X,\sh{F})
      \ar[r]_{\mu_{x,0}}
    &
    \HH^n(X,\sh{F})
  }
\]
is commutative.
The interpretation of $\mathrm{R}_k$ as the image of
\[
  \HH^n(X,\mathfrak{J}^{k+1}\sh{F})
  \longrightarrow
  \HH^n(X,\sh{F})
\]
therefore proves \eqref{III.4.1.7.10}.
It also shows that the graded $S$-module
\[
  \mathrm{R}=\bigoplus_{k\geq0}\mathrm{R}_k
\]
is a quotient of the sub-$S$-module
\[
  \mathrm{M}
  =
  \bigoplus_{k\geq0}
  \HH^n(X,\mathfrak{J}^{k+1}\sh{F})
\]
of $\mathrm{E}$.
The preceding observation thus shows that $\mathrm{R}$ is an
$S$-module of finite type.
This is equivalent to the assertion of the lemma
(Bourbaki, \emph{Alg.\ comm.}, Chap.~III, \S3, no.~1, Th.~1).
\end{proof}

\begin{env}[4.1.7.12]
\label{III.4.1.7.12}
Consider now the graded $S$-module
\[
  \mathrm{N}
  =
  \bigoplus_{k\geq0}
  \HH^{n+1}(X,\mathfrak{J}^{k+1}\sh{F}),
\]
defined as in \eqref{III.4.1.7.8}.
It too is an $S$-module of finite type by \sref{III.3.3.2}.
By \eqref{III.4.1.7.5}, we have
$\mathrm{Q}_k\subset\mathrm{N}_k$ for every $k$, and diagram
\eqref{III.4.1.7.11}, with $n$ replaced by $n+1$, shows that
\[
  S_m\mathrm{Q}_k
  =
  \mathfrak{J}^m\mathrm{Q}_k
  \subset
  \mathrm{Q}_{k+m}.
\]
In other words, $\mathrm{Q}$ is a graded sub-$S$-module of
$\mathrm{N}$ and is therefore of finite type.
\end{env}

\begin{env}[4.1.7.13]
\label{III.4.1.7.13}
Let $\alpha_m$ denote the canonical injection
\[
  \mathfrak{J}^m\longrightarrow A,
\]
which may also be written $S_m\to S_0$.
Since $\mathfrak{J}^{k+1}\sh{F}_k=0$, the $A$-module
$\HH^n(X,\sh{F}_k)$ is annihilated by $\mathfrak{J}^{k+1}$.
As $\mathrm{Q}_k$ is the image of the $A$-homomorphism
\[
  \HH^n(X,\sh{F}_k)
  \longrightarrow
  \HH^{n+1}(X,\mathfrak{J}^{k+1}\sh{F}),
\]
$\mathrm{Q}_k$ is likewise annihilated by $\mathfrak{J}^{k+1}$.
Equivalently, in the $S$-module $\mathrm{Q}$,
\[
  \alpha_{k+1}(S_{k+1})\mathrm{Q}_k=0.
  \tag{4.1.7.14}
  \label{III.4.1.7.14}
\]
Since $\mathrm{Q}$ is an $S$-module of finite type, there exist
integers $k_0$ and $h$ such that
\[
  \mathrm{Q}_{k+h}=S_h\mathrm{Q}_k
  \qquad(k\geq k_0)
\]
\sref[II]{II.2.1.6}[ii].
This relation and \eqref{III.4.1.7.14} imply that there exists an
integer $r>0$ such that
\[
  \alpha_r(S_r)\mathrm{Q}=0.
  \tag{4.1.7.15}
  \label{III.4.1.7.15}
\]
\end{env}

\begin{env}[4.1.7.16]
\label{III.4.1.7.16}
The canonical injection
\[
  \mathfrak{J}^{k+m}\sh{F}
  \longrightarrow
  \mathfrak{J}^k\sh{F}
\]
induces in cohomology an $A$-homomorphism
\[
  \nu_m:
  \HH^{n+1}(X,\mathfrak{J}^{k+m}\sh{F})
  \longrightarrow
  \HH^{n+1}(X,\mathfrak{J}^k\sh{F}).
  \tag{4.1.7.17}
  \label{III.4.1.7.17}
\]
For every $x\in\mathfrak{J}^m$, we plainly have the factorization
\[
  \xymatrix@C=4.5em{
    \HH^{n+1}(X,\mathfrak{J}^k\sh{F})
      \ar[r]^{\mu_{x,m}}
    &
    \HH^{n+1}(X,\mathfrak{J}^{k+m}\sh{F})
      \ar[r]^{\nu_m}
    &
    \HH^{n+1}(X,\mathfrak{J}^k\sh{F}),
  }
  \tag{4.1.7.18}
  \label{III.4.1.7.18}
\]
whose composite is $\mu_{x,0}$.
\end{env}

\oldpage[III]{128}

It follows that, for every sub-$A$-module $\mathrm{P}$ of
$\HH^{n+1}(X,\mathfrak{J}^k\sh{F})$, we have in the $S$-module
$\mathrm{N}$
\[
  \nu_m(S_m\mathrm{P})
  =
  \alpha_m(S_m)\mathrm{P}.
  \tag{4.1.7.19}
  \label{III.4.1.7.19}
\]

\begin{lemma}[4.1.7.20]
\label{III.4.1.7.20}
There exists an integer $m>0$ such that
\[
  \nu_m(\mathrm{Q}_{k+m})=0
\]
for every $k\geq k_0$.
\end{lemma}

\begin{proof}
Take $m$ to be a multiple of $h$ with $m\geq r$.
Since
\[
  \mathrm{Q}_{k+m}=S_m\mathrm{Q}_k
  \qquad(k\geq k_0),
\]
relations \eqref{III.4.1.7.19} and \eqref{III.4.1.7.15} give
\[
  \nu_m(\mathrm{Q}_{k+m})
  =
  \alpha_m(S_m)\mathrm{Q}_k
  \subset
  \alpha_r(S_r)\mathrm{Q}_k
  =
  0.
\]
\end{proof}

\begin{env}[4.1.7.21]
\label{III.4.1.7.21}
The commutative diagram
\[
  \xymatrix@C=2.5em{
    \HH^n(X,\sh{F})
      \ar[r]
    &
    \HH^n(X,\sh{F}_k)
      \ar[r]^-{\partial}
    &
    \HH^{n+1}(X,\mathfrak{J}^{k+1}\sh{F})
      \ar[r]
    &
    \HH^{n+1}(X,\sh{F})
    \\
    \HH^n(X,\sh{F})
      \ar[r]\ar[u]
    &
    \HH^n(X,\sh{F}_{k+m})
      \ar[r]^-{\partial}\ar[u]
    &
    \HH^{n+1}(X,\mathfrak{J}^{k+m+1}\sh{F})
      \ar[r]\ar[u]
    &
    \HH^{n+1}(X,\sh{F})
      \ar[u]
  }
\]
itself comes from the commutative diagram
\[
  \xymatrix{
    0\ar[r]
    &
    \mathfrak{J}^{k+1}\sh{F}\ar[r]
    &
    \sh{F}\ar[r]
    &
    \sh{F}_k\ar[r]
    &
    0
    \\
    0\ar[r]
    &
    \mathfrak{J}^{k+m+1}\sh{F}\ar[r]\ar[u]
    &
    \sh{F}\ar[r]\ar[u]
    &
    \sh{F}_{k+m}\ar[r]\ar[u]
    &
    0.
  }
\]
Here the vertical arrows are the canonical maps.
We therefore obtain the commutative diagram
\[
  \xymatrix@C=2.2em{
    0\ar[r]
    &
    \mathrm{R}_k\ar[r]
    &
    \mathrm{H}\ar[r]
    &
    \mathrm{H}_k\ar[r]
    &
    \mathrm{Q}_k\ar[r]
    &
    0
    \\
    0\ar[r]
    &
    \mathrm{R}_{k+m}\ar[r]\ar[u]
    &
    \mathrm{H}\ar[r]\ar@{=}[u]
    &
    \mathrm{H}_{k+m}\ar[r]\ar[u]
    &
    \mathrm{Q}_{k+m}\ar[r]\ar[u]^{\nu_m}
    &
    0,
  }
\]
whose rows are exact.
For $k\geq k_0$, the last vertical arrow is zero by
\sref{III.4.1.7.20}.
Thus the image of $\mathrm{H}_{k+m}$ in $\mathrm{H}_k$ is contained
in
\[
  \Ker(\mathrm{H}_k\longrightarrow\mathrm{Q}_k)
  =
  \Im(\mathrm{H}\longrightarrow\mathrm{H}_k).
\]
Conversely, commutativity shows that it contains
$\Im(\mathrm{H}\to\mathrm{H}_k)$, and hence it is equal to that
image.
The same is therefore true of the images in $\mathrm{H}_k$ of
$\mathrm{H}_{k'}$ for $k'\geq k+m$.
This proves condition \emph{(ML)} for the projective system
$(\mathrm{H}_k)_{k\geq0}$.

Moreover, for every affine open subset $U$ of $X$,
\[
  \HH^i(U,\sh{F}_k)=0
  \qquad(i>0)
\]
by \sref{III.1.3.1}, and, for $m>0$, the map
\[
  \HH^0(U,\sh{F}_{k+m})
  \longrightarrow
  \HH^0(U,\sh{F}_k)
\]
is surjective \sref[I]{I.1.3.9}.
We may therefore

\oldpage[III]{129}

apply \sref[0]{0.13.3.1}; the canonical homomorphism
\[
  \HH^n(\widehat{X},\widehat{\sh{F}})
  \longrightarrow
  \varprojlim_k\HH^n(X,\sh{F}_k)
\]
is bijective for every $n\geq0$.

Since the projective system
$(\mathrm{H}/\mathrm{R}_k)_{k\geq0}$ is strict, we may pass to the
projective limit in the exact sequences
\[
  0
  \longrightarrow
  \mathrm{H}/\mathrm{R}_k
  \longrightarrow
  \mathrm{H}_k
  \longrightarrow
  \mathrm{Q}_k
  \longrightarrow
  0
  \tag{4.1.7.22}
  \label{III.4.1.7.22}
\]
by \sref[0]{0.13.2.2}.
Since $\nu_m(\mathrm{Q}_{k+m})=0$, we have
\[
  \varprojlim_k\mathrm{Q}_k=0,
\]
and hence a topological isomorphism
\[
  \varprojlim_k(\mathrm{H}/\mathrm{R}_k)
  \isoto
  \varprojlim_k\mathrm{H}_k.
\]
But the filtration $(\mathrm{R}_k)$ of $\mathrm{H}$ is
$\mathfrak{J}$-good, and therefore defines on $\mathrm{H}$ the
$\mathfrak{J}$-preadic topology.
Thus $\varprojlim_k(\mathrm{H}/\mathrm{R}_k)$ is the separated
completion of $\mathrm{H}$ for this topology.
This completes the proof of \sref{III.4.1.7}.
\end{env}
\end{proof}

\begin{env}[4.1.8]
\label{III.4.1.8}
We can now complete the proof of \sref{III.4.1.5}.
For every affine open subset $V$ of $Y$,
\[
  \Gamma\bigl(
    V,\widehat{\RR^n f_*(\sh{F})}
  \bigr)
\]
is the separated completion of
$\Gamma(V,\RR^n f_*(\sh{F}))$ for the
$\mathfrak{J}$-preadic topology, if
$\sh{J}|V=\widetilde{\mathfrak{J}}$, because
$\RR^n f_*(\sh{F})$ is a coherent $\sh{O}_Y$-module
\sref[I]{I.10.8.4}.
Furthermore,
\[
  \Gamma\left(
    V,\varprojlim_k\RR^n f_*(\sh{F}_k)
  \right)
  =
  \varprojlim_k
  \Gamma\bigl(V,\RR^n f_*(\sh{F}_k)\bigr)
\]
by \sref[0]{0.3.2.6}.
It follows from \sref{III.4.1.7} and \sref{III.1.4.11} that
$\vphi_n$ is a topological isomorphism.
Likewise, again by \sref{III.1.4.11}, \sref{III.4.1.7} shows that the
homomorphism $\psi_{n,V}$ of \eqref{III.4.1.3.5} is an isomorphism.
By the definition of
$\RR^n\widehat{f}_*(\widehat{\sh{F}})$, it follows that $\psi_n$ is
an isomorphism.
\end{env}

\begin{corollary}[4.1.9]
\label{III.4.1.9}
Under the hypotheses of \sref{III.4.1.4}, for every affine open
subset $V$ of $Y$, the canonical homomorphism
\[
  \HH^n\bigl(
    \widehat{X}\cap f^{-1}(V),
    \widehat{\sh{F}}
  \bigr)
  \longrightarrow
  \Gamma\bigl(
    \widehat{Y}\cap V,
    \RR^n\widehat{f}_*(\widehat{\sh{F}})
  \bigr)
\]
is bijective.
\end{corollary}

\begin{env}[4.1.10]
\label{III.4.1.10}
\emph{Remark.}
Let $f:X\to Y$ be a morphism of finite type between Noetherian usual
preschemes, and let $\sh{F}$ be a coherent $\sh{O}_X$-module whose
support is proper over $Y$ \sref[II]{II.5.4.10}.
We then know \sref{III.3.2.4} that $\RR^n f_*(\sh{F})$ is a coherent
$\sh{O}_Y$-module for every $n\geq0$.
Moreover, we may always suppose that
\[
  \sh{F}=u_*(\sh{G}),
  \qquad
  \sh{G}=u^*(\sh{F}),
\]
where $\sh{G}$ is a coherent $\sh{O}_Z$-module, $Z$ is a suitable
closed subprescheme of $X$ whose underlying space is
$\Supp(\sh{F})$, and $u:Z\to X$ is the canonical injection
\sref[I]{I.9.3.5}.
If
\[
  \sh{G}_k
  =
  \sh{G}\otimes_{\sh{O}_Y}
  (\sh{O}_Y/\sh{J}^{k+1}),
\]
then
\[
  \sh{G}_k=u^*(\sh{F}_k),
  \qquad
  \RR^n f_*(\sh{F}_k)
  =
  \RR^n(f\circ u)_*(\sh{G}_k)
\]
and
\[
  \RR^n f_*(\sh{F})
  =
  \RR^n(f\circ u)_*(\sh{G})
\]
by \sref{III.1.3.4}; and, taking \sref[I]{I.10.9.5} into account,
\[
  \RR^n\widehat{f}_*(\widehat{\sh{F}})
  =
  \RR^n\widehat{(f\circ u)}_*(\widehat{\sh{G}}).
\]
We may therefore apply \sref{III.4.1.5} to $\sh{G}$ and to the proper
morphism $f\circ u$.
It follows that, under these hypotheses, the results of
\sref{III.4.1.5} remain valid for $\sh{F}$ and $f$.
\end{env}


\subsection{Particular cases and variants}
\label{subsection:III.4.2}

The most useful form of the comparison theorem \sref{III.4.1.5} is
the following.

\begin{proposition}[4.2.1]
\label{III.4.2.1}
Let $Y$ be a locally Noetherian prescheme, let $f:X\to Y$ be a
proper morphism, and let $\sh{F}$ be a coherent $\sh{O}_X$-module.
Then, for every $y\in Y$ and every $p$,
\[
  \bigl(\RR^p f_*(\sh{F})\bigr)_y
\]

\oldpage[III]{130}

is an $\sh{O}_y$-module of finite type and is therefore separated
for the $\mathfrak{m}_y$-preadic topology.
Moreover, there is a canonical topological isomorphism
\[
  \widehat{\bigl(\RR^p f_*(\sh{F})\bigr)_y}
  \isoto
  \varprojlim_k
  \HH^p\left(
    f^{-1}(y),
    \sh{F}\otimes_{\sh{O}_Y}
    (\sh{O}_y/\mathfrak{m}_y^k)
  \right).
  \tag{4.2.1.1}
  \label{III.4.2.1.1}
\]
Here the first term is the completion of
$\bigl(\RR^p f_*(\sh{F})\bigr)_y$ for the
$\mathfrak{m}_y$-preadic topology.
In the second term, for every $k\geq0$, $f^{-1}(y)$ is regarded as
the underlying space of the prescheme
\[
  X\times_Y\Spec(\sh{O}_y/\mathfrak{m}_y^k)
\]
\sref[I]{I.3.6.1}.
\end{proposition}

\begin{proof}
Since $\sh{O}_y$ is a Noetherian local ring and
$\bigl(\RR^p f_*(\sh{F})\bigr)_y$ is an $\sh{O}_y$-module of finite
type \sref{III.3.2.1}, the $\mathfrak{m}_y$-preadic topology on this
module is separated \sref[0]{0.7.3.5}.

When $Y$ is Noetherian and $y$ is closed, the other assertions follow
from \sref{III.4.1.7}, after replacing $Y$ by an affine neighbourhood
of $y$ and taking $Y'=\{y\}$; here we use (G,~II,~4.9.1).

In general, put
\[
  Y_1=\Spec(\sh{O}_y),
  \qquad
  X_1=X\times_Y Y_1,
  \qquad
  \sh{F}_1
  =
  \sh{F}\otimes_{\sh{O}_Y}\sh{O}_{Y_1},
\]
and let
\[
  f_1=f\times 1_{Y_1}:X_1\longrightarrow Y_1.
\]
The prescheme $Y_1$ is Noetherian, the morphism $f_1$ is proper
\sref[II]{II.5.4.2}[iii], and $\sh{F}_1$ is coherent
\sref[0]{0.5.3.11}.
Let $y_1$ be the unique closed point of $Y_1$.
The proposition already proved applies to $f_1$, $\sh{F}_1$, and
$y_1$.
We have
\[
  \sh{O}_{y_1}=\sh{O}_y,
  \qquad
  f_1^{-1}(y_1)=f^{-1}(y)
\]
\sref[I]{I.3.6.5}.
Furthermore, the preschemes
\[
  X\times_Y\Spec(\sh{O}_y/\mathfrak{m}_y^k)
  \quad\text{and}\quad
  X_1\times_{Y_1}\Spec(\sh{O}_y/\mathfrak{m}_y^k)
\]
are canonically identified \sref[I]{I.3.3.9}, as are
\[
  \sh{F}_1
  \otimes_{\sh{O}_{Y_1}}
  (\sh{O}_y/\mathfrak{m}_y^k)
  \quad\text{and}\quad
  \sh{F}
  \otimes_{\sh{O}_Y}
  (\sh{O}_y/\mathfrak{m}_y^k)
\]
\sref[I]{I.9.1.6}.
It remains only to observe that
\[
  \RR^p f_{1*}(\sh{F}_1)
  \simeq
  \RR^p f_*(\sh{F})
  \otimes_{\sh{O}_Y}
  \sh{O}_{Y_1}.
\]
This follows from \sref{III.1.4.15}, since the local morphism
$\Spec(\sh{O}_y)\to Y$ is flat
\sref[0]{0.6.7.1}, \sref[I]{I.2.4.2}.
\end{proof}

The following corollary uses the terminology of dimension theory
(Chapter~IV) and will not be applied before that chapter.

\begin{corollary}[4.2.2]
\label{III.4.2.2}
Let $Y$ be a locally Noetherian prescheme, let $f:X\to Y$ be a
proper morphism, let $y\in Y$, and let $r$ be the dimension of
$f^{-1}(y)$.
Then, for every coherent $\sh{O}_X$-module $\sh{F}$, the sheaves
$\RR^p f_*(\sh{F})$ vanish in a neighbourhood of $y$ for every
$p>r$.
\end{corollary}

\begin{proof}
For every $k$, we have
\[
  \HH^p\left(
    f^{-1}(y),
    \sh{F}\otimes_{\sh{O}_Y}
    (\sh{O}_y/\mathfrak{m}_y^k)
  \right)
  =
  0
\]
by (G,~II,~4.15.2).
It follows from \sref{III.4.2.1} that the separated completion of
$\bigl(\RR^p f_*(\sh{F})\bigr)_y$ for the
$\mathfrak{m}_y$-preadic topology is zero.
Since this topology is separated, we also have
\[
  \bigl(\RR^p f_*(\sh{F})\bigr)_y=0.
\]
The conclusion follows because $\RR^p f_*(\sh{F})$ is coherent
\sref[0]{0.5.2.2}.
\end{proof}

\begin{env}[4.2.3]
\label{III.4.2.3}
The result \sref{III.4.2.1} is used especially when $p=0$.
We obtain the following corollary.
\end{env}

\begin{corollary}[4.2.4]
\label{III.4.2.4}
Under the hypotheses of \sref{III.4.2.1}, there is a canonical
topological isomorphism
\[
  \widehat{\bigl(f_*(\sh{F})\bigr)_y}
  \isoto
  \varprojlim_k
  \Gamma\left(
    f^{-1}(y),
    \sh{F}_y/\mathfrak{m}_y^k\sh{F}_y
  \right).
\]
\end{corollary}


\subsection{Zariski's connectedness theorem}
\label{subsection:III.4.3}
\label{III.4.3}

The results of this section and the next generalize well-known
theorems of Zariski, and may all be deduced from
\sref{III.4.2.4}.  They are consequences of the following theorem.

\begin{theorem}[4.3.1]
\label{III.4.3.1}
\textnormal{(Connectedness theorem).}
Let $Y$ be a locally Noetherian prescheme and let $f:X\to Y$ be a
proper morphism.  Then
\[
  \sh{A}(X)=f_*(\sh{O}_X)
\]
is a coherent $\sh{O}_Y$-algebra.

\oldpage[III]{131}

Let $Y'$ be the finite $Y$-scheme such that
$\sh{A}(Y')=\sh{A}(X)$; it is determined up to a $Y$-isomorphism
\sref[II]{II.1.3.1}, \sref[II]{II.6.1.3}.
If $f'=\sh{A}(e)$ is the $Y$-morphism $X\to Y'$ deduced from the
identity isomorphism
\[
  e:\sh{A}(Y')\isoto\sh{A}(X)
\]
\sref[II]{II.1.2.7}, then $f'$ is proper,
$f'_*(\sh{O}_X)$ is isomorphic to $\sh{O}_{Y'}$, and every fibre
$f'^{-1}(y')$, $y'\in Y'$, is connected and nonempty.
\end{theorem}

\begin{proof}
Let $g:Y'\to Y$ be the structural morphism.  To prove that the
homomorphism
\[
  \theta:\sh{O}_{Y'}\longrightarrow f'_*(\sh{O}_X)
\]
entering into the definition of $f'$ is bijective, it is enough,
since $Y'$ is affine over $Y$, to prove that
\[
  g_*(\theta):
  g_*(\sh{O}_{Y'})
  \longrightarrow
  g_*\bigl(f'_*(\sh{O}_X)\bigr)
  =
  f_*(\sh{O}_X)
\]
is the identity \sref[II]{II.1.4.2}.  This follows from the
definitions, since
$g_*(\sh{O}_{Y'})=\sh{A}(Y')$ and
$f_*(\sh{O}_X)=\sh{A}(X)$.
The coherence of $\sh{A}(X)$ is a special case of the finiteness
theorem \sref{III.3.2.1}.  Since $f$ is proper and $g$ is separated,
$f'$ is proper \sref[II]{II.5.4.3}[i].  It remains only to prove
\sref{III.4.3.2}.
\end{proof}

\begin{corollary}[4.3.2]
\label{III.4.3.2}
Under the hypotheses of \sref{III.4.3.1}, suppose in addition that
$f_*(\sh{O}_X)$ is isomorphic to $\sh{O}_Y$.  Then every fibre
$f^{-1}(y)$ of $f$ is connected and nonempty.
\end{corollary}

\begin{proof}
The hypothesis that $f_*(\sh{O}_X)$ is isomorphic to $\sh{O}_Y$
already implies that $f$ is dominant, hence surjective because $f$
is a closed map.  As in \sref{III.4.2.1}, we may reduce to the case
where $y$ is closed in $Y$.  Since $f^{-1}(y)$ is a Noetherian
space, it has finitely many connected components; these are also the
connected components of the underlying space of the completion
$\widehat{X}$ along $f^{-1}(y)$.

If $Z_i$, $1\leq i\leq n$, are these connected components, then
\[
  \Gamma(\widehat{X},\sh{O}_{\widehat{X}})
  =
  \prod_{i=1}^{n}\Gamma(Z_i,\sh{O}_{\widehat{X}}),
\]
and none of the rings on the right is zero, since the unit section is
nonzero at every point of $\widehat{X}$.  Applying
\sref{III.4.1.5} to $\sh{F}=\sh{O}_X$, whose completion along
$f^{-1}(y)$ is $\sh{O}_{\widehat{X}}$, shows that
$\Gamma(\widehat{X},\sh{O}_{\widehat{X}})$ is isomorphic to the
separated $\mathfrak{m}_y$-adic completion
$\widehat{\sh{O}}_y$ of the local ring $\sh{O}_y$.  It is therefore
a local ring and cannot be a direct product of several nonzero rings.
Thus $n=1$.
\end{proof}

\begin{corollary}[4.3.3]
\label{III.4.3.3}
Under the hypotheses of \sref{III.4.3.1}, for every $y\in Y$ the set
of connected components of the fibre $f^{-1}(y)$ is in bijective
correspondence with the finite set of points of the fibre
$g^{-1}(y)$, where $g:Y'\to Y$ is the structural morphism; in other
words, it corresponds to the set of maximal ideals of
$\bigl(f_*(\sh{O}_X)\bigr)_y$.
\end{corollary}

\begin{proof}
Since $Y'$ is finite over $Y$, the space $g^{-1}(y)$ is finite and
discrete \sref[II]{II.6.1.7}.  As
\[
  f^{-1}(y)=f'^{-1}\bigl(g^{-1}(y)\bigr),
\]
the assertion follows from this observation and
\sref{III.4.3.1}.
\end{proof}

This gives a remarkable interpretation of the $Y$-prescheme $Y'$
defined in \sref{III.4.3.1}.  The factorization
$f=g\circ f'$ of the proper morphism $f$ is analogous to the
factorization obtained by K.~Stein for holomorphic maps of analytic
spaces, and will henceforth be called the \emph{Stein factorization}
of $f$.

\begin{remark}[4.3.4]
\label{III.4.3.4}
Let $k$ be an extension of the field $\kres(y)$.  If
\[
  f^{-1}(y)\otimes_{\kres(y)}k
  =
  X\times_Y\Spec(k)
\]
is connected, then $f^{-1}(y)$ is connected as well, since it is the
image of the former under a projection morphism
\sref[I]{I.3.4.7}.  For a morphism $f:X\to Y$ of preschemes and a
point $y\in Y$, we shall say that the fibre $f^{-1}(y)$ is
\emph{geometrically connected} if, for every extension $k$ of
$\kres(y)$, the prescheme
$f^{-1}(y)\otimes_{\kres(y)}k=X\times_Y\Spec(k)$ is connected.

\oldpage[III]{132}

Under the hypotheses of \sref{III.4.3.2}, the conclusion can be
strengthened: the fibres $f^{-1}(y)$ are geometrically connected.
Indeed, for every extension $k$ of $\kres(y)$ there exist a local
Noetherian ring $A$ and a local homomorphism
\[
  \varphi:\sh{O}_y\longrightarrow A
\]
such that $A$ is a flat $\sh{O}_y$-module and the residue field of
$A$ is $\kres(y)$-isomorphic to $k$ \sref[0]{0.10.3.1}.
Let $Y_1=\Spec(A)$, and let $h:Y_1\to Y$ be the local morphism
corresponding to $\varphi$, taking the unique closed point $y_1$ of
$Y_1$ to $y$ \sref[I]{I.2.4.1}.  Set
\[
  X_1=X\times_Y Y_1,
  \qquad
  f_1=f\times 1_{Y_1}.
\]
The prescheme $Y_1$ is Noetherian, $f_1$ is proper
\sref[II]{II.5.4.2}, and $f_1^{-1}(y_1)$ is a
$\kres(y_1)$-prescheme isomorphic to $X\times_Y\Spec(k)$.  It is
therefore enough to prove that
\[
  f_{1*}(\sh{O}_{X_1})=\sh{O}_{Y_1},
\]
so that \sref{III.4.3.2} may be applied to $f_1$.  Now $h$ is flat
by \sref[I]{I.2.4.2}; hence, by \sref{III.1.4.15} with $q=0$,
\[
  f_{1*}(\sh{O}_{X_1})
  =
  h^*\bigl(f_*(\sh{O}_X)\bigr)
  =
  h^*(\sh{O}_Y)
  =
  \sh{O}_{Y_1}.
\]

In the general case of \sref{III.4.3.1}, the same argument shows,
with the notation of that theorem, that
\[
  f_{1*}(\sh{O}_{X_1})
  =
  h^*\bigl(g_*(\sh{O}_{Y'})\bigr),
\]
and that the Stein factorization $f_1=g_1\circ f'_1$ of $f_1$
satisfies
\[
  g_1=g\times 1_{Y_1}
\]
\sref[II]{II.1.5.2}, the corresponding finite $Y_1$-scheme being
$Y'_1=Y'\times_Y Y_1$.  Taking the transitivity of fibres into
account \sref[I]{I.3.6.4}, it follows from
\sref{III.4.3.3} that the number of connected components of
$f_1^{-1}(y_1)$ is the number of elements of
\[
  g_1^{-1}(y_1)
  =
  g^{-1}(y)\otimes_{\kres(y)}k.
\]
If $k$ is chosen algebraically closed, this number is independent of
the algebraically closed extension chosen and is equal to the
\emph{geometric number of points} of $g^{-1}(y)$
\sref[I]{I.6.4.7}, or equivalently to
\[
  \sum_i[\kres(y'_i):\kres(y)]_s,
\]
where the $y'_i$ range over the finite set $g^{-1}(y)$.  This number
is called the \emph{geometric number of connected components} of
$f^{-1}(y)$.  The fields $\kres(y'_i)$ are precisely the residue
fields of the semilocal ring $\bigl(f_*(\sh{O}_X)\bigr)_y$.
\end{remark}

\begin{proposition}[4.3.5]
\label{III.4.3.5}
Let $X$ and $Y$ be locally Noetherian integral preschemes, and let
$f:X\to Y$ be a proper dominant morphism.  For every $y\in Y$, the
number of connected components of $f^{-1}(y)$ is at most the number
of maximal ideals of the integral closure $\sh{O}'_y$ of
$\sh{O}_y$ in the field of rational functions $\mathrm{R}(X)$.
\end{proposition}

\begin{proof}
For every open set $U$ of $Y$,
\[
  \Gamma\bigl(U,f_*(\sh{O}_X)\bigr)
  =
  \Gamma\bigl(f^{-1}(U),\sh{O}_X\bigr)
\]
is the intersection of the local rings $\sh{O}_x$ for
$x\in f^{-1}(U)$ \sref[I]{I.8.2.1.1}.  It follows at once that
$\bigl(f_*(\sh{O}_X)\bigr)_y$ is a subring of $\mathrm{R}(X)$ containing
$\sh{O}_y$.  Since $f_*(\sh{O}_X)$ is a coherent
$\sh{O}_Y$-module, $\bigl(f_*(\sh{O}_X)\bigr)_y$ is a finite
$\sh{O}_y$-module, and hence is contained in $\sh{O}'_y$.
Every maximal ideal of such a ring is the intersection of that ring
with a maximal ideal of $\sh{O}'_y$
\cite[vol.~I, pp.~257 and 259]{I-13}, proving the proposition.
\end{proof}

\begin{definition}[4.3.6]
\label{III.4.3.6}
An integral local ring is said to be \emph{unibranch} if its integral
closure is again a local ring.  A point $y$ of an integral prescheme
$Y$ is said to be \emph{unibranch} if the local ring $\sh{O}_y$ is
unibranch; this is in particular the case when $Y$ is normal at
$y$.
\end{definition}

Let $A$ be an integral local ring and let $K$ be its field of
fractions.  The ring $A$ is unibranch if and only if every subring
$A_1$ of $K$ that contains $A$ and is a finite $A$-algebra is
local.  Indeed, let $A'$ be the integral closure of $A$.  By the
first Cohen--Seidenberg theorem (Bourbaki, \emph{Alg\`ebre
commutative}, Chapter~V, \S2, no.~1, Theorem~1), every maximal ideal
of $A_1$ is the contraction to $A_1$ of a maximal ideal of $A'$;
thus, if $A'$ is local, so is $A_1$.

\oldpage[III]{133}

Conversely, $A'$ is the inductive limit of the directed family of
finite $A$-subalgebras $A_\alpha$ of $A'$.  If each $A_\alpha$ is
local, then for $A_\alpha\subset A_\beta$ the maximal ideal of
$A_\alpha$ is the contraction of the maximal ideal of $A_\beta$, by
the same argument; hence $A'$ is local
\sref[0]{0.10.3.1.3}.

Notice also that if the completion of a Noetherian local ring $A$ is
integral---a property expressed by saying that $A$ is
\emph{analytically integral}---then $A$ is unibranch.  Let
$\mathfrak{m}$ be the maximal ideal of $A$, let $K$ be its field of
fractions, and let $K'$ be the field of fractions of $\widehat{A}$;
then
\[
  K'=K\otimes_A\widehat{A}.
\]
Let $A'_F$ be a finite $A$-subalgebra of $K$.  The subring $B_F$ of
$K'$ generated by $\widehat{A}$ and $A'_F$ is isomorphic to
\[
  A'_F\otimes_A\widehat{A};
\]
it is a finite $\widehat{A}$-module and is the completion of $A'_F$
for the $\mathfrak{m}$-adic topology
\sref[0]{0.7.3.3}, \sref[0]{0.7.3.6}.  Since $A'_F$ is semilocal
(Bourbaki, \emph{Alg\`ebre commutative}, Chapter~IV, \S2, no.~5,
Corollary~3 to Proposition~9), while its completion is integral,
$A'_F$ can have only one maximal ideal $\mathfrak{m}'_F$, and
$\mathfrak{m}'_F\cap A=\mathfrak{m}$.  This proves the assertion.

\begin{corollary}[4.3.7]
\label{III.4.3.7}
Under the hypotheses of \sref{III.4.3.5}, suppose that the algebraic
closure of $\mathrm{R}(Y)$ in $\mathrm{R}(X)$ has separable degree $n$ and that
$y\in Y$ is unibranch.  Then the fibre $f^{-1}(y)$ has at most $n$
connected components.  In particular, if the algebraic closure of
$\mathrm{R}(Y)$ in $\mathrm{R}(X)$ is radicial over $\mathrm{R}(Y)$, then $f^{-1}(y)$ is
connected.
\end{corollary}

\begin{proof}
Let $\sh{O}'_y$ be the integral closure of $\sh{O}_y$, and let
$\sh{O}''_y$ be the integral closure of $\sh{O}_y$ in $\mathrm{R}(X)$.
The latter is also the integral closure of $\sh{O}'_y$ in $\mathrm{R}(X)$.
If $\sh{O}'_y$ is local, then $\sh{O}''_y$ is a semilocal ring
having at most $n$ maximal ideals
\cite[vol.~I, p.~289, Theorem~22]{I-13}.

This corollary is essentially the form in which Zariski states his
``connectedness theorem'' for algebraic schemes.
\end{proof}

\begin{remark}[4.3.8]
\label{III.4.3.8}
If, in addition to the hypotheses of \sref{III.4.3.7}, $Y$ is
normal at $y$, then the fibre $f^{-1}(y)$ is geometrically
connected: with the notation of \sref{III.4.3.4}, $g^{-1}(y)$ is a
single point $y'$, and $\kres(y')$ is radicial over $\kres(y)$.
\end{remark}

\begin{definition}[4.3.9]
\label{III.4.3.9}
Let $Y$ be a locally Noetherian prescheme.  A morphism of finite type
$f:X\to Y$ is said to be \emph{universally open} if, for every
irreducible locally Noetherian prescheme $Y'$ and every dominant
morphism $g:Y'\to Y$, every irreducible component of
$X'=X\times_Y Y'$ dominates $Y'$.
\end{definition}

If $Y$ is irreducible, this says that if $\eta$ and $\eta'$ are the
generic points of $Y$ and $Y'$, respectively, and if
$f'=f\times 1_{Y'}$, every irreducible component of $X'$ has its
generic point in $f'^{-1}(\eta')$.  Consequently, for every open set
$U$ of $Y$, the restricted morphism $f^{-1}(U)\to U$ is universally
open.

\begin{corollary}[4.3.10]
\label{III.4.3.10}
Let $X$ and $Y$ be locally Noetherian integral preschemes, and let
$f:X\to Y$ be proper, dominant, and universally open.  If the
algebraic closure of $\mathrm{R}(Y)$ in $\mathrm{R}(X)$ is radicial over
$\mathrm{R}(Y)$, then every fibre $f^{-1}(y)$, $y\in Y$, is geometrically
connected.
\end{corollary}

\begin{proof}
We may reduce to the case $Y=\Spec(B)$, where $B$ is an integral
Noetherian ring.  By \sref[II]{II.7.1.7}, there exists an integrally
closed Noetherian local ring $A$ that dominates $\sh{O}_y$ and has
$\mathrm{R}(Y)$ as its field of fractions.  Put $Y'=\Spec(A)$, and let
$h:Y'\to Y$ be the morphism corresponding to the canonical
injection $B\to A$.  It is birational, hence dominant; moreover, if
$y'$ is the unique closed point of $Y'$, then $h(y')=y$.

Set
\[
  X'=X\times_Y Y',
  \qquad
  f'=f\times 1_{Y'}.
\]
Let $\eta$, $\eta'$, and $\xi$ be the generic points of $Y$, $Y'$,
and $X$, respectively.  Then $f(\xi)=\eta$ and
$h^{-1}(\eta)=\{\eta'\}$; moreover
$\kres(\eta)=\kres(\eta')=\mathrm{R}(Y)$.  Thus $f'^{-1}(\eta')$ is
isomorphic to $f^{-1}(\eta)$ \sref[I]{I.3.6.4}, and it has a single
generic point, since $\xi$ is the generic point of
$f^{-1}(\eta)$ \sref[0]{0.2.1.8}.  By universal openness, every
irreducible component of $X'$ has its generic point in
$f'^{-1}(\eta')$.  Hence $X'$ is irreducible; its generic point
$\xi'$ is that of $f'^{-1}(\eta')$, and
$\kres(\xi')=\kres(\xi)$.

\oldpage[III]{134}

Put $X''=X'_{\mathrm{red}}$ and $f''=f'_{\mathrm{red}}$.  Then
$X''$ is integral and Noetherian, $f''$ is proper
\sref[II]{II.5.4.6}, and the underlying spaces of the fibres
$f'^{-1}(y')$ and $f''^{-1}(y')$ are the same.  Moreover,
\[
  \mathrm{R}(X'')=\kres(\xi')=\mathrm{R}(X).
\]
Thus $f''$ satisfies the hypotheses of \sref{III.4.3.8}, and
$f''^{-1}(y')$ is geometrically connected.

Let $k$ now be any extension of $\kres(y)$.  There exists an
extension $k_1$ of $\kres(y)$ in which both $\kres(y')$ and $k$ may
be regarded as subextensions (Bourbaki, \emph{Alg\`ebre},
Chapter~V, \S4, Proposition~2).  By the preceding paragraph,
\[
  f''^{-1}(y')\times_{Y'}\Spec(k_1)
\]
is connected.  It has the same associated reduced prescheme as
$f'^{-1}(y')\times_{Y'}\Spec(k_1)$
\sref[I]{I.5.1.8}, so the latter is connected; it is isomorphic to
$f^{-1}(y)\times_Y\Spec(k_1)$ \sref[I]{I.3.6.4}.  Therefore
$f^{-1}(y)\times_Y\Spec(k)$ is connected, and the observation at the
beginning of \sref{III.4.3.4} completes the proof.
\end{proof}

\begin{remark}[4.3.11]
\label{III.4.3.11}
\emph{(i)}
The preceding argument is due in substance to Zariski
\cite{I-20}, except that in his setting one may take for $A$ the
integral closure of $\sh{O}_y$, since this is a Noetherian ring for
the local rings of classical algebraic geometry.  Zariski also
proves that if $Y$ is the Chow variety of a projective space
$\mathbf{P}_k^r$ over a field $k$, and $X$ is the closed subset of
$\mathbf{P}_k^r\times_kY$ defining the Chow correspondence between
$\mathbf{P}_k^r$ and $Y$, then the projection $X\to Y$ is
universally open (\emph{loc.\ cit.}, lemma on p.~82).
It appears that this is the only formal property of ``Chow
coordinates'' used in certain applications.  In such a situation it
is therefore useful to replace the language of specialization of
cycles in projective space by that of fibres of a proper morphism,
possibly assumed universally open or subject to analogous local
regularity conditions.

\emph{(ii)}
In Chapter~IV we shall see that a universally open morphism
$f:X\to Y$ may equivalently be defined by requiring that, for every
morphism $Z\to Y$, the base-changed morphism
$f_{(Z)}:X_{(Z)}\to Z$ be open.  One can also show that if $f$
satisfies the hypotheses of \sref{III.4.3.10}, and if $y$ is a
specialization of $y'$, then the geometric number of connected
components of $f^{-1}(y)$ is at most that of $f^{-1}(y')$.
\end{remark}

\begin{corollary}[4.3.12]
\label{III.4.3.12}
Under the hypotheses of \sref{III.4.3.5}, suppose in addition that
$\mathrm{R}(Y)$ is algebraically closed in $\mathrm{R}(X)$, and let $y$ be a
normal point of $Y$.  Then $f^{-1}(y)$ is geometrically connected,
and there exists an open neighbourhood $U$ of $y$ in $Y$ such that
\[
  f_*\bigl(\sh{O}_X|_{f^{-1}(U)}\bigr)
  \simeq
  \sh{O}_Y|_U.
\]
In particular, if $Y$ is normal and $\mathrm{R}(Y)$ is algebraically closed
in $\mathrm{R}(X)$, then $f_*(\sh{O}_X)\simeq\sh{O}_Y$.
\end{corollary}

\begin{proof}
The assertion concerning $f^{-1}(y)$ is a special case of
\sref{III.4.3.8}.

\oldpage[III]{135}

If
\[
  X\xrightarrow{f'}Y'\xrightarrow{g}Y
\]
is the Stein factorization of $f$ \sref{III.4.3.3}, it follows that
$g^{-1}(y)$ consists of a single point $y'$.  Moreover,
\[
  \sh{O}_y
  \subset
  \sh{O}_{y'}
  =
  \bigl(f_*(\sh{O}_X)\bigr)_y
  \subset
  \mathrm{R}(X).
\]
Since $\sh{O}_{y'}$ is finite over $\sh{O}_y$, and hence
\emph{a fortiori} algebraic over $\mathrm{R}(Y)$ inside the relevant
fraction field, the hypothesis forces it to lie in $\mathrm{R}(Y)$.  Since
$y$ is normal, necessarily $\sh{O}_{y'}=\sh{O}_y$.  Thus $g$ is a
local isomorphism at $y'$ \sref[I]{I.6.5.4}, which proves the first
part.  The second assertion follows from the first: the additional
hypothesis makes $g$ bijective and a local isomorphism in a
neighbourhood of every point of $Y'$, hence an isomorphism.
\end{proof}

The fact that \sref{III.4.3.7} is available for schemes yields, for
example, the following application.

\begin{proposition}[4.3.13]
\label{III.4.3.13}
Let $A$ be a unibranch Noetherian local ring, let $\mathfrak{a}$ be
an ideal of definition of $A$, let
\[
  A_0=A/\mathfrak{a},
  \qquad
  S=\operatorname{gr}_{\mathfrak{a}}(A)
\]
be the graded ring associated to the $\mathfrak{a}$-preadic
filtration.  Suppose that $S$ is a graded $A_0$-algebra generated by
$S_1$, where $S_1$ is a finite $A_0$-module.  Then $\Proj(S)$ is a
connected $A_0$-scheme.
\end{proposition}

\begin{proof}
Let $\mathfrak{m}$ be the maximal ideal of $A$.  The scheme
$Y=\Spec(A)$ is integral, and the point $y$ corresponding to
$\mathfrak{m}$ is its unique closed point.  For some integer $p$,
\[
  \mathfrak{m}^p\subset\mathfrak{a}\subset\mathfrak{m},
\]
so $V(\mathfrak{a})=\{\mathfrak{m}\}$.  Put
\[
  S'=\bigoplus_{n\geq0}\mathfrak{a}^n
  \qquad\text{and}\qquad
  X=\Proj(S').
\]
This is the $Y$-scheme obtained by blowing up the ideal
$\mathfrak{a}$.  It is integral, and the structural morphism
$f:X\to Y$ is birational \sref[II]{II.8.1.4} and projective.
Therefore \sref{III.4.3.7} applies and shows that $f^{-1}(y)$ is
connected.  But the underlying space of $f^{-1}(y)$ is that of
\[
  \Proj(S'\otimes_A A_0)
\]
\sref[I]{I.3.6.1}, \sref[II]{II.2.8.10}; since
$S'\otimes_AA_0=S$ by definition, the proposition follows.
\end{proof}


\subsection{Zariski's ``main theorem''}
\label{subsection:III.4.4}

\begin{proposition}[4.4.1]
\label{III.4.4.1}
Let $Y$ be a locally Noetherian prescheme and let $f:X\to Y$ be a
proper morphism.  Let $X'$ be the set of points $x\in X$ that are
isolated in their fibre $f^{-1}(f(x))$.  Then $X'$ is open in $X$.
If
\[
  f=g\circ f'
\]
is the Stein factorization of $f$ \sref{III.4.3.3}, the restriction
of $f'$ to $X'$ is an isomorphism of $X'$ onto the subprescheme
induced on an open subset $U$ of $Y'$, and
$X'=f'^{-1}(U)$.
\end{proposition}

\begin{proof}
Since $g^{-1}(f(x))$ is finite and discrete
\sref{III.4.3.3}, \sref[II]{II.6.1.7}, the point $x$ is isolated in
$f^{-1}(f(x))$ if and only if it is isolated in
$f'^{-1}(f'(x))$.  We may therefore reduce to the case $f'=f$, so
that $f_*(\sh{O}_X)=\sh{O}_Y$.  If $x\in X'$, the fibre
$f^{-1}(f(x))$, being connected by \sref{III.4.3.2}, must then
consist of the single point $x$.

Since $f$ is closed, for every open neighbourhood $V$ of $x$ in
$X$, the closed set $f(X-V)$ does not contain $y=f(x)$.  If $U$ is
its complement in $Y$, then $f^{-1}(U)\subset V$.  Thus the inverse
images under $f$ of a fundamental system of open neighbourhoods of
$y$ form a fundamental system of open neighbourhoods of $x$.
The equality $f_*(\sh{O}_X)=\sh{O}_Y$ and the definition of the
direct image of a sheaf \sref[0]{0.3.4.1}, \sref{III.4.2.1} now
imply that, if $f=(\psi,\theta)$, the homomorphism
\[
  \theta_y^\#:\sh{O}_y\longrightarrow\sh{O}_x
\]
is an isomorphism.  Hence there are open neighbourhoods $V$ of $x$
and $U$ of $y$ such that the restriction of $f$ is an isomorphism
$V\isoto U$ \sref[I]{I.6.5.4}.

\oldpage[III]{136}

By the preceding paragraph we may moreover arrange that
$f^{-1}(U)=V$.  It follows immediately that $V\subset X'$, which
completes the proof.
\end{proof}

The following proposition was proved by Chevalley for algebraic
schemes.

\begin{proposition}[4.4.2]
\label{III.4.4.2}
Let $Y$ be a locally Noetherian prescheme and let $f:X\to Y$ be a
morphism.  The following conditions are equivalent:
\begin{enumerate}
  \item[{\rm a)}] $f$ is finite;
  \item[{\rm b)}] $f$ is affine and proper;
  \item[{\rm c)}] $f$ is proper and, for every $y\in Y$, the set
    $f^{-1}(y)$ is finite.
\end{enumerate}
\end{proposition}

\begin{proof}
Condition~a) implies b) by \sref[II]{II.6.1.2} and
\sref[II]{II.6.1.11}.  If $f$ is proper and affine, then so is
\[
  f^{-1}(y)\longrightarrow\Spec(\kres(y))
\]
\sref[II]{II.1.6.2}, \sref[II]{II.5.4.2}.  The finiteness theorem
\sref{III.3.2.1}, applied to the structural sheaf of $f^{-1}(y)$,
shows that $f^{-1}(y)=\Spec(A)$ for a finite
$\kres(y)$-algebra $A$.  Thus $f^{-1}(y)$ is a finite set
\sref[II]{II.6.1.7}, and b) implies c).

Finally, if c) holds, the fibre $f^{-1}(y)$ is an algebraic
prescheme over $\kres(y)$ whose underlying set is finite, and is
therefore discrete \sref[I]{I.6.4.4}.  With the notation of
\sref{III.4.4.1}, we have $X'=X$, so
$f':X\to Y'$ is an isomorphism.  Since $g$ is finite, c) implies
a).
\end{proof}

\begin{theorem}[4.4.3]
\label{III.4.4.3}
\textnormal{(Zariski's ``main theorem'').}
Let $Y$ be a Noetherian prescheme, let $f:X\to Y$ be a
quasi-projective morphism, and let $X'$ be the set of points $x\in X$
that are isolated in their fibre $f^{-1}(f(x))$.  Then $X'$ is open
in $X$, and the subprescheme induced by $X$ on $X'$ is isomorphic to
a subprescheme induced on an open subset of a $Y$-prescheme $Y'$
finite over $Y$.
\end{theorem}

\begin{proof}
There is a projective $Y$-prescheme $Z$ such that $X$ is
$Y$-isomorphic to the subprescheme induced on an open subset of $Z$
\sref[II]{II.5.3.2}, \sref[II]{II.5.5.1}.  We are thus reduced to
the case where $f$ is projective, hence proper
\sref[II]{II.5.5.3}, and the theorem follows immediately from
\sref{III.4.4.1}.
\end{proof}

\begin{remark}[4.4.4]
\label{III.4.4.4}
If $X$ is reduced (respectively, if $X$ is irreducible and $X'$ is
nonempty), one may require in \sref{III.4.4.3} that $Y'$ be reduced
(respectively, irreducible).  Indeed, $Y'$ may be replaced by the
schematic closure $\overline{X'}$ of $X'$ in $Y'$
\sref[I]{I.9.5.11}, \sref[II]{II.6.1.5}.  If $X'$ is reduced, so is
$\overline{X'}$ \sref[I]{I.9.5.9}; if $X'$ is nonempty and
irreducible, then $\overline{X'}$ is irreducible as well.
\end{remark}

\begin{corollary}[4.4.5]
\label{III.4.4.5}
Let $Y$ be a locally Noetherian scheme, let $f:X\to Y$ be a morphism
of finite type, and let $x$ be a point of $X$ isolated in its fibre
$f^{-1}(f(x))$.  Then there is an open neighbourhood of $x$ in $X$
which is isomorphic to an open subset of a $Y$-prescheme finite over
$Y$.
\end{corollary}

\begin{proof}
Set $y=f(x)$, let $U$ be an affine open neighbourhood of $y$ in $Y$,
and let $V$ be an affine open neighbourhood of $x$ in $X$, contained
in $f^{-1}(U)$.  Since $Y$ is separated, the inclusion $U\to Y$ is
affine \sref[II]{II.1.6.3}; and since $V$ is affine over $U$
(\emph{ibid.}), the restriction of $f$ to $V$ is an affine morphism
$V\to Y$ \sref[II]{II.1.6.2}.  A fortiori, this restriction is
quasi-projective, since it is of finite type
\sref[I]{I.6.3.5}, \sref[II]{II.5.3.4}.  It remains only to apply
\sref{III.4.4.3} to this restriction.
\end{proof}

Corollary \sref{III.4.4.5} may be stated in the language of
commutative algebra as follows.

\begin{corollary}[4.4.6]
\label{III.4.4.6}
Let $A$ be a Noetherian ring, let $B$ be an $A$-algebra of finite
type, let $\mathfrak q$ be a prime ideal of $B$, and let $\mathfrak p$
be its inverse image in $A$.  Suppose that $\mathfrak q$ is both
maximal and minimal in the set of prime ideals of $B$ whose inverse
image is $\mathfrak p$.  Then there exist an element
$g\in B-\mathfrak q$, a finite $A$-algebra $A'$, and an element
$f'\in A'$ such that the $A$-algebras $B_g$ and $A'_{f'}$ are
isomorphic.
\end{corollary}

\begin{proof}
Apply \sref{III.4.4.5} to $Y=\Spec(A)$ and $X=\Spec(B)$.  The
hypothesis on $\mathfrak q$ says precisely that $\mathfrak q$ is
isolated in its fibre \sref[I]{I.1.1.7}.
\end{proof}

The following consequence appears less general.

\oldpage[III]{137}

\begin{corollary}[4.4.7]
\label{III.4.4.7}
Let $A$ be a Noetherian local ring, let $B$ be an $A$-algebra of
finite type, and let $\mathfrak n$ be a prime ideal of $B$ whose
inverse image in $A$ is the maximal ideal $\mathfrak m$.  Suppose
that $\mathfrak n$ is maximal in $B$ and minimal in the set of prime
ideals of $B$ whose inverse image is $\mathfrak m$ (which also means
that $\mathfrak mB$ is $\mathfrak n$-primary).  Then there exist a
finite $A$-algebra $A'$ and a maximal ideal $\mathfrak m'$ of $A'$
(whose inverse image in $A$ is $\mathfrak m$) such that $B_{\mathfrak
n}$ is isomorphic, as an $A$-algebra, to $A'_{\mathfrak m'}$.
\end{corollary}

The following special case of \sref{III.4.4.7} is also sometimes
called the ``main theorem''.

\begin{corollary}[4.4.8]
\label{III.4.4.8}
Under the hypotheses of \sref{III.4.4.7}, suppose moreover that $A$
and $B$ are integral domains with the same field of fractions $K$.
If $A$ is integrally closed, then $A=B$.
\end{corollary}

\begin{proof}
Remark \sref{III.4.4.4} shows that, in applying
\sref{III.4.4.7}, we may suppose that $A'$ is an integral domain
with field of fractions $K$.  The hypothesis on $A$ then gives
$A'=A$, hence $B_{\mathfrak n}=A$.  Since
$A\subset B\subset B_{\mathfrak n}$, it follows that $A=B$.
\end{proof}

Statement \sref{III.4.4.8} is the form in which Zariski gave his
``main theorem'' (extended here to arbitrary Noetherian integral
local rings).

The preceding corollaries were local variants of
\sref{III.4.4.3}, which is a global result.  Here is another
consequence of a global nature.

\begin{corollary}[4.4.9]
\label{III.4.4.9}
Let $X$ and $Y$ be integral locally Noetherian preschemes, and let
$f:X\to Y$ be a separated morphism locally of finite type which is
birational.  Suppose moreover that $Y$ is normal and that every
fibre $f^{-1}(y)$, for $y\in Y$, is finite.  Then $f$ is an open
immersion.  If in addition $f$ is closed (in particular, if $f$ is
proper), then $f$ is an isomorphism.
\end{corollary}

\begin{proof}
Let $x\in X$ and put $y=f(x)$.  Since $f^{-1}(y)$ is a scheme
locally of finite type over $\kres(y)$, the assumption that it is finite implies
that it is discrete \sref[I]{I.6.4.4}.  Moreover, $\sh{O}_y$ is
integrally closed, and $\sh{O}_x$ and $\sh{O}_y$ have the same field
of fractions \sref[I]{I.7.1.5}.  We may therefore apply
\sref{III.4.4.8}; if $f=(\psi,\theta)$, the homomorphism
\[
  \theta_y^\#:\sh{O}_y\longrightarrow\sh{O}_x
\]
is bijective.  It follows from \sref[I]{I.6.5.4} that $f$ is a local
isomorphism.  Since $f$ is separated and $X$ is integral, $f$ is an
open immersion \sref[I]{I.8.2.8}.  The last assertion follows from
the fact that $f$ is dominant.
\end{proof}

\begin{proposition}[4.4.10]
\label{III.4.4.10}
Let $Y$ be a locally Noetherian prescheme, and let $f:X\to Y$ be a
morphism locally of finite type.  The set $X'$ of points $x\in X$
which are isolated in their fibre $f^{-1}(f(x))$ is open in $X$.
\end{proposition}

\begin{proof}
The question is local on $X$ and $Y$, so we may suppose that $X$ and
$Y$ are affine Noetherian and that $f$ is of finite type.  The
morphism $f$ is then affine and of finite type, hence quasi-projective
\sref[II]{II.5.3.4}; it suffices to apply \sref{III.4.4.3}.
\end{proof}

\oldpage[III]{138}

\begin{corollary}[4.4.11]
\label{III.4.4.11}
Let $Y$ be a locally Noetherian prescheme, and let $f:X\to Y$ be a
proper morphism.  The set $U$ of points $y\in Y$ for which
$f^{-1}(y)$ is discrete is open in $Y$, and the restricted morphism
$f^{-1}(U)\to U$ is finite.  In particular, every proper
quasi-finite morphism $X\to Y$ is finite.
\end{corollary}

\begin{proof}
The complement of $U$ in $Y$ is the image under $f$ of $X-X'$, which
is closed in $X$ by \sref{III.4.4.10}; since $f$ is a closed map,
$U$ is open.  Furthermore, \sref[II]{II.6.2.2} shows that
$f^{-1}(y)$ is finite for every $y\in U$.  Since the restricted
morphism $f^{-1}(U)\to U$ is proper \sref[II]{II.5.4.1}, it is finite
by \sref{III.4.4.2}.
\end{proof}

\begin{remark}[4.4.12]
\label{III.4.4.12}
\emph{(i)}
As announced in \sref[II]{II.6.2.7}, we shall show in Chapter~V that
if $Y$ is locally Noetherian, every separated quasi-finite morphism
$f:X\to Y$ is quasi-affine, hence quasi-projective.  It will follow
that the conclusion of the main theorem \sref{III.4.4.3} remains
valid when $f$ is assumed only to be separated and of finite type.
Indeed, \sref{III.4.4.10} shows that $X'$ is open in $X$; and since
$X$ is locally Noetherian, the restriction of $f$ to $X'$ is still
of finite type \sref[I]{I.6.3.5}, hence quasi-finite by the definition
of $X'$, and clearly separated.  We may therefore apply
\sref{III.4.4.3} to this restriction.

\emph{(ii)}
In Chapter~IV we shall give a more elementary proof of
\sref{III.4.4.10}, using dimension theory.
\end{remark}


\subsection{Completions of modules of homomorphisms}
\label{subsection:III.4.5}

\begin{proposition}[4.5.1]
\label{III.4.5.1}
Let $A$ be a Noetherian ring, let $\mathfrak J$ be an ideal of $A$,
let $X$ be an $A$-prescheme of finite type, and let $\sh{F}$ and
$\sh{G}$ be two coherent $\sh{O}_X$-modules whose supports have
proper intersection over $Y=\Spec(A)$ \sref[II]{II.5.4.10}.  Then,
for every integer $n\geq0$,
\[
  \Ext^n_{\sh{O}_X}(X;\sh{F},\sh{G})
\]
is a finite $A$-module, and its separated completion for the
$\mathfrak J$-preadic topology is canonically identified, with the
notation of \sref{III.4.1.7}, with
\[
  \Ext^n_{\sh{O}_{\hat{X}}}
    (\hat{X};\hat{\sh{F}},\hat{\sh{G}}).
\]
\end{proposition}

\begin{proof}
We know from (T, 4.2) that there is a biregular spectral sequence
$\mathrm{E}(\sh{F},\sh{G})$ whose abutment is
\[
  \Ext_{\sh{O}_X}(X;\sh{F},\sh{G})
\]
and whose $\mathrm{E}_2$-terms are
\[
  \mathrm{E}_2^{pq}
  =\HH^p\bigl(X,\shExt^q_{\sh{O}_X}(\sh{F},\sh{G})\bigr).
\]
The $\sh{O}_X$-module
$\shExt^q_{\sh{O}_X}(\sh{F},\sh{G})$ is coherent
\sref[0]{0.12.3.3}; its support is contained in the intersection of
the supports of $\sh{F}$ and $\sh{G}$ (T, 4.2.2), and is consequently
proper over $Y$ \sref[II]{II.5.4.10}.  It follows from
\sref{III.3.2.4} that the $\mathrm{E}_2^{pq}$ are finite
$A$-modules, and then from \sref[0]{0.11.1.8} that the same is true
of every term $\mathrm{E}_r^{pq}$ of the spectral sequence and of
its abutment.

On the other hand, if $i:\hat{X}\to X$ is the canonical morphism,
then $\hat{\sh{F}}$ and $\hat{\sh{G}}$ are canonically identified
with $i^*(\sh{F})$ and $i^*(\sh{G})$, and $i$ is flat
\sref[I]{I.10.8.8}, \sref[I]{I.10.8.9}.  For every $q\geq0$ there
is therefore a canonical $i$-morphism
\[
  u_q:\shExt^q_{\sh{O}_X}(\sh{F},\sh{G})
  \longrightarrow
  \shExt^q_{\sh{O}_{\hat{X}}}
    (\hat{\sh{F}},\hat{\sh{G}})
\]
\sref[0]{0.12.3.4}, and the corresponding
$\sh{O}_{\hat{X}}$-homomorphism
\[
  u_q^\#:
  i^*\bigl(\shExt^q_{\sh{O}_X}(\sh{F},\sh{G})\bigr)
  \longrightarrow
  \shExt^q_{\sh{O}_{\hat{X}}}
    (\hat{\sh{F}},\hat{\sh{G}})
\]
is an isomorphism \sref[0]{0.12.3.5}.  In other words,
\sref[I]{I.10.8.8} identifies
\[
  \shExt^q_{\sh{O}_{\hat{X}}}
    (\hat{\sh{F}},\hat{\sh{G}})
\]
canonically with the completion of
$\shExt^q_{\sh{O}_X}(\sh{F},\sh{G})$, in the notation of
\sref{III.4.1.7}.

The comparison theorem \sref{III.4.1.10} now shows, for every
$p\geq0$, that
\[
  \HH^p\bigl(\hat{X},
    \shExt^q_{\sh{O}_{\hat{X}}}
      (\hat{\sh{F}},\hat{\sh{G}})\bigr)
\]
is canonically identified with the separated completion, for the
$\mathfrak J$-preadic topology, of
\[
  \HH^p\bigl(X,\shExt^q_{\sh{O}_X}(\sh{F},\sh{G})\bigr).
\]

\oldpage[III]{139}

Let $\mathrm{E}(\hat{\sh{F}},\hat{\sh{G}})$ denote the biregular
spectral sequence of (T, 4.2) associated with $\hat{\sh{F}}$ and
$\hat{\sh{G}}$.  If $\hat{A}$ is the separated completion of $A$ for
the $\mathfrak J$-preadic topology, then, up to canonical
isomorphism,
\[
  \mathrm{E}_2^{pq}(\hat{\sh{F}},\hat{\sh{G}})
  =
  \mathrm{E}_2^{pq}(\sh{F},\sh{G})\otimes_A\hat{A}
\]
\sref[0]{0.7.3.3}.

The flat morphism $i$ defines a canonical homomorphism of spectral
sequences
\[
  \varphi:\mathrm{E}(\sh{F},\sh{G})
  \longrightarrow
  \mathrm{E}(\hat{\sh{F}},\hat{\sh{G}})
  =
  \mathrm{E}(i^*(\sh{F}),i^*(\sh{G})).
\]
On the $\mathrm{E}_2$-terms and on the abutment, respectively, this
is the homomorphism
\[
  \varphi_2^{pq}:
  \HH^p\bigl(X,\shExt^q_{\sh{O}_X}(\sh{F},\sh{G})\bigr)
  \longrightarrow
  \HH^p\bigl(\hat{X},
    \shExt^q_{\sh{O}_{\hat{X}}}
      (\hat{\sh{F}},\hat{\sh{G}})\bigr)
\]
and the homomorphism
\[
  \varphi^n:
  \Ext^n_{\sh{O}_X}(X;\sh{F},\sh{G})
  \longrightarrow
  \Ext^n_{\sh{O}_{\hat{X}}}
    (\hat{X};\hat{\sh{F}},\hat{\sh{G}}),
\]
deduced from $u_q$ and $u_0$, respectively, by functoriality
\sref[0]{0.12.3.4}.

After tensoring with $\hat{A}$, the $\varphi_r^{pq}$ and
$\varphi^n$ give homomorphisms of $\hat{A}$-modules
\[
  \psi_r^{pq}:
  \mathrm{E}_r^{pq}(\sh{F},\sh{G})\otimes_A\hat{A}
  \longrightarrow
  \mathrm{E}_r^{pq}(\hat{\sh{F}},\hat{\sh{G}})
\]
and
\[
  \psi^n:
  \Ext^n_{\sh{O}_X}(X;\sh{F},\sh{G})\otimes_A\hat{A}
  \longrightarrow
  \Ext^n_{\sh{O}_{\hat{X}}}
    (\hat{X};\hat{\sh{F}},\hat{\sh{G}}).
\]
Since $\hat{A}$ is a flat $A$-module \sref[0]{0.7.3.3}, the
$\hat{A}$-modules
\[
  \mathrm{E}_r^{pq}(\sh{F},\sh{G})\otimes_A\hat{A}
\]
form a biregular spectral sequence with abutment
\[
  \Ext^n_{\sh{O}_X}(X;\sh{F},\sh{G})\otimes_A\hat{A},
\]
and the $\psi_r^{pq}$ and $\psi^n$ form a morphism of spectral
sequences.  Since the $\psi_2^{pq}$ are isomorphisms, the same is true
of the $\psi^n$ \sref[0]{0.11.1.5}.
\end{proof}

\begin{corollary}[4.5.2]
\label{III.4.5.2}
Under the hypotheses of \sref{III.4.5.1}, suppose in addition that
$A$ is a Noetherian $\mathfrak J$-adic ring.  Then, for every integer
$n\geq0$,
\[
  \Ext^n_{\sh{O}_X}(X;\sh{F},\sh{G})
  \isoto
  \Ext^n_{\sh{O}_{\hat{X}}}
    (\hat{X};\hat{\sh{F}},\hat{\sh{G}})
\]
canonically.
\end{corollary}

\begin{proof}
The finite $A$-module
$\Ext^n_{\sh{O}_X}(X;\sh{F},\sh{G})$ is separated and complete for
the $\mathfrak J$-preadic topology \sref[0]{0.7.3.6}.
\end{proof}

The special case $n=0$ of \sref{III.4.5.1} may be stated as follows.

\begin{corollary}[4.5.3]
\label{III.4.5.3}
Under the hypotheses of \sref{III.4.5.1}, for every homomorphism
$u:\sh{F}\to\sh{G}$ let $\hat{u}:\hat{\sh{F}}\to\hat{\sh{G}}$
denote the completed homomorphism \sref[I]{I.10.8.4}.  There is a
canonical isomorphism
\[
\label{III.4.5.3.1}
  \left(\Hom_{\sh{O}_X}(\sh{F},\sh{G})\right)^\wedge
  \isoto
  \Hom_{\sh{O}_{\hat{X}}}(\hat{\sh{F}},\hat{\sh{G}}),
  \tag{4.5.3.1}
\]
where the left-hand side is the separated completion, for the
$\mathfrak J$-preadic topology, of the $A$-module
$\Hom_{\sh{O}_X}(\sh{F},\sh{G})$.  This isomorphism is obtained by
passing to separated completions in the homomorphism
$u\mapsto\hat{u}$.
\end{corollary}


\subsection{Relations between formal morphisms and ordinary morphisms}
\label{subsection:III.4.6}

\begin{proposition}[4.6.1]
\label{III.4.6.1}
Let $Y$ be a locally Noetherian prescheme, let $f:X\to Y$ be a
proper morphism, let $\sh{F}$ be a coherent $\sh{O}_X$-module which
is flat over $Y$, and let $y\in Y$.  Suppose that, for some integer
$n$,
\[
  \HH^n\bigl(f^{-1}(y),
    \sh{F}\otimes_{\sh{O}_Y}\kres(y)\bigr)=0.
\]
Then there is a neighbourhood $U$ of $y$ in $Y$ such that
\[
  \RR^n f_*(\sh{F})|U=0,
\]
and, for every integer $p\geq0$, the canonical homomorphism
\[
  \bigl(\RR^{n-1}f_*(\sh{F})\bigr)_y
  \longrightarrow
  \HH^{n-1}\bigl(f^{-1}(y),
    \sh{F}\otimes_{\sh{O}_Y}
      (\sh{O}_y/\mathfrak m_y^{p+1})\bigr)
\]
\eqref{III.4.2.1.1} is surjective.
\end{proposition}

\begin{proof}
The $\sh{O}_Y$-module $\RR^n f_*(\sh{F})$ is coherent
\sref{III.3.2.1}; hence the first assertion will follow once we prove
\[
  \bigl(\RR^n f_*(\sh{F})\bigr)_y=0
\]
\sref[0]{0.5.2.2}.  By \sref{III.4.2.1}, it is enough to prove that
\[
  \HH^n\bigl(f^{-1}(y),
    \sh{F}\otimes_{\sh{O}_Y}
      (\sh{O}_y/\mathfrak m_y^{p+1})\bigr)=0
\]
for every $p$.  This is true by hypothesis for $p=0$; we proceed by
induction on $p$.

Put
\[
  X_p=X\times_Y
    \Spec(\sh{O}_y/\mathfrak m_y^{p+1}).
\]
Then $X_{p-1}$ is a closed subprescheme of $X_p$ with the same
underlying space \sref[I]{I.3.6.1}.  The induction hypothesis
\[
  \HH^n\bigl(X_{p-1},
    \sh{F}\otimes_{\sh{O}_Y}
      (\sh{O}_y/\mathfrak m_y^p)\bigr)=0
\]
therefore gives
\[
  \HH^n\bigl(X_p,
    \sh{F}\otimes_{\sh{O}_Y}
      (\sh{O}_y/\mathfrak m_y^p)\bigr)=0.
\]
On the other hand, the exact sequence of $\sh{O}_{X_p}$-modules
\[
  0\longrightarrow
  \mathfrak m_y^p\sh{F}/\mathfrak m_y^{p+1}\sh{F}
  \longrightarrow
  \sh{F}/\mathfrak m_y^{p+1}\sh{F}
  \longrightarrow
  \sh{F}/\mathfrak m_y^p\sh{F}
  \longrightarrow0
\]
gives the exact cohomology sequence
\[
  \HH^n\bigl(X_p,
    \mathfrak m_y^p\sh{F}/\mathfrak m_y^{p+1}\sh{F}\bigr)
  \longrightarrow
  \HH^n\bigl(X_p,\sh{F}/\mathfrak m_y^{p+1}\sh{F}\bigr)
  \longrightarrow
  \HH^n\bigl(X_p,\sh{F}/\mathfrak m_y^p\sh{F}\bigr).
\]
It is thus enough to show that
\[
\label{III.4.6.1.1}
  \HH^n\bigl(X_p,
    \mathfrak m_y^p\sh{F}/\mathfrak m_y^{p+1}\sh{F}\bigr)=0.
  \tag{4.6.1.1}
\]
Indeed, the middle term will then be a submodule of the last term,
which is zero by the induction hypothesis.

\oldpage[III]{140}

The fibre
\[
  Z=f^{-1}(y)=X\times_Y\Spec(\kres(y))
\]
is a closed subprescheme of $X_p$, and
$\mathfrak m_y^p\sh{F}/\mathfrak m_y^{p+1}\sh{F}$ is annihilated by
$\mathfrak m_y$.  It may therefore be regarded as an
$\sh{O}_Z$-module, and
\[
  \HH^n\bigl(Z,
    \mathfrak m_y^p\sh{F}/\mathfrak m_y^{p+1}\sh{F}\bigr)
  =
  \HH^n\bigl(X_p,
    \mathfrak m_y^p\sh{F}/\mathfrak m_y^{p+1}\sh{F}\bigr).
\]
We shall now prove that the canonical $\sh{O}_Z$-homomorphism
\[
\label{III.4.6.1.2}
  (\sh{F}/\mathfrak m_y\sh{F})
    \otimes_{\kres(y)}
      (\mathfrak m_y^p/\mathfrak m_y^{p+1})
  \longrightarrow
  \mathfrak m_y^p\sh{F}/\mathfrak m_y^{p+1}\sh{F}
  \tag{4.6.1.2}
\]
is bijective.  Once this is established, the fact that
$\mathfrak m_y^p/\mathfrak m_y^{p+1}$ is a free
$\kres(y)$-module gives
\[
\begin{split}
  \HH^n\bigl(Z,
    \mathfrak m_y^p\sh{F}/\mathfrak m_y^{p+1}\sh{F}\bigr)
  &=
  \HH^n\bigl(Z,\sh{F}/\mathfrak m_y\sh{F}\bigr)
    \otimes_{\kres(y)}
      (\mathfrak m_y^p/\mathfrak m_y^{p+1})\\
  &=0
\end{split}
\]
\sref[0]{0.12.2.3}, since
$\HH^n(Z,\sh{F}/\mathfrak m_y\sh{F})=0$ by hypothesis.  This proves
\eqref{III.4.6.1.1}.

To finish the proof of the first assertion it remains only to show
that \eqref{III.4.6.1.2} is bijective.  The question is pointwise on
$X$, and $\sh{F}_x$ is a flat $\sh{O}_y$-module by hypothesis for
every $x\in f^{-1}(y)$.  It is therefore enough to apply
\sref[0]{0.10.2.1}, c), since
$\sh{F}_x/\mathfrak m_y\sh{F}_x$ is flat over the field
$\kres(y)=\sh{O}_y/\mathfrak m_y$.

For the second assertion of \sref{III.4.6.1}, we reduce at once, as
in \sref{III.4.2.1}, to the case in which $Y$ is affine and $y$ is
closed.  The same argument shows, for every $k>0$, that
\[
\label{III.4.6.1.3}
  \HH^n\bigl(X_{k+1},
    \mathfrak m_y^k\sh{F}/
      \mathfrak m_y^{k+p+1}\sh{F}\bigr)=0.
  \tag{4.6.1.3}
\]

\oldpage[III]{141}

It follows from \sref{III.4.2.1} that
\[
\label{III.4.6.1.4}
  \bigl(\RR^n f_*(\mathfrak m_y^k\sh{F})\bigr)_y=0.
  \tag{4.6.1.4}
\]
The exact cohomology sequence now yields the exactness of
\[
  \bigl(\RR^{n-1}f_*(\sh{F})\bigr)_y
  \longrightarrow
  \bigl(\RR^{n-1}f_*
    (\sh{F}/\mathfrak m_y^p\sh{F})\bigr)_y
  \longrightarrow
  \bigl(\RR^n f_*(\mathfrak m_y^p\sh{F})\bigr)_y=0.
\]
Since $y$ is closed and $Y$ is affine, we have
\sref{III.1.4.11}
\[
  \RR^{n-1}f_*(\sh{F}/\mathfrak m_y^p\sh{F})
  =
  \bigl(\HH^{n-1}
    (X,\sh{F}/\mathfrak m_y^p\sh{F})\bigr)^{\sim}
  =
  \bigl(\HH^{n-1}
    (f^{-1}(y),\sh{F}/\mathfrak m_y^p\sh{F})\bigr)^{\sim}
\]
(G, II, 4.9.1).  The module
\[
  \HH^{n-1}
    (f^{-1}(y),\sh{F}/\mathfrak m_y^p\sh{F})
\]
is an $\sh{O}_y/\mathfrak m_y^p$-module, and hence
\[
  \bigl(\RR^{n-1}f_*
    (\sh{F}/\mathfrak m_y^p\sh{F})\bigr)_y
  =
  \HH^{n-1}
    (f^{-1}(y),\sh{F}/\mathfrak m_y^p\sh{F}).
\]
This completes the proof.
\end{proof}

\begin{corollary}[4.6.2]
\label{III.4.6.2}
Let $Y$ be a locally Noetherian prescheme, let $f:X\to Y$ be a
proper flat morphism, let $\sh{F}$ and $\sh{G}$ be two locally free
$\sh{O}_X$-modules, and let $y$ be a point of $Y$.  Put
\[
  X_y=f^{-1}(y)=X\times_Y\Spec(\kres(y)),\qquad
  \sh{F}_y=\sh{F}\otimes_{\sh{O}_Y}\kres(y),\qquad
  \sh{G}_y=\sh{G}\otimes_{\sh{O}_Y}\kres(y),
\]
and suppose that
\[
\label{III.4.6.2.1}
  \HH^1\bigl(
    X_y,\shHom_{\sh{O}_{X_y}}(\sh{F}_y,\sh{G}_y)
  \bigr)=0.
  \tag{4.6.2.1}
\]
Then, for every homomorphism $u_0:\sh{F}_y\to\sh{G}_y$, there exist
an open neighbourhood $U$ of $y$ in $Y$ and a homomorphism
\[
  u:\sh{F}|_{f^{-1}(U)}\longrightarrow
    \sh{G}|_{f^{-1}(U)}
\]
such that $u_0=u\otimes1$.
\end{corollary}

\begin{proof}
The hypothesis permits us to apply \sref{III.4.6.1}, with $n=1$
and $p=0$, to the coherent $\sh{O}_X$-module
\[
  \sh{H}=\shHom_{\sh{O}_X}(\sh{F},\sh{G}).
\]
Indeed, $\sh{H}$ is locally free and hence, a fortiori, flat over
$Y$, while the $\sh{O}_{X_y}$-module
$\sh{H}\otimes_{\sh{O}_Y}\kres(y)$ identifies with
$\shHom_{\sh{O}_{X_y}}(\sh{F}_y,\sh{G}_y)$
\sref[0]{0.6.2.2}.  We may suppose that $Y=\Spec(A)$ is affine.
Then \sref{III.1.4.11} gives
\[
  \RR^0f_*(\sh{H})
  =
  \bigl(\Hom_{\sh{O}_X}(\sh{F},\sh{G})\bigr)^{\sim},
\]
and consequently
\[
  \bigl(\RR^0f_*(\sh{H})\bigr)_y
  =
  \Hom_{\sh{O}_X}(\sh{F},\sh{G})
    \otimes_A\sh{O}_y.
\]
By \sref{III.4.6.1}, the canonical homomorphism
\[
  \Hom_{\sh{O}_X}(\sh{F},\sh{G})
    \otimes_A\sh{O}_y
  \longrightarrow
  \Hom_{\sh{O}_{X_y}}(\sh{F}_y,\sh{G}_y)
\]
is surjective.  This proves the corollary, since every element of
the module on the left may be written as $u\otimes(1/s)$, where
$s\in A\setminus\mathfrak m_y$.
\end{proof}

\begin{corollary}[4.6.3]
\label{III.4.6.3}
Under the hypotheses of \sref{III.4.6.2}, if $u_0$ is injective
(respectively surjective, bijective), one may choose $u$ to be
injective (respectively surjective, bijective).
\end{corollary}

\begin{proof}
We may restrict ourselves to the case $U=Y$.  It is enough to prove
that, if $u_0$ is injective (respectively surjective), then
\[
  \Ker u_x=0
  \quad\text{(respectively }\Coker u_x=0\text{)}
\]
for every $x\in f^{-1}(y)$.  Indeed, $\Ker u$ and $\Coker u$ are
coherent $\sh{O}_X$-modules \sref[0]{0.5.3.4}; hence there is a
neighbourhood $V$ of $f^{-1}(y)$ in $X$ on which the restriction of
$\Ker u$ (respectively $\Coker u$) is zero \sref[0]{0.5.2.2}.
Since $f$ is closed, there is an open neighbourhood
$U'\subset U$ of $y$ such that $f^{-1}(U')\subset V$, which proves
the assertion after replacing $U$ by $U'$.

\oldpage[III]{142}

By hypothesis,
\[
  u_x\otimes1:
  \sh{F}_x\otimes_{\sh{O}_y}\kres(y)
  \longrightarrow
  \sh{G}_x\otimes_{\sh{O}_y}\kres(y)
\]
is injective (respectively surjective), $\sh{F}_x$ and $\sh{G}_x$
are finite free $\sh{O}_x$-modules, and $\sh{O}_x$ is flat over
$\sh{O}_y$.  If $u_x\otimes1$ is injective, the injectivity of $u_x$
follows from \sref[0]{0.10.2.4}.  If $u_x\otimes1$ is surjective,
then, a fortiori, the induced homomorphism
\[
  \sh{F}_x/\mathfrak m_x\sh{F}_x
  \longrightarrow
  \sh{G}_x/\mathfrak m_x\sh{G}_x
\]
is surjective.  Since $\sh{G}_x$ is a finite $\sh{O}_x$-module and
$\sh{O}_x$ is a local ring with maximal ideal $\mathfrak m_x$, the
conclusion follows from Nakayama's lemma (Bourbaki,
\emph{Alg\`ebre}, chap.~VIII, \S\,6, no.~3, cor.~4 to prop.~6).
\end{proof}

\begin{corollary}[4.6.4]
\label{III.4.6.4}
Let $Y$ be a locally Noetherian prescheme, let $f:X\to Y$ be a
proper flat morphism, let $y$ be a point of $Y$, and put
$X_y=X\times_Y\Spec(\kres(y))$.  Let $\sh{E}_0$ be a locally free
$\sh{O}_{X_y}$-module such that
\[
\label{III.4.6.4.1}
  \HH^1\bigl(
    X_y,\shHom_{\sh{O}_{X_y}}(\sh{E}_0,\sh{E}_0)
  \bigr)=0.
  \tag{4.6.4.1}
\]
Let $\sh{F}$ and $\sh{G}$ be two locally free $\sh{O}_X$-modules
such that $\sh{F}_y$ and $\sh{G}_y$, with the notation of
\sref{III.4.6.2}, are isomorphic to $\sh{E}_0$.  Then there is an
open neighbourhood $U$ of $y$ such that
$\sh{F}|_{f^{-1}(U)}$ and $\sh{G}|_{f^{-1}(U)}$ are isomorphic.
\end{corollary}

\begin{corollary}[4.6.5]
\label{III.4.6.5}
Under the hypotheses of \sref{III.4.6.4} on $f$, $X$, and $Y$,
suppose that
\[
  \HH^1(X_y,\sh{O}_{X_y})=0.
\]
If $\sh{F}$ and $\sh{G}$ are two invertible $\sh{O}_X$-modules such
that $\sh{F}_y$ and $\sh{G}_y$ are isomorphic, then there is an open
neighbourhood $U$ of $y$ such that $\sh{F}|_{f^{-1}(U)}$ and
$\sh{G}|_{f^{-1}(U)}$ are isomorphic.
\end{corollary}

\begin{proof}
It is enough to apply \sref{III.4.6.4} to the modules
$\sh{F}^{-1}\otimes\sh{G}$ and $\sh{O}_X$.
\end{proof}

\begin{remarks}[4.6.6]
\label{III.4.6.6}
\begin{enumerate}
\item[(i)] Using \sref{III.4.6.5}, we shall establish in Chapter~V
  the classification of invertible sheaves on a projective bundle
  announced in \sref[II]{II.4.2.7}.
\item[(ii)] The result of \sref{III.4.6.1} will appear in \S\,7 as
  a consequence of more general propositions.
\end{enumerate}
\end{remarks}

\begin{proposition}[4.6.7]
\label{III.4.6.7}
Let $Z$ be a locally Noetherian prescheme, and let $X$ and $Y$ be
two $Z$-preschemes whose structural morphisms
\[
  g:X\longrightarrow Z,\qquad h:Y\longrightarrow Z
\]
are proper.  Let $f:X\to Y$ be a $Z$-morphism, let $z$ be a point
of $Z$, and put
\[
  f_z=f\times_Z1:
  X\times_Z\Spec(\kres(z))
  \longrightarrow
  Y\times_Z\Spec(\kres(z)).
\]
\begin{enumerate}
\item[(i)] If $f_z$ is finite (respectively a closed immersion),
  there is an open neighbourhood $U$ of $z$ such that the morphism
  \[
    g^{-1}(U)\longrightarrow h^{-1}(U)
  \]
  induced by $f$ is finite (respectively a closed immersion).
\item[(ii)] Suppose in addition that $g$ is flat.  If $f_z$ is an
  isomorphism, there is an open neighbourhood $U$ of $z$ such that
  the morphism $g^{-1}(U)\to h^{-1}(U)$ induced by $f$ is an
  isomorphism.
\end{enumerate}
\end{proposition}

\begin{proof}
In either case it is enough to prove that, for every
$y\in h^{-1}(z)$, there is a neighbourhood $V_y$ of $y$ such that
the restriction
\[
  f^{-1}(V_y)\longrightarrow V_y
\]
of $f$ is finite (respectively a closed immersion, an isomorphism).
Indeed, if $V$ is the union of the $V_y$, then
$f^{-1}(V)\to V$ is finite (respectively a closed immersion, an
isomorphism) \sref[II]{II.6.1.1}, \sref[I]{I.4.2.4}.  Since $h$ is
closed, there is an open neighbourhood $U$ of $z$ such that
$h^{-1}(U)\subset V$, and the proposition follows.

\noindent\textit{(i)} First note that $f$ is proper
\sref[II]{II.5.4.3}.  If
$f_z$ is finite, the required neighbourhood $V_y$, for each
$y\in h^{-1}(z)$, exists by \sref{III.4.4.11}.

\oldpage[III]{143}

To treat the case in which $f_z$ is a closed immersion, we may
therefore already suppose that $f$ is finite.  Thus
$X=\Spec(\sh{B})$, where $\sh{B}$ is a coherent
$\sh{O}_Y$-algebra, and $f$ corresponds \sref[II]{II.1.2.7} to the
canonical homomorphism
\[
  u:\sh{O}_Y\longrightarrow\sh{B}.
\]
If, for every $y\in h^{-1}(z)$, the homomorphism
$u_y:\sh{O}_y\to\sh{B}_y$ is surjective, then $u$ is surjective on
some neighbourhood $V_y$ of $y$, since $\Coker u$ is coherent
\sref[0]{0.5.3.4}, \sref[0]{0.5.2.2}.  Now the finite morphism
$f_z$ corresponds to
\[
  v=u\otimes1:
  \sh{O}_Y\otimes_{\sh{O}_Z}\kres(z)
  \longrightarrow
  \sh{B}\otimes_{\sh{O}_Z}\kres(z),
\]
and the hypothesis that $f_z$ is a closed immersion implies that
\[
  u_y\otimes1:
  \sh{O}_y\otimes_{\sh{O}_z}\kres(z)
  \longrightarrow
  \sh{B}_y\otimes_{\sh{O}_z}\kres(z)
\]
is surjective.  Since $\sh{B}_y$ is a finite $\sh{O}_y$-module and
$\sh{O}_y$ is a local Noetherian ring, the conclusion follows, as
in \sref{III.4.6.3}, from Nakayama's lemma.

\noindent\textit{(ii)} The same reasoning shows that it is now
enough to prove
that $u_y$ is bijective, knowing that $u_y\otimes1$ is bijective.
The required injectivity follows from the lemma below; surjectivity
follows from the preceding argument.
\end{proof}

\begin{lemma}[4.6.7.1]
\label{III.4.6.7.1}
Let $A$ and $B$ be local Noetherian rings, let
$\rho:A\to B$ be a local homomorphism, and let $u:N\to M$ be a
homomorphism of $B$-modules.  Suppose that $M$ is flat over $A$,
that $N$ is a finite $B$-module, and that
\[
  u\otimes1:
  N\otimes_A k\longrightarrow M\otimes_A k,
\]
where $k$ is the residue field of $A$, is injective.  Then $N$ is
flat over $A$ and $u$ is injective.
\end{lemma}

\begin{proof}
To prove the first assertion, we must show that, for every pair of
finite $A$-modules $P,Q$ and every injective $A$-homomorphism
$v:P\to Q$, the homomorphism
\[
  1_N\otimes v:N\otimes_A P\longrightarrow N\otimes_A Q
\]
is injective.  Consider the commutative diagram
\[
  \xymatrix{
    N\otimes_A P
      \ar[r]^-{1_N\otimes v}
      \ar[d]_{u\otimes1_P}
      &
    N\otimes_A Q
      \ar[d]^{u\otimes1_Q}
    \\
    M\otimes_A P
      \ar[r]_{1_M\otimes v}
      &
    M\otimes_A Q .
  }
\]
Since $1_M\otimes v$ is injective by the flatness of $M$, it is
enough to prove that $u\otimes1_P$ is injective.

Let $\mathfrak m$ be the maximal ideal of $A$.  The
$\mathfrak m$-adic filtration on the $A$-module $N\otimes_A P$ is
also its $\mathfrak mB$-adic filtration as a $B$-module.  The
topology defined by this filtration is separated: $B$ is
Noetherian, $\mathfrak mB$ is contained in the radical of $B$, and
$N\otimes_A P$ is a finite $B$-module, since $N$ is finite over $B$
and $P$ is finite over $A$ \sref[0]{0.7.3.5}.  It is therefore
enough to prove that
\[
  \operatorname{gr}(u\otimes1_P):
  \operatorname{gr}_\bullet(N\otimes_A P)
  \longrightarrow
  \operatorname{gr}_\bullet(M\otimes_A P)
\]
is injective, where the associated graded modules are taken with
respect to the $\mathfrak m$-adic filtrations (Bourbaki,
\emph{Alg\`ebre commutative}, chap.~III, \S\,2, no.~8, cor.~1 to
th.~1).

Since $M$ is flat over $A$, the homomorphisms
\[
  M\otimes_A(\mathfrak m^nP)
  \longrightarrow
  \mathfrak m^n(M\otimes_A P)
\]
are bijective.  Consequently the canonical homomorphism
\[
  \vphi_M:
  \operatorname{gr}_0(M)\otimes_A\operatorname{gr}_\bullet(P)
  \longrightarrow
  \operatorname{gr}_\bullet(M\otimes_A P)
\]
is bijective.

\oldpage[III]{144}

We have the commutative diagram
\[
  \xymatrix{
    \operatorname{gr}_0(N)\otimes_A\operatorname{gr}_\bullet(P)
      \ar[r]^-{\operatorname{gr}_0(u)\otimes1}
      \ar[d]_{\vphi_N}
      &
    \operatorname{gr}_0(M)\otimes_A\operatorname{gr}_\bullet(P)
      \ar[d]^{\vphi_M}
    \\
    \operatorname{gr}_\bullet(N\otimes_A P)
      \ar[r]_{\operatorname{gr}(u\otimes1)}
      &
    \operatorname{gr}_\bullet(M\otimes_A P).
  }
\]
Here $\vphi_M$ is bijective and $\vphi_N$ is surjective.
Moreover, $\operatorname{gr}_0(u)$ is injective by hypothesis, and
\[
\begin{split}
  \operatorname{gr}_0(N)\otimes_A\operatorname{gr}_\bullet(P)
    &=
  \operatorname{gr}_0(N)\otimes_k\operatorname{gr}_\bullet(P),\\
  \operatorname{gr}_0(M)\otimes_A\operatorname{gr}_\bullet(P)
    &=
  \operatorname{gr}_0(M)\otimes_k\operatorname{gr}_\bullet(P).
\end{split}
\]
Thus $\operatorname{gr}_0(u)\otimes1$ is also injective.  It
follows from the diagram that $\operatorname{gr}(u\otimes1_P)$ is
injective, proving the first assertion.  The second follows from
the same argument by taking $P=A$.
\end{proof}

\begin{proposition}[4.6.8]
\label{III.4.6.8}
Let $Z$ be a locally Noetherian prescheme, and let $X$ and $Y$ be
two $Z$-preschemes whose structural morphisms
$g:X\to Z$ and $h:Y\to Z$ are proper.  Let $Z'$ be a closed subset
of $Z$, put
\[
  X'=g^{-1}(Z'),\qquad Y'=h^{-1}(Z'),
\]
and denote by
\[
  \widehat X=X_{/X'},\qquad
  \widehat Y=Y_{/Y'},\qquad
  \widehat Z=Z_{/Z'}
\]
the formal completions along these closed subsets.  Let $f:X\to Y$
be a $Z$-morphism and let
$\widehat f:\widehat X\to\widehat Y$ be its extension to the formal
completions.  In order that $\widehat f$ be an isomorphism
(respectively a closed immersion), it is necessary and sufficient
that there exist an open neighbourhood $U$ of $Z'$ such that the
morphism
\[
  g^{-1}(U)\longrightarrow h^{-1}(U)
\]
induced by $f$ is an isomorphism (respectively a closed immersion).
\end{proposition}

\begin{proof}
The sufficiency of the condition is immediate
\sref[I]{I.10.14.7}.  To prove its necessity, it is enough, by the
same reasoning as in \sref{III.4.6.7}, to show that for every
$y\in Y'$ there is an open neighbourhood $V_y$ of $y$ such that
$f^{-1}(V_y)\to V_y$ is an isomorphism (respectively a closed
immersion).  We are thus reduced to the case
\[
  Y=Z=\Spec(A),
\]
where $A$ is Noetherian.

By hypothesis \sref[I]{I.10.9.1}, \sref[I]{I.10.14.2}, the fibre
$f^{-1}(y)$ consists of a single point for every $y\in Y'$.  Since
$f$ is proper \sref[II]{II.5.4.3}, \sref{III.4.4.11} gives an open
neighbourhood $U$ of $y$ for which $f^{-1}(U)\to U$ is finite.  We
may therefore suppose that $f$ is finite, so that
\[
  X=\Spec(B),
\]
where $B$ is a finite $A$-algebra.

If $Y'=V(\mathfrak J)$, then
\[
  \widehat Y=\Spf(\widehat A),
  \qquad
  \widehat X=\Spf(\widehat B),
\]
where $\widehat A$ is the separated completion of $A$ for the
$\mathfrak J$-adic topology, and $\widehat B$ is the separated
completion of $B$ for the $\mathfrak JB$-adic topology, or,
equivalently, the separated completion of the $A$-module $B$ for
the $\mathfrak J$-adic topology.  Furthermore, $\widehat f$
corresponds to the continuous extension
\[
  \widehat{\vphi}:\widehat A\longrightarrow\widehat B
\]
of the canonical ring homomorphism $\vphi:A\to B$, and the
hypothesis is that $\widehat{\vphi}$ is surjective (respectively
bijective) \sref[I]{I.10.14.2}.  But $\widehat{\vphi}$ is also the
continuous extension of $\vphi$ regarded as a homomorphism of
$A$-modules.  It follows from \sref[I]{I.10.8.14} that there is an
open neighbourhood $U$ of $Y'$ such that the restriction to $U$ of
the homomorphism of $\sh{O}_Y$-modules
\[
  \widetilde{\vphi}:\widetilde A\longrightarrow\widetilde B
\]
is surjective (respectively bijective), which completes the proof.
\end{proof}


\subsection{An ampleness criterion}
\label{subsection:III.4.7}

\begin{theorem}[4.7.1]
\label{III.4.7.1}
Let $Y$ be a locally Noetherian prescheme, let $f:X\to Y$ be a
proper morphism, let $\sh{L}$ be an invertible $\sh{O}_X$-module,
and let $y$ be a point of $Y$.  Put
\[
  X_y=X\times_Y\Spec(\kres(y))=f^{-1}(y),
\]
and let $g:X_y\to X$ be the canonical morphism.  If
\[
  \sh{L}_y=g^*(\sh{L})
  =\sh{L}\otimes_{\sh{O}_Y}\kres(y)
\]
is ample on $X_y$, then there is an open neighbourhood $U$ of $y$
in $Y$ such that $\sh{L}|_{f^{-1}(U)}$ is ample for the restriction
of $f$ to $f^{-1}(U)$.
\end{theorem}

\begin{proof}
\oldpage[III]{145}

\noindent\textit{I)} Put
\[
  Y'=\Spec(\sh{O}_y),\qquad
  X'=X\times_Y Y',\qquad
  \sh{L}'=\sh{L}\otimes_{\sh{O}_Y}\sh{O}_y.
\]
We first prove that $\sh{L}'$ is ample for $f'=f_{(Y')}$.  We have
the commutative diagram
\[
  \xymatrix{
    X\ar[d]_f
      &
    X'\ar[l]\ar[d]_{f'}
      &
    X_y\ar[l]\ar[d]_{f_y}
    \\
    Y
      &
    Y'\ar[l]
      &
    \Spec(\kres(y))\ar[l].
  }
\]
Since $f'$ is proper \sref[II]{II.5.4.2}, (iii), and $\sh{O}_y$ is
Noetherian, \sref[II]{II.4.6.6} shows that we may reduce to the case
\[
  Y=Y'=\Spec(\sh{O}_y),\qquad X=X',
\]
assume that $\sh{L}_y$ is ample for $f_y$, and prove that $\sh{L}$
is ample for $f$.  We apply the criterion
\sref{III.2.6.1}, c), and prove that for every coherent
$\sh{O}_X$-module $\sh{F}$ there is an integer $N$ such that
\[
  \HH^1(X,\sh{F}(n))=0
  \quad\text{for all }n\geq N,
  \qquad
  \sh{F}(n)=\sh{F}\otimes\sh{L}^{\otimes n}.
\]

The point $y$ is closed in $Y$ and corresponds to the maximal ideal
$\mathfrak m$ of $\sh{O}_y$.  Hence $X_y$ is the closed
subprescheme of $X$ defined by the coherent ideal
\[
  \sh{J}
  =
  f^*(\widetilde{\mathfrak m})\sh{O}_X
  =
  \mathfrak m\sh{O}_X
\]
\sref[I]{I.4.4.5}, and $g$ is its canonical injection.  Consider
the graded $\kres(y)$-algebra
\[
  S=\operatorname{gr}(\sh{O}_y)
   =\bigoplus_{j\geq0}
      \mathfrak m^j/\mathfrak m^{j+1}.
\]
It is of finite type because $\sh{O}_y$ is Noetherian.  The
$\sh{O}_X$-algebra
\[
  \sh{S}=f^*(\widetilde S)
\]
is quasi-coherent and of finite type, and is annihilated by
$\sh{J}$.  Thus, if $\sh{S}_y=g^*(\sh{S})$, then $\sh{S}_y$ is a
quasi-coherent $\sh{O}_{X_y}$-algebra of finite type and
\[
  \sh{S}=g_*(\sh{S}_y).
\]

Put
\[
  \sh{M}_j
  =
  \mathfrak m^j\sh{F}/\mathfrak m^{j+1}\sh{F},
  \qquad
  \sh{M}
  =
  \bigoplus_{j\geq0}\sh{M}_j
  =
  \operatorname{gr}(\sh{F}).
\]
Since $\sh{F}$ is coherent, $\sh{M}$ is a quasi-coherent
$\sh{S}$-module of finite type \sref[0]{0.10.1.1}; it too is
annihilated by $\sh{J}$.  Consequently, if
\[
  \sh{M}'_j=g^*(\sh{M}_j),
  \qquad
  \sh{M}'=g^*(\sh{M})
    =\bigoplus_{j\geq0}\sh{M}'_j,
\]
then $\sh{M}'$ is a graded quasi-coherent $\sh{S}_y$-module of
finite type and $\sh{M}=g_*(\sh{M}')$.  Moreover, on putting
\[
  \sh{M}'_j(n)
  =
  \sh{M}'_j\otimes\sh{L}_y^{\otimes n},
\]
we have $\sh{M}'_j(n)=g^*(\sh{M}_j(n))$.

The morphism $f_y$ is proper \sref[II]{II.5.4.2}, (iii), and
$\sh{L}_y$ is ample; hence $f_y$ is projective
\sref[II]{II.5.5.4}, \sref[II]{II.4.6.11}.  We may therefore apply
\sref{III.2.4.1}, (ii), to
$\Spec(\kres(y))$, $f_y$, $\sh{S}_y$, $\sh{L}_y$, and $\sh{M}'$.
There is an integer $N$ such that, for $n\geq N$,
\[
  \HH^q(X_y,\sh{M}'_j(n))=0
\]
for every $q>0$ and every $j$.  It follows that
\[
  \HH^q(X,\sh{M}_j(n))=0
\]
for every $q>0$ and every $j$ (G, II, 4.9.1).

Now put
\[
  \sh{F}(n)_j
  =
  \sh{F}(n)/\mathfrak m^{j+1}\sh{F}(n).
\]
For $j\geq1$ we have
\[
  \sh{F}(n)_{j-1}
  =
  \sh{F}(n)_j/\sh{M}_j(n),
  \qquad
  \sh{F}(n)_0=\sh{M}_0(n).
\]
Thus $\HH^1(X,\sh{F}(n)_0)=0$, and the exact cohomology sequence
gives, for every $j\geq1$,
\[
  \HH^1(X,\sh{F}(n)_j)
  =
  \HH^1(X,\sh{F}(n)_{j-1}).
\]
Hence $\HH^1(X,\sh{F}(n)_j)=0$ for every $j\geq0$.  By
\sref{III.4.2.1}, it follows that
\[
  \HH^1(X,\sh{F}(n))=0,
\]
which proves the assertion.

\oldpage[III]{146}

\noindent\textit{II)} Return to the notation at the beginning of
the proof.  We may suppose that $Y=\Spec(A)$ is affine.  Since
$f'$ is of finite type and $\sh{L}'$ is ample for $f'$, there is an
integer $m>0$ such that $\sh{L}'^{\otimes m}$ is very ample for
$f'$ \sref[II]{II.4.6.11}.  Replacing $\sh{L}$ by
$\sh{L}^{\otimes m}$, it is enough to consider the case in which
$\sh{L}'$ is very ample for $f'$ and to prove that
$\sh{L}|_{f^{-1}(U)}$ is very ample for the restriction of $f$.

Since $f'$ is proper, there is, for a suitable integer $r>0$, a
closed $Y'$-immersion
\[
  j:X'\longrightarrow P=\bb{P}_{Y'}^r
\]
such that $\sh{L}'\simeq j^*(\sh{O}_P(1))$
\sref[II]{II.5.5.4}, (ii).  This closed immersion corresponds
canonically to a surjective $\sh{O}_{X'}$-homomorphism
\[
  u:\sh{O}_{X'}^{r+1}\longrightarrow\sh{L}'
\]
\sref[II]{II.4.2.3}, or equivalently \sref[0]{0.5.1.1} to
$r+1$ sections $s'_i$ $(0\leq i\leq r)$ of $\sh{L}'$ over $X'$
which generate this $\sh{O}_{X'}$-module.

By definition these are also sections of $f'_*(\sh{L}')$ over
$Y'$, and
\[
  f'_*(\sh{L}')
  =
  f_*(\sh{L})\otimes_{\sh{O}_Y}\sh{O}_{Y'}
\]
\sref[0]{0.5.4.10}.  Since $Y$ is affine and $\sh{O}_y$ is the
local ring of $A$ at the prime ideal $\mathfrak j_y$, each $s'_i$
has the form
\[
  s'_i=s''_i/t_i,
\]
where $s''_i$ is a section of $f_*(\sh{L})$ over $Y$ and
$t_i\in A\setminus\mathfrak j_y$.  It follows that there are an
open neighbourhood $V$ of $y$ in $Y$ and sections $s_i$ of
$f_*(\sh{L})|_V$ such that $s'_i=s_i/1$ (the space $Y'$ is
contained in $V$; see \sref[I]{I.2.4.2}).

The sections $s_i$ are sections of $\sh{L}$ over $f^{-1}(V)$ and
therefore define a homomorphism
\[
  v:
  \bigl(\sh{O}_X|_{f^{-1}(V)}\bigr)^{r+1}
  \longrightarrow
  \sh{L}|_{f^{-1}(V)}.
\]
By construction it is surjective at every point of $f^{-1}(y)$.
Since $\Coker(v)$ is coherent \sref[0]{0.5.3.4}, its support is
closed \sref[0]{0.5.2.2}; hence there is an open neighbourhood
$W\subset f^{-1}(V)$ of $f^{-1}(y)$ on which $v$ is surjective.
Because $f$ is closed, after shrinking around $y$ we may choose an
open neighbourhood $U\subset V$ with $f^{-1}(U)\subset W$.
The conclusion now follows from \sref[II]{II.4.2.3}.
\end{proof}


\subsection{Finite morphisms of formal preschemes}
\label{subsection:III.4.8}

\begin{proposition}[4.8.1]
\label{III.4.8.1}
Let $\mathfrak{Y}$ be a locally Noetherian formal prescheme,
$\sh{K}$ an ideal of definition of $\mathfrak{Y}$, and
$f:\mathfrak{X}\to\mathfrak{Y}$ a morphism of formal preschemes.
The following conditions are equivalent:
\begin{enumerate}
  \item[\rm{a)}] $\mathfrak{X}$ is locally Noetherian, $f$ is an
    adic morphism \sref[I]{I.10.12.1}, and, if
    $\sh{J}=f^*(\sh{K})\sh{O}_{\mathfrak{X}}$, the morphism
    \[
      f_0:
      (\mathfrak{X},\sh{O}_{\mathfrak{X}}/\sh{J})
      \longrightarrow
      (\mathfrak{Y},\sh{O}_{\mathfrak{Y}}/\sh{K})
    \]
    induced by $f$ is finite.
  \item[\rm{b)}] $\mathfrak{X}$ is locally Noetherian and is the
    inductive limit of an adic inductive $(Y_n)$-system $(X_n)$
    such that the morphism $X_0\to Y_0$ is finite.
  \item[\rm{c)}] Every point of $\mathfrak{Y}$ has a Noetherian
    formal affine open neighbourhood $V$ such that $f^{-1}(V)$ is
    a formal affine open and
    $\Gamma(f^{-1}(V),\sh{O}_{\mathfrak{X}})$ is a
    $\Gamma(V,\sh{O}_{\mathfrak{Y}})$-module of finite type.
\end{enumerate}
\end{proposition}

\begin{proof}
It is immediate from \sref[I]{I.10.12.3} that \emph{a)} implies
\emph{b)}.  To see that \emph{b)} implies \emph{c)}, we may suppose
that
\[
  \mathfrak{Y}=\Spf(B),\qquad \mathfrak{X}=\Spf(A),
\]
where $B$ is an adic Noetherian ring and $\mathfrak{R}$ is an ideal
of definition of $B$.  By hypothesis, $X_0$ is an affine scheme
whose ring $A_0$ is a $B/\mathfrak{R}$-module of finite type
\sref[II]{II.6.1.3}.  By \sref[I]{I.5.1.9}, each $X_n$ is an affine
scheme; if $A_n$ denotes its ring, condition \emph{b)} implies that,
for $m\leq n$,
\[
  A_m\simeq A_n/\mathfrak{R}^{m+1}A_n.
\]
It follows that $\mathfrak{X}\simeq\Spf(A)$, where
$A=\varprojlim A_n$, and the conclusion follows from
\sref[0]{0.7.2.9}.

Finally, to prove that \emph{c)} implies \emph{a)}, we may again
restrict to the case
\[
  \mathfrak{Y}=\Spf(B),\qquad \mathfrak{X}=\Spf(A),
\]
where $A$ is a finite $B$-algebra.  Since
$A/\mathfrak{R}A$ is then a finite
$B/\mathfrak{R}$-algebra, \sref[I]{I.10.10.9} shows that the
conditions in \emph{a)} are satisfied.
\end{proof}

\begin{definition}[4.8.2]
\label{III.4.8.2}
When the equivalent conditions \emph{a)}, \emph{b)}, and \emph{c)}
of \sref{III.4.8.1} are satisfied, the morphism $f$ is said to be
\emph{finite}; one also says that $\mathfrak{X}$ is a
\emph{finite formal $\mathfrak{Y}$-prescheme}, or a formal
prescheme \emph{finite over $\mathfrak{Y}$}.
\end{definition}

\begin{proposition}[4.8.3]
\label{III.4.8.3}
\begin{enumerate}
  \item[\rm{(i)}] A closed immersion of locally Noetherian formal
    preschemes is a finite morphism.
  \item[\rm{(ii)}] The composite of two finite morphisms of locally
    Noetherian formal preschemes is a finite morphism.
  \item[\rm{(iii)}] Let $\mathfrak{X}$, $\mathfrak{Y}$, and
    $\mathfrak{S}$ be locally Noetherian formal preschemes,
    $f:\mathfrak{X}\to\mathfrak{S}$ a finite morphism, and
    $g:\mathfrak{Y}\to\mathfrak{S}$ a morphism.  Then the projection
    \[
      \mathfrak{X}\times_{\mathfrak{S}}\mathfrak{Y}
      \longrightarrow\mathfrak{Y}
    \]
    is finite.
  \item[\rm{(iv)}] Let $\mathfrak{S}$ be a locally Noetherian formal
    prescheme, and let $\mathfrak{X}'$ and $\mathfrak{Y}'$ be
    locally Noetherian formal preschemes such that
    $\mathfrak{X}'\times_{\mathfrak{S}}\mathfrak{Y}'$ is locally
    Noetherian.  If $\mathfrak{X}$ and $\mathfrak{Y}$ are locally
    Noetherian formal $\mathfrak{S}$-preschemes and
    \[
      f:\mathfrak{X}\to\mathfrak{X}',\qquad
      g:\mathfrak{Y}\to\mathfrak{Y}'
    \]
    are finite $\mathfrak{S}$-morphisms, then
    $f\times_{\mathfrak{S}}g$ is finite.
  \item[\rm{(v)}] Let
    \[
      f:\mathfrak{X}\to\mathfrak{Y},\qquad
      g:\mathfrak{Y}\to\mathfrak{Z}
    \]
    be morphisms of locally Noetherian formal preschemes such that
    $g$ is of finite type and separated.  If $g\circ f$ is finite,
    then $f$ is finite.
\end{enumerate}
\end{proposition}

\begin{proof}
Assertion \emph{(i)} is immediate.  Using criterion \emph{a)} of
\sref{III.4.8.1}, the other assertions reduce at once to the
corresponding propositions for morphisms of ordinary preschemes
\sref[II]{II.6.1.5}.  We leave the details to the reader, following
the model of \sref[I]{I.10.13.5}.
\end{proof}

\begin{corollary}[4.8.4]
\label{III.4.8.4}
Under the hypotheses of \sref[I]{I.10.9.9}, if $f$ is a finite
morphism, then so is its extension $\widehat{f}$ to the
completions.
\end{corollary}

\begin{corollary}[4.8.5]
\label{III.4.8.5}
Let $\mathfrak{X}$ be a formal prescheme finite over
$\mathfrak{Y}$, and let
$f:\mathfrak{X}\to\mathfrak{Y}$ be the structural morphism.  For
every open subset $U\subset\mathfrak{Y}$, the formal prescheme
$f^{-1}(U)$ is finite over $U$.
\end{corollary}

\begin{proposition}[4.8.6]
\label{III.4.8.6}
If $f:\mathfrak{X}\to\mathfrak{Y}$ is a finite morphism of locally
Noetherian formal preschemes, then
$f_*(\sh{O}_{\mathfrak{X}})$ is a coherent
$\sh{O}_{\mathfrak{Y}}$-algebra.
\end{proposition}

\begin{proof}
We may regard $f$ as the inductive limit of an inductive system
$(f_n)$ of morphisms $f_n:X_n\to Y_n$.  Let us show that the
$f_n$ are finite and that $f_*(\sh{O}_{\mathfrak{X}})$ is
isomorphic to the projective limit of the
$(f_n)_*(\sh{O}_{X_n})$; this will prove the assertion
\sref[I]{I.10.10.5}.  It is enough to restrict to the case
\[
  \mathfrak{Y}=\Spf(B),\qquad \mathfrak{X}=\Spf(A),
\]
and to observe that, if $\mathfrak{R}$ is an ideal of definition
of $B$ and $A$ is a $B$-module of finite type, then
$A/\mathfrak{R}^{n+1}A$ is a module of finite type over
$B/\mathfrak{R}^{n+1}B$, and $A$ is the projective limit of the
$A/\mathfrak{R}^{n+1}A$.
\end{proof}

Conversely:

\begin{proposition}[4.8.7]
\label{III.4.8.7}
Let $\mathfrak{Y}$ be a locally Noetherian formal prescheme and
$\sh{A}$ a coherent $\sh{O}_{\mathfrak{Y}}$-algebra.  There exists
a formal prescheme $\mathfrak{X}$ finite over $\mathfrak{Y}$,
unique up to a unique $\mathfrak{Y}$-isomorphism, such that
\[
  f_*(\sh{O}_{\mathfrak{X}})=\sh{A},
\]
where $f:\mathfrak{X}\to\mathfrak{Y}$ is the structural morphism.
\end{proposition}

\begin{proof}
\oldpage[III]{148}
Let $\sh{K}$ be an ideal of definition of $\mathfrak{Y}$, and set
\[
  Y_n=(\mathfrak{Y},
       \sh{O}_{\mathfrak{Y}}/\sh{K}^{n+1}),
  \qquad
  \sh{A}_n=\sh{A}/\sh{K}^{n+1}\sh{A}.
\]
Clearly $\sh{A}_n$ is a finite $\sh{O}_{Y_n}$-algebra, and hence
defines a finite $Y_n$-prescheme
\[
  X_n=\Spec(\sh{A}_n)
\]
\sref[II]{II.6.1.3}.  For $m\leq n$, the canonical surjective
homomorphism
\[
  h_{mn}:\sh{A}_n\longrightarrow\sh{A}_m
\]
defines a morphism $u_{mn}:X_m\to X_n$ for which the diagram
\[
  \xymatrix{
    X_n\ar[d]_{f_n} &
    X_m\ar[l]_{u_{mn}}\ar[d]^{f_m}\\
    Y_n &
    Y_m\ar[l]
  }
\]
is commutative, where $f_n$ denotes the structural morphism; it
also identifies $X_m$ with
$X_n\times_{Y_n}Y_m$, as follows at once from
\sref[II]{II.1.4.6}.  The formal prescheme $\mathfrak{X}$ obtained
as the inductive limit of the inductive system $(X_n)$ is therefore
locally Noetherian, and its structural morphism
\[
  f:\mathfrak{X}\longrightarrow\mathfrak{Y},
\]
the inductive limit of the system $(f_n)$, is finite
\sref{III.4.8.1} \sref[I]{I.10.12.3.1}.  Moreover, the proof of
\sref{III.4.8.6} shows that
$f_*(\sh{O}_{\mathfrak{X}})$ is the projective limit of the
$\sh{A}_n$, and is therefore equal to $\sh{A}$
\sref[I]{I.10.10.6}.  The uniqueness assertion follows from the
more general result below.
\end{proof}

\begin{proposition}[4.8.8]
\label{III.4.8.8}
Let $\mathfrak{Y}$ be a locally Noetherian formal prescheme, and
let $\mathfrak{X}$ and $\mathfrak{X}'$ be formal
$\mathfrak{Y}$-preschemes finite over $\mathfrak{Y}$, with
structural morphisms
\[
  f:\mathfrak{X}\to\mathfrak{Y},\qquad
  f':\mathfrak{X}'\to\mathfrak{Y}.
\]
There is a canonical bijection
\[
  \Hom_{\mathfrak{Y}}(\mathfrak{X},\mathfrak{X}')
  \xrightarrow{\ \sim\ }
  \Hom_{\sh{O}_{\mathfrak{Y}}}
  \bigl(f'_*(\sh{O}_{\mathfrak{X}'}),
        f_*(\sh{O}_{\mathfrak{X}})\bigr)\footnotemark.
\]
\footnotetext{The latter expression denotes the set of homomorphisms of
$\sh{O}_{\mathfrak{Y}}$-algebras
$f'_*(\sh{O}_{\mathfrak{X}'})\to f_*(\sh{O}_{\mathfrak{X}})$.}
\end{proposition}

\begin{proof}
The map $h\mapsto\sh{A}(h)$ is defined as in
\sref[II]{II.1.1.2}.  To prove that it is bijective, we reduce at
once to the case in which
$\mathfrak{Y}=\Spf(B)$ is a Noetherian formal affine prescheme.
Then
\[
  \mathfrak{X}=\Spf(A),\qquad
  \mathfrak{X}'=\Spf(A'),
\]
where $A$ and $A'$ are finite $B$-algebras, and
\[
  f_*(\sh{O}_{\mathfrak{X}})=A^\Delta,\qquad
  f'_*(\sh{O}_{\mathfrak{X}'})=A'^\Delta.
\]
The conclusion follows from the bijective correspondence between
$\mathfrak{Y}$-morphisms
$\mathfrak{X}\to\mathfrak{X}'$ and continuous $B$-algebra
homomorphisms $A'\to A$ \sref[I]{I.10.2.2}, together with the
bijective correspondence between homomorphisms of $B$-modules
$A'\to A$ and homomorphisms of
$\sh{O}_{\mathfrak{Y}}$-modules
$A'^\Delta\to A^\Delta$ \sref[I]{I.10.10.2.3}.
\end{proof}

\begin{corollary}[4.8.9]
\label{III.4.8.9}
Under the canonical bijection of \sref{III.4.8.8}, closed
immersions $\mathfrak{X}\to\mathfrak{X}'$ correspond to
surjective homomorphisms of $\sh{O}_{\mathfrak{Y}}$-algebras
\[
  f'_*(\sh{O}_{\mathfrak{X}'})
  \longrightarrow
  f_*(\sh{O}_{\mathfrak{X}}).
\]
\end{corollary}

\begin{proof}
The question is local on $\mathfrak{Y}$, and hence reduces to the
definition of closed immersions of locally Noetherian formal
preschemes \sref[I]{I.10.14.2}.
\end{proof}

\begin{corollary}[4.8.10]
\label{III.4.8.10}
With the notation and hypotheses of \sref{III.4.8.1}, an adic
morphism $f$ is a closed immersion if and only if $f_0$ is a
closed immersion of ordinary preschemes.
\end{corollary}

\begin{proof}
This follows at once from \sref{III.4.8.9} and the surjectivity
criterion for a homomorphism of coherent
$\sh{O}_{\mathfrak{Y}}$-modules \sref[I]{I.10.11.5}.
\end{proof}

\begin{proposition}[4.8.11]
\label{III.4.8.11}
A morphism $f:\mathfrak{X}\to\mathfrak{Y}$ of locally Noetherian
formal preschemes
\oldpage[III]{149}
is finite if and only if it is proper and all its fibres
$f^{-1}(y)$, for $y\in\mathfrak{Y}$, are finite.
\end{proposition}

\begin{proof}
By the definitions \sref{III.3.4.1} and \sref{III.4.8.2}, the
assertion reduces immediately to the same assertion for $f_0$
(with the notation of \sref{III.4.8.1}), which is precisely
\sref{III.4.4.2}.
\end{proof}

\section{An existence theorem for coherent algebraic sheaves}
\label{section:III.5}


\subsection{Statement of the theorem}
\label{subsection:III.5.1}

\begin{env}[5.1.1]
\label{III.5.1.1}
Let $A$ be an adic Noetherian ring and $\mathfrak{J}$ an ideal of
definition of $A$, so that $A$ is separated and complete for the
$\mathfrak{J}$-adic topology.  If $Y=\Spec(A)$, the affine formal
prescheme $\Spf(A)$ is identified with the completion
$\widehat{Y}$ of $Y$ along the closed subset
$Y'=\mathrm{V}(\mathfrak{J})$ \sref[I]{I.10.10.1}.  Let $X$ be an
ordinary $Y$-prescheme of finite type, with structural morphism
$f:X\to Y$.  We denote by $\widehat{X}$ the completion of $X$ along
the closed subset $X'=f^{-1}(Y')$, or equivalently the formal
$\widehat{Y}$-prescheme $X\times_Y\widehat{Y}$, and by
$\widehat{f}:\widehat{X}\to\widehat{Y}$ the extension of $f$ to the
completions.  Finally, for every coherent $\sh{O}_X$-module
$\sh{F}$, we denote by $\widehat{\sh{F}}$ its completion
$\sh{F}_{/X'}$, which is a coherent
$\sh{O}_{\widehat{X}}$-module.
\end{env}

\begin{proposition}[5.1.2]
\label{III.5.1.2}
With the hypotheses and notation of \sref{III.5.1.1}, let $\sh{F}$
be a coherent $\sh{O}_X$-module whose support is proper over $Y$
\sref[II]{II.5.4.10}.  The canonical homomorphisms
\sref{III.4.1.4}
\[
  \rho_i:
  \HH^i(X,\sh{F})
  \longrightarrow
  \HH^i(\widehat{X},\widehat{\sh{F}})
\]
are isomorphisms.
\end{proposition}

\begin{proof}
The $A$-module $\HH^i(X,\sh{F})$ is of finite type
\sref{III.3.2.4}, and is therefore identical with its separated
completion for the $\mathfrak{J}$-preadic topology
\sref[0]{0.7.3.6}.  The proposition is consequently a special case
of \sref{III.4.1.10}.
\end{proof}

Recall that the canonical isomorphisms $\rho_i$ commute with the
connecting homomorphisms associated with every exact sequence
\[
  0\longrightarrow\sh{F}'
  \longrightarrow\sh{F}
  \longrightarrow\sh{F}''
  \longrightarrow0
\]
of coherent $\sh{O}_X$-modules
\sref[0]{0.12.1.6} \sref[I]{I.10.8.9}.

\begin{corollary}[5.1.3]
\label{III.5.1.3}
Let $\sh{F}$ and $\sh{G}$ be coherent $\sh{O}_X$-modules such that
the intersection of their supports is proper over $Y$.  The
canonical homomorphism
\[
  \Hom_{\sh{O}_X}(\sh{F},\sh{G})
  \longrightarrow
  \Hom_{\sh{O}_{\widehat{X}}}
  (\widehat{\sh{F}},\widehat{\sh{G}})
  \tag{5.1.3.1}
  \label{III.5.1.3.1.eq}
\]
which sends a homomorphism $u:\sh{F}\to\sh{G}$ to its completion
$\widehat{u}:\widehat{\sh{F}}\to\widehat{\sh{G}}$, is an
isomorphism.  Moreover, when $f$ is closed,
$\widehat{u}$ is injective (respectively, surjective) if and only
if $u$ is injective (respectively, surjective).
\end{corollary}

\begin{proof}
The first assertion is a special case of \sref{III.4.5.3}; it also
follows from the fact that the source of \eqref{III.5.1.3.1.eq} is
an $A$-module of finite type and is therefore identical with its
separated completion.  For the second assertion, note that, by
\sref[I]{I.10.8.14}, $\widehat{u}$ is injective (respectively,
surjective) if and only if there is a neighbourhood of $X'$ in
$X$ on which $u$ is injective (respectively, surjective).  The
conclusion therefore follows from the lemma below.
\end{proof}

\oldpage[III]{150}

\begin{lemma}[5.1.3.1]
\label{III.5.1.3.1}
Under the hypotheses of \sref{III.5.1.1}, suppose in addition that
$f$ is closed.  Every neighbourhood of $X'$ in $X$ is then equal
to $X$.
\end{lemma}

\begin{proof}
We first reduce to the case $f(X)=Y$.  Indeed, by hypothesis,
$f(X)$ is a closed subset $Z$ of $Y$.  We may replace $f$ by
$f_{\mathrm{red}}$ \sref[I]{I.6.3.4}, and hence suppose that $X$
and $Y$ are reduced; we may then replace $Y$ by the reduced closed
subprescheme having $Z$ as its underlying space
\sref[I]{I.5.2.2}, since every ideal of $A$ is closed and every
quotient ring of $A$ is therefore adic and Noetherian.

We then have $f(X')=Y'$.  If $V$ is an open neighbourhood of $X'$
in $X$, the set $f(X-V)$ is closed in $Y$ and does not meet $Y'$.
This is impossible unless $X-V$ is empty, because $\mathfrak{J}$
is contained in the radical of $A$
\sref[I]{I.1.1.15} \sref[0]{0.7.1.10}.  Thus $V=X$.
\end{proof}

If we restrict attention to coherent $\sh{O}_X$-modules whose
support is proper over $Y$, \sref{III.5.1.3} says, in categorical
language, that the functor
\[
  \sh{F}\longmapsto\widehat{\sh{F}}
\]
is fully faithful from the category of such $\sh{O}_X$-modules to
the category of coherent $\sh{O}_{\widehat{X}}$-modules.  It thus
gives an equivalence of the former category with a full
subcategory of the latter \sref[0]{0.8.1.6}.  The existence theorem
will show that, when $X$ is proper over $Y$, this subcategory is in
fact the category of all coherent
$\sh{O}_{\widehat{X}}$-modules.  More precisely:

\begin{theorem}[5.1.4]
\label{III.5.1.4}
Let $A$ be an adic Noetherian ring, let $Y=\Spec(A)$, let
$\mathfrak{J}$ be an ideal of definition of $A$, and put
$Y'=\mathrm{V}(\mathfrak{J})$.  Let $f:X\to Y$ be a separated
morphism of finite type and put $X'=f^{-1}(Y')$.  Let
\[
  \widehat{Y}=Y_{/Y'}=\Spf(A),\qquad
  \widehat{X}=X_{/X'}
\]
be the completions of $Y$ and $X$ along $Y'$ and $X'$, and let
$\widehat{f}:\widehat{X}\to\widehat{Y}$ be the extension of $f$ to
the completions.  Then the functor
\[
  \sh{F}\longmapsto\sh{F}_{/X'}=\widehat{\sh{F}}
\]
is an equivalence from the category of coherent
$\sh{O}_X$-modules whose support is proper over $\Spec(A)$ to the
category of coherent $\sh{O}_{\widehat{X}}$-modules whose support
is proper over $\Spf(A)$.
\end{theorem}

In other words, in view of \sref{III.5.1.3}:

\begin{corollary}[5.1.5]
\label{III.5.1.5}
An $\sh{O}_{\widehat{X}}$-module is isomorphic to the completion of
a coherent $\sh{O}_X$-module whose support is proper over
$\Spec(A)$ if and only if it is coherent and its support is proper
over $\Spf(A)$.
\end{corollary}

The most important case is the following:

\begin{corollary}[5.1.6]
\label{III.5.1.6}
Suppose that $X$ is proper over $Y=\Spec(A)$.  Then the functor
$\sh{F}\mapsto\widehat{\sh{F}}$ is an equivalence between the
category of coherent $\sh{O}_X$-modules and the category of
coherent $\sh{O}_{\widehat{X}}$-modules.
\end{corollary}

\begin{env}[5.1.7]
\label{III.5.1.7}
\emph{Scholium.}
Taking account of the characterization of coherent sheaves on
formal preschemes \sref[I]{I.10.11.3}, we see that, under the
conditions of \sref{III.5.1.1}, giving a coherent
$\sh{O}_X$-module whose support is proper over $\Spec(A)$ is
equivalent---if
\[
  Y_n=\Spec(A/\mathfrak{J}^{n+1}),\qquad
  X_n=X\times_Y Y_n
\]
---to giving a projective system $(\sh{F}_n)$ of coherent
$\sh{O}_{X_n}$-modules such that, for $m\leq n$,
\[
  \sh{F}_m
  =
  \sh{F}_n\otimes_{\sh{O}_{Y_n}}\sh{O}_{Y_m},
  \qquad\text{or equivalently}\qquad
  \sh{F}_m
  =
  \sh{F}_n/\mathfrak{J}^{m+1}\sh{F}_n,
\]
and such that the support of $\sh{F}_0$ is a subset of $X_0$
proper over $Y_0$.  By \sref[I]{I.10.11.4}, homomorphisms of
coherent $\sh{O}_X$-modules are interpreted similarly as
homomorphisms of projective systems of coherent
$\sh{O}_{X_n}$-modules.

\oldpage[III]{151}

In all known applications, $A$ is in fact a local adic Noetherian
ring.  Thus the $Y_n$ are spectra of local Artinian rings, and the
results of this subsection and of the preceding ones reduce, to a
large extent, algebraic geometry over a local adic Noetherian ring
to algebraic geometry over local Artinian rings.
\end{env}

\begin{corollary}[5.1.8]
\label{III.5.1.8}
Under the hypotheses of \sref{III.5.1.4}, the map
\[
  Z\longmapsto\widehat{Z}=Z_{/(Z\cap X')}
\]
is a bijection from the set of closed subpreschemes $Z$ of $X$
which are proper over $Y$ onto the set of closed formal
subpreschemes of $\widehat{X}$ which are proper over
$\widehat{Y}$.
\end{corollary}

\begin{proof}
A closed formal subprescheme of $\widehat{X}$ has the form
\[
  \bigl(T,
    (\sh{O}_{\widehat{X}}/\sh{A})|_T\bigr),
\]
where $\sh{A}$ is a coherent ideal of
$\sh{O}_{\widehat{X}}$ \sref[I]{I.10.14.2}.  If $T$ is proper over
$\widehat{Y}$, \sref{III.5.1.4} shows that
$\sh{O}_{\widehat{X}}/\sh{A}$ is isomorphic to an
$\sh{O}_{\widehat{X}}$-module of the form
$\widehat{\sh{F}}$, where $\sh{F}$ is a coherent
$\sh{O}_X$-module whose support is proper over $Y$.  Moreover,
\sref{III.5.1.3} shows that the canonical homomorphism
\[
  \sh{O}_{\widehat{X}}
  \longrightarrow
  \sh{O}_{\widehat{X}}/\sh{A}
\]
has the form $\widehat{u}$, where
$u:\sh{O}_X\to\sh{F}$ is a surjective homomorphism of
$\sh{O}_X$-modules.  Consequently
$\sh{F}=\sh{O}_X/\sh{N}$ for a coherent ideal $\sh{N}$ of
$\sh{O}_X$, and
$\sh{A}=\widehat{\sh{N}}$ \sref[I]{I.10.8.8}.  The conclusion now
follows from \sref[I]{I.10.14.7}.
\end{proof}


\subsection{Proof of the existence theorem: the projective and quasi-projective case}
\label{subsection:III.5.2}

\begin{env}[5.2.1]
\label{III.5.2.1}
Under the conditions of \sref{III.5.1.4}, we shall say
provisionally that a coherent $\sh{O}_{\widehat{X}}$-module is
\emph{algebraizable} if it is isomorphic to the completion
$\widehat{\sh{F}}$ of a coherent $\sh{O}_X$-module $\sh{F}$ whose
support is proper over $Y$.
\end{env}

\begin{lemma}[5.2.2]
\label{III.5.2.2}
Let $\sh{F}'$ and $\sh{G}'$ be algebraizable
$\sh{O}_{\widehat{X}}$-modules.  For every homomorphism
$u:\sh{F}'\to\sh{G}'$, the modules
$\Ker(u)$, $\Im(u)$, and $\Coker(u)$ are algebraizable.
\end{lemma}

\begin{proof}
Write
\[
  \sh{F}'=\widehat{\sh{F}},\qquad
  \sh{G}'=\widehat{\sh{G}},
\]
where $\sh{F}$ and $\sh{G}$ are coherent $\sh{O}_X$-modules whose
supports are proper over $Y$.  By \sref{III.5.1.3},
$u=\widehat{v}$ for a homomorphism
$v:\sh{F}\to\sh{G}$.  Since completion is exact,
$\Ker(\widehat{v})\simeq\widehat{\Ker(v)}$; and because the support
of $\Ker(v)$ is contained in that of $\sh{F}$, it follows that
$\Ker(u)$ is algebraizable.  The arguments for $\Im(u)$ and
$\Coker(u)$ are the same.
\end{proof}

\begin{proposition}[5.2.3]
\label{III.5.2.3}
Let $A$ be an adic Noetherian ring, $\mathfrak{J}$ an ideal of
definition of $A$, $\mathfrak{Y}=\Spf(A)$, and
$f:\mathfrak{X}\to\mathfrak{Y}$ a proper morphism of formal
preschemes.  Put
\[
  Y_k=\Spec(A/\mathfrak{J}^{k+1}),\qquad
  X_k=\mathfrak{X}\times_{\mathfrak{Y}}Y_k,
\]
and, for every $\sh{O}_{\mathfrak{X}}$-module $\sh{F}$, put
\[
  \sh{F}_k
  =
  \sh{F}\otimes_{\sh{O}_{\mathfrak{X}}}\sh{O}_{X_k}
  =
  \sh{F}/\mathfrak{J}^{k+1}\sh{F}.
\]
Let $\sh{L}$ be an invertible $\sh{O}_{\mathfrak{X}}$-module, and
suppose that
$\sh{L}_0=\sh{L}/\mathfrak{J}\sh{L}$ is an ample
$\sh{O}_{X_0}$-module.  For every
$\sh{O}_{\mathfrak{X}}$-module $\sh{F}$ and every integer $n$, set
\[
  \sh{F}(n)=\sh{F}\otimes\sh{L}^{\otimes n}.
\]
Then, for every coherent $\sh{O}_{\mathfrak{X}}$-module $\sh{F}$,
there is an integer $n_0$ such that, for every $n\geq n_0$:
\begin{enumerate}
  \item[\rm{(i)}] the canonical homomorphism
    \[
      \HH^0(\mathfrak{X},\sh{F}(n))
      \longrightarrow
      \HH^0(\mathfrak{X},\sh{F}_k(n))
    \]
    is surjective for every $k\geq0$;
  \item[\rm{(ii)}]
    $\HH^q(\mathfrak{X},\sh{F}(n))=0$ for every $q>0$.
\end{enumerate}
\end{proposition}

\begin{proof}
The underlying spaces of $\mathfrak{X}$ and $X_0$ are the same.
The sheaves
\[
  \sh{M}_k
  =
  \mathfrak{J}^k\sh{F}/\mathfrak{J}^{k+1}\sh{F},
\]
being annihilated by $\mathfrak{J}$, may be regarded as coherent
$\sh{O}_{X_0}$-modules \sref[0]{0.5.3.10}.  Moreover, if
\[
  \sh{M}_k(n)
  =
  \sh{M}_k\otimes_{\sh{O}_{X_0}}\sh{L}_0^{\otimes n},
\]
then
\[
  \sh{M}_k(n)
  =
  \mathfrak{J}^k\sh{F}(n)/
  \mathfrak{J}^{k+1}\sh{F}(n).
\]
Since $\sh{L}_0$ is ample for $f_0:X_0\to Y_0$ and $f_0$ is
proper,
\oldpage[III]{152}
$f_0$ is projective \sref[II]{II.5.5.4}.  Let
\[
  S=\bigoplus_{k\geq0}
  \mathfrak{J}^k/\mathfrak{J}^{k+1}
\]
be the graded $A$-algebra associated with the
$\mathfrak{J}$-adic filtration of $A$.  It is of finite type
because $A$ is Noetherian.  Put
\[
  \sh{S}'=f_0^*(\widetilde{S}),\qquad
  \sh{M}=\bigoplus_{k\geq0}\sh{M}_k.
\]
Then $\sh{S}'$ is a quasi-coherent $\sh{O}_{X_0}$-algebra of
finite type, and $\sh{M}$ is a quasi-coherent graded
$\sh{S}'$-module of finite type, since $\sh{F}_0$ is coherent and
generates $\sh{M}$ as an $\sh{S}'$-module.  Theorem
\sref{III.2.4.1},~\emph{(ii)}, therefore gives an integer $n_0$
such that, for $n\geq n_0$ and every $k$,
\[
  \HH^q(X_0,\sh{M}_k(n))=0
  \qquad(q>0).
  \tag{5.2.3.1}
  \label{III.5.2.3.1}
\]
Regarded now as an $\sh{O}_{\mathfrak{X}}$-module,
$\sh{M}_k(n)$ consequently also satisfies
\[
  \HH^q(\mathfrak{X},\sh{M}_k(n))=0
  \qquad(q>0,\ n\geq n_0).
\]

Apply the exact cohomology sequence to
\[
  0\longrightarrow
  \frac{\mathfrak{J}^k\sh{F}(n)}
       {\mathfrak{J}^{k+1}\sh{F}(n)}
  \longrightarrow
  \frac{\mathfrak{J}^h\sh{F}(n)}
       {\mathfrak{J}^{k+1}\sh{F}(n)}
  \longrightarrow
  \frac{\mathfrak{J}^h\sh{F}(n)}
       {\mathfrak{J}^{k}\sh{F}(n)}
  \longrightarrow0.
\]
For $0\leq h<k$, $n\geq n_0$, and $q>0$, induction on $k-h$
gives
\[
  \HH^q\left(
    \mathfrak{X},
    \frac{\mathfrak{J}^h\sh{F}(n)}
         {\mathfrak{J}^{k}\sh{F}(n)}
  \right)=0,
  \tag{5.2.3.2}
  \label{III.5.2.3.2}
\]
and in particular, for $h=0$,
\[
  \HH^q(\mathfrak{X},\sh{F}_k(n))=0
  \qquad(n\geq n_0,\ k\geq0,\ q>0).
  \tag{5.2.3.3}
  \label{III.5.2.3.3}
\]
Another part of the exact cohomology sequence, with $h=0$, gives
\[
  \HH^0(\mathfrak{X},\sh{F}_{k+1}(n))
  \longrightarrow
  \HH^0(\mathfrak{X},\sh{F}_k(n))
  \longrightarrow
  \HH^1\left(
    \mathfrak{X},
    \frac{\mathfrak{J}^k\sh{F}(n)}
         {\mathfrak{J}^{k+1}\sh{F}(n)}
  \right)
  =0.
  \tag{5.2.3.4}
  \label{III.5.2.3.4}
\]
It follows that, for $h\leq k$, the canonical map
\[
  \HH^0(\mathfrak{X},\sh{F}_k(n))
  \longrightarrow
  \HH^0(\mathfrak{X},\sh{F}_h(n))
  \tag{5.2.3.5}
  \label{III.5.2.3.5}
\]
is surjective.  Thus, for every $q$, the projective system
\[
  \bigl(\HH^q(\mathfrak{X},\sh{F}_k(n))\bigr)_{k\geq0}
\]
satisfies condition (ML) when $n\geq n_0$.

Furthermore, every formal affine open $U$ of $\mathfrak{X}$ is an
affine open in each $X_k$ \sref[I]{I.10.5.2}.  Hence
$\HH^q(U,\sh{F}_k(n))=0$ for $q>0$ \sref[I]{I.1.3.1}, and
\[
  \HH^0(U,\sh{F}_k(n))
  \longrightarrow
  \HH^0(U,\sh{F}_h(n))
\]
is surjective for $h\leq k$ \sref[I]{I.1.3.9}.  The hypotheses of
\sref[0]{0.13.3.1} are therefore satisfied.  For $n\geq n_0$ we
obtain:
\begin{enumerate}
  \item[$1^\circ$] for every $q>0$, the map
    \[
      \HH^q(\mathfrak{X},\sh{F}(n))
      \longrightarrow
      \varprojlim_k
      \HH^q(\mathfrak{X},\sh{F}_k(n))
    \]
    is bijective; hence, by \eqref{III.5.2.3.3},
    $\HH^q(\mathfrak{X},\sh{F}(n))=0$;
  \item[$2^\circ$] the homomorphism
    \[
      \HH^0(\mathfrak{X},\sh{F}(n))
      \longrightarrow
      \varprojlim_k
      \HH^0(\mathfrak{X},\sh{F}_k(n))
    \]
    is bijective.  Since the homomorphisms
    \eqref{III.5.2.3.5} are surjective, so is each homomorphism
    \[
      \varprojlim_k
      \HH^0(\mathfrak{X},\sh{F}_k(n))
      \longrightarrow
      \HH^0(\mathfrak{X},\sh{F}_h(n)).
    \]
\end{enumerate}
This proves both assertions.
\end{proof}

\begin{corollary}[5.2.4]
\label{III.5.2.4}
Under the hypotheses of \sref{III.5.2.3}, for every coherent
$\sh{O}_{\mathfrak{X}}$-module $\sh{F}$ there is an integer $N$
such that, for $n\geq N$, $\sh{F}(n)$ is generated by its sections
over
\oldpage[III]{153}
$\mathfrak{X}$.  Equivalently, $\sh{F}$ is isomorphic to a quotient
of an $\sh{O}_{\mathfrak{X}}$-module of the form
$(\sh{O}_{\mathfrak{X}}(-n))^k$.
\end{corollary}

\begin{proof}
Since $X_0$ is Noetherian, the hypothesis on $\sh{L}_0$ and
\sref[II]{II.4.5.5} give an integer $n_0$ such that
$\sh{F}_0(n)$ is generated by its sections over $\mathfrak{X}$ for
$n\geq n_0$.  We may take $n_0$ large enough that
\[
  \Gamma(\mathfrak{X},\sh{F}(n))
  \longrightarrow
  \Gamma(\mathfrak{X},\sh{F}_0(n))
\]
is also surjective for $n\geq n_0$ \sref{III.5.2.3}.  There are
therefore finitely many sections
$s_i\in\Gamma(\mathfrak{X},\sh{F}(n))$ whose images generate
$\sh{F}_0(n)$ \sref[0]{0.5.2.3}.  Since $\mathfrak{J}$ is
contained in the maximal ideal of the local ring at every point
of $\mathfrak{X}$, Nakayama's lemma applied to those local rings
shows that the $s_i$ generate $\sh{F}(n)$
\sref[0]{0.5.1.1}.
\end{proof}

\begin{env}[5.2.5]
\label{III.5.2.5}
\emph{Proof of the existence theorem: the projective case.}
With the notation of \sref{III.5.1.4}, suppose that $f$ is
projective, so that there is an ample $\sh{O}_X$-module $\sh{L}$
\sref[I]{I.5.5.4}.  By definition,
\[
  X_n=\widehat{X}\times_{\widehat{Y}}Y_n
\]
is the closed subprescheme
$X\times_Y Y_n=f^{-1}(Y_n)$ of $X$.  If
\[
  \sh{L}'=\sh{L}\otimes_{\sh{O}_X}\sh{O}_{\widehat{X}}
\]
is the completion of $\sh{L}$, then
$\sh{L}'_0=\sh{L}/\mathfrak{J}\sh{L}$, regarded as an
$\sh{O}_{X_0}$-module; this module is ample
\sref[II]{II.4.6.13}[(i~\emph{bis})].

We may therefore apply \sref{III.5.2.4} to $\sh{L}'$ and to any
coherent $\sh{O}_{\widehat{X}}$-module $\sh{F}$.  It follows that
$\sh{F}$ is isomorphic to a quotient of
\[
  \sh{G}=
  \bigl(\sh{L}'^{\otimes(-n)}\bigr)^k
\]
for suitable integers $n>0$ and $k>0$.  The module $\sh{G}$ is the
completion of $(\sh{L}^{\otimes(-n)})^k$
\sref[I]{I.10.8.10}, and is therefore algebraizable.  Let
$u:\sh{G}\to\sh{F}$ be the canonical surjection, and put
$\sh{H}=\Ker(u)$; this is a coherent
$\sh{O}_{\widehat{X}}$-module \sref[0]{0.5.3.4}.  Applying the same
argument to $\sh{H}$, we find an algebraizable
$\sh{O}_{\widehat{X}}$-module $\sh{K}$ and a homomorphism
$v:\sh{K}\to\sh{G}$ such that $\sh{H}=\Im(v)$.  Thus
$\sh{F}=\Coker(v)$, and $\sh{F}$ is algebraizable by
\sref{III.5.2.2}.
\end{env}

\begin{env}[5.2.6]
\label{III.5.2.6}
\emph{Proof of the existence theorem: the quasi-projective case.}
Retain the notation of \sref{III.5.1.4}, and now suppose that $f$
is quasi-projective.  There is a projective morphism
$g:Z\to Y$ such that $X$ is identified with the $Y$-prescheme
induced on an open subset of $Z$ \sref[II]{II.5.3.2}.  If
$Z'=g^{-1}(Y')$, then $X'=X\cap Z'$.  Consequently the completion
$\widehat{X}=X_{/X'}$ is identified with the formal prescheme
induced by $\widehat{Z}=Z_{/Z'}$ on the open subset $X\cap Z'$ of
$\widehat{Z}$ \sref[I]{I.10.8.5}.

Let $\sh{F}'$ be a coherent $\sh{O}_{\widehat{X}}$-module whose
support $T'$ is proper over $\widehat{Y}$.  By definition, there
is a closed subprescheme of $X'$ with underlying space
$T'\subset X'$ such that the restriction $T'\to Y'$ of $f$ is
proper.  Hence $T'$ is proper over $Y$, and is therefore closed in
$Z'$ \sref[II]{II.5.4.10}.  It follows that $\sh{F}'$ is induced
on $\widehat{X}$ by the $\sh{O}_{\widehat{Z}}$-module $\sh{G}'$
obtained by gluing $\sh{F}'$, on the open $\widehat{X}$, with the
zero sheaf on the open $\widehat{Z}-T'$; the two sheaves agree on
their intersection $\widehat{X}-T'$.

The module $\sh{G}'$ is coherent.  By \sref{III.5.2.5}, there is a
coherent $\sh{O}_Z$-module $\sh{G}$ such that
$\sh{G}'=\widehat{\sh{G}}$.  Let $T$ be the support of $\sh{G}$;
then $T'=T\cap Z'$ \sref[I]{I.10.8.12}.  If $h$ is the restriction
of $g$ to the reduced closed subprescheme of $Z$ having $T$ as
its underlying space, then
\[
  T'=h^{-1}(Y')=T\cap g^{-1}(Y').
\]
Thus $X\cap T$ is an open subset of $T$ containing $T'$.

\oldpage[III]{154}
Since $h$ is proper \sref[II]{II.5.4.2}, and hence closed,
\sref{III.5.1.3.1} gives $X\cap T=T$.  Thus $T\subset X$; and
since $T$ is closed in $Z$, it is proper over $Y$.  If $\sh{F}$ is
the sheaf induced on $X$ by $\sh{G}$, its completion
$\widehat{\sh{F}}$ is induced on $\widehat{X}$ by
$\widehat{\sh{G}}$ \sref[I]{I.10.8.4}, and is therefore equal to
$\sh{F}'$.  This completes the proof.
\end{env}


\subsection{Proof of the existence theorem: the general case}
\label{subsection:III.5.3}

\begin{lemma}[5.3.1]
\label{III.5.3.1}
Under the conditions of \sref{III.5.1.4}, let
\[
  0\longrightarrow\sh{H}\longrightarrow\sh{F}
  \longrightarrow\sh{G}\longrightarrow0
\]
be an exact sequence of coherent
$\sh{O}_{\widehat{X}}$-modules.  If $\sh{G}$ and $\sh{H}$ are
algebraizable, then $\sh{F}$ is algebraizable.
\end{lemma}

\begin{proof}
Suppose that
\[
  \sh{G}=\widehat{\sh{B}},
  \qquad
  \sh{H}=\widehat{\sh{C}},
\]
where $\sh{B}$ and $\sh{C}$ are coherent $\sh{O}_X$-modules whose
supports are proper over $Y$.  The module $\sh{F}$ canonically
defines an element of the $A$-module
\[
  \Ext^1_{\sh{O}_{\widehat{X}}}
  (\widehat{X};\widehat{\sh{B}},\widehat{\sh{C}})
\]
\sref[0]{0.12.3.2}.  The hypotheses imply that this $A$-module is
canonically isomorphic to
\[
  \Ext^1_{\sh{O}_X}(X;\sh{B},\sh{C})
\]
\sref{III.4.5.2}.  There is therefore an exact sequence
\[
  0\longrightarrow\sh{C}\longrightarrow\sh{A}
  \longrightarrow\sh{B}\longrightarrow0
\]
of coherent $\sh{O}_X$-modules whose canonical class maps to the
class corresponding to $\sh{F}$.  By definition, taking account of
\sref[I]{I.10.8.8}, (ii), this says that $\sh{F}$ is isomorphic to
$\widehat{\sh{A}}$.  The lemma follows, since
$\Supp(\sh{A})$ is contained in
$\Supp(\sh{B})\cup\Supp(\sh{C})$ and is therefore proper over $Y$.
\end{proof}

\begin{corollary}[5.3.2]
\label{III.5.3.2}
Under the conditions of \sref{III.5.1.1}, let
$u:\sh{F}\to\sh{G}$ be a homomorphism of coherent
$\sh{O}_{\widehat{X}}$-modules.  If $\sh{G}$, $\Ker(u)$, and
$\Coker(u)$ are algebraizable, then $\sh{F}$ is algebraizable.
\end{corollary}

\begin{proof}
Lemma \sref{III.5.2.2}, applied to
$\sh{G}\to\Coker(u)$, shows that $\Im(u)$ is algebraizable.  It then
suffices to apply Lemma \sref{III.5.3.1} to the exact sequence
\[
  0\longrightarrow\Ker(u)\longrightarrow\sh{F}
  \longrightarrow\Im(u)\longrightarrow0.
\]
\end{proof}

\begin{lemma}[5.3.3]
\label{III.5.3.3}
Under the conditions of \sref{III.5.1.1}, let $h:Z\to Y$ be a
morphism of finite type, let $\widehat{Z}$ be the completion of $Z$
along $Z'=h^{-1}(Y')$, let $g:Z\to X$ be a proper $Y$-morphism,
and let $\widehat{g}:\widehat{Z}\to\widehat{X}$ be its extension to
the completions.  For every algebraizable
$\sh{O}_{\widehat{Z}}$-module $\sh{F}'$, the module
$\widehat{g}_*(\sh{F}')$ is an algebraizable
$\sh{O}_{\widehat{X}}$-module.
\end{lemma}

\begin{proof}
If $\sh{F}$ is a coherent $\sh{O}_Z$-module such that
$\sh{F}'=\widehat{\sh{F}}$, the first comparison theorem
\sref{III.4.1.5} identifies
$\widehat{g}_*(\widehat{\sh{F}})$ with the completion of
$g_*(\sh{F})$.
\end{proof}

\begin{lemma}[5.3.4]
\label{III.5.3.4}
Let $X$ be an ordinary Noetherian scheme, let $X'$ be a closed
subset of $X$, let $f:Z\to X$ be a proper morphism, and put
\[
  Z'=f^{-1}(X'),\qquad
  \widehat{X}=X_{/X'},\qquad
  \widehat{Z}=Z_{/Z'}.
\]
Let $\widehat{f}:\widehat{Z}\to\widehat{X}$ be the extension of $f$
to the completions.  Let $\sh{M}$ be a coherent ideal of
$\sh{O}_X$ such that, if
\[
  U=X-\Supp(\sh{O}_X/\sh{M}),
\]
the restriction $f^{-1}(U)\to U$ of $f$ is an isomorphism.  Then,
for every coherent $\sh{O}_{\widehat{X}}$-module $\sh{F}$, there is
an integer $n>0$ such that the kernel and cokernel of the canonical
homomorphism
\[
  \rho_{\sh{F}}:
  \sh{F}\longrightarrow
  \widehat{f}_*\bigl(\widehat{f}^{\,*}(\sh{F})\bigr)
\]
\sref[0]{0.4.4.3} are annihilated by
$\widehat{\sh{M}}^{\,n}$.
\end{lemma}

\begin{proof}
We may restrict to the case $X=\Spec(B)$, where $B$ is a
Noetherian ring; then $X'=\mathrm{V}(\mathfrak{K})$ for an ideal
$\mathfrak{K}$ of $B$.  We first reduce to the case in which $B$ is
adic Noetherian and $\mathfrak{K}$ is an ideal of definition.

Let $B_1$ be the separated completion of $B$ for the
$\mathfrak{K}$-preadic topology.
\oldpage[III]{155}
If $\mathfrak{K}_1=\mathfrak{K}B_1$, then $B_1$ is an adic
Noetherian ring and $\mathfrak{K}_1$ is an ideal of definition.
Put $X_1=\Spec(B_1)$, and let $h:X_1\to X$ be the morphism
corresponding to the canonical homomorphism $B\to B_1$.  If
$X'_1=h^{-1}(X')$, then $X'_1=\mathrm{V}(\mathfrak{K}_1)$.  Finally
put
\[
  Z_1=Z\times_X X_1,\qquad f_1:Z_1\longrightarrow X_1.
\]
The morphism $f_1$ is proper \sref[II]{II.5.4.2}.  Denote by
$\widehat{X}_1$ the completion of $X_1$ along $X'_1$, by
\[
  \widehat{Z}_1=
  Z_1\times_{X_1}\widehat{X}_1
\]
the completion of $Z_1$ along
$Z'_1=f_1^{-1}(X'_1)$, and by
$\widehat{f}_1$ the extension of $f_1$ to the completions.

The extension
$\widehat{h}:\widehat{X}_1\to\widehat{X}$ corresponds to the
identity map of $B_1$ and is therefore an isomorphism
\sref[I]{I.10.9.1}.  The corresponding morphism
$\widehat{Z}_1\to\widehat{Z}$ is consequently also an isomorphism,
and these isomorphisms identify $\widehat{f}_1$ with
$\widehat{f}$.  Moreover,
$\sh{M}_1=h^*(\sh{M})$ is a coherent ideal of $\sh{O}_{X_1}$ and
\[
  \Supp(\sh{O}_{X_1}/\sh{M}_1)
  =h^{-1}\bigl(\Supp(\sh{O}_X/\sh{M})\bigr)
\]
\sref[I]{I.9.1.13}.  Thus, if
$U_1=X_1-\Supp(\sh{O}_{X_1}/\sh{M}_1)$, then
$U_1=h^{-1}(U)$, and the restriction
$f_1^{-1}(U_1)\to U_1$ is an isomorphism
\sref[I]{I.3.2.7}.  Finally, $\widehat{\sh{M}}$ and
$\widehat{\sh{M}}_1$ are identified by $\widehat{h}$
\sref[I]{I.10.9.5}.  All the hypotheses of the lemma are therefore
satisfied by $X_1$, $X'_1$, $f_1$, and $\sh{M}_1$; henceforth we
may suppose that $B$ is adic Noetherian and that $\mathfrak{K}$ is
an ideal of definition.

We then have $\widehat{X}=\Spf(B)$ and
$\sh{F}=N^\Delta$, where $N$ is a finite-type $B$-module.  Thus
$\sh{F}=\widehat{\sh{G}}$, where $\sh{G}$ is the coherent
$\sh{O}_X$-module $\widetilde{N}$ \sref[I]{I.10.10.5}, and
\[
  \widehat{f}^{\,*}(\sh{F})
  =\widehat{f^*(\sh{G})}
\]
\sref[I]{I.10.9.5}.  By the first comparison theorem
\sref{III.4.1.5},
\[
  \widehat{f}_*
  \bigl(\widehat{f^*(\sh{G})}\bigr)
  \simeq
  \widehat{f_*\bigl(f^*(\sh{G})\bigr)},
\]
and, by \sref{III.5.1.3}, the canonical homomorphism
$\rho_{\sh{F}}$ is precisely $\widehat{\rho}_{\sh{G}}$.

The kernel $\sh{P}$ and cokernel $\sh{R}$ of
\[
  \rho_{\sh{G}}:
  \sh{G}\longrightarrow f_*\bigl(f^*(\sh{G})\bigr)
\]
are coherent $\sh{O}_X$-modules, and their restrictions to $U$ are
zero.  There is therefore an integer $n>0$ such that
\[
  \sh{M}^n\sh{P}=\sh{M}^n\sh{R}=0
\]
\sref[I]{I.9.3.4}.  It follows that
\[
  \widehat{\sh{M}}^{\,n}\widehat{\sh{P}}
  =
  \widehat{\sh{M}}^{\,n}\widehat{\sh{R}}
  =0
\]
\sref[I]{I.10.8.8} \sref[I]{I.10.8.10}, which proves the lemma.
\end{proof}

\begin{env}[5.3.5]
\label{III.5.3.5}
\emph{End of the proof of the existence theorem.}
Retain the hypotheses of \sref{III.5.1.4}.  We use Noetherian
induction \sref[0]{0.2.2.2}, assuming the theorem known for every
closed subprescheme $T$ of $X$ whose underlying space is distinct
from $X$; here $\widehat{T}$ denotes the completion of $T$ along
$T\cap X'$.  We may suppose that $X$ is nonempty.

Since $f$ is separated and of finite type, Chow's lemma
\sref[II]{II.5.6.1} gives a $Y$-scheme $Z$ and a $Y$-morphism
$g:Z\to X$ such that the structural morphism $h:Z\to Y$ is
quasi-projective, $g$ is projective and surjective, and there is a
nonempty open subset $U$ of $X$ for which
$g^{-1}(U)\to U$ is an isomorphism.  Let $\sh{M}$ be a coherent
ideal of $\sh{O}_X$ defining a closed subprescheme with underlying
space $X-U$ \sref[I]{I.5.2.2}, and let $\sh{F}$ be a coherent
$\sh{O}_{\widehat{X}}$-module whose support $E$ is proper over $Y$.
Denote by $\widehat{Z}$ the completion of $Z$ along
$h^{-1}(Y')$, and by
$\widehat{g}:\widehat{Z}\to\widehat{X}$ the extension of $g$ to
the completions.

The module $\widehat{g}^{\,*}(\sh{F})$ is a coherent
$\sh{O}_{\widehat{Z}}$-module whose support is contained in
$g^{-1}(E)$ and is therefore proper over $Y$, since $g$ is
projective and hence proper \sref[II]{II.5.4.6}.  Since $h$ is
quasi-projective, $\widehat{g}^{\,*}(\sh{F})$ is algebraizable by
\sref{III.5.2.6}.
\oldpage[III]{156}
It follows from \sref{III.5.3.3}, since $g$ is proper, that
\[
  \widehat{g}_*
  \bigl(\widehat{g}^{\,*}(\sh{F})\bigr)
\]
is an algebraizable $\sh{O}_{\widehat{X}}$-module.

Apply \sref{III.5.3.4} to $\sh{F}$ and $g$.  The kernel $\sh{P}$
and cokernel $\sh{R}$ of
\[
  \rho_{\sh{F}}:
  \sh{F}\longrightarrow
  \widehat{g}_*\bigl(\widehat{g}^{\,*}(\sh{F})\bigr)
\]
are annihilated by a power $\widehat{\sh{M}}^{\,n}$.  Let $T$ be
the closed subprescheme of $X$ defined by $\sh{M}^n$, whose
underlying space is $X-U$, and let $j:T\to X$ be the canonical
injection.  Its extension
$\widehat{j}:\widehat{T}\to\widehat{X}$ is the canonical injection
\sref[I]{I.10.14.7}, and we may write
\[
  \sh{P}=
  \widehat{j}_*\bigl(\widehat{j}^{\,*}(\sh{P})\bigr),
  \qquad
  \sh{R}=
  \widehat{j}_*\bigl(\widehat{j}^{\,*}(\sh{R})\bigr).
\]
Since $U$ is nonempty, the induction hypothesis shows that
$\widehat{j}^{\,*}(\sh{P})$ and
$\widehat{j}^{\,*}(\sh{R})$ are algebraizable
$\sh{O}_{\widehat{T}}$-modules.  Lemma \sref{III.5.3.3} then shows
that $\sh{P}$ and $\sh{R}$ are algebraizable, and
\sref{III.5.3.2} finally proves that $\sh{F}$ is algebraizable.
\end{env}


\subsection{Application: comparison between morphisms of ordinary preschemes and morphisms of formal preschemes. Algebraizable formal preschemes}
\label{subsection:III.5.4}

\begin{theorem}[5.4.1]
\label{III.5.4.1}
Let $A$ be an adic Noetherian ring, let $\mathfrak{J}$ be an ideal
of definition of $A$, and put
\[
  S=\Spec(A),\qquad S'=\mathrm{V}(\mathfrak{J}).
\]
Let $u:X\to S$ be a proper morphism and let $v:Y\to S$ be a
separated morphism of finite type.  Let
$\widehat{S}$, $\widehat{X}$, and $\widehat{Y}$ be the completions
of $S$, $X$, and $Y$ along $S'$, $u^{-1}(S')$, and $v^{-1}(S')$,
respectively.  For every $S$-morphism $f:X\to Y$, let
$\widehat{f}:\widehat{X}\to\widehat{Y}$ be its extension to the
completions.  Then the map $f\mapsto\widehat{f}$ is a bijection
\[
  \Hom_S(X,Y)
  \xrightarrow{\ \sim\ }
  \Hom_{\widehat{S}}(\widehat{X},\widehat{Y}).
\]
\end{theorem}

\begin{proof}
We first prove injectivity.  Suppose that $f,g:X\to Y$ are
$S$-morphisms such that $\widehat{f}=\widehat{g}$.  There is then
an open neighbourhood $V$ of $X'=u^{-1}(S')$ on which $f$ and $g$
coincide \sref[I]{I.10.9.4}.  Since $u$ is closed,
\sref{III.5.1.3.1} gives $V=X$, and hence $f=g$.

We now prove surjectivity.  Let
$h:\widehat{X}\to\widehat{Y}$ be an
$\widehat{S}$-morphism.  Put
\[
  Z=X\times_S Y,
\]
and denote the canonical projections by $p:Z\to X$ and
$q:Z\to Y$.  The prescheme $Z$ is of finite type over $S$
\sref[I]{I.6.3.4}, and is therefore Noetherian.  Let
$\widehat{Z}$ be its completion along
\[
  Z'=p^{-1}\bigl(u^{-1}(S')\bigr).
\]
The formal prescheme $\widehat{Z}$ is canonically identified with
$\widehat{X}\times_{\widehat{S}}\widehat{Y}$, and its projections
onto $\widehat{X}$ and $\widehat{Y}$ are identified with
$\widehat{p}$ and $\widehat{q}$ \sref[I]{I.10.9.7}.

Since $Y$ is separated over $S$, $\widehat{Y}$ is separated over
$\widehat{S}$ \sref[I]{I.10.15.7}.  The graph morphism
\[
  \Gamma_h=(1_{\widehat{X}},h):
  \widehat{X}\longrightarrow\widehat{Z}
\]
is consequently a closed immersion \sref[I]{I.10.15.4}.  Let
$\mathfrak{T}$ be the closed formal subprescheme of
$\widehat{Z}$ associated with this immersion, and let
$j:\mathfrak{T}\to\widehat{Z}$ be the canonical injection.  Thus
\[
  \Gamma_h=j\circ w,
\]
where $w:\widehat{X}\to\mathfrak{T}$ is an isomorphism
\sref[I]{I.10.14.3} whose inverse is
$\widehat{p}\circ j$.  Moreover, $\mathfrak{T}$ is proper over
$\widehat{S}$ because $\widehat{X}$ is.  It follows from
\sref{III.5.1.8} that there is a closed subprescheme $T$ of $Z$
such that
\[
  \mathfrak{T}
  =\widehat{T}
  =T_{/(T\cap Z')},
\]
and that $j=\widehat{i}$, where $i:T\to Z$ is the canonical
injection \sref[I]{I.10.14.7}.

The morphism $p\circ i:T\to X$ is an isomorphism: its extension to
the completions is the isomorphism
$\widehat{p}\circ\widehat{i}$, and it suffices to apply
\sref{III.4.6.8}, noting as above that $S$ is the only
neighbourhood of $S'$ in $S$.
\oldpage[III]{157}
Let $g:X\to T$ be the inverse of $p\circ i$, and put
\[
  f=q\circ i\circ g.
\]
Then $f:X\to Y$ has graph $\Gamma_f=i\circ g$.  Since
$\widehat{g}$ is the inverse of
$\widehat{p\circ i}$, we have
\[
  \widehat{\Gamma_f}
  =\widehat{i}\circ\widehat{g}
  =j\circ w
  =\Gamma_h.
\]
But $\widehat{\Gamma_f}=\Gamma_{\widehat{f}}$
\sref[I]{I.10.9.8}; hence $h=\widehat{f}$, which completes the
proof.
\end{proof}

In categorical language, the functor
$X\mapsto\widehat{X}$ is therefore fully faithful
\sref[0]{0.8.1.6} from the category of preschemes proper over
$\Spec(A)$ to the category of formal preschemes proper over
$\Spf(A)$, for every adic Noetherian ring $A$.  It consequently
gives an equivalence between the first category and a full
subcategory of the second.  The objects of this subcategory will
henceforth be called \emph{algebraizable formal preschemes}.  For
such a formal prescheme $\mathfrak{X}$, there is an ordinary
prescheme $X$, proper over $\Spec(A)$ and determined up to unique
isomorphism, such that $\mathfrak{X}\simeq\widehat{X}$.

\begin{env}[5.4.2]
\label{III.5.4.2}
\emph{Scholium.}
With the notation of \sref{III.5.4.1}, put
\[
  S_n=\Spec(A/\mathfrak{J}^{n+1}),\qquad
  X_n=X\times_S S_n,\qquad
  Y_n=Y\times_S S_n.
\]
It follows from \sref{III.5.4.1} and \sref[I]{I.10.12.3} that
giving an $S$-morphism $f:X\to Y$ is equivalent to giving an adic
inductive $(S_n)$-system \sref[I]{I.10.12.2} of
$S_n$-morphisms $f_n:X_n\to Y_n$.
\end{env}

\begin{remark}[5.4.3]
\label{III.5.4.3}
Contrary to what the existence theorem \sref{III.5.1.6} might
suggest, there are formal preschemes proper over $\Spf(A)$ which
are \emph{not algebraizable}, just as there are compact analytic
spaces which do not arise from complex algebraic varieties.  We
shall encounter such preschemes later in the ``theory of moduli,''
which, when the base field is $\mathbb{C}$, deals precisely with
infinitesimal variations of the complex structure of a complete
algebraic variety; it is well known that such variations can give
rise to analytic varieties which are not algebraic.
\end{remark}

\begin{proposition}[5.4.4]
\label{III.5.4.4}
Let $A$ be an adic Noetherian ring, put
$\mathfrak{S}=\Spf(A)$, and let
\[
  g:\mathfrak{X}\to\mathfrak{S},
  \qquad
  h:\mathfrak{Y}\to\mathfrak{S}
\]
be two proper morphisms of formal preschemes.  Let
$f:\mathfrak{X}\to\mathfrak{Y}$ be an
$\mathfrak{S}$-morphism.  If $f$ is finite and
$\mathfrak{Y}$ is algebraizable, then $\mathfrak{X}$ is
algebraizable.
\end{proposition}

\begin{proof}
The hypotheses on $g$ and $h$ already imply that $f$ is proper
\sref{III.3.4.1}; for $f$ to be finite, it is enough that the fibre
$f^{-1}(y)$ be finite for every $y\in\mathfrak{Y}$
\sref[0]{0.4.8.11}.  By hypothesis,
\[
  \sh{B}=f_*(\sh{O}_{\mathfrak{X}})
\]
is a coherent $\sh{O}_{\mathfrak{Y}}$-algebra
\sref[0]{0.4.8.6}.  Write
$\mathfrak{Y}=\widehat{Y}$ and $h=\widehat{w}$, where
$w:Y\to\Spec(A)$ is a proper morphism of ordinary preschemes.  The
existence theorem gives a coherent $\sh{O}_Y$-algebra $\sh{C}$
such that $\sh{B}=\widehat{\sh{C}}$.  Put
\[
  X=\Spec(\sh{C}),
\]
and let $u:X\to Y$ be the structural morphism.  The definition of
$\mathfrak{X}$ from $\sh{B}$ \sref[0]{0.4.8.7} then shows that
$\mathfrak{X}$ is canonically isomorphic to $\widehat{X}$ and
that $f$ is identified with $\widehat{u}$; it is enough to check
this when $Y$ is affine.
\end{proof}

Proposition \sref{III.5.1.8} is a special case of
\sref{III.5.4.4}.

\begin{theorem}[5.4.5]
\label{III.5.4.5}
Let $A$ be an adic Noetherian ring, let $\mathfrak{J}$ be an ideal
of definition of $A$, put
\[
  S=\Spec(A),\qquad
  \mathfrak{S}=\widehat{S}=\Spf(A),
\]
and let $f:\mathfrak{X}\to\mathfrak{S}$ be a proper morphism of
formal preschemes.  Put
\[
  S_k=\Spec(A/\mathfrak{J}^{k+1}),\qquad
  X_k=\mathfrak{X}\times_{\mathfrak{S}}S_k,
\]
and, for every $\sh{O}_{\mathfrak{X}}$-module $\sh{F}$, put
\[
  \sh{F}_k
  =\sh{F}\otimes_{\sh{O}_{\mathfrak{X}}}\sh{O}_{X_k}
  =\sh{F}/\mathfrak{J}^{k+1}\sh{F}.
\]
\oldpage[III]{158}
Let $\sh{L}$ be an invertible
$\sh{O}_{\mathfrak{X}}$-module, and suppose that
\[
  \sh{L}_0=\sh{L}/\mathfrak{J}\sh{L}
\]
is an ample $\sh{O}_{X_0}$-module.  Then $\mathfrak{X}$ is
algebraizable.  Moreover, if $X$ is a proper $S$-prescheme such
that $\mathfrak{X}\simeq\widehat{X}$, there is an ample
$\sh{O}_X$-module $\sh{M}$ such that
$\sh{L}\simeq\widehat{\sh{M}}$; in particular, $X$ is projective
over $S$.
\end{theorem}

\begin{proof}
Apply \sref{III.5.2.3} with
$\sh{F}=\sh{O}_{\mathfrak{X}}$.  There is an integer $n_0$ such
that, for $n\geq n_0$, the canonical homomorphism
\[
  \Gamma(\mathfrak{X},\sh{L}^{\otimes n})
  \longrightarrow
  \Gamma(X_0,\sh{L}_0^{\otimes n})
\]
is surjective.  Choose $n\geq n_0$ large enough that
$\sh{L}_0^{\otimes n}$ is very ample for $S_0$
\sref[II]{II.4.5.10}.  Since
$f_0:X_0\to S_0$ is proper,
$\Gamma(X_0,\sh{L}_0^{\otimes n})$ is a finite-type $A$-module
\sref{III.3.2.1}.  There is consequently a finite-type
$A$-submodule
\[
  E\subset\Gamma(\mathfrak{X},\sh{L}^{\otimes n})
\]
whose image in $\Gamma(X_0,\sh{L}_0^{\otimes n})$ is the whole
module.

For every $k\geq0$, consider the homomorphism
\[
  u_k:
  E/\mathfrak{J}^{k+1}E
  \longrightarrow
  \Gamma(X_k,\sh{L}_k^{\otimes n})
\]
induced by the canonical injection
$E\to\Gamma(\mathfrak{X},\sh{L}^{\otimes n})$.  The module
$(f_k)_*(\sh{L}_k^{\otimes n})$ is quasi-coherent, and
\[
  \Gamma\bigl(S_k,(f_k)_*(\sh{L}_k^{\otimes n})\bigr)
  =
  \Gamma(X_k,\sh{L}_k^{\otimes n}).
\]
Thus $u_k$ defines a homomorphism
\[
  \widetilde{u}_k:
  \widetilde{E/\mathfrak{J}^{k+1}E}
  \longrightarrow
  (f_k)_*(\sh{L}_k^{\otimes n}),
\]
and hence, by adjunction, a homomorphism
\[
  \widetilde{u}_k^{\,\sharp}:
  f_k^*\bigl(\widetilde{E/\mathfrak{J}^{k+1}E}\bigr)
  \longrightarrow
  \sh{L}_k^{\otimes n}.
\]
Put
\[
  \sh{G}_k=
  f_k^*\bigl(\widetilde{E/\mathfrak{J}^{k+1}E}\bigr).
\]
We have
$\sh{G}_0=\sh{G}_k/\mathfrak{J}\sh{G}_k$
\sref[I]{I.9.1.5}, and
\[
  \widetilde{u}_0^{\,\sharp}:
  \sh{G}_0/\mathfrak{J}\sh{G}_0
  \longrightarrow
  \sh{L}_0^{\otimes n}/
  \mathfrak{J}\sh{L}_0^{\otimes n}
\]
is obtained from $\widetilde{u}_k^{\,\sharp}$ by passage to the
quotients.

By the definition of $E$, $\widetilde{u}_0^{\,\sharp}$ is the
canonical homomorphism
\[
  \sigma:
  f_0^*\bigl((f_0)_*(\sh{L}_0^{\otimes n})\bigr)
  \longrightarrow
  \sh{L}_0^{\otimes n}.
\]
Since $\sh{L}_0^{\otimes n}$ is very ample,
$\widetilde{u}_0^{\,\sharp}$ is surjective
\sref[II]{II.4.4.3}.  Nakayama's lemma then shows that every
$\widetilde{u}_k^{\,\sharp}$ is surjective.

Each $\widetilde{u}_k^{\,\sharp}$ therefore defines
\sref[II]{II.4.2.2} an $S_k$-morphism
\[
  g_k:X_k\longrightarrow
  P_k=\mathbf{P}(E/\mathfrak{J}^{k+1}E).
\]
For $h\leq k$,
\[
  P_h=P_k\times_{S_k}S_h
\]
by \sref[II]{II.4.1.3}.  The relations among the
$\widetilde{u}_k^{\,\sharp}$ and \sref[II]{II.4.2.10} show that
$(g_k)$ is an adic inductive $(S_k)$-system
\sref[I]{I.10.12.2}.  The $g_k$ consequently define an
$\mathfrak{S}$-morphism of formal preschemes
\[
  g:\mathfrak{X}\longrightarrow\mathfrak{P},
\]
where $\mathfrak{P}$ is the inductive limit of the system $(P_k)$,
or equivalently the completion $\widehat{P}$ of
$P=\mathbf{P}(E)$.  The very ampleness of
$\sh{L}_0^{\otimes n}$ implies that $g_0$ is a closed immersion
\sref[II]{II.4.4.3}.  Hence $g$ is a closed immersion of formal
preschemes \sref[0]{0.4.8.10}, and $\mathfrak{X}$ is algebraizable
by \sref{III.5.1.8}.

The existence theorem \sref{III.5.1.6} now shows that $\sh{L}$ is
the completion $\widehat{\sh{M}}$ of an invertible
$\sh{O}_X$-module $\sh{M}$.  Moreover,
$\sh{L}^{\otimes n}$ is the completion of
$\sh{M}^{\otimes n}$ \sref[I]{I.10.8.10}, and the homomorphisms
$\widetilde{u}_k^{\,\sharp}$ define a unique homomorphism
\[
  v:f^*(\widetilde{E})
  \longrightarrow
  \sh{M}^{\otimes n}
\]
\sref{III.5.1.7}.  Since the
$\widetilde{u}_k^{\,\sharp}$ are surjective,
$\widehat{v}$ is surjective \sref[I]{I.10.11.5}, and therefore so
is $v$ \sref{III.5.1.3}.  Let $r:X\to P$ be the morphism defined
by $v$ \sref[II]{II.4.2.2}.  Its extension to the completions is
$g$.  Since $g$ is a closed immersion, so is $r$, by
\sref{III.5.1.8} and \sref{III.5.4.1}.  It follows that
$\sh{M}^{\otimes n}$ is very ample \sref[II]{II.4.4.6}, and hence
that $\sh{M}$ is ample \sref[II]{II.4.5.10}.
\end{proof}

\begin{remark}[5.4.6]
\label{III.5.4.6}
Let $A$ be an adic Noetherian ring, put
$\mathfrak{S}=\Spf(A)$, and let
$f:\mathfrak{X}\to\mathfrak{S}$ be a proper morphism of formal
preschemes.  Let $\sh{N}$ be the coherent ideal of
$\sh{O}_{\mathfrak{X}}$ such that, for every affine formal open
subset $U$ of $\mathfrak{X}$,
\[
  \Gamma(U,\sh{N})
\]
is the nilradical of
$\Gamma(U,\sh{O}_{\mathfrak{X}})$.  The existence of this ideal
follows readily from \sref[I]{I.10.10.2} and from the fact that
every ring homomorphism $B\to C$ sends the nilradical of $B$ into
that of $C$.  Let $\mathfrak{X}'$ be the closed formal
subprescheme of $\mathfrak{X}$ defined by $\sh{N}$
\sref[I]{I.10.14.2}.  It would be interesting to know whether
algebraizability of $\mathfrak{X}'$ implies algebraizability of
$\mathfrak{X}$ itself.

One would doubtless arrive at a solution of this problem if one
knew how to classify, for example by invariants
\oldpage[III]{159}
of a cohomological nature, the \emph{extensions} of a structural
sheaf $\sh{O}_{\mathfrak{X}}$---for an ordinary or a formal
prescheme---by a square-zero ideal.  Equivalently, these are the
$\sh{O}_{\mathfrak{X}}$-algebras $\sh{A}$ such that
$\sh{O}_{\mathfrak{X}}\simeq\sh{A}/\sh{J}$, where $\sh{J}$ is a
square-zero ideal of $\sh{A}$.
\end{remark}


\subsection{A decomposition of certain preschemes}
\label{subsection:III.5.5}

\begin{proposition}[5.5.1]
\label{III.5.5.1}
Let $A$ be an adic Noetherian ring, let $\mathfrak{J}$ be an ideal
of definition of $A$, and put $Y=\Spec(A)$.  Let $f:X\to Y$ be a
separated morphism of finite type, and put
\[
  Y_0=\Spec(A/\mathfrak{J}),\qquad
  X_0=X\times_Y Y_0=f^{-1}(Y_0).
\]
Let $Z_0$ be an open subset of $X_0$ which is proper over $Y_0$.
Then there is an open and closed subset $Z$ of $X$, proper over
$Y$, such that
\[
  Z\cap X_0=Z_0.
\]
\end{proposition}

\begin{proof}
By hypothesis, there is an open subset $T$ of $X$ such that
$T\cap X_0=Z_0$.  Let $\widehat{T}$ be the completion along $Z_0$
of the prescheme induced by $X$ on $T$.  The support of
$\sh{O}_{\widehat{T}}$ is $Z_0$, which is proper over $Y_0$;
hence $\widehat{T}$ is proper over
$\widehat{Y}=\Spf(A)$ \sref{III.3.4.1}.

It follows from \sref{III.5.1.8} that there is a closed
subprescheme $Z$ of $T$, proper over $Y$, such that, if
$i:Z\to T$ is the canonical injection, then
\[
  \widehat{i}:\widehat{Z}\longrightarrow\widehat{T}
\]
is an isomorphism; here $\widehat{Z}$ is the completion of $Z$
along $Z_0$.  By \sref{III.4.6.8}, there is an open neighbourhood
$V$ of $Z_0$ in $T$ such that
$i^{-1}(V)\to V$ is an isomorphism.  But $i^{-1}(V)$ is a
neighbourhood of $Z_0$ in $Z$, and is therefore equal to $Z$ by
\sref{III.5.1.3.1}.  Thus $Z$ is open in $T$, and hence in $X$.
It is also closed in $X$, since $Z$ is proper over $Y$ and $X$ is
separated over $Y$.  This proves the proposition.
\end{proof}

\begin{corollary}[5.5.2]
\label{III.5.5.2}
If $X_0$ is proper over $Y_0$, then $X$ is the union of two
disjoint open subsets $Z$ and $Z'$ such that $Z$ is proper over
$Y$ and contains $X_0$.  Moreover, every closed subset $P$ of $X$
which is proper over $Y$ is contained in $Z$.
\end{corollary}

\begin{proof}
Apply \sref{III.5.5.1} with $Z_0=X_0$, and put $Z'=X-Z$.  For the
last assertion, $P\cap Z'$ is closed in $P$ and is therefore proper
over $Y$.  If it were nonempty, $f(P\cap Z')$ would be a nonempty
closed subset of $Y$, and hence would meet $Y_0$
\sref{III.5.1.3.1}.  This contradicts
$Z'\cap X_0=\varnothing$, and proves that $P\subset Z$.
\end{proof}

\bigskip
\noindent\emph{(To be continued.)}

% Printed page 504 is blank.
\clearpage
\thispagestyle{empty}
\null
\clearpage
\thispagestyle{plain}

\oldpage[III]{161}
\section*{BIBLIOGRAPHY (continued)}

\begin{flushleft}
[27] P. Gabriel, \emph{Des catégories abéliennes}, thesis, Paris (1961).

\medskip
[28] J. E. Roos, \emph{Sur les foncteurs dérivés de $\varprojlim$.
Applications}, C. R. Acad. Sci. Paris, vol. CCLII, 1961, pp. 3702--3704.

\medskip
[29] H. Cartan, \emph{Séminaire de l'École Normale Supérieure},
13th year (1960--1961), lecture no. 11.
\end{flushleft}

\clearpage
\oldpage[III]{162}
\section*{INDEX OF NOTATION}

\begin{flushleft}
$\Set$: 0, 8.1.1.

$h_X(Y)$, $h_X(u)$ ($X$, $Y$ objects of a category $\C$, $u$ a morphism
of $\C$): 0, 8.1.1.

$h_w(Y)$ ($Y$ an object and $w$ a morphism of $\C$): 0, 8.1.2.

$Z_k(E_r^{pq})$, $B_k(E_r^{pq})$ ($k\geq r$): 0, 11.1.1.

$Z_\infty(E_2^{pq})$, $B_\infty(E_2^{pq})$: 0, 11.1.1.

$K^\bullet$, $K_\bullet$ (complexes): 0, 11.2.1.

$K^{\bullet\bullet}$, $K_{\bullet\bullet}$ (bicomplexes): 0, 11.2.1.

$B^i(K^\bullet)$, $Z^i(K^\bullet)$, $H^i(K^\bullet)$,
$H^\bullet(K^\bullet)$: 0, 11.2.1.

$B_i(K_\bullet)$, $Z_i(K_\bullet)$, $H_i(K_\bullet)$,
$H_\bullet(K_\bullet)$: 0, 11.2.1.

$T(K^\bullet)$, $T(K_\bullet)$ ($T$ a functor): 0, 11.2.1.

$E(K^\bullet)$ ($K^\bullet$ a filtered complex): 0, 11.2.2.

$K^{i,\bullet}$, $K^{\bullet,j}$, $Z_{\mathrm{II}}^p(K^{i,\bullet})$,
$B_{\mathrm{II}}^p(K^{i,\bullet})$, $H_{\mathrm{II}}^p(K^{i,\bullet})$,
$Z_{\mathrm I}^p(K^{\bullet,j})$, $B_{\mathrm I}^p(K^{\bullet,j})$,
$H_{\mathrm I}^p(K^{\bullet,j})$: 0, 11.3.1.

$H_{\mathrm{II}}^p(K^{\bullet\bullet})$,
$Z_{\mathrm I}^q(H_{\mathrm{II}}^p(K^{\bullet\bullet}))$,
$B_{\mathrm I}^q(H_{\mathrm{II}}^p(K^{\bullet\bullet}))$,
$H_{\mathrm I}^q(H_{\mathrm{II}}^p(K^{\bullet\bullet}))$,
$H^n(K^{\bullet\bullet})$: 0, 11.3.1.

$F_{\mathrm I}^p(K^{\bullet\bullet})$,
$F_{\mathrm{II}}^p(K^{\bullet\bullet})$, ${}'E(K^{\bullet\bullet})$,
${}''E(K^{\bullet\bullet})$: 0, 11.3.2.

$K_{i,\bullet}$, $K_{\bullet,j}$, $H_p^{\mathrm{II}}(K_{i,\bullet})$,
$H_p^{\mathrm I}(K_{\bullet,j})$, $H_p^{\mathrm{II}}(K_{\bullet\bullet})$,
$H_p^{\mathrm I}(K_{\bullet\bullet})$,
$H_q^{\mathrm I}(H_p^{\mathrm{II}}(K_{\bullet\bullet}))$,
$H_q^{\mathrm{II}}(H_p^{\mathrm I}(K_{\bullet\bullet}))$,
$H_n(K_{\bullet\bullet})$: 0, 11.3.5.

$\RR^\bullet T(K^\bullet)$ ($T$ an additive covariant functor): 0, 11.4.3.

$\RR^\bullet T(K^\bullet,K'^{\bullet})$ ($T$ an additive covariant
bifunctor): 0, 11.4.6.

$\LL_\bullet T(K_\bullet)$ ($T$ an additive covariant functor): 0, 11.6.2.

$\LL_\bullet T(K_\bullet,K'_\bullet)$ ($T$ an additive covariant
bifunctor): 0, 11.6.6.

$\LL_\bullet T(K_{\bullet\bullet})$ ($T$ an additive covariant functor):
0, 11.7.2.

$|\sigma|$, $\Sigma(A)$, $\mathrm C_\bullet(A)$, $\mathrm L_\bullet(A)$
($A$ a finite set, $\sigma$ a simplex of $A$): 0, 11.8.1.

$\mathrm C^\bullet(A,B;\sh{S})$, $\mathrm L^\bullet(A,B;\sh{S})$
($A$, $B$ finite sets, $\sh{S}$ a coefficient system): 0, 11.8.4.

$\mathrm P^\bullet(A,B;\sh{S})$ ($A$, $B$ finite sets, $\sh{S}$ a
coefficient system): 0, 11.8.6.

$\mathrm C^\bullet(A,B;\sh{S}^\bullet)$,
$\mathrm L^\bullet(A,B;\sh{S}^\bullet)$ ($A$, $B$ finite sets,
$\sh{S}^\bullet$ a complex of coefficient systems): 0, 11.8.8.

$\mathrm P^\bullet(A,B;\sh{S}^\bullet)$ ($A$, $B$ finite sets,
$\sh{S}^\bullet$ a complex of coefficient systems): 0, 11.8.10.

$\chi(M^\bullet)$ ($M^\bullet$ a finite complex of finite-length modules):
0, 11.10.1.

$\check C^\bullet(\mathfrak U,\sh{F})$ ($\sh{F}$ a sheaf of abelian groups
on $X$, $\mathfrak U$ an open cover of $X$): 0, 12.1.2.

$\RR^p f_*(\sh{F})$ ($f:X\to Y$ a morphism of ringed spaces,
$\sh{F}$ a sheaf of $\sh{O}_X$-Modules): 0, 12.2.1.

$\sh{H}^p(f,\sh{K}^\bullet)$,
$\sh{H}^p_f(\sh{K}^\bullet)$, ${}'E(f,\sh{K}^\bullet)$,
${}''E(f,\sh{K}^\bullet)$ ($f:X\to Y$ a morphism of ringed spaces,
$\sh{K}^\bullet$ a complex of $\sh{O}_X$-Modules): 0, 12.4.1.

$\mathbf H^\bullet(X,\sh{K}^\bullet)$,
$\mathbf H^\bullet(V,\sh{K}^\bullet)$ ($X$ a ringed space,
$\sh{K}^\bullet$ a complex of $\sh{O}_X$-Modules, $V$ open in $X$):
0, 12.4.2.

$\operatorname{gr}^\bullet(A)$ ($A$ a projective system satisfying (ML)):
0, 13.4.2.

$E(A)$ ($A$ a filtered object): 0, 13.6.4.

$K_\bullet(\mathbf f)$, $i_{\mathbf f}$ ($\mathbf f$ a finite family of
elements of a ring): III, 1.1.1.

$K^\bullet(\mathbf f,M)$, $K_\bullet(\mathbf f,M)$ ($\mathbf f$ a finite
family of elements of a ring $A$, $M$ an $A$-module): III, 1.1.2.

$H^\bullet(\mathbf f,M)$, $H_\bullet(\mathbf f,M)$ ($\mathbf f$ a finite
family of elements of a ring $A$, $M$ an $A$-module): III, 1.1.3.

$C^\bullet((\mathbf f),M)$, $H^\bullet((\mathbf f),M)$ ($\mathbf f$ a
finite family of elements of a ring $A$, $M$ an $A$-module): III, 1.1.6.

$H^\bullet(U,\sh{F}^{(*)})$,
$H^\bullet(\mathfrak U,\sh{F}^{(*)})$ ($\sh{F}$ a quasi-coherent
$\sh{O}_X$-Module, $U$ open in $X$, $\mathfrak U$ an open cover of $X$):
III, 2.1.1.
\end{flushleft}

\clearpage
\oldpage[III]{163}
\section*{TERMINOLOGICAL INDEX}

\begin{flushleft}
Abutment of a spectral sequence: 0, 11.1.1.

Algebraizable ($\sh{O}_X$-Module): III, 5.2.1.

Algebraizable (formal scheme): III, 5.4.2.

Analytically integral (ring): III, 4.3.6.

Quasi-compact map: 0, 9.1.1.

Augmentation of a resolution: 0, 11.4.1.

Euler--Poincaré characteristic of a complex: 0, 11.10.1.

Euler--Poincaré characteristic of a coherent $\sh{O}_X$-Module
($X$ projective over $\Spec(A)$, $A$ an Artinian ring): III, 2.5.1.

Bi-alternating cochain: 0, 11.8.4.

Complex defined by a bicomplex: 0, 11.3.1.

Exterior algebra complex: III, 1.1.1.

Condition (ML): 0, 13.1.2.

Constructible (subset, set): 0, 9.1.2.

Constructible (function): 0, 9.3.1.

Cup product: 0, 12.1.2.

Di-homomorphism for algebraic structures on categories: 0, 8.2.1.

Exact (subset) in an abelian category: III, 3.1.1.

Essentially constant (projective system): 0, 13.4.2.

Stein factorization: III, 4.3.3.

Co-discrete, co-separated, discrete, exhaustive, finite, separated
(filtration): 0, 11.1.3.

Final (object) of a category: 0, 8.1.10.

Finite (morphism) of formal preschemes: III, 4.8.2.

Finite (formal prescheme) over $\mathfrak Y$,
$\mathfrak Y$-finite (formal prescheme): III, 4.8.2.

Canonical covariant functor $\C\to\CHom(\C\op,\Set)$: 0, 8.1.2.

Arithmetic genus: III, 2.5.1.

Geometrically connected (fibre): III, 4.3.4.

Homomorphism of algebraic structures on categories: 0, 8.2.1.

Hypercohomology of a functor with respect to a complex $K^\bullet$:
0, 11.4.3.

Hypercohomology of a bifunctor with respect to two complexes
$K^\bullet$, $K'^{\bullet}$: 0, 11.4.6.

Hyperhomology of a functor with respect to a complex $K_\bullet$:
0, 11.6.2.

Hyperhomology of a bifunctor with respect to two complexes $K_\bullet$,
$K'_\bullet$: 0, 11.6.6.

Hyperhomology of a functor with respect to a bicomplex: 0, 11.7.2.

Locally constructible (set): 0, 9.1.11.

External (internal) composition law on the objects of a category: 0, 8.2.1.

\oldpage[III]{164}

Filtered $S$-$\C$-module, graded $\operatorname{gr}^\bullet(S)$-$\C$-module,
bigraded $\operatorname{gr}^{\bullet\bullet}(S)$-$\C$-module: 0, 13.6.6.

Morphism of spectral sequences: 0, 11.1.2.

Geometric number of connected components of a fibre: III, 4.3.4.

Group object in $\C$, $\C$-group, $\C$-ring, $\C$-module: 0, 8.2.3.

Object of boundaries, object of coboundaries, object of cocycles, object of
cycles: 0, 11.2.1.

Object of universal images: 0, 13.1.1.

Graded object associated with a projective system satisfying (ML): 0, 13.4.2.

Graded object bounded below (above): 0, 11.2.1.

Proper subset (over $\mathfrak Y$) of a formal prescheme: III, 3.4.1.

Full (subcategory): 0, 8.1.5.

Fully faithful (functor): 0, 8.1.5.

Hilbert polynomial: III, 2.5.3.

Proper (morphism) of formal preschemes: III, 3.4.1.

Representable (functor): 0, 8.1.8.

Cohomological, right, left, homological, injective, free, flat, projective
(resolution): 0, 11.4.1.

Right, injective Cartan--Eilenberg resolution: 0, 11.4.2.

Left, projective Cartan--Eilenberg resolution: 0, 11.6.1.

Retrocompact (set): 0, 9.1.1.

Strict (projective system): 0, 13.4.2.

Spectral sequence in an abelian category: 0, 11.1.1.

Weakly convergent, regular, coregular, biregular (spectral sequence):
0, 11.1.3.

Degenerate spectral sequence: 0, 11.1.6.

Spectral sequence of a functor relative to a filtered object: 0, 13.6.4.

Spectral sequences of a bicomplex: 0, 11.3.2 and 11.3.5.

Hypercohomology spectral sequences of a functor with respect to a complex:
0, 11.4.3.

Hyperhomology spectral sequences of a functor with respect to a complex:
0, 11.6.2.

Hyperhomology spectral sequences of a functor with respect to a bicomplex:
0, 11.7.2.

Coefficient system: 0, 11.8.4.

Unibranch (ring, point): III, 4.3.6.

Universally open (morphism): III, 4.3.9.
\end{flushleft}

\clearpage
\oldpage[III]{165}
\section*{TABLE OF CONTENTS}

\begin{flushleft}
\textsc{Chapter 0.} --- \textbf{Preliminaries (continued)} \dotfill 5

\medskip
\S~8. Representable functors \dotfill 5

8.1. Representable functors \dotfill 5

8.2. Algebraic structures in categories \dotfill 9

\medskip
\S~9. Constructible sets \dotfill 12

9.1. Constructible sets \dotfill 12

9.2. Constructible subsets of Noetherian spaces \dotfill 14

9.3. Constructible functions \dotfill 16

\medskip
\S~10. Supplement on flat modules \dotfill 17

10.1. Relations between flat modules and free modules \dotfill 17

10.2. Local flatness criteria \dotfill 18

10.3. Existence of flat extensions of local rings \dotfill 20

\medskip
\S~11. Supplement on homological algebra \dotfill 23

11.1. Review of spectral sequences \dotfill 23

11.2. The spectral sequence of a filtered complex \dotfill 27

11.3. The spectral sequences of a bicomplex \dotfill 29

11.4. Hypercohomology of a functor with respect to a complex $K^\bullet$
\dotfill 32

11.5. Passage to the inductive limit in hypercohomology \dotfill 35

11.6. Hyperhomology of a functor with respect to a complex $K_\bullet$
\dotfill 39

11.7. Hyperhomology of a functor with respect to a bicomplex
$K_{\bullet\bullet}$ \dotfill 41

11.8. Supplement on the cohomology of simplicial complexes \dotfill 43

11.9. A lemma on complexes of finite type \dotfill 46

11.10. Euler--Poincaré characteristic of a complex of finite-length modules
\dotfill 48

\medskip
\S~12. Supplement on sheaf cohomology \dotfill 49

12.1. Cohomology of sheaves of modules on ringed spaces \dotfill 49

12.2. Higher direct images \dotfill 57

12.3. Supplements on the Ext functors of sheaves \dotfill 60

12.4. Hypercohomology of the direct image functor \dotfill 62

\oldpage[III]{166}

\S~13. Projective limits in homological algebra \dotfill 64

13.1. The Mittag--Leffler condition \dotfill 64

13.2. The Mittag--Leffler condition for abelian groups \dotfill 65

13.3. Application: cohomology of a projective limit of sheaves \dotfill 68

13.4. Mittag--Leffler condition and graded objects associated with projective
systems \dotfill 69

13.5. Projective limits of spectral sequences of filtered complexes
\dotfill 71

13.6. Spectral sequence of a functor relative to an object equipped with a
finite filtration \dotfill 73

13.7. Derived functors of a projective limit of arguments \dotfill 75

\medskip
\textsc{Chapter III.} --- \textbf{Cohomological study of coherent sheaves}
\dotfill 81

\medskip
\S~1. Cohomology of affine schemes \dotfill 82

1.1. Review of the exterior algebra complex \dotfill 82

1.2. \v Cech cohomology of an open cover \dotfill 85

1.3. Cohomology of an affine scheme \dotfill 88

1.4. Application to the cohomology of arbitrary preschemes \dotfill 89

\medskip
\S~2. Cohomological study of projective morphisms \dotfill 95

2.1. Explicit calculations of certain cohomology groups \dotfill 95

2.2. The fundamental theorem of projective morphisms \dotfill 100

2.3. Application to graded sheaves of algebras and modules \dotfill 102

2.4. A generalisation of the fundamental theorem \dotfill 107

2.5. Euler--Poincaré characteristic and the Hilbert polynomial \dotfill 109

2.6. Application: ampleness criteria \dotfill 111

\medskip
\S~3. Finiteness theorem for proper morphisms \dotfill 115

3.1. The dévissage lemma \dotfill 115

3.2. The finiteness theorem: the case of usual schemes \dotfill 116

3.3. Generalization of the finiteness theorem (usual schemes) \dotfill 118

3.4. Finiteness theorem: the case of formal schemes \dotfill 119

\medskip
\S~4. The fundamental theorem of proper morphisms. Applications \dotfill 122

4.1. The fundamental theorem \dotfill 122

4.2. Particular cases and variants \dotfill 129

4.3. Zariski's connectedness theorem \dotfill 130

4.4. Zariski's ``main theorem'' \dotfill 135

4.5. Completions of modules of homomorphisms \dotfill 138

4.6. Relations between formal morphisms and ordinary morphisms \dotfill 139

4.7. An ampleness criterion \dotfill 145

4.8. Finite morphisms of formal preschemes \dotfill 146

\oldpage[III]{167}

\S~5. An existence theorem for coherent algebraic sheaves \dotfill 149

5.1. Statement of the theorem \dotfill 149

5.2. Proof of the existence theorem: the projective and quasi-projective case
\dotfill 151

5.3. Proof of the existence theorem: the general case \dotfill 154

5.4. Application: comparison between morphisms of ordinary preschemes and
morphisms of formal preschemes. Algebraizable formal preschemes \dotfill 156

5.5. A decomposition of certain preschemes \dotfill 159

\medskip
\textsc{Bibliography (continued)} \dotfill 161

\textsc{Index of notation} \dotfill 162

\textsc{Terminological index} \dotfill 163

\begin{flushright}
\emph{Received June 28, 1961.}
\end{flushright}
\end{flushleft}


\clearpage
\phantomsection
\addcontentsline{toc}{section}{Consolidated references for this edition}
\section*{Consolidated references for this edition}
\bibliography{the}
\bibliographystyle{amsalpha}

\end{document}
